\documentclass{article}
\usepackage{amsmath, amssymb, amsthm}
\usepackage{hyperref}
\usepackage{geometry}
\usepackage{titlesec}
\geometry{margin=1in}

\title{Erd\H{o}s Open Problems\\[0.5em]\large A Collection of All Currently Unsolved Erd\H{o}s Problems}
\date{Generated: 2026-06-09}
\author{Source: \url{https://www.erdosproblems.com}}

\begin{document}
\maketitle
\tableofcontents
\newpage


%% ============================================================
\section{Problem \#1}
\label{prob:1}

\noindent\textbf{Status:} OPEN \quad|\quad \textbf{Prize:} \$500 \quad|\quad \textbf{Updated:} 2025-08-31

\noindent\textbf{Tags:} number theory, additive combinatorics

\medskip

\noindent\textbf{Statement:}

If $A\subseteq \{1,\ldots,N\}$ with $\lvert A\rvert=n$ is such that the subset sums $\sum_{a\in S}a$ are distinct for all $S\subseteq A$ then\[N \gg 2^{n}.\]




Erd\H{o}s called this 'perhaps my first serious problem' (in \cite{Er98} he dates it to 1931). The powers of $2$ show that $2^n$ would be best possible here: this provides an upper bound for the minimal such $N$ of $N\leq 2^{n-1}$. This was improved by Conway and Guy \cite{CoGu68} (see also \cite{Gu82}) to $N\leq 2^{n-2}$ (for all large $n$). The best known upper bound is\[N\leq 0.22002\cdot 2^n,\]due to Bohman. The trivial lower bound is $N \gg 2^{n}/n$, since all $2^n$ distinct subset sums must lie in $[0,Nn)$. Erd\H{o}s and Moser \cite{Er56} proved\[ N\geq (\tfrac{1}{4}-o(1))\frac{2^n}{\sqrt{n}}.\](In \cite{Er85c} Erd\H{o}s offered \$100 for any improvement of the constant $1/4$ here.) A number of improvements of the constant have been given (see \cite{St23} for a history), with the current record $\sqrt{2/\pi}$ first proved in unpublished work of Elkies and Gleason. Two proofs achieving this constant are provided by Dubroff, Fox, and Xu \cite{DFX21}, who in fact prove the exact bound $N\geq \binom{n}{\lfloor n/2\rfloor}$. An equivalent formulation is to ask for the maximal size of a set of integers in $[1,x]$ which is dissociated (all subset sums are distinct). If $F(x)$ is the size of such a set then this problem is equivalent to\[F(x) &lt;\log_2x+O(1).\]Conway and Guy (see \cite{Gu82}) conjectured that $F(2^k)=k+2$ for all large $k$, but Erd\H{o}s \cite{Er80} wrote he had 'no opinion'. In \cite{Er73} and \cite{ErGr80} the generalisation where $A\subseteq (0,N]$ is a set of real numbers such that the subset sums all differ by at least $1$ is proposed, with the same conjectured bound. (The second proof of \cite{DFX21} applies also to this generalisation.) This generalisation seems to have first appeared in \cite{Gr71}. This problem appears in Erd\H{o}s' book with Spencer \cite{ErSp74} in the final chapter titled 'The kitchen sink'. As Ruzsa writes in \cite{Ru99} "it is a rich kitchen where such things go to the sink". The sequence of minimal $N$ for a given $n$ is A276661 in the OEIS. See also [350] . This is discussed in problem C8 of Guy's collection \cite{Gu04}.




References

[CoGu68] J. H. Conway and R. K. Guy, Sets of natural numbers with distinct sums . Notices Amer. Math. Soc. (1968), 345.

[DFX21] Dubroff, Q. and Fox, J. and Xu, M. W., A note on the Erd\H{o}s distinct subset sums problem . SIAM Journal on Discrete Mathematics (2021), 322-324.

[Er56] Erd\H{o}s, P., Problems and results in additive number theory . Colloque sur la Th\'{e}orie des Nombres, Bruxelles, 1955 (1956), 127-137.

[Er73] Erd\H{o}s, P., Problems and results on combinatorial number theory . A survey of combinatorial theory (Proc. Internat. Sympos., Colorado State Univ., Fort Collins, Colo., 1971) (1973), 117-138.

[Er80] Erd\H{o}s, Paul, A survey of problems in combinatorial number theory . Ann. Discrete Math. (1980), 89-115.

[Er85c] Erd\H{o}s, P., On some of my problems in number theory I would most like to see solved . Number theory (Ootacamund, 1984) (1985), 74-84.

[Er98] Erd\H{o}s, Paul, Some of my new and almost new problems and results in combinatorial number theory . Number theory (Eger, 1996) (1998), 169-180.

[ErGr80] Erd\H{o}s, P. and Graham, R., Old and new problems and results in combinatorial number theory . Monographies de L'Enseignement Mathematique (1980).

[ErSp74] Erd\H{o}s, Paul and Spencer, Joel, Probabilistic methods in combinatorics . Akad\'{e}miai Kiad\'{o} (1974).

[Gr71] Graham, R. L., On sums of integers taken from a fixed sequence . (1971), 22--40.

[Gu04] Guy, Richard K., Unsolved problems in number theory . (2004), xviii+437.

[Gu82] Guy, Richard K., Sets of integers whose subsets have distinct sums . (1982), 141--154.

[Ru99] Ruzsa, I., Erd\H{o}s and the Integers . Journal of Number Theory (1999), 115-163.

[St23] Steinerberger, S., Some remarks on the Erd\H{o}s distinct subset sums problem . arXiv:2208.12182 (2023).

\noindent\rule{\linewidth}{0.4pt}


%% ============================================================
\section{Problem \#3}
\label{prob:3}

\noindent\textbf{Status:} OPEN \quad|\quad \textbf{Prize:} \$5000 \quad|\quad \textbf{Updated:} 2025-08-31

\noindent\textbf{Tags:} number theory, additive combinatorics, arithmetic progressions

\medskip

\noindent\textbf{Statement:}

If $A\subseteq \mathbb{N}$ has $\sum_{n\in A}\frac{1}{n}=\infty$ then must $A$ contain arbitrarily long arithmetic progressions?




This is essentially asking for good bounds on $r_k(N)$, the size of the largest subset of $\{1,\ldots,N\}$ without a non-trivial $k$-term arithmetic progression. For example, a bound like\[r_k(N) \ll_k \frac{N}{(\log N)(\log\log N)^2}\]would be sufficient. Even the case $k=3$ is non-trivial, but was proved by Bloom and Sisask \cite{BlSi20}. Much better bounds for $r_3(N)$ were subsequently proved by Kelley and Meka \cite{KeMe23}. Green and Tao \cite{GrTa17} proved $r_4(N)\ll N/(\log N)^{c}$ for some small constant $c&gt;0$. Gowers \cite{Go01} proved\[r_k(N) \ll \frac{N}{(\log\log N)^{c_k}},\]where $c_k&gt;0$ is a small constant depending on $k$. The current best bounds for general $k$ are due to Leng, Sah, and Sawhney \cite{LSS24}, who show that\[r_k(N) \ll \frac{N}{\exp((\log\log N)^{c_k})}\]for some constant $c_k&gt;0$ depending on $k$. Curiously, Erd\H{o}s \cite{Er83c} thought this conjecture was the 'only way to approach' the conjecture that there are arbitrarily long arithmetic progressions of prime numbers, now a theorem due to Green and Tao \cite{GrTa08} (see [219] ). In \cite{Er81} Erd\H{o}s makes the stronger conjecture that\[r_k(N) \ll_C\frac{N}{(\log N)^C}\]for every $C&gt;0$ (now known for $k=3$ due to Kelley and Meka \cite{KeMe23}) - see [140] . See also [139] and [142] . This is discussed in problem A5 of Guy's collection \cite{Gu04}.




References

[BlSi20] Bloom, T.F. and Sisask, O., Breaking the logarithmic barrier in Roth's theorem on arithmetic progressions . arXiv:2007.03528 (2020).

[Er81] Erd\H{o}s, P., On the combinatorial problems which I would most like to see solved . Combinatorica (1981), 25-42.

[Er83c] Erd\H{o}s, Paul, Combinatorial problems in geometry . Math. Chronicle (1983), 35-54.

[Go01] Gowers, W. T., A new proof of Szemer\'{e}di's theorem . Geom. Funct. Anal. (2001), 465-588.

[GrTa08] Green, Ben and Tao, Terence, The primes contain arbitrarily long arithmetic progressions . Ann. of Math. (2) (2008), 481-547.

[GrTa17] Green, Ben and Tao, Terence, New bounds for Szemer\'{e}di's theorem, III: a polylogarithmic bound for $r_4(N)$ . Mathematika (2017), 944-1040.

[Gu04] Guy, Richard K., Unsolved problems in number theory . (2004), xviii+437.

[KeMe23] Kelley, Z. and Meka, R., Strong Bounds for 3-Progressions . arXiv:2302.05537 (2023).

[LSS24] Leng, J., Sah, A. and Sawhney, M., Improved bounds for Szemer\'{e}di's theorem . arXiv:2402.17995 (2024).

\noindent\rule{\linewidth}{0.4pt}


%% ============================================================
\section{Problem \#5}
\label{prob:5}

\noindent\textbf{Status:} OPEN \quad|\quad \textbf{Prize:} none \quad|\quad \textbf{Updated:} 2025-08-31

\noindent\textbf{Tags:} number theory, primes

\medskip

\noindent\textbf{Statement:}

Let $C\geq 0$. Is there an infinite sequence of $n_i$ such that\[\lim_{i\to \infty}\frac{p_{n_i+1}-p_{n_i}}{\log n_i}=C?\]




Let $S$ be the set of limit points of $(p_{n+1}-p_n)/\log n$. This problem asks whether $S=[0,\infty]$. Although this conjecture remains unproven, a lot is known about $S$. Some highlights: {UL} {LI}$\infty\in S$ by Westzynthius' result \cite{We31} on large prime gaps,{/LI} {LI}$0\in S$ by the work of Goldston, Pintz, and Yildirim \cite{GPY09} on small prime gaps,{/LI} {LI}Erd\H{o}s \cite{Er55} and Ricci \cite{Ri56} independently showed that $S$ has positive Lebesgue measure,{/LI} {LI} Hildebrand and Maier \cite{HiMa88} showed that $S$ contains arbitrarily large (finite) numbers,{/LI} {LI} Pintz \cite{Pi16} showed that there exists some small constant $c&gt;0$ such that $[0,c]\subset S$,{/LI} {LI} Banks, Freiberg, and Maynard \cite{BFM16} showed that at least $12.5\%$ of $[0,\infty)$ belongs to $S$,{/LI} {LI} Merikoski \cite{Me20} showed that at least $1/3$ of $[0,\infty)$ belongs to $S$, and that $S$ has bounded gaps.{/LI} {/UL} In \cite{Er65b}, \cite{Er85c}, and \cite{Er97c} Erd\H{o}s asks whether $S$ is everywhere dense (but Weisenberg notes that clearly $S$ is closed so this is equivalent to asking whether $S=[0,\infty]$). See also [234] .




References

[BFM16] Banks, William D. and Freiberg, Tristan and Maynard, James, On limit points of the sequence of normalized prime gaps . Proc. Lond. Math. Soc. (3) (2016), 515-539.

[Er55] Erd\H{o}s, Paul, Some remarks on number theory . Riveon Lematematika (1955), 45-48.

[Er65b] Erd\H{o}s, Paul, Some recent advances and current problems in number theory . Lectures on Modern Mathematics, Vol. III (1965), 196-244.

[Er85c] Erd\H{o}s, P., On some of my problems in number theory I would most like to see solved . Number theory (Ootacamund, 1984) (1985), 74-84.

[Er97c] Erd\H{o}s, Paul, Some of my favorite problems and results . The mathematics of Paul Erd\H{o}s, I (1997), 47-67.

[GPY09] Goldston, Daniel A. and Pintz, J\'{a}nos and Y\i ld\i r\i m, Cem Y., Primes in tuples. I . Ann. of Math. (2) (2009), 819-862.

[HiMa88] Hildebrand, Adolf and Maier, Helmut, Gaps between prime numbers . Proc. Amer. Math. Soc. (1988), 1-9.

[Me20] Merikoski, Jori, Limit points of normalized prime gaps . J. Lond. Math. Soc. (2) (2020), 99-124.

[Pi16] Pintz, J\'{a}nos, Polignac numbers, conjectures of Erd\H{o}s on gaps between primes, arithmetic progressions in primes, and the bounded gap conjecture . From arithmetic to zeta-functions (2016), 367-384.

[Ri56] Ricci, Giovanni, Recherches sur l'allure de la suite $\{p_{n+1}-p_n/\log p_n\}$ . Colloque sur la Th\'{e}orie des Nombres, Bruxelles, 1955 (1956), 93-106.

[We31] Westzynthius, E., \"{U}ber die Verteilung der Zahlen, die zu den n ersten Primzahlen teilerfremd sind . Commentat. Phys. Math. (1931), 1-37.

\noindent\rule{\linewidth}{0.4pt}


%% ============================================================
\section{Problem \#9}
\label{prob:9}

\noindent\textbf{Status:} OPEN \quad|\quad \textbf{Prize:} none \quad|\quad \textbf{Updated:} 2025-08-31

\noindent\textbf{Tags:} number theory, additive basis, primes

\medskip

\noindent\textbf{Statement:}

Let $A$ be the set of all odd integers $\geq 1$ not of the form $p+2^{k}+2^l$ (where $k,l\geq 0$ and $p$ is prime). Is the upper density of $A$ positive?




Crocker \cite{Cr71} proved that there are infinitely many odd integers not of this form; his proof in fact proves there are $\gg\log\log N$ such integers in $\{1,\ldots,N\}$. Pan \cite{Pa11} improved this to $\gg_\epsilon N^{1-\epsilon}$ for any $\epsilon&gt;0$. Erd\H{o}s believed this cannot be proved by covering systems, i.e. integers of the form $p+2^k+2^l$ exist in every infinite arithmetic progression. The sequence of such numbers is A006286 in the OEIS. In \cite{Er80} Erd\H{o}s conjectured 'with some trepidation' that for any finite set of primes $P$ all large integers $n$ can be written as $n=m+2^k+2^l$ where $m$ is a multiple of one of the primes in $P$. See also [10] , [11] , and [16] . This is discussed in problem A19 of Guy's collection \cite{Gu04}.




References

[Cr71] Crocker, Roger, On the sum of a prime and of two powers of two . Pacific J. Math. (1971), 103-107.

[Er80] Erd\H{o}s, Paul, A survey of problems in combinatorial number theory . Ann. Discrete Math. (1980), 89-115.

[Gu04] Guy, Richard K., Unsolved problems in number theory . (2004), xviii+437.

[Pa11] Pan, Hao, On the integers not of the form {$p+2^a+2^b$} . Acta Arith. (2011), 55-61.

\noindent\rule{\linewidth}{0.4pt}


%% ============================================================
\section{Problem \#10}
\label{prob:10}

\noindent\textbf{Status:} OPEN \quad|\quad \textbf{Prize:} none \quad|\quad \textbf{Updated:} 2025-08-31

\noindent\textbf{Tags:} number theory, additive basis, primes

\medskip

\noindent\textbf{Statement:}

Is there some $k$ such that every large integer is the sum of a prime and at most $k$ powers of 2?




Erd\H{o}s described this as 'probably unattackable'. In \cite{ErGr80} Erd\H{o}s and Graham suggest that no such $k$ exists, although in \cite{Er80} Erd\H{o}s conjectured with 'trepidation' that such a $k$ does exist. Gallagher \cite{Ga75} has shown that for any $\epsilon&gt;0$ there exists $k(\epsilon)$ such that the set of integers which are the sum of a prime and at most $k(\epsilon)$ many powers of 2 has lower density at least $1-\epsilon$. Granville and Soundararajan \cite{GrSo98} have conjectured that at most $3$ powers of 2 suffice for all odd integers, and hence at most $4$ powers of $2$ suffice for all even integers. (The restriction to odd integers is important here - for example, Bogdan Grechuk has observed that $1117175146$ is not the sum of a prime and at most $3$ powers of $2$, and pointed out that parity considerations, coupled with the fact that there are many integers not the sum of a prime and $2$ powers of $2$ (see [9] ) suggest that there exist infinitely many even integers which are not the sum of a prime and at most $3$ powers of $2$). See also [9] , [11] , and [16] .




References

[Er80] Erd\H{o}s, Paul, A survey of problems in combinatorial number theory . Ann. Discrete Math. (1980), 89-115.

[ErGr80] Erd\H{o}s, P. and Graham, R., Old and new problems and results in combinatorial number theory . Monographies de L'Enseignement Mathematique (1980).

[Ga75] Gallagher, P. X., Primes and powers of 2 . Invent. Math. (1975), 125-142.

[GrSo98] Granville, A. and Soundararajan, K., A Binary Additive Problem of Erd\H{o}s and the Order of $2$ mod $p^2$ . The Ramanujan Journal (1998), 283-298.

\noindent\rule{\linewidth}{0.4pt}


%% ============================================================
\section{Problem \#11}
\label{prob:11}

\noindent\textbf{Status:} OPEN \quad|\quad \textbf{Prize:} none \quad|\quad \textbf{Updated:} 2026-03-14

\noindent\textbf{Tags:} number theory, additive basis

\medskip

\noindent\textbf{Statement:}

Is every large odd integer $n$ the sum of a squarefree number and a power of 2?




Odlyzko has checked this up to $10^7$. Hercher \cite{He24b} has verified this is true for all odd integers up to $2^{50}\approx 1.12\times 10^{15}$. Granville and Soundararajan \cite{GrSo98} have proved that this is very related to the problem of finding Wieferich primes , which are $p$ for which $2^{p-1}\equiv 1\pmod{p^2}$ - for example, if every odd integer is the sum of a squarefree number and a power of $2$ then a positive proportion of primes are non-Wieferich primes. Erd\H{o}s often asked this under the weaker assumption that $n$ is not divisible by $4$. Erd\H{o}s thought that proving this with two powers of 2 is perhaps easy, and could prove that it is true (with a single power of two) for almost all $n$. See also [9] , [10] , and [16] . This is mentioned in problem A19 of Guy's collection \cite{Gu04}.




References

[GrSo98] Granville, A. and Soundararajan, K., A Binary Additive Problem of Erd\H{o}s and the Order of $2$ mod $p^2$ . The Ramanujan Journal (1998), 283-298.

[Gu04] Guy, Richard K., Unsolved problems in number theory . (2004), xviii+437.

[He24b] C. Hercher, On the Sum of Squarefree Integers and a Power of Two . arXiv:2411.01964 (2024).

\noindent\rule{\linewidth}{0.4pt}


%% ============================================================
\section{Problem \#12}
\label{prob:12}

\noindent\textbf{Status:} OPEN \quad|\quad \textbf{Prize:} none \quad|\quad \textbf{Updated:} 2025-08-31

\noindent\textbf{Tags:} number theory

\medskip

\noindent\textbf{Statement:}

Let $A$ be an infinite set such that there are no distinct $a,b,c\in A$ such that $a\mid (b+c)$ and $b,c&#62;a$. Is there such an $A$ with\[\liminf \frac{\lvert A\cap\{1,\ldots,N\}\rvert}{N^{1/2}}&#62;0?\]Does there exist some absolute constant $c&#62;0$ such that there are always infinitely many $N$ with\[\lvert A\cap\{1,\ldots,N\}\rvert&#60;N^{1-c}?\]Is it true that\[\sum_{n\in A}\frac{1}{n}&#60;\infty?\]




Asked by Erd\H{o}s and S\'{a}rk\"{o}zy \cite{ErSa70}, who proved that $A$ must have density $0$. They also prove that this is essentially best possible, in that given any function $f(x)\to \infty$ as $x\to \infty$ there exists a set $A$ with this property and infinitely many $N$ such that\[\lvert A\cap\{1,\ldots,N\}\rvert&gt;\frac{N}{f(N)}.\](Their example is given by all integers in $(y_i,\frac{3}{2}y_i)$ congruent to $1$ modulo $(2y_{i-1})!$, where $y_i$ is some sufficiently quickly growing sequence.) An example of an $A$ with this property where\[\liminf \frac{\lvert A\cap\{1,\ldots,N\}\rvert}{N^{1/2}}\log N&gt;0\]is given by the set of $p^2$, where $p\equiv 3\pmod{4}$ is prime. Elsholtz and Planitzer \cite{ElPl17} have constructed such an $A$ with\[\lvert A\cap\{1,\ldots,N\}\rvert\gg \frac{N^{1/2}}{(\log N)^{1/2}(\log\log N)^2(\log\log\log N)^2}.\]Schoen \cite{Sc01} proved that if all elements in $A$ are pairwise coprime then\[\lvert A\cap\{1,\ldots,N\}\rvert \ll N^{2/3}\]for infinitely many $N$. Baier \cite{Ba04} has improved this to $\ll N^{2/3}/\log N$. DeepMind has provided a construction that shows the answer to the second question is no (and hence the answer to the first question is yes). This construction has been simplified and improved in the comments; it is now known that there exists an $A$ with this property such that, for all large $N$,\[\lvert A\cap \{1,\ldots,N\}\rvert \geq \frac{N}{(\log N)^{O(\log\log\log N)}}.\]It is unknown whether there exists such an $A$ with $\sum_{n\in A}\frac{1}{n}=\infty$. For the finite version see [13] .




References

[Ba04] Baier, Stephan, A note on {$\scr P$}-sets . Integers (2004), A13, 6.

[ElPl17] Elsholtz, Christian and Planitzer, Stefan, On Erd\H{o}s and S\'{a}rk\"ozy's sequences with Property P . Monatsh. Math. (2017), 565--575.

[ErSa70] Erd\H{o}s, P. and S\'{a}rk\"ozi, A., On the divisibility properties of sequences of integers . Proc. London Math. Soc. (3) (1970), 97-101.

[Sc01] Schoen, Tomasz, On a problem of Erd\H{o}s and S\'{a}rk\"ozy . J. Combin. Theory Ser. A (2001), 191--195.

\noindent\rule{\linewidth}{0.4pt}


%% ============================================================
\section{Problem \#14}
\label{prob:14}

\noindent\textbf{Status:} OPEN \quad|\quad \textbf{Prize:} none \quad|\quad \textbf{Updated:} 2025-08-31

\noindent\textbf{Tags:} number theory, sidon sets, additive combinatorics

\medskip

\noindent\textbf{Statement:}

Let $A\subseteq \mathbb{N}$. Let $B\subseteq \mathbb{N}$ be the set of integers which are representable in exactly one way as the sum of two elements from $A$. Is it true that for all $\epsilon&#62;0$ and large $N$\[\lvert \{1,\ldots,N\}\backslash B\rvert \gg_\epsilon N^{1/2-\epsilon}?\]Is it possible that\[\lvert \{1,\ldots,N\}\backslash B\rvert =o(N^{1/2})?\]




Apparently originally considered by Erd\H{o}s and Nathanson, although later Erd\H{o}s attributes this to Erd\H{o}s, S\'{a}rk\"{o}zy, and Szemer\'{e}di (but gives no reference), and claims a construction of an $A$ such that for all $\epsilon&#62;0$ and all large $N$\[\lvert \{1,\ldots,N\}\backslash B\rvert \ll_\epsilon N^{1/2+\epsilon},\]and yet there for all $\epsilon&#62;0$ there exist infinitely many $N$ where\[\lvert \{1,\ldots,N\}\backslash B\rvert \gg_\epsilon N^{1/3-\epsilon}.\]Erd\"{o}s and Freud investigated the finite analogue in \cite{ErFr91}, proving that there exists $A\subseteq \{1,\ldots,N\}$ such that the number of integers not representable in exactly one way as the sum of two elements from $A$ is $&#60;2^{3/2}N^{1/2}$, and suggest the constant $2^{3/2}$ is perhaps best possible.




References

[ErFr91] Erd\H{o}s, P. and Freud, R., On sums of a Sidon-sequence . J. Number Theory (1991), 196--205.

\noindent\rule{\linewidth}{0.4pt}


%% ============================================================
\section{Problem \#15}
\label{prob:15}

\noindent\textbf{Status:} OPEN \quad|\quad \textbf{Prize:} none \quad|\quad \textbf{Updated:} 2025-08-31

\noindent\textbf{Tags:} number theory, primes

\medskip

\noindent\textbf{Statement:}

Is it true that\[\sum_{n=1}^\infty(-1)^n\frac{n}{p_n}\]converges, where $p_n$ is the sequence of primes?




Erd\H{o}s suggested that a computer could be used to explore this, and did not see any other method to attack this. Tao \cite{Ta23} has proved that this series does converge assuming a strong form of the Hardy-Littlewood prime tuples conjecture. In \cite{Er98} Erd\H{o}s further conjectures that\[\sum_{n=1}^\infty (-1)^n \frac{1}{n(p_{n+1}-p_n)}\]converges and\[\sum_{n=1}^\infty (-1)^n \frac{1}{p_{n+1}-p_n}\]diverges. Weisenberg notes that the existence of infinitely many bounded gaps between primes (as proved by Zhang \cite{Zh14}) implies the latter series does not converge. Weisenberg also has an argument which shows that, assuming the Hardy-Littlewood prime $k$-tuples conjecture, the series is unbounded in at least one direction (positive or negative). Erd\H{o}s further conjectured that\[\sum_{n=1}^\infty (-1)^n \frac{1}{n(p_{n+1}-p_n)(\log\log n)^c}\]converges for every $c&#62;0$, and reports that he and Nathanson can prove that this series converges absolutely for $c&#62;2$ (and can show, conditional on 'hopeless' conjectures about the primes, that this sum does not converge absolutely for $c=2$). Sawhney has provided the following proof that this series converges absolutely for $c&#62;2$: note that, whenever $c&#62;1$, the contribution to the sum from gaps $p_{n+1}-p_n\geq \log n$ is convergent, so it suffices to consider only small gaps. The number of $n\leq X$ such that $p_{n+1}-p_n\in [\epsilon\log n,2\epsilon \log n)$ is bounded above by $\ll \epsilon X$ (this can be proved via the Selberg sieve). In particular, applying this bound for $\frac{1}{\log n}\leq \epsilon \leq 1$ of the shape $2^{-j}$ (of which there are at most $\log\log n$ possibilities) shows the desired convergence, since\[\sum \frac{1}{n(\log n)(\log\log n)^{c-1}}\]converges.




References

[Er98] Erd\H{o}s, Paul, Some of my new and almost new problems and results in combinatorial number theory . Number theory (Eger, 1996) (1998), 169-180.

[Ta23] Tao, T., The convergence of an alternating series of Erd\H{o}s, assuming the Hardy-Littlewood prime tuples conjecture . arXiv:2308.07205 (2023).

[Zh14] Zhang, Yitang, Bounded gaps between primes . Ann. of Math. (2) (2014), 1121--1174.

\noindent\rule{\linewidth}{0.4pt}


%% ============================================================
\section{Problem \#17}
\label{prob:17}

\noindent\textbf{Status:} OPEN \quad|\quad \textbf{Prize:} none \quad|\quad \textbf{Updated:} 2025-08-31

\noindent\textbf{Tags:} number theory, primes

\medskip
\noindent\textit{cluster primes}


\noindent\textbf{Statement:}

Are there infinitely many primes $p$ such that every even number $n\leq p-3$ can be written as a difference of primes $n=q_1-q_2$ where $q_1,q_2\leq p$?




The first prime without this property is $97$. The sequence of such primes is A038133 in the OEIS. These are called cluster primes. Blecksmith, Erd\H{o}s, and Selfridge \cite{BES99} proved that the number of such primes is\[\ll_A \frac{x}{(\log x)^A}\]for every $A&gt;0$, and Elsholtz \cite{El03} improved this to\[\ll x\exp(-c(\log\log x)^2)\]for every $c&lt;1/8$. This is discussed in problem C1 of Guy's collection \cite{Gu04}.




References

[BES99] Blecksmith, Richard and Erd\H{o}s, Paul and Selfridge, J. L., Cluster primes . Amer. Math. Monthly (1999), 43--48.

[El03] Elsholtz, Christian, On cluster primes . Acta Arith. (2003), 281--284.

[Gu04] Guy, Richard K., Unsolved problems in number theory . (2004), xviii+437.

\noindent\rule{\linewidth}{0.4pt}


%% ============================================================
\section{Problem \#18}
\label{prob:18}

\noindent\textbf{Status:} OPEN \quad|\quad \textbf{Prize:} none \quad|\quad \textbf{Updated:} 2025-08-31

\noindent\textbf{Tags:} number theory, divisors, factorials

\medskip
\noindent\textit{practical numbers}


\noindent\textbf{Statement:}

We call $m$ practical if every integer $1\leq n&#60;m$ is the sum of distinct divisors of $m$. If $m$ is practical then let $h(m)$ be such that $h(m)$ many divisors always suffice. Are there infinitely many practical $m$ such that\[h(m) &#60; (\log\log m)^{O(1)}?\]Is it true that $h(n!)&#60;n^{o(1)}$? Or perhaps even $h(n!)&#60;(\log n)^{O(1)}$?




It is easy to see that almost all numbers are not practical. Erd\H{o}s originally showed that $h(n!) &lt;n$. Vose \cite{Vo85} proved the existence of infinitely many practical $m$ such that $h(m)\ll (\log m)^{1/2}$. The sequence of practical numbers is A005153 in the OEIS. The reward of \$250 is offered in \cite{Er81h} for a proof or disproof of whether\[h(n) &lt;(\log \log n)^{O(1)}\]for infinitely many practical $n$. See also [304] and [825] .




References

[Er81h] Erd\H{o}s, P., Some problems and results on additive and multiplicative number theory . Analytic number theory (Philadelphia, Pa., 1980) (1981), 171-182.

[Vo85] Vose, Michael D., Egyptian fractions . Bull. London Math. Soc. (1985), 21-24.

\noindent\rule{\linewidth}{0.4pt}


%% ============================================================
\section{Problem \#20}
\label{prob:20}

\noindent\textbf{Status:} OPEN \quad|\quad \textbf{Prize:} \$1000 \quad|\quad \textbf{Updated:} 2025-08-31

\noindent\textbf{Tags:} combinatorics

\medskip
\noindent\textit{sunflower conjecture}


\noindent\textbf{Statement:}

Let $f(n,k)$ be minimal such that every family $\mathcal{F}$ of $n$-uniform sets with $\lvert \mathcal{F}\rvert \geq f(n,k)$ contains a $k$-sunflower. Is it true that\[f(n,k) &lt; c_k^n\]for some constant $c_k&gt;0$?




Erd\H{o}s and Rado \cite{ErRa60} originally proved $f(n,k)\leq (k-1)^nn!$. Kostochka \cite{Ko97} improved this slightly (in particular establishing an upper bound of $o(n!)$, for which Erd\H{o}s awarded him the consolation prize of \$100), but the bound stood at $n^{(1+o(1))n}$ for a long time until Alweiss, Lovett, Wu, and Zhang \cite{ALWZ20} proved\[f(n,k) &lt; (Ck\log n\log\log n)^n\]for some constant $C&gt;1$. This was refined slightly, independently by Rao \cite{Ra20}, Frankston, Kahn, Narayanan, and Park \cite{FKNP19}, and Bell, Chueluecha, and Warnke \cite{BCW21}, leading to the current record of\[f(n,k) &lt; (Ck\log n)^n\]for some constant $C&gt;1$. This proof was streamlined by Hu . A constant of $C=64$ was achieved in the presentation by Stoeckl . In \cite{Er81} Erd\H{o}s offered \$1000 for a proof or disproof even just in the special case when $k=3$, which he expected 'contains the whole difficulty'. He also wrote 'I really do not see why this question is so difficult'. The usual focus is on the regime where $k=O(1)$ is fixed (say $k=3$) and $n$ is large, although for the opposite regime Kostochka, R\"{o}dl, and Talysheva \cite{KRT99} have shown\[f(n,k)=(1+O_n(k^{-1/2^n}))k^n.\]




References

[ALWZ20] Alweiss, R. and Lovett, S. and Wu, K. and Zhang, J., Improved bounds for the sunflower lemma . (2020).

[BCW21] Bell, T. and Chueluecha, S. and Warnke, L., Note on sunflowers . Discret. Math. (2021).

[Er81] Erd\H{o}s, P., On the combinatorial problems which I would most like to see solved . Combinatorica (1981), 25-42.

[ErRa60] Erd\H{o}s, P. and Rado, R., Intersection theorems for systems of sets . J. London Math. Soc. (1960), 85-90.

[FKNP19] Frankston, K. and Kahn, J. and Narayanan, B. and Park, J., Thresholds versus fractional expectation-thresholds . CoRR (2019).

[KRT99] Kostochka, A. V. and R\"{o}dl, V. and Talysheva, L. A., On systems of small sets with no large $\Delta$-subsystems . Combin. Probab. Comput. (1999), 265-268.

[Ko97] Kostochka, A., A bound on the cardinality of families not containing $\Delta$-systems . (1997).

[Ra20] Rao, A., Coding for sunflowers . Discrete Analysis (2020).

\noindent\rule{\linewidth}{0.4pt}


%% ============================================================
\section{Problem \#25}
\label{prob:25}

\noindent\textbf{Status:} OPEN \quad|\quad \textbf{Prize:} none \quad|\quad \textbf{Updated:} 2025-08-31

\noindent\textbf{Tags:} number theory

\medskip

\noindent\textbf{Statement:}

Let $1\leq n_1&#60;n_2&#60;\cdots$ be an arbitrary sequence of integers, each with an associated residue class $a_i\pmod{n_i}$. Let $A$ be the set of integers $n$ such that for every $i$ either $n&#60;n_i$ or $n\not\equiv a_i\pmod{n_i}$. Must the logarithmic density of $A$ exist?




This is a special case of [486] .

\noindent\rule{\linewidth}{0.4pt}


%% ============================================================
\section{Problem \#28}
\label{prob:28}

\noindent\textbf{Status:} OPEN \quad|\quad \textbf{Prize:} \$500 \quad|\quad \textbf{Updated:} 2025-08-31

\noindent\textbf{Tags:} number theory, additive basis

\medskip

\noindent\textbf{Statement:}

If $A\subseteq \mathbb{N}$ is such that $A+A$ contains all but finitely many integers then $\limsup 1_A\ast 1_A(n)=\infty$.




Conjectured by Erd\H{o}s and Tur\'{a}n. They also suggest the stronger conjecture that $\limsup 1_A\ast 1_A(n)/\log n&gt;0$. Another stronger conjecture would be that the hypothesis $\lvert A\cap [1,N]\rvert \gg N^{1/2}$ for all large $N$ suffices. See also [40] , and [1145] for a stronger generalisation. This is discussed in problem C9 of Guy's collection \cite{Gu04}.




References

[Gu04] Guy, Richard K., Unsolved problems in number theory . (2004), xviii+437.

\noindent\rule{\linewidth}{0.4pt}


%% ============================================================
\section{Problem \#30}
\label{prob:30}

\noindent\textbf{Status:} OPEN \quad|\quad \textbf{Prize:} \$1000 \quad|\quad \textbf{Updated:} 2025-08-31

\noindent\textbf{Tags:} number theory, sidon sets, additive combinatorics

\medskip

\noindent\textbf{Statement:}

Let $h(N)$ be the maximum size of a Sidon set in $\{1,\ldots,N\}$. Is it true that, for every $\epsilon&#62;0$,\[h(N) = N^{1/2}+O_\epsilon(N^\epsilon)?\]




A problem of Erd\H{o}s and Tur\'{a}n. It may even be true that $h(N)=N^{1/2}+O(1)$, but Erd\H{o}s remarks this is perhaps too optimistic. Erd\H{o}s and Tur\'{a}n \cite{ErTu41} proved an upper bound of $N^{1/2}+O(N^{1/4})$, with an alternative proof by Lindstr\"{o}m \cite{Li69}. Both proofs in fact give\[h(N) \leq N^{1/2}+N^{1/4}+1.\]Balogh, F\"{u}redi, and Roy \cite{BFR21} improved the bound in the error term to $0.998N^{1/4}$. This was further optimised by O'Bryant \cite{OB22}. The current record is\[h(N)\leq N^{1/2}+0.98183N^{1/4}+O(1),\]due to Carter, Hunter, and O'Bryant \cite{CHO25}. Singer \cite{Si38} was the first to show that $h(N)\geq (1-o(1))N^{1/2}$ for all $N$. For a detailed survey of the literature we refer to \cite{OB04}. See also [241] and [840] . This problem is Problem 31 on Green's open problems list . This is discussed in problem C9 of Guy's collection \cite{Gu04}.




References

[BFR21] Balogh, J. and F\"{u}redi, Z. and Roy, S., An upper bound on the size of Sidon sets . arXiv:2103.15850 (2021).

[CHO25] Carter, D. and Hunter, Z. and O'Bryant, K., On the diameter of finite Sidon sets . Acta Math. Hungar. (2025), 108--126.

[ErTu41] Erd\H{o}s, P. and Tur\'{a}n, P., On a problem of Sidon in additive number theory, and on some related problems . J. London Math. Soc. (1941), 212-215.

[Gu04] Guy, Richard K., Unsolved problems in number theory . (2004), xviii+437.

[Li69] Lindstr\"{o}m, B., An inequality for $B_2$-sequences . J. Combinatorial Theory (1969), 211-212.

[OB04] O'Bryant, Kevin, A complete annotated bibliography of work related to Sidon sequences . Electron. J. Combin. (2004), 39.

[OB22] O'Bryant, K., On the size of finite Sidon sets . arXiv:2207.07800 (2022).

[Si38] Singer, James, A theorem in finite projective geometry and some applications to number theory . Trans. Amer. Math. Soc. (1938), 377--385.

\noindent\rule{\linewidth}{0.4pt}


%% ============================================================
\section{Problem \#32}
\label{prob:32}

\noindent\textbf{Status:} OPEN \quad|\quad \textbf{Prize:} none \quad|\quad \textbf{Updated:} 2025-08-31

\noindent\textbf{Tags:} number theory, additive basis

\medskip

\noindent\textbf{Statement:}

Is there a set $A\subset\mathbb{N}$ such that\[\lvert A\cap\{1,\ldots,N\}\rvert = o((\log N)^2)\]and such that every large integer can be written as $p+a$ for some prime $p$ and $a\in A$? Can the bound $O(\log N)$ be achieved? Must such an $A$ satisfy\[\liminf \frac{\lvert A\cap\{1,\ldots,N\}\rvert}{\log N}&#62; 1?\]




Such a set is called an additive complement to the primes. Erd\H{o}s \cite{Er54} proved that such a set $A$ exists with $\lvert A\cap\{1,\ldots,N\}\rvert\ll (\log N)^2$ (improving a previous result of Lorentz \cite{Lo54} who achieved $\ll (\log N)^3$). Wolke \cite{Wo96} has shown that such a bound is almost true, in that we can achieve $\ll (\log N)^{1+o(1)}$ if we only ask for almost all integers to be representable. Kolountzakis \cite{Ko96} improved this to $\ll (\log N)(\log\log N)$, and Ruzsa \cite{Ru98c} further improved this to $\ll \omega(N)\log N$ for any $\omega\to \infty$. The answer to the third question is yes: Ruzsa \cite{Ru98c} has shown that we must have\[\liminf \frac{\lvert A\cap\{1,\ldots,N\}\rvert}{\log N}\geq e^\gamma\approx 1.781.\]This is discussed in problem E1 of Guy's collection \cite{Gu04}, where it is stated that Erd\H{o}s offered \$50 for determining whether $O(\log N)$ can be achieved.




References

[Er54] Erd\H{o}s, Paul, Some results on additive number theory . Proc. Amer. Math. Soc. (1954), 847-853.

[Gu04] Guy, Richard K., Unsolved problems in number theory . (2004), xviii+437.

[Ko96] Kolountzakis, Mihail N., On the additive complements of the primes and sets of similar growth . Acta Arith. (1996), 1--8.

[Lo54] Lorentz, G. G., On a problem of additive number theory . Proc. Amer. Math. Soc. (1954), 838-841.

[Ru98c] Ruzsa, Imre Z., On the additive completion of primes . Acta Arith. (1998), 269-275.

[Wo96] Wolke, Dieter, On a problem of Erd\H{o}s in additive number theory . J. Number Theory (1996), 209-213.

\noindent\rule{\linewidth}{0.4pt}


%% ============================================================
\section{Problem \#33}
\label{prob:33}

\noindent\textbf{Status:} OPEN \quad|\quad \textbf{Prize:} none \quad|\quad \textbf{Updated:} 2025-08-31

\noindent\textbf{Tags:} number theory, additive basis

\medskip

\noindent\textbf{Statement:}

Let $A\subset\mathbb{N}$ be such that every large integer can be written as $n^2+a$ for some $a\in A$ and $n\geq 0$. What is the smallest possible value of\[\limsup \frac{\lvert A\cap\{1,\ldots,N\}\rvert}{N^{1/2}}?\]Is\[\liminf \frac{\lvert A\cap\{1,\ldots,N\}\rvert}{N^{1/2}}&#62;1?\]




Such a set $A$ is called an additive complement of the set of squares. Erd\H{o}s observed that there exist $A$ for which the $\limsup$ is finite and $&gt;1$. Moser \cite{Mo65} proved that, for any such $A$,\[\liminf \frac{\lvert A\cap\{1,\ldots,N\}\rvert}{N^{1/2}}&gt;1.06.\]The best-known lower bound is\[\liminf \frac{\lvert A\cap\{1,\ldots,N\}\rvert}{N^{1/2}}\geq\frac{4}{\pi}\approx 1.273\]proved by Cilleruelo \cite{Ci93}, Habsieger \cite{Ha95}, and Balasubramanian and Ramana \cite{BaRa01}. The problem of minimising the $\limsup$ appears to have been much less studied. van Doorn has a construction of such an $A$ in which, for all $N$,\[\frac{\lvert A\cap\{1,\ldots,N\}\rvert}{N^{1/2}}&lt; 2\phi^{5/2}\approx 6.66,\]where $\phi=\frac{1+\sqrt{5}}{2}$ is the golden ratio.




References

[BaRa01] Balasubramanian, R. and Ramana, D. S., Additive complements of the squares . C. R. Math. Acad. Sci. Soc. R. Can. (2001), 6--11.

[Ci93] Cilleruelo, Javier, The additive completion of {$k$}th-powers . J. Number Theory (1993), 237--243.

[Ha95] Habsieger, Laurent, On the additive completion of polynomial sets . J. Number Theory (1995), 130--135.

[Mo65] Moser, Leo, On the additive completion of sets of integers . (1965), 175--180.

\noindent\rule{\linewidth}{0.4pt}


%% ============================================================
\section{Problem \#36}
\label{prob:36}

\noindent\textbf{Status:} OPEN \quad|\quad \textbf{Prize:} none \quad|\quad \textbf{Updated:} 2025-08-31

\noindent\textbf{Tags:} number theory, additive combinatorics

\medskip
\noindent\textit{minimum overlap problem}


\noindent\textbf{Statement:}

Find the optimal constant $c&#62;0$ such that the following holds. For all sufficiently large $N$, if $A\sqcup B=\{1,\ldots,2N\}$ is a partition into two equal parts, so that $\lvert A\rvert=\lvert B\rvert=N$, then there is some $x$ such that the number of solutions to $a-b=x$ with $a\in A$ and $b\in B$ is at least $cN$.




The minimum overlap problem . The example (with $N$ even) $A=\{N/2+1,\ldots,3N/2\}$ shows that $c\leq 1/2$ (indeed, Erd\H{o}s initially conjectured that $c=1/2$). The lower bound of $c\geq 1/4$ is trivial, and Scherk improved this to $1-1/\sqrt{2}=0.29\cdots$. The current records are\[0.379005 &lt; c &lt; 0.380876,\]the lower bound due to White \cite{Wh22} and the upper bound due to the TTT-Discover LLM \cite{YKLBMWKCZGS26}, improving slightly on earlier bounds due to AlphaEvolve \cite{GGTW25} and Haugland \cite{Ha16}. This is discussed in problem C17 of Guy's collection \cite{Gu04}.




References

[GGTW25] B. Georgiev, J. G\'{o}mez-Serrano, T. Tao, and A. Wagner, Mathematical exploration and discovery at scale . arXiv:2511.02864 (2025).

[Gu04] Guy, Richard K., Unsolved problems in number theory . (2004), xviii+437.

[Ha16] Haugland, J. K., The minimum overlap problem revisited . arXiv:1609.08000 (2016).

[Wh22] White, E. P., Erd\H{o}s' minimum overlap problem . arXiv:2201.05704 (2022).

[YKLBMWKCZGS26] M. Yuksekgonul, D. Koceja, X. Li, F. Bianchi, J. McCaleb, X. Wang, J. Kautz, Y. Choi, J. Zou, C. Guestrin, and Y. Sun, Learning to Discover at Test Time . https://test-time-training.github.io/discover.pdf (2026).

\noindent\rule{\linewidth}{0.4pt}


%% ============================================================
\section{Problem \#39}
\label{prob:39}

\noindent\textbf{Status:} OPEN \quad|\quad \textbf{Prize:} \$500 \quad|\quad \textbf{Updated:} 2025-08-31

\noindent\textbf{Tags:} number theory, sidon sets, additive combinatorics

\medskip

\noindent\textbf{Statement:}

Is there an infinite Sidon set $A\subset \mathbb{N}$ such that\[\lvert A\cap \{1\ldots,N\}\rvert \gg_\epsilon N^{1/2-\epsilon}\]for all $\epsilon&#62;0$?




The trivial greedy construction achieves $\gg N^{1/3}$. The first improvement on this was achieved by Ajtai, Koml\'{o}s, and Szemer\'{e}di \cite{AKS81b}, who found an infinite Sidon set with growth rate $\gg (N\log N)^{1/3}$. The current best bound of $\gg N^{\sqrt{2}-1+o(1)}$ is due to Ruzsa \cite{Ru98}. Erd\H{o}s \cite{Er73} had offered \$25 for any construction which achieves $N^{c}$ for some $c&#62;1/3$. Later he (\cite{Er77c} and \cite{Er80}) offered \$100 for a construction which achieves $\omega(N)N^{1/3}$ for some $\omega(N)\to \infty$. Erd\H{o}s proved that for every infinite Sidon set $A$ we have\[\liminf \frac{\lvert A\cap \{1,\ldots,N\}\rvert}{N^{1/2}}=0.\]Erd\H{o}s and R\'{e}nyi have constructed, for any $\epsilon&#62;0$, a set $A$ such that\[\lvert A\cap \{1\ldots,N\}\rvert \gg_\epsilon N^{1/2-\epsilon}\]for all large $N$ and $1_A\ast 1_A(n)\ll_\epsilon 1$ for all $n$. This is discussed in problem C9 of Guy's collection \cite{Gu04}.




References

[AKS81b] Ajtai, Mikl\'os and Koml\'os, J\'{a}nos and Szemer\'{e}di, Endre, A dense infinite Sidon sequence . European J. Combin. (1981), 1--11.

[Er73] Erd\H{o}s, P., Problems and results on combinatorial number theory . A survey of combinatorial theory (Proc. Internat. Sympos., Colorado State Univ., Fort Collins, Colo., 1971) (1973), 117-138.

[Er77c] Erd\H{o}s, Paul, Problems and results on combinatorial number theory. III . Number theory day (Proc. Conf., Rockefeller Univ., New York, 1976) (1977), 43-72.

[Er80] Erd\H{o}s, Paul, A survey of problems in combinatorial number theory . Ann. Discrete Math. (1980), 89-115.

[Gu04] Guy, Richard K., Unsolved problems in number theory . (2004), xviii+437.

[Ru98] Ruzsa, Imre Z., An infinite Sidon sequence . J. Number Theory (1998), 63-71.

\noindent\rule{\linewidth}{0.4pt}


%% ============================================================
\section{Problem \#40}
\label{prob:40}

\noindent\textbf{Status:} OPEN \quad|\quad \textbf{Prize:} \$500 \quad|\quad \textbf{Updated:} 2025-08-31

\noindent\textbf{Tags:} number theory, additive basis

\medskip

\noindent\textbf{Statement:}

For what functions $g(N)\to \infty$ is it true that\[\lvert A\cap \{1,\ldots,N\}\rvert \gg \frac{N^{1/2}}{g(N)}\]implies $\limsup 1_A\ast 1_A(n)=\infty$?




This is a stronger form of the Erd\H{o}s-Tur\'{a}n conjecture [28] (since establishing this for any function $g(N)\to \infty$ would imply a positive solution to [28] ).

\noindent\rule{\linewidth}{0.4pt}


%% ============================================================
\section{Problem \#41}
\label{prob:41}

\noindent\textbf{Status:} OPEN \quad|\quad \textbf{Prize:} \$500 \quad|\quad \textbf{Updated:} 2025-08-31

\noindent\textbf{Tags:} number theory, sidon sets, additive combinatorics

\medskip

\noindent\textbf{Statement:}

Let $A\subset\mathbb{N}$ be an infinite set such that the triple sums $a+b+c$ are all distinct for $a,b,c\in A$ (aside from the trivial coincidences). Is it true that\[\liminf \frac{\lvert A\cap \{1,\ldots,N\}\rvert}{N^{1/3}}=0?\]




Erd\H{o}s proved that if the pairwise sums $a+b$ are all distinct aside from the trivial coincidences then\[\liminf \frac{\lvert A\cap \{1,\ldots,N\}\rvert}{N^{1/2}}=0.\]This is discussed in problem C11 of Guy's collection \cite{Gu04}, in which Guy says Erd\H{o}s offered \$500 for the general problem of whether, for all $h\geq 2$,\[\liminf \frac{\lvert A\cap \{1,\ldots,N\}\rvert}{N^{1/h}}=0\]whenever the sum of $h$ terms in $A$ are distinct. This was proved for $h=4$ by Nash \cite{Na89} and for all even $h$ by Chen \cite{Ch96b}.




References

[Ch96b] Chen, Sheng, A note on {$B_{2k}$} sequences . J. Number Theory (1996), 1--3.

[Gu04] Guy, Richard K., Unsolved problems in number theory . (2004), xviii+437.

[Na89] Nash, John C. M., On {$B_4$}-sequences . Canad. Math. Bull. (1989), 446--449.

\noindent\rule{\linewidth}{0.4pt}


%% ============================================================
\section{Problem \#44}
\label{prob:44}

\noindent\textbf{Status:} OPEN \quad|\quad \textbf{Prize:} none \quad|\quad \textbf{Updated:} 2025-08-31

\noindent\textbf{Tags:} number theory, sidon sets, additive combinatorics

\medskip

\noindent\textbf{Statement:}

Let $N\geq 1$ and $A\subset \{1,\ldots,N\}$ be a Sidon set. Is it true that, for any $\epsilon&#62;0$, there exist $M$ and $B\subset \{N+1,\ldots,M\}$ (which may depend on $N,A,\epsilon$) such that $A\cup B\subset \{1,\ldots,M\}$ is a Sidon set of size at least $(1-\epsilon)M^{1/2}$?




See also [329] and [707] (indeed a positive solution to [707] implies a positive solution to this problem, which in turn implies a positive solution to [329] ). This is discussed in problem C9 of Guy's collection \cite{Gu04}.




References

[Gu04] Guy, Richard K., Unsolved problems in number theory . (2004), xviii+437.

\noindent\rule{\linewidth}{0.4pt}


%% ============================================================
\section{Problem \#50}
\label{prob:50}

\noindent\textbf{Status:} OPEN \quad|\quad \textbf{Prize:} \$250 \quad|\quad \textbf{Updated:} 2025-08-31

\noindent\textbf{Tags:} number theory

\medskip

\noindent\textbf{Statement:}

Schoenberg proved that for every $c\in [0,1]$ the density of\[\{ n\in \mathbb{N} : \phi(n)&#60;cn\}\]exists. Let this density be denoted by $f(c)$. Is it true that there are no $x$ such that $f'(x)$ exists and is positive?




Erd\H{o}s \cite{Er95} could prove the distribution function is purely singular.




References

[Er95] Erd\H{o}s, Paul, Some of my favourite problems in number theory, combinatorics, and geometry . Resenhas (1995), 165-186.

\noindent\rule{\linewidth}{0.4pt}


%% ============================================================
\section{Problem \#51}
\label{prob:51}

\noindent\textbf{Status:} OPEN \quad|\quad \textbf{Prize:} none \quad|\quad \textbf{Updated:} 2025-08-31

\noindent\textbf{Tags:} number theory

\medskip

\noindent\textbf{Statement:}

Is there an infinite set $A\subset \mathbb{N}$ such that for every $a\in A$ there is an integer $n$ such that $\phi(n)=a$, and yet if $n_a$ is the smallest such integer then $n_a/a\to \infty$ as $a\to\infty$?




Carmichael has asked whether there is an integer $t$ for which $\phi(n)=t$ has exactly one solution. Erd\H{o}s has proved that if such a $t$ exists then there must be infinitely many such $t$. See also [694] . This is discussed in problems B36 and B39 of Guy's collection \cite{Gu04}.




References

[Gu04] Guy, Richard K., Unsolved problems in number theory . (2004), xviii+437.

\noindent\rule{\linewidth}{0.4pt}


%% ============================================================
\section{Problem \#52}
\label{prob:52}

\noindent\textbf{Status:} OPEN \quad|\quad \textbf{Prize:} \$250 \quad|\quad \textbf{Updated:} 2026-05-28

\noindent\textbf{Tags:} number theory, additive combinatorics

\medskip
\noindent\textit{sum-product problem}


\noindent\textbf{Statement:}

Let $A$ be a finite set of integers. Is it true that for every $\epsilon&#62;0$\[\max( \lvert A+A\rvert,\lvert AA\rvert)\gg_\epsilon \lvert A\rvert^{2-\epsilon}?\]




The sum-product problem. Erd\H{o}s and Szemer\'{e}di \cite{ErSz83} proved a lower bound of $\lvert A\rvert^{1+c}$ for some constant $c&gt;0$, and an upper bound of\[\lvert A\rvert^2 \exp\left(-c\frac{\log\lvert A\rvert}{\log\log \lvert A\rvert}\right)\]for some constant $c&gt;0$. The lower bound has been improved a number of times. The current record is\[\max( \lvert A+A\rvert,\lvert AA\rvert)\gg\lvert A\rvert^{\frac{1962}{1469}-o(1)}\]due to Cushman \cite{Cu25} (note $1962/1469=1.3356\cdots$). A complete history of sum-product bounds can be found at this webpage . There is likely nothing special about the integers in this question, and indeed Erd\H{o}s and Szemer\'{e}di also ask a similar question about finite sets of real or complex numbers. The current best lower bound for sets of reals is the same bound of Cushman above; the original conjecture is false for sets of reals, see below. The best bound for complex numbers is\[\max( \lvert A+A\rvert,\lvert AA\rvert)\gg\lvert A\rvert^{\frac{4}{3}+c}\]for some absolute constant $c&gt;0$, due to Basit and Lund \cite{BaLu19}. One can in general ask this question in any setting where addition and multiplication are defined (once one avoids any trivial obstructions such as zero divisors or finite subfields). For example, it makes sense for subsets of finite fields. The current record is that there exists $c&gt;0$ such that if $A\subseteq \mathbb{F}_p$ with $\lvert A\rvert &lt;p^{c}$ then\[\max( \lvert A+A\rvert,\lvert AA\rvert)\gg\lvert A\rvert^{\frac{5}{4}+o(1)},\]due to Mohammadi and Stevens \cite{MoSt23}. There is also a natural generalisation to higher-fold sum and product sets. For example, in \cite{ErSz83} (and in \cite{Er91}) Erd\H{o}s and Szemer\'{e}di also conjecture that for any $m\geq 2$ and finite set of integers $A$\[\max( \lvert mA\rvert,\lvert A^m\rvert)\gg \lvert A\rvert^{m-o(1)}.\]The conjecture is false for $A\subset \mathbb{R}$, as proved by Bloom, Sawin, Schildkraut, and Zhelezov \cite{BSSZ26}, who construct arbitrarily large $A$ such that\[\max(\lvert A+A\rvert, \lvert AA\rvert) \leq \lvert A\rvert^{2-c}\]for some absolute constant $c&gt;0$. A similar construction also disproves this for small subsets of $\mathbb{F}_p$, and disproves the above higher-fold variant for sets of real numbers. See [53] for more on this generalisation and [808] for a stronger form of the original conjecture. See also [818] for a special case.




References

[BSSZ26] T. F. Bloom, W. Sawin, C. Schildkraut, D. Zhelezov, The sum-product conjecture is false for real numbers . arXiv:2605.28781 (2026).

[BaLu19] Basit, Abdul and Lund, Ben, An improved sum-product bound for quaternions . SIAM J. Discrete Math. (2019), 1044--1060.

[Cu25] A. Cushman, A Note on the Sum-Product Problem and the Convex Sumset Problem . arXiv:2512.13849 (2025).

[Er91] Erd\H{o}s, P., Problems and results in combinatorial analysis and combinatorial number theory . Graph theory, combinatorics, and applications, Vol. 1 (Kalamazoo, MI, 1988) (1991), 397-406.

[ErSz83] Erd\H{o}s, P. and Szemer\'{e}di, E., On sums and products of integers . Studies in pure mathematics (1983), 213-218.

[MoSt23] Mohammadi, Ali and Stevens, Sophie, Attaining the exponent 5/4 for the sum-product problem in finite fields . Int. Math. Res. Not. IMRN (2023), 3516--3532.

\noindent\rule{\linewidth}{0.4pt}


%% ============================================================
\section{Problem \#60}
\label{prob:60}

\noindent\textbf{Status:} OPEN \quad|\quad \textbf{Prize:} none \quad|\quad \textbf{Updated:} 2025-08-31

\noindent\textbf{Tags:} graph theory, cycles

\medskip

\noindent\textbf{Statement:}

Does every graph on $n$ vertices with $&#62;\mathrm{ex}(n;C_4)$ edges contain $\gg n^{1/2}$ many copies of $C_4$?




Conjectured by Erd\H{o}s and Simonovits, who could not even prove that at least $2$ copies of $C_4$ are guaranteed. The behaviour of $\mathrm{ex}(n;C_4)$ is the subject of [765] . He, Ma, and Yang \cite{HeMaYa21} have proved this conjecture when $n=q^2+q+1$ for some even integer $q$.




References

[HeMaYa21] He, J. and Ma, J. and Yang, T., Some extremal results on 4-cycles . Journal of Combinatorial Theory B (2021).

\noindent\rule{\linewidth}{0.4pt}


%% ============================================================
\section{Problem \#61}
\label{prob:61}

\noindent\textbf{Status:} OPEN \quad|\quad \textbf{Prize:} none \quad|\quad \textbf{Updated:} 2025-08-31

\noindent\textbf{Tags:} graph theory

\medskip

\noindent\textbf{Statement:}

For any graph $H$ is there some $c=c(H)&#62;0$ such that every graph $G$ on $n$ vertices that does not contain $H$ as an induced subgraph contains either a complete graph or independent set on $\geq n^c$ vertices?




Conjectured by Erd\H{o}s and Hajnal \cite{ErHa89}, who proved that a complete graph or independent set must exist on\[\geq \exp(c_H\sqrt{\log n})\]many vertices, where $c_H&gt;0$ is some constant. This was improved by Buci\'{c}, Nguyen, Scott, and Seymour \cite{BNSS23} to\[\geq \exp(c_H\sqrt{\log n\log\log n}).\]The conjecture has been proved: {UL} {LI} for $H$ with $\leq 4$ vertices by Erd\H{o}s and Hajnal \cite{ErHa89}.{/LI} {LI} for the 'bull' graph on $5$ vertices (a $P_4$ with a new vertex adjacent to the two middle vertices) by Chudnovsky and Safra \cite{ChSa08}.{/LI} {LI} for $C_5$ by Chudnovsky, Scott, Seymour, and Spirkl \cite{CSSS23}.{/LI} {LI} for $P_5$ by Nguyen, Scott, and Seymour \cite{NSS26}.{/LI} {/UL} Alon, Pach, and Solymosi \cite{APS01} have proved that the family of $H$ for which the conjecture holds is closed under vertex substitution; together with the three final results above this means the conjecture is now known for all $H$ on $5$ vertices. Nguyen, Scott, and Seymour \cite{NSS24} have proved that if $H$ is a path then all $H$-free graphs on $n$ vertices contain either a complete graph or independent set on\[\geq 2^{(\log n)^{1-o(1)}}\]many vertices. A detailed account of this problem, including the proof of some special cases, is given in the PhD thesis of Nguyen. This problem is #80 in Extremal Graph Theory in the graphs problem collection.




References

[APS01] Alon, Noga and Pach, J\'{a}nos and Solymosi, J\'ozsef, Ramsey-type theorems with forbidden subgraphs . Combinatorica (2001), 155--170.

[BNSS23] Buci\'C, M. and Nguyen, T. and Scott, A. and Seymour, P., A loglog step towards Erdos-Hajnal . arXiv:2301.10147 (2023).

[CSSS23] Chudnovsky, Maria and Scott, Alex and Seymour, Paul and Spirkl, Sophie, Erd\H{o}s-{H}ajnal for graphs with no 5-hole . Proc. Lond. Math. Soc. (3) (2023), 997--1014.

[ChSa08] Chudnovsky, Maria and Safra, Shmuel, The {E}rd\H os-{H}ajnal conjecture for bull-free graphs . J. Combin. Theory Ser. B (2008), 1301--1310.

[ErHa89] Erd\H{o}s, P. and Hajnal, A., Ramsey-type theorems . Discrete Appl. Math. (1989), 37-52.

[NSS24] Nguyen, Tung and Scott, Alex and Seymour, Paul, On a problem of {E}l-{Z}ahar and {E}rd\H{o}s . J. Combin. Theory Ser. B (2024), 211--222.

[NSS26] Nguyen, Tung and Scott, Alex and Seymour, Paul, Induced subgraph density. {VII}. The five-vertex path . Proc. Lond. Math. Soc. (3) (2026), Paper No. e70133.

\noindent\rule{\linewidth}{0.4pt}


%% ============================================================
\section{Problem \#62}
\label{prob:62}

\noindent\textbf{Status:} OPEN \quad|\quad \textbf{Prize:} none \quad|\quad \textbf{Updated:} 2025-08-31

\noindent\textbf{Tags:} graph theory

\medskip

\noindent\textbf{Statement:}

If $G_1,G_2$ are two graphs with chromatic number $\aleph_1$ then must there exist a graph $G$ whose chromatic number is $4$ (or even $\aleph_0$) which is a subgraph of both $G_1$ and $G_2$?




Erd\H{o}s also asked \cite{Er87} about finding a common subgraph $H$ (with chromatic number either $4$ or $\aleph_0$) in any finite collection of graphs with chromatic number $\aleph_1$. Every graph with chromatic number $\aleph_1$ contains all sufficiently large odd cycles (which have chromatic number $3$), see [594] . This was proved by Erd\H{o}s, Hajnal, and Shelah \cite{EHS74}. Erd\H{o}s wrote \cite{Er87} that 'probably' every graph with chromatic number $\aleph_1$ contains as subgraphs all graphs with chromatic number $4$ with sufficiently large girth.




References

[EHS74] Erd\H{o}s, P. and Hajnal, A. and Shelah, S., On some general properties of chromatic numbers . Topics in topology (Proc. Colloq., Keszthely, 1972) (1974), 243-255.

[Er87] Erd\H{o}s, P., Some problems on finite and infinite graphs . Logic and combinatorics (Arcata, Calif., 1985) (1987), 223-228.

\noindent\rule{\linewidth}{0.4pt}


%% ============================================================
\section{Problem \#65}
\label{prob:65}

\noindent\textbf{Status:} OPEN \quad|\quad \textbf{Prize:} none \quad|\quad \textbf{Updated:} 2025-08-31

\noindent\textbf{Tags:} graph theory, cycles

\medskip

\noindent\textbf{Statement:}

Let $G$ be a graph with $n$ vertices and $kn$ edges, and $a_1&#60;a_2&#60;\cdots $ be the lengths of cycles in $G$. Is it true that\[\sum\frac{1}{a_i}\gg \log k?\]Is the sum $\sum\frac{1}{a_i}$ minimised when $G$ is a complete bipartite graph?




A problem of Erd\H{o}s and Hajnal. Gy\'{a}rf\'{a}s, Koml\'{o}s, and Szemer\'{e}di \cite{GKS84} have proved that this sum is $\gg \log k$, so that only the second question remains. Liu and Montgomery \cite{LiMo20} have proved the asymptotically sharp lower bound of $\geq (\tfrac{1}{2}-o(1))\log k$. Montgomery has written a survey including this problem, in which he mentions forthcoming work of himself, Milojevi\'{c}, Pokrovskiy, and Sudakov which proves that, if $k$ is sufficiently large, then $\sum\frac{1}{a_i}$ is maximised when $G$ is a complete bipartite graph. This problem is #65 in Extremal Graph Theory in the graphs problem collection. See also [57] .




References

[GKS84] Gy\'{a}rf\'{a}s, A. and Koml\'{o}s, J. and Szemer\'{e}di, E., On the distribution of cycle lengths in graphs . J. Graph Theory (1984), 441-462.

[LiMo20] Liu, Hong and Montgomery, Richard, A solution to Erd\H{o}s and Hajnal's odd cycle problem . arXiv:2010.15802 (2020).

\noindent\rule{\linewidth}{0.4pt}


%% ============================================================
\section{Problem \#66}
\label{prob:66}

\noindent\textbf{Status:} OPEN \quad|\quad \textbf{Prize:} \$500 \quad|\quad \textbf{Updated:} 2025-08-31

\noindent\textbf{Tags:} number theory, additive basis

\medskip

\noindent\textbf{Statement:}

Is there $A\subseteq \mathbb{N}$ such that\[\lim_{n\to \infty}\frac{1_A\ast 1_A(n)}{\log n}\]exists and is $\neq 0$?




A suitably constructed random set has this property if we are allowed to ignore an exceptional set of density zero. The challenge is obtaining this with no exceptional set. Erd\H{o}s believed the answer should be no. In \cite{Er80} he explicitly asked whether there exists such a set where the limit is $1$ (and did not believe there existed such a sequence). Erd\H{o}s and S\'{a}rk\"{o}zy proved that\[\frac{\lvert 1_A\ast 1_A(n)-\log n\rvert}{\sqrt{\log n}}\to 0\]is impossible. Erd\H{o}s suggests it may even be true that the $\liminf$ and $\limsup$ of $1_A\ast 1_A(n)/\log n$ are always separated by some absolute constant. Horv\'{a}th \cite{Ho07} proved that\[\lvert 1_A\ast 1_A(n)-\log n\rvert \leq (1-\epsilon)\sqrt{\log n}\]cannot hold for all large $n$.




References

[Er80] Erd\H{o}s, Paul, A survey of problems in combinatorial number theory . Ann. Discrete Math. (1980), 89-115.

[Ho07] G. Horv\'{a}th, An improvement of a theorem of Erd\H{o}s and S\'{a}rk\"{o}zy . Pollack Periodica (2007), 155-161.

\noindent\rule{\linewidth}{0.4pt}


%% ============================================================
\section{Problem \#68}
\label{prob:68}

\noindent\textbf{Status:} OPEN \quad|\quad \textbf{Prize:} none \quad|\quad \textbf{Updated:} 2025-08-31

\noindent\textbf{Tags:} number theory, irrationality

\medskip

\noindent\textbf{Statement:}

Is\[\sum_{n\geq 2}\frac{1}{n!-1}\]irrational?




The decimal expansion is A331373 in the OEIS. Weisenberg has observed that this sum can also be written as\[\sum_{k\geq 1}\sum_{n\geq 2}\frac{1}{(n!)^k}.\]Erd\H{o}s \cite{Er88c} notes that $\sum \frac{1}{n!+t}$ should be transcendental for every integer $t$.




References

[Er88c] Erd\H{o}s, P., On the irrationality of certain series: problems and results . New advances in transcendence theory (Durham, 1986) (1988), 102-109.

\noindent\rule{\linewidth}{0.4pt}


%% ============================================================
\section{Problem \#70}
\label{prob:70}

\noindent\textbf{Status:} OPEN \quad|\quad \textbf{Prize:} none \quad|\quad \textbf{Updated:} 2025-08-31

\noindent\textbf{Tags:} graph theory, ramsey theory, set theory

\medskip

\noindent\textbf{Statement:}

Let $\mathfrak{c}$ be the ordinal of the real numbers, $\beta$ be any countable ordinal, and $2\leq n&#60;\omega$. Is it true that $\mathfrak{c}\to (\beta, n)_2^3$?




Erd\H{o}s and Rado proved that $\mathfrak{c}\to (\omega+n,4)_2^3$ for any $2\leq n&#60;\omega$.

\noindent\rule{\linewidth}{0.4pt}


%% ============================================================
\section{Problem \#74}
\label{prob:74}

\noindent\textbf{Status:} OPEN \quad|\quad \textbf{Prize:} \$500 \quad|\quad \textbf{Updated:} 2025-08-31

\noindent\textbf{Tags:} graph theory, chromatic number, cycles

\medskip

\noindent\textbf{Statement:}

Let $f(n)\to \infty$ (possibly very slowly). Is there a graph of infinite chromatic number such that every finite subgraph on $n$ vertices can be made bipartite by deleting at most $f(n)$ edges?




Conjectured by Erd\H{o}s, Hajnal, and Szemer\'{e}di \cite{EHS82}. R\"{o}dl \cite{Ro82} has proved this for hypergraphs, and also proved there is such a graph (with chromatic number $\aleph_0$) if $f(n)=\epsilon n$ for any fixed constant $\epsilon&gt;0$. It is open even for $f(n)=\sqrt{n}$. Erd\H{o}s offered \$500 for a proof but only \$250 for a counterexample. This fails (even with $f(n)\gg n$) if the graph has chromatic number $\aleph_1$ (see [111] ).




References

[EHS82] Erd\H{o}s, P. and Hajnal, A. and Szemer\'{e}di, E., On almost bipartite large chromatic graphs . Theory and practice of combinatorics (1982), 117-123.

[Ro82] R\"{o}dl, Vojt\vEch, Nearly bipartite graphs with large chromatic number . Combinatorica (1982), 377-383.

\noindent\rule{\linewidth}{0.4pt}


%% ============================================================
\section{Problem \#75}
\label{prob:75}

\noindent\textbf{Status:} OPEN \quad|\quad \textbf{Prize:} none \quad|\quad \textbf{Updated:} 2025-08-31

\noindent\textbf{Tags:} graph theory, chromatic number

\medskip

\noindent\textbf{Statement:}

Is there a graph of chromatic number $\aleph_1$ with $\aleph_1$ vertices such that for all $\epsilon&#62;0$ if $n$ is sufficiently large and $H$ is a subgraph on $n$ vertices then $H$ contains an independent set of size $&#62;n^{1-\epsilon}$? What about an independent set of size $\gg n$?




Conjectured by Erdős, Hajnal, and Szemerédi \cite{EHS82}. In \cite{Er95} Erd\H{o}s asks this without the condition that the graph also have $\aleph_1$ vertices, but this is an oversight, since already in \cite{EHS82} they provide such a construction. In \cite{Er95d} Erd\H{o}s offers \$1000 for a complete solution to all problems of this type (for example including also [74] ), and a 'generous reward for any significant partial results'. See also [750] .




References

[EHS82] Erd\H{o}s, P. and Hajnal, A. and Szemer\'{e}di, E., On almost bipartite large chromatic graphs . Theory and practice of combinatorics (1982), 117-123.

[Er95] Erd\H{o}s, Paul, Some of my favourite problems in number theory, combinatorics, and geometry . Resenhas (1995), 165-186.

[Er95d] Erd\H{o}s, Paul, On some problems in combinatorial set theory . Publ. Inst. Math. (Beograd) (N.S.) (1995), 61-65.

\noindent\rule{\linewidth}{0.4pt}


%% ============================================================
\section{Problem \#77}
\label{prob:77}

\noindent\textbf{Status:} OPEN \quad|\quad \textbf{Prize:} \$250 \quad|\quad \textbf{Updated:} 2025-08-31

\noindent\textbf{Tags:} graph theory, ramsey theory

\medskip

\noindent\textbf{Statement:}

If $R(k)$ is the Ramsey number for $K_k$, the minimal $n$ such that every $2$-colouring of the edges of $K_n$ contains a monochromatic copy of $K_k$, then find the value of\[\lim_{k\to \infty}R(k)^{1/k}.\]




Erd\H{o}s offered \$100 for just a proof of the existence of this constant, without determining its value. He also offered \$1000 for a proof that the limit does not exist, but says 'this is really a joke as [it] certainly exists'. (In \cite{Er88} he raises this prize to \$10000). Erd\H{o}s proved\[\sqrt{2}\leq \liminf_{k\to \infty}R(k)^{1/k}\leq \limsup_{k\to \infty}R(k)^{1/k}\leq 4.\]The upper bound has been improved to $4-\tfrac{1}{128}$ by Campos, Griffiths, Morris, and Sahasrabudhe \cite{CGMS23}. This was improved to $3.7992\cdots$ by Gupta, Ndiaye, Norin, and Wei \cite{GNNW24}. A shorter and simpler proof of an upper bound of the strength $4-c$ for some constant $c&gt;0$ (and a generalisation to the case of more than two colours) was given by Balister, Bollob\'{a}s, Campos, Griffiths, Hurley, Morris, Sahasrabudhe, and Tiba \cite{BBCGHMST24}. In \cite{Er93} Erd\H{o}s writes 'I have no idea what the value of $\lim R(k)^{1/k}$ should be, perhaps it is $2$ but we have no real evidence for this.' This problem is #3 in Ramsey Theory in the graphs problem collection. See also [1029] for a problem concerning a lower bound for $R(k)$ and discussion of lower bounds in general. The limit in this question is closely related to the limit in [627] . A famous quote of Erd\H{o}s concerns the difficulty of finding exact values for $R(k)$. This is often repeated in the words of Spencer, who phrased it as an alien attacking race. The earliest such quote in a paper of Erd\H{o}s I have found is in \cite{Er93}, where he writes: 'Sometime ago, I made the following joke. If an evil spirit would appear and say "unless you give me the value of $R(5)$ within a year, I will exterminate humanity", then our best bet would be perhaps to get all our computers working on $R(5)$ and we probably would get its value in a year. If he would ask for $R(6)$, the best strategy probably would be to destroy it before it can destroy us. If we would be so clever that we could give the answer by mathematics, we would just tell him: "if you try to do something you will see what will happent to you...". I think we are strong enugh now and the only evil spirit we have to feel is the one which is in ourselves (quoting somebody: I have seen the enemy and them are us). Now enough of the idle talk and back to Mathematics.'




References

[BBCGHMST24] Balister, P. and Bollob\'{a}s, B. and Campos, M. and Griffiths, S. and Hurley, E. and Morris, R. and Sahasrabudhe, J. and Tiba, M., Upper bounds for multicolour Ramsey numbers . arXiv:2410.17197 (2024).

[CGMS23] Campos, Marcelo and Griffiths, Simon and Morris, Robert and Sahasrabudhe, Julian, An exponential improvement for diagonal Ramsey . arXiv:2303.09521 (2023).

[Er88] Erd\H{o}s, P, Problems and results in combinatorial analysis and graph theory . Discrete Math. (1988), 81-92.

[Er93] Erd\H{o}s, Paul, Some of my favorite solved and unsolved problems in graph theory . Quaestiones Math. (1993), 333-350.

[GNNW24] Gupta, P. and Ndiaye, N. and Norin, S. and Wei, L., Optimizing the CGMS upper bound on Ramsey numbers . arXiv:2407.19026 (2024).

\noindent\rule{\linewidth}{0.4pt}


%% ============================================================
\section{Problem \#78}
\label{prob:78}

\noindent\textbf{Status:} OPEN \quad|\quad \textbf{Prize:} \$100 \quad|\quad \textbf{Updated:} 2025-08-31

\noindent\textbf{Tags:} graph theory, ramsey theory

\medskip

\noindent\textbf{Statement:}

Let $R(k)$ be the Ramsey number for $K_k$, the minimal $n$ such that every $2$-colouring of the edges of $K_n$ contains a monochromatic copy of $K_k$. Give a constructive proof that $R(k)&#62;C^k$ for some constant $C&#62;1$.




Erd\H{o}s gave a simple probabilistic proof that $R(k) \gg k2^{k/2}$. Equivalently, this question asks for an explicit construction of a graph on $n$ vertices which does not contain any clique or independent set of size $\geq c\log n$ for some constant $c&gt;0$. In \cite{Er69b} Erd\H{o}s asks for even a construction whose largest clique or independent set has size $o(n^{1/2})$, which is now known. Cohen \cite{Co15} (see the introduction for further history) constructed a graph on $n$ vertices which does not contain any clique or independent set of size\[\geq 2^{(\log\log n)^{C}}\]for some constant $C&gt;0$. Li \cite{Li23b} has recently improved this to\[\geq (\log n)^{C}\]for some constant $C&gt;0$. This problem is #4 in Ramsey Theory in the graphs problem collection.




References

[Co15] Gil Cohen, Two-Source Dispersers for Polylogarithmic Entropy and Improved Ramsey Graphs . Electronic Colloquium on Computational Complexity (2015).

[Er69b] Erd\H{o}s, P., Problems and results in chromatic graph theory . Proof Techniques in Graph Theory (Proc. Second Ann Arbor Graph Theory Conf., Ann Arbor, Mich., 1968) (1969), 27-35.

[Li23b] Li, X., Two Source Extractors for Asymptotically Optimal Entropy, and (Many) More . arXiv:2303.06802 (2023).

\noindent\rule{\linewidth}{0.4pt}


%% ============================================================
\section{Problem \#80}
\label{prob:80}

\noindent\textbf{Status:} OPEN \quad|\quad \textbf{Prize:} none \quad|\quad \textbf{Updated:} 2025-08-31

\noindent\textbf{Tags:} graph theory, ramsey theory

\medskip

\noindent\textbf{Statement:}

Let $c&#62;0$ and let $f_c(n)$ be the maximal $m$ such that every graph $G$ with $n$ vertices and at least $cn^2$ edges, where each edge is contained in at least one triangle, must contain a book of size $m$, that is, an edge shared by at least $m$ different triangles. Estimate $f_c(n)$. In particular, is it true that $f_c(n)&#62;n^{\epsilon}$ for some $\epsilon&#62;0$? Or $f_c(n)\gg \log n$?




A problem of Erd\H{o}s and Rothschild. Alon and Trotter showed that, provided $c&lt;1/4$, $f_c(n)\ll_c n^{1/2}$. Szemer\'{e}di observed that his regularity lemma implies that $f_c(n)\to \infty$. Edwards (unpublished) and Khadzhiivanov and Nikiforov \cite{KhNi79} proved independently that $f_c(n) \geq n/6$ when $c&gt;1/4$ (see [905] ). Fox and Loh \cite{FoLo12} proved that\[f_c(n) \leq n^{O(1/\log\log n)}\]for all $c&lt;1/4$, disproving the first conjecture of Erd\H{o}s. The best known lower bounds for $f_c(n)$ are those from Szemer\'{e}di's regularity lemma, and as such remain very poor. See also [600] and the entry in the graphs problem collection .




References

[FoLo12] Fox, Jacob and Loh, Po-Shen, On a problem of Erd\H{o}s and {R}othschild on edges in triangles . Combinatorica (2012), 619--628.

[KhNi79] Khadzhiivanov, N. G. and Nikiforov, S. V., Solution of a problem of {P}. Erd\H{o}s about the maximum number of triangles with a common edge in a graph . C. R. Acad. Bulgare Sci. (1979), 1315--1318.

\noindent\rule{\linewidth}{0.4pt}


%% ============================================================
\section{Problem \#81}
\label{prob:81}

\noindent\textbf{Status:} OPEN \quad|\quad \textbf{Prize:} none \quad|\quad \textbf{Updated:} 2025-08-31

\noindent\textbf{Tags:} graph theory

\medskip

\noindent\textbf{Statement:}

Let $G$ be a chordal graph on $n$ vertices - that is, $G$ has no induced cycles of length greater than $3$. Can the edges of $G$ be partitioned into $n^2/6+O(n)$ many cliques?




Asked by Erd\H{o}s, Ordman, and Zalcstein \cite{EOZ93}, who proved an upper bound of $(1/4-\epsilon)n^2$ many cliques (for some very small $\epsilon&gt;0$). The example of all edges between a complete graph on $n/3$ vertices and an empty graph on $2n/3$ vertices show that $n^2/6+O(n)$ is sometimes necessary. A split graph is one where the vertices can be split into a clique and an independent set. Every split graph is chordal. Chen, Erd\H{o}s, and Ordman \cite{CEO94} have shown that any split graph can be partitioned into $\frac{3}{16}n^2+O(n)$ many cliques. See also [1017] .




References

[CEO94] Chen, Guan-Tao and Erd\H{o}s, Paul and Ordman, Edward T., Clique partitions of split graphs . Combinatorics, graph theory, algorithms and applications (Beijing, 1993) (1994), 21-30.

[EOZ93] Erd\H{o}s, Paul and Ordman, Edward T. and Zalcstein, Yechezkel, Clique partitions of chordal graphs . Combin. Probab. Comput. (1993), 409-415.

\noindent\rule{\linewidth}{0.4pt}


%% ============================================================
\section{Problem \#82}
\label{prob:82}

\noindent\textbf{Status:} OPEN \quad|\quad \textbf{Prize:} none \quad|\quad \textbf{Updated:} 2025-08-31

\noindent\textbf{Tags:} graph theory

\medskip

\noindent\textbf{Statement:}

Let $F(n)$ be maximal such that every graph on $n$ vertices contains a regular induced subgraph on at least $F(n)$ vertices. Prove that $F(n)/\log n\to \infty$.




Conjectured by Erd\H{o}s, Fajtlowicz, and Stanton. It is known that $F(5)=3$ and $F(7)=4$. Ramsey's theorem implies that $F(n)\gg \log n$. Bollob\'{a}s observed that $F(n)\ll n^{1/2+o(1)}$. Alon, Krivelevich, and Sudakov \cite{AKS07} have improved this to $n^{1/2}(\log n)^{O(1)}$. Dyson and McKay \cite{DyMc26} have improved this further to $F(n) \ll n^{1/2}$. In \cite{Er93} Erd\H{o}s asks whether, if $t(n)$ is the largest trivial (either empty or complete) subgraph which a graph on $n$ vertices must contain (so that $t(n) \gg \log n$ by Ramsey's theorem), then is it true that\[F(n)-t(n)\to \infty?\]Equivalently, and in analogue with the definition of Ramsey numbers, one can define $G(n)$ to be the minimal $m$ such that every graph on $m$ vertices contains a regular induced subgraph on at least $n$ vertices. This problem can be rephrased as asking whether $G(n) \leq 2^{o(n)}$. Fajtlowicz, McColgan, Reid, and Staton \cite{FMRS95} showed that $G(1)=1$, $G(2)=2$, $G(3)=5$, $G(4)=7$, and $G(5)\geq 12$. Boris Alexeev and Brendan McKay (see the comments and this site ) have computed $G(5)=17$, $G(6)\geq 21$, and $G(7)\geq 29$. Dyson and McKay \cite{DyMc26} have proved $G(7)\geq 30$, and $G(k)\geq \frac{9}{163}k^2$ for all large $k$. See also [1031] for another question regarding induced regular subgraphs.




References

[AKS07] Alon, N. and Krivelevich, M. and Sudakov, B., Large nearly regular induced subgraphs . arXiv:0710.2106 (2007).

[DyMc26] P. Dyson and B. McKay, Ramsey numbers for regular induced subgraphs . arXiv:2604.08215 (2026).

[Er93] Erd\H{o}s, Paul, Some of my favorite solved and unsolved problems in graph theory . Quaestiones Math. (1993), 333-350.

[FMRS95] Fajtlowicz, Siemion and McColgan, Tamara and Reid, Talmage and Staton, William, Ramsey numbers for induced regular subgraphs . Ars Combin. (1995), 149--154.

\noindent\rule{\linewidth}{0.4pt}


%% ============================================================
\section{Problem \#84}
\label{prob:84}

\noindent\textbf{Status:} OPEN \quad|\quad \textbf{Prize:} none \quad|\quad \textbf{Updated:} 2025-08-31

\noindent\textbf{Tags:} graph theory, cycles

\medskip

\noindent\textbf{Statement:}

The cycle set of a graph $G$ on $n$ vertices is a set $A\subseteq \{3,\ldots,n\}$ such that there is a cycle in $G$ of length $\ell$ if and only if $\ell \in A$. Let $f(n)$ count the number of possible such $A$. Prove that $f(n)=o(2^n)$. Prove that $f(n)/2^{n/2}\to \infty$.




Conjectured by Erd\H{o}s and Faudree, who showed that $2^{n/2}&#60;f(n) \leq 2^{n-2}$. The first problem was solved by Verstra\"{e}te \cite{Ve04}, who proved\[f(n)\ll 2^{n-n^{1/10}}.\]This was improved by Nenadov \cite{Ne25} to\[f(n) \ll 2^{n-n^{1/2-o(1)}}.\]One can also ask about the existence and value of $\lim f(n)^{1/n}$.




References

[Ne25] R. Nenadov, Improved bound on the number of cycle sets . arXiv:2501.09904 (2025).

[Ve04] Verstra\"{e}te, Jacques, On the number of sets of cycle lengths . Combinatorica (2004), 719-730.

\noindent\rule{\linewidth}{0.4pt}


%% ============================================================
\section{Problem \#85}
\label{prob:85}

\noindent\textbf{Status:} OPEN \quad|\quad \textbf{Prize:} none \quad|\quad \textbf{Updated:} 2026-03-14

\noindent\textbf{Tags:} graph theory

\medskip

\noindent\textbf{Statement:}

Let $n\geq 4$ and $f(n)$ be minimal such that every graph on $n$ vertices with minimal degree $\geq f(n)$ contains a $C_4$. Is it true that, for all large $n$, $f(n+1)\geq f(n)$?




The function $f(n)$ is a reformulation of the Ramsey number $R(C_4,K_{1,n})$, in that\[R(C_4,K_{1,n})=\min\{ m : f(m)\leq m-n\}\]and\[f(n)=\min\{ m : m\geq R(C_4, K_{1,n-m})\}.\]The behaviour of this Ramsey number more generally is [552] . A weaker version of the conjecture asks for some constant $c$ such that $f(m)&gt;f(n)-c$ for all $m&gt;n$. This question can be asked for other graphs than $C_4$. The bounds in [552] imply in particular that $f(n)&lt;\sqrt{n}+1$ and\[f(n)=(1+o(1))\sqrt{n}.\]It is easy to check that $f(4)=2$.

\noindent\rule{\linewidth}{0.4pt}


%% ============================================================
\section{Problem \#86}
\label{prob:86}

\noindent\textbf{Status:} OPEN \quad|\quad \textbf{Prize:} \$100 \quad|\quad \textbf{Updated:} 2025-08-31

\noindent\textbf{Tags:} graph theory

\medskip

\noindent\textbf{Statement:}

Let $Q_n$ be the $n$-dimensional hypercube graph (so that $Q_n$ has $2^n$ vertices and $n2^{n-1}$ edges). Is it true that every subgraph of $Q_n$ with\[\geq \left(\frac{1}{2}+o(1)\right)n2^{n-1}\]many edges contains a $C_4$?




Let $f(n)$ be the maximum number of edges in a subgraph of $Q_n$ without a $C_4$, so that this conjecture is that $f(n)\leq (\frac{1}{2}+o(1))n2^{n-1}$. Erd\H{o}s \cite{Er91} showed that\[f(n) \geq \left(\frac{1}{2}+\frac{c}{n}\right)n2^{n-1}\]for some constant $c&gt;0$, and wrote it is 'perhaps not hopeless' to determine $f(n)$ exactly. Brass, Harborth, and Nienborg \cite{BHN95} improved this to\[f(n) \geq \left(\frac{1}{2}+\frac{c}{\sqrt{n}}\right)n2^{n-1}\]for some constant $c&gt;0$. Balogh, Hu, Lidicky, and Liu \cite{BHLL14} proved that $f(n)\leq 0.6068 n2^{n-1}$. This was improved to $\leq 0.60318 n2^{n-1}$ by Baber \cite{Ba12b}. A similar question can be asked for other even cycles. See also [666] and the entry in the graphs problem collection .




References

[BHLL14] Balogh, J\'{o}zsef and Hu, Ping and Lidick\'{y}, Bernard and Liu, Hong, Upper bounds on the size of 4- and 6-cycle-free subgraphs of the hypercube . European J. Combin. (2014), 75-85.

[BHN95] Brass, Peter and Harborth, Heiko and Nienborg, Hauke, On the maximum number of edges in a {$C_4$}-free subgraph of {$Q_n$} . J. Graph Theory (1995), 17--23.

[Ba12b] R. Baber, Tur\'{a}n densities of hypercubes . arXiv:1201.3587 (2012).

[Er91] Erd\H{o}s, P., Problems and results in combinatorial analysis and combinatorial number theory . Graph theory, combinatorics, and applications, Vol. 1 (Kalamazoo, MI, 1988) (1991), 397-406.

\noindent\rule{\linewidth}{0.4pt}


%% ============================================================
\section{Problem \#87}
\label{prob:87}

\noindent\textbf{Status:} OPEN \quad|\quad \textbf{Prize:} none \quad|\quad \textbf{Updated:} 2025-08-31

\noindent\textbf{Tags:} graph theory, ramsey theory

\medskip

\noindent\textbf{Statement:}

Let $\epsilon &#62;0$. Is it true that, if $k$ is sufficiently large, then\[R(G)&#62;(1-\epsilon)^kR(k)\]for every graph $G$ with chromatic number $\chi(G)=k$? Even stronger, is there some $c&#62;0$ such that, for all large $k$, $R(G)&#62;cR(k)$ for every graph $G$ with chromatic number $\chi(G)=k$?




Erd\H{o}s originally conjectured that $R(G)\geq R(k)$, which is trivial for $k=3$, but fails already for $k=4$, as Faudree and McKay \cite{FaMc93} showed that $R(W)=17$ for the pentagonal wheel $W$. Since $R(k)\leq 4^k$ this is trivial for $\epsilon\geq 3/4$. Yuval Wigderson points out that $R(G)\gg 2^{k/2}$ for any $G$ with chromatic number $k$ (via a random colouring), which asymptotically matches the best-known lower bounds for $R(k)$. This problem is #12 and #13 in Ramsey Theory in the graphs problem collection.




References

[FaMc93] Faudree, R. J. and McKay, B., A conjecture of Erd\H{o}s and the Ramsey number $r(W_6)$ . J. Combinatorial Math. and Combinatorial Computing (1993), 23-31.

\noindent\rule{\linewidth}{0.4pt}


%% ============================================================
\section{Problem \#89}
\label{prob:89}

\noindent\textbf{Status:} OPEN \quad|\quad \textbf{Prize:} \$500 \quad|\quad \textbf{Updated:} 2025-08-31

\noindent\textbf{Tags:} geometry, distances

\medskip
\noindent\textit{Erdős distance problem}


\noindent\textbf{Statement:}

Does every set of $n$ distinct points in $\mathbb{R}^2$ determine $\gg n/\sqrt{\log n}$ many distinct distances?




A $\sqrt{n}\times\sqrt{n}$ integer grid shows that this would be the best possible. Nearly solved by Guth and Katz \cite{GuKa15} who proved that there are always $\gg n/\log n$ many distinct distances. A stronger form (see [604] ) may be true: is there a single point which determines $\gg n/\sqrt{\log n}$ distinct distances, or even $\gg n$ many such points, or even that this is true averaged over all points - for example, if $d(x)$ counts the number of distinct distances from $x$ then in \cite{Er75f} Erd\H{o}s conjectured\[\sum_{x\in A}d(x) \gg \frac{n^2}{\sqrt{\log n}},\]where $A\subset \mathbb{R}^2$ is any set of $n$ points. See also [661] , and [1083] for the generalisation to higher dimensions.




References

[Er75f] Erd\H{o}s, Paul, On some problems of elementary and combinatorial geometry . Ann. Mat. Pura Appl. (4) (1975), 99-108.

[GuKa15] Guth, Larry and Katz, Nets Hawk, On the Erd\H{o}s distinct distances problem in the plane . Ann. of Math. (2) (2015), 155-190.

\noindent\rule{\linewidth}{0.4pt}


%% ============================================================
\section{Problem \#91}
\label{prob:91}

\noindent\textbf{Status:} OPEN \quad|\quad \textbf{Prize:} none \quad|\quad \textbf{Updated:} 2025-08-31

\noindent\textbf{Tags:} geometry, distances

\medskip

\noindent\textbf{Statement:}

Let $n$ be a sufficiently large integer. Suppose $A\subset \mathbb{R}^2$ has $\lvert A\rvert=n$ and minimises the number of distinct distances between points in $A$. Prove that there are at least two (and probably many) such $A$ which are non-similar.




For $n=3$ the equilateral triangle is the only such set. For $n=4$ the square or two equilateral triangles sharing an edge give two non-similar examples. For $n=5$ the regular pentagon is the unique such set (which has two distinct distances). Erd\H{o}s mysteriously remarks in \cite{Er90} this was proved by 'a colleague'. (In \cite{Er87b} this is described as 'a colleague from Zagreb (unfortunately I do not have his letter)'.) A published proof of this fact is provided by Kov\'{a}cs \cite{Ko24c}. In \cite{Er87b} Erd\H{o}s says that there are at least two non-similar examples for $6\leq n\leq 9$. The minimal possible number of distinct distances is the subject of [89] .




References

[Er87b] Erd\H{o}s, P., Some combinatorial and metric problems in geometry . Intuitive geometry (Si\'{o}fok, 1985) (1987), 167-177.

[Er90] Erd\H{o}s, Paul, Some of my favourite unsolved problems . A tribute to Paul Erd\H{o}s (1990), 467-478.

[Ko24c] Z. Kov\'{a}cs, A note on Erd\H{o}s's mysterious remark . arXiv:2412.05190 (2024).

\noindent\rule{\linewidth}{0.4pt}


%% ============================================================
\section{Problem \#96}
\label{prob:96}

\noindent\textbf{Status:} OPEN \quad|\quad \textbf{Prize:} none \quad|\quad \textbf{Updated:} 2025-08-31

\noindent\textbf{Tags:} geometry, distances, convex

\medskip

\noindent\textbf{Statement:}

If $n$ points in $\mathbb{R}^2$ form a convex polygon then there are $O(n)$ many pairs which are distance $1$ apart.




Conjectured by Erd\H{o}s and Moser. In \cite{Er92e} Erd\H{o}s credits the conjecture that the true upper bound is $2n$ to himself and Fishburn. F\"{u}redi \cite{Fu90} proved an upper bound of $O(n\log n)$. A short proof of this bound was given by Brass and Pach \cite{BrPa01}. The best known upper bound is\[\leq n\log_2n+4n,\]due to Aggarwal \cite{Ag15}. Edelsbrunner and Hajnal \cite{EdHa91} have constructed $n$ such points with $2n-7$ pairs distance $1$ apart. (This disproved an early stronger conjecture of Erd\H{o}s and Moser, that the true answer was $\frac{5}{3}n+O(1)$.) A positive answer would follow from [97] . See also [90] . In \cite{Er92e} Erd\H{o}s makes the stronger conjecture that, if $g(x)$ counts the largest number of points equidistant from $x$ in $A$, then\[\sum_{x\in A}g(x)&lt; 4n.\]He notes that the example of Edelsbrunner and Hajnal shows that $\sum_{x\in A}g(x)&gt;4n-O(1)$ is possible.




References

[Ag15] Aggarwal, Amol, On unit distances in a convex polygon . Discrete Math. (2015), 88-92.

[BrPa01] Brass , Peter and Pach, J\'{a}nos, The maximum number of times the same distance can occur among the vertices of a convex {$n$}-gon is {$O(n\log n)$} . J. Combin. Theory Ser. A (2001), 178-179.

[EdHa91] Edelsbrunner, Herbert and Hajnal, P\'{e}ter, A lower bound on the number of unit distances between the vertices of a convex polygon . J. Combin. Theory Ser. A (1991), 312-316.

[Er92e] Erd\H{o}s, P\'{a}l, Some Unsolved problems in Geometry, Number Theory and Combinatorics . Eureka (1992), 44-48.

[Fu90] F\"{u}redi, Zolt\'{a}n, The maximum number of unit distances in a convex {$n$}-gon . J. Combin. Theory Ser. A (1990), 316-320.

\noindent\rule{\linewidth}{0.4pt}


%% ============================================================
\section{Problem \#98}
\label{prob:98}

\noindent\textbf{Status:} OPEN \quad|\quad \textbf{Prize:} none \quad|\quad \textbf{Updated:} 2025-08-31

\noindent\textbf{Tags:} geometry, distances

\medskip

\noindent\textbf{Statement:}

Let $h(n)$ be such that any $n$ points in $\mathbb{R}^2$, with no three on a line and no four on a circle, determine at least $h(n)$ distinct distances. Does $h(n)/n\to \infty$?




Erd\H{o}s could not even prove $h(n)\geq n$. Pach has shown $h(n)&lt;n^{\log_23}$. Erd\H{o}s, F\"{u}redi, and Pach \cite{EFPR93} have improved this to\[h(n) &lt; n\exp(c\sqrt{\log n})\]for some constant $c&gt;0$.




References

[EFPR93] Erd\H{o}s, Paul and F\"{u}redi, Zolt\'{a}n and Pach, J\'{a}nos and Ruzsa, Imre Z., The grid revisited . Discrete Math. (1993), 189--196.

\noindent\rule{\linewidth}{0.4pt}


%% ============================================================
\section{Problem \#99}
\label{prob:99}

\noindent\textbf{Status:} OPEN \quad|\quad \textbf{Prize:} \$100 \quad|\quad \textbf{Updated:} 2025-08-31

\noindent\textbf{Tags:} geometry, distances

\medskip

\noindent\textbf{Statement:}

Let $A\subseteq\mathbb{R}^2$ be a set of $n$ points with minimum distance equal to 1, chosen to minimise the diameter of $A$. If $n$ is sufficiently large then must there be three points in $A$ which form an equilateral triangle of size 1?




Thue proved that the minimal such diameter is achieved (asymptotically) by the points in a triangular lattice intersected with a circle. In general Erd\H{o}s believed such a set must have very large intersection with the triangular lattice (perhaps as many as $(1-o(1))n$). Erd\H{o}s \cite{Er94b} wrote 'I could not prove it but felt that it should not be hard. To my great surprise both B. H. Sendov and M. Simonovits doubted the truth of this conjecture.' In \cite{Er94b} he offers \$100 for a counterexample but only \$50 for a proof. The stated problem is false for $n=4$, for example taking the points to be vertices of a square. The behaviour of such sets for small $n$ is explored by Bezdek and Fodor \cite{BeFo99}. See also [103] .




References

[BeFo99] Bezdek, Andr\'{a}s and Fodor, Ferenc, Minimal diameter of certain sets in the plane . J. Combin. Theory Ser. A (1999), 105-111.

[Er94b] Erd\H{o}s, Paul, Some problems in number theory, combinatorics and combinatorial geometry . Math. Pannon. (1994), 261-269.

\noindent\rule{\linewidth}{0.4pt}


%% ============================================================
\section{Problem \#100}
\label{prob:100}

\noindent\textbf{Status:} OPEN \quad|\quad \textbf{Prize:} none \quad|\quad \textbf{Updated:} 2025-08-31

\noindent\textbf{Tags:} geometry, distances

\medskip

\noindent\textbf{Statement:}

Let $A$ be a set of $n$ points in $\mathbb{R}^2$ such that all pairwise distances are at least $1$ and if two distinct distances differ then they differ by at least $1$. Is the diameter of $A$ $\gg n$?




Perhaps the diameter is even $\geq n-1$ for sufficiently large $n$. Piepmeyer has an example of $9$ such points with diameter $&lt;5$. Kanold proved the diameter is $\geq n^{3/4}$. The bounds on the distinct distance problem [89] proved by Guth and Katz \cite{GuKa15} imply a lower bound of $\gg n/\log n$.




References

[GuKa15] Guth, Larry and Katz, Nets Hawk, On the Erd\H{o}s distinct distances problem in the plane . Ann. of Math. (2) (2015), 155-190.

\noindent\rule{\linewidth}{0.4pt}


%% ============================================================
\section{Problem \#101}
\label{prob:101}

\noindent\textbf{Status:} OPEN \quad|\quad \textbf{Prize:} \$100 \quad|\quad \textbf{Updated:} 2025-08-31

\noindent\textbf{Tags:} geometry

\medskip

\noindent\textbf{Statement:}

Given $n$ points in $\mathbb{R}^2$, no five of which are on a line, the number of lines containing four points is $o(n^2)$.




There are examples of sets of $n$ points with $\sim n^2/6$ many collinear triples and no four points on a line. Such constructions are given by Burr, Gr\"{u}nbaum, and Sloane \cite{BGS74} and F\"{u}redi and Pal\'{a}sti \cite{FuPa84}. Gr\"{u}nbaum \cite{Gr76} constructed an example with $\gg n^{3/2}$ such lines. Erd\H{o}s speculated this may be the correct order of magnitude. This is false: Solymosi and Stojakovi\'{c} \cite{SoSt13} have constructed a set with no five on a line and at least\[n^{2-O(1/\sqrt{\log n})}\]many lines containing exactly four points. See also [102] and [669] . A generalisation of this problem is asked in [588] . This problem is Problem 71 on Green's open problems list .




References

[BGS74] Burr, Stefan A. and Gr\"{u}nbaum, Branko and Sloane, N. J. A., The orchard problem . Geometriae Dedicata (1974), 397-424.

[FuPa84] F\"{u}redi, Z. and Pal\'{a}sti, I., Arrangements of lines with a large number of triangles . Proc. Amer. Math. Soc. (1984), 561-566.

[Gr76] Gr\"{u}nbaum, Branko, New views on some old questions of combinatorial geometry . Colloquio Internazionale sulle Teorie Combinatorie (Roma, 1973), Tomo I (1976), 451-468.

[SoSt13] Solymosi, J\'{o}zsef and Stojakovi\'C, Milo\vS, Many collinear {$k$}-tuples with no {$k+1$} collinear points . Discrete Comput. Geom. (2013), 811-820.

\noindent\rule{\linewidth}{0.4pt}


%% ============================================================
\section{Problem \#102}
\label{prob:102}

\noindent\textbf{Status:} OPEN \quad|\quad \textbf{Prize:} none \quad|\quad \textbf{Updated:} 2025-08-31

\noindent\textbf{Tags:} geometry

\medskip

\noindent\textbf{Statement:}

Let $c&#62;0$ and $h_c(n)$ be such that for any $n$ points in $\mathbb{R}^2$ such that there are $\geq cn^2$ lines each containing more than three points, there must be some line containing $h_c(n)$ many points. Estimate $h_c(n)$. Is it true that, for fixed $c&#62;0$, we have $h_c(n)\to \infty$?




A problem of Erd\H{o}s and Purdy. It is not even known if $h_c(n)\geq 5$ (see [101] ). It is easy to see that $h_c(n) \ll_c n^{1/2}$, and Erd\H{o}s at one point \cite{Er95} suggested that perhaps a similar lower bound $h_c(n)\gg_c n^{1/2}$ holds. Zach Hunter has pointed out that this is false, even replacing $&gt;3$ points on each line with $&gt;k$ points: consider the set of points in $\{1,\ldots,m\}^d$ where $n\approx m^d$. These intersect any line in $\ll_d n^{1/d}$ points, and have $\gg_d n^2$ many pairs of points each of which determine a line with at least $k$ points. This is a construction in $\mathbb{R}^d$, but a random projection into $\mathbb{R}^2$ preserves the relevant properties. This construction shows that $h_c(n) \ll n^{1/\log(1/c)}$.




References

[Er95] Erd\H{o}s, Paul, Some of my favourite problems in number theory, combinatorics, and geometry . Resenhas (1995), 165-186.

\noindent\rule{\linewidth}{0.4pt}


%% ============================================================
\section{Problem \#103}
\label{prob:103}

\noindent\textbf{Status:} OPEN \quad|\quad \textbf{Prize:} none \quad|\quad \textbf{Updated:} 2025-08-31

\noindent\textbf{Tags:} geometry, distances

\medskip

\noindent\textbf{Statement:}

Let $h(n)$ count the number of incongruent sets of $n$ points in $\mathbb{R}^2$ which minimise the diameter subject to the constraint that $d(x,y)\geq 1$ for all points $x\neq y$. Is it true that $h(n)\to \infty$?




It is not even known whether $h(n)\geq 2$ for all large $n$. See also [99] .

\noindent\rule{\linewidth}{0.4pt}


%% ============================================================
\section{Problem \#104}
\label{prob:104}

\noindent\textbf{Status:} OPEN \quad|\quad \textbf{Prize:} \$100 \quad|\quad \textbf{Updated:} 2025-08-31

\noindent\textbf{Tags:} geometry

\medskip

\noindent\textbf{Statement:}

Given $n$ points in $\mathbb{R}^2$ the number of distinct unit circles containing at least three points is $o(n^2)$.




In \cite{Er81d} Erd\H{o}s proved that $\gg n$ many circles is possible, and that there cannot be more than $O(n^2)$ many circles. The argument is very simple: every pair of points determines at most $2$ unit circles, and the claimed bound follows from double counting. Erd\H{o}s claims in a number of places this produces the upper bound $n(n-1)$, but Harborth and Mengerson \cite{HaMe86} note that in fact this delivers an upper bound of $\frac{n(n-1)}{3}$. Elekes \cite{El84} has a simple construction of a set with $\gg n^{3/2}$ such circles. This may be the correct order of magnitude. In \cite{Er75h} and \cite{Er92e} Erd\H{o}s also asks how many such unit circles there must be if the points are in general position. In \cite{Er92e} Erd\H{o}s offered £100 for a proof or disproof that the answer is $O(n^{3/2})$. The maximal number of unit circles achieved by $n$ points is A003829 in the OEIS. See also [506] and [831] .




References

[El84] Elekes, G., {$n$} points in the plane can determine $n^{3/2}$ unit circles . Combinatorica (1984), 131.

[Er75h] Erd\H{o}s, P., Some problems on elementary geometry . Austral. Math. Soc. Gaz. (1975), 2-3.

[Er81d] Erd\H{o}s, P., Some applications of graph theory and combinatorial methods to number theory and geometry . Algebraic methods in graph theory, Vol. I, II (Szeged, 1978) (1981), 137-148.

[Er92e] Erd\H{o}s, P\'{a}l, Some Unsolved problems in Geometry, Number Theory and Combinatorics . Eureka (1992), 44-48.

[HaMe86] Harborth, Heiko and Mengersen, Ingrid, Point sets with many unit circles . Discrete Math. (1986), 193--197.

\noindent\rule{\linewidth}{0.4pt}


%% ============================================================
\section{Problem \#108}
\label{prob:108}

\noindent\textbf{Status:} OPEN \quad|\quad \textbf{Prize:} none \quad|\quad \textbf{Updated:} 2025-08-31

\noindent\textbf{Tags:} graph theory, chromatic number, cycles

\medskip

\noindent\textbf{Statement:}

For every $r\geq 4$ and $k\geq 2$ is there some finite $f(k,r)$ such that every graph of chromatic number $\geq f(k,r)$ contains a subgraph of girth $\geq r$ and chromatic number $\geq k$?




Conjectured by Erd\H{o}s and Hajnal. R\"{o}dl \cite{Ro77} has proved the $r=4$ case (see [923] ). The infinite version (whether every graph of infinite chromatic number contains a subgraph of infinite chromatic number whose girth is $&gt;k$) is also open. In \cite{Er79b} Erd\H{o}s also asks whether\[\lim_{k\to \infty}\frac{f(k,r+1)}{f(k,r)}=\infty.\]See also the entry in the graphs problem collection and [740] for the infinitary version.




References

[Er79b] Erd\H{o}s, Paul, Problems and results in graph theory and combinatorial analysis . Graph theory and related topics (Proc. Conf., Univ. Waterloo, Waterloo, Ont., 1977) (1979), 153-163.

[Ro77] R\"{o}dl, V., On the chromatic number of subgraphs of a given graph . Proc. Amer. Math. Soc. (1977), 370-371.

\noindent\rule{\linewidth}{0.4pt}


%% ============================================================
\section{Problem \#111}
\label{prob:111}

\noindent\textbf{Status:} OPEN \quad|\quad \textbf{Prize:} none \quad|\quad \textbf{Updated:} 2025-08-31

\noindent\textbf{Tags:} graph theory, chromatic number, set theory

\medskip

\noindent\textbf{Statement:}

If $G$ is a graph let $h_G(n)$ be defined such that any subgraph of $G$ on $n$ vertices can be made bipartite after deleting at most $h_G(n)$ edges. What is the behaviour of $h_G(n)$? Is it true that $h_G(n)/n\to \infty$ for every graph $G$ with chromatic number $\aleph_1$?




A problem of Erd\H{o}s, Hajnal, and Szemer\'{e}di \cite{EHS82}. Every $G$ with chromatic number $\aleph_1$ must have $h_G(n)\gg n$ since $G$ must contain, for some $r$, $\aleph_1$ many vertex disjoint odd cycles of length $2r+1$. On the other hand, Erd\H{o}s, Hajnal, and Szemer\'{e}di proved that there is a $G$ with chromatic number $\aleph_1$ such that $h_G(n)\ll n^{3/2}$. In \cite{Er81} Erd\H{o}s conjectured that this can be improved to $\ll n^{1+\epsilon}$ for every $\epsilon&gt;0$. See also [74] .




References

[EHS82] Erd\H{o}s, P. and Hajnal, A. and Szemer\'{e}di, E., On almost bipartite large chromatic graphs . Theory and practice of combinatorics (1982), 117-123.

[Er81] Erd\H{o}s, P., On the combinatorial problems which I would most like to see solved . Combinatorica (1981), 25-42.

\noindent\rule{\linewidth}{0.4pt}


%% ============================================================
\section{Problem \#112}
\label{prob:112}

\noindent\textbf{Status:} OPEN \quad|\quad \textbf{Prize:} none \quad|\quad \textbf{Updated:} 2025-08-31

\noindent\textbf{Tags:} graph theory, ramsey theory

\medskip

\noindent\textbf{Statement:}

Let $k=k(n,m)$ be minimal such that any directed graph on $k$ vertices must contain either an independent set of size $n$ or a transitive tournament of size $m$. Determine $k(n,m)$.




A problem of Erd\H{o}s and Rado \cite{ErRa67}, who showed $k(n,m) \ll_m n^{m-1}$, or more precisely,\[k(n,m) \leq \frac{2^{m-1}(n-1)^m+n-2}{2n-3}.\]Larson and Mitchell \cite{LaMi97} improved the dependence on $m$, establishing in particular that $k(n,3)\leq n^{2}$. Zach Hunter has observed that\[R(n,m) \leq k(n,m)\leq R(n,m,m),\]which in particular proves the upper bound $k(n,m)\leq 3^{n+2m}$. See also the entry in the graphs problem collection - on this site the problem replaces transitive tournament with directed path, but Zach Hunter and Raphael Steiner have a simple argument that proves, for this alternative definition, that $k(n,m)=(n-1)(m-1)$.




References

[ErRa67] Erd\H{o}s, P. and Rado, R., Partition relations and transitivity domains of binary relations . J. London Math. Soc. (1967), 624-633.

[LaMi97] Larson, Jean A. and Mitchell, William J., On a problem of Erd\H{o}s and Rado . Ann. Comb. (1997), 245-252.

\noindent\rule{\linewidth}{0.4pt}


%% ============================================================
\section{Problem \#117}
\label{prob:117}

\noindent\textbf{Status:} OPEN \quad|\quad \textbf{Prize:} none \quad|\quad \textbf{Updated:} 2025-08-31

\noindent\textbf{Tags:} group theory

\medskip

\noindent\textbf{Statement:}

Let $h(n)$ be minimal such that any group $G$ with the property that any subset of $&#62;n$ elements contains some $x\neq y$ such that $xy=yx$ can be covered by at most $h(n)$ many Abelian subgroups. Estimate $h(n)$ as well as possible.




Pyber \cite{Py87} has proved there exist constants $c_2&#62;c_1&#62;1$ such that $c_1^n&#60;h(n)&#60;c_2^n$. Erd\H{o}s \cite{Er97f} writes that the lower bound was already known to Isaacs.




References

[Er97f] Erd\H{o}s, Paul, Some unsolved problems . Combinatorics, geometry and probability (Cambridge, 1993) (1997), 1-10.

[Py87] Pyber, L., The number of pairwise noncommuting elements and the index of the centre in a finite group . J. London Math. Soc. (2) (1987), 287-295.

\noindent\rule{\linewidth}{0.4pt}


%% ============================================================
\section{Problem \#119}
\label{prob:119}

\noindent\textbf{Status:} OPEN \quad|\quad \textbf{Prize:} \$100 \quad|\quad \textbf{Updated:} 2025-08-31

\noindent\textbf{Tags:} analysis, polynomials

\medskip

\noindent\textbf{Statement:}

Let $z_i$ be an infinite sequence of complex numbers such that $\lvert z_i\rvert=1$ for all $i\geq 1$, and for $n\geq 1$ let\[p_n(z)=\prod_{i\leq n} (z-z_i).\]Let $M_n=\max_{\lvert z\rvert=1}\lvert p_n(z)\rvert$. Is it true that $\limsup M_n=\infty$? Is it true that there exists $c&#62;0$ such that for infinitely many $n$ we have $M_n &#62; n^c$? Is it true that there exists $c&#62;0$ such that, for all large $n$,\[\sum_{k\leq n}M_k &#62; n^{1+c}?\]




This is Problem 4.1 in \cite{Ha74} where it is attributed to Erd\H{o}s. The weaker conjecture that $\limsup M_n=\infty$ was proved by Wagner \cite{Wa80}, who show that there is some $c&#62;0$ with $M_n&#62;(\log n)^c$ infinitely often. The second question was answered by Beck \cite{Be91}, who proved that there exists some $c&#62;0$ such that\[\max_{n\leq N} M_n &#62; N^c.\]Erd\H{o}s (e.g. see \cite{Ha74}) gave a construction of a sequence with $M_n\leq n+1$ for all $n$. Linden \cite{Li77} improved this to give a sequence with $M_n\ll n^{1-c}$ for some $c&#62;0$. The third question seems to remain open.




References

[Be91] Beck, J., The modulus of polynomials with zeros on the unit circle: A problem of Erd\H{o}s . Annals of Math. (1991), 609-651.

[Ha74] Hayman, W. K., Research problems in function theory: new problems . (1974), 155--180.

[Li77] Linden, C. N., The modulus of polynomials with zeros on the unit circle . Bull. London Math. Soc. (1977), 65--69.

[Wa80] Wagner, Gerold, On a problem of {E}rd\H{o}s in {D}iophantine approximation . Bull. London Math. Soc. (1980), 81--88.

\noindent\rule{\linewidth}{0.4pt}


%% ============================================================
\section{Problem \#120}
\label{prob:120}

\noindent\textbf{Status:} OPEN \quad|\quad \textbf{Prize:} \$100 \quad|\quad \textbf{Updated:} 2025-08-31

\noindent\textbf{Tags:} combinatorics

\medskip
\noindent\textit{Erdős similarity problem}


\noindent\textbf{Statement:}

Let $A\subseteq\mathbb{R}$ be an infinite set. Must there be a set $E\subset \mathbb{R}$ of positive measure which does not contain any set of the shape $aA+b$ for some $a,b\in\mathbb{R}$ and $a\neq 0$?




The Erd\H{o}s similarity problem. This is true if $A$ is unbounded or dense in some interval. It therefore suffices to prove this when $A=\{a_1&gt;a_2&gt;\cdots\}$ is a countable strictly monotone sequence which converges to $0$. Steinhaus \cite{St20} has proved this is false whenever $A$ is a finite set. This conjecture is known in many special cases (but, for example, it is open when $A=\{1,1/2,1/4,\ldots\}$, which is Problem 94 on Green's open problems list ). For an overview of progress we recommend a nice survey by Svetic \cite{Sv00} on this problem. A survey of more recent progress was written by Jung, Lai, and Mooroogen \cite{JLM24}.




References

[JLM24] Y. Jung and C.-K. Lai and Y. Mooroogen, Some recent progress on the Erd\H{o}s similarity conjecture . arXiv:2412.11062 (2024).

[St20] Steinhaus, Hugo, Sur les distances des points dans les ensembles de measure positive . Fund. Math. (1920), 93-104.

[Sv00] Svetic, R. E., The Erd\H{o}s similarity problem: a survey . Real Anal. Exchange (2000/01), 525-539.

\noindent\rule{\linewidth}{0.4pt}


%% ============================================================
\section{Problem \#122}
\label{prob:122}

\noindent\textbf{Status:} OPEN \quad|\quad \textbf{Prize:} none \quad|\quad \textbf{Updated:} 2025-08-31

\noindent\textbf{Tags:} number theory

\medskip

\noindent\textbf{Statement:}

For which number theoretic functions $f$ is it true that, for any $F(n)$ such that $F(n)/f(n)\to 0$ for almost all $n$, there are infinitely many $x$ such that\[\frac{\#\{ n\in \mathbb{N} : n+f(n)\in (x,x+F(x))\}}{F(x)}\to \infty?\]




In \cite{Er97} Erd\H{o}s is only considering number theoretic functions which grow 'slowly' (i.e. slower than $(\log n)^{1-c}$ for some $c&#62;0$). This is considered by Erd\H{o}s, Pomerance, and S\'{a}rk\"{o}zy \cite{EPS97}, who prove in particular that if $\omega(n)$ counts the number of distinct prime divisors of $n$, then for all large $x$ there are intervals $I,J\subset [1,x]$ of width\[\lvert I\rvert\asymp \left(\frac{\log x}{\log\log x}\right)^{1/2}\]and\[\lvert J\rvert \asymp (\log\log x)^{1/2}\]such that if $n\in I$ then $n+\omega(n)\in J$. (Note the normal order of $\omega$ is $\log\log x$, and hence $(\log\log n)^{1/2}/\omega(n)\to 0$ for almost all $n$.) In \cite{Er97} and \cite{Er97e} Erd\H{o}s reports that he, Pomerance, and S\'{a}rkz\"{o}zy can prove the more general claim above for $f$ being $\tau(n)$, the divisor function, or $\omega(n)$, and states it 'probably fails' for $\phi(n)$ or $\sigma(n)$.




References

[EPS97] Erd\H{o}s, Paul and Pomerance, Carl and S\'{a}rk\"{o}zy, Andr\'{a}s, On locally repeated values of certain arithmetic functions. IV . Ramanujan J. (1997), 227-241.

[Er97] Erd\H{o}s, Paul, Problems in number theory . New Zealand J. Math. (1997), 155-160.

[Er97e] Erd\H{o}s, Paul, Some of my favourite unsolved problems . Math. Japon. (1997), 527-537.

\noindent\rule{\linewidth}{0.4pt}


%% ============================================================
\section{Problem \#123}
\label{prob:123}

\noindent\textbf{Status:} OPEN \quad|\quad \textbf{Prize:} \$250 \quad|\quad \textbf{Updated:} 2025-08-31

\noindent\textbf{Tags:} number theory

\medskip

\noindent\textbf{Statement:}

Let $a,b,c\geq 1$ be three integers which are pairwise coprime. Is every large integer the sum of distinct integers of the form $a^kb^lc^m$ ($k,l,m\geq 0$), none of which divide any other?




A sequence is said to be $d$-complete if every large integer is the sum of distinct integers from the sequence, none of which divide any other. This particular case of $d$-completeness was conjectured by Erd\H{o}s and Lewin \cite{ErLe96}, who (among other related results) prove this when $a=3$, $b=5$, and $c=7$. As a partial record of progress so far, the sequence $\{a^kb^lc^m\}$ is known to be $d$-complete when: {UL} {LI}$a=3$, $b=5$, $c=7$ (Erd\H{o}s and Lewin \cite{ErLe96}).{/LI} {LI}$a=2$, $b=5$, $c\in \{7,11,13,17,19\}$ (Erd\H{o}s and Lewin \cite{ErLe96}).{/LI} {LI}$a=2$, $b=5$, $c\in \{9,21,23,27,29,31\}$ - more generally, $a=2$, $b=5$, and any $c&gt;6$ with $(c,10)=1$ such that there exists $N$ where every integer in $(N,25cN)$ is the sum of distinct elements of $\{2^k3^lc^m\}$, none of which divide any other (Ma and Chen \cite{MaCh16}).{/LI} {LI} $a=2$, $b=5$, $3\leq c\leq 87$ with $(c,10)=1$, or $a=2$, $b=7$, $3\leq c\leq 33$ with $(c,14)=1$, or $a=3$, $b=5$, $2\leq c\leq 14$ with $(c,15)=1$ (Chen and Yu \cite{ChYu23b}).{/LI} {/UL} In \cite{Er92b} Erd\H{o}s makes the stronger conjecture (for $a=2$, $b=3$, and $c=5$) that, for any $\epsilon&gt;0$, all large integers $n$ can be written as the sum of distinct integers $b_1&lt;\cdots &lt;b_t$ of the form $2^k3^l5^m$ where $b_t&lt;(1+\epsilon)b_1$. See also [845] , and [1110] for the case of two powers.




References

[ChYu23b] Chen, Yong-Gao and Yu, Wang-Xing, On {$d$}-complete sequences of integers, {II} . Acta Arith. (2023), 161--181.

[Er92b] Erd\H{o}s, Paul, Some of my favourite problems in various branches of combinatorics . Matematiche (Catania) (1992), 231-240.

[ErLe96] Erd\H{o}s, P. and Lewin, Mordechai, $d$-complete sequences of integers . Math. Comp. (1996), 837-840.

[MaCh16] Ma, Mi-Mi and Chen, Yong-Gao, On {$d$}-complete sequences of integers . J. Number Theory (2016), 1--12.

\noindent\rule{\linewidth}{0.4pt}


%% ============================================================
\section{Problem \#124}
\label{prob:124}

\noindent\textbf{Status:} OPEN \quad|\quad \textbf{Prize:} none \quad|\quad \textbf{Updated:} 2025-08-31

\noindent\textbf{Tags:} number theory, base representations

\medskip

\noindent\textbf{Statement:}

For any $d\geq 1$ and $k\geq 0$ let $P(d,k)$ be the set of integers which are the sum of distinct powers $d^i$ with $i\geq k$. Let $3\leq d_1&#60;d_2&#60;\cdots &#60;d_r$ be integers such that\[\sum_{1\leq i\leq r}\frac{1}{d_r-1}\geq 1.\]Can all sufficiently large integers be written as a sum of the shape $\sum_i c_ia_i$ where $c_i\in \{0,1\}$ and $a_i\in P(d_i,0)$? If we further have $\mathrm{gcd}(d_1,\ldots,d_r)=1$ then, for any $k\geq 1$, can all sufficiently large integers be written as a sum of the shape $\sum_i c_ia_i$ where $c_i\in \{0,1\}$ and $a_i\in P(d_i,k)$?




The second question was conjectured by Burr, Erd\H{o}s, Graham, and Li \cite{BEGL96}, who proved it for $\{3,4,7\}$. The first question was asked separately by Erd\H{o}s in \cite{Er97} and \cite{Er97e} (although there is some ambiguity over whether he intended $P(d,0)$ or $P(d,1)$ - certainly he mentions no gcd condition). A simple positive proof of the first question was provided (and formalised in Lean) by Aristotle thanks to Alexeev; see the comments for details. In \cite{BEGL96} they record that Pomerance observed that the condition $\sum 1/(d_i-1)\geq 1$ is necessary (for both questions), but give no details. Tao has sketched an explanation in the comments. It is trivial that $\mathrm{gcd}(d_1,\ldots,d_r)=1$ is a necessary condition in the second question. Melfi \cite{Me04} gives a construction, for any $\epsilon&gt;0$, of an infinite set of $d_i$ for which every sufficiently large integer can be written as a finite sum of the shape $\sum_i c_ia_i$ where $c_i\in \{0,1\}$ and $a_i\in P(d_i,0)$ and yet $\sum_{i}\frac{1}{d_i-1}&lt;\epsilon$. See also [125] .




References

[BEGL96] Burr, S. A. and Erd\H{o}s, P. and Graham, R. L. and Li, W. Wen-Ching, Complete sequences of sets of integer powers . Acta Arith. (1996), 133-138.

[Er97] Erd\H{o}s, Paul, Problems in number theory . New Zealand J. Math. (1997), 155-160.

[Er97e] Erd\H{o}s, Paul, Some of my favourite unsolved problems . Math. Japon. (1997), 527-537.

[Me04] Melfi, Giuseppe, On certain positive integer sequences . Riv. Mat. Univ. Parma (7) (2004), 253--260.

\noindent\rule{\linewidth}{0.4pt}


%% ============================================================
\section{Problem \#126}
\label{prob:126}

\noindent\textbf{Status:} OPEN \quad|\quad \textbf{Prize:} \$250 \quad|\quad \textbf{Updated:} 2025-08-31

\noindent\textbf{Tags:} number theory

\medskip

\noindent\textbf{Statement:}

Let $f(n)$ be maximal such that if $A\subseteq\mathbb{N}$ has $\lvert A\rvert=n$ then $\prod_{a\neq b\in A}(a+b)$ has at least $f(n)$ distinct prime factors. Is it true that $f(n)/\log n\to\infty$?




Investigated by Erd\H{o}s and Tur\'{a}n \cite{ErTu34} (prompted by a question of L\'{a}z\'{a}r and Gr\"{u}nwald) in their first joint paper, where they proved that\[\log n \ll f(n) \ll n/\log n\](the upper bound is trivial, taking $A=\{1,\ldots,n\}$). Erd\H{o}s says that $f(n)=o(n/\log n)$ has never been proved, but perhaps never seriously attacked. This problem has been formalised in Lean as part of the Google DeepMind Formal Conjectures project .




References

[ErTu34] Erd\H{o}s, Paul and Turan, Paul, On a Problem in the Elementary Theory of Numbers . Amer. Math. Monthly (1934), 608-611.

\noindent\rule{\linewidth}{0.4pt}


%% ============================================================
\section{Problem \#129}
\label{prob:129}

\noindent\textbf{Status:} OPEN \quad|\quad \textbf{Prize:} none \quad|\quad \textbf{Updated:} 2025-08-31

\noindent\textbf{Tags:} graph theory, ramsey theory

\medskip
\noindent\textit{ambiguous statement}


\noindent\textbf{Statement:}

Let $R(n;k,r)$ be the smallest $N$ such that if the edges of $K_N$ are $r$-coloured then there is a set of $n$ vertices which does not contain a copy of $K_k$ in at least one of the $r$ colours. Prove that there is a constant $C=C(r)&#62;1$ such that\[R(n;3,r) &#60; C^{\sqrt{n}}.\]




Conjectured by Erd\H{o}s and Gy\'{a}rf\'{a}s, who proved the existence of some $C&#62;1$ such that $R(n;3,r)&#62;C^{\sqrt{n}}$. Note that when $r=k=2$ we recover the classic Ramsey numbers. Erd\H{o}s thought it likely that for all $r,k\geq 2$ there exists some $C_1,C_2&#62;1$ (depending only on $r$) such that\[ C_1^{n^{1/k-1}}&lt; R(n;k,r) &lt; C_2^{n^{1/k-1}}.\]Antonio Girao has pointed out that this problem as written is easily disproved, and indeed $R(n;3,2) \geq C^{n}$: The obvious probabilistic construction (randomly colour the edges red/blue independently uniformly at random) yields a 2-colouring of the edges of $K_N$ such every set on $n$ vertices contains a red triangle and a blue triangle (using that every set of $n$ vertices contains $\gg n^2$ edge-disjoint triangles), provided $N \leq C^n$ for some absolute constant $C&gt;1$. This implies $R(n;3,2) \geq C^{n}$, contradicting the conjecture. Perhaps Erd\H{o}s had a different problem in mind, but it is not clear what that might be. It would presumably be one where the natural probabilistic argument would deliver a bound like $C^{\sqrt{n}}$ as Erd\H{o}s and Gy\'{a}rf\'{a}s claim to have achieved via the probabilistic method.

\noindent\rule{\linewidth}{0.4pt}


%% ============================================================
\section{Problem \#130}
\label{prob:130}

\noindent\textbf{Status:} OPEN \quad|\quad \textbf{Prize:} none \quad|\quad \textbf{Updated:} 2025-08-31

\noindent\textbf{Tags:} graph theory, chromatic number

\medskip

\noindent\textbf{Statement:}

Let $A\subset\mathbb{R}^2$ be an infinite set which contains no three points on a line and no four points on a circle. Consider the graph with vertices the points in $A$, where two vertices are joined by an edge if and only if they are an integer distance apart. How large can the chromatic number and clique number of this graph be? In particular, can the chromatic number be infinite?




Asked by Andr\'{a}sfai and Erd\H{o}s. Erd\H{o}s \cite{Er97b} also asked where such a graph could contain an infinite complete graph, but this is impossible by an earlier result of Anning and Erd\H{o}s \cite{AnEr45}. See also [213] .




References

[AnEr45] Anning, Norman H. and Erd\H{o}s, Paul, Integral distances . Bull. Amer. Math. Soc. (1945), 598-600.

[Er97b] Erd\H{o}s, Paul, Some old and new problems in various branches of combinatorics . Discrete Math. (1997), 227-231.

\noindent\rule{\linewidth}{0.4pt}


%% ============================================================
\section{Problem \#131}
\label{prob:131}

\noindent\textbf{Status:} OPEN \quad|\quad \textbf{Prize:} none \quad|\quad \textbf{Updated:} 2025-08-31

\noindent\textbf{Tags:} number theory

\medskip

\noindent\textbf{Statement:}

Let $F(N)$ be the maximal size of $A\subseteq\{1,\ldots,N\}$ such that no $a\in A$ divides the sum of any distinct elements of $A\backslash\{a\}$. Estimate $F(N)$. In particular, is it true that\[F(N) &#62; N^{1/2-o(1)}?\]




This was studied by Erd\H{o}s, Lev, Rauzy, S\'{a}ndor, and S\'{a}rk\"{o}zy \cite{ELRSS99}, where they call such a property 'non-dividing', and prove the explicit bound\[F(N)&lt;3N^{1/2}+1.\]In \cite{Er97b} Erd\H{o}s credits Csaba with a construction that proves $F(N) \gg N^{1/5}$. Such a construction was also given in \cite{ELRSS99}, where it is linked to the problem of non-averaging sets (see [186] ). Indeed, every such set is non-averaging, and hence the result of Pham and Zakharov \cite{PhZa24} implies\[F(N) \leq N^{1/4+o(1)}.\]This shows the answer to the original question is no, but the general question of the correct growth of $F(N)$ remains open. In \cite{Er75b} Erd\H{o}s writes that he originally thought $F(N) &lt;(\log N)^{O(1)}$, but that Straus proved that\[F(N) &gt; \exp((\sqrt{\tfrac{2}{\log 2}}+o(1))\sqrt{\log N}).\]See also [13] . This is discussed in problem C16 of Guy's collection \cite{Gu04}.




References

[ELRSS99] Erd\H{o}s, P. and Lev, V. and Rauzy, G. and S\'{a}ndor, C. and S\'{a}rk\"ozy, A., Greedy algorithm, arithmetic progressions, subset sums and divisibility . Discrete Math. (1999), 119--135.

[Er75b] Erd\H{o}s, Paul, Problems and results in combinatorial number theory . Journ\'{e}es Arithm\'{e}tiques de Bordeaux (Conf., Univ. Bordeaux, Bordeaux, 1974) (1975), 295-310.

[Er97b] Erd\H{o}s, Paul, Some old and new problems in various branches of combinatorics . Discrete Math. (1997), 227-231.

[Gu04] Guy, Richard K., Unsolved problems in number theory . (2004), xviii+437.

[PhZa24] Pham, H. T. and Zakharov, D., Sharp bound for the Erd\H{o}s-Straus non-averaging set problem . arXiv:2410.14624 (2024).

\noindent\rule{\linewidth}{0.4pt}


%% ============================================================
\section{Problem \#132}
\label{prob:132}

\noindent\textbf{Status:} OPEN \quad|\quad \textbf{Prize:} \$100 \quad|\quad \textbf{Updated:} 2025-08-31

\noindent\textbf{Tags:} distances

\medskip

\noindent\textbf{Statement:}

Let $A\subset \mathbb{R}^2$ be a set of $n$ points. Must there be two distances which occur at least once but between at most $n$ pairs of points? Must the number of such distances $\to \infty$ as $n\to \infty$?




Asked by Erd\H{o}s and Pach. Hopf and Pannowitz \cite{HoPa34} proved that the largest distance between points of $A$ can occur at most $n$ times, but it is unknown whether a second such distance must occur. It may be true that there are at least $n^{1-o(1)}$ many such distances. In \cite{Er97e} Erd\H{o}s offers \$100 for 'any nontrivial result'. Erd\H{o}s \cite{Er84c} believed that for $n\geq 5$ there must always exist at least two such distances. This is false for $n=4$, as witnessed by two equilateral triangles of the same side-length glued together. Erd\H{o}s and Fishburn \cite{ErFi95} proved this is true for $n=5$ and $n=6$. Clemen, Dumitrescu, and Liu \cite{CDL25} have proved that there always at least two such distances if $A$ is in convex position (that is, no point lies inside the convex hull of the others). They also prove it is true if the set $A$ is 'not too convex', in a specific technical sense. See also [223] , [756] , and [957] .




References

[CDL25] F. Clemen, A. Dumitrescu, and D. Liu, On multiplicities of interpoint distances . arXiv:2505.04283 (2025).

[Er84c] Erd\H{o}s, Paul, Some old and new problems in combinatorial geometry . Convexity and graph theory (Jerusalem, 1981) (1984), 129-136.

[Er97e] Erd\H{o}s, Paul, Some of my favourite unsolved problems . Math. Japon. (1997), 527-537.

[ErFi95] Erd\H{o}s, Paul and Fishburn, Peter C., Multiplicities of interpoint distances in finite planar sets . Discrete Appl. Math. (1995), 141--147.

[HoPa34] Hopf, H. and Pannwitz, E., Aufgabe 167 . Jber. Deutsch. Math. Verein. (1934), 114.

\noindent\rule{\linewidth}{0.4pt}


%% ============================================================
\section{Problem \#137}
\label{prob:137}

\noindent\textbf{Status:} OPEN \quad|\quad \textbf{Prize:} none \quad|\quad \textbf{Updated:} 2025-08-31

\noindent\textbf{Tags:} number theory

\medskip

\noindent\textbf{Statement:}

We say that $N$ is powerful if whenever $p\mid N$ we also have $p^2\mid N$. Let $k\geq 3$. Can the product of any $k$ consecutive positive integers ever be powerful?




Conjectured by Erd\H{o}s and Selfridge. There are infinitely many $n$ such that $n(n+1)$ is powerful (see [364] ). Erd\H{o}s and Selfridge \cite{ErSe75} proved that the product of $k\geq 3$ consecutive positive integers can never be a perfect power. Erd\H{o}s remarked that this 'seems hopeless at present'. In \cite{Er82c} he further conjectures that, if $k$ is fixed and $n$ is sufficiently large, then, for all $m$, there must be at least $k$ distinct primes $p$ such that\[p\mid m(m+1)\cdots (m+n)\]and yet $p^2$ does not divide the right-hand side. See also [364] .




References

[Er82c] Erd\H{o}s, P., Miscellaneous problems in number theory . Congr. Numer. (1982), 25-45.

[ErSe75] Erd\H{o}s, P. and Selfridge, J. L., The product of consecutive integers is never a power . Illinois J. Math. (1975), 292-301.

\noindent\rule{\linewidth}{0.4pt}


%% ============================================================
\section{Problem \#138}
\label{prob:138}

\noindent\textbf{Status:} OPEN \quad|\quad \textbf{Prize:} \$500 \quad|\quad \textbf{Updated:} 2025-08-31

\noindent\textbf{Tags:} additive combinatorics

\medskip

\noindent\textbf{Statement:}

Let the van der Waerden number $W(k)$ be such that whenever $N\geq W(k)$ and $\{1,\ldots,N\}$ is $2$-coloured there must exist a monochromatic $k$-term arithmetic progression. Improve the bounds for $W(k)$ - for example, prove that $W(k)^{1/k}\to \infty$.




When $p$ is prime Berlekamp \cite{Be68} has proved $W(p+1)\geq p2^p$. Gowers \cite{Go01} has proved\[W(k) \leq 2^{2^{2^{2^{2^{k+9}}}}}.\]The best general lower bound is $W(k)\gg 2^k$, due to Kozik and Shabanov \cite{KoSh16}. In \cite{Er81} Erd\H{o}s further asks whether $W(k+1)/W(k)\to \infty$, or $W(k+1)-W(k)\to \infty$. DeepMind has proved (see the comments) that\[W(k+1)\geq W(k)+k,\]answering the second question. In \cite{Er80} Erd\H{o}s asks whether $W(k)/2^k\to \infty$, and offers \$500 for a proof or disproof of $W(k)^{1/k}\to \infty$. The analogous question for $r\geq 3$ colours was resolved by Fox and Hunter \cite{FoHu26}, who proved\[W_3(k)^{1/k}\geq C^{\log_*k}\]for some constant $C&#62;1$.




References

[Be68] Berlekamp, E. R., A construction for partitions which avoid long arithmetic progressions . Canad. Math. Bull. (1968), 409-414.

[Er80] Erd\H{o}s, Paul, A survey of problems in combinatorial number theory . Ann. Discrete Math. (1980), 89-115.

[Er81] Erd\H{o}s, P., On the combinatorial problems which I would most like to see solved . Combinatorica (1981), 25-42.

[FoHu26] J. Fox and Z. Hunter, Three-color van der Waerden numbers grow super-exponentially . arXiv:2606.02541 (2026).

[Go01] Gowers, W. T., A new proof of Szemer\'{e}di's theorem . Geom. Funct. Anal. (2001), 465-588.

[KoSh16] Kozik, Jakub and Shabanov, Dmitry, Improved algorithms for colorings of simple hypergraphs and applications . J. Combin. Theory Ser. B (2016), 312--332.

\noindent\rule{\linewidth}{0.4pt}


%% ============================================================
\section{Problem \#141}
\label{prob:141}

\noindent\textbf{Status:} OPEN \quad|\quad \textbf{Prize:} none \quad|\quad \textbf{Updated:} 2025-08-31

\noindent\textbf{Tags:} additive combinatorics, primes, arithmetic progressions

\medskip

\noindent\textbf{Statement:}

Let $k\geq 3$. Are there $k$ consecutive primes in arithmetic progression?




Green and Tao \cite{GrTa08} have proved that there must always exist some $k$ primes in arithmetic progression, but these need not be consecutive. Erd\H{o}s called this conjecture 'completely hopeless at present'. The existence of such progressions for small $k$ has been verified for $k\leq 10$, see the Wikipedia page . It is open, even for $k=3$, whether there are infinitely many such progressions. See also [219] . This is discussed in problem A6 of Guy's collection \cite{Gu04}.




References

[GrTa08] Green, Ben and Tao, Terence, The primes contain arbitrarily long arithmetic progressions . Ann. of Math. (2) (2008), 481-547.

[Gu04] Guy, Richard K., Unsolved problems in number theory . (2004), xviii+437.

\noindent\rule{\linewidth}{0.4pt}


%% ============================================================
\section{Problem \#142}
\label{prob:142}

\noindent\textbf{Status:} OPEN \quad|\quad \textbf{Prize:} \$10000 \quad|\quad \textbf{Updated:} 2025-08-31

\noindent\textbf{Tags:} additive combinatorics, arithmetic progressions

\medskip

\noindent\textbf{Statement:}

Let $r_k(N)$ be the largest possible size of a subset of $\{1,\ldots,N\}$ that does not contain any non-trivial $k$-term arithmetic progression. Prove an asymptotic formula for $r_k(N)$.




Erd\H{o}s remarked this is 'probably unattackable at present'. In \cite{Er97c} Erd\H{o}s offered \$1000, but given that he elsewhere offered \$5000 just for (essentially) showing that $r_k(N)=o_k(N/\log N)$ (see [3] ), that value seems odd. In \cite{Er81} he offers \$10000, stating it is 'probably enormously difficult'. The best known upper bounds for $r_k(N)$ are due to Kelley and Meka \cite{KeMe23} for $k=3$, Green and Tao \cite{GrTa17} for $k=4$, and Leng, Sah, and Sawhney \cite{LSS24} for $k\geq 5$. An asymptotic formula is still far out of reach, even for $k=3$. In \cite{Er80} and \cite{Va99} he asks (much more reasonably) for the order of magnitude of $r_k(N)$. In \cite{Er80} he remarks that we do not even know whether $r_k(n)/r_{k+1}(n)\to 0$ for any $k\geq 3$. See also [3] and [139] .




References

[Er80] Erd\H{o}s, Paul, A survey of problems in combinatorial number theory . Ann. Discrete Math. (1980), 89-115.

[Er81] Erd\H{o}s, P., On the combinatorial problems which I would most like to see solved . Combinatorica (1981), 25-42.

[Er97c] Erd\H{o}s, Paul, Some of my favorite problems and results . The mathematics of Paul Erd\H{o}s, I (1997), 47-67.

[GrTa17] Green, Ben and Tao, Terence, New bounds for Szemer\'{e}di's theorem, III: a polylogarithmic bound for $r_4(N)$ . Mathematika (2017), 944-1040.

[KeMe23] Kelley, Z. and Meka, R., Strong Bounds for 3-Progressions . arXiv:2302.05537 (2023).

[LSS24] Leng, J., Sah, A. and Sawhney, M., Improved bounds for Szemer\'{e}di's theorem . arXiv:2402.17995 (2024).

[Va99] Various, Some of Paul's favorite problems . Booklet produced for the conference "Paul Erd\H{o}s and his mathematics", Budapest, July 1999 (1999).

\noindent\rule{\linewidth}{0.4pt}


%% ============================================================
\section{Problem \#143}
\label{prob:143}

\noindent\textbf{Status:} OPEN \quad|\quad \textbf{Prize:} \$500 \quad|\quad \textbf{Updated:} 2025-08-31

\noindent\textbf{Tags:} primitive sets

\medskip

\noindent\textbf{Statement:}

Let $A\subset (1,\infty)$ be a countably infinite set such that for all $x\neq y\in A$ and integers $k\geq 1$ we have\[ \lvert kx -y\rvert \geq 1.\]Does this imply that $A$ is sparse? In particular, does this imply that\[\sum_{x\in A}\frac{1}{x\log x}&#60;\infty\]or\[\sum_{\substack{x &#60;n\\ x\in A}}\frac{1}{x}=o(\log n)?\]




Note that if $A$ is a set of integers then the condition implies that $A$ is a primitive set (that is, no element of $A$ is divisible by any other), for which the convergence of $\sum_{n\in A}\frac{1}{n\log n}$ was proved by Erd\H{o}s \cite{Er35}, and the upper bound\[\sum_{\substack{x &lt;n\\ x\in A}}\frac{1}{n}\ll \frac{\log x}{\sqrt{\log\log x}}\]was proved by Behrend \cite{Be35}. This $O(\cdot)$ bound was improved to a $o(\cdot)$ bound by Erd\H{o}s, S\'{a}rk\H{o}zy, and Szemer\'{e}di \cite{ESS67}. In \cite{Er73} and \cite{Er77c} Erd\H{o}s mentions an unpublished proof of Haight that\[\lim \frac{\lvert A\cap [1,x]\rvert}{x}=0\]holds if the elements of $A$ are independent over $\mathbb{Q}$. Over the years Erd\H{o}s asked for various different quantitative estimates, for example\[\liminf \frac{\lvert A\cap [1,x]\rvert}{x}=0\]or even (motivated by Behrend's bound)\[\sum_{\substack{x &lt;n\\ x\in A}}\frac{1}{x}\ll \frac{\log x}{\sqrt{\log\log x}}.\]In \cite{Er97c} he offers \$500 for resolving the questions in the main problem statement above. This was partially resolved by Koukoulopoulos, Lamzouri, and Lichtman \cite{KLL25}, who proved that we must have\[\sum_{\substack{x &lt;n\\ x\in A}}\frac{1}{x}=o(\log n).\]See also [858] .




References

[Be35] Behrend, F., On sequences of numbers not divisible by another . London Math. Soc. Journal (1935), 42-45.

[ESS67] Erd\H{o}s, P. and S\'{a}rk\"ozy, A. and Szemer\'{e}di, E., On a theorem of Behrend . J. Austral. Math. Soc. (1967), 9--16.

[Er35] Erd\H{o}s, Paul, Note on Sequences of Integers No One of Which is Divisible By Any Other . J. London Math. Soc. (1935), 126-128.

[Er73] Erd\H{o}s, P., Problems and results on combinatorial number theory . A survey of combinatorial theory (Proc. Internat. Sympos., Colorado State Univ., Fort Collins, Colo., 1971) (1973), 117-138.

[Er77c] Erd\H{o}s, Paul, Problems and results on combinatorial number theory. III . Number theory day (Proc. Conf., Rockefeller Univ., New York, 1976) (1977), 43-72.

[Er97c] Erd\H{o}s, Paul, Some of my favorite problems and results . The mathematics of Paul Erd\H{o}s, I (1997), 47-67.

[KLL25] D. Koukoulopoulos, Y. Lamzouri, and J. D. Lichtman, Erd\H{o}s's integer dilation approximation problem and GCD graphs . arXiv:2502.09539 (2025).

\noindent\rule{\linewidth}{0.4pt}


%% ============================================================
\section{Problem \#145}
\label{prob:145}

\noindent\textbf{Status:} OPEN \quad|\quad \textbf{Prize:} none \quad|\quad \textbf{Updated:} 2025-08-31

\noindent\textbf{Tags:} number theory

\medskip

\noindent\textbf{Statement:}

Let $s_1&#60;s_2&#60;\cdots$ be the sequence of squarefree numbers. Is it true that, for any $\alpha \geq 0$,\[\lim_{x\to \infty}\frac{1}{x}\sum_{s_n\leq x}(s_{n+1}-s_n)^\alpha\]exists?




Erd\H{o}s \cite{Er51} proved this for all $0\leq \alpha \leq 2$, and Hooley \cite{Ho73} extended this to all $\alpha \leq 3$. Greaves, Harman, and Huxley showed (in Chapter 11 of \cite{GHH97}) that this is true for $\alpha \leq 11/3$. Chan \cite{Ch23c} has extended this to $\alpha \leq 3.75$. Granville \cite{Gr98} proved that this follows (for all $\alpha \geq 0$) from the ABC conjecture . See also [208] .




References

[Ch23c] Chan, Tsz Ho, On moments of gaps between consecutive square-free numbers . Mosc. J. Comb. Number Theory (2023), 287--295.

[Er51] Erd\H{o}s, P., Some problems and results in elementary number theory . Publ. Math. Debrecen (1951), 103-109.

[GHH97] Greaves, G. R. H. and Harman, G. and Huxley, M. N., Sieve Methods, Exponential Sums, and their Applications in Number Theory . (1997).

[Gr98] Granville, Andrew, {$ABC$} allows us to count squarefrees . Internat. Math. Res. Notices (1998), 991--1009.

[Ho73] Hooley, Christopher, On the intervals between consecutive terms of sequences . (1973), 129-140.

\noindent\rule{\linewidth}{0.4pt}


%% ============================================================
\section{Problem \#146}
\label{prob:146}

\noindent\textbf{Status:} OPEN \quad|\quad \textbf{Prize:} \$500 \quad|\quad \textbf{Updated:} 2025-08-31

\noindent\textbf{Tags:} graph theory, turan number

\medskip

\noindent\textbf{Statement:}

If $H$ is bipartite and is $r$-degenerate, that is, every induced subgraph of $H$ has minimum degree $\leq r$, then\[\mathrm{ex}(n;H) \ll n^{2-1/r}.\]




Conjectured by Erd\H{o}s and Simonovits \cite{ErSi84}. Open even for $r=2$. Alon, Krivelevich, and Sudakov \cite{AKS03} have proved\[\mathrm{ex}(n;H) \ll n^{2-1/4r}.\]They also prove the full Erd\H{o}s-Simonovits conjectured bound if $H$ is bipartite and the maximum degree in one side of the bipartition is $r$. See also [113] and [147] . This problem is #43 in Extremal Graph Theory in the graphs problem collection.




References

[AKS03] Alon, Noga and Krivelevich, Michael and Sudakov, Benny, Tur\'{a}n numbers of bipartite graphs and related Ramsey-type questions . Combin. Probab. Comput. (2003), 477-494.

[ErSi84] Erd\H{o}s, P. and Simonovits, M., Cube-supersaturated graphs and related problems . Progress in graph theory (Waterloo, Ont., 1982) (1984), 203-218.

\noindent\rule{\linewidth}{0.4pt}


%% ============================================================
\section{Problem \#148}
\label{prob:148}

\noindent\textbf{Status:} OPEN \quad|\quad \textbf{Prize:} none \quad|\quad \textbf{Updated:} 2025-08-31

\noindent\textbf{Tags:} number theory, unit fractions

\medskip

\noindent\textbf{Statement:}

Let $F(k)$ be the number of solutions to\[ 1= \frac{1}{n_1}+\cdots+\frac{1}{n_k},\]where $1\leq n_1&#60;\cdots&#60;n_k$ are distinct integers. Find good estimates for $F(k)$.




The current best bounds known are\[2^{c^{\frac{k}{\log k}}}\leq F(k) \leq c_0^{(\frac{1}{5}+o(1))2^k},\]where $c&#62;0$ is some absolute constant and $c_0=1.26408\cdots$ is the 'Vardi constant'. The lower bound is due to Konyagin \cite{Ko14} and the upper bound to Elsholtz and Planitzer \cite{ElPl21}.




References

[ElPl21] Elsholtz, Christian and Planitzer, Stefan, Sums of four and more unit fractions and approximate parametrizations . Bull. Lond. Math. Soc. (2021), 695-709.

[Ko14] Konyagin, S. V., Double exponential lower bound for the number of representations of unity by Egyptian fractions . Math. Notes (2014), 277-281.

\noindent\rule{\linewidth}{0.4pt}


%% ============================================================
\section{Problem \#149}
\label{prob:149}

\noindent\textbf{Status:} OPEN \quad|\quad \textbf{Prize:} none \quad|\quad \textbf{Updated:} 2025-08-31

\noindent\textbf{Tags:} graph theory

\medskip

\noindent\textbf{Statement:}

The strong chromatic index of a graph $G$, denoted by $\mathrm{sq}(G)$, is the minimum $k$ such that the edges of $G$ can be partitioned into $k$ sets of 'strongly independent' edges, that is, such that the subgraph of $G$ induced by each set is the union of vertex-disjoint edges. Is it true that, for any graph $G$ with maximum degree $\Delta$,\[\mathrm{sq}(G)\leq\frac{5}{4}\Delta^2?\]




Asked by Erd\H{o}s and Ne\v{s}et\v{r}il in 1985 (see \cite{FGST89}). This is equivalent to asking whether the chromatic number of the square of the line graph $L(G)^2$ is at most $\frac{5}{4}\Delta^2$. This bound would be the best possible, as witnessed by a blowup of $C_5$, although improvements may be possible for odd $\Delta$. It is easy to see that $\mathrm{sq}(G)\leq 5$ when $\Delta\leq 2$: in general, the trivial bound of\[\mathrm{sq}(G)\leq 2\Delta^2-2\Delta+1\]follows from the observations that the number of neighbours of any edge of $G$ is $\leq 2\Delta^2-2\Delta$. This was improved to $\mathrm{sq}(G)\leq 1.998\Delta^2$ (for $\Delta$ sufficiently large) by Molloy and Reed \cite{MoRe97}. This was improved to $1.93\Delta^2$ by Bruhn and Joos \cite{BrJo18} and to $1.835\Delta^2$ by Bonamy, Perrett, and Postle \cite{BPP22}. The best bound currently available is\[1.772\Delta^2,\]proved by Hurley, de Joannis de Verclos, and Kang \cite{HJK22}. Mahdian has, in their Masters' thesis , proved an upper bound of $(2+o(1))\frac{\Delta^2}{\log \Delta}$ under the additional assumption that $G$ is $C_4$-free. Andersen \cite{92} and Hor\'{a}k, He, and Trotter \cite{HHT93} have independently shown that $\mathrm{sq}(G)\leq 10$ when $\Delta\leq 3$. (This is best possible, as shown e.g. by a $C_8$ with all four diagonals). Huang, Santana, and Yu \cite{HSY18} proved that $\mathrm{sq}(G)\leq 21$ when $\Delta\leq 4$. Erd\H{o}s and Ne\v{s}et\v{r}il also asked the easier problem of whether $G$ containing at least $\tfrac{5}{4}\Delta^2$ many edges implies $G$ containing two strongly independent edges. This was proved by Chung, Gy\'{a}rf\'{a}s, Tuza, and Trotter \cite{CGTT90}. It is still open even whether the clique number of $L(G)^2$ at most $\frac{5}{4}\Delta^2$. Let $\omega=\omega(L(G)^2)$ be this clique number. \'{S}leszy\'{n}ska-Nowak \cite{Sl16} proved $\omega \leq \frac{3}{2}\Delta^2$. Faron and Postle \cite{FaPo19} proved $\omega\leq \frac{4}{3}\Delta^2$. Cames van Batenburg, Kang, and Pirot \cite{CKP20} have proved $\omega\leq \frac{5}{4}\Delta^2$ under the additional assumption that $G$ is triangle-free (and $\omega\leq \Delta^2$ if $G$ is $C_5$-free).




References

[92] No reference found.


[BPP22] Bonamy, Marthe and Perrett, Thomas and Postle, Luke, Colouring graphs with sparse neighbourhoods: bounds and applications . J. Combin. Theory Ser. B (2022), 278-317.

[BrJo18] Bruhn, Henning and Joos, Felix, A stronger bound for the strong chromatic index . Combin. Probab. Comput. (2018), 21-43.

[CGTT90] Chung, F. R. K. and Gy\'{a}rf\'{a}s, A. and Tuza, Z. and Trotter, W. T., The maximum number of edges in {$2K_2$}-free graphs of bounded degree . Discrete Math. (1990), 129--135.

[CKP20] Cames van Batenburg, Wouter and Kang, Ross J. and Pirot, Fran\c cois, Strong cliques and forbidden cycles . Indag. Math. (N.S.) (2020), 64--82.

[FGST89] Faudree, R. J. and Gy\'{a}rf\'{a}s, A. and Schelp, R. H. and Tuza, Zs., Induced matchings in bipartite graphs . Discrete Math. (1989), 83-87.

[FaPo19] Faron, Maxime and Postle, Luke, On the clique number of the square of a line graph and its relation to maximum degree of the line graph . J. Graph Theory (2019), 261--274.

[HHT93] Hor\'{a}k, Peter and He, Qing and Trotter, William T., Induced matchings in cubic graphs . J. Graph Theory (1993), 151--160.

[HJK22] Hurley, Eoin and de Joannis de Verclos, R\'{e}mi and Kang, Ross J., An improved procedure for colouring graphs of bounded local density . Adv. Comb. (2022), Paper No. 7, 33.

[HSY18] Huang, Mingfang and Santana, Michael and Yu, Gexin, Strong chromatic index of graphs with maximum degree four . Electron. J. Combin. (2018), Paper No. 3.31, 24.

[MoRe97] Molloy, Michael and Reed, Bruce, A bound on the strong chromatic index of a graph . J. Combin. Theory Ser. B (1997), 103-109.

[Sl16] \'Sleszy\'nska-Nowak, Ma\l gorzata, Clique number of the square of a line graph . Discrete Math. (2016), 1551--1556.

\noindent\rule{\linewidth}{0.4pt}


%% ============================================================
\section{Problem \#151}
\label{prob:151}

\noindent\textbf{Status:} OPEN \quad|\quad \textbf{Prize:} none \quad|\quad \textbf{Updated:} 2025-08-31

\noindent\textbf{Tags:} graph theory

\medskip

\noindent\textbf{Statement:}

For a graph $G$ let $\tau(G)$ denote the minimal number of vertices that include at least one from each maximal clique of $G$ on at least two vertices (sometimes called the clique transversal number). Let $H(n)$ be maximal such that every triangle-free graph on $n$ vertices contains an independent set on $H(n)$ vertices. If $G$ is a graph on $n$ vertices then is\[\tau(G)\leq n-H(n)?\]




It is easy to see that $\tau(G) \leq n-\sqrt{n}$. Note also that if $G$ is triangle-free then trivially $\tau(G)\leq n-H(n)$. This is listed in \cite{Er88} as a problem of Erd\H{o}s and Gallai, who were unable to make progress even assuming $G$ is $K_4$-free. There Erd\H{o}s remarked that this conjecture is 'perhaps completely wrongheaded'. It later appeared as Problem 1 in \cite{EGT92}. The general behaviour of $\tau(G)$ is the subject of [610] .




References

[EGT92] Erd\H{o}s, Paul and Gallai, Tibor and Tuza, Zsolt, Covering the cliques of a graph with vertices . Discrete Math. (1992), 279-289.

[Er88] Erd\H{o}s, P, Problems and results in combinatorial analysis and graph theory . Discrete Math. (1988), 81-92.

\noindent\rule{\linewidth}{0.4pt}


%% ============================================================
\section{Problem \#153}
\label{prob:153}

\noindent\textbf{Status:} OPEN \quad|\quad \textbf{Prize:} none \quad|\quad \textbf{Updated:} 2025-08-31

\noindent\textbf{Tags:} sidon sets

\medskip

\noindent\textbf{Statement:}

Let $A$ be a finite Sidon set and $A+A=\{s_1&#60;\cdots&#60;s_t\}$. Is it true that\[\frac{1}{t}\sum_{1\leq i&#60;t}(s_{i+1}-s_i)^2 \to \infty\]as $\lvert A\rvert\to \infty$?




A similar problem can be asked for infinite Sidon sets.

\noindent\rule{\linewidth}{0.4pt}


%% ============================================================
\section{Problem \#155}
\label{prob:155}

\noindent\textbf{Status:} OPEN \quad|\quad \textbf{Prize:} none \quad|\quad \textbf{Updated:} 2025-08-31

\noindent\textbf{Tags:} additive combinatorics, sidon sets

\medskip

\noindent\textbf{Statement:}

Let $F(N)$ be the size of the largest Sidon subset of $\{1,\ldots,N\}$. Is it true that for every $k\geq 1$ we have\[F(N+k)\leq F(N)+1\]for all sufficiently large $N$?




This may even hold with $k\approx \epsilon N^{1/2}$.

\noindent\rule{\linewidth}{0.4pt}


%% ============================================================
\section{Problem \#156}
\label{prob:156}

\noindent\textbf{Status:} OPEN \quad|\quad \textbf{Prize:} none \quad|\quad \textbf{Updated:} 2025-08-31

\noindent\textbf{Tags:} sidon sets

\medskip

\noindent\textbf{Statement:}

Does there exist a maximal Sidon set $A\subset \{1,\ldots,N\}$ of size $O(N^{1/3})$?




A question of Erd\H{o}s, S\'{a}rk\"{o}zy, and S\'{o}s \cite{ESS94}. It is easy to prove that the greedy construction of a maximal Sidon set in $\{1,\ldots,N\}$ has size $\gg N^{1/3}$. Ruzsa \cite{Ru98b} constructed a maximal Sidon set of size $\ll (N\log N)^{1/3}$. See also [340] .




References

[ESS94] Erd\H{o}s, P. and S\'{a}rk\"{o}zy, A. and S\'{o}s, T., On Sum Sets of Sidon Sets, I . Journal of Number Theory (1994), 329-347.

[Ru98b] Ruzsa, Imre Z., A small maximal Sidon set . Ramanujan J. (1998), 55-58.

\noindent\rule{\linewidth}{0.4pt}


%% ============================================================
\section{Problem \#158}
\label{prob:158}

\noindent\textbf{Status:} OPEN \quad|\quad \textbf{Prize:} none \quad|\quad \textbf{Updated:} 2025-08-31

\noindent\textbf{Tags:} sidon sets

\medskip

\noindent\textbf{Statement:}

Let $A\subset \mathbb{N}$ be an infinite set such that, for any $n$, there are most $2$ solutions to $a+b=n$ with $a\leq b$. Must\[\liminf_{N\to\infty}\frac{\lvert A\cap \{1,\ldots,N\}\rvert}{N^{1/2}}=0?\]




If we replace $2$ by $1$ then $A$ is a Sidon set, for which Erd\H{o}s proved this is true.

\noindent\rule{\linewidth}{0.4pt}


%% ============================================================
\section{Problem \#159}
\label{prob:159}

\noindent\textbf{Status:} OPEN \quad|\quad \textbf{Prize:} none \quad|\quad \textbf{Updated:} 2025-08-31

\noindent\textbf{Tags:} graph theory, ramsey theory

\medskip

\noindent\textbf{Statement:}

There exists some constant $c&#62;0$ such that $$R(C_4,K_n) \ll n^{2-c}.$$




The prize of \$100 is offered in \cite{Er78} for a proof or disproof. The current bounds are\[ \frac{n^{3/2}}{(\log n)^{3/2}}\ll R(C_4,K_n)\ll \frac{n^2}{(\log n)^2}.\]The upper bound is due to Szemer\'{e}di (mentioned in \cite{EFRS78}), and the lower bound is due to Spencer \cite{Sp77}. This problem is #17 in Ramsey Theory in the graphs problem collection.




References

[EFRS78] Erd\H{o}s, Paul and Faudree, R. J. and Rousseau, C. C. and Schelp, R. H., On cycle-complete graph Ramsey numbers . J. Graph Theory (1978), 53-64.

[Er78] Erd\H{o}s, Paul, Problems and results in combinatorial analysis and combinatorial number theory . Proceedings of the Ninth Southeastern Conference on Combinatorics, Graph Theory, and Computing (Florida Atlantic Univ., Boca Raton, Fla., 1978) (1978), 29-40.

[Sp77] Spencer, J., Asymptotic lower bounds for Ramsey functions . Discrete Math. (1977), 69-76.

\noindent\rule{\linewidth}{0.4pt}


%% ============================================================
\section{Problem \#160}
\label{prob:160}

\noindent\textbf{Status:} OPEN \quad|\quad \textbf{Prize:} none \quad|\quad \textbf{Updated:} 2025-08-31

\noindent\textbf{Tags:} additive combinatorics, arithmetic progressions

\medskip

\noindent\textbf{Statement:}

Let $h(N)$ be the smallest $k$ such that $\{1,\ldots,N\}$ can be coloured with $k$ colours so that every four-term arithmetic progression must contain at least three distinct colours. Estimate $h(N)$.




Investigated by Erd\H{o}s and Freud. This has been discussed on MathOverflow , where LeechLattice shows\[h(N) \ll N^{2/3}.\]In the comments of this site Hunter improves this to\[h(N) \ll N^{\frac{\log 3}{\log 22}+o(1)}\](note $\frac{\log 3}{\log 22}\approx 0.355$). The observation of Zach Hunter in that question coupled with recent progress on the size of subsets without three-term arithmetic progression (see \cite{BlSi23} which improves slightly on the bounds due to Kelley and Meka \cite{KeMe23}) imply that\[h(N) \gg \exp(c(\log N)^{1/9})\]for some $c&gt;0$.




References

[BlSi23] T. F. Bloom and O. Sisask, An improvement to the Kelley-Meka bounds on three-term arithmetic progressions . arXiv:2309.02353 (2023).

[KeMe23] Kelley, Z. and Meka, R., Strong Bounds for 3-Progressions . arXiv:2302.05537 (2023).

\noindent\rule{\linewidth}{0.4pt}


%% ============================================================
\section{Problem \#161}
\label{prob:161}

\noindent\textbf{Status:} OPEN \quad|\quad \textbf{Prize:} \$500 \quad|\quad \textbf{Updated:} 2025-08-31

\noindent\textbf{Tags:} combinatorics, ramsey theory, discrepancy

\medskip

\noindent\textbf{Statement:}

Let $\alpha\in[0,1/2)$ and $n,t\geq 1$. Let $F^{(t)}(n,\alpha)$ be the smallest $m$ such that we can $2$-colour the edges of the complete $t$-uniform hypergraph on $n$ vertices such that if $X\subseteq [n]$ with $\lvert X\rvert \geq m$ then there are at least $\alpha \binom{\lvert X\rvert}{t}$ many $t$-subsets of $X$ of each colour. For fixed $n,t$ as we change $\alpha$ from $0$ to $1/2$ does $F^{(t)}(n,\alpha)$ increase continuously or are there jumps? Only one jump?




For $\alpha=0$ this is the usual Ramsey function. A conjecture of Erd\H{o}s, Hajnal, and Rado (see [562] ) implies that\[ F^{(t)}(n,0)\asymp \log_{t-1} n\]and results of Erd\H{o}s and Spencer imply that\[F^{(t)}(n,\alpha) \gg_\alpha (\log n)^{\frac{1}{t-1}}\]for all $\alpha&gt;0$, and a similar upper bound holds for $\alpha$ close to $1/2$. Erd\H{o}s said in \cite{Er90b}: 'If I can hazard a guess completely unsupported by evidence, I am afraid that the jump occurs all in one step at $0$. It would be much more interesting if my conjecture would be wrong and perhaps there is some hope for this for $t&gt;3$. I know nothing and offer \$500 to anybody who can clear up this mystery.' Conlon, Fox, and Sudakov \cite{CFS11} have proved that, for any fixed $\alpha&gt;0$,\[F^{(3)}(n,\alpha) \ll_\alpha \sqrt{\log n}.\]Coupled with the lower bound above, this implies that there is only one jump for fixed $\alpha$ when $t=3$, at $\alpha=0$. For all $\alpha&gt;0$ it is known that\[F^{(t)}(n,\alpha)\gg_t (\log n)^{c_\alpha}.\]See also [563] for more on the case $t=2$. This problem is #40 in Ramsey Theory in the graphs problem collection.




References

[CFS11] Conlon, David and Fox, Jacob and Sudakov, Benny, Large almost monochromatic subsets in hypergraphs . Israel J. Math. (2011), 423--432.

[Er90b] Erd\H{o}s, Paul, Problems and results on graphs and hypergraphs: similarities and differences . Mathematics of Ramsey theory (1990), 12-28.

\noindent\rule{\linewidth}{0.4pt}


%% ============================================================
\section{Problem \#162}
\label{prob:162}

\noindent\textbf{Status:} OPEN \quad|\quad \textbf{Prize:} none \quad|\quad \textbf{Updated:} 2025-08-31

\noindent\textbf{Tags:} combinatorics, ramsey theory, discrepancy

\medskip

\noindent\textbf{Statement:}

Let $\alpha&#62;0$ and $n\geq 1$. Let $F(n,\alpha)$ be the largest $k$ such that there exists some 2-colouring of the edges of $K_n$ in which any induced subgraph $H$ on at least $k$ vertices contains more than $\alpha\binom{\lvert H\rvert}{2}$ many edges of each colour. Prove that for every fixed $0\leq \alpha \leq 1/2$, as $n\to\infty$,\[F(n,\alpha)\sim c_\alpha \log n\]for some constant $c_\alpha$.




It is easy to show with the probabilistic method that there exist $c_1(\alpha),c_2(\alpha)$ such that\[c_1(\alpha)\log n &#60; F(n,\alpha) &#60; c_2(\alpha)\log n.\]

\noindent\rule{\linewidth}{0.4pt}


%% ============================================================
\section{Problem \#165}
\label{prob:165}

\noindent\textbf{Status:} OPEN \quad|\quad \textbf{Prize:} \$250 \quad|\quad \textbf{Updated:} 2025-08-31

\noindent\textbf{Tags:} graph theory, ramsey theory

\medskip

\noindent\textbf{Statement:}

Give an asymptotic formula for $R(3,k)$.




It is known that there exists some constant $c&gt;0$ such that for large $k$\[(c+o(1))\frac{k^2}{\log k}\leq R(3,k) \leq (1+o(1))\frac{k^2}{\log k}.\]The lower bound is due to Kim \cite{Ki95}, the upper bound is due to Shearer \cite{Sh83}, improving an earlier bound of Ajtai, Koml\'{o}s, and Szemer\'{e}di \cite{AKS80}. The value of $c$ in the lower bound has seen a number of improvements. Kim's original proof gave $c\geq 1/162$. The bound $c\geq 1/4$ was proved independently by Bohman and Keevash \cite{BoKe21} and Pontiveros, Griffiths and Morris \cite{PGM20}. The latter collection of authors conjecture that this lower bound is the true order of magnitude. This was, however, improved by Campos, Jenssen, Michelen, and Sahasrabudhe \cite{CJMS25} to $c\geq 1/3$, and further by Hefty, Horn, King, and Pfender \cite{HHKP25} to $c\geq 1/2$. Both of these papers conjecture that $c=1/2$ is the correct asymptotic. See also [544] , and [986] for the general case. See [1013] for a related function.




References

[AKS80] Ajtai, Mikl\'{o}s and Koml\'{o}s, J\'{a}nos and Szemer\'{e}di, Endre, A note on Ramsey numbers . J. Combin. Theory Ser. A (1980), 354-360.

[BoKe21] Bohman, Tom and Keevash, Peter, Dynamic concentration of the triangle-free process . Random Structures Algorithms (2021), 221-293.

[CJMS25] M. Campos, M. Jenssen, M. Michelen, and J. Sahasrabudhe, A new lower bound for the Ramsey numbers $R(3,k)$ . arXiv:2505.13371 (2025).

[HHKP25] Z. Hefty, P. Horn, D. King, and F. Pfender, Improving $R(3,k)$ in just two bites . arXiv:2510.19718 (2025).

[Ki95] Kim, J. H., The Ramsey number $R(3,t)$ has order of magnitude $t^2/\log t$ . Random Structures and Algorithms (1995), 173-207.

[PGM20] Fiz Pontiveros, Gonzalo and Griffiths, Simon and Morris, Robert, The triangle-free process and the Ramsey number $R(3,k)$ . Mem. Amer. Math. Soc. (2020), v+125.

[Sh83] Shearer J., A note on the independence number of triangle-free graphs . Discrete Math. (1983), 83-87.

\noindent\rule{\linewidth}{0.4pt}


%% ============================================================
\section{Problem \#168}
\label{prob:168}

\noindent\textbf{Status:} OPEN \quad|\quad \textbf{Prize:} none \quad|\quad \textbf{Updated:} 2025-08-31

\noindent\textbf{Tags:} additive combinatorics

\medskip

\noindent\textbf{Statement:}

Let $F(N)$ be the size of the largest subset of $\{1,\ldots,N\}$ which does not contain any set of the form $\{n,2n,3n\}$. What is\[ \lim_{N\to \infty}\frac{F(N)}{N}?\]Is this limit irrational?




This limit was proved to exist by Graham, Spencer, and Witsenhausen \cite{GSW77}, who showed it is equal to\[\frac{1}{3}\sum_{k\in K}\frac{1}{d_k},\]where $d_1&lt;d_2&lt;\cdots $are the $3$-smooth numbers and $K$ is the set of $k$ for which $f(k)&gt;f(k-1)$, where $f$ counts the largest subset of $\{d_1,\ldots,d_k\}$ that avoids $\{n,2n,3n\}$. Similar questions can be asked for the density or upper density of infinite sets without such configurations. The limit can be estimated by elementary arguments (see the comments). Eberhard has used the formula of \cite{GSW77} mentioned above to calculate the value of the limit as\[0.800965\cdots.\]




References

[GSW77] Graham, R. and Spencer, J. and Witsenhausen, H., On Extremal Density Theorems for Linear Forms . Number Theory and Algebra (1977).

\noindent\rule{\linewidth}{0.4pt}


%% ============================================================
\section{Problem \#169}
\label{prob:169}

\noindent\textbf{Status:} OPEN \quad|\quad \textbf{Prize:} none \quad|\quad \textbf{Updated:} 2025-08-31

\noindent\textbf{Tags:} additive combinatorics, arithmetic progressions

\medskip

\noindent\textbf{Statement:}

Let $k\geq 3$ and $f(k)$ be the supremum of $\sum_{n\in A}\frac{1}{n}$ as $A$ ranges over all sets of positive integers which do not contain a $k$-term arithmetic progression. Estimate $f(k)$. Is\[\lim_{k\to \infty}\frac{f(k)}{\log W(k)}=\infty\]where $W(k)$ is the van der Waerden number?




Berlekamp \cite{Be68} proved $f(k) \geq \frac{\log 2}{2}k$. Gerver \cite{Ge77} proved\[f(k) \geq (1-o(1))k\log k.\]It is trivial that\[\frac{f(k)}{\log W(k)}\geq \frac{1}{2},\]but improving the right-hand side to any constant $&gt;1/2$ is open. Gerver also proved (see the comments for an alternative argument of Tao) that [3] is equivalent to $f(k)$ being finite for all $k$. The current record for $f(3)$ is $f(3)\geq 3.00849$, due to Wr\'{o}blewski \cite{Wr84}. Walker \cite{Wa25} proved $f(4)\geq 4.43975$. Walker \cite{Wa25} has shown that it suffices to consider Kempner sets (that is, sets of integers defined as all those whose base $b$ digits are contained in some $S\subset \{0,\ldots,b-1\}$ for fixed $b$ and $S$), in the sense that for any $k\geq 3$ and $\epsilon&gt;0$ there is a Kempner set $A$ lacking $k$-term arithmetic progressions such that\[\sum_{n\in A}\frac{1}{n}\geq f(k)-\epsilon.\]In \cite{Er80} Erd\H{o}s asks whether for every $\epsilon&gt;0$ and $k\geq 3$ if $A$ is a set of integers without a $k$-term arithmetic progression such that $\min(A)$ is sufficiently large in terms of $\epsilon$ and $k$ then $\sum_{n\in A}\frac{1}{n}&lt;\epsilon$.




References

[Be68] Berlekamp, E. R., A construction for partitions which avoid long arithmetic progressions . Canad. Math. Bull. (1968), 409-414.

[Er80] Erd\H{o}s, Paul, A survey of problems in combinatorial number theory . Ann. Discrete Math. (1980), 89-115.

[Ge77] Gerver, Joseph L., The sum of the reciprocals of a set of integers with no arithmetic progression of {$k$} terms . Proc. Amer. Math. Soc. (1977), 211--214.

[Wa25] A. Walker, Integer sets of large harmonic sum which avoid long arithmetic progressions . arXiv:2203.06045 (2025).

[Wr84] Wr\'oblewski, J., A nonaveraging set of integers with a large sum of reciprocals . Math. Comp. (1984), 261--262.

\noindent\rule{\linewidth}{0.4pt}


%% ============================================================
\section{Problem \#170}
\label{prob:170}

\noindent\textbf{Status:} OPEN \quad|\quad \textbf{Prize:} none \quad|\quad \textbf{Updated:} 2025-08-31

\noindent\textbf{Tags:} additive combinatorics

\medskip
\noindent\textit{sparse ruler problem}


\noindent\textbf{Statement:}

Let $F(N)$ be the smallest possible size of $A\subset \{0,1,\ldots,N\}$ such that $\{0,1,\ldots,N\}\subset A-A$. Find the value of\[\lim_{N\to \infty}\frac{F(N)}{N^{1/2}}.\]




The Sparse Ruler problem. R\'{e}dei asked whether this limit exists, which was proved by Erd\H{o}s and G\'{a}l \cite{ErGa48}. Bounds on the limit were improved by Leech \cite{Le56}. The limit is known to be in the interval $[1.56,\sqrt{3}]$. The lower bound is due to Leech \cite{Le56}, the upper bound is due to Wichmann \cite{Wi63}. Computational evidence by Pegg \cite{Pe20} suggests that the upper bound is the truth. A similar question can be asked without the restriction $A\subset \{0,1,\ldots,N\}$.




References

[ErGa48] Erd\H{o}s, P. and G\'{a}l, I., On the representation of $1,2,\ldots,N$ by differences . Nederl. Akad. Wetensch., Proc. (1948), 1155-1158.

[Le56] Leech, J., On the representation of $1,2,\ldots,n$ by differences . J. London Math. Soc. (1956), 160-169.

[Pe20] Pegg, E., Hitting All the Marks: Exploring New Bounds for Sparse Rulers and a Wolfram Language Proof . https://blog.wolfram.com/2020/02/12/hitting-all-the-marks-exploring-new-bounds-for-sparse-rulers-and-a-wolfram-language-proof/ (2020).

[Wi63] Wichmann, B., A note on restricted difference bases . J. London Math. Soc. (1963), 465-466.

\noindent\rule{\linewidth}{0.4pt}


%% ============================================================
\section{Problem \#172}
\label{prob:172}

\noindent\textbf{Status:} OPEN \quad|\quad \textbf{Prize:} none \quad|\quad \textbf{Updated:} 2025-08-31

\noindent\textbf{Tags:} additive combinatorics, ramsey theory

\medskip

\noindent\textbf{Statement:}

Is it true that in any finite colouring of $\mathbb{N}$ there exist arbitrarily large finite $A$ such that all sums and products of distinct elements in $A$ are the same colour?




First asked by Hindman. Hindman \cite{Hi80} has proved this is false (with 7 colours) if we ask for an infinite $A$. In \cite{Er77c} Erd\H{o}s asks about the case for an infinite $A$ with just $2$ colours (see [1198] ). Moreira \cite{Mo17} has proved that in any finite colouring of $\mathbb{N}$ there exist $x,y$ such that $\{x,x+y,xy\}$ are all the same colour. Alweiss \cite{Al23} has proved that in any finite colouring of $\mathbb{Q}\backslash \{0\}$ there exist arbitrarily large finite $A$ such that all sums and products of distinct elements in $A$ are the same colour. Bowen and Sabok \cite{BoSa22} had proved this earlier for the first non-trivial case of $\lvert A\rvert=2$.




References

[Al23] R. Alweiss, Hindman's conjecture over the rationals . arXiv:2307.08901 (2023).

[BoSa22] M. Bowen and M. Sabok, Monochromatic Sums and Products in the Rationals . arXiv:2210.12290 (2022).

[Er77c] Erd\H{o}s, Paul, Problems and results on combinatorial number theory. III . Number theory day (Proc. Conf., Rockefeller Univ., New York, 1976) (1977), 43-72.

[Hi80] Hindman, Neil, Partitions and sums and products-two counterexamples . J. Combin. Theory Ser. A (1980), 113-120.

[Mo17] Moreira, J., Monochromatic sums and products in $\mathbbN$ . Ann. Math. (2017), 1069-1090.

\noindent\rule{\linewidth}{0.4pt}


%% ============================================================
\section{Problem \#173}
\label{prob:173}

\noindent\textbf{Status:} OPEN \quad|\quad \textbf{Prize:} none \quad|\quad \textbf{Updated:} 2025-08-31

\noindent\textbf{Tags:} geometry, ramsey theory

\medskip

\noindent\textbf{Statement:}

In any $2$-colouring of $\mathbb{R}^2$, for all but at most one triangle $T$, there is a monochromatic congruent copy of $T$.




For some colourings a single equilateral triangle has to be excluded, considering the colouring by alternating strips. Shader \cite{Sh76} has proved this is true if we just consider a single right-angled triangle.




References

[Sh76] Shader, L., All right triangles are Ramsey in $\mathbbE^2$! . J. Comb. Th. A (1976), 385-389.

\noindent\rule{\linewidth}{0.4pt}


%% ============================================================
\section{Problem \#174}
\label{prob:174}

\noindent\textbf{Status:} OPEN \quad|\quad \textbf{Prize:} none \quad|\quad \textbf{Updated:} 2025-08-31

\noindent\textbf{Tags:} geometry, ramsey theory

\medskip

\noindent\textbf{Statement:}

A finite set $A\subset \mathbb{R}^n$ is called Ramsey if, for any $k\geq 1$, there exists some $d=d(A,k)$ such that in any $k$-colouring of $\mathbb{R}^d$ there exists a monochromatic copy of $A$. Characterise the Ramsey sets in $\mathbb{R}^n$.




Erd\H{o}s, Graham, Montgomery, Rothschild, Spencer, and Straus \cite{EGMRSS73} proved that every Ramsey set is 'spherical': it lies on the surface of some sphere. Graham has conjectured that every spherical set is Ramsey. Leader, Russell, and Walters \cite{LRW12} have alternatively conjectured that a set is Ramsey if and only if it is 'subtransitive': it can be embedded in some higher-dimensional set on which rotations act transitively. Sets known to be Ramsey include vertices of $k$-dimensional rectangles \cite{EGMRSS73}, non-degenerate simplices \cite{FrRo90}, trapezoids \cite{Kr92}, and regular polygons/polyhedra \cite{Kr91}.




References

[EGMRSS73] Erd\H{o}s, P. and Graham, R. L. and Montgomery, P. and Rothschild, B. L. and Spencer, J. and Straus, E. G., Euclidean Ramsey Theorems I . J. Comb. Th. A (1973), 341-363.

[FrRo90] Frankl, P. and R\"{o}dl, V., A partition property of simplices in Euclidean space . J. Amer. Math. Soc. (1990), 1-7.

[Kr91] K\v{r}\'{\i}\v{z}, Igor, Permutation groups in Euclidean Ramsey theory . Proc. Amer. Math. Soc. (1991), 899-907.

[Kr92] K\v{r}\'{\i}\v{z}, Igor, All trapezoids are Ramsey . Discrete Math. (1992), 59-62.

[LRW12] Leader, Imre and Russell, Paul A. and Walters, Mark, Transitive sets in Euclidean Ramsey theory . J. Combin. Theory Ser. A (2012), 382-396.

\noindent\rule{\linewidth}{0.4pt}


%% ============================================================
\section{Problem \#176}
\label{prob:176}

\noindent\textbf{Status:} OPEN \quad|\quad \textbf{Prize:} none \quad|\quad \textbf{Updated:} 2025-08-31

\noindent\textbf{Tags:} additive combinatorics, arithmetic progressions, discrepancy

\medskip

\noindent\textbf{Statement:}

Let $N(k,\ell)$ be the minimal $N$ such that for any $f:\{1,\ldots,N\}\to\{-1,1\}$ there must exist a $k$-term arithmetic progression $P$ such that\[ \left\lvert \sum_{n\in P}f(n)\right\rvert\geq \ell.\]Find good upper bounds for $N(k,\ell)$. Is it true that for any $c&#62;0$ there exists some $C&#62;1$ such that\[N(k,ck)\leq C^k?\]What about\[N(k,2)\leq C^k\]or\[N(k,\sqrt{k})\leq C^k?\]




When $\ell=k$ this is the van der Waerden number $W(k)$ (see [138] ). Spencer \cite{Sp73} has proved that if $k=2^tm$ with $m$ odd then\[N(k,1)=2^t(k-1)+1.\]Erd\H{o}s and Graham write that 'no decent bound' is known even for $N(k,2)$. Erd\H{o}s \cite{Er63d} proved that, for every $c&gt;0$,\[N(k,ck)&gt; (1+\alpha_c)^k\]where $\alpha_c\to 0$ as $c\to 0$ and $\alpha_c\to \sqrt{2}-1$ as $c\to 1$. Hunter in the comment section observes that the local lemma implies an improved bound of\[N(k,ck) \gg \frac{2^k}{k^{O(1)}\sum_{i&gt;\frac{1+c}{2}k}\binom{k}{i}},\]so that in particular as $c\to 1$ we have, for all large $k$, $N(k,ck)\geq (2-o(1))^k$ (where the $o(1)$ term $\to 0$ as $c\to 1$).




References

[Er63d] Erd\H{o}s, P\'{a}l, On combinatorial questions connected with a theorem of {R}amsey and van der {W}aerden . Mat. Lapok (1963), 29--37.

[Sp73] J. Spencer, Problems 185 . Bull. Canad. Math. Soc. (1973), 185.

\noindent\rule{\linewidth}{0.4pt}


%% ============================================================
\section{Problem \#177}
\label{prob:177}

\noindent\textbf{Status:} OPEN \quad|\quad \textbf{Prize:} none \quad|\quad \textbf{Updated:} 2025-08-31

\noindent\textbf{Tags:} discrepancy, arithmetic progressions

\medskip

\noindent\textbf{Statement:}

Find the smallest $h(d)$ such that the following holds. There exists a function $f:\mathbb{N}\to\{-1,1\}$ such that, for every $d\geq 1$,\[\max_{P_d}\left\lvert \sum_{n\in P_d}f(n)\right\rvert\leq h(d),\]where $P_d$ ranges over all finite arithmetic progressions with common difference $d$.




Cantor, Erd\H{o}s, Schreiber, and Straus \cite{Er66} proved that $h(d)\ll d!$ is possible. Van der Waerden's theorem implies that $h(d)\to \infty$. Beck \cite{Be17} has shown that $h(d) \leq d^{8+\epsilon}$ is possible for every $\epsilon&#62;0$. Roth's famous discrepancy lower bound \cite{Ro64} implies that $h(d)\gg d^{1/2}$.




References

[Be17] Beck, J\'{o}zsef, A discrepancy problem: balancing infinite dimensional vectors . Number theory-Diophantine problems, uniform distribution and applications (2017), 61-82.

[Er66] Erd\H{o}s, P\'{a}l, Remarks on number theory. {V}. {E}xtremal problems in number theory. {II} . Mat. Lapok (1966), 135--155.

[Ro64] Roth, K. F., Remark concerning integer sequences . Acta Arith. (1964), 257-260.

\noindent\rule{\linewidth}{0.4pt}


%% ============================================================
\section{Problem \#180}
\label{prob:180}

\noindent\textbf{Status:} OPEN \quad|\quad \textbf{Prize:} none \quad|\quad \textbf{Updated:} 2025-08-31

\noindent\textbf{Tags:} graph theory, turan number

\medskip

\noindent\textbf{Statement:}

If $\mathcal{F}$ is a finite set of finite graphs then $\mathrm{ex}(n;\mathcal{F})$ is the maximum number of edges a graph on $n$ vertices can have without containing any subgraphs from $\mathcal{F}$. Note that it is trivial that $\mathrm{ex}(n;\mathcal{F})\leq \mathrm{ex}(n;G)$ for every $G\in\mathcal{F}$. Is it true that, for every $\mathcal{F}$, there exists $G\in\mathcal{F}$ such that\[\mathrm{ex}(n;G)\ll_{\mathcal{F}}\mathrm{ex}(n;\mathcal{F})?\]




A problem of Erd\H{o}s and Simonovits. This is trivially true if $\mathcal{F}$ does not contain any bipartite graphs, since by the Erd\H{o}s-Stone theorem if $H\in\mathcal{F}$ has minimal chromatic number $r\geq 2$ then\[\mathrm{ex}(n;H)=\mathrm{ex}(n;\mathcal{F})=\left(\frac{r-2}{r-1}+o(1)\right)\binom{n}{2}.\]Erd\H{o}s and Simonovits observe that this is false for infinite families $\mathcal{F}$, e.g. the family of all cycles. Hunter has provided the following 'folklore counterexample': if $\mathcal{F}=\{H_1,H_2\}$ where $H_1$ is a star and $H_2$ is a matching, both with at least two edges, then $\mathrm{ex}(n;\mathcal{F})\ll 1$, but $\mathrm{ex}(n;H_i)\asymp n$ for $1\leq i\leq 2$. This conjecture may still hold for all other $\mathcal{F}$. See also [575] . This problem is #47 in Extremal Graph Theory in the graphs problem collection.

\noindent\rule{\linewidth}{0.4pt}


%% ============================================================
\section{Problem \#181}
\label{prob:181}

\noindent\textbf{Status:} OPEN \quad|\quad \textbf{Prize:} none \quad|\quad \textbf{Updated:} 2025-08-31

\noindent\textbf{Tags:} graph theory, ramsey theory

\medskip

\noindent\textbf{Statement:}

Let $Q_n$ be the $n$-dimensional hypercube graph (so that $Q_n$ has $2^n$ vertices and $n2^{n-1}$ edges). Prove that\[R(Q_n) \ll 2^n.\]




Conjectured by Burr and Erd\H{o}s, althouhg in \cite{Er93} Erd\H{o}s says the behaviour of $R(Q_n)$ was considered by himself and S\'{o}s, who could not decide whether $R(Q_n)/2^n\to \infty$ or not. The trivial bound is\[R(Q_n) \leq R(K_{2^n})\leq C^{2^n}\]for some constant $C&gt;1$. This was improved a number of times; the current best bound due to Tikhomirov \cite{Ti22} is\[R(Q_n)\ll 2^{(2-c)n}\]for some small constant $c&gt;0$. (In fact $c\approx 0.03656$ is permissible.) This problem is #20 in Ramsey Theory in the graphs problem collection.




References

[Er93] Erd\H{o}s, Paul, Some of my favorite solved and unsolved problems in graph theory . Quaestiones Math. (1993), 333-350.

[Ti22] Tikhomirov, K., A remark on the Ramsey number of the hypercube . arXiv:2208.14568 (2022).

\noindent\rule{\linewidth}{0.4pt}


%% ============================================================
\section{Problem \#183}
\label{prob:183}

\noindent\textbf{Status:} OPEN \quad|\quad \textbf{Prize:} \$250 \quad|\quad \textbf{Updated:} 2025-08-31

\noindent\textbf{Tags:} graph theory, ramsey theory

\medskip

\noindent\textbf{Statement:}

Let $R(3;k)$ be the minimal $n$ such that if the edges of $K_n$ are coloured with $k$ colours then there must exist a monochromatic triangle. Determine\[\lim_{k\to \infty}R(3;k)^{1/k}.\]




Erd\H{o}s offers \$100 for showing that this limit is finite. An easy pigeonhole argument shows that\[R(3;k)\leq 2+k(R(3;k-1)-1),\]from which $R(3;k)\leq \lceil e k!\rceil$ immediately follows. The best-known upper bound is\[R(3;k)\leq (e-\tfrac{1}{6}) k!+1,\]due to Xu, Xie, and Chen \cite{XXC02} (improving previous bounds of Wan \cite{Wa97} and Whitehead \cite{Wh73}). See Eliahou \cite{El19} for more on the upper bound. These arise from the same type of inductive relationship and computational bounds for $R(3;k)$ for small $k$. The best-known lower bound (coming from lower bounds for Schur numbers) is\[R(3,k)\geq (380)^{k/5}-O(1),\]due to Ageron, Casteras, Pellerin, Portella, Rimmel, and Tomasik \cite{ACPPRT21} (improving previous bounds of Exoo \cite{Ex94} and Fredricksen and Sweet \cite{FrSw00}). Note that $380^{1/5}\approx 3.2806$. See also [483] . This problem is #21 in Ramsey Theory in the graphs problem collection.




References

[ACPPRT21] R. Ageron, P. Casteras, T. Pellerin, Y. Portella, A. Rimmel, and J. Tomasik, New lower bounds for Schur and weak Schur numbers . arXiv:2112.03175 (2021).

[El19] No reference found.


[Ex94] Exoo, G., A lower bound for Schur numbers and multicolor Ramsey numbers . Electronic J. of Combinatorics (1994).

[FrSw00] Fredricksen, Harold and Sweet, Melvin M., Symmetric sum-free partitions and lower bounds for Schur numbers . Electron. J. Combin. (2000), Research Paper 32, 9.

[Wa97] Wan, Honghui, Upper bounds for {R}amsey numbers {$R(3,3,\cdots,3)$} and Schur numbers . J. Graph Theory (1997), 119--122.

[Wh73] Whitehead, Jr., Earl Glen, The {R}amsey number {$N(3,\,3,\,3,\,3;\,2)$} . Discrete Math. (1973), 389--396.

[XXC02] Xu, Xiao Dong and Xie, Zheng and Chen, Zhi, Upper bounds for {R}amsey numbers {$R_n(3)$} and Schur numbers . Math. Econ. (2002), 81--84.

\noindent\rule{\linewidth}{0.4pt}


%% ============================================================
\section{Problem \#184}
\label{prob:184}

\noindent\textbf{Status:} OPEN \quad|\quad \textbf{Prize:} none \quad|\quad \textbf{Updated:} 2025-08-31

\noindent\textbf{Tags:} graph theory, cycles

\medskip

\noindent\textbf{Statement:}

Any graph on $n$ vertices can be decomposed into $O(n)$ many edge-disjoint cycles and edges.




Conjectured by Erd\H{o}s and Gallai, who proved that $O(n\log n)$ many cycles and edges suffices (see Section 5 of \cite{EGP66}). The graph $K_{3,n-3}$ shows that at least $(1+c)n$ many cycles and edges are required, for some constant $c&gt;0$. In \cite{Er71} Erd\H{o}s suggests that only $n-1$ many cycles and edges are required if we do not require them to be edge-disjoint. The best bound available is due to Buci\'{c} and Montgomery \cite{BM22}, who prove that $O(n\log^*n)$ many cycles and edges suffice, where $\log^*$ is the iterated logarithm function. Conlon, Fox, and Sudakov \cite{CFS14} proved that $O_\epsilon(n)$ cycles and edges suffice if $G$ has minimum degree at least $\epsilon n$, for any $\epsilon&gt;0$. See also [583] for an analogous problem decomposing into paths, and [1017] for decomposing into complete graphs.




References

[BM22] Buci\'C, M. and Montgomery, R., Towards the Erd\H{o}s-Gallai Cycle Decomposition Conjecture . arXiv:2211.07689 (2022).

[CFS14] Conlon, David and Fox, Jacob and Sudakov, Benny, Cycle packing . Random Structures Algorithms (2014), 608-626.

[EGP66] Erd\H{o}s, Paul and Goodman, A. W. and P\'{o}sa, Lajos, The representation of a graph by set intersections . Canadian J. Math. (1966), 106-112.

[Er71] Erd\H{o}s, P., Some unsolved problems in graph theory and combinatorial analysis . Combinatorial Mathematics and its Applications (Proc. Conf., Oxford, 1969) (1971), 97-109.

\noindent\rule{\linewidth}{0.4pt}


%% ============================================================
\section{Problem \#187}
\label{prob:187}

\noindent\textbf{Status:} OPEN \quad|\quad \textbf{Prize:} none \quad|\quad \textbf{Updated:} 2025-08-31

\noindent\textbf{Tags:} additive combinatorics, ramsey theory

\medskip

\noindent\textbf{Statement:}

Find the best function $f(d)$ such that, in any 2-colouring of the integers, at least one colour class contains an arithmetic progression with common difference $d$ of length $f(d)$ for infinitely many $d$.




Originally asked by Cohen. Erd\H{o}s observed that colouring according to whether $\{ \sqrt{2}n\}&#60;1/2$ or not implies $f(d) \ll d$ (using the fact that $\|\sqrt{2}q\| \gg 1/q$ for all $q$, where $\|x\|$ is the distance to the nearest integer). In \cite{Er80} he states that Petruska and Szemer\'{e}di have proved that $f(d) \ll d^{1/2}$, and expected $f(d)\leq d^{o(1)}$. Beck \cite{Be80} has improved this using the probabilistic method, constructing a colouring that shows $f(d)\leq (1+o(1))\log_2 d$. Van der Waerden's theorem implies $f(d)\to \infty$ is necessary.




References

[Be80] Beck, J\'{o}zsef, A remark concerning arithmetic progressions . J. Combin. Theory Ser. A (1980), 376-379.

[Er80] Erd\H{o}s, Paul, A survey of problems in combinatorial number theory . Ann. Discrete Math. (1980), 89-115.

\noindent\rule{\linewidth}{0.4pt}


%% ============================================================
\section{Problem \#188}
\label{prob:188}

\noindent\textbf{Status:} OPEN \quad|\quad \textbf{Prize:} none \quad|\quad \textbf{Updated:} 2025-08-31

\noindent\textbf{Tags:} geometry, ramsey theory

\medskip

\noindent\textbf{Statement:}

What is the smallest $k$ such that $\mathbb{R}^2$ can be red/blue coloured with no pair of red points unit distance apart, and no $k$-term arithmetic progression of blue points with distance $1$?




Erd\H{o}s, Graham, Montgomery, Rothschild, Spencer, and Straus \cite{EGMRSS75} proved $k\geq 5$. Tsaturian \cite{Ts17} improved this to $k\geq 6$. Erd\H{o}s and Graham claim that $k\leq 10000000$ ('more or less'), but give no proof. Erd\H{o}s and Graham asked this with just any $k$-term arithmetic progression in blue (not necessarily with distance $1$), but Alon has pointed out that in fact no such $k$ exists: in any red/blue colouring of the integer points on a line either there are two red points distance $1$ apart, or else the set of blue points and the same set shifted by $1$ cover all integers, and hence by van der Waerden's theorem there are arbitrarily long blue arithmetic progressions. It seems most likely, from context, that Erd\H{o}s and Graham intended to restrict the blue arithmetic progression to have distance $1$ (although they do not write this restriction in their papers).




References

[EGMRSS75] Erd\H{o}s, P. and Graham, R. L. and Montgomery, P. and Rothschild, B. L. and Spencer, J. and Straus, E. G., Euclidean {R}amsey theorems. {II} . (1975), 529--557.

[Ts17] Tsaturian, Sergei, A {E}uclidean {R}amsey result in the plane . Electron. J. Combin. (2017), Paper No. 4.35, 9.

\noindent\rule{\linewidth}{0.4pt}


%% ============================================================
\section{Problem \#193}
\label{prob:193}

\noindent\textbf{Status:} OPEN \quad|\quad \textbf{Prize:} none \quad|\quad \textbf{Updated:} 2025-08-31

\noindent\textbf{Tags:} geometry

\medskip

\noindent\textbf{Statement:}

Let $S\subseteq \mathbb{Z}^3$ be a finite set and let $A=\{a_1,a_2,\ldots,\}\subset \mathbb{Z}^3$ be an infinite $S$-walk, so that $a_{i+1}-a_i\in S$ for all $i$. Must $A$ contain three collinear points?




Originally conjectured by Gerver and Ramsey \cite{GeRa79}, who showed that the answer is yes for $\mathbb{Z}^2$, and for $\mathbb{Z}^3$ that the largest number of collinear points can be bounded.




References

[GeRa79] Gerver, Joseph L. and Ramsey, L. Thomas, On certain sequences of lattice points . Pacific J. Math. (1979), 357-363.

\noindent\rule{\linewidth}{0.4pt}


%% ============================================================
\section{Problem \#195}
\label{prob:195}

\noindent\textbf{Status:} OPEN \quad|\quad \textbf{Prize:} none \quad|\quad \textbf{Updated:} 2025-08-31

\noindent\textbf{Tags:} arithmetic progressions

\medskip

\noindent\textbf{Statement:}

What is the largest $k$ such that in any permutation of $\mathbb{Z}$ there must exist a monotone $k$-term arithmetic progression $x_1&#60;\cdots&#60;x_k$?




Geneson \cite{Ge19} proved that $k\leq 5$. Adenwalla \cite{Ad22} proved that $k\leq 4$. See also [194] and [196] .




References

[Ad22] Adenwalla, S., Avoiding Monotone Arithmetic Progressions in Permutations of Integers . arXiv:2211.04451 (2022).

[Ge19] Geneson, Jesse, Forbidden arithmetic progressions in permutations of subsets of the integers . Discrete Math. (2019), 1489-1491.

\noindent\rule{\linewidth}{0.4pt}


%% ============================================================
\section{Problem \#196}
\label{prob:196}

\noindent\textbf{Status:} OPEN \quad|\quad \textbf{Prize:} none \quad|\quad \textbf{Updated:} 2025-08-31

\noindent\textbf{Tags:} arithmetic progressions

\medskip

\noindent\textbf{Statement:}

Must every permutation of $\mathbb{N}$ contain a monotone 4-term arithmetic progression? In other words, given a permutation $x$ of $\mathbb{N}$ must there be indices with either $i&lt;j&lt;k&lt;l$ or $i&gt;j&#62;k&#62;l$ such that $x_i,x_j,x_k,x_l$ are an arithmetic progression?




Davis, Entringer, Graham, and Simmons \cite{DEGS77} have shown that there must exist a monotone 3-term arithmetic progression and need not contain a 5-term arithmetic progression. See also [194] and [195] .




References

[DEGS77] Davis, J. A. and Entringer, R. C. and Graham, R. L. and Simmons, G. J., On permutations containing no long arithmetic progressions . Acta Arith. (1977/78), 81-90.

\noindent\rule{\linewidth}{0.4pt}


%% ============================================================
\section{Problem \#197}
\label{prob:197}

\noindent\textbf{Status:} OPEN \quad|\quad \textbf{Prize:} none \quad|\quad \textbf{Updated:} 2025-08-31

\noindent\textbf{Tags:} arithmetic progressions

\medskip

\noindent\textbf{Statement:}

Can $\mathbb{N}$ be partitioned into two sets, each of which can be permuted to avoid monotone 3-term arithmetic progressions?




If three sets are allowed then this is possible.

\noindent\rule{\linewidth}{0.4pt}


%% ============================================================
\section{Problem \#200}
\label{prob:200}

\noindent\textbf{Status:} OPEN \quad|\quad \textbf{Prize:} none \quad|\quad \textbf{Updated:} 2025-08-31

\noindent\textbf{Tags:} primes, arithmetic progressions

\medskip

\noindent\textbf{Statement:}

Does the longest arithmetic progression of primes in $\{1,\ldots,N\}$ have length $o(\log N)$?




It follows from the prime number theorem that such a progression has length $\leq(1+o(1))\log N$.

\noindent\rule{\linewidth}{0.4pt}


%% ============================================================
\section{Problem \#201}
\label{prob:201}

\noindent\textbf{Status:} OPEN \quad|\quad \textbf{Prize:} none \quad|\quad \textbf{Updated:} 2025-08-31

\noindent\textbf{Tags:} additive combinatorics, arithmetic progressions

\medskip

\noindent\textbf{Statement:}

Let $G_k(N)$ be such that any set of $N$ integers contains a subset of size at least $G_k(N)$ which does not contain a $k$-term arithmetic progression. Determine the size of $G_k(N)$. How does it relate to $R_k(N)$, the size of the largest subset of $\{1,\ldots,N\}$ without a $k$-term arithmetic progression? Is it true that\[\lim_{N\to \infty}\frac{R_3(N)}{G_3(N)}=1?\]




First asked and investigated by Riddell \cite{Ri69}. It is trivial that $G_k(N)\leq R_k(N)$, and it is possible that $G_k(N) &#60;R_k(N)$ (for example $G_3(5)=3$ and $R_3(5)=4$, and $G_3(14)\leq 7$ and $R_3(14)=8$). Koml\'{o}s, Sulyok, and Szemer\'{e}di \cite{KSS75} have shown that $R_k(N) \ll_k G_k(N)$.




References

[KSS75] Koml\'{o}s, J. and Sulyok, M. and Szemeredi, E., Linear problems in combinatorial number theory . Acta Math. Acad. Sci. Hungar. (1975), 113-121.

[Ri69] Riddell, J., On sets of numbers containing no $l$ terms in arithmetic progression . Nieuw Arch. Wisk. (3) (1969), 204-209.

\noindent\rule{\linewidth}{0.4pt}


%% ============================================================
\section{Problem \#203}
\label{prob:203}

\noindent\textbf{Status:} OPEN \quad|\quad \textbf{Prize:} none \quad|\quad \textbf{Updated:} 2025-08-31

\noindent\textbf{Tags:} primes, covering systems

\medskip

\noindent\textbf{Statement:}

Is there an integer $m\geq 1$ with $(m,6)=1$ such that none of $2^k3^\ell m+1$ are prime, for any $k,\ell\geq 0$?




Positive odd integers $m$ such that none of $2^km+1$ are prime are called Sierpinski numbers - see [1113] for more details. Erd\H{o}s and Graham also ask more generally about $p_1^{k_1}\cdots p_r^{k_r}m+1$ for distinct primes $p_i$, or $q_1\cdots q_rm+1$ where the $q_i$ are primes congruent to $1\pmod{4}$. (Dogmachine has noted in the comments this latter question has the trivial answer $m=1$ - perhaps some condition such as $m$ even is meant.)

\noindent\rule{\linewidth}{0.4pt}


%% ============================================================
\section{Problem \#208}
\label{prob:208}

\noindent\textbf{Status:} OPEN \quad|\quad \textbf{Prize:} none \quad|\quad \textbf{Updated:} 2025-08-31

\noindent\textbf{Tags:} number theory

\medskip

\noindent\textbf{Statement:}

Let $s_1&lt;s_2&lt;\cdots$ be the sequence of squarefree numbers. Is it true that, for any $\epsilon&gt;0$ and large $n$,\[s_{n+1}-s_n \ll_\epsilon s_n^{\epsilon}?\]Is it true that\[s_{n+1}-s_n \leq (1+o(1))\frac{\pi^2}{6}\frac{\log s_n}{\log\log s_n}?\]




Erd\H{o}s \cite{Er51} showed that there are infinitely many $n$ such that\[s_{n+1}-s_n &gt; (1+o(1))\frac{\pi^2}{6}\frac{\log s_n}{\log\log s_n},\]so this bound would be the best possible. In \cite{Er79} Erd\H{o}s says perhaps $s_{n+1}-s_n \ll \log s_n$, but he is 'very doubtful'. Filaseta and Trifonov \cite{FiTr92} proved an upper bound of $s_n^{1/5+o(1)}$. Pandey \cite{Pa24} has improved this exponent to $1/5-c$ for some constant $c&gt;0$. Granville \cite{Gr98} showed that $s_{n+1}-s_n\ll_\epsilon s_n^\epsilon$ for all $\epsilon&gt;0$ follows from the ABC conjecture . See also [489] and [145] . A more general form of this problem is given in [1101] .




References

[Er51] Erd\H{o}s, P., Some problems and results in elementary number theory . Publ. Math. Debrecen (1951), 103-109.

[Er79] Erd\H{o}s, Paul, Some unconventional problems in number theory . Math. Mag. (1979), 67-70.

[FiTr92] Filaseta, M. and Trifonov, O., On gaps between squarefree numbers II . J. London Math. Soc. (1992), 215-221.

[Gr98] Granville, Andrew, {$ABC$} allows us to count squarefrees . Internat. Math. Res. Notices (1998), 991--1009.

[Pa24] Pandey, M., Squarefree numbers in short intervals . arXiv:2401.13981 (2024).

\noindent\rule{\linewidth}{0.4pt}


%% ============================================================
\section{Problem \#212}
\label{prob:212}

\noindent\textbf{Status:} OPEN \quad|\quad \textbf{Prize:} none \quad|\quad \textbf{Updated:} 2025-08-31

\noindent\textbf{Tags:} geometry, distances

\medskip

\noindent\textbf{Statement:}

Is there a dense subset of $\mathbb{R}^2$ such that all pairwise distances are rational?




Conjectured by Ulam. Erd\H{o}s believed there cannot be such a set. This problem is discussed in a blogpost by Terence Tao, in which he shows that there cannot be such a set, assuming the Bombieri-Lang conjecture. The same conclusion was independently obtained by Shaffaf \cite{Sh18}. Indeed, Shaffaf and Tao actually proved that such a rational distance set must be contained in a finite union of real algebraic curves. Solymosi and de Zeeuw \cite{SdZ10} then proved (unconditionally) that a rational distance set contained in a real algebraic curve must be finite, unless the curve contains a line or a circle. Ascher, Braune, and Turchet \cite{ABT20} observed that, combined, these facts imply that a rational distance set in general position must be finite (conditional on the Bombieri-Lang conjecture). In \cite{Er87b} Erd\H{o}s mentions that Besicovitch conjectured that the limit points of a rational distance set cannot contain arbitrarily large convex sets.




References

[ABT20] Ascher, K. and Braune, L. and Turchet, A., The Erd\H{o}s-Ulam problem, Lang's conjecture, and uniformity . arXiv:1901.02616 (2020).

[Er87b] Erd\H{o}s, P., Some combinatorial and metric problems in geometry . Intuitive geometry (Si\'{o}fok, 1985) (1987), 167-177.

[SdZ10] Solymosi, Jozsef and de Zeeuw, Frank, On a question of Erd\H{o}s and Ulam . Discrete Comput. Geom. (2010), 393-401.

[Sh18] Shaffaf, Jafar, A solution of the Erd\H{o}s-Ulam problem on rational distance sets assuming the Bombieri-Lang conjecture . Discrete Comput. Geom. (2018), 283-293.

\noindent\rule{\linewidth}{0.4pt}


%% ============================================================
\section{Problem \#213}
\label{prob:213}

\noindent\textbf{Status:} OPEN \quad|\quad \textbf{Prize:} none \quad|\quad \textbf{Updated:} 2025-08-31

\noindent\textbf{Tags:} geometry, distances

\medskip

\noindent\textbf{Statement:}

Let $n\geq 4$. Are there $n$ points in $\mathbb{R}^2$, no three on a line and no four on a circle, such that all pairwise distances are integers?




Anning and Erd\H{o}s \cite{AnEr45} proved there cannot exist an infinite such set. Harborth constructed such a set when $n=5$. The best construction to date, due to Kreisel and Kurz \cite{KK08}, has $n=7$. Ascher, Braune, and Turchet \cite{ABT20} have shown that there is a uniform upper bound on the size of such a set, conditional on the Bombieri-Lang conjecture. Greenfeld, Iliopoulou, and Peluse \cite{GIP24} have shown (unconditionally) that any such set must be very sparse, in that if $S\subseteq [-N,N]^2$ has no three on a line and no four on a circle, and all pairwise distances integers, then\[\lvert S\rvert \ll (\log N)^{O(1)}.\]See also [130] .




References

[ABT20] Ascher, K. and Braune, L. and Turchet, A., The Erd\H{o}s-Ulam problem, Lang's conjecture, and uniformity . arXiv:1901.02616 (2020).

[AnEr45] Anning, Norman H. and Erd\H{o}s, Paul, Integral distances . Bull. Amer. Math. Soc. (1945), 598-600.

[GIP24] Greenfeld, R. and Iliopoulou, M. and Peluse, S., On integer distance sets . arXiv:2401.10821 (2024).

[KK08] Kreisel, Tobias and Kurz, Sascha, There are integral heptagons, no three points on a line, on four on a circle . Discrete Comput. Geom. (2008), 786-790.

\noindent\rule{\linewidth}{0.4pt}


%% ============================================================
\section{Problem \#217}
\label{prob:217}

\noindent\textbf{Status:} OPEN \quad|\quad \textbf{Prize:} none \quad|\quad \textbf{Updated:} 2025-08-31

\noindent\textbf{Tags:} geometry, distances

\medskip

\noindent\textbf{Statement:}

For which $n$ are there $n$ points in $\mathbb{R}^2$, no three on a line and no four on a circle, which determine $n-1$ distinct distances and so that (in some ordering of the distances) the $i$th distance occurs $i$ times?




An example with $n=4$ is an isosceles triangle with the point in the centre. Erd\H{o}s originally believed this was impossible for $n\geq 5$, but Pomerance constructed a set with $n=5$ (see \cite{Er83c} for a description). Pal\'{a}sti \cite{Pa87} constructed such a set with $n=7$. According to \cite{Pa89} when Erd\H{o}s saw this he asked whether there exists such a set of $5$ points with no equilateral triangles; Pal\'{a}sti \cite{Pa89} {IMAGE=217a,constructed a set} of $n=6$ points, satisfying these restrictions, with no equilateral triangles. Pal\'{a}sti \cite{Pa89b} has {IMAGE=217b,constructed a set} with $n=8$. Erd\H{o}s believed this is impossible for all sufficiently large $n$. This would follow from $h(n)\geq n$ for sufficiently large $n$, where $h(n)$ is as in [98] .




References

[Er83c] Erd\H{o}s, Paul, Combinatorial problems in geometry . Math. Chronicle (1983), 35-54.

[Pa87] Pal\'{a}sti, Ilona, On the seven points problem of {P}. {E}rd\H os . Studia Sci. Math. Hungar. (1987), 447--448.

[Pa89] Pal\'{a}sti, Ilona, A distance problem of {P}. {E}rd\H os with some further restrictions . Discrete Math. (1989), 155--156.

[Pa89b] Pal\'{a}sti, I., Lattice-point examples for a question of {E}rd\H os . Period. Math. Hungar. (1989), 231--235.

\noindent\rule{\linewidth}{0.4pt}


%% ============================================================
\section{Problem \#218}
\label{prob:218}

\noindent\textbf{Status:} OPEN \quad|\quad \textbf{Prize:} none \quad|\quad \textbf{Updated:} 2025-08-31

\noindent\textbf{Tags:} number theory, primes

\medskip

\noindent\textbf{Statement:}

Let $d_n=p_{n+1}-p_n$. The set of $n$ such that $d_{n+1}\geq d_n$ has density $1/2$, and similarly for $d_{n+1}\leq d_n$. Furthermore, there are infinitely many $n$ such that $d_{n+1}=d_n$.




In \cite{Er85c} Erd\H{o}s also conjectures that $d_n=d_{n+1}=\cdots=d_{n+k}$ is solvable for every $k$ (which is equivalent to $k$ consecutive primes in arithmetic progression, see [141] ). Banks \cite{Ba23} has given a heuristic argument, assuming a quantitative form of the prime tuples conjecture, that the first statement is true, and in fact for any $c\geq 0$ and $\epsilon&gt;0$ the number of $n$ such that $p_n\leq x$ and $d_{n+1}\geq cd_n$ is\[\frac{\pi(x)}{c+1}+O\left(\frac{x}{(\log x)^{3/2+\epsilon}}\right).\]




References

[Ba23] Banks, William D., On ratios of consecutive prime gaps . Integers (2023), Paper No. A50, 13.

[Er85c] Erd\H{o}s, P., On some of my problems in number theory I would most like to see solved . Number theory (Ootacamund, 1984) (1985), 74-84.

\noindent\rule{\linewidth}{0.4pt}


%% ============================================================
\section{Problem \#222}
\label{prob:222}

\noindent\textbf{Status:} OPEN \quad|\quad \textbf{Prize:} none \quad|\quad \textbf{Updated:} 2025-08-31

\noindent\textbf{Tags:} number theory, squares

\medskip

\noindent\textbf{Statement:}

Let $n_1&#60;n_2&#60;\cdots$ be the sequence of integers which are the sum of two squares. Explore the behaviour of (i.e. find good upper and lower bounds for) the consecutive differences $n_{k+1}-n_k$.




Erd\H{o}s \cite{Er51} proved that, for infinitely many $k$,\[ n_{k+1}-n_k \gg \frac{\log n_k}{\sqrt{\log\log n_k}}.\]Richards \cite{Ri82} improved this to\[\limsup_{k\to \infty} \frac{n_{k+1}-n_k}{\log n_k} \geq 1/4.\]The constant $1/4$ here has been improved, most lately to $0.868\cdots$ by Dietmann, Elsholtz, Kalmynin, Konyagin, and Maynard \cite{DEKKM22}. The best known upper bound is due to Bambah and Chowla \cite{BaCh47}, who proved that\[n_{k+1}-n_k \ll n_k^{1/4}.\]The differences are listed at A256435 on the OEIS.




References

[BaCh47] Bambah, R. P. and Chowla, S., On numbers which can be expressed as a sum of two squares . Proc. Nat. Inst. Sci. India (1947), 101-103.

[DEKKM22] Dietmann, R. and Elsholtz, C. and Kaymynin, A. and Konyagin, S. and Maynard, J., Longer Gaps Between Values of Binary Quadratic Forms . International Mathematics Research Notices (2023), 10313–10349.

[Er51] Erd\H{o}s, P., Some problems and results in elementary number theory . Publ. Math. Debrecen (1951), 103-109.

[Ri82] Richards, Ian, On the gaps between numbers which are sums of two squares . Adv. in Math. (1982), 1-2.

\noindent\rule{\linewidth}{0.4pt}


%% ============================================================
\section{Problem \#233}
\label{prob:233}

\noindent\textbf{Status:} OPEN \quad|\quad \textbf{Prize:} none \quad|\quad \textbf{Updated:} 2025-08-31

\noindent\textbf{Tags:} number theory, primes

\medskip

\noindent\textbf{Statement:}

Let $d_n=p_{n+1}-p_n$, where $p_n$ is the $n$th prime. Prove that\[\sum_{1\leq n\leq N}d_n^2 \ll N(\log N)^2.\]




Cramer \cite{Cr36} proved an upper bound of $O(N(\log N)^4)$ conditional on the Riemann hypothesis. Selberg \cite{Se43} improved this slightly (still assuming the Riemann hypothesis) to\[\sum_{1\leq n\leq N}\frac{d_n^2}{n}\ll (\log N)^4.\]The prime number theorem immediately implies a lower bound of\[\sum_{1\leq n\leq N}d_n^2\gg N(\log N)^2.\]This would imply in particular that $d_n\ll n^{1/2}\log n$ for all $n$, which is known only the assumption of the Riemann Hypothesis. The values of the sum are listed at A074741 on the OEIS. This is discussed in problem A8 of Guy's collection \cite{Gu04}.




References

[Cr36] Cram\'{e}r, Harald, On the order of magnitude of the difference between consecutive prime numbers . Acta Arithmetica (1936), 23--46.

[Gu04] Guy, Richard K., Unsolved problems in number theory . (2004), xviii+437.

[Se43] Selberg, Atle, On the normal density of primes in small intervals, and the difference between consecutive primes . Arch. Math. Naturvid. (1943), 87--105.

\noindent\rule{\linewidth}{0.4pt}


%% ============================================================
\section{Problem \#234}
\label{prob:234}

\noindent\textbf{Status:} OPEN \quad|\quad \textbf{Prize:} none \quad|\quad \textbf{Updated:} 2025-08-31

\noindent\textbf{Tags:} number theory, primes

\medskip

\noindent\textbf{Statement:}

For every $c\geq 0$ the density $f(c)$ of integers for which\[\frac{p_{n+1}-p_n}{\log n}&#60; c\]exists and is a continuous function of $c$.




See also [5] .

\noindent\rule{\linewidth}{0.4pt}


%% ============================================================
\section{Problem \#236}
\label{prob:236}

\noindent\textbf{Status:} OPEN \quad|\quad \textbf{Prize:} none \quad|\quad \textbf{Updated:} 2025-08-31

\noindent\textbf{Tags:} number theory, primes

\medskip

\noindent\textbf{Statement:}

Let $f(n)$ count the number of solutions to $n=p+2^k$ for prime $p$ and $k\geq 0$. Is it true that $f(n)=o(\log n)$?




Erd\H{o}s \cite{Er50} proved that there are infinitely many $n$ such that $f(n)\gg \log\log n$. Erd\H{o}s could not even prove that there do not exist infinitely many integers $n$ such that for all $1&lt; 2^k&lt;n$ the number $n-2^k$ is prime - see [1142] . The sequence of values of $f(n)$ is A109925 on the OEIS. See also [237] .




References

[Er50] Erd\H{o}s, P., On integers of the form $2^k+p$ and some related problems . Summa Brasil. Math. (1950), 113-123.

\noindent\rule{\linewidth}{0.4pt}


%% ============================================================
\section{Problem \#238}
\label{prob:238}

\noindent\textbf{Status:} OPEN \quad|\quad \textbf{Prize:} none \quad|\quad \textbf{Updated:} 2025-08-31

\noindent\textbf{Tags:} number theory, primes

\medskip

\noindent\textbf{Statement:}

Let $c_1,c_2&#62;0$. Is it true that, for any sufficiently large $x$, there exist more than $c_1\log x$ many consecutive primes $\leq x$ such that the difference between any two is $&#62;c_2$?




Erd\H{o}s \cite{Er49c} proved this is true for any $c_2&#62;0$ if $c_1&#62;0$ is sufficiently small (depending on $c_1$).




References

[Er49c] Erd\H{o}s, P., On some applications of {B}run's method . Acta Univ. Szeged. Sect. Sci. Math. (1949), 57--63.

\noindent\rule{\linewidth}{0.4pt}


%% ============================================================
\section{Problem \#241}
\label{prob:241}

\noindent\textbf{Status:} OPEN \quad|\quad \textbf{Prize:} \$100 \quad|\quad \textbf{Updated:} 2025-08-31

\noindent\textbf{Tags:} additive combinatorics, sidon sets

\medskip

\noindent\textbf{Statement:}

Let $f(N)$ be the maximum size of $A\subseteq \{1,\ldots,N\}$ such that the sums $a+b+c$ with $a,b,c\in A$ are all distinct (aside from the trivial coincidences). Is it true that\[ f(N)\sim N^{1/3}?\]




Originally asked to Erd\H{o}s by Bose. Bose and Chowla \cite{BoCh62} provided a construction proving one half of this, namely\[(1+o(1))N^{1/3}\leq f(N).\]The best upper bound known to date is due to Green \cite{Gr01},\[f(N) \leq ((7/2)^{1/3}+o(1))N^{1/3}\](note that $(7/2)^{1/3}\approx 1.519$). More generally, Bose and Chowla conjectured that the maximum size of $A\subseteq \{1,\ldots,N\}$ with all $r$-fold sums distinct (aside from the trivial coincidences) then\[\lvert A\rvert \sim N^{1/r}.\]This is known only for $r=2$ (see [30] ). This is discussed in problem C11 of Guy's collection \cite{Gu04}.




References

[BoCh62] Bose, R. C. and Chowla, S., Theorems in the additive theory of numbers . Comment. Math. Helv. (1962/63), 141-147.

[Gr01] Green, Ben, The number of squares and {$B_h[g]$} sets . Acta Arith. (2001), 365-390.

[Gu04] Guy, Richard K., Unsolved problems in number theory . (2004), xviii+437.

\noindent\rule{\linewidth}{0.4pt}


%% ============================================================
\section{Problem \#243}
\label{prob:243}

\noindent\textbf{Status:} OPEN \quad|\quad \textbf{Prize:} none \quad|\quad \textbf{Updated:} 2025-08-31

\noindent\textbf{Tags:} number theory

\medskip

\noindent\textbf{Statement:}

Let $1\leq a_1&#60;a_2&#60;\cdots$ be a sequence of integers such that\[\lim_{n\to \infty}\frac{a_n}{a_{n-1}^2}=1\]and $\sum\frac{1}{a_n}\in \mathbb{Q}$. Then, for all sufficiently large $n\geq 1$,\[ a_n = a_{n-1}^2-a_{n-1}+1.\]




Erd\H{o}s and Straus \cite{ErSt64} proved that if $\lim a_n/a_{n-1}^2=1$ and $\sum \frac{1}{a_n}$ is rational, and $a_n$ does not satisfy the recurrence, then\[\limsup_{n\to \infty} \frac{[a_1,\ldots,a_n]}{a_{n+1}}\left(\frac{a_n^2}{a_{n+1}}-1\right)&gt;0.\]A sequence satisfying the reucrrence $a_n = a_{n-1}^2-a_{n-1}+1$ is known as Sylvester's sequence . Duverney \cite{Du01} proved a weaker version of this problem: if\[\sum_{n\geq 0}\left(\frac{a_{n+1}}{a_n^2}-1\right)\]converges then $\sum \frac{1}{a_n}$ is rational if and only if\[a_{n}=a_{n-1}^2-a_{n-1}+1\]for all large $n$.




References

[Du01] Duverney, Daniel, Irrationality of fast converging series of rational numbers . J. Math. Sci. Univ. Tokyo (2001), 275--316.

[ErSt64] Erd\H{o}s, P. and Straus, E. G., On the irrationality of certain {A}hmes series . J. Indian Math. Soc. (N.S.) (1964), 129--133.

\noindent\rule{\linewidth}{0.4pt}


%% ============================================================
\section{Problem \#244}
\label{prob:244}

\noindent\textbf{Status:} OPEN \quad|\quad \textbf{Prize:} none \quad|\quad \textbf{Updated:} 2025-08-31

\noindent\textbf{Tags:} number theory, primes

\medskip

\noindent\textbf{Statement:}

Let $C&#62;1$. Does the set of integers of the form $p+\lfloor C^k\rfloor$, for some prime $p$ and $k\geq 0$, have density $&#62;0$?




Originally asked to Erd\H{o}s by Kalm\'{a}r. Erd\H{o}s believed the answer is yes. Romanoff \cite{Ro34} proved that the answer is yes if $C$ is an integer. Ding \cite{Di25} has proved that this is true for almost all $C&#62;1$.




References

[Di25] Y. Ding, On a Romanoff type problem of Erd\H{o}s and Kalm\'{a}r . arXiv:2503.22700 (2025).

[Ro34] Romanoff, N. P., \"{U}ber einige S\"Atze der additiven Zahlentheorie . Math. Ann. (1934), 668-678.

\noindent\rule{\linewidth}{0.4pt}


%% ============================================================
\section{Problem \#247}
\label{prob:247}

\noindent\textbf{Status:} OPEN \quad|\quad \textbf{Prize:} none \quad|\quad \textbf{Updated:} 2025-08-31

\noindent\textbf{Tags:} number theory, irrationality

\medskip

\noindent\textbf{Statement:}

Let $1\leq a_1&#60;a_2&#60;\cdots$ be a sequence of integers such that\[\limsup \frac{a_n}{n}=\infty.\]Is\[\sum_{n=1}^\infty \frac{1}{2^{a_n}}\]transcendental?




Erd\H{o}s \cite{Er75c} proved the answer is yes under the stronger condition that $\limsup n_k/k^t=\infty$ for all $t\geq 1$. Erd\H{o}s \cite{Er88c} says 'many of these problems seem hopeless at present, but perhaps one can prove that if $a_n&gt;cn^2$ then $\sum_{n=1}^\infty \frac{1}{2^{a_n}}$ is not the root of any quadratic polynomial'. This problem has been formalised in Lean as part of the Google DeepMind Formal Conjectures project .




References

[Er75c] Erd\H{o}s, P., Some problems and results on the irrationality of the sum of infinite series . J. Math. Sci. (1975), 1-7 (1976).

[Er88c] Erd\H{o}s, P., On the irrationality of certain series: problems and results . New advances in transcendence theory (Durham, 1986) (1988), 102-109.

\noindent\rule{\linewidth}{0.4pt}


%% ============================================================
\section{Problem \#249}
\label{prob:249}

\noindent\textbf{Status:} OPEN \quad|\quad \textbf{Prize:} none \quad|\quad \textbf{Updated:} 2025-08-31

\noindent\textbf{Tags:} number theory, irrationality

\medskip

\noindent\textbf{Statement:}

Is\[\sum_n \frac{\phi(n)}{2^n}\]irrational? Here $\phi$ is the Euler totient function .




The decimal expansion of this sum is A256936 on the OEIS.

\noindent\rule{\linewidth}{0.4pt}


%% ============================================================
\section{Problem \#251}
\label{prob:251}

\noindent\textbf{Status:} OPEN \quad|\quad \textbf{Prize:} none \quad|\quad \textbf{Updated:} 2025-08-31

\noindent\textbf{Tags:} number theory, irrationality

\medskip

\noindent\textbf{Statement:}

Is\[\sum \frac{p_n}{2^n}\]irrational? (Here $p_n$ is the $n$th prime.)




Erd\H{o}s \cite{Er58b} proved that $\sum \frac{p_n^k}{n!}$ is irrational for every $k\geq 1$. In \cite{Er88c} he further conjectures that $\sum \frac{p_n^k}{2^n}$ is irrational for every $k$, and that if $g_n\geq 2$ and $g_n=o(p_n)$ then\[\sum_{n=1}^\infty \frac{p_n}{g_1\cdots g_n}\]is irrational. (The example $g_n=p_n+1$ shows that some condition on the growth of the $g_n$ is necessary here.) The decimal expansion of this sum is A098990 on the OEIS.




References

[Er58b] Erd\H{o}s, Paul, Sur certaines s\'{e}ries \`a{} valeur irrationnelle . Enseign. Math. (2) (1958), 93--100.

[Er88c] Erd\H{o}s, P., On the irrationality of certain series: problems and results . New advances in transcendence theory (Durham, 1986) (1988), 102-109.

\noindent\rule{\linewidth}{0.4pt}


%% ============================================================
\section{Problem \#252}
\label{prob:252}

\noindent\textbf{Status:} OPEN \quad|\quad \textbf{Prize:} none \quad|\quad \textbf{Updated:} 2025-08-31

\noindent\textbf{Tags:} number theory, irrationality

\medskip

\noindent\textbf{Statement:}

Let $k\geq 1$ and $\sigma_k(n)=\sum_{d\mid n}d^k$. Is\[\sum \frac{\sigma_k(n)}{n!}\]irrational?




This is known now for $1\leq k\leq 4$. The cases $k=1,2$ are reasonably straightforward, as observed by Erd\H{o}s \cite{Er52}. The case $k=3$ was proved independently by Schlage-Puchta \cite{ScPu06} and Friedlander, Luca, and Stoiciu \cite{FLC07}. The case $k=4$ was proved by Pratt \cite{Pr22}. It is known that this sum is irrational for all $k\geq 1$ conditional on either Schinzel's conjecture (Schlage-Puchta \cite{ScPu06}) or Dickson's conjecture (Friedlander, Luca, and Stoiciu \cite{FLC07}). This is discussed in problem B14 of Guy's collection \cite{Gu04}.




References

[Er52] Erd\H{o}s, P., Problem 4493 . Amer. Math. Monthly (1952), 557-558.

[FLC07] Friedlander, J. B. and Luca, F. and Stoiciu, M., On the irrationality of a divisor function series . Integers (2007).

[Gu04] Guy, Richard K., Unsolved problems in number theory . (2004), xviii+437.

[Pr22] Pratt, K., The irrationality of a divisor function series of Erd\H{o}s and Kac . arXiv:2209.11124 (2022).

[ScPu06] Schlage-Puchta, J. C., The irrationality of a number theoretical series . Ramanujan J. (2006), 455-460.

\noindent\rule{\linewidth}{0.4pt}


%% ============================================================
\section{Problem \#254}
\label{prob:254}

\noindent\textbf{Status:} OPEN \quad|\quad \textbf{Prize:} none \quad|\quad \textbf{Updated:} 2025-08-31

\noindent\textbf{Tags:} number theory

\medskip

\noindent\textbf{Statement:}

Let $A\subseteq \mathbb{N}$ be such that\[\lvert A\cap [1,2x]\rvert -\lvert A\cap [1,x]\rvert \to \infty\textrm{ as }x\to \infty\]and\[\sum_{n\in A} \{ \theta n\}=\infty\]for every $\theta\in (0,1)$, where $\{x\}$ is the distance of $x$ from the nearest integer. Then every sufficiently large integer is the sum of distinct elements of $A$.




Cassels \cite{Ca60} proved this under the alternative hypotheses\[\lim \frac{\lvert A\cap [1,2x]\rvert -\lvert A\cap [1,x]\rvert}{\log\log x}=\infty\]and\[\sum_{n\in A} \{ \theta n\}^2=\infty\]for every $\theta\in (0,1)$.




References

[Ca60] Cassels, J. W. S., On the representation of integers as the sums of distinct summands taken from a fixed set . Acta Sci. Math. (Szeged) (1960), 111-124.

\noindent\rule{\linewidth}{0.4pt}


%% ============================================================
\section{Problem \#256}
\label{prob:256}

\noindent\textbf{Status:} OPEN \quad|\quad \textbf{Prize:} none \quad|\quad \textbf{Updated:} 2025-08-31

\noindent\textbf{Tags:} analysis

\medskip

\noindent\textbf{Statement:}

Let $n\geq 1$ and $f(n)$ be maximal such that for any integers $1\leq a_1\leq \cdots \leq a_n$ we have\[\max_{\lvert z\rvert=1}\left\lvert \prod_{i}(1-z^{a_i})\right\rvert\geq f(n).\]Estimate $f(n)$ - in particular, is it true that there exists some constant $c&#62;0$ such that\[\log f(n) \gg n^c?\]




Erd\H{o}s and Szekeres \cite{ErSz59} proved that $\lim f(n)^{1/n}=1$ and $f(n)&gt;\sqrt{2n}$. Erd\H{o}s proved an upper bound of $\log f(n) \ll n^{1-c}$ for some constant $c&gt;0$ with probabilistic methods. Atkinson \cite{At61} showed that $\log f(n) \ll n^{1/2}\log n$. This was improved to\[\log f(n) \ll n^{1/3}(\log n)^{4/3}\]by Odlyzko \cite{Od82}. If we denote by $f^*(n)$ the analogous quantity with the assumption that $a_1&lt;\cdots&lt;a_n$ then Bourgain and Chang \cite{BoCh18} prove that\[\log f^*(n)\ll (n\log n)^{1/2}\log\log n.\]Atkinson \cite{At61} noted this is related to the Chowla cosine problem [510] , in that if for any set of $n$ integers $A$ there exists $\theta$ such that $\sum_{n\in A}\cos(n\theta) &lt; -M_n$ then\[\log f^*(n) \ll M_n \log n.\]The answer to the specific question asked is no - Belov and Konyagin \cite{BeKo96} proved that\[\log f(n) \ll (\log n)^4.\]




References

[At61] Atkinson, F. V., On a problem of Erd\H{o}s and Szekeres . Canad. Math. Bull. (1961), 7-12.

[BeKo96] Belov, A. S. and Konyagin, S. V., An estimate for the free term of a nonnegative trigonometric polynomial with integer coefficients . Mat. Zametki (1996), 627--629.

[BoCh18] Bourgain, J. and Chang, Mei-Chu, On a paper of Erd\H{o}s and Szekeres . J. Anal. Math. (2018), 253-271.

[ErSz59] Erd\H{o}s, P. and Szekeres, G., On the product $\Pi^n_{k=1}(1-z^ak)$ . Acad. Serbe Sci. Publ. Inst. Math. (1959), 29-34.

[Od82] Odlyzko, A. M., Minima of cosine sums and maxima of polynomials on the unit circle . J. London Math. Soc. (2) (1982), 412-420.

\noindent\rule{\linewidth}{0.4pt}


%% ============================================================
\section{Problem \#257}
\label{prob:257}

\noindent\textbf{Status:} OPEN \quad|\quad \textbf{Prize:} none \quad|\quad \textbf{Updated:} 2025-08-31

\noindent\textbf{Tags:} irrationality

\medskip

\noindent\textbf{Statement:}

Let $A\subseteq \mathbb{N}$ be an infinite set. Is\[\sum_{n\in A}\frac{1}{2^n-1}\]irrational?




If $A=\mathbb{N}$ then this series is $\sum_{n}\frac{\tau(n)}{2^n}$, where $\tau(n)$ is the number of divisors of $n$, which Erd\H{o}s \cite{Er48} proved is irrational. In general, if $f_A(n)$ counts the number of divisors of $n$ which are elements of $A$ then\[\sum_{n\in A}\frac{1}{2^n-1}=\sum_n \frac{f_A(n)}{2^n}.\]The case when $A$ is the set of primes is [69] . This case (and when $A$ is the set of prime powers) was settled in the affirmative by Tao and Ter\"{a}v\"{a}inen \cite{TaTe25}. Erd\H{o}s \cite{Er68d} proved this sum is irrational whenever $(a,b)=1$ for all $a\neq b\in A$ and $\sum_{n\in A}\frac{1}{n}&lt;\infty$ (and thought that the condition $(a,b)=1$ could be dropped by complicating his proof). There is nothing special about $2$ here, and this sum is likely irrational with $2$ replaced by any integer $t\geq 2$. In \cite{Er88c} Erd\H{o}s goes further and speculates that $\sum_{n\in A}\frac{1}{2^n-t_n}$ is irrational for every infinite set $A$ and bounded sequence $t_n$ (presumably of integers, and presumably excluding the case when $t_n=0$ for all $n$). This was disproved by Kova\v{c} and Tao \cite{KoTa24}, and in the comments Kova\v{c} has sketched a proof that there exists some choice of $t_n$ with $1\leq t_n\leq 6$ such that this sum is rational.




References

[Er48] Erd\H{o}s, P., On arithmetical properties of Lambert series . J. Indian Math. Soc. (N.S.) (1948), 63-66.

[Er68d] Erd\H{o}s, P., On the irrationality of certain series . Math. Student (1968), 222--226.

[Er88c] Erd\H{o}s, P., On the irrationality of certain series: problems and results . New advances in transcendence theory (Durham, 1986) (1988), 102-109.

[KoTa24] Kova\vC, V. and Tao T., On several irrationality problems for Ahmes series . arXiv:2406.17593 (2024).

[TaTe25] T. Tao and J. Ter\"{a}v\"{a}inen, Quantitative correlations and some problems on prime factors of consecutive integers . arXiv:2512.01739 (2025).

\noindent\rule{\linewidth}{0.4pt}


%% ============================================================
\section{Problem \#260}
\label{prob:260}

\noindent\textbf{Status:} OPEN \quad|\quad \textbf{Prize:} none \quad|\quad \textbf{Updated:} 2025-08-31

\noindent\textbf{Tags:} irrationality

\medskip

\noindent\textbf{Statement:}

Let $a_1&#60;a_2&#60;\cdots$ be an increasing sequence such that $a_n/n\to \infty$. Is the sum\[\sum_n \frac{a_n}{2^{a_n}}\]irrational?




Erd\H{o}s \cite{Er81l} proved this is true under either of the stronger assumptions that {UL} {LI} $a_{n+1}-a_n\to \infty$ or {/LI} {LI} $a_n \gg n\sqrt{\log n\log\log n}$.{/LI} {/UL} Erd\H{o}s and Graham speculate that the condition $\limsup a_{n+1}-a_n=\infty$ is not sufficient, but know of no example. In \cite{Va99} this is asked with $a_n/n\to \infty$ replaced by $a_{n+1}-a_n\to \infty$, which as above was proved by Erd\H{o}s in \cite{Er81l}.




References

[Er81l] Erd\H{o}s, Paul, Sur l'irrationalit\'{e}{} d'une certaine s\'{e}rie . C. R. Acad. Sci. Paris S\'{e}r. I Math. (1981), 765--768.

[Va99] Various, Some of Paul's favorite problems . Booklet produced for the conference "Paul Erd\H{o}s and his mathematics", Budapest, July 1999 (1999).

\noindent\rule{\linewidth}{0.4pt}


%% ============================================================
\section{Problem \#261}
\label{prob:261}

\noindent\textbf{Status:} OPEN \quad|\quad \textbf{Prize:} none \quad|\quad \textbf{Updated:} 2025-08-31

\noindent\textbf{Tags:} number theory

\medskip

\noindent\textbf{Statement:}

Are there infinitely many $n$ such that there exists some $t\geq 2$ and distinct integers $a_1,\ldots,a_t\geq 1$ such that\[\frac{n}{2^n}=\sum_{1\leq k\leq t}\frac{a_k}{2^{a_k}}?\]Is this true for all $n$? Is there a rational $x$ such that\[x = \sum_{k=1}^\infty \frac{a_k}{2^{a_k}}\]has at least $2^{\aleph_0}$ solutions?




Related to [260] . In \cite{Er88c} Erd\H{o}s notes that Cusick had a simple proof that there do exist infinitely many such $n$. Erd\H{o}s does not record what this was, but a later paper by Borwein and Loring \cite{BoLo90} provides the following proof: for every positive integer $m$ and $n=2^{m+1}-m-2$ we have\[\frac{n}{2^n}=\sum_{n&lt;k\leq n+m}\frac{k}{2^k}.\]Tengely, Ulas, and Zygadlo \cite{TUZ20} have verified that all $n\leq 10000$ have the required property. In \cite{Er88c} Erd\H{o}s weakens the second question to asking for the existence of a rational $x$ which has two solutions.




References

[BoLo90] Borwein, Peter and Loring, Terry A., Some questions of {E}rd\H{o}s and {G}raham on numbers of the form {$\sum g_n/2^{g_n}$} . Math. Comp. (1990), 377--394.

[Er88c] Erd\H{o}s, P., On the irrationality of certain series: problems and results . New advances in transcendence theory (Durham, 1986) (1988), 102-109.

[TUZ20] Tengely, Szabolcs and Ulas, Maciej and Zygad\l o, Jakub, On a {D}iophantine equation of {E}rd\H{o}s and {G}raham . J. Number Theory (2020), 445--459.

\noindent\rule{\linewidth}{0.4pt}


%% ============================================================
\section{Problem \#263}
\label{prob:263}

\noindent\textbf{Status:} OPEN \quad|\quad \textbf{Prize:} none \quad|\quad \textbf{Updated:} 2025-08-31

\noindent\textbf{Tags:} irrationality

\medskip

\noindent\textbf{Statement:}

Let $a_n$ be an increasing sequence of positive integers such that for every sequence of positive integers $b_n$ with $b_n/a_n\to 1$ the sum\[\sum\frac{1}{b_n}\]is irrational. Is $a_n=2^{2^n}$ such a sequence? Must such a sequence satisfy $a_n^{1/n}\to \infty$?




One possible definition of an 'irrationality sequence' (see also [262] and [264] ). A folklore result states that $\sum \frac{1}{a_n}$ is irrational whenever $\lim a_n^{1/2^n}=\infty$. Kova\v{c} and Tao \cite{KoTa24} have proved that any strictly increasing sequence such that $\sum \frac{1}{a_n}$ converges and $\lim a_{n+1}/a_n^2=0$ is not such an irrationality sequence. On the other hand, if\[\liminf \frac{a_{n+1}}{a_n^{2+\epsilon}}&gt;0\]for some $\epsilon&gt;0$ then the above folklore result implies that $a_n$ is such an irrationality sequence. This problem was originally wrongly stated without the assumption that the sequence be increasing; DeepMind has given a counterexample without this assumption (see the comments). Koizumi \cite{Ko25c} has proved that $a_n=\lfloor \alpha^{2^n}\rfloor$ is an irrationality sequence of this type for all but countably many $\alpha&gt;1$.




References

[Ko25c] J. Koizumi, Irrationality of the reciprocal sum of doubly exponential sequences . arXiv:2504.05933 (2025).

[KoTa24] Kova\vC, V. and Tao T., On several irrationality problems for Ahmes series . arXiv:2406.17593 (2024).

\noindent\rule{\linewidth}{0.4pt}


%% ============================================================
\section{Problem \#264}
\label{prob:264}

\noindent\textbf{Status:} OPEN \quad|\quad \textbf{Prize:} none \quad|\quad \textbf{Updated:} 2025-08-31

\noindent\textbf{Tags:} irrationality

\medskip

\noindent\textbf{Statement:}

Let $a_n$ be a sequence of positive integers such that for every bounded sequence of integers $b_n$ (with $a_n+b_n\neq 0$ and $b_n\neq 0$ for all $n$) the sum\[\sum \frac{1}{a_n+b_n}\]is irrational. Are $a_n=2^n$ or $a_n=n!$ examples of such a sequence?




A possible definition of an 'irrationality sequence' (see also [262] and [263] ). One example is $a_n=2^{2^n}$. In \cite{ErGr80} they also ask whether such a sequence can have polynomial growth, but Erd\H{o}s later retracted this in \cite{Er88c}, claiming 'It is not hard to show that it cannot increase slower than exponentially'. Kova\v{c} and Tao \cite{KoTa24} have proved that $2^n$ is not such an irrationality sequence. More generally, they prove that any strictly increasing sequence of positive integers such that $\sum\frac{1}{a_n}$ converges and\[\liminf \left(a_n^2\sum_{k&gt;n}\frac{1}{a_k^2}\right) &gt;0 \]is not such an irrationality sequence. In particular, any strictly increasing sequence with $\limsup a_{n+1}/a_n &lt;\infty$ is not such an irrationality sequence. On the other hand, Kova\v{c} and Tao do prove that for any function $F$ with $\lim F(n+1)/F(n)=\infty$ there exists such an irrationality sequence with $a_n\sim F(n)$.




References

[Er88c] Erd\H{o}s, P., On the irrationality of certain series: problems and results . New advances in transcendence theory (Durham, 1986) (1988), 102-109.

[ErGr80] Erd\H{o}s, P. and Graham, R., Old and new problems and results in combinatorial number theory . Monographies de L'Enseignement Mathematique (1980).

[KoTa24] Kova\vC, V. and Tao T., On several irrationality problems for Ahmes series . arXiv:2406.17593 (2024).

\noindent\rule{\linewidth}{0.4pt}


%% ============================================================
\section{Problem \#265}
\label{prob:265}

\noindent\textbf{Status:} OPEN \quad|\quad \textbf{Prize:} none \quad|\quad \textbf{Updated:} 2025-08-31

\noindent\textbf{Tags:} irrationality

\medskip

\noindent\textbf{Statement:}

Let $1\leq a_1&#60;a_2&#60;\cdots$ be an increasing sequence of integers. How fast can $a_n\to \infty$ grow if\[\sum\frac{1}{a_n}\quad\textrm{and}\quad\sum\frac{1}{a_n-1}\]are both rational?




Cantor observed that $a_n=\binom{n}{2}$ is such a sequence. If we replace $-1$ by a different constant then higher degree polynomials can be used - for example if we consider $\sum_{n\geq 2}\frac{1}{a_n}$ and $\sum_{n\geq 2}\frac{1}{a_n-12}$ then $a_n=n^3+6n^2+5n$ is an example of both series being rational. Erd\H{o}s believed that $a_n^{1/n}\to \infty$ is possible, but $a_n^{1/2^n}\to 1$ is necessary. This has been almost completely solved by Kova\v{c} and Tao \cite{KoTa24}, who prove that such a sequence can grow doubly exponentially. More precisely, there exists such a sequence such that $a_n^{1/\beta^n}\to \infty$ for some $\beta &#62;1$. It remains open whether one can achieve\[\limsup a_n^{1/2^n}&#62;1.\]A folklore result states that $\sum \frac{1}{a_n}$ is irrational whenever $\lim a_n^{1/2^n}=\infty$, and hence such a sequence cannot grow faster than doubly exponentially - the remaining question is the precise exponent possible.




References

[KoTa24] Kova\vC, V. and Tao T., On several irrationality problems for Ahmes series . arXiv:2406.17593 (2024).

\noindent\rule{\linewidth}{0.4pt}


%% ============================================================
\section{Problem \#267}
\label{prob:267}

\noindent\textbf{Status:} OPEN \quad|\quad \textbf{Prize:} none \quad|\quad \textbf{Updated:} 2025-08-31

\noindent\textbf{Tags:} irrationality

\medskip

\noindent\textbf{Statement:}

Let $F_1=F_2=1$ and $F_{n+1}=F_n+F_{n-1}$ be the Fibonacci sequence. Let $n_1&lt;n_2&lt;\cdots $ be an infinite sequence with $n_{k+1}/n_k \geq c&gt;1$. Must\[\sum_k\frac{1}{F_{n_k}}\]be irrational?




It may be sufficient to have $n_k/k\to \infty$. Good \cite{Go74} and Bicknell and Hoggatt \cite{BiHo76} have shown that $\sum \frac{1}{F_{2^n}}$ is irrational - in fact,\[\sum \frac{1}{F_{2^n}}=\frac{7-\sqrt{5}}{2}.\]Badea \cite{Ba87} proved that $\sum \frac{1}{F_{2^n+1}}$ is irrational. The sum $\sum \frac{1}{F_n}$ itself was proved to be irrational by Andr\'{e}-Jeannin \cite{An89}. The main problem has been proved for $c\geq 2$ by Badea \cite{Ba93}. It remains open for $1&#60;c&#60;2$.




References

[An89] Andr\'{e}-Jeannin, Richard, Irrationalit\'{e}{} de la somme des inverses de certaines suites r\'{e}currentes . C. R. Acad. Sci. Paris S\'{e}r. I Math. (1989), 539--541.

[Ba87] Badea, C., The irrationality of certain infinite series . Glasgow Math. J. (1987), 221--228.

[Ba93] Badea, C., A theorem on irrationality of infinite series and applications . Acta Arith. (1993), 313--323.

[BiHo76] Hoggatt, Jr., V. E. and Bicknell, Marjorie, A reciprocal series of Fibonacci numbers with subscripts $2^nk$ . Fibonacci Quart. (1976), 453-455.

[Go74] Good, I. J., A reciprocal series of Fibonacci numbers . Fibonacci Quart. (1974), 346.

\noindent\rule{\linewidth}{0.4pt}


%% ============================================================
\section{Problem \#269}
\label{prob:269}

\noindent\textbf{Status:} OPEN \quad|\quad \textbf{Prize:} none \quad|\quad \textbf{Updated:} 2025-08-31

\noindent\textbf{Tags:} irrationality

\medskip

\noindent\textbf{Statement:}

Let $P$ be a finite set of primes with $\lvert P\rvert \geq 2$ and let $\{a_1&#60;a_2&#60;\cdots\}=\{ n\in \mathbb{N} : \textrm{if }p\mid n\textrm{ then }p\in P\}$. Is the sum\[\sum_{n=1}^\infty \frac{1}{[a_1,\ldots,a_n]},\]where $[a_1,\ldots,a_n]$ is the lowest common multiple of $a_1,\ldots,a_n$, irrational?




If $P$ is infinite this sum is always irrational (in \cite{Er88c} Erd\H{o}s says this is a 'simple exercise'). This problem was asked by Erd\H{o}s in a letter to the editor written January 1st 1973 in issue 12 of the Fibonacci Quarterly, 1974, p. 335. In that letter he says that he can prove the sum is irrational if duplicate summands are removed.




References

[Er88c] Erd\H{o}s, P., On the irrationality of certain series: problems and results . New advances in transcendence theory (Durham, 1986) (1988), 102-109.

\noindent\rule{\linewidth}{0.4pt}


%% ============================================================
\section{Problem \#271}
\label{prob:271}

\noindent\textbf{Status:} OPEN \quad|\quad \textbf{Prize:} none \quad|\quad \textbf{Updated:} 2025-08-31

\noindent\textbf{Tags:} additive combinatorics, arithmetic progressions

\medskip
\noindent\textit{Stanley sequences}


\noindent\textbf{Statement:}

Let $A(n)=\{a_0&#60;a_1&#60;\cdots\}$ be the sequence defined by $a_0=0$ and $a_1=n$, and for $k\geq 1$ define $a_{k+1}$ as the least positive integer such that there is no three-term arithmetic progression in $\{a_0,\ldots,a_{k+1}\}$. Can the $a_k$ be explicitly determined? How fast do they grow?




It is easy to see that $A(1)$ is the set of integers which have no 2 in their base 3 expansion. Odlyzko and Stanley \cite{OdSt78} have found similar characterisations are known for $A(3^k)$ and $A(2\cdot 3^k)$ for any $k\geq 0$ and conjectured in general that such a sequence always eventually either satisfies\[a_k\asymp k^{\log_23}\]or\[a_k \asymp \frac{k^2}{\log k}.\]There is no known sequence which satisfies the second growth rate, but Lindhurst \cite{Li90} gives data which suggests that $A(4)$ has such growth ($A(4)$ is given as A005487 in the OEIS). Moy \cite{Mo11} has proved that, for all such sequences, for all $\epsilon&gt;0$, $a_k\leq (\frac{1}{2}+\epsilon)k^2$ for all sufficiently large $k$. van Doorn and Sothanaphan have noted in the comment section that Moy's proof can be upgraded to give a fully explicit result of\[a_k\leq \frac{(k-1)(k+2)}{2}+n\]for all $k\geq 0$. In general, sequences which begin with some initial segment and thereafter are continued in a greedy fashion to avoid three-term arithmetic progressions are known as Stanley sequences.




References

[Li90] S. Lindhurst, An investigation of several interesting sets of numbers generated by the greedy algorithm . Senior thesis at Princeton University (1990).

[Mo11] Moy, Richard A., On the growth of the counting function of Stanley sequences . Discrete Math. (2011), 560-562.

[OdSt78] A. Odlyzko and R. Stanley, Some curious sequences constructed with the greedy algorithm . Bell Laboratories internal memorandum (1978).

\noindent\rule{\linewidth}{0.4pt}


%% ============================================================
\section{Problem \#272}
\label{prob:272}

\noindent\textbf{Status:} OPEN \quad|\quad \textbf{Prize:} none \quad|\quad \textbf{Updated:} 2025-08-31

\noindent\textbf{Tags:} additive combinatorics, arithmetic progressions

\medskip

\noindent\textbf{Statement:}

Let $N\geq 1$. What is the largest $t$ such that there are $A_1,\ldots,A_t\subseteq \{1,\ldots,N\}$ with $A_i\cap A_j$ a non-empty arithmetic progression for all $i\neq j$?




Simonovits and S\'{o}s \cite{SiSo81} have shown that $t\ll N^2$. Erd\H{o}s and Graham asked whether the maximal $t$ is achieved when we take the $A_i$ to be all arithmetic progressions in $\{1,\ldots,N\}$ containing some fixed element, 'presumably the integer $\lfloor N/2\rfloor$'. This was disproved by Simonovits and S\'{o}s \cite{SiSo81}, who observed that taking all sets containing at most $3$ elements, containing some fixed element, produces $\binom{N}{2}+1$ many such sets, which is asymptotically greater than the number of arithmetic progressions containing a fixed element, which is $\sim \frac{\pi^2}{24}N^2$. If we drop the non-empty requirement then Graham, Simonovits, and S\'{o}s \cite{GSS80} have shown that\[t\leq \binom{N}{3}+\binom{N}{2}+\binom{N}{1}+1\]and this is best possible. Szabo \cite{Sz99} proved that the maximal such $t$ is equal to\[\frac{N^2}{2}+O(N^{5/3}(\log N)^3),\]resolving the asymptotic question. On the other hand, Szabo showed that the conjecture of Simonovits and S\'{o}s that $\binom{n}{2}+1$ is best possible is false, giving a construction which yields\[t \geq \binom{N}{2}+\left\lfloor\frac{N-1}{4}\right\rfloor+1.\]Szabo conjectures that the asymptotic $t=\binom{N}{2}+O(N)$ holds, and that in any extremal example there is an integer contained in all sets.




References

[GSS80] Graham, R. L. and Simonovits, M. and S\'{o}s, V. T., A note on the intersection properties of subsets of integers . J. Combin. Theory Ser. A (1980), 106-110.

[SiSo81] Simonovits, Mikl\'{o}s and S\'{o}s, Vera T., Intersection properties of subsets of integers . European J. Combin. (1981), 363-372.

[Sz99] Szab\'o, Tibor, Intersection properties of subsets of integers . European J. Combin. (1999), 429--444.

\noindent\rule{\linewidth}{0.4pt}


%% ============================================================
\section{Problem \#273}
\label{prob:273}

\noindent\textbf{Status:} OPEN \quad|\quad \textbf{Prize:} none \quad|\quad \textbf{Updated:} 2025-08-31

\noindent\textbf{Tags:} number theory, covering systems

\medskip

\noindent\textbf{Statement:}

Is there a covering system all of whose moduli are of the form $p-1$ for some primes $p\geq 5$?




Selfridge has found an example using divisors of $360$ if $p=3$ is allowed.

\noindent\rule{\linewidth}{0.4pt}


%% ============================================================
\section{Problem \#274}
\label{prob:274}

\noindent\textbf{Status:} OPEN \quad|\quad \textbf{Prize:} none \quad|\quad \textbf{Updated:} 2025-08-31

\noindent\textbf{Tags:} group theory, covering systems

\medskip
\noindent\textit{Herzog-Schönheim conjecture}


\noindent\textbf{Statement:}

If $G$ is a group then can there exist an exact covering of $G$ by more than one cosets of different sizes? (i.e. each element is contained in exactly one of the cosets)




A question of Herzog and Sch\"{o nheim}, who conjectured more generally that if $G$ is any (not necessarily finite) group and $a_1G_1,\ldots,a_kG_k$ are finitely many cosets of subgroups of $G$ with distinct indices $[G:G_i]$ then the $a_iG_i$ cannot form a partition of $G$. This conjecture was proved in the case when all the $G_i$ are subnormal in $G$ by Sun \cite{Su04}. In particular if $G$ is abelian (which was the special case asked about in \cite{Er77c} and \cite{ErGr80}) the answer to the original question is no. Margolis and Schnabel \cite{MaSc19} proved this conjecture for all groups $G$ of size $&lt;1440$.




References

[Er77c] Erd\H{o}s, Paul, Problems and results on combinatorial number theory. III . Number theory day (Proc. Conf., Rockefeller Univ., New York, 1976) (1977), 43-72.

[ErGr80] Erd\H{o}s, P. and Graham, R., Old and new problems and results in combinatorial number theory . Monographies de L'Enseignement Mathematique (1980).

[MaSc19] Margolis, Leo and Schnabel, Ofir, The {H}erzog-Sch\"onheim conjecture for small groups and harmonic subgroups . Beitr. Algebra Geom. (2019), 399--418.

[Su04] Sun, Zhi-Wei, On the {H}erzog-Sch\"onheim conjecture for uniform covers of groups . J. Algebra (2004), 153--175.

\noindent\rule{\linewidth}{0.4pt}


%% ============================================================
\section{Problem \#276}
\label{prob:276}

\noindent\textbf{Status:} OPEN \quad|\quad \textbf{Prize:} none \quad|\quad \textbf{Updated:} 2025-08-31

\noindent\textbf{Tags:} number theory, covering systems

\medskip

\noindent\textbf{Statement:}

Is there an infinite Lucas sequence $a_0,a_1,\ldots$ where $a_{n+2}=a_{n+1}+a_n$ for $n\geq 0$ such that all $a_k$ are composite, and yet no integer has a common factor with every term of the sequence?




Whether such a composite Lucas sequence even exists was open for a while, but using covering systems Graham \cite{Gr64} showed that\[a_0 = 1786772701928802632268715130455793\]and\[a_1 = 1059683225053915111058165141686995\]generate such a sequence. This problem asks whether one can have a composite Lucas sequence without 'an underlying system of covering congruences responsible'. This problem has been 'conjecturally solved' by Ismailescu and Son \cite{IsSo14}, in that they provide an explicit infinite Lucas sequence in which all the terms are composite, and believe that no covering system is responsible for this. See the comment by van Doorn below for more details. See also [1113] for another problem in which the question is whether covering systems are always responsible.




References

[Gr64] Graham, R. L., A Fibonacci-Like Sequence of Composite Numbers . Math. Mag. (1964), 322-324.

[IsSo14] Ismailescu, Dan and Son, Jaesung, A new kind of {F}ibonacci-like sequence of composite numbers . J. Integer Seq. (2014), Article 14.8.2, 9.

\noindent\rule{\linewidth}{0.4pt}


%% ============================================================
\section{Problem \#278}
\label{prob:278}

\noindent\textbf{Status:} OPEN \quad|\quad \textbf{Prize:} none \quad|\quad \textbf{Updated:} 2025-08-31

\noindent\textbf{Tags:} number theory, covering systems

\medskip

\noindent\textbf{Statement:}

Let $A=\{n_1&#60;\cdots&#60;n_r\}$ be a finite set of positive integers. What is the maximum density of integers covered by a suitable choice of congruences $a_i\pmod{n_i}$? Is the minimum density achieved when all the $a_i$ are equal?




Simpson \cite{Si86} has observed that the density of integers covered is at least\[\sum_i \frac{1}{n_i}-\sum_{i&#60;j}\frac{1}{[n_i,n_j]}+\sum_{i&#60;j&#60;k}\frac{1}{[n_i,n_j,n_k]}-\cdot\](where $[\cdots]$ denotes the least common multiple) which is achieved when all $a_i$ are equal, settling the second question.




References

[Si86] Simpson, R. J., Exact coverings of the integers by arithmetic progressions . Discrete Math. (1986), 181--190.

\noindent\rule{\linewidth}{0.4pt}


%% ============================================================
\section{Problem \#279}
\label{prob:279}

\noindent\textbf{Status:} OPEN \quad|\quad \textbf{Prize:} none \quad|\quad \textbf{Updated:} 2025-08-31

\noindent\textbf{Tags:} number theory, covering systems, primes

\medskip

\noindent\textbf{Statement:}

Let $k\geq 3$. Is there a choice of congruence classes $a_p\pmod{p}$ for every prime $p$ such that all sufficiently large integers can be written as $a_p+tp$ for some prime $p$ and integer $t\geq k$?




Even the case $k=3$ seems difficult. This may be true with the primes replaced by any set $A\subseteq \mathbb{N}$ such that\[\lvert A\cap [1,N]\rvert \gg N/\log N\]and\[\sum_{\substack{n\in A\\ n\leq N}}\frac{1}{n} -\log\log N\to \infty\]as $N\to \infty$. If we weaken the question to ask for only almost all integers to be represented in this fashion then, for any $k\geq 2$, any set $A$ such that $\sum_{n\in A}\frac{1}{n}=\infty$ has this property.

\noindent\rule{\linewidth}{0.4pt}


%% ============================================================
\section{Problem \#282}
\label{prob:282}

\noindent\textbf{Status:} OPEN \quad|\quad \textbf{Prize:} none \quad|\quad \textbf{Updated:} 2025-08-31

\noindent\textbf{Tags:} number theory, unit fractions

\medskip

\noindent\textbf{Statement:}

Let $A\subseteq \mathbb{N}$ be an infinite set and consider the following greedy algorithm for a rational $x\in (0,1)$: choose the minimal $n\in A$ such that $n\geq 1/x$ and repeat with $x$ replaced by $x-\frac{1}{n}$. If this terminates after finitely many steps then this produces a representation of $x$ as the sum of distinct unit fractions with denominators from $A$. Does this process always terminate if $x$ has odd denominator and $A$ is the set of odd numbers? More generally, for which pairs $x$ and $A$ does this process terminate?




In 1202 Fibonacci observed that this process terminates for any $x$ when $A=\mathbb{N}$. The problem when $A$ is the set of odd numbers is due to Stein. Graham \cite{Gr64b} has shown that $\frac{m}{n}$ is the sum of distinct unit fractions with denominators $\equiv a\pmod{d}$ if and only if\[\left(\frac{n}{(n,(a,d))},\frac{d}{(a,d)}\right)=1.\]Does the greedy algorithm always terminate in such cases? Graham \cite{Gr64c} has also shown that $x$ is the sum of distinct unit fractions with square denominators if and only if $x\in [0,\pi^2/6-1)\cup [1,\pi^2/6)$. Does the greedy algorithm for this always terminate? Erd\H{o}s and Graham believe not - indeed, perhaps it fails to terminate almost always. See also [206] .




References

[Gr64b] Graham, R. L., On finite sums of unit fractions . Proc. London Math. Soc. (3) (1964), 193-207.

[Gr64c] Graham, R. L., On finite sums of reciprocals of distinct nth powers . Pacific J. Math. (1964), 85-92.

\noindent\rule{\linewidth}{0.4pt}


%% ============================================================
\section{Problem \#288}
\label{prob:288}

\noindent\textbf{Status:} OPEN \quad|\quad \textbf{Prize:} none \quad|\quad \textbf{Updated:} 2025-08-31

\noindent\textbf{Tags:} number theory, unit fractions

\medskip

\noindent\textbf{Statement:}

Is it true that there are only finitely many pairs of intervals $I_1,I_2$ such that\[\sum_{n_1\in I_1}\frac{1}{n_1}+\sum_{n_2\in I_2}\frac{1}{n_2}\in \mathbb{N}?\]




For example,\[\frac{1}{3}+\frac{1}{4}+\frac{1}{5}+\frac{1}{6}+\frac{1}{20}=1.\]This is still open even if $\lvert I_2\rvert=1$. It is perhaps true with two intervals replaced by any $k$ intervals.

\noindent\rule{\linewidth}{0.4pt}


%% ============================================================
\section{Problem \#289}
\label{prob:289}

\noindent\textbf{Status:} OPEN \quad|\quad \textbf{Prize:} none \quad|\quad \textbf{Updated:} 2025-08-31

\noindent\textbf{Tags:} number theory, unit fractions

\medskip

\noindent\textbf{Statement:}

Is it true that, for all sufficiently large $k$, there exist finite intervals $I_1,\ldots,I_k\subset \mathbb{N}$, distinct, not overlapping or adjacent, with $\lvert I_i\rvert \geq 2$ for $1\leq i\leq k$ such that\[1=\sum_{i=1}^k \sum_{n\in I_i}\frac{1}{n}?\]




Erd\H{o}s and Graham posed this in \cite{ErGr80} without the stipulation the intervals be distinct, non-overlapping, or adjacent, but Kovac in the comments has provided a simple argument showing that it is easily possible without this restriction, and likely \cite{ErGr80} just forgot to mention this natural restriction. As an example representing $2$ rather than $1$, Hickerson and Montgomery, in the solution to AMS Monthly problem E2689 proposed by Hahn, found\[2=\sum_{i=1}^5 \sum_{n\in I_i}\frac{1}{n}\]where $I_1=[2,7]$, $I_2=[9,10]$, $I_3=[17,18]$, $I_4=[34,35]$, and $I_5=[84,85]$.




References

[ErGr80] Erd\H{o}s, P. and Graham, R., Old and new problems and results in combinatorial number theory . Monographies de L'Enseignement Mathematique (1980).

\noindent\rule{\linewidth}{0.4pt}


%% ============================================================
\section{Problem \#291}
\label{prob:291}

\noindent\textbf{Status:} OPEN \quad|\quad \textbf{Prize:} none \quad|\quad \textbf{Updated:} 2025-08-31

\noindent\textbf{Tags:} number theory, unit fractions

\medskip

\noindent\textbf{Statement:}

Let $n\geq 1$ and define $L_n$ to be the least common multiple of $\{1,\ldots,n\}$ and $a_n$ by\[\sum_{1\leq k\leq n}\frac{1}{k}=\frac{a_n}{L_n}.\]Is it true that $(a_n,L_n)=1$ and $(a_n,L_n)&#62;1$ both occur for infinitely many $n$?




Steinerberger has observed that the answer to the second question is trivially yes: for example, any $n$ which begins with a $2$ in base $3$ has $3\mid (a_n,L_n)$. More generally, if the leading digit of $n$ in base $p$ is $p-1$ then $p\mid (a_n,L_n)$. There is in fact a necessary and sufficient condition: a prime $p\leq n$ divides $(a_n,L_n)$ if and only if $p$ divides the numerator of $1+\cdots+\frac{1}{k}$, where $k$ is the leading digit of $n$ in base $p$. This can be seen by writing\[a_n = \frac{L_n}{1}+\cdots+\frac{L_n}{n}\]and observing that the right-hand side is congruent to $1+\cdots+1/k$ modulo $p$. (The previous claim about $p-1$ follows immediately from Wolstenholme's theorem .) This leads to a heuristic prediction (see for example a preprint of Shiu \cite{Sh16}) of $\asymp\frac{x}{\log x}$ for the number of $n\in [1,x]$ such that $(a_n,L_n)=1$. In particular, there should be infinitely many $n$, but the set of such $n$ should have density zero. Unfortunately this heuristic is difficult to turn into a proof. Wu and Yan \cite{WuYa22} have proved, conditional on $\frac{1}{\log p}$ being linearly independent over $\mathbb{Q}$ for any finite collection of primes $p$ (itself a consequence of Schanuel's conjecture), that the set of $n$ for which $(a_n,L_n)&gt;1$ has upper density $1$.




References

[Sh16] P. Shiu, The denominators of harmonic numbers . arXiv:1607.02863 (2016).

[WuYa22] Wu, Bing-Ling and Yan, Xiao-Hui, On the denominators of harmonic numbers. {IV} . C. R. Math. Acad. Sci. Paris (2022), 53--57.

\noindent\rule{\linewidth}{0.4pt}


%% ============================================================
\section{Problem \#293}
\label{prob:293}

\noindent\textbf{Status:} OPEN \quad|\quad \textbf{Prize:} none \quad|\quad \textbf{Updated:} 2025-08-31

\noindent\textbf{Tags:} number theory, unit fractions

\medskip

\noindent\textbf{Statement:}

Let $k\geq 1$ and let $v(k)$ be the minimal integer which does not appear as some $n_i$ in a solution to\[1=\frac{1}{n_1}+\cdots+\frac{1}{n_k}\]with $1\leq n_1&#60;\cdots &#60;n_k$. Estimate the growth of $v(k)$.




Results of Bleicher and Erd\H{o}s \cite{BlEr75} imply $v(k) \gg k!$. It may be that $v(k)$ grows doubly exponentially in $\sqrt{k}$ or even $k$. An elementary inductive argument shows that $n_k\leq ku_k$ where $u_1=1$ and $u_{i+1}=u_i(u_i+1)$, and hence\[v(k) \leq kc_0^{2^k},\]where\[c_0=\lim_n u_n^{1/2^n}=1.26408\cdots\]is the 'Vardi constant' (small improvements on this are possible as in [148] ). van Doorn and Tang \cite{vDTa25b} have proved that\[v(k)\geq e^{ck^2}\]for some constant $c&gt;0$, and noted a close connection to [304] . In particular, if $N(b)\ll \log\log b$ as in [304] then it is likely the methods of \cite{vDTa25b} prove $v(k) \geq e^{e^{ck}}$ for some $c&gt;0$.




References

[BlEr75] Bleicher, M. N. and Erd\H{o}s, P., The number of distinct subsums of $\sum \sb{1}\spN\,1/i$ . Math. Comp. (1975), 29-42.

[vDTa25b] W. van Doorn and Q. Tang, The smallest denominator not contained in a unit fraction decomposition of 1 with fixed length . arXiv:2512.22083 (2025).

\noindent\rule{\linewidth}{0.4pt}


%% ============================================================
\section{Problem \#295}
\label{prob:295}

\noindent\textbf{Status:} OPEN \quad|\quad \textbf{Prize:} none \quad|\quad \textbf{Updated:} 2025-08-31

\noindent\textbf{Tags:} number theory, unit fractions

\medskip

\noindent\textbf{Statement:}

Let $N\geq 1$ and let $k(N)$ denote the smallest $k$ such that there exist $N\leq n_1&#60;\cdots &#60;n_k$ with\[1=\frac{1}{n_1}+\cdots+\frac{1}{n_k}.\]Is it true that\[\lim_{N\to \infty} k(N)-(e-1)N=\infty?\]




Erd\H{o}s and Straus \cite{ErSt71b} have proved the existence of some constant $c&#62;0$ such that\[-c &#60; k(N)-(e-1)N \ll \frac{N}{\log N}.\]




References

[ErSt71b] Erd\H{o}s, P. and Straus, E. G., Solution to Problem . Amer. Math. Monthly (1971), 302-303.

\noindent\rule{\linewidth}{0.4pt}


%% ============================================================
\section{Problem \#301}
\label{prob:301}

\noindent\textbf{Status:} OPEN \quad|\quad \textbf{Prize:} none \quad|\quad \textbf{Updated:} 2025-08-31

\noindent\textbf{Tags:} number theory, unit fractions

\medskip

\noindent\textbf{Statement:}

Let $f(N)$ be the size of the largest $A\subseteq \{1,\ldots,N\}$ such that there are no solutions to\[\frac{1}{a}= \frac{1}{b_1}+\cdots+\frac{1}{b_k}\]with distinct $a,b_1,\ldots,b_k\in A$? Estimate $f(N)$. In particular, is it true that $f(N)=(\tfrac{1}{2}+o(1))N$?




The example $A=(N/2,N]\cap \mathbb{N}$ shows that $f(N)\geq N/2$. Wouter van Doorn has given an elementary argument that proves\[f(N)\leq (25/28+o(1))N.\]Indeed, consider the sets $S_a=\{2a,3a,4a,6a,12a\}\cap [1,N]$ as $a$ ranges over all integers of the form $8^b9^cd$ with $(d,6)=1$. All such $S_a$ are disjoint and, if $A$ has no solutions to the given equation, then $A$ must omit at least two elements of $S_a$ when $a\leq N/12$ and at least one element of $S_a$ when $N/12&lt;a\leq N/6$, and an elementary calculation concludes the proof. Stijn Cambie and Wouter van Doorn have noted that, if we allow solutions to this equation with non-distinct $b_i$, then the size of the maximal set is at most $N/2$. Indeed, this is the classical threshold for the existence of some distinct $a,b\in A$ such that $a\mid b$. See also [302] and [327] .

\noindent\rule{\linewidth}{0.4pt}


%% ============================================================
\section{Problem \#302}
\label{prob:302}

\noindent\textbf{Status:} OPEN \quad|\quad \textbf{Prize:} none \quad|\quad \textbf{Updated:} 2025-08-31

\noindent\textbf{Tags:} number theory, unit fractions

\medskip

\noindent\textbf{Statement:}

Let $f(N)$ be the size of the largest $A\subseteq \{1,\ldots,N\}$ such that there are no solutions to\[\frac{1}{a}= \frac{1}{b}+\frac{1}{c}\]with distinct $a,b,c\in A$? Estimate $f(N)$. In particular, is $f(N)=(\tfrac{1}{2}+o(1))N$?




The colouring version of this is [303] , which was solved by Brown and R\"{o}dl \cite{BrRo91}. One can take either $A$ to be all odd integers in $[1,N]$ or all integers in $[N/2,N]$ to show $f(N)\geq (1/2+o(1))N$. Wouter van Doorn has proved (see this note ) that\[f(N) \leq (9/10+o(1))N.\]Stijn Cambie has observed that\[f(N)\geq (5/8+o(1))N,\]taking $A$ to be all odd integers $\leq N/4$ and all integers in $[N/2,N]$. Stijn Cambie has also observed that, if we allow $b=c$, then there is a solution to this equation when $\lvert A\rvert \geq (\tfrac{2}{3}+o(1))N$, since then there must exist some $n,2n\in A$. See also [301] and [327] .




References

[BrRo91] Brown, Tom C. and R\"{o}dl, Voijtech, Monochromatic solutions to equations with unit fractions . Bull. Austral. Math. Soc. (1991), 387-392.

\noindent\rule{\linewidth}{0.4pt}


%% ============================================================
\section{Problem \#304}
\label{prob:304}

\noindent\textbf{Status:} OPEN \quad|\quad \textbf{Prize:} none \quad|\quad \textbf{Updated:} 2025-08-31

\noindent\textbf{Tags:} number theory, unit fractions

\medskip

\noindent\textbf{Statement:}

For integers $1\leq a&#60;b$ let $N(a,b)$ denote the minimal $k$ such that there exist integers $1&#60;n_1&#60;\cdots&#60;n_k$ with\[\frac{a}{b}=\frac{1}{n_1}+\cdots+\frac{1}{n_k}.\]Estimate $N(b)=\max_{1\leq a&#60;b}N(a,b)$. Is it true that $N(b) \ll \log\log b$?




Erd\H{o}s \cite{Er50c} proved that\[\log\log b \ll N(b) \ll \frac{\log b}{\log\log b}.\]The upper bound was improved by Vose \cite{Vo85} to\[N(b) \ll \sqrt{\log b}.\]One can also investigate the average of $N(a,b)$ for fixed $b$, and it is known that\[\frac{1}{b}\sum_{1\leq a&lt;b}N(a,b) \gg \log\log b.\]Related to [18] . There is also a close connection to [293] (particularly with $N(b-1,b)$), as elucidated by van Doorn and Tang \cite{vDTa25b}. This problem has been formalised in Lean as part of the Google DeepMind Formal Conjectures project .




References

[Er50c] Erd\H{o}s, P., Az $1/x_1 + 1/x_2 + \ldots + 1/x_n =A/B$ egyenlet eg\'{e}sz sz\'{a}m\'{u} megold\'{a}sair\'{o}l . Mat. Lapok (1950), 192-210.

[Vo85] Vose, Michael D., Egyptian fractions . Bull. London Math. Soc. (1985), 21-24.

[vDTa25b] W. van Doorn and Q. Tang, The smallest denominator not contained in a unit fraction decomposition of 1 with fixed length . arXiv:2512.22083 (2025).

\noindent\rule{\linewidth}{0.4pt}


%% ============================================================
\section{Problem \#306}
\label{prob:306}

\noindent\textbf{Status:} OPEN \quad|\quad \textbf{Prize:} none \quad|\quad \textbf{Updated:} 2025-08-31

\noindent\textbf{Tags:} number theory, unit fractions

\medskip

\noindent\textbf{Statement:}

Let $a/b\in \mathbb{Q}_{&#62;0}$ with $b$ squarefree. Are there integers $1&#60;n_1&#60;\cdots&#60;n_k$, each the product of two distinct primes, such that\[\frac{a}{b}=\frac{1}{n_1}+\cdots+\frac{1}{n_k}?\]




For $n_i$ the product of three distinct primes, this is true when $b=1$, as proved by Butler, Erd\H{o}s and Graham \cite{BEG15} (this paper is perhaps Erd\H{o}s' last paper, appearing 19 years after his death).




References

[BEG15] Butler, Steve and Erd\H{o}s, Paul and Graham, Ron, Egyptian fractions with each denominator having three distinct prime divisors . Integers (2015), Paper No. A51, 9.

\noindent\rule{\linewidth}{0.4pt}


%% ============================================================
\section{Problem \#311}
\label{prob:311}

\noindent\textbf{Status:} OPEN \quad|\quad \textbf{Prize:} none \quad|\quad \textbf{Updated:} 2025-08-31

\noindent\textbf{Tags:} number theory, unit fractions

\medskip

\noindent\textbf{Statement:}

Let $\delta(N)$ be the minimal non-zero value of $\lvert 1-\sum_{n\in A}\frac{1}{n}\rvert$ as $A$ ranges over all subsets of $\{1,\ldots,N\}$. Is it true that\[\delta(N)=e^{-(c+o(1))N}\]for some constant $c\in (0,1)$?




It is trivial that\[\delta(N)\geq \frac{1}{[1,\ldots,N]}=e^{-(1+o(1))N},\]where $[1,\ldots,N]$ is the least common multiple of $\{1,\ldots,N\}$. The formulation in \cite{ErGr80} has the additional condition that $A$ contain no $S$ such that $\sum_{n\in S}\frac{1}{n}=1$, but Kovac in the comments has shown that the simpler formulation above is equivalent. Tang has shown that\[\delta(N) \leq \exp\left(-c\frac{N}{(\log N\log\log N)^3}\right)\]for some constant $c&gt;0$.




References

[ErGr80] Erd\H{o}s, P. and Graham, R., Old and new problems and results in combinatorial number theory . Monographies de L'Enseignement Mathematique (1980).

\noindent\rule{\linewidth}{0.4pt}


%% ============================================================
\section{Problem \#312}
\label{prob:312}

\noindent\textbf{Status:} OPEN \quad|\quad \textbf{Prize:} none \quad|\quad \textbf{Updated:} 2025-08-31

\noindent\textbf{Tags:} number theory, unit fractions

\medskip

\noindent\textbf{Statement:}

Does there exist some $c&#62;0$ such that, for any $K&#62;1$, whenever $A$ is a sufficiently large finite multiset of positive integers with $\sum_{n\in A}\frac{1}{n}&#62;K$ there exists some $S\subseteq A$ such that\[1-e^{-cK} &#60; \sum_{n\in S}\frac{1}{n}\leq 1?\]




Erd\H{o}s and Graham knew this with $e^{-cK}$ replaced by $c/K^2$.

\noindent\rule{\linewidth}{0.4pt}


%% ============================================================
\section{Problem \#313}
\label{prob:313}

\noindent\textbf{Status:} OPEN \quad|\quad \textbf{Prize:} none \quad|\quad \textbf{Updated:} 2025-08-31

\noindent\textbf{Tags:} number theory, unit fractions

\medskip

\noindent\textbf{Statement:}

Are there infinitely many solutions to\[\frac{1}{p_1}+\cdots+\frac{1}{p_k}=1-\frac{1}{m},\]where $m\geq 2$ is an integer and $p_1&#60;\cdots&#60;p_k$ are distinct primes?




For example,\[\frac{1}{2}+\frac{1}{3}=1-\frac{1}{6}\]and\[\frac{1}{2}+\frac{1}{3}+\frac{1}{7}=1-\frac{1}{42}.\]It is clear that we must have $m=p_1\cdots p_k$, and hence in particular there is at most one solution for each $m$. The integers $m$ for which there is such a solution are known as primary pseudoperfect numbers , and there are $8$ known, listed in A054377 at the OEIS.

\noindent\rule{\linewidth}{0.4pt}


%% ============================================================
\section{Problem \#317}
\label{prob:317}

\noindent\textbf{Status:} OPEN \quad|\quad \textbf{Prize:} none \quad|\quad \textbf{Updated:} 2025-08-31

\noindent\textbf{Tags:} number theory, unit fractions

\medskip

\noindent\textbf{Statement:}

Is there some constant $c&#62;0$ such that for every $n\geq 1$ there exists some $\delta_k\in \{-1,0,1\}$ for $1\leq k\leq n$ with\[0&lt; \left\lvert \sum_{1\leq k\leq n}\frac{\delta_k}{k}\right\rvert &lt; \frac{c}{2^n}?\]Is it true that for sufficiently large $n$, for any $\delta_k\in \{-1,0,1\}$,\[\left\lvert \sum_{1\leq k\leq n}\frac{\delta_k}{k}\right\rvert &gt; \frac{1}{[1,\ldots,n]}\]whenever the left-hand side is not zero?




Inequality is obvious for the second claim, the problem is strict inequality. This fails for small $n$, for example\[\frac{1}{2}-\frac{1}{3}-\frac{1}{4}=-\frac{1}{12}.\]Arguments of Kovac and van Doorn in the comment section prove a weak version of the first question, with an upper bound of\[2^{-n\frac{(\log\log\log n)^{1+o(1)}}{\log n}},\]and van Doorn gives a heuristic that suggests this may be the true order of magnitude.

\noindent\rule{\linewidth}{0.4pt}


%% ============================================================
\section{Problem \#319}
\label{prob:319}

\noindent\textbf{Status:} OPEN \quad|\quad \textbf{Prize:} none \quad|\quad \textbf{Updated:} 2025-08-31

\noindent\textbf{Tags:} number theory, unit fractions

\medskip

\noindent\textbf{Statement:}

What is the size of the largest $A\subseteq \{1,\ldots,N\}$ such that there is a function $\delta:A\to \{-1,1\}$ such that\[\sum_{n\in A}\frac{\delta_n}{n}=0\]and\[\sum_{n\in A'}\frac{\delta_n}{n}\neq 0\]for all non-empty $A'\subsetneq A$?




Adenwalla has observed that a lower bound of\[\lvert A\rvert\geq (1-\tfrac{1}{e}+o(1))N\]follows from the main result of Croot \cite{Cr01}, which states that there exists some set of integers $B\subset [(\frac{1}{e}-o(1))N,N]$ such that $\sum_{b\in B}\frac{1}{b}=1$. Since the sum of $\frac{1}{m}$ for $m\in [c_1N,c_2N]$ is asymptotic to $\log(c_2/c_1)$ we must have $\lvert B\rvert \geq (1-\tfrac{1}{e}+o(1))N$. We may then let $A=B\cup\{1\}$ and choose $\delta(n)=-1$ for all $n\in B$ and $\delta(1)=1$. This problem has been formalised in Lean as part of the Google DeepMind Formal Conjectures project .




References

[Cr01] Croot, III, Ernest S., On unit fractions with denominators in short intervals . Acta Arith. (2001), 99-114.

\noindent\rule{\linewidth}{0.4pt}


%% ============================================================
\section{Problem \#320}
\label{prob:320}

\noindent\textbf{Status:} OPEN \quad|\quad \textbf{Prize:} none \quad|\quad \textbf{Updated:} 2025-08-31

\noindent\textbf{Tags:} number theory, unit fractions

\medskip

\noindent\textbf{Statement:}

Let $S(N)$ count the number of distinct sums of the form $\sum_{n\in A}\frac{1}{n}$ for $A\subseteq \{1,\ldots,N\}$. Estimate $S(N)$.




Bleicher and Erd\H{o}s \cite{BlEr75} proved the lower bound\[\log S(N)\geq \frac{N}{\log N}\left(\log 2\prod_{i=3}^k\log_iN\right),\]valid for $k\geq 4$ and $\log_kN\geq k$, and also \cite{BlEr76b} proved the upper bound\[\log S(N)\leq \frac{N}{\log N}\left(\log_r N \prod_{i=3}^r \log_iN\right),\]valid for $r\geq 1$ and $\log_{2r}N\geq 1$. (In these bounds $\log_in$ denotes the $i$-fold iterated logarithm.) Bettin, Greni\'{e}, Molteni, and Sanna \cite{BGMS25} improved the lower bound to\[\log S(N) \geq \frac{N}{\log N}\left(2\log 2\left(1-\frac{3/2}{\log_kN}\right)\prod_{i=3}^k\log_iN\right),\]valid for $k\geq 4$ and $\log_kN\geq 3/2$. (In particular this goes to infinity faster than the lower bound of Bleicher and Erd\H{o}s.) See also [321] .




References

[BGMS25] S. Bettin, L. Greni\'{e}, G. Molteni, and C. Sanna, A lower bound for the number of Egyptian fractions . arXiv:2509.10030 (2025).

[BlEr75] Bleicher, M. N. and Erd\H{o}s, P., The number of distinct subsums of $\sum \sb{1}\spN\,1/i$ . Math. Comp. (1975), 29-42.

[BlEr76b] Bleicher, Michael N. and Erd\H{o}s, Paul, Denominators of Egyptian fractions. II . Illinois J. Math. (1976), 598-613.

\noindent\rule{\linewidth}{0.4pt}


%% ============================================================
\section{Problem \#321}
\label{prob:321}

\noindent\textbf{Status:} OPEN \quad|\quad \textbf{Prize:} none \quad|\quad \textbf{Updated:} 2025-08-31

\noindent\textbf{Tags:} number theory, unit fractions

\medskip

\noindent\textbf{Statement:}

What is the size of the largest $A\subseteq \{1,\ldots,N\}$ such that all sums $\sum_{n\in S}\frac{1}{n}$ are distinct for $S\subseteq A$?




Let $R(N)$ be the maximal such size. Results of Bleicher and Erd\H{o}s from \cite{BlEr75} and \cite{BlEr76b} imply that\[\frac{N}{\log N}\prod_{i=3}^k\log_iN\leq R(N)\leq \frac{1}{\log 2}\log_r N\left(\frac{N}{\log N} \prod_{i=3}^r \log_iN\right),\]valid for any $k\geq 4$ with $\log_kN\geq k$ and any $r\geq 1$ with $\log_{2r}N\geq 1$. (In these bounds $\log_in$ denotes the $i$-fold iterated logarithm.) See also [320] .




References

[BlEr75] Bleicher, M. N. and Erd\H{o}s, P., The number of distinct subsums of $\sum \sb{1}\spN\,1/i$ . Math. Comp. (1975), 29-42.

[BlEr76b] Bleicher, Michael N. and Erd\H{o}s, Paul, Denominators of Egyptian fractions. II . Illinois J. Math. (1976), 598-613.

\noindent\rule{\linewidth}{0.4pt}


%% ============================================================
\section{Problem \#322}
\label{prob:322}

\noindent\textbf{Status:} OPEN \quad|\quad \textbf{Prize:} none \quad|\quad \textbf{Updated:} 2025-08-31

\noindent\textbf{Tags:} number theory, powers

\medskip

\noindent\textbf{Statement:}

Let $k\geq 3$ and $A\subset \mathbb{N}$ be the set of $k$th powers. What is the order of growth of $1_A^{(k)}(n)$, i.e. the number of representations of $n$ as the sum of $k$ many $k$th powers? Does there exist some $c&#62;0$ and infinitely many $n$ such that\[1_A^{(k)}(n) &#62;n^c?\]




Connected to Waring's problem. The famous Hypothesis $K$ of Hardy and Littlewood was that $1_A^{(k)}(n)\leq n^{o(1)}$, but this was disproved by Mahler \cite{Ma36} for $k=3$, who constructed infinitely many $n$ such that\[1_A^{(3)}(n)\gg n^{1/12}\](where $A$ is the set of cubes). Erd\H{o}s believed Hypothesis $K$ fails for all $k\geq 4$, but this is unknown. Hardy and Littlewood made the weaker Hypothesis $K^*$ that for all $N$ and $\epsilon&#62;0$\[\sum_{n\leq N}1_A^{(k)}(n)^2 \ll_\epsilon N^{1+\epsilon}.\]Erd\H{o}s and Graham remark: 'This is probably true but no doubt very deep. However, it would suffice for most applications.' Independently Erd\H{o}s \cite{Er36} and Chowla proved that for all $k\geq 3$ and infinitely many $n$\[1_A^{(k)}(n) \gg n^{c/\log\log n}\]for some constant $c&#62;0$ (depending on $k$). In \cite{Er65b} Erd\H{o}s claims an unpublished proof that, if $B$ is the set of $k$th powers of any set of positive density, then\[\limsup 1_B^{(k)}(n)=\infty.\]This is discussed in problem D4 of Guy's collection \cite{Gu04}.




References

[Er36] Erd\H{o}s, Paul, On the Representation of an Integer as the Sum of k k-th Powers . J. London Math. Soc. (1936), 133-136.

[Er65b] Erd\H{o}s, Paul, Some recent advances and current problems in number theory . Lectures on Modern Mathematics, Vol. III (1965), 196-244.

[Gu04] Guy, Richard K., Unsolved problems in number theory . (2004), xviii+437.

[Ma36] Mahler, Kurt, Note on Hypothesis K of Hardy and Littlewood . J. London Math. Soc. (1936), 136-138.

\noindent\rule{\linewidth}{0.4pt}


%% ============================================================
\section{Problem \#323}
\label{prob:323}

\noindent\textbf{Status:} OPEN \quad|\quad \textbf{Prize:} none \quad|\quad \textbf{Updated:} 2025-08-31

\noindent\textbf{Tags:} number theory, powers

\medskip

\noindent\textbf{Statement:}

Let $1\leq m\leq k$ and $f_{k,m}(x)$ denote the number of integers $\leq x$ which are the sum of $m$ many nonnegative $k$th powers. Is it true that\[f_{k,k}(x) \gg_\epsilon x^{1-\epsilon}\]for all $\epsilon&#62;0$? Is it true that if $m&#60;k$ then\[f_{k,m}(x) \gg x^{m/k}\]for sufficiently large $x$?




This would have significant applications to Waring's problem. Erd\H{o}s and Graham describe this as 'unattackable by the methods at our disposal'. The case $k=2$ was resolved by Landau, who showed\[f_{2,2}(x) \sim \frac{cx}{\sqrt{\log x}}\]for some constant $c&#62;0$. For $k&#62;2$ it is not known if $f_{k,k}(x)=o(x)$.

\noindent\rule{\linewidth}{0.4pt}


%% ============================================================
\section{Problem \#324}
\label{prob:324}

\noindent\textbf{Status:} OPEN \quad|\quad \textbf{Prize:} none \quad|\quad \textbf{Updated:} 2025-08-31

\noindent\textbf{Tags:} number theory, powers

\medskip

\noindent\textbf{Statement:}

Does there exist a polynomial $f(x)\in\mathbb{Z}[x]$ such that all the sums $f(a)+f(b)$ with $a&#60;b$ nonnegative integers are distinct?




In other words, the set $\{ f(n) : n\geq 1\}$ should be a Sidon set. Erd\H{o}s and Graham describe this problem as 'very annoying'. Probably $f(x)=x^5$ should work. The Lander, Parkin, and Selfridge conjecture would imply that $f(x)=x^n$ has this property for all $n\geq 5$. Ruzsa \cite{Ru01b} has proved there exists $c\in [0,1]$ and $n_0\geq 1$ such that\[\{ n^5 +\lfloor cn^4\rfloor : n \geq n_0\}\]is a Sidon set. It is easy to check that a quadratic $f$ cannot have this property. Dubickas and Novikas \cite{DuNo21} have shown that a cubic $f$ cannot have this property. It is a classical fact that $x^4$ cannot have this property. This is discussed in problems C9, D1, and F30 of Guy's collection \cite{Gu04}.




References

[DuNo21] Dubickas, Arturas and Novikas, Aivaras, No cubic integer polynomial generates a Sidon sequence . Math. Nachr. (2021), 1859--1865.

[Gu04] Guy, Richard K., Unsolved problems in number theory . (2004), xviii+437.

[Ru01b] Ruzsa, I. Z., An almost polynomial Sidon sequence . Studia Sci. Math. Hungar. (2001), 367--375.

\noindent\rule{\linewidth}{0.4pt}


%% ============================================================
\section{Problem \#325}
\label{prob:325}

\noindent\textbf{Status:} OPEN \quad|\quad \textbf{Prize:} none \quad|\quad \textbf{Updated:} 2025-08-31

\noindent\textbf{Tags:} number theory, powers

\medskip

\noindent\textbf{Statement:}

Let $k\geq 3$ and $f_{k,3}(x)$ denote the number of integers $\leq x$ which are the sum of three nonnegative $k$th powers. Is it true that\[f_{k,3}(x) \gg x^{3/k}\]or even $\gg_\epsilon x^{3/k-\epsilon}$?




Mahler and Erd\H{o}s \cite{ErMa38} proved that $f_{k,2}(x) \gg x^{2/k}$. For $k=3$ the best known is due to Wooley \cite{Wo15},\[f_{3,3}(x) \gg x^{0.917\cdots}.\]This problem has been formalised in Lean as part of the Google DeepMind Formal Conjectures project .




References

[ErMa38] Erd\H{o}s, P\'{a}l and Mahler, Kurt, On the number of integers which can be represented by a binary form . Doc. Math. (2019), 475-481.

[Wo15] Wooley, Trevor D., Sums of three cubes, II . Acta Arith. (2015), 73-100.

\noindent\rule{\linewidth}{0.4pt}


%% ============================================================
\section{Problem \#326}
\label{prob:326}

\noindent\textbf{Status:} OPEN \quad|\quad \textbf{Prize:} none \quad|\quad \textbf{Updated:} 2025-08-31

\noindent\textbf{Tags:} number theory, additive basis

\medskip

\noindent\textbf{Statement:}

Does there exist $A=\{a_1&#60;a_2&#60;\cdots\}\subset \mathbb{N}$ which is a minimal basis of order $2$ (i.e. every large integer is the sum of $2$ elements from $A$, and no proper subset of $A$ has this property), such that\[\lim_{k\to \infty}\frac{a_k}{k^2}=c\]for some $c\neq 0$?




Erd\H{o}s originally asked this for any basis (not necessarily minimal); such a basis was constructed by Cassels \cite{Ca57}. Erd\H{o}s and Graham conjectured a negative answer to this question.




References

[Ca57] Cassels, J. W. S., \"{U}ber Basen der nat\"{u}rlichen Zahlenreihe . Abh. Math. Sem. Univ. Hamburg (1957), 247-257.

\noindent\rule{\linewidth}{0.4pt}


%% ============================================================
\section{Problem \#327}
\label{prob:327}

\noindent\textbf{Status:} OPEN \quad|\quad \textbf{Prize:} none \quad|\quad \textbf{Updated:} 2025-08-31

\noindent\textbf{Tags:} number theory, unit fractions

\medskip

\noindent\textbf{Statement:}

Suppose $A\subseteq \{1,\ldots,N\}$ is such that if $a,b\in A$ and $a\neq b$ then $a+b\nmid ab$. Can $A$ be 'substantially more' than the odd numbers? What if $a,b\in A$ with $a\neq b$ implies $a+b\nmid 2ab$? Must $\lvert A\rvert=o(N)$?




The connection to unit fractions comes from the observation that $\frac{1}{a}+\frac{1}{b}$ is a unit fraction if and only if $a+b\mid ab$. Wouter van Doorn has given an elementary argument that proves that if $A\subseteq \{1,\ldots,N\}$ has $\lvert A\rvert \geq (25/28+o(1))N$ then $A$ must contain $a\neq b$ with $a+b\mid ab$ (see the discussion in [301] ). See also [302] .

\noindent\rule{\linewidth}{0.4pt}


%% ============================================================
\section{Problem \#329}
\label{prob:329}

\noindent\textbf{Status:} OPEN \quad|\quad \textbf{Prize:} none \quad|\quad \textbf{Updated:} 2025-08-31

\noindent\textbf{Tags:} number theory

\medskip

\noindent\textbf{Statement:}

Suppose $A\subseteq \mathbb{N}$ is a Sidon set. How large can\[\limsup_{N\to \infty}\frac{\lvert A\cap \{1,\ldots,N\}\rvert}{N^{1/2}}\]be?




Erd\H{o}s proved that $1/2$ is possible and Kr\"{u}ckeberg \cite{Kr61} proved $1/\sqrt{2}$ is possible. Erd\H{o}s and Tur\'{a}n \cite{ErTu41} have proved this $\limsup$ is always $\leq 1$. In \cite{Er80} he writes that he and Kr\"{u}ckeberg conjecture that $1$ is possible. The fact that $1$ is possible would follow if any finite Sidon set is a subset of a perfect difference set (see [44] and [707] ). This question can also be asked for $B_2[g]$ sequences (i.e. where the number of solutions to $n=a_1+a_2$ with $a_1\leq a_2$ is at most $g$ for all $n$, so that a $B_2[1]$ set is a Sidon set). Kolountzakis \cite{Ko96} constructed a $B_2[2]$ sequence where the $\limsup$ is $1$, and for larger $g$ constructions were provided by Cilleruelo and Trujillo \cite{CiTr01}.




References

[CiTr01] Cilleruelo, Javier and Trujillo, Carlos, Infinite {$B_2[g]$} sequences . Israel J. Math. (2001), 263--267.

[Er80] Erd\H{o}s, Paul, A survey of problems in combinatorial number theory . Ann. Discrete Math. (1980), 89-115.

[ErTu41] Erd\H{o}s, P. and Tur\'{a}n, P., On a problem of Sidon in additive number theory, and on some related problems . J. London Math. Soc. (1941), 212-215.

[Ko96] Kolountzakis, Mihail N., On the additive complements of the primes and sets of similar growth . Acta Arith. (1996), 1--8.

[Kr61] Kr\"{u}ckeberg, Fritz, $B\sb{2}$-Folgen und verwandte Zahlenfolgen . J. Reine Angew. Math. (1961), 53-60.

\noindent\rule{\linewidth}{0.4pt}


%% ============================================================
\section{Problem \#332}
\label{prob:332}

\noindent\textbf{Status:} OPEN \quad|\quad \textbf{Prize:} none \quad|\quad \textbf{Updated:} 2025-08-31

\noindent\textbf{Tags:} number theory

\medskip

\noindent\textbf{Statement:}

Let $A\subseteq \mathbb{N}$ and $D(A)$ be the set of those numbers which occur infinitely often as $a_1-a_2$ with $a_1,a_2\in A$. What conditions on $A$ are sufficient to ensure $D(A)$ has bounded gaps?




Prikry, Tijdeman, Stewart, and others (see the survey articles \cite{St78} and \cite{Ti79}) have shown that a sufficient condition is that $A$ has positive density. One can also ask what conditions are sufficient for $D(A)$ to have positive density, or for $\sum_{d\in D(A)}\frac{1}{d}=\infty$, or even just $D(A)\neq\emptyset$.




References

[St78] Stewart, Cam L., On difference sets of sets of integers . S\'{e}minaire Delange-Pisot-Poitou, 19e ann\'{e}e: 1977/78, Th\'{e}orie des nombres, Fasc. 1 (1978), Exp. No. 5, 8.

[Ti79] Tijdeman, R., Distance sets of sequences of integers . Proceedings, Bicentennial Congress Wiskundig Genootschap (Vrije Univ., Amsterdam, 1978), Part II (1979), 405-415.

\noindent\rule{\linewidth}{0.4pt}


%% ============================================================
\section{Problem \#334}
\label{prob:334}

\noindent\textbf{Status:} OPEN \quad|\quad \textbf{Prize:} none \quad|\quad \textbf{Updated:} 2025-08-31

\noindent\textbf{Tags:} number theory

\medskip

\noindent\textbf{Statement:}

Find the best function $f(n)$ such that every $n$ can be written as $n=a+b$ where both $a,b$ are $f(n)$-smooth (that is, are not divisible by any prime $p&#62;f(n)$.)




Erd\H{o}s originally asked if even $f(n)\leq n^{1/3}$ is true. This is known, and the best bound is due to Balog \cite{Ba89} who proved that\[f(n) \ll_\epsilon n^{\frac{4}{9\sqrt{e}}+\epsilon}\]for all $\epsilon&gt;0$. (Note $\frac{4}{9\sqrt{e}}=0.2695\ldots$.) It is likely that $f(n)\leq n^{o(1)}$. See also Problem 59 on Green's open problems list.




References

[Ba89] Balog, A., On additive representation of integers . Acta Math. Hungar. (1989), 297-301.

\noindent\rule{\linewidth}{0.4pt}


%% ============================================================
\section{Problem \#335}
\label{prob:335}

\noindent\textbf{Status:} OPEN \quad|\quad \textbf{Prize:} none \quad|\quad \textbf{Updated:} 2025-08-31

\noindent\textbf{Tags:} number theory, additive combinatorics

\medskip

\noindent\textbf{Statement:}

Let $d(A)$ denote the density of $A\subseteq \mathbb{N}$. Characterise those $A,B\subseteq \mathbb{N}$ with positive density such that\[d(A+B)=d(A)+d(B).\]




One way this can happen is if there exists $\theta&#62;0$ such that\[A=\{ n&#62;0 : \{ n\theta\} \in X_A\}\textrm{ and }B=\{ n&#62;0 : \{n\theta\} \in X_B\}\]where $\{x\}$ denotes the fractional part of $x$ and $X_A,X_B\subseteq \mathbb{R}/\mathbb{Z}$ are such that $\mu(X_A+X_B)=\mu(X_A)+\mu(X_B)$. Are all possible $A$ and $B$ generated in a similar way (using other groups)? This has been partially resolved by Ackelsburg and Richter \cite{AcRi26}, under the additional assumption that one of the sets meets every residue class in $\mathbb{N}$. Roughly speaking, if $d(A)&#62;0$ and $d(A+B)=d(A)+d(B)&#60;1$ and $B$ meets every residue class then either both $A$ and $B$ arise from a rotation on the product of the circle group with a finite cyclic group, or $A$ is in a single residue class modulo some $h$ and $B$ 'looks like' the union of all but one residue class modulo $h$. A full characterisation without the assumption that one of the sets meets every residue class seems hopeless; for example, taking $A=B$ to be a random subset of the even integers with density $1/4$ we have $d(A+A)=1/2$.




References

[AcRi26] E. Ackelsburg and F. Richter, An inverse theorem for sumsets of sets of positive density in the integers . arXiv:2604.12864 (2026).

\noindent\rule{\linewidth}{0.4pt}


%% ============================================================
\section{Problem \#336}
\label{prob:336}

\noindent\textbf{Status:} OPEN \quad|\quad \textbf{Prize:} none \quad|\quad \textbf{Updated:} 2025-08-31

\noindent\textbf{Tags:} number theory, additive basis

\medskip

\noindent\textbf{Statement:}

For $r\geq 2$ let $h(r)$ be the maximal finite $k$ such that there exists a basis $A\subseteq \mathbb{N}$ of order $r$ (so every large integer is the sum of at most $r$ integers from $A$) and exact order $k$ (so every large integer is the sum of exactly $k$ integers from $A$). Find the value of\[\lim_r \frac{h(r)}{r^2}.\]




A simple example of the order of a basis differing from the exact order is given by $A=\cup_{k\geq 0}(2^{2k},2^{2k+1}]$, which has order $2$ but exact order $3$. Erd\H{o}s and Graham \cite{ErGr80b} have shown that a basis $A$ has an exact order if and only if $a_2-a_1,a_3-a_2,a_4-a_3,\ldots$ are coprime. They also proved that\[\frac{1}{4}\leq \lim_r \frac{h(r)}{r^2}\leq \frac{5}{4}.\]The best bounds known for the limit are\[\frac{1}{3}\leq \lim_r \frac{h(r)}{r^2}\leq \frac{1}{2},\]the lower bound originally due to Grekos \cite{Gr88} and the upper bound to Nash \cite{Na93}. Improved bounds in the lower order terms were given by Plagne \cite{Pl04}. Erd\H{o}s and Graham \cite{ErGr80b} showed $h(2)=4$. Nash \cite{Na93} showed $h(3)=7$. The value of $h(4)$ is unknown, but it is known \cite{Pl04} that $10\leq h(4)\leq 11$.




References

[ErGr80b] Erd\H{o}s, P. and Graham, R. L., On bases with an exact order . Acta Arith. (1980), 201-207.

[Gr88] Grekos, Georges, Sur l'ordre d'une base additive . ([1988?]), Exp. No. 31, 13.

[Na93] Nash, John C. M., Some applications of a theorem of {M}. {K}neser . J. Number Theory (1993), 1--8.

[Pl04] Plagne, Alain, \`A{} propos de la fonction {$X$} d'{E}rd\H{o}s et {G}raham . Ann. Inst. Fourier (Grenoble) (2004), 1717--1767.

\noindent\rule{\linewidth}{0.4pt}


%% ============================================================
\section{Problem \#338}
\label{prob:338}

\noindent\textbf{Status:} OPEN \quad|\quad \textbf{Prize:} none \quad|\quad \textbf{Updated:} 2025-08-31

\noindent\textbf{Tags:} number theory, additive basis

\medskip

\noindent\textbf{Statement:}

The restricted order of a basis is the least integer $t$ (if it exists) such that every large integer is the sum of at most $t$ distinct summands from $A$. What are necessary and sufficient conditions that this exists? Can it be bounded (when it exists) in terms of the order of the basis? What are necessary and sufficient conditions that this is equal to the order of the basis?




Bateman has observed that for $h\geq 3$ there is a basis of order $h$ with no restricted order, taking\[A=\{1\}\cup \{x&#62;0 : h\mid x\}.\]Kelly \cite{Ke57} has shown that any basis of order $2$ has restricted order at most $4$ and conjectured it always has restricted order at most $3$ (which he proved under the additional assumption that the basis has positive lower density). Kelly's conjecture was disproved by Hennecart \cite{He05}, who constructed a basis of order $2$ with restricted order $4$. The set of squares has order $4$ and restricted order $5$ (see \cite{Pa33}) and the set of triangular numbers has order $3$ and restricted order $3$ (see \cite{Sc54}). Is it true that if $A\backslash F$ is a basis for all finite sets $F$ then $A$ must have a restricted order? What if they are all bases of the same order? Hegyv\'{a}ri, Hennecart, and Plagne \cite{HHP07} have shown that for all $k\geq2$ there exists a basis of order $k$ which has restricted order at least\[2^{k-2}+k-1.\]




References

[HHP07] Hegyv\'{a}ri, Norbert and Hennecart, Fran\c cois and Plagne, Alain, Answer to a question by {B}urr and {E}rd\H{o}s on restricted addition, and related results . Combin. Probab. Comput. (2007), 747--756.

[He05] Hennecart, Fran\c cois, On the restricted order of asymptotic bases of order two . Ramanujan J. (2005), 123--130.

[Ke57] Kelly, John B., Restricted bases . Amer. J. Math. (1957), 258-264.

[Pa33] Pall, Gordon, On Sums of Squares . Amer. Math. Monthly (1933), 10-18.

[Sc54] Schinzel, A., Sur la d\'{e}composition des nombres naturels en sommes de nombres triangulaires distincts . Bull. Acad. Polon. Sci. Cl. III. (1954), 409-410.

\noindent\rule{\linewidth}{0.4pt}


%% ============================================================
\section{Problem \#340}
\label{prob:340}

\noindent\textbf{Status:} OPEN \quad|\quad \textbf{Prize:} none \quad|\quad \textbf{Updated:} 2025-08-31

\noindent\textbf{Tags:} number theory, additive combinatorics, sidon sets

\medskip

\noindent\textbf{Statement:}

Let $A=\{1,2,4,8,13,21,31,45,66,81,97,\ldots\}$ be the greedy Sidon sequence: we begin with $1$ and iteratively include the next smallest integer that preserves the Sidon property (i.e. there are no non-trivial solutions to $a+b=c+d$). What is the order of growth of $A$? Is it true that\[\lvert A\cap \{1,\ldots,N\}\rvert \gg N^{1/2-\epsilon}\]for all $\epsilon&#62;0$ and large $N$?




This sequence is sometimes called the Mian-Chowla sequence. It is trivial that this sequence grows at least like $\gg N^{1/3}$. Erd\H{o}s and Graham \cite{ErGr80} also asked about the difference set $A-A$, whether this has positive density, and whether this contains $22$. It does contain $22$, since $a_{15}-a_{14}=204-182=22$. The smallest integer which is unknown to be in $A-A$ is $33$ (see A080200 ). It may be true that all or almost all integers are in $A-A$. This sequence is at OEIS A005282 . See also [156] .




References

[ErGr80] Erd\H{o}s, P. and Graham, R., Old and new problems and results in combinatorial number theory . Monographies de L'Enseignement Mathematique (1980).

\noindent\rule{\linewidth}{0.4pt}


%% ============================================================
\section{Problem \#341}
\label{prob:341}

\noindent\textbf{Status:} OPEN \quad|\quad \textbf{Prize:} none \quad|\quad \textbf{Updated:} 2025-08-31

\noindent\textbf{Tags:} number theory

\medskip

\noindent\textbf{Statement:}

Let $A=\{a_1&#60;\cdots&#60;a_k\}$ be a finite set of positive integers and extend it to an infinite sequence $\overline{A}=\{a_1&#60;a_2&#60;\cdots \}$ by defining $a_{n+1}$ for $n\geq k$ to be the least integer exceeding $a_n$ which is not of the form $a_i+a_j$ with $i,j\leq n$. Is it true that the sequence of differences $a_{m+1}-a_m$ is eventually periodic?




An old problem of Dickson. Even a starting set as small as $\{1,4,9,16,25\}$ requires thousands of terms before periodicity occurs. This problem is discussed under Problem 7 on Green's open problems list .

\noindent\rule{\linewidth}{0.4pt}


%% ============================================================
\section{Problem \#342}
\label{prob:342}

\noindent\textbf{Status:} OPEN \quad|\quad \textbf{Prize:} none \quad|\quad \textbf{Updated:} 2025-08-31

\noindent\textbf{Tags:} number theory

\medskip

\noindent\textbf{Statement:}

With $a_1=1$ and $a_2=2$ let $a_{n+1}$ for $n\geq 2$ be the least integer $&#62;a_n$ which can be expressed uniquely as $a_i+a_j$ for $i&#60;j\leq n$. What can be said about this sequence? Do infinitely many pairs $a,a+2$ occur? Does this sequence eventually have periodic differences? Is the density $0$?




A problem of Ulam. The sequence is\[1,2,3,4,6,8,11,13,16,18,26,28,\ldots\]at OEIS A002858 . See also Problem 7 of Green's open problems list . This is problem C4 in Guy's collection \cite{Gu04}.




References

[Gu04] Guy, Richard K., Unsolved problems in number theory . (2004), xviii+437.

\noindent\rule{\linewidth}{0.4pt}


%% ============================================================
\section{Problem \#345}
\label{prob:345}

\noindent\textbf{Status:} OPEN \quad|\quad \textbf{Prize:} none \quad|\quad \textbf{Updated:} 2025-08-31

\noindent\textbf{Tags:} number theory, complete sequences

\medskip

\noindent\textbf{Statement:}

Let $A\subseteq \mathbb{N}$ be a complete sequence, and define the threshold of completeness $T(A)$ to be the least integer $m$ such that all $n\geq m$ are in\[P(A) = \left\{\sum_{n\in B}n : B\subseteq A\textrm{ finite }\right\}\](the existence of $T(A)$ is guaranteed by completeness). Is it true that there are infinitely many $k$ such that $T(n^k)&#62;T(n^{k+1})$?




Erd\H{o}s and Graham \cite{ErGr80} remark that very little is known about $T(A)$ in general. It is known that\[T(n)=1, T(n^2)=128, T(n^3)=12758,\]\[T(n^4)=5134240,\textrm{ and }T(n^5)=67898771.\]Erd\H{o}s and Graham remark that a good candidate for the $n$ in the question are $k=2^t$ for large $t$, perhaps even $t=3$, because of the highly restricted values of $n^{2^t}$ modulo $2^{t+1}$.




References

[ErGr80] Erd\H{o}s, P. and Graham, R., Old and new problems and results in combinatorial number theory . Monographies de L'Enseignement Mathematique (1980).

\noindent\rule{\linewidth}{0.4pt}


%% ============================================================
\section{Problem \#346}
\label{prob:346}

\noindent\textbf{Status:} OPEN \quad|\quad \textbf{Prize:} none \quad|\quad \textbf{Updated:} 2025-08-31

\noindent\textbf{Tags:} number theory, complete sequences

\medskip

\noindent\textbf{Statement:}

Let $A=\{1\leq a_1&lt; a_2&lt;\cdots\}$ be a set of integers such that {UL} {LI} $A\backslash B$ is complete for any finite subset $B$ and {/LI} {LI} $A\backslash B$ is not complete for any infinite subset $B$.{/LI} {/UL} (Here 'complete' means all sufficiently large integers can be written as a sum of distinct members of the sequence.) Is it true that if $a_{n+1}/a_n \geq 1+\epsilon$ for some $\epsilon&gt;0$ and all $n$ then\[\lim_n \frac{a_{n+1}}{a_n}=\frac{1+\sqrt{5}}{2}?\]




Graham \cite{Gr64d} has shown that the sequence $a_n=F_n-(-1)^{n}$, where $F_n$ is the $n$th Fibonacci number, has these properties. Erd\H{o}s and Graham \cite{ErGr80} remark that it is easy to see that if $a_{n+1}/a_n&#62;\frac{1+\sqrt{5}}{2}$ then the second property is automatically satisfied, and that it is not hard to construct very irregular sequences satisfying both properties.




References

[ErGr80] Erd\H{o}s, P. and Graham, R., Old and new problems and results in combinatorial number theory . Monographies de L'Enseignement Mathematique (1980).

[Gr64d] Graham, R. L., A property of Fibonacci numbers . Fibonacci Quart. (1964), 1-10.

\noindent\rule{\linewidth}{0.4pt}


%% ============================================================
\section{Problem \#348}
\label{prob:348}

\noindent\textbf{Status:} OPEN \quad|\quad \textbf{Prize:} none \quad|\quad \textbf{Updated:} 2025-08-31

\noindent\textbf{Tags:} number theory, complete sequences

\medskip

\noindent\textbf{Statement:}

For what values of $0\leq m&#60;n$ is there a complete sequence $A=\{a_1\leq a_2\leq \cdots\}$ of integers such that {UL} {LI} $A$ remains complete after removing any $m$ elements, but {/LI} {LI} $A$ is not complete after removing any $n$ elements? {/LI} {/UL}




The Fibonacci sequence $1,1,2,3,5,\ldots$ shows that $m=1$ and $n=2$ is possible. The sequence of powers of $2$ shows that $m=0$ and $n=1$ is possible. The case $m=2$ and $n=3$ is not known. van Doorn has shown that no such sequence exists for $2\leq m&lt;n$ if we interpret complete in the strong sense that\[\left\{ \sum_{n\in B}n : \textrm{ for all finite }B\subset A\right\}=\mathbb{N}.\]Erd\H{o}s and Graham most likely meant the weaker notion of completeness which allows finitely many exceptions, however.

\noindent\rule{\linewidth}{0.4pt}


%% ============================================================
\section{Problem \#349}
\label{prob:349}

\noindent\textbf{Status:} OPEN \quad|\quad \textbf{Prize:} none \quad|\quad \textbf{Updated:} 2025-08-31

\noindent\textbf{Tags:} number theory, complete sequences

\medskip

\noindent\textbf{Statement:}

For what values of $t,\alpha \in (0,\infty)$ is the sequence $\lfloor t\alpha^n\rfloor$ complete (that is, all sufficiently large integers are the sum of distinct integers of the form $\lfloor t\alpha^n\rfloor$)?




Even in the range $t\in (0,1)$ and $\alpha\in (1,2)$ the behaviour is surprisingly complex. For example, Graham \cite{Gr64e} has shown that for any $k$ there exists some $t_k\in (0,1)$ such that the set of $\alpha$ such that the sequence is complete consists of at least $k$ disjoint line segments. It seems likely that the sequence is complete for all $t&#62;0$ and all $1&#60;\alpha &#60; \frac{1+\sqrt{5}}{2}$. Proving this seems very difficult, since we do not even know whether $\lfloor (3/2)^n\rfloor$ is odd or even infinitely often.




References

[Gr64e] Graham, R. L., On a conjecture of Erd\H{o}s in additive number theory . Acta Arith. (1964/65), 63-70.

\noindent\rule{\linewidth}{0.4pt}


%% ============================================================
\section{Problem \#352}
\label{prob:352}

\noindent\textbf{Status:} OPEN \quad|\quad \textbf{Prize:} none \quad|\quad \textbf{Updated:} 2025-08-31

\noindent\textbf{Tags:} geometry

\medskip

\noindent\textbf{Statement:}

Is there some $c&#62;0$ such that every measurable $A\subseteq \mathbb{R}^2$ of measure $\geq c$ contains the vertices of a triangle of area 1?




Erd\H{o}s (unpublished) proved that this is true if $A$ has infinite measure, or if $A$ is an unbounded set of positive measure (stating in \cite{Er78d} and \cite{Er83d} it 'follows easily from the Lebesgue density theorem'). In \cite{Er78d} and \cite{Er83d} he speculated that perhaps $C=4\pi/\sqrt{27}\approx 2.418$ works, which would be the best possible, as witnessed by a circle of radius $&lt;2\cdot 3^{-3/4}$. Further evidence for this is given by a result of Freiling and Mauldin \cite{Ma02}, who proved that if $A$ has outer measure $&gt;4\pi/\sqrt{27}$ then $A$ contains the vertices of a triangle with area $&#62;1$. This also proves the same threshold for the original problem under the assumption that $A$ is a compact convex set. Mauldin also discusses this problem in \cite{Ma13}, in which he notes that it suffices to prove this under the assumption that $A$ is the union of the interiors of $n&#60;\infty$ many compact convex sets. Freiling and Mauldin (see \cite{Ma13}) have proved this conjecture if $1\leq n\leq 3$.




References

[Er78d] Erd\H{o}s, P., Set-theoretic, measure-theoretic, combinatorial, and number-theoretic problems concerning point sets in Euclidean space . Real Anal. Exchange (1978/79), 113-138.

[Er83d] Erd\H{o}s, Paul, Some combinatorial, geometric and set theoretic problems in measure theory . Measure Theory, Oberwolfach 1983: Proceedings of the Conference held at Oberwolfach, June 26-July 2, 1983 (1984), 321-327.

[Ma02] Mauldin, R. D., Some problems in set theory, analysis and geometry . (2002), 493--506.

[Ma13] Mauldin, R. Daniel, Some problems and ideas of {E}rd\H{o}s in analysis and geometry . (2013), 365--376.

\noindent\rule{\linewidth}{0.4pt}


%% ============================================================
\section{Problem \#354}
\label{prob:354}

\noindent\textbf{Status:} OPEN \quad|\quad \textbf{Prize:} none \quad|\quad \textbf{Updated:} 2025-08-31

\noindent\textbf{Tags:} number theory

\medskip

\noindent\textbf{Statement:}

Let $\alpha,\beta\in \mathbb{R}_{&#62;0}$ such that $\alpha/\beta$ is irrational. Is the multiset\[\{ \lfloor \alpha\rfloor,\lfloor 2\alpha\rfloor,\lfloor 4\alpha\rfloor,\ldots\}\cup \{ \lfloor \beta\rfloor,\lfloor 2\beta\rfloor,\lfloor 4\beta\rfloor,\ldots\}\]complete? That is, can all sufficiently large natural numbers $n$ be written as\[n=\sum_{s\in S}\lfloor 2^s\alpha\rfloor+\sum_{t\in T}\lfloor 2^t\beta\rfloor\]for some finite $S,T\subset \mathbb{N}$? What if $2$ is replaced by some $\gamma\in(1,2)$?




This question was first mentioned by Graham \cite{Gr71}. Hegyv\'{a}ri \cite{He89} proved that this holds if $\alpha=m/2^n$ is a dyadic rational and $\beta$ is not. He later \cite{He91} proved that, for any fixed $\alpha&#62;0$, the set of $\beta$ for which this holds either has measure $0$ or infinite measure. In \cite{He94} he proved that the set of $(\alpha,\beta)$ for which the corresponding set of sums does not contain an infinite arithmetic progression has cardinality continuum. Hegyv\'{a}ri \cite{He89} proved that the sequence is not complete if $\alpha\geq 2$ and $\beta =2^k\alpha$ for some $k\geq 0$. Jiang and Ma \cite{JiMa24} and Fang and He \cite{FaHe25} prove that the sequence is not complete if $1&#60;\alpha&#60;2$ and $\beta=2^k\alpha$ for some sufficiently large $k$. It is likely (and Hegyv\'{a}ri conjectures) that the assumption $\alpha/\beta$ irrational can be weakened to $\alpha/\beta \neq 2^k$ and either $\alpha$ or $\beta$ not a dyadic rational. In the comments van Doorn proves the sequence is complete if $\alpha &#60; 2&#60;\beta&#60;3$, and also proves that if either $\alpha$ or $\beta$ is not a dyadic rational then the corresponding sequence with ceiling functions replacing the floor functions is complete.




References

[FaHe25] Fang, J.-H. and He, J.-Y., On a problem of {E}rd\H{o}s and {G}raham . Acta Math. Hungar. (2025), 532--542.

[Gr71] Graham, R. L., On sums of integers taken from a fixed sequence . (1971), 22--40.

[He89] Hegyv\'{a}ri, N., Some remarks on a problem of {E}rd\H{o}s and {G}raham . Acta Math. Hungar. (1989), 149--154.

[He91] Hegyv\'{a}ri, N., On complete sequences . Ann. Univ. Sci. Budapest. E\"otv\"os Sect. Math. (1991), 7--10.

[He94] Hegyv\'{a}ri, Norbert, On sumset of certain sets . Publ. Math. Debrecen (1994), 115--122.

[JiMa24] Jiang, Xing-Wang and Ma, Wu-Xia, A conjecture of {H}egyv\'{a}ri . Int. J. Number Theory (2024), 915--933.

\noindent\rule{\linewidth}{0.4pt}


%% ============================================================
\section{Problem \#357}
\label{prob:357}

\noindent\textbf{Status:} OPEN \quad|\quad \textbf{Prize:} none \quad|\quad \textbf{Updated:} 2025-08-31

\noindent\textbf{Tags:} number theory

\medskip

\noindent\textbf{Statement:}

Let $1\leq a_1&#60;\cdots &#60;a_k\leq n$ be integers such that all sums of the shape $\sum_{u\leq i\leq v}a_i$ are distinct. Let $f(n)$ be the maximal such $k$. How does $f(n)$ grow? Is $f(n)=o(n)$?




Asked by Erd\H{o}s and Harzheim. In \cite{Er77c} Erd\H{o}s asks about an infinite such set of integers, and whether such a set must have density $0$. He notes that a simple averaging process implies $a_k \gg k\log k$ for infinitely many $k$, and so the lower density is $0$. He also asks whether $\sum\frac{1}{a_k}$ must converge. Weisenberg in the comments observes that any set which satisfies [874] also has this property, which implies $f(n)\geq (2+o(1))n^{1/2}$. If $g(n)$ is the maximal $k$ such that there are $1\leq a_1,\ldots,a_k\leq n$ with all consecutive sums distinct (i.e. we drop the monotonicity assumption in the definition of $f$) then Hegyv\'{a}ri \cite{He86} has proved that\[\left(\frac{1}{3}+o(1)\right) n\leq g(n)\leq \left(\frac{2}{3}+o(1)\right)n.\]The upper bound of Coppersmith and Phillips in [867] implies\[g(n) \leq \left(\frac{2}{3}-\frac{1}{512}+o(1)\right)n.\]A similar question can be asked if we replace strict monotonicity with weak monotonicity (i.e. we allow $a_i=a_j$). Erd\H{o}s and Harzheim also ask what is the least $m$ which is not a sum of the given form? Can it be much larger than $n$? Erd\H{o}s and Harzheim can show that $\sum_{x&lt;a_i&lt;x^2}\frac{1}{a_i}\ll 1$. Is it true that $\sum_i \frac{1}{a_i}\ll 1$? See also [34] , [356] , [670] , and [867] . The multiplicative analogue is [421] .




References

[Er77c] Erd\H{o}s, Paul, Problems and results on combinatorial number theory. III . Number theory day (Proc. Conf., Rockefeller Univ., New York, 1976) (1977), 43-72.

[He86] Hegyv\'{a}ri, N., On consecutive sums in sequences . Acta Math. Hungar. (1986), 193--200.

\noindent\rule{\linewidth}{0.4pt}


%% ============================================================
\section{Problem \#359}
\label{prob:359}

\noindent\textbf{Status:} OPEN \quad|\quad \textbf{Prize:} none \quad|\quad \textbf{Updated:} 2025-08-31

\noindent\textbf{Tags:} number theory

\medskip
\noindent\textit{segmented numbers}


\noindent\textbf{Statement:}

Let $a_1&lt;a_2&lt;\cdots$ be an infinite sequence of integers such that $a_1=n$ and $a_{i+1}$ is the least integer which is not a sum of consecutive earlier $a_j$s. What can be said about the density of this sequence? In particular, in the case $n=1$, can one prove that $a_k/k\to \infty$ and $a_k/k^{1+c}\to 0$ for any $c&gt;0$?




A problem of MacMahon, studied by Andrews \cite{An75}. When $n=1$ this sequence begins\[1,2,4,5,8,10,14,15,\ldots.\]This sequence is A002048 in the OEIS. Andrews conjectures\[a_k\sim \frac{k\log k}{\log\log k}.\]Porubsky \cite{Po77} proved that, for any $\epsilon&gt;0$, there are infinitely many $k$ such that\[a_k &lt; (\log k)^\epsilon \frac{k\log k}{\log\log k},\]and also that if $A(x)$ counts the number of $a_i\leq x$ then\[\limsup \frac{A(x)}{\pi(x)}\geq \frac{1}{\log 2}\]where $\pi(x)$ counts the number of primes $\leq x$. See also [839] .




References

[An75] Andrews, George E., Research Problems: Mac Mahon's Prime Numbers of Measurement . Amer. Math. Monthly (1975), 922-923.

[Po77] Porubsk\'y, \v S., On {M}ac{M}ahon's segmented numbers and related sequences . Nieuw Arch. Wisk. (3) (1977), 403--408.

\noindent\rule{\linewidth}{0.4pt}


%% ============================================================
\section{Problem \#361}
\label{prob:361}

\noindent\textbf{Status:} OPEN \quad|\quad \textbf{Prize:} none \quad|\quad \textbf{Updated:} 2025-08-31

\noindent\textbf{Tags:} number theory

\medskip

\noindent\textbf{Statement:}

Let $c&#62;0$ and $n$ be some large integer. What is the size of the largest $A\subseteq \{1,\ldots,\lfloor cn\rfloor\}$ such that $n$ is not a sum of a subset of $A$? Does this depend on $n$ in an irregular way?

\noindent\rule{\linewidth}{0.4pt}


%% ============================================================
\section{Problem \#365}
\label{prob:365}

\noindent\textbf{Status:} OPEN \quad|\quad \textbf{Prize:} none \quad|\quad \textbf{Updated:} 2025-09-20

\noindent\textbf{Tags:} number theory

\medskip

\noindent\textbf{Statement:}

Do all pairs of consecutive powerful numbers $n$ and $n+1$ come from solutions to Pell equations ? In other words, must either $n$ or $n+1$ be a square? Is the number of such $n\leq x$ bounded by $(\log x)^{O(1)}$?




Erd\H{o}s originally asked Mahler whether there are infinitely many pairs of consecutive powerful numbers, but Mahler immediately observed that the answer is yes from the infinitely many solutions to the Pell equation $x^2=2^3y^2+1$. The list of $n$ such that $n$ and $n+1$ are both powerful is A060355 in the OEIS. The answer to the first question is no: Golomb \cite{Go70} observed that both $12167=23^3$ and $12168=2^33^213^2$ are powerful. Walker \cite{Wa76} proved that the equation\[7^3x^2=3^3y^2+1\]has infinitely many solutions, giving infinitely many counterexamples. See also [364] . This is discussed in problem B16 of Guy's collection \cite{Gu04}.




References

[Go70] Golomb, S. W., Powerful numbers . Amer. Math. Monthly (1970), 848-855.

[Gu04] Guy, Richard K., Unsolved problems in number theory . (2004), xviii+437.

[Wa76] Walker, David T., Consecutive integer pairs of powerful numbers and related Diophantine equations . Fibonacci Quart. (1976), 111-116.

\noindent\rule{\linewidth}{0.4pt}


%% ============================================================
\section{Problem \#367}
\label{prob:367}

\noindent\textbf{Status:} OPEN \quad|\quad \textbf{Prize:} none \quad|\quad \textbf{Updated:} 2025-08-31

\noindent\textbf{Tags:} number theory

\medskip

\noindent\textbf{Statement:}

Let $B_2(n)$ be the $2$-full part of $n$ (that is, $B_2(n)=n/n'$ where $n'$ is the product of all primes that divide $n$ exactly once). Is it true that, for every fixed $k\geq 1$,\[\prod_{n\leq m&#60;n+k}B_2(m) \ll n^{2+o(1)}?\]Or perhaps even $\ll_k n^2$?




It would also be interesting to find upper and lower bounds for the analogous product with $B_r$ for $r\geq 3$, where $B_r(n)$ is the $r$-full part of $n$ (that is, the product of prime powers $p^a \mid n$ such that $p^{a+1}\nmid n$ and $a\geq r$). Is it true that, for every fixed $r,k\geq 2$ and $\epsilon&gt;0$,\[\limsup \frac{\prod_{n\leq m&lt;n+k}B_r(m) }{n^{1+\epsilon}}\to\infty?\]van Doorn notes in the comments that for $k\leq 2$ we trivially have\[\prod_{n\leq m&lt;n+k}B_2(m) \ll n^{2},\]but that this fails for all $k\geq 3$, and in fact\[\prod_{n\leq m&lt;n+3}B_2(m) \gg n^{2}\log n\]infinitely often. This question is equivalent (up to constants) to [935] .

\noindent\rule{\linewidth}{0.4pt}


%% ============================================================
\section{Problem \#368}
\label{prob:368}

\noindent\textbf{Status:} OPEN \quad|\quad \textbf{Prize:} none \quad|\quad \textbf{Updated:} 2025-08-31

\noindent\textbf{Tags:} number theory

\medskip

\noindent\textbf{Statement:}

How large is the largest prime factor of $n(n+1)$?




Let $F(n)$ be the prime in question. P\'{o}lya \cite{Po18} proved that $F(n)\to \infty$ as $n\to\infty$. Mahler \cite{Ma35} showed that $F(n)\gg \log\log n$. Schinzel \cite{Sc67b} observed that for infinitely many $n$ we have $F(n)\leq n^{O(1/\log\log\log n)}$. The truth is probably $F(n)\gg (\log n)^2$ for all $n$. Erd\H{o}s \cite{Er76d} conjectured that, for every $\epsilon&gt;0$, there are infinitely many $n$ such that $F(n) &lt;(\log n)^{2+\epsilon}$. Pasten \cite{Pa24b} has proved that\[F(n) \gg \frac{(\log\log n)^2}{\log\log\log n}.\]The largest prime factors of $n(n+1)$ are listed as A074399 in the OEIS.




References

[Er76d] Erd\H{o}s, P., Problems and results on number theoretic properties of consecutive integers and related questions . Proceedings of the Fifth Manitoba Conference on Numerical Mathematics (Univ. Manitoba, Winnipeg, Man., 1975) (1976), 25-44.

[Ma35] Mahler, Kurt, \"{U}ber den gr\"{o}ssten Primteiler spezieller Polynome zweiten Grades . Archiv f\"{u}r math. og naturvid (1935).

[Pa24b] Pasten, Hector, The largest prime factor of {$n^2+1$} and improvements on subexponential {$ABC$} . Invent. Math. (2024), 373--385.

[Po18] P\'{o}lya, Georg, Zur arithmetischen {U}ntersuchung der {P}olynome . Math. Z. (1918), 143--148.

[Sc67b] Schinzel, A., On two theorems of Gelfond and some of their applications . Acta Arith. (1967/68), 177-236.

\noindent\rule{\linewidth}{0.4pt}


%% ============================================================
\section{Problem \#371}
\label{prob:371}

\noindent\textbf{Status:} OPEN \quad|\quad \textbf{Prize:} none \quad|\quad \textbf{Updated:} 2025-08-31

\noindent\textbf{Tags:} number theory

\medskip

\noindent\textbf{Statement:}

Let $P(n)$ denote the largest prime factor of $n$. Show that the set of $n$ with $P(n)&#60;P(n+1)$ has density $1/2$.




Conjectured by Erd\H{o}s and Pomerance \cite{ErPo78}, who proved that this set and its complement both have positive upper density. The best unconditional lower bound available is due to L\"{u} and Wang \cite{LuWa25}, who prove that\[\#\{ n&lt;x :P(n)&lt;P(n+1)\} &gt; (0.2017-o(1))x,\]and the same lower bound for the complement. In \cite{Er79e} Erd\H{o}s also asks whether, for every $\alpha$, the density of the set of $n$ where\[P(n+1)&gt;P(n)n^\alpha\]exists. Ter\"{a}v\"{a}inen \cite{Te18} has proved that the logarithmic density of the set of $n$ for which $P(n)&lt;P(n+1)$ is $1/2$. Tao and Ter\"{a}v\"{a}inen \cite{TaTe19} have proved that the asymptotic density is $1/2$ at 'almost all scales'. More generally, for any $0\leq \alpha \leq1$, Ter\"{a}v\"{a}inen \cite{Te18} proved that the logarithmic density of the set of $n$ for which $P(n+1)&gt;P(n)n^\alpha$ exists and is equal to\[\int_{[0,1]^2}1_{y\geq x+\alpha}u(x)u(y)\mathrm{d}x\mathrm{d}y\]where $u(x)=x^{-1}\rho(x^{-1}-1)$ and $\rho$ is the Dickman function . Wang \cite{Wa21} has proved the same value holds for the asymptotic density (and in particular provided an affirmative answer to the original question) conditional on the Elliott-Halberstam conjecture for friable integers. The sequence of such $n$ is A070089 in the OEIS. See also [372] and [928] .




References

[Er79e] Erd\H{o}s, Paul, Some unconventional problems in number theory . Ast\'{e}risque (1979), 73-82.

[ErPo78] Erd\H{o}s, Paul and Pomerance, Carl, On the largest prime factors of {$n$} and {$n+1$} . Aequationes Math. (1978), 311-321.

[LuWa25] L\"u, Xiaodong and Wang, Zhiwei, On the largest prime factors of consecutive integers . Monatsh. Math. (2025), 403--418.

[TaTe19] Tao, Terence and Ter\"{a}v\"{a}inen, Joni, The structure of correlations of multiplicative functions at almost all scales, with applications to the {C}howla and {E}lliott conjectures . Algebra Number Theory (2019), 2103--2150.

[Te18] Ter\"{a}v\"{a}inen, Joni, On binary correlations of multiplicative functions . Forum Math. Sigma (2018), Paper No. e10, 41.

[Wa21] Wang, Zhiwei, Three conjectures on {$P^+(n)$} and {$P^+(n+1)$} hold under the {E}lliott-{H}alberstam conjecture for friable integers . J. Number Theory (2021), 1--11.

\noindent\rule{\linewidth}{0.4pt}


%% ============================================================
\section{Problem \#373}
\label{prob:373}

\noindent\textbf{Status:} OPEN \quad|\quad \textbf{Prize:} none \quad|\quad \textbf{Updated:} 2025-08-31

\noindent\textbf{Tags:} number theory, factorials

\medskip

\noindent\textbf{Statement:}

Show that the equation\[n! = a_1!a_2!\cdots a_k!,\]with $n-1&#62;a_1\geq a_2\geq \cdots \geq a_k\geq 2$, has only finitely many solutions.




This would follow if $P(n(n+1))/\log n\to \infty$, where $P(m)$ denotes the largest prime factor of $m$ (see Problem [368] ). Erd\H{o}s \cite{Er76d} proved that this problem would also follow from showing that $P(n(n-1))&gt;4\log n$. The condition $a_1&lt;n-1$ is necessary to rule out the trivial solutions when $n=a_2!\cdots a_k!$. Sur\'{a}nyi was the first to conjecture that the only non-trivial solution to $a!b!=n!$ is $6!7!=10!$. More generally, Hickerson (as reported in \cite{Er76d}) conjectured that the only non-trivial solutions to the equation in the problem statement are\[9!=2!3!3!7!,\]\[10!=6!7!,\]\[10!=3!5!7!,\]and\[16!=14!5!2!.\]Luca \cite{Lu07b} has shown that there are only finitely many solutions, conditional on the ABC conjecture , and proved unconditionally that the number of $n\leq x$ which admit a non-trivial solution is\[\leq \exp \bigg(f(x)\frac{\log (x)}{\log\log (x)}\bigg)\]for any function $f(x)$ which tends to infinity. This is discussed in problem B23 of Guy's collection \cite{Gu04}. In the case when $k=2$, Erd\H{o}s \cite{Er93} proved that if $n!=a_1!a_2!$ with $n-1&gt;a_1\geq a_2$ then\[a_1\geq n-5\log\log n,\]and says it 'would be nice' to prove $a_1\geq n-o(\log\log n)$. Bhat and Ramachandra \cite{BhRa10} replace the $5$ with $(1+o(1))\frac{1}{\log 2}$, and also prove that the same bound holds for arbitrary $k\geq 2$. Numerical investigations on solutions to $n!=a_1!a_2!$ have been carried out by Caldwell \cite{Ca94} and Habsieger \cite{Ha19}, and it is known that there are no solutions aside from $10!=6!7!$ for $n\leq 10^{3000}$.




References

[BhRa10] Bhat, K. Dzh. and Ramachandra, K., A remark on factorials that are products of factorials . Mat. Zametki (2010), 350--354.

[Ca94] C. Caldwell, The Diophantine equation $A!B!=C!$ . J. Recreat. Math. (1994), 128-133.

[Er76d] Erd\H{o}s, P., Problems and results on number theoretic properties of consecutive integers and related questions . Proceedings of the Fifth Manitoba Conference on Numerical Mathematics (Univ. Manitoba, Winnipeg, Man., 1975) (1976), 25-44.

[Er93] Erd\H{o}s, Paul, Some of my favorite solved and unsolved problems in graph theory . Quaestiones Math. (1993), 333-350.

[Gu04] Guy, Richard K., Unsolved problems in number theory . (2004), xviii+437.

[Ha19] Habsieger, Laurent, Explicit bounds for the {D}iophantine equation {$A!B!=C!$} . Fibonacci Quart. (2019), 21--28.

[Lu07b] Luca, Florian, On factorials which are products of factorials . Math. Proc. Cambridge Philos. Soc. (2007), 533--542.

\noindent\rule{\linewidth}{0.4pt}


%% ============================================================
\section{Problem \#374}
\label{prob:374}

\noindent\textbf{Status:} OPEN \quad|\quad \textbf{Prize:} none \quad|\quad \textbf{Updated:} 2025-08-31

\noindent\textbf{Tags:} number theory

\medskip

\noindent\textbf{Statement:}

For any $m\in \mathbb{N}$, let $F(m)$ be the minimal $k\geq 2$ (if it exists) such that there are $a_1&#60;\cdots &#60;a_k=m$ with $a_1!\cdots a_k!$ a square. Let $D_k=\{ m : F(m)=k\}$. What is the order of growth of $\lvert D_k\cap\{1,\ldots,n\}\rvert$ for $3\leq k\leq 6$? For example, is it true that $\lvert D_6\cap \{1,\ldots,n\}\rvert \gg n$?




Studied by Erd\H{o}s and Graham \cite{ErGr76} (see also \cite{LSS14}). It is known, for example, that: {UL} {LI}no $D_k$ contains a prime,{/LI} {LI}$D_2=\{ n^2 : n&#62;1\}$,{/LI} {LI} $\lvert D_3\cap \{1,\ldots,n\}\rvert = o(\lvert D_4\cap \{1,\ldots,n\}\rvert)$,{/LI} {LI} the least element of $D_6$ is $527$, and{/LI} {LI} $D_k=\emptyset$ for $k&#62;6$.{/LI} {/UL}




References

[ErGr76] Erd\H{o}s, P. and Graham, R. L., On products of factorials . Bull. Inst. Math. Acad. Sinica (1976), 337-355.

[LSS14] Luca, F. and Saradha, N. and Shorey, T. N., Squares and factorials in products of factorials . Monatsh. Math. (2014), 385-400.

\noindent\rule{\linewidth}{0.4pt}


%% ============================================================
\section{Problem \#376}
\label{prob:376}

\noindent\textbf{Status:} OPEN \quad|\quad \textbf{Prize:} none \quad|\quad \textbf{Updated:} 2025-08-31

\noindent\textbf{Tags:} number theory, binomial coefficients, base representations

\medskip

\noindent\textbf{Statement:}

Are there infinitely many $n$ such that $\binom{2n}{n}$ is coprime to $105$?




Erd\H{o}s, Graham, Ruzsa, and Straus \cite{EGRS75} have shown that, for any two odd primes $p$ and $q$, there are infinitely many $n$ such that $\binom{2n}{n}$ is coprime to $pq$. This is equivalent (via Kummer's theorem ) to whether there are infinitely many $n$ which have only digits $0,1$ in base $3$, digits $0,1,2$ in base $5$, and digits $0,1,2,3$ in base $7$. The sequence of such $n$ is A030979 in the OEIS. The best result in this direction is due to Bloom and Croot \cite{BlCr25}, who proved that, if $p_1,p_2,p_3$ are sufficiently large primes, then there are infinitely many $n$ such that almost all of the base $p_i$ digits are $&lt;p_i/2$. In other words, for all $\epsilon&gt;0$, there are infinitely many $n$ such that $\binom{2n}{n}$ is coprime to $p_1p_2p_3$, except for a factor of size $\leq n^\epsilon$. This is mentioned in problem B33 of Guy's collection \cite{Gu04}. It is also discussed in an article of Pomerance \cite{Po15c}. Graham offered \$1000 for a solution to this problem (as mentioned in \cite{Gu04} and \cite{BeHa98}).




References

[BeHa98] Berend, Daniel and Harmse, J\o rgen E., On some arithmetical properties of middle binomial coefficients . Acta Arith. (1998), 31--41.

[BlCr25] T. F. Bloom and E. Croot, Integers with small digits in multiple bases . arXiv:2509.02835 (2025).

[EGRS75] Erd\H{o}s, P. and Graham, R. L. and Ruzsa, I. Z. and Straus, E. G., On the prime factors of $(\sp{2n}\sb{n})$ . Math. Comp. (1975), 83-92.

[Gu04] Guy, Richard K., Unsolved problems in number theory . (2004), xviii+437.

[Po15c] Pomerance, Carl, Divisors of the middle binomial coefficient . Amer. Math. Monthly (2015), 636--644.

\noindent\rule{\linewidth}{0.4pt}


%% ============================================================
\section{Problem \#377}
\label{prob:377}

\noindent\textbf{Status:} OPEN \quad|\quad \textbf{Prize:} none \quad|\quad \textbf{Updated:} 2025-08-31

\noindent\textbf{Tags:} number theory, binomial coefficients

\medskip

\noindent\textbf{Statement:}

Is there some absolute constant $C&#62;0$ such that\[\sum_{p\leq n}1_{p\nmid \binom{2n}{n}}\frac{1}{p}\leq C\]for all $n$ (where the summation is restricted to primes $p\leq n$)?




A question of Erd\H{o}s, Graham, Ruzsa, and Straus \cite{EGRS75}, who proved that if $f(n)$ is the sum in question then\[\lim_{x\to \infty}\frac{1}{x}\sum_{n\leq x}f(n) = \sum_{k=2}^\infty \frac{\log k}{2^k}=\gamma_0\]and\[\lim_{x\to \infty}\frac{1}{x}\sum_{n\leq x}f(n)^2 = \gamma_0^2,\]so that for almost all integers $f(m)=\gamma_0+o(1)$. They also prove that, for all large $n$,\[f(n) \leq c\log\log n\]for some constant $c&#60;1$. (It is trivial from Mertens estimates that $f(n)\leq (1+o(1))\log\log n$.) A positive answer would imply that\[\sum_{p\leq n}1_{p\mid \binom{2n}{n}}\frac{1}{p}=(1-o(1))\log\log n,\]and Erd\H{o}s, Graham, Ruzsa, and Straus say there is 'no doubt' this latter claim is true. This is mentioned in problem B33 of Guy's collection \cite{Gu04}.




References

[EGRS75] Erd\H{o}s, P. and Graham, R. L. and Ruzsa, I. Z. and Straus, E. G., On the prime factors of $(\sp{2n}\sb{n})$ . Math. Comp. (1975), 83-92.

[Gu04] Guy, Richard K., Unsolved problems in number theory . (2004), xviii+437.

\noindent\rule{\linewidth}{0.4pt}


%% ============================================================
\section{Problem \#382}
\label{prob:382}

\noindent\textbf{Status:} OPEN \quad|\quad \textbf{Prize:} none \quad|\quad \textbf{Updated:} 2025-08-31

\noindent\textbf{Tags:} number theory

\medskip

\noindent\textbf{Statement:}

Let $u\leq v$ be such that the largest prime dividing $\prod_{u\leq m\leq v}m$ appears with exponent at least $2$. Is it true that $v-u=v^{o(1)}$? Can $v-u$ be arbitrarily large?




Erd\H{o}s and Graham report it follows from results of Ramachandra that $v-u\leq v^{1/2+o(1)}$. Cambie has observed that the first question boils down to some old conjectures on prime gaps. By Cram\'{e r's conjecture}, for every $\epsilon&gt;0,$ for every $u$ sufficiently large there is a prime between $u$ and $u+u^\epsilon$. Thus for $u+u^\epsilon&lt;v$, the largest prime divisor of \( \prod_{u \leq m \leq v} m \) appears with exponent $1$. Since this is not the case in the question, \( v - u = v^{o(1)} \). Cambie also gives the following heuristic for the second question. The 'probability' that the largest prime divisor of $n$ is $&lt;n^{1/2}$ is $1-\log 2&gt;0$. For any fixed $k$, there is therefore a positive 'probability' that there are $k$ consecutive integers around $q^2$ (for a prime $q$) all of whose prime divisors are bounded above by $q$, when $v-u\geq k$. See [383] for a conjecture along these lines. A similar argument applies if we replace multiplicity $2$ with multiplicity $r$, for any fixed $r\geq 2$. See also [380] .

\noindent\rule{\linewidth}{0.4pt}


%% ============================================================
\section{Problem \#383}
\label{prob:383}

\noindent\textbf{Status:} OPEN \quad|\quad \textbf{Prize:} none \quad|\quad \textbf{Updated:} 2025-08-31

\noindent\textbf{Tags:} number theory

\medskip

\noindent\textbf{Statement:}

Is it true that for every $k$ there are infinitely many primes $p$ such that the largest prime divisor of\[\prod_{0\leq i\leq k}(p^2+i)\]is $p$?




A positive answer to this would give an answer to the second part of [382] . Heuristically, the 'probability' that $n$ has no prime divisors $\geq n^{1/2}$ is $1-\log 2&gt;0$, so standard heuristics predict the answer to this is yes.

\noindent\rule{\linewidth}{0.4pt}


%% ============================================================
\section{Problem \#385}
\label{prob:385}

\noindent\textbf{Status:} OPEN \quad|\quad \textbf{Prize:} none \quad|\quad \textbf{Updated:} 2025-08-31

\noindent\textbf{Tags:} number theory

\medskip

\noindent\textbf{Statement:}

Let\[F(n) = \max_{\substack{m&lt;n\\ m\textrm{ composite}}} m+p(m),\]where $p(m)$ is the least prime divisor of $m$. Is it true that $F(n)&gt;n$ for all sufficiently large $n$? Does $F(n)-n\to \infty$ as $n\to\infty$?




A question of Erd\H{o}s, Eggleton, and Selfridge, who write that 'plausible conjectures on primes' imply that $F(n)\leq n$ for only finitely many $n$, and in fact it is possible that this quantity is always at least $n+(1-o(1))\sqrt{n}$ (note that it is trivially $\leq n+\sqrt{n}$). Tao has discussed this problem in a blog post . Sarosh Adenwalla has observed that the first question is equivalent to [430] . Indeed, if $n$ is large and $a_i$ is the sequence defined in the latter problem, then [430] implies that there is a composite $a_j$ such that $a_j-p(a_j)&gt;n$ and hence $F(n)&gt;n$. See also [463] .

\noindent\rule{\linewidth}{0.4pt}


%% ============================================================
\section{Problem \#386}
\label{prob:386}

\noindent\textbf{Status:} OPEN \quad|\quad \textbf{Prize:} none \quad|\quad \textbf{Updated:} 2025-08-31

\noindent\textbf{Tags:} number theory, binomial coefficients

\medskip

\noindent\textbf{Statement:}

Let $2\leq k\leq n-2$. Can $\binom{n}{k}$ be the product of consecutive primes infinitely often? For example\[\binom{21}{2}=2\cdot 3\cdot 5\cdot 7.\]




Erd\H{o}s and Graham write that 'a proof that this cannot happen infinitely often for $\binom{n}{2}$ seems hopeless; probably this can never happen for $\binom{n}{k}$ if $3\leq k\leq n-3$.' Weisenberg has provided four easy examples that show Erd\H{o}s and Graham were too optimistic here:\[\binom{7}{3}=5\cdot 7,\]\[\binom{10}{4}= 2\cdot 3\cdot 5\cdot 7,\]\[\binom{14}{4} = 7\cdot 11\cdot 13,\]and\[\binom{15}{6}=5\cdot 7\cdot 11\cdot 13.\]The known values of $n$ for which $\binom{n}{2}$ is the product of consecutive primes are $4,6,15,21,715$ (see A280992 ).

\noindent\rule{\linewidth}{0.4pt}


%% ============================================================
\section{Problem \#387}
\label{prob:387}

\noindent\textbf{Status:} OPEN \quad|\quad \textbf{Prize:} none \quad|\quad \textbf{Updated:} 2025-08-31

\noindent\textbf{Tags:} number theory, binomial coefficients

\medskip

\noindent\textbf{Statement:}

Is there an absolute constant $c&#62;0$ such that, for all $1\leq k&#60; n$, the binomial coefficient $\binom{n}{k}$ has a divisor in $(cn,n]$?




Erd\H{o}s once conjectured that $\binom{n}{k}$ must always have a divisor in $(n-k,n]$, but this was disproved by Schinzel and Erd\H{o}s \cite{Sc58}. A counterexample is given by $n=99215$ and $k=15$. Schinzel conjectured (see problem B34 of \cite{Gu04}) that, for all sufficiently large $k$ which are not prime powers, there exists an $n$ such that $\binom{n}{k}$ is not divisible by any integer in $(n-k,n]$. In \cite{Er78g} Erd\H{o}s thought that this question had a negative answer. It is easy to see that $\binom{n}{k}$ always has a divisor in $[n/k,n]$. Faulkner \cite{Fa66} proved that, if $p$ is the least prime $&#62;2k$ and $n\geq p$, then $\binom{n}{k}$ has a prime divisor $\geq p$ (except $\binom{9}{2}$ and $\binom{10}{3}$). This is discussed in problems B33 and B34 of Guy's collection \cite{Gu04}, who says that Erd\H{o}s conjectured this is true for any $c&#60;1$ (if $n$ is sufficiently large).




References

[Er78g] Erd\H{o}s, P\'{a}l, On prime factors of binomial coefficients. {II} . Mat. Lapok (1978/82), 307--316.

[Fa66] Faulkner, M., On a theorem of Sylvester and Schur . J. London Math. Soc. (1966), 107--110.

[Gu04] Guy, Richard K., Unsolved problems in number theory . (2004), xviii+437.

[Sc58] Schinzel, A., Sur un probl\`eme de {P}. {E}rd\H{o}s . Colloq. Math. (1958), 198--204.

\noindent\rule{\linewidth}{0.4pt}


%% ============================================================
\section{Problem \#388}
\label{prob:388}

\noindent\textbf{Status:} OPEN \quad|\quad \textbf{Prize:} none \quad|\quad \textbf{Updated:} 2025-08-31

\noindent\textbf{Tags:} number theory

\medskip

\noindent\textbf{Statement:}

Can one classify all solutions of\[\prod_{1\leq i\leq k_1}(m_1+i)=\prod_{1\leq j\leq k_2}(m_2+j)\]where $k_1,k_2&#62;3$ and $m_1+k_1\leq m_2$? Are there only finitely many solutions?




More generally, if $k_1&gt;2$ then for fixed $a$ and $b$\[a\prod_{1\leq i\leq k_1}(m_1+i)=b\prod_{1\leq j\leq k_2}(m_2+j)\]should have only a finite number of solutions. See also [363] , [686] , and [931] .

\noindent\rule{\linewidth}{0.4pt}


%% ============================================================
\section{Problem \#389}
\label{prob:389}

\noindent\textbf{Status:} OPEN \quad|\quad \textbf{Prize:} none \quad|\quad \textbf{Updated:} 2025-08-31

\noindent\textbf{Tags:} number theory

\medskip

\noindent\textbf{Statement:}

Is it true that for every $n\geq 1$ there is a $k$ such that\[n(n+1)\cdots(n+k-1)\mid (n+k)\cdots (n+2k-1)?\]




Asked by Erd\H{o}s and Straus. For example when $n=2$ we have $k=5$:\[2\times 3 \times 4 \times 5\times 6 \mid 7 \times 8 \times 9\times 10\times 11.\]and when $n=3$ we have $k=4$:\[3\times 4\times 5\times 6 \mid 7\times 8\times 9\times 10.\]Bhavik Mehta has computed the minimal such $k$ for $1\leq n\leq 18$ (now available as A375071 on the OEIS).

\noindent\rule{\linewidth}{0.4pt}


%% ============================================================
\section{Problem \#390}
\label{prob:390}

\noindent\textbf{Status:} OPEN \quad|\quad \textbf{Prize:} none \quad|\quad \textbf{Updated:} 2025-08-31

\noindent\textbf{Tags:} number theory, factorials

\medskip

\noindent\textbf{Statement:}

Let $f(n)$ be the minimal $m$ such that\[n! = a_1\cdots a_k\]with $n&#60; a_1&#60;\cdots &#60;a_k=m$. Is there (and what is it) a constant $c$ such that\[f(n)-2n \sim c\frac{n}{\log n}?\]




Erd\H{o}s, Guy, and Selfridge \cite{EGS82} have shown that\[f(n)-2n \asymp \frac{n}{\log n}.\]




References

[EGS82] Erd\H{o}s, P. and Guy, R. K. and Selfridge, J. L., Another property of {$239$} and some related questions . Congr. Numer. (1982), 243-257.

\noindent\rule{\linewidth}{0.4pt}


%% ============================================================
\section{Problem \#393}
\label{prob:393}

\noindent\textbf{Status:} OPEN \quad|\quad \textbf{Prize:} none \quad|\quad \textbf{Updated:} 2025-08-31

\noindent\textbf{Tags:} number theory, factorials

\medskip

\noindent\textbf{Statement:}

Let $f(n)$ denote the minimal $m\geq 1$ such that\[n! = a_1\cdots a_t\]with $a_1&#60;\cdots &#60;a_t=a_1+m$. What is the behaviour of $f(n)$?




Erd\H{o}s and Graham write that they do not even know whether $f(n)=1$ infinitely often (i.e. whether a factorial is the product of two consecutive integers infinitely often). Let $F_m(N)$ count the number of $n\leq N$ such that $f(n)=m$. Berend and Osgood \cite{BeOs92} proved that, for each fixed $m$, $F_m(N)=o(N)$. Bui, Pratt, and Zaharescu \cite{BPZ23} have shown that\[F_m(N)\ll_m N^{33/34}.\]A result of Luca \cite{Lu02} implies that $f(n)\to \infty$ as $n\to \infty$, conditional on the ABC conjecture .




References

[BPZ23] Bui, Hung M. and Pratt, Kyle and Zaharescu, Alexandru, Power savings for counting solutions to polynomial-factorial equations . Adv. Math. (2023), Paper No. 109021, 32.

[BeOs92] Berend, Daniel and Osgood, Charles F., On the equation {$P(x)=n!$} and a question of {E}rd\H{o}s . J. Number Theory (1992), 189--193.

[Lu02] Luca, Florian, The {D}iophantine equation {$P(x)=n!$} and a result of {M}. {O}verholt . Glas. Mat. Ser. III (2002), 269--273.

\noindent\rule{\linewidth}{0.4pt}


%% ============================================================
\section{Problem \#394}
\label{prob:394}

\noindent\textbf{Status:} OPEN \quad|\quad \textbf{Prize:} none \quad|\quad \textbf{Updated:} 2025-10-28

\noindent\textbf{Tags:} number theory

\medskip

\noindent\textbf{Statement:}

Let $t_k(n)$ denote the least $m$ such that\[n\mid m(m+1)(m+2)\cdots (m+k-1).\]Is it true that\[\sum_{n\leq x}t_2(n)\ll \frac{x^2}{(\log x)^c}\]for some $c&#62;0$? Is it true that, for $k\geq 2$,\[\sum_{n\leq x}t_{k+1}(n) =o\left(\sum_{n\leq x}t_k(n)\right)?\]




In \cite{ErGr80} they mention a conjecture of Erd\H{o}s that the sum is $o(x^2)$. This was proved by Erd\H{o}s and Hall \cite{ErHa78}, who proved that in fact\[\sum_{n\leq x}t_2(n)\ll \frac{\log\log\log x}{\log\log x}x^2.\]Erd\H{o}s and Hall conjecture that the sum is $o(x^2/(\log x)^c)$ for any $c&#60;\log 2$. Since $t_2(p)=p-1$ for prime $p$ it is trivial that\[\sum_{n\leq x}t_2(n)\gg \frac{x^2}{\log x}.\]Erd\H{o}s and Hall \cite{ErHa78} also note that $t_{n-1}(n!)=2$ and $t_{n-2}(n!)\ll n$, which $n=2^r$ shows is the best possible. They ask about the behaviour of $t_{n-3}(n!)$ and also ask ask whether, for infinitely many $n$,\[t_k(n!)&#60; t_{k-1}(n!)-1\]for all $1\leq k&#60;n$. They proved (with Selfridge) that this holds for $n=10$.




References

[ErGr80] Erd\H{o}s, P. and Graham, R., Old and new problems and results in combinatorial number theory . Monographies de L'Enseignement Mathematique (1980).

[ErHa78] Erd\H{o}s, P. and Hall, R. R., On some unconventional problems on the divisors of integers . J. Austral. Math. Soc. Ser. A (1978), 479--485.

\noindent\rule{\linewidth}{0.4pt}


%% ============================================================
\section{Problem \#396}
\label{prob:396}

\noindent\textbf{Status:} OPEN \quad|\quad \textbf{Prize:} none \quad|\quad \textbf{Updated:} 2025-08-31

\noindent\textbf{Tags:} number theory, binomial coefficients

\medskip

\noindent\textbf{Statement:}

Is it true that for every $k$ there exists $n$ such that\[\prod_{0\leq i\leq k}(n-i) \mid \binom{2n}{n}?\]




Erd\H{o}s and Graham write that $n+1$ always divides $\binom{2n}{n}$ (indeed $\frac{1}{n+1}\binom{2n}{n}$ is the $n$th Catalan number ), but it is quite rare that $n$ divides $\binom{2n}{n}$. Pomerance \cite{Po14} has shown that for any $k\geq 0$ there are infinitely many $n$ such that $n-k\mid\binom{2n}{n}$, although the set of such $n$ has upper density $&lt;1/3$. Pomerance also shows that the set of $n$ such that\[\prod_{1\leq i\leq k}(n+i)\mid \binom{2n}{n}\]has density $1$. The smallest $n$ for each $k$ are listed as A375077 on the OEIS.




References

[Po14] Pomerance, C., Divisors of the middle binomial coefficient . American Mathematical Monthly (2014).

\noindent\rule{\linewidth}{0.4pt}


%% ============================================================
\section{Problem \#400}
\label{prob:400}

\noindent\textbf{Status:} OPEN \quad|\quad \textbf{Prize:} none \quad|\quad \textbf{Updated:} 2025-08-31

\noindent\textbf{Tags:} number theory, factorials

\medskip

\noindent\textbf{Statement:}

For any $k\geq 2$ let $g_k(n)$ denote the maximum value of\[(a_1+\cdots+a_k)-n\]where $a_1,\ldots,a_k$ are integers such that $a_1!\cdots a_k! \mid n!$. Can one show that\[\sum_{n\leq x}g_k(n) \sim c_k x\log x\]for some constant $c_k$? Is it true that there is a constant $c_k$ such that for almost all $n&#60;x$ we have\[g_k(n)=c_k\log x+o(\log x)?\]




Erd\H{o}s and Graham write that it is easy to show that $g_k(n) \ll_k \log n$ always, but the best possible constant is unknown. See also [401] .

\noindent\rule{\linewidth}{0.4pt}


%% ============================================================
\section{Problem \#404}
\label{prob:404}

\noindent\textbf{Status:} OPEN \quad|\quad \textbf{Prize:} none \quad|\quad \textbf{Updated:} 2025-08-31

\noindent\textbf{Tags:} number theory, factorials

\medskip

\noindent\textbf{Statement:}

For which integers $a\geq 1$ and primes $p$ is there a finite upper bound on those $k$ such that there are $a=a_1&#60;\cdots&#60;a_n$ with\[p^k \mid (a_1!+\cdots+a_n!)?\]If $f(a,p)$ is the greatest such $k$, how does this function behave? Is there a prime $p$ and an infinite sequence $a_1&#60;a_2&#60;\cdots$ such that if $p^{m_k}$ is the highest power of $p$ dividing $\sum_{i\leq k}a_i!$ then $m_k\to \infty$?




See also [403] . Lin \cite{Li76} has shown that $f(2,2) \leq 254$.




References

[Li76] Lin, S., On two problems of Erd\H{o}s concerning sums of distinct factorials . Bell Laboratories internal memorandum (1960).

\noindent\rule{\linewidth}{0.4pt}


%% ============================================================
\section{Problem \#406}
\label{prob:406}

\noindent\textbf{Status:} OPEN \quad|\quad \textbf{Prize:} none \quad|\quad \textbf{Updated:} 2025-08-31

\noindent\textbf{Tags:} number theory, base representations

\medskip

\noindent\textbf{Statement:}

Is it true that there are only finitely many powers of $2$ which have only the digits $0$ and $1$ when written in base $3$?




The only examples seem to be $1$, $4=1+3$, and $256=1+3+3^2+3^5$. If we only allow the digits $1$ and $2$ then $2^{15}$ seems to be the largest such power of $2$. This would imply via Kummer's theorem that\[3\mid \binom{2^{k+1}}{2^k}\]for all large $k$. Saye \cite{Sa22} has computed that $2^n$ contains every possible ternary digit for $16\leq n \leq 5.9\times 10^{21}$. Let $N(x)$ count the number of $n\leq x$ such that $2^n$ has only the digits $0$ and $1$ in base $3$. Narkiewicz \cite{Na80} proved\[N(x)\leq 1.62 x^{\log_32}\]This is mentioned in problem B33 of Guy's collection \cite{Gu04}. There are many generalisations possible - see, for example, \cite{AbLa14} and \cite{La09}.




References

[AbLa14] Abram, William C. and Lagarias, Jeffrey C., Intersections of multiplicative translates of 3-adic {C}antor sets . J. Fractal Geom. (2014), 349--390.

[Gu04] Guy, Richard K., Unsolved problems in number theory . (2004), xviii+437.

[La09] Lagarias, Jeffrey C., Ternary expansions of powers of 2 . J. Lond. Math. Soc. (2) (2009), 562--588.

[Na80] Narkiewicz, W., A note on a paper of {H}. {G}upta concerning powers of two and three: ``{P}owers of {$2$}\ and sums of distinct powers of {$3$}''\ [{U}niv. {B}eograd. {P}ubl. {E}lektrotehn. {F}ak. Ser. {M}at. {F}iz. {N}o. 602-633 (1978), 151--158 (1979);\ {MR} 81g:10016] . Univ. Beograd. Publ. Elektrotehn. Fak. Ser. Mat. Fiz. (1980), 173--174.

[Sa22] Saye, Robert I., On two conjectures concerning the ternary digits of powers of two . J. Integer Seq. (2022), Art. 22.3.4, 9.

\noindent\rule{\linewidth}{0.4pt}


%% ============================================================
\section{Problem \#408}
\label{prob:408}

\noindent\textbf{Status:} OPEN \quad|\quad \textbf{Prize:} none \quad|\quad \textbf{Updated:} 2025-08-31

\noindent\textbf{Tags:} number theory, iterated functions

\medskip

\noindent\textbf{Statement:}

Let $\phi(n)$ be the Euler totient function and $\phi_k(n)$ be the iterated $\phi$ function, so that $\phi_1(n)=\phi(n)$ and $\phi_k(n)=\phi(\phi_{k-1}(n))$. Let\[f(n) = \min \{ k : \phi_k(n)=1\}.\]Does $f(n)/\log n$ have a distribution function? Is $f(n)/\log n$ almost always constant? What can be said about the largest prime factor of $\phi_k(n)$ when, say, $k=\log\log n$?




Pillai \cite{Pi29} was the first to investigate this function, and proved\[\log_3 n &#60; f(n) &#60; \log_2 n\]for all large $n$. Shapiro \cite{Sh50} proved that $f(n)$ is essentially multiplicative. Erd\H{o}s, Granville, Pomerance, and Spiro \cite{EGPS90} have proved that the answer to the first two questions is yes, conditional on a form of the Elliott-Halberstam conjecture. It is likely true that, if $k\to \infty$ however slowly with $n$, then for almost all $n$ the largest prime factor of $\phi_k(n)$ is $\leq n^{o(1)}$. This is discussed in problem B41 of Guy's collection \cite{Gu04}.




References

[EGPS90] Erd\H{o}s, P. and Granville, A. and Pomerance, C. and Spiro, C., On the normal behavior of the iterates of some arithmetic functions . Analytic number theory (Allerton Park, IL, 1989) (1990), 165-204.

[Gu04] Guy, Richard K., Unsolved problems in number theory . (2004), xviii+437.

[Pi29] Pillai, S. Sivasankaranarayana, On some functions connected with {$\phi(n)$} . Bull. Amer. Math. Soc. (1929), 832-836.

[Sh50] Shapiro, Harold N., On the iterates of a certain class of arithmetic functions . Comm. Pure Appl. Math. (1950), 259-272.

\noindent\rule{\linewidth}{0.4pt}


%% ============================================================
\section{Problem \#409}
\label{prob:409}

\noindent\textbf{Status:} OPEN \quad|\quad \textbf{Prize:} none \quad|\quad \textbf{Updated:} 2025-08-31

\noindent\textbf{Tags:} number theory, iterated functions

\medskip

\noindent\textbf{Statement:}

How many iterations of $n\mapsto \phi(n)+1$ are needed before a prime is reached? Can infinitely many $n$ reach the same prime? What is the density of $n$ which reach any fixed prime?




A problem of Finucane. One can also ask similar questions about $n\mapsto \sigma(n)-1$: do iterates of this always reach a prime? If so, how soon? (It is easily seen that iterates of this cannot reach the same prime infinitely often, since they are non-decreasing.) This problem is somewhat ambiguously phrased. Let $F(n)$ count the number of iterations of $n\mapsto \phi(n)+1$ before reaching a prime. The number of iterations required is A039651 in the OEIS. Cambie notes in the comments that $F(n)=o(n)$ is trivial, and $F(n)=1$ infinitely often. Presumably the intended question is to find 'good' upper bounds for $F(n)$. This is discussed in problem B41 of Guy's collection \cite{Gu04}.




References

[Gu04] Guy, Richard K., Unsolved problems in number theory . (2004), xviii+437.

\noindent\rule{\linewidth}{0.4pt}


%% ============================================================
\section{Problem \#410}
\label{prob:410}

\noindent\textbf{Status:} OPEN \quad|\quad \textbf{Prize:} none \quad|\quad \textbf{Updated:} 2025-08-31

\noindent\textbf{Tags:} number theory, iterated functions

\medskip

\noindent\textbf{Statement:}

Let $\sigma_1(n)=\sigma(n)$, the sum of divisors function, and $\sigma_k(n)=\sigma(\sigma_{k-1}(n))$. Is it true that for all $n\geq 2$\[\lim_{k\to \infty} \sigma_k(n)^{1/k}=\infty?\]




This is discussed in problem B9 of Guy's collection \cite{Gu04}.




References

[Gu04] Guy, Richard K., Unsolved problems in number theory . (2004), xviii+437.

\noindent\rule{\linewidth}{0.4pt}


%% ============================================================
\section{Problem \#411}
\label{prob:411}

\noindent\textbf{Status:} OPEN \quad|\quad \textbf{Prize:} none \quad|\quad \textbf{Updated:} 2025-08-31

\noindent\textbf{Tags:} number theory, iterated functions

\medskip

\noindent\textbf{Statement:}

Let $g_1=g(n)=n+\phi(n)$ and $g_k(n)=g(g_{k-1}(n))$. For which $n$ and $r$ is it true that $g_{k+r}(n)=2g_k(n)$ for all large $k$?




The known solutions to $g_{k+2}(n)=2g_k(n)$ are $n=10$ and $n=94$. Selfridge and Weintraub found solutions to $g_{k+9}(n)=9g_k(n)$ and Weintraub found\[g_{k+25}(3114)=729g_k(3114)\]for all $k\geq 6$. Steinerberger \cite{St25} has observed that, for $r=2$, this problem is equivalent to asking for solutions to\[\phi(n)+\phi(n+\phi(n))=n,\]and has shown that if this holds then either the odd part of $n$ is in $\{1,3,5,7,35,47\}$, or is equal to $8m+7$ or $6m+5$, where $8m+7\geq 10^{10}$ is a prime number and $\phi(6m+5)=4m+4$. Whether there are infinitely many such $m$ is related to the question of whether\[\phi(n)=\frac{2}{3}(n+1)\]has infinitely many solutions. Cambie conjectures that the only solutions have $r=2$ and $n=2^lp$ for some $l\geq 1$ and $p\in \{2,3,5,7,35,47\}$. Cambie has shown this problem is reducible to the question of which integers $r,t\geq 1$ and primes $p\equiv 7\pmod{8}$ satisfy $g_k(2p^t)=4p^t$, and conjectures there are no solutions to this except when $t=1$ and $p\in \{7,47\}$. Cambie has also observed that\[g_{k+4}(738)=3g_k(738),\]\[g_{k+4}(148646)=4g_k(148646),\]and\[g_{k+4}(4325798)=4g_{k}(4325798)\]for all $k\geq 1$.




References

[St25] S. Steinerberger, On an iterated arithmetic function problem of Erd\H{o}s and Graham . arXiv:2504.08023 (2025).

\noindent\rule{\linewidth}{0.4pt}


%% ============================================================
\section{Problem \#412}
\label{prob:412}

\noindent\textbf{Status:} OPEN \quad|\quad \textbf{Prize:} none \quad|\quad \textbf{Updated:} 2025-08-31

\noindent\textbf{Tags:} number theory, iterated functions

\medskip

\noindent\textbf{Statement:}

Let $\sigma_1(n)=\sigma(n)$, the sum of divisors function, and $\sigma_k(n)=\sigma(\sigma_{k-1}(n))$. Is it true that, for every $m,n\geq 2$, there exist some $i,j$ such that $\sigma_i(m)=\sigma_j(n)$?




In \cite{Er79d} Erd\H{o}s attributes this conjecture to van Wijngaarden, who told it to Erd\H{o}s in the 1950s. That is, there is (eventually) only one possible sequence that the iterated sum of divisors function can settle on. Selfridge reports numerical evidence which suggests the answer is no, but Erd\H{o}s and Graham write 'it seems unlikely that anything can be proved about this in the near future'. See also [413] and [414] .




References

[Er79d] Erd\H{o}s, P., Some unconventional problems in number theory . Acta Math. Acad. Sci. Hungar. (1979), 71-80.

\noindent\rule{\linewidth}{0.4pt}


%% ============================================================
\section{Problem \#413}
\label{prob:413}

\noindent\textbf{Status:} OPEN \quad|\quad \textbf{Prize:} none \quad|\quad \textbf{Updated:} 2025-08-31

\noindent\textbf{Tags:} number theory, iterated functions

\medskip

\noindent\textbf{Statement:}

Let $\omega(n)$ count the number of distinct primes dividing $n$. Are there infinitely many $n$ such that, for all $m&lt;n$, we have $m+\omega(m) \leq n$? Can one show that there exists an $\epsilon&gt;0$ such that there are infinitely many $n$ where $m+\epsilon \omega(m)\leq n$ for all $m&#60;n$?




In \cite{Er79} Erd\H{o}s calls such an $n$ a 'barrier' for $\omega$. Some other natural number theoretic functions (such as $\phi$ and $\sigma$) have no barriers because they increase too rapidly. Erd\H{o}s believed that $\omega$ should have infinitely many barriers. In \cite{Er79d} he proves that $F(n)=\prod k_i$, where $n=\prod p_i^{k_i}$, has infinitely many barriers (in fact the set of barriers has positive density in the integers). Erd\H{o}s also believed that $\Omega$, the count of the number of prime factors with multiplicity), should have infinitely many barriers. Selfridge found the largest barrier for $\Omega$ which is $&lt;10^5$ is $99840$. In \cite{ErGr80} this problem is suggested as a way of showing that the iterated behaviour of $n\mapsto n+\omega(n)$ eventually settles into a single sequence, regardless of the starting value of $n$ (see also [412] and [414] ). Erd\H{o}s and Graham report it could be attacked by sieve methods, but 'at present these methods are not strong enough'. Lau \cite{La26} (see [679] ) has provided a positive answer to the second question, and proved a weaker version of the first: there exists a constant $C$ such that, for infinitely many $n$,\[m+\omega(m) \leq n\]for all $1\leq m\leq n-C$. The sequence of barriers for $\omega$ is A005236 in the OEIS. This is discussed in problem B8 of Guy's collection \cite{Gu04}. See also [647] and [679]




References

[Er79] Erd\H{o}s, Paul, Some unconventional problems in number theory . Math. Mag. (1979), 67-70.

[Er79d] Erd\H{o}s, P., Some unconventional problems in number theory . Acta Math. Acad. Sci. Hungar. (1979), 71-80.

[ErGr80] Erd\H{o}s, P. and Graham, R., Old and new problems and results in combinatorial number theory . Monographies de L'Enseignement Mathematique (1980).

[Gu04] Guy, Richard K., Unsolved problems in number theory . (2004), xviii+437.

[La26] C. F. Lau, On the number of prime factors of consecutive integers . arXiv:2604.15042 (2026).

\noindent\rule{\linewidth}{0.4pt}


%% ============================================================
\section{Problem \#414}
\label{prob:414}

\noindent\textbf{Status:} OPEN \quad|\quad \textbf{Prize:} none \quad|\quad \textbf{Updated:} 2025-08-31

\noindent\textbf{Tags:} number theory, iterated functions

\medskip

\noindent\textbf{Statement:}

Let $h_1(n)=h(n)=n+\tau(n)$ (where $\tau(n)$ counts the number of divisors of $n$) and $h_k(n)=h(h_{k-1}(n))$. Is it true, for any $m,n$, there exist $i$ and $j$ such that $h_i(m)=h_j(n)$?




Asked by Spiro. That is, there is (eventually) only one possible sequence that the iterations of $n\mapsto h(n)$ can settle on. Erd\H{o}s and Graham believed the answer is yes. Similar questions can be asked by the iterates of many other functions. See also [412] and [413] .

\noindent\rule{\linewidth}{0.4pt}


%% ============================================================
\section{Problem \#415}
\label{prob:415}

\noindent\textbf{Status:} OPEN \quad|\quad \textbf{Prize:} none \quad|\quad \textbf{Updated:} 2025-08-31

\noindent\textbf{Tags:} number theory

\medskip

\noindent\textbf{Statement:}

For any $n$ let $F(n)$ be the largest $k$ such that any of the $k!$ possible ordering patterns appears in some sequence of $\phi(m+1),\ldots,\phi(m+k)$ with $m+k\leq n$. Is it true that\[F(n)=(c+o(1))\log\log\log n\]for some constant $c$? Is the first pattern which fails to appear always\[\phi(m+1)&#62;\phi(m+2)&#62;\cdots &#62;\phi(m+k)?\]Is it true that the 'natural' ordering which mimics what happens to $\phi(1),\ldots,\phi(k)$ is the most likely to appear?




In \cite{ErGr80} it is left ambiguous whether an 'ordering pattern' involves strict inequalities or allows equality. Weisenberg has observed that the final question only makes sense if we allow equality. In \cite{ErGr80} it is claimed that Erd\H{o}s \cite{Er36b} proved that\[F(n)\asymp \log\log\log n,\]and similarly if we replace $\phi$ with $\sigma$ or $\tau$ or $\nu$ or any 'decent' additive or multiplicative function, but this does not seem to be present in \cite{Er36b}, which only shows that the number of $m\leq n$ such that $\phi(m)&lt;\phi(m+1)$ is $\sim n/2$, and similarly for $\phi(m)&gt;\phi(m+1)$. The following asymptotic shows that such a result is false regardless. Pollack, Pomerance, and Trevi\~{n}o \cite{PPT13} have proved that, if $G(n)$ is the maximum $k$ such that\[\phi(m+1)&gt; \phi(m+2)&gt;\cdots &gt; \phi(m+k)\]for some $m$ with $m+k\leq n$, then\[G(n)\sim \frac{\log\log\log n}{\log\log\log\log\log\log n}.\]The same asymptotic holds if we replace $&gt;$ with $&lt;,\leq,\geq$. Since $F(n)\leq G(n)$ trivially, this answers the first question in the negative. Chojecki and GPT-5.4 have sketched how the proof of \cite{PPT13} can be adapted to obtain the same asymptotic for an arbitrary pattern of (strict) inequalities. Note that a positive answer to the final question would imply that there exist infinitely many $n$ with $\phi(n)=\phi(n+1)$, which is the subject of [1003] .




References

[Er36b] Erd\H{o}s, P., On a problem of Chowla and some related problems . Proc. Cambridge Philos. Soc. (1936), 530-540.

[ErGr80] Erd\H{o}s, P. and Graham, R., Old and new problems and results in combinatorial number theory . Monographies de L'Enseignement Mathematique (1980).

[PPT13] Pollack, Paul and Pomerance, Carl and Trevi\~no, Enrique, Sets of monotonicity for {E}uler's totient function . Ramanujan J. (2013), 379--398.

\noindent\rule{\linewidth}{0.4pt}


%% ============================================================
\section{Problem \#416}
\label{prob:416}

\noindent\textbf{Status:} OPEN \quad|\quad \textbf{Prize:} none \quad|\quad \textbf{Updated:} 2025-08-31

\noindent\textbf{Tags:} number theory

\medskip

\noindent\textbf{Statement:}

Let $V(x)$ count the number of $n\leq x$ such that $\phi(m)=n$ is solvable. Does $V(2x)/V(x)\to 2$? Is there an asymptotic formula for $V(x)$?




Pillai \cite{Pi29} proved $V(x)=o(x)$. Erd\H{o}s \cite{Er35b} proved $V(x)=x(\log x)^{-1+o(1)}$. The behaviour of $V(x)$ is now almost completely understood. Maier and Pomerance \cite{MaPo88} proved\[V(x)=\frac{x}{\log x}e^{(C+o(1))(\log\log\log x)^2},\]for some explicit constant $C&gt;0$. Ford \cite{Fo98} improved this to\[V(x)\asymp\frac{x}{\log x}e^{C_1(\log\log\log x-\log\log\log\log x)^2+C_2\log\log\log x-C_3\log\log\log\log x}\]for some explicit constants $C_1,C_2,C_3&gt;0$. Unfortunately this falls just short of an asymptotic formula for $V(x)$ and determining whether $V(2x)/V(x)\to 2$. In \cite{Er79e} Erd\H{o}s asks further to estimate the number of $n\leq x$ such that the smallest solution to $\phi(m)=n$ satisfies $kx&lt;m\leq (k+1)x$. See also [417] and [821] . This is discussed in problem B36 of Guy's collection \cite{Gu04}.




References

[Er35b] Erd\H{o}s, P., On the normal number of prime factors of $p-1$ and some related problems concerning Euler's $\varphi$-function . Quart. J. Math. (1935), 205-213.

[Er79e] Erd\H{o}s, Paul, Some unconventional problems in number theory . Ast\'{e}risque (1979), 73-82.

[Fo98] Ford, Kevin, The distribution of totients . Ramanujan J. (1998), 67-151.

[Gu04] Guy, Richard K., Unsolved problems in number theory . (2004), xviii+437.

[MaPo88] Maier, Helmut and Pomerance, Carl, On the number of distinct values of Euler's $\phi$-function . Acta Arith. (1988), 263-275.

[Pi29] Pillai, S. Sivasankaranarayana, On some functions connected with {$\phi(n)$} . Bull. Amer. Math. Soc. (1929), 832-836.

\noindent\rule{\linewidth}{0.4pt}


%% ============================================================
\section{Problem \#417}
\label{prob:417}

\noindent\textbf{Status:} OPEN \quad|\quad \textbf{Prize:} none \quad|\quad \textbf{Updated:} 2025-08-31

\noindent\textbf{Tags:} number theory

\medskip

\noindent\textbf{Statement:}

Let\[V'(x)=\#\{\phi(m) : 1\leq m\leq x\}\]and\[V(x)=\#\{\phi(m) \leq x : 1\leq m\}.\]Does $\lim V(x)/V'(x)$ exist? Is it $&#62;1$?




It is trivial that $V'(x) \leq V(x)$. In \cite{Er98} Erd\H{o}s suggests the limit may be infinite. See also [416] .




References

[Er98] Erd\H{o}s, Paul, Some of my new and almost new problems and results in combinatorial number theory . Number theory (Eger, 1996) (1998), 169-180.

\noindent\rule{\linewidth}{0.4pt}


%% ============================================================
\section{Problem \#420}
\label{prob:420}

\noindent\textbf{Status:} OPEN \quad|\quad \textbf{Prize:} none \quad|\quad \textbf{Updated:} 2025-08-31

\noindent\textbf{Tags:} number theory

\medskip

\noindent\textbf{Statement:}

If $\tau(n)$ counts the number of divisors of $n$ then let\[F(f,n)=\frac{\tau((n+\lfloor f(n)\rfloor)!)}{\tau(n!)}.\]Is it true that\[\lim_{n\to \infty}F((\log n)^C,n)=\infty\]for large $C$? Is it true that $F(\log n,n)$ is everywhere dense in $(1,\infty)$? More generally, if $f(n)\leq \log n$ is a monotonic function such that $f(n)\to \infty$ as $n\to \infty$, then is $F(f,n)$ everywhere dense?




Erd\H{o}s and Graham write that it is easy to show that $\lim F(n^{1/2},n)=\infty$, and in fact the $n^{1/2}$ can be replaced by $n^{1/2-c}$ for some small constant $c&#62;0$. Erd\H{o}s, Graham, Ivi\'{c}, and Pomerance \cite{EGIP96} have proved that\[\liminf F(c\log n, n) = 1\]for any $c&#62;0$, and\[\lim F(n^{4/9},n)=\infty.\](The exponent $4/9$ can be improved slightly.) They also prove that, if $f(n)=o((\log n)^2)$, then for almost all $n$\[F(f,n)\sim 1.\]van Doorn notes in the comments that the existence of infinitely many bounded prime gaps implies\[\limsup_{n\to \infty}F(g(n),n)=\infty\]for any $g(n)\to \infty$, and that Cram\'{e}r's conjecture implies\[\lim F(g(n)(\log n)^2, n)=\infty\]for any $g(n)\to \infty$&#62;




References

[EGIP96] Erd\H{o}s, Paul and Graham, S. W. and Ivi\'c, Aleksandar and Pomerance, Carl, On the number of divisors of {$n!$} . (1996), 337--355.

\noindent\rule{\linewidth}{0.4pt}


%% ============================================================
\section{Problem \#421}
\label{prob:421}

\noindent\textbf{Status:} OPEN \quad|\quad \textbf{Prize:} none \quad|\quad \textbf{Updated:} 2025-08-31

\noindent\textbf{Tags:} number theory

\medskip

\noindent\textbf{Statement:}

Is there a sequence $1\leq d_1&#60;d_2&#60;\cdots$ with density $1$ such that all products $\prod_{u\leq i\leq v}d_i$ are distinct?




A construction of Selfridge (see [786] ) shows that there exists such a sequence of density $&gt;1/e-\epsilon$ for any $\epsilon&gt;0$. See also [786] .

\noindent\rule{\linewidth}{0.4pt}


%% ============================================================
\section{Problem \#422}
\label{prob:422}

\noindent\textbf{Status:} OPEN \quad|\quad \textbf{Prize:} none \quad|\quad \textbf{Updated:} 2025-08-31

\noindent\textbf{Tags:} number theory

\medskip

\noindent\textbf{Statement:}

Let $f(1)=f(2)=1$ and for $n&#62;2$\[f(n) = f(n-f(n-1))+f(n-f(n-2)).\]Does $f(n)$ miss infinitely many integers? What is its behaviour?




Asked by Hofstadter. The sequence begins $1,1,2,3,3,4,\ldots$ and is A005185 in the OEIS. It is not even known whether $f(n)$ is well-defined for all $n$.

\noindent\rule{\linewidth}{0.4pt}


%% ============================================================
\section{Problem \#423}
\label{prob:423}

\noindent\textbf{Status:} OPEN \quad|\quad \textbf{Prize:} none \quad|\quad \textbf{Updated:} 2025-08-31

\noindent\textbf{Tags:} number theory

\medskip

\noindent\textbf{Statement:}

Let $a_1=1$ and $a_2=2$ and for $k\geq 3$ choose $a_k$ to be the least integer $&#62;a_{k-1}$ which is the sum of at least two consecutive terms of the sequence. What is the asymptotic behaviour of this sequence?




Asked by Hofstadter (in \cite{Er77c} Erd\H{o}s says Hofstadter was inspired by a similar question of Ulam). The sequence begins\[1,2,3,5,6,8,10,11,\ldots\]and is A005243 in the OEIS. Bolan and Tang \cite{Ta26} have independently proved that there are infinitely many integers which do not appear in this sequence. In fact, the sequence $a_n-n$ is nondecreasing and unbounded. Tang \cite{Ta26} further proved that $a_n \ll n^{\frac{1}{c-1}+o(1)}$, where $c$ is any exponent such that $\lvert A-A\rvert \geq \lvert A\rvert^{c-o(1)}$ for all convex sets $A$. In particular, the latest bounds of Cushman \cite{Cu25} yield\[a_n \ll n^{\frac{688}{413}+o(1)}\leq n^{1.6659+o(1)},\]and the conjecture of Erd\H{o}s and Hegyv\'{a}ri that one can take $c=2$ in the convex set problem would imply $a_n\leq n^{1+o(1)}$. It seems likely that $a_n=n+o(n)$ is the truth. Tang has further proved an improved lower bound of\[a_n=n+\Omega(\log\log n).\]




References

[Cu25] A. Cushman, A Note on the Sum-Product Problem and the Convex Sumset Problem . arXiv:2512.13849 (2025).

[Er77c] Erd\H{o}s, Paul, Problems and results on combinatorial number theory. III . Number theory day (Proc. Conf., Rockefeller Univ., New York, 1976) (1977), 43-72.

[Ta26] Q. Tang, The Hofstadter consecutive-sum sequence omits infinitely many positive integers . arXiv:2603.09939 (2026).

\noindent\rule{\linewidth}{0.4pt}


%% ============================================================
\section{Problem \#424}
\label{prob:424}

\noindent\textbf{Status:} OPEN \quad|\quad \textbf{Prize:} none \quad|\quad \textbf{Updated:} 2025-08-31

\noindent\textbf{Tags:} number theory

\medskip

\noindent\textbf{Statement:}

Let $a_1=2$ and $a_2=3$ and continue the sequence by appending to $a_1,\ldots,a_n$ all possible values of $a_ia_j-1$ with $i\neq j$. Is it true that the set of integers which eventually appear has positive density?




Asked by Hofstadter. The sequence begins $2,3,5,9,14,17,26,\ldots$ and is A005244 in the OEIS. This problem is also discussed in section E31 of Guy's book Unsolved Problems in Number Theory. In \cite{ErGr80} (and in Guy's book) this problem as written is asking for whether almost all integers appear in this sequence, but the answer to this is trivially no (as pointed out to me by Steinerberger): no integer $\equiv 1\pmod{3}$ is ever in the sequence, so the set of integers which appear has density at most $2/3$. This is easily seen by induction, and the fact that if $a,b\in \{0,2\}\pmod{3}$ then $ab-1\in \{0,2\}\pmod{3}$. Presumably it is the weaker question of whether a positive density of integers appear (as correctly asked in \cite{Er77c}) that was also intended in \cite{ErGr80}. As with many of Erd\H{o}s' questions, by 'positive density' he most likely meant 'positive lower density' - in other words, does there exist $c&gt;0$ such that for all large $x$ the number of values of $a_i$ in $[1,x]$ is at least $cx$? See also Problem 63 of Green's open problems list .




References

[Er77c] Erd\H{o}s, Paul, Problems and results on combinatorial number theory. III . Number theory day (Proc. Conf., Rockefeller Univ., New York, 1976) (1977), 43-72.

[ErGr80] Erd\H{o}s, P. and Graham, R., Old and new problems and results in combinatorial number theory . Monographies de L'Enseignement Mathematique (1980).

\noindent\rule{\linewidth}{0.4pt}


%% ============================================================
\section{Problem \#425}
\label{prob:425}

\noindent\textbf{Status:} OPEN \quad|\quad \textbf{Prize:} none \quad|\quad \textbf{Updated:} 2025-08-31

\noindent\textbf{Tags:} number theory, sidon sets

\medskip

\noindent\textbf{Statement:}

Let $F(n)$ be the maximum possible size of a subset $A\subseteq\{1,\ldots,N\}$ such that the products $ab$ are distinct for all $a&#60;b$. Is there a constant $c$ such that\[F(n)=\pi(n)+(c+o(1))n^{3/4}(\log n)^{-3/2}?\]If $A\subseteq \{1,\ldots,n\}$ is such that all products $a_1\cdots a_r$ are distinct for $a_1&#60;\cdots &#60;a_r$ then is it true that\[\lvert A\rvert \leq \pi(n)+O(n^{\frac{r+1}{2r}})?\]




Erd\H{o}s \cite{Er68} proved that there exist some constants $0&lt;c_1\leq c_2$ such that\[\pi(n)+c_1 n^{3/4}(\log n)^{-3/2}\leq F(n)\leq \pi(n)+c_2 n^{3/4}(\log n)^{-3/2}.\]This problem can also be considered in the real numbers: that is, what is the size of the the largest $A\subset [1,x]$ such that for any distinct $a,b,c,d\in A$ we have $\lvert ab-cd\rvert \geq 1$? Erd\H{o}s had conjectured (see \cite{Er73} and \cite{Er77c}) that $\lvert A\rvert=o(x)$. In \cite{ErGr80} Erd\H{o}s and Graham report that Alexander had given a construction disproving this conjecture, establishing that $\lvert A\rvert\gg x$ is possible. Alexander's construction is given in \cite{Er80}, and we sketch a simplified version. Let $B\subseteq [1,X^2]$ be a Sidon set of integers of size $\gg X$ and\[A=\{ X e^{b/X^2} : b\in B\}.\]It is easy to check that\[\lvert ab-cd\rvert \geq X^2\lvert 1-e^{1/X^2}\rvert \gg 1\]for distinct $a,b,c,d\in A$, and (after rescaling $A$ by some constant factor) this produces a set of size $\gg X$ with the desired property in the interval $[X,O(X)]$. Furthermore, a simple modification allows for $A$ to also be $1$-separated. In \cite{Er77c} Erd\H{o}s considers a similar generalisation for sets of complex numbers or complex integers. See also [490] , [793] , and [796] .




References

[Er68] Erd\H{o}s, P., On some applications of graph theory to number theoretic problems . Publ. Ramanujan Inst. (1968/69), 131-136.

[Er73] Erd\H{o}s, P., Problems and results on combinatorial number theory . A survey of combinatorial theory (Proc. Internat. Sympos., Colorado State Univ., Fort Collins, Colo., 1971) (1973), 117-138.

[Er77c] Erd\H{o}s, Paul, Problems and results on combinatorial number theory. III . Number theory day (Proc. Conf., Rockefeller Univ., New York, 1976) (1977), 43-72.

[Er80] Erd\H{o}s, Paul, A survey of problems in combinatorial number theory . Ann. Discrete Math. (1980), 89-115.

[ErGr80] Erd\H{o}s, P. and Graham, R., Old and new problems and results in combinatorial number theory . Monographies de L'Enseignement Mathematique (1980).

\noindent\rule{\linewidth}{0.4pt}


%% ============================================================
\section{Problem \#428}
\label{prob:428}

\noindent\textbf{Status:} OPEN \quad|\quad \textbf{Prize:} none \quad|\quad \textbf{Updated:} 2025-08-31

\noindent\textbf{Tags:} number theory, primes

\medskip

\noindent\textbf{Statement:}

Is there a set $A\subseteq \mathbb{N}$ such that, for infinitely many $n$, all of $n-a$ are prime for all $a\in A$ with $0&lt;a&lt;n$ and\[\liminf\frac{\lvert A\cap [1,x]\rvert}{\pi(x)}&gt;0?\]




Erd\H{o}s and Graham could show this is true (assuming the prime $k$-tuple conjecture) if we replace $\liminf$ by $\limsup$.

\noindent\rule{\linewidth}{0.4pt}


%% ============================================================
\section{Problem \#430}
\label{prob:430}

\noindent\textbf{Status:} OPEN \quad|\quad \textbf{Prize:} none \quad|\quad \textbf{Updated:} 2025-08-31

\noindent\textbf{Tags:} number theory

\medskip

\noindent\textbf{Statement:}

Fix some integer $n$ and define a decreasing sequence in $[1,n)$ by $a_1=n-1$ and, for $k\geq 2$, letting $a_k$ be the greatest integer in $[1,a_{k-1})$ such that all of the prime factors of $a_k$ are $&#62;n-a_k$. Is it true that, for sufficiently large $n$, not all of this sequence can be prime?




Erd\H{o}s and Graham write 'preliminary calculations made by Selfridge indicate that this is the case but no proof is in sight'. For example if $n=8$ we have $a_1=7$ and $a_2=5$ and then must stop. Sarosh Adenwalla has observed that this problem is equivalent to (the first part of) [385] . Indeed, assuming a positive answer to that, for all large $n$, there exists a composite $m&lt;n$ such that all primes dividing $m$ are $&gt;n-m$. It follows that such an $m$ is equal to some $a_i$ in the sequence defined for $[1,n)$, and $m$ is composite by assumption.

\noindent\rule{\linewidth}{0.4pt}


%% ============================================================
\section{Problem \#431}
\label{prob:431}

\noindent\textbf{Status:} OPEN \quad|\quad \textbf{Prize:} none \quad|\quad \textbf{Updated:} 2025-08-31

\noindent\textbf{Tags:} number theory, primes

\medskip
\noindent\textit{inverse Goldbach problem}


\noindent\textbf{Statement:}

Are there two infinite sets $A$ and $B$ such that $A+B$ agrees with the set of prime numbers up to finitely many exceptions?




A problem of Ostmann, sometimes known as the 'inverse Goldbach problem'. (In \cite{Er80} Erd\H{o}s dates this to abou 25 years old, putting it at about 1955.) The answer is surely no. The best result in this direction is due to Elsholtz and Harper \cite{ElHa15}, who showed that if $A,B$ are such sets then for all large $x$ we must have\[\frac{x^{1/2}}{\log x\log\log x} \ll \lvert A \cap [1,x]\rvert \ll x^{1/2}\log\log x\]and similarly for $B$. Elsholtz \cite{El01} has proved there are no sets $A,B,C$ (all of size at least $2$) such that $A+B+C$ agrees with the set of prime numbers up to finitely many exceptions. Granville \cite{Gr90} proved, conditional on the prime $k$-tuples conjecture, that there are infinite sets $B$ and $C$ such that\[\{ \tfrac{b+c}{2}: b\in B, c\in C\}\]is a subset of the primes. Tao and Ziegler \cite{TaZi23} gave an unconditional proof that there are infinite sets $B=\{b_1&lt;\cdots\}$ and $C=\{c_1&lt;\cdots\}$ such that\[\{ b_i+c_j : b_i\in B, c_j\in C, i&lt;j\}\]is a subset of the primes. See also [429] and [432] .




References

[El01] Elsholtz, Christian, The inverse Goldbach problem . Mathematika (2001), 151-158.

[ElHa15] Elsholtz, Christian and Harper, Adam J., Additive decompositions of sets with restricted prime factors . Trans. Amer. Math. Soc. (2015), 7403-7427.

[Er80] Erd\H{o}s, Paul, A survey of problems in combinatorial number theory . Ann. Discrete Math. (1980), 89-115.

[Gr90] Granville, Andrew, A note on sums of primes . Canad. Math. Bull. (1990), 452--454.

[TaZi23] Tao, Terence and Ziegler, Tamar, Infinite partial sumsets in the primes . J. Anal. Math. (2023), 375--389.

\noindent\rule{\linewidth}{0.4pt}


%% ============================================================
\section{Problem \#432}
\label{prob:432}

\noindent\textbf{Status:} OPEN \quad|\quad \textbf{Prize:} none \quad|\quad \textbf{Updated:} 2025-08-31

\noindent\textbf{Tags:} number theory

\medskip

\noindent\textbf{Statement:}

Let $A,B\subseteq \mathbb{N}$ be two infinite sets. How dense can $A+B$ be if all elements of $A+B$ are pairwise relatively prime?




Asked by Straus, inspired by a problem of Ostmann (see [431] ).

\noindent\rule{\linewidth}{0.4pt}


%% ============================================================
\section{Problem \#436}
\label{prob:436}

\noindent\textbf{Status:} OPEN \quad|\quad \textbf{Prize:} none \quad|\quad \textbf{Updated:} 2025-08-31

\noindent\textbf{Tags:} number theory

\medskip

\noindent\textbf{Statement:}

If $p$ is a prime and $k,m\geq 2$ then let $r(k,m,p)$ be the minimal $r$ such that $r,r+1,\ldots,r+m-1$ are all $k$th power residues modulo $p$. Let\[\Lambda(k,m)=\limsup_{p\to \infty} r(k,m,p).\]Is it true that $\Lambda(k,2)$ is finite for all $k$? Is $\Lambda(k,3)$ finite for all odd $k$? How large are they?




Asked by Lehmer and Lehmer \cite{LeLe62}, who note that for example $\Lambda(2,2)=9$ - indeed, $9$ is always a quadratic residue, and if $10$ isn't then either $2$ or $5$ is, and hence at least one of $1,2$ or $4,5$ or $9,10$ is a consecutive pair of quadratic residues (and similarly there are infinitely many $p$ for which there are no consecutive quadratic residues below $9,10$). A similar argument of Dunton \cite{Du65} proves $\Lambda(3,2)=77$, and Bierstedt and Mills \cite{BiMi63} proved $\Lambda(4,2)=1224$. Lehmer and Lehmer proved that $\Lambda(k,3)=\infty$ for all even $k$ and $\Lambda(k,4)=\infty$ for all $k\leq 1048909$. Lehmer, Lehmer, and Mills \cite{LLM63} proved $\Lambda(5,2)=7888$ and $\Lambda(6,2)=202124$. Brillhart, Lehmer, and Lehmer \cite{BLL64} proved $\Lambda(7,2)=1649375$. Lehmer, Lehmer, Mills, and Selfridge \cite{LLMS62} proved that $\Lambda(3,3)=23532$. Graham \cite{Gr64g} proved that $\Lambda(k,l)=\infty$ for all $k\geq 2$ and $l\geq 4$. Hildebrand \cite{Hi91} resolved the first question, proving that $\Lambda(k,2)$ is finite for all $k$: in other words, for any $k\geq 2$, if $p$ is sufficiently large then there exists a pair of consecutive $k$th power residues modulo $p$ in $[1,O_k(1)]$. The remaining questions are to examine whether $\Lambda(k,3)$ is finite for all odd $k\geq 5$, and the growth rate of $\Lambda(k,2)$ and $\Lambda(k,3)$ as functions of $k$.




References

[BLL64] Brillhart, John and Lehmer, D. H. and Lehmer, Emma, Bounds for pairs of consecutive seventh and higher power residues . Math. Comp. (1964), 397--407.

[BiMi63] Bierstedt, R. G. and Mills, W. H., On the bound for a pair of consecutive quartic residues of a prime . Proc. Amer. Math. Soc. (1963), 628--632.

[Du65] Dunton, M., Bounds for pairs of cubic residues . Proc. Amer. Math. Soc. (1965), 330--332.

[Gr64g] Graham, R. L., On quadruples of consecutive {$k$}th power residues . Proc. Amer. Math. Soc. (1964), 196--197.

[Hi91] Hildebrand, Adolf, On consecutive {$k$}th power residues. II . Michigan Math. J. (1991), 241-253.

[LLM63] Lehmer, D. H. and Lehmer, Emma and Mills, W. H., Pairs of consecutive power residues . Canadian J. Math. (1963), 172--177.

[LLMS62] Lehmer, D. H. and Lehmer, E. and Mills, W. H. and Selfridge, J. L., Machine proof of a theorem on cubic residues . Math. Comp. (1962), 407--415.

[LeLe62] Lehmer, D. H. and Lehmer, Emma, On runs of residues . Proc. Amer. Math. Soc. (1962), 102-106.

\noindent\rule{\linewidth}{0.4pt}


%% ============================================================
\section{Problem \#445}
\label{prob:445}

\noindent\textbf{Status:} OPEN \quad|\quad \textbf{Prize:} none \quad|\quad \textbf{Updated:} 2025-08-31

\noindent\textbf{Tags:} number theory

\medskip

\noindent\textbf{Statement:}

Is it true that, for any $c&#62;1/2$, if $p$ is a sufficiently large prime then, for any $n\geq 0$, there exist $a,b\in(n,n+p^c)$ such that $ab\equiv 1\pmod{p}$?




Heilbronn (unpublished) proved this for $c$ sufficiently close to $1$. Heath-Brown \cite{He00} used Kloosterman sums to prove this for all $c&gt;3/4$. This is discussed in this MathOverflow question .




References

[He00] Heath-Brown, D. R., Arithmetic applications of {K}loosterman sums . Nieuw Arch. Wiskd. (5) (2000), 380--384.

\noindent\rule{\linewidth}{0.4pt}


%% ============================================================
\section{Problem \#450}
\label{prob:450}

\noindent\textbf{Status:} OPEN \quad|\quad \textbf{Prize:} none \quad|\quad \textbf{Updated:} 2025-08-31

\noindent\textbf{Tags:} number theory, divisors

\medskip

\noindent\textbf{Statement:}

How large must $y=y(\epsilon,n)$ be such that the number of integers in $(x,x+y)$ with a divisor in $(n,2n)$ is at most $\epsilon y$?




It is not clear what the intended quantifier on $x$ is. Cambie has observed that if this is intended to hold for all $x$ then, provided\[\epsilon(\log n)^\delta (\log\log n)^{3/2}\to \infty\]as $n\to \infty$, where $\delta=0.086\cdots$, there is no such $y$, which follows from an averaging argument and the work of Ford \cite{Fo08}. On the other hand, Cambie has observed that if $\epsilon\ll 1/n$ then $y(\epsilon,n)\sim 2n$: indeed, if $y&#60;2n$ then this is impossible taking $x+n$ to be a multiple of the lowest common multiple of $\{n+1,\ldots,2n-1\}$. On the other hand, for every fixed $\delta\in (0,1)$ and $n$ large every $2(1+\delta)n$ consecutive elements contains many elements which are a multiple of an element in $(n,2n)$.




References

[Fo08] Ford, Kevin, The distribution of integers with a divisor in a given interval . Ann. of Math. (2) (2008), 367-433.

\noindent\rule{\linewidth}{0.4pt}


%% ============================================================
\section{Problem \#451}
\label{prob:451}

\noindent\textbf{Status:} OPEN \quad|\quad \textbf{Prize:} none \quad|\quad \textbf{Updated:} 2025-08-31

\noindent\textbf{Tags:} number theory

\medskip

\noindent\textbf{Statement:}

Estimate $n_k$, the smallest integer $&#62;2k$ such that $\prod_{1\leq i\leq k}(n_k-i)$ has no prime factor in $(k,2k)$.




Erd\H{o}s and Graham write 'we can prove $n_k&#62;k^{1+c}$ but no doubt much more is true'. In \cite{Er79d} Erd\H{o}s writes that probably $n_k&lt;e^{o(k)}$ but $n_k&gt;k^d$ for all constant $d$. Adenwalla observes that an easy upper bound is $n_k\leq \prod_{k&#60;p&#60;2k}p=e^{O(k)}$.




References

[Er79d] Erd\H{o}s, P., Some unconventional problems in number theory . Acta Math. Acad. Sci. Hungar. (1979), 71-80.

\noindent\rule{\linewidth}{0.4pt}


%% ============================================================
\section{Problem \#452}
\label{prob:452}

\noindent\textbf{Status:} OPEN \quad|\quad \textbf{Prize:} none \quad|\quad \textbf{Updated:} 2025-08-31

\noindent\textbf{Tags:} number theory

\medskip

\noindent\textbf{Statement:}

Let $\omega(n)$ count the number of distinct prime factors of $n$. What is the size of the largest interval $I\subseteq [x,2x]$ such that $\omega(n)&#62;\log\log n$ for all $n\in I$?




Erd\H{o}s \cite{Er37} proved that the density of integers $n$ with $\omega(n)&#62;\log\log n$ is $1/2$. The Chinese remainder theorem implies that there is such an interval with\[\lvert I\rvert \geq (1+o(1))\frac{\log x}{(\log\log x)^2}.\]It could be true that there is such an interval of length $(\log x)^{k}$ for arbitrarily large $k$.




References

[Er37] Erd\H{o}s, Paul, Note on the number of prime divisors of integers . J. London Math. Soc. (1937), 308-314.

\noindent\rule{\linewidth}{0.4pt}


%% ============================================================
\section{Problem \#454}
\label{prob:454}

\noindent\textbf{Status:} OPEN \quad|\quad \textbf{Prize:} none \quad|\quad \textbf{Updated:} 2025-08-31

\noindent\textbf{Tags:} number theory, primes

\medskip

\noindent\textbf{Statement:}

Let\[f(n) = \min_{i&#60;n} (p_{n+i}+p_{n-i}),\]where $p_k$ is the $k$th prime. Is it true that\[\limsup_n (f(n)-2p_n)=\infty?\]




Pomerance \cite{Po79} has proved the $\limsup$ is at least $2$.




References

[Po79] Pomerance, Carl, The prime number graph . Math. Comp. (1979), 399-408.

\noindent\rule{\linewidth}{0.4pt}


%% ============================================================
\section{Problem \#455}
\label{prob:455}

\noindent\textbf{Status:} OPEN \quad|\quad \textbf{Prize:} none \quad|\quad \textbf{Updated:} 2025-08-31

\noindent\textbf{Tags:} number theory

\medskip

\noindent\textbf{Statement:}

Let $q_1&#60;q_2&#60;\cdots$ be a sequence of primes such that\[q_{n+1}-q_n\geq q_n-q_{n-1}.\]Must\[\lim_n \frac{q_n}{n^2}=\infty?\]




Richter \cite{Ri76} proved that\[\liminf_n \frac{q_n}{n^2}&#62;0.352\cdots.\]




References

[Ri76] Richter, Bernd, \"{U}ber die Monotonie von Differenzenfolgen . Acta Arith. (1976), 225-227.

\noindent\rule{\linewidth}{0.4pt}


%% ============================================================
\section{Problem \#456}
\label{prob:456}

\noindent\textbf{Status:} OPEN \quad|\quad \textbf{Prize:} none \quad|\quad \textbf{Updated:} 2025-08-31

\noindent\textbf{Tags:} number theory

\medskip

\noindent\textbf{Statement:}

Let $p_n$ be the smallest prime $\equiv 1\pmod{n}$ and let $m_n$ be the smallest integer such that $n\mid \phi(m_n)$. Is it true that $m_n&#60;p_n$ for almost all $n$? Does $p_n/m_n\to \infty$ for almost all $n$? Are there infinitely many primes $p$ such that $p-1$ is the only $n$ for which $m_n=p$?




Linnik's theorem implies that $p_n\leq n^{O(1)}$. It is trivial that $m_n\leq p_n$ always. If $n=q-1$ for some prime $q$ then $m_n=p_n$. Erd\H{o}s \cite{Er79e} writes it is 'easy to show' that for infinitely many $n$ we have $m_n &lt;p_n$, and that $m_n/n\to \infty$ for almost all $n$. van Doorn in the comments has noted that if $n=2^{2k+1}$ then $m_n\leq 2n$ and $p_n\geq 2n+1$.




References

[Er79e] Erd\H{o}s, Paul, Some unconventional problems in number theory . Ast\'{e}risque (1979), 73-82.

\noindent\rule{\linewidth}{0.4pt}


%% ============================================================
\section{Problem \#460}
\label{prob:460}

\noindent\textbf{Status:} OPEN \quad|\quad \textbf{Prize:} none \quad|\quad \textbf{Updated:} 2025-08-31

\noindent\textbf{Tags:} number theory

\medskip
\noindent\textit{ambiguous statement}


\noindent\textbf{Statement:}

Let $a_0=0$ and $a_1=1$, and in general define $a_k$ to be the least integer $&#62;a_{k-1}$ for which $(n-a_k,n-a_i)=1$ for all $0\leq i&#60;k$. Does\[\sum_{0&#60;a_i&#60; n}\frac{1}{a_i}\to \infty\]as $n\to \infty$? What about if we restrict the sum to those $i$ such that $n-a_j$ is divisible by some prime $\leq a_j$, or the complement of such $i$?




This question arose in work of Eggleton, Erd\H{o}s, and Selfridge, who could prove that $a_k &lt;k^{2+o(1)}$ for $k$ large enough depending on $n$, but conjectured that in fact $a_k\ll k\log k$ is true. The problem above is from \cite{Er77c}. This question is stated slightly differently in \cite{ErGr80}, which has $a_0=n$ instead of $a_0=0$ and $1\leq i&lt;k$ instead of $0\leq i&lt;k$. The material difference this makes is that the main formulation above has $(a_k,n)=1$ for all $k\geq 1$, which is not imposed in the second formulation. Furthermore, in both \cite{Er77c} and \cite{ErGr80} this is stated without the restriction to $a_k&lt;n$ in the sum, although perhaps this was implicitly intended. Without this condition the sum is infinite since the sequence of $a_k$ contains $n+p$ for all primes $p&gt;n$ (this observation was made by Svyable using ChatGPT). Unfortunately, although both sources mention a forthcoming paper of Eggleton, Erd\H{o}s, and Selfridge, I cannot find a candidate paper of theirs with this problem in, and hence the motivation behind this problem, and what the precise problem intended was, is unclear. Chojecki has noted in the comments that a positive solution to the main problem would follow if\[f(n) = \sum_{a&lt;n}1_{P^-(n-a)&gt;a}\frac{1}{a}\to \infty,\]where $P^-(\cdot)$ is the least prime factor. Standard estimates on rough numbers show that $\frac{1}{N}\sum_{n\leq N}f(n)\gg \log\log N$, so $f(n)$ does diverge on average, but it is unclear whether $f(n)\to \infty$ for all $n$.




References

[Er77c] Erd\H{o}s, Paul, Problems and results on combinatorial number theory. III . Number theory day (Proc. Conf., Rockefeller Univ., New York, 1976) (1977), 43-72.

[ErGr80] Erd\H{o}s, P. and Graham, R., Old and new problems and results in combinatorial number theory . Monographies de L'Enseignement Mathematique (1980).

\noindent\rule{\linewidth}{0.4pt}


%% ============================================================
\section{Problem \#461}
\label{prob:461}

\noindent\textbf{Status:} OPEN \quad|\quad \textbf{Prize:} none \quad|\quad \textbf{Updated:} 2025-08-31

\noindent\textbf{Tags:} number theory, primes

\medskip

\noindent\textbf{Statement:}

Let $s_t(n)$ be the $t$-smooth component of $n$ - that is, the product of all primes $p$ (with multiplicity) dividing $n$ such that $p&#60;t$. Let $f(n,t)$ count the number of distinct possible values for $s_t(m)$ for $m\in [n+1,n+t]$. Is it true that\[f(n,t)\gg t\](uniformly, for all $t$ and $n$)?




Erd\H{o}s and Graham report they can show\[f(n,t) \gg \frac{t}{\log t}.\]

\noindent\rule{\linewidth}{0.4pt}


%% ============================================================
\section{Problem \#462}
\label{prob:462}

\noindent\textbf{Status:} OPEN \quad|\quad \textbf{Prize:} none \quad|\quad \textbf{Updated:} 2025-08-31

\noindent\textbf{Tags:} number theory, primes

\medskip

\noindent\textbf{Statement:}

Let $p(n)$ denote the least prime factor of $n$. There is a constant $c&#62;0$ such that\[\sum_{\substack{n&lt;x\\ n\textrm{ not prime}}}\frac{p(n)}{n}\sim c\frac{x^{1/2}}{(\log x)^2}.\]Is it true that there exists a constant $C&gt;0$ such that\[\sum_{x\leq n\leq x+Cx^{1/2}(\log x)^2}\frac{p(n)}{n} \gg 1\]for all large $x$?

\noindent\rule{\linewidth}{0.4pt}


%% ============================================================
\section{Problem \#463}
\label{prob:463}

\noindent\textbf{Status:} OPEN \quad|\quad \textbf{Prize:} none \quad|\quad \textbf{Updated:} 2025-08-31

\noindent\textbf{Tags:} number theory, primes

\medskip

\noindent\textbf{Statement:}

Is there a function $f$ with $f(n)\to \infty$ as $n\to \infty$ such that, for all large $n$, there is a composite number $m$ such that\[n+f(n)&#60;m&#60;n+p(m)?\](Here $p(m)$ is the least prime factor of $m$.)




In \cite{Er92e} Erd\H{o}s asks about\[F(n)=\min_{m&gt;n}(m-p(m)),\]and whether $n-F(n)\sim cn^{1/2}$ for some $c&gt;0$. See also [385] .




References

[Er92e] Erd\H{o}s, P\'{a}l, Some Unsolved problems in Geometry, Number Theory and Combinatorics . Eureka (1992), 44-48.

\noindent\rule{\linewidth}{0.4pt}


%% ============================================================
\section{Problem \#467}
\label{prob:467}

\noindent\textbf{Status:} OPEN \quad|\quad \textbf{Prize:} none \quad|\quad \textbf{Updated:} 2025-08-31

\noindent\textbf{Tags:} number theory

\medskip
\noindent\textit{ambiguous statement}


\noindent\textbf{Statement:}

Prove the following for all large $x$: there is a choice of congruence classes $a_p$ for all primes $p\leq x$ and a decomposition $\{p\leq x\}=A\sqcup B$ into two non-empty sets such that, for all $n&#60;x$, there exist some $p\in A$ and $q\in B$ such that $n\equiv a_p\pmod{p}$ and $n\equiv a_q\pmod{q}$.




This is what I assume the intended problem is, although the presentation in \cite{ErGr80} is missing some crucial quantifiers, so I may have misinterpreted it.




References

[ErGr80] Erd\H{o}s, P. and Graham, R., Old and new problems and results in combinatorial number theory . Monographies de L'Enseignement Mathematique (1980).

\noindent\rule{\linewidth}{0.4pt}


%% ============================================================
\section{Problem \#468}
\label{prob:468}

\noindent\textbf{Status:} OPEN \quad|\quad \textbf{Prize:} none \quad|\quad \textbf{Updated:} 2025-08-31

\noindent\textbf{Tags:} number theory, divisors

\medskip

\noindent\textbf{Statement:}

For any $n$ let $D_n$ be the set of sums of the shape $d_1,d_1+d_2,d_1+d_2+d_3,\ldots$ where $1&#60;d_1&#60;d_2&#60;\cdots$ are the divisors of $n$. What is the size of $D_n\backslash \cup_{m&#60;n}D_m$? If $f(N)$ is the minimal $n$ such that $N\in D_n$ then is it true that $f(N)=o(N)$? Perhaps just for almost all $N$?

\noindent\rule{\linewidth}{0.4pt}


%% ============================================================
\section{Problem \#469}
\label{prob:469}

\noindent\textbf{Status:} OPEN \quad|\quad \textbf{Prize:} none \quad|\quad \textbf{Updated:} 2025-08-31

\noindent\textbf{Tags:} number theory, divisors

\medskip

\noindent\textbf{Statement:}

Let $A$ be the set of all $n$ such that $n=d_1+\cdots+d_k$ with $d_i$ distinct proper divisors of $n$, but this is not true for any $m\mid n$ with $m&#60;n$. Does\[\sum_{n\in A}\frac{1}{n}\]converge?




The integers in $A$ are also known as primitive pseudoperfect numbers and are listed as A006036 in the OEIS. The same question can be asked for those $n$ which do not have distinct sums of sets of divisors, but any proper divisor of $n$ does (which are listed as A119425 in the OEIS). Benkoski and Erd\H{o}s \cite{BeEr74} ask about these two sets, and also about the set of $n$ that have a divisor expressible as a distinct sum of other divisors of $n$, but where no proper divisor of $n$ has this property.




References

[BeEr74] Benkoski, S. J. and Erd\H{o}s, P., On weird and pseudoperfect numbers . Math. Comp. (1974), 617-623.

\noindent\rule{\linewidth}{0.4pt}


%% ============================================================
\section{Problem \#470}
\label{prob:470}

\noindent\textbf{Status:} OPEN \quad|\quad \textbf{Prize:} \$10 \quad|\quad \textbf{Updated:} 2025-08-31

\noindent\textbf{Tags:} number theory, divisors

\medskip

\noindent\textbf{Statement:}

Call $n$ weird if $\sigma(n)\geq 2n$ and $n$ is not pseudoperfect, that is, it is not the sum of any set of its divisors. Are there any odd weird numbers? Are there infinitely many primitive weird numbers, i.e. those such that no proper divisor of $n$ is weird?




Weird numbers were investigated by Benkoski and Erd\H{o}s \cite{BeEr74}, who proved that the set of weird numbers has positive density. The smallest weird number is $70$. Melfi \cite{Me15} has proved that there are infinitely many primitive weird numbers, conditional on the fact that $p_{n+1}-p_n&lt;\frac{1}{10}p_n^{1/2}$ for all large $n$, which in turn would follow from well-known conjectures concerning prime gaps. The sequence of weird numbers is A006037 in the OEIS. Fang \cite{Fa22} has shown there are no odd weird numbers below $10^{21}$, and Liddy and Riedl \cite{LiRi18} have shown that an odd weird number must have at least 6 prime divisors. If there are no odd weird numbers then every weird number has abundancy index $&lt;4$ (see [825] ). This is problem B2 in Guy's collection \cite{Gu04} (the \$10 is reported by Guy, offered by Erd\H{o}s for a solution to the question of whether any odd weird numbers exist).




References

[BeEr74] Benkoski, S. J. and Erd\H{o}s, P., On weird and pseudoperfect numbers . Math. Comp. (1974), 617-623.

[Fa22] Searching on the boundary of abundance for odd weird numbers, W. Fang . arXiv:2207.12906 (2022).

[Gu04] Guy, Richard K., Unsolved problems in number theory . (2004), xviii+437.

[LiRi18] J. Liddy and J. Riedl, An algorithm to determine all odd primitive abundant numbers with $d$ prime divisors . Honors Research Projects. 728 (2018).

[Me15] Melfi, Giuseppe, On the conditional infiniteness of primitive weird numbers . J. Number Theory (2015), 508-514.

\noindent\rule{\linewidth}{0.4pt}


%% ============================================================
\section{Problem \#472}
\label{prob:472}

\noindent\textbf{Status:} OPEN \quad|\quad \textbf{Prize:} none \quad|\quad \textbf{Updated:} 2025-08-31

\noindent\textbf{Tags:} number theory

\medskip

\noindent\textbf{Statement:}

Given some initial finite sequence of primes $q_1&#60;\cdots&#60;q_m$ extend it so that $q_{n+1}$ is the smallest prime of the form $q_n+q_i-1$ for $n\geq m$. Is there an initial starting sequence so that the resulting sequence is infinite?




A problem due to Ulam. For example if we begin with $3,5$ then the sequence continues $3,5,7,11,13,17,\ldots$. It is possible that this sequence is infinite.

\noindent\rule{\linewidth}{0.4pt}


%% ============================================================
\section{Problem \#477}
\label{prob:477}

\noindent\textbf{Status:} OPEN \quad|\quad \textbf{Prize:} none \quad|\quad \textbf{Updated:} 2025-08-31

\noindent\textbf{Tags:} number theory

\medskip

\noindent\textbf{Statement:}

Is there a polynomial $f:\mathbb{Z}\to \mathbb{Z}$ of degree at least $2$ and a set $A\subset \mathbb{Z}$ such that for any $n\in \mathbb{Z}$ there is exactly one $a\in A$ and $b\in \{ f(k) : k\in\mathbb{Z}\}$ such that $n=a+b$?




A question of Erd\H{o}s and Graham, who thought the answer was negative. Clearly any such $A$ must be infinite. A simple proof that this is impossible for $f$ of degree $2$ is given in the comments (with the combined efforts of AlphaProof and Adenwalla): let $f(x)=c_2x^2+c_1x+c_0$ with $c_1\neq 0$. For any infinite $A$ there exist $a,b\in A$ such that $a-b=2c_1k$ for some $k\geq 1$, whence\[0\neq a-b=f(k)-f(-k)\]and so $n=a+f(-k)=b+f(k)$ has two distinct solutions. If $c_1=0$ then we can argue similarly finding $a,b\in A$ with $a-b=4c_2k$ for some $k\geq 1$, whence\[0\neq a-b=f(k+1)-f(k-1).\]The same argument works for any $f$ such that $f(\mathbb{Z})-f(\mathbb{Z})$ contains all multiples of some fixed $q\geq 1$, such as any polynomial of the shape\[f(x)=g((x-k)^2)+c_1x+c_0\]for some $g\in \mathbb{Z}[x]$ and $k,c_1,c_0\in \mathbb{Z}$ with $c_1\neq 0$.

\noindent\rule{\linewidth}{0.4pt}


%% ============================================================
\section{Problem \#478}
\label{prob:478}

\noindent\textbf{Status:} OPEN \quad|\quad \textbf{Prize:} none \quad|\quad \textbf{Updated:} 2025-08-31

\noindent\textbf{Tags:} number theory, factorials

\medskip

\noindent\textbf{Statement:}

Let $p$ be a prime and\[A_p = \{ k! \pmod{p} : 1\leq k&#60;p\}.\]Is it true that\[\lvert A_p\rvert \sim (1-\tfrac{1}{e})p?\]




Since $A_p/A_p=\{1,\ldots,p-1\}$ it follows that $\lvert A_p\rvert \gg p^{1/2}$. The best known lower bound is due to Grebennikov, Sagdeev, Semchankau, and Vasilevskii \cite{GSSV24},\[\lvert A_p\rvert \geq (\sqrt{2}-o(1))p^{1/2},\]which follows from proving that $\lvert A_pA_p\rvert=(1+o(1))p$. Wilson's theorem implies $(p-2)!\equiv 1\pmod{p}$, and hence $\lvert A_p\rvert\leq p-2$. Results on $\lvert A_p\rvert$ on average were obtained by Klurman and Munsch \cite{KlMu17}. It is even open whether there are infinitely many primes $p$ such that $\lvert A_p\rvert=p-2$: in other words, such that $2!,\ldots,(p-1)!$ are all distinct modulo $p$ (such primes are sometimes called socialist primes). Rokowska and Schinzel \cite{RoSc60} have proved that any socialist prime is $5\pmod{8}$, and must satisfy $-(\frac{5}{p})=(\frac{-23}{p})=1$. Moreover, the congruence class missing from $2!,\ldots,(p-1)!$ must be $-(\frac{p-1}{2})!$. The only known socialist prime is $p=5$. Trudgian \cite{Tr13} has shown there are no other socialist primes $&#60;10^9$. This was extended to $10^{11}$ by Andreji\'{c} and Tatarevic \cite{AnTa16}. This is also discussed in problems A2 and F11 of Guy's collection \cite{Gu04}.




References

[AnTa16] V. Andreji\'{c} and M. Tatarevic, On distinct residues of factorials . arXiv:1603.04086 (2016).

[GSSV24] Grebennikov, Alexandr and Sagdeev, Arsenii and Semchankau, Aliaksei and Vasilevskii, Aliaksei, On the sequence {$n! \bmod p$} . Rev. Mat. Iberoam. (2024), 637--648.

[Gu04] Guy, Richard K., Unsolved problems in number theory . (2004), xviii+437.

[KlMu17] Klurman, Oleksiy and Munsch, Marc, Distribution of factorials modulo {$p$} . J. Th\'{e}or. Nombres Bordeaux (2017), 169--177.

[RoSc60] Rokowska, B. and Schinzel, A., Sur un probl\`eme de {M}. {E}rd\H os . Elem. Math. (1960), 84--85.

[Tr13] T. Trudgian, There are no socialist primes less than $10^9$ . arXiv:1310.6403 (2013).

\noindent\rule{\linewidth}{0.4pt}


%% ============================================================
\section{Problem \#479}
\label{prob:479}

\noindent\textbf{Status:} OPEN \quad|\quad \textbf{Prize:} none \quad|\quad \textbf{Updated:} 2025-08-31

\noindent\textbf{Tags:} number theory

\medskip

\noindent\textbf{Statement:}

Is it true that, for all $k\neq 1$, there are infinitely many $n$ such that $2^n\equiv k\pmod{n}$?




A conjecture of Graham. It is easy to see that $2^n\not\equiv 1\mod{n}$ for all $n&gt;1$, so the restriction $k\neq 1$ is necessary. Erd\H{o}s and Graham report that Graham, Lehmer, and Lehmer have proved this for $k=2^i$ for $i\geq 1$, or if $k=-1$, but I cannot find such a paper. Tang has written a short note giving a proof for this case. As an indication of the difficulty, when $k=3$ the smallest $n$ such that $2^n\equiv 3\pmod{n}$ is $n=4700063497$. The minimal such $n$ for each $k$ is A036236 in the OEIS.

\noindent\rule{\linewidth}{0.4pt}


%% ============================================================
\section{Problem \#483}
\label{prob:483}

\noindent\textbf{Status:} OPEN \quad|\quad \textbf{Prize:} none \quad|\quad \textbf{Updated:} 2025-08-31

\noindent\textbf{Tags:} number theory, additive combinatorics, ramsey theory

\medskip

\noindent\textbf{Statement:}

Let $f(k)$ be the minimal $N$ such that if $\{1,\ldots,N\}$ is $k$-coloured then there is a monochromatic solution to $a+b=c$. Estimate $f(k)$. In particular, is it true that $f(k) &lt; c^k$ for some constant $c&gt;0$?




The values of $f(k)$ are known as Schur numbers. The best-known bounds for large $k$ are\[(380)^{k/5}-O(1)\leq f(k) \leq (e-\tfrac{1}{6}) k!.\]Note that $380^{1/5}\approx 3.2806$. The lower bound is due to Ageron, Casteras, Pellerin, Portella, Rimmel, and Tomasik \cite{ACPPRT21} (improving previous bounds of Exoo \cite{Ex94} and Fredricksen and Sweet \cite{FrSw00}). The upper bound is due to Xu, Xie, and Chen \cite{XXC02} (improving previous bounds of Wan \cite{Wa97} and Whitehead \cite{Wh73}). See Eliahou \cite{El19} for more on the upper bound. The known values of $f$ are $f(1)=2$, $f(2)=5$, $f(3)=14$, $f(4)=45$, and $f(5)=161$ (see A030126 ). (The equality $f(5)=161$ was established by Heule \cite{He17}). See also [183] (in particular a folklore observation gives $f(k)\leq R(3;k)-1$).




References

[ACPPRT21] R. Ageron, P. Casteras, T. Pellerin, Y. Portella, A. Rimmel, and J. Tomasik, New lower bounds for Schur and weak Schur numbers . arXiv:2112.03175 (2021).

[El19] No reference found.


[Ex94] Exoo, G., A lower bound for Schur numbers and multicolor Ramsey numbers . Electronic J. of Combinatorics (1994).

[FrSw00] Fredricksen, Harold and Sweet, Melvin M., Symmetric sum-free partitions and lower bounds for Schur numbers . Electron. J. Combin. (2000), Research Paper 32, 9.

[He17] M. Heuele, Schur Number Five . arXiv:1711.08076 (2017).

[Wa97] Wan, Honghui, Upper bounds for {R}amsey numbers {$R(3,3,\cdots,3)$} and Schur numbers . J. Graph Theory (1997), 119--122.

[Wh73] Whitehead, Jr., Earl Glen, The {R}amsey number {$N(3,\,3,\,3,\,3;\,2)$} . Discrete Math. (1973), 389--396.

[XXC02] Xu, Xiao Dong and Xie, Zheng and Chen, Zhi, Upper bounds for {R}amsey numbers {$R_n(3)$} and Schur numbers . Math. Econ. (2002), 81--84.

\noindent\rule{\linewidth}{0.4pt}


%% ============================================================
\section{Problem \#486}
\label{prob:486}

\noindent\textbf{Status:} OPEN \quad|\quad \textbf{Prize:} none \quad|\quad \textbf{Updated:} 2025-08-31

\noindent\textbf{Tags:} number theory, primitive sets

\medskip

\noindent\textbf{Statement:}

Let $A\subseteq \mathbb{N}$, and for each $n\in A$ choose some $X_n\subseteq \mathbb{Z}/n\mathbb{Z}$. Let\[B = \{ m\in \mathbb{N} : m\not\in X_n\pmod{n}\textrm{ for all }n\in A\textrm{ with }m&#62;n\}.\]Must $B$ have a logarithmic density, i.e. is it true that\[\lim_{x\to \infty} \frac{1}{\log x}\sum_{\substack{m\in B\\ m&#60;x}}\frac{1}{m}\]exists?




Davenport and Erd\H{o}s \cite{DaEr36} proved that the answer is yes when $X_n=\{0\}$ for all $n\in A$. An alternative elementary proof was later given by Davenport and Erd\H{o}s in \cite{DaEr51}. Erd\H{o}s \cite{Er80} wrote 'perhaps this question is not very difficult as far as I know it has not been attacked really seriously'. The problem considers logarithmic density since Besicovitch \cite{Be34} showed examples exist without a natural density, even when $X_n=\{0\}$ for all $n\in A$. This is a generalisation of [25] (which is the case when $\lvert X_n\rvert=1$ for all $n\in A$).




References

[Be34] Besicovitch, A., On the density of certain sequences of integers . Math. Annalen (1934), 336-341.

[DaEr36] Davenport, H. and Erd\H{o}s, P., On sequences of positive integers . Acta Arithmetica (1936), 147-151.

[DaEr51] Davenport, H. and Erd\H{o}s, P., On sequences of positive integers . J. Indian Math. Soc. (N.S.) (1951), 19--24.

[Er80] Erd\H{o}s, Paul, A survey of problems in combinatorial number theory . Ann. Discrete Math. (1980), 89-115.

\noindent\rule{\linewidth}{0.4pt}


%% ============================================================
\section{Problem \#489}
\label{prob:489}

\noindent\textbf{Status:} OPEN \quad|\quad \textbf{Prize:} none \quad|\quad \textbf{Updated:} 2025-08-31

\noindent\textbf{Tags:} number theory

\medskip

\noindent\textbf{Statement:}

Let $A\subseteq \mathbb{N}$ be a set such that $\lvert A\cap [1,x]\rvert=o(x^{1/2})$. Let\[B=\{ n\geq 1 : a\nmid n\textrm{ for all }a\in A\}.\]If $B=\{b_1&#60;b_2&#60;\cdots\}$ then is it true that\[\lim \frac{1}{x}\sum_{b_i&#60;x}(b_{i+1}-b_i)^2\]exists (and is finite)?




For example, when $A=\{p^2: p\textrm{ prime}\}$ then $B$ is the set of squarefree numbers, and the existence of this limit was proved by Erd\H{o}s. See also [208] .

\noindent\rule{\linewidth}{0.4pt}


%% ============================================================
\section{Problem \#495}
\label{prob:495}

\noindent\textbf{Status:} OPEN \quad|\quad \textbf{Prize:} none \quad|\quad \textbf{Updated:} 2025-08-31

\noindent\textbf{Tags:} diophantine approximation, number theory

\medskip
\noindent\textit{Littlewood conjecture}


\noindent\textbf{Statement:}

Let $\alpha,\beta \in \mathbb{R}$. Is it true that\[\liminf_{n\to \infty} n \| n\alpha \| \| n\beta\| =0\]where $\|x\|$ is the distance from $x$ to the nearest integer?




The infamous Littlewood conjecture .

\noindent\rule{\linewidth}{0.4pt}


%% ============================================================
\section{Problem \#500}
\label{prob:500}

\noindent\textbf{Status:} OPEN \quad|\quad \textbf{Prize:} \$500 \quad|\quad \textbf{Updated:} 2025-08-31

\noindent\textbf{Tags:} graph theory, hypergraphs, turan number

\medskip

\noindent\textbf{Statement:}

What is $\mathrm{ex}_3(n,K_4^3)$? That is, the largest number of $3$-edges which can placed on $n$ vertices so that there exists no $K_4^3$, a set of 4 vertices which is covered by all 4 possible $3$-edges.




A problem of Tur\'{a}n. Tur\'{a}n observed that dividing the vertices into three equal parts $X_1,X_2,X_3$, and taking the edges to be those triples that either have exactly one vertex in each part or two vertices in $X_i$ and one vertex in $X_{i+1}$ (where $X_4=X_1$) shows that\[\mathrm{ex}_3(n,K_4^3)\geq\left(\frac{5}{9}+o(1)\right)\binom{n}{3}.\]This is probably the truth. The current best upper bound is\[\mathrm{ex}_3(n,K_4^3)\leq 0.5611666\binom{n}{3},\]due to Razborov \cite{Ra10}. See also [712] for the general case.




References

[Ra10] Razborov, Alexander A., On 3-hypergraphs with forbidden 4-vertex configurations . SIAM J. Discrete Math. (2010), 946-963.

\noindent\rule{\linewidth}{0.4pt}


%% ============================================================
\section{Problem \#501}
\label{prob:501}

\noindent\textbf{Status:} OPEN \quad|\quad \textbf{Prize:} none \quad|\quad \textbf{Updated:} 2025-08-31

\noindent\textbf{Tags:} combinatorics, set theory

\medskip

\noindent\textbf{Statement:}

For every $x\in\mathbb{R}$ let $A_x\subset \mathbb{R}$ be a bounded set with outer measure $&#60;1$. Must there exist an infinite independent set, that is, some infinite $X\subseteq \mathbb{R}$ such that $x\not\in A_y$ for all $x\neq y\in X$? If the sets $A_x$ are closed and have measure $&#60;1$, then must there exist an independent set of size $3$?




Erd\H{o}s and Hajnal \cite{ErHa60} proved the existence of arbitrarily large finite independent sets (under the assumptions in the first problem). Gladysz \cite{Gl62} proved the existence of an independent set of size $2$ under the assumptions of the second question. Hechler \cite{He72} has shown the answer to the first question is no, assuming the continuum hypothesis. Newelski, Pawlikowski, and Seredy\'{n}ski \cite{NPS87} proved that if all the $A_x$ are closed with measure $&#60;1$ then there is an infinite independent set, giving a strong affirmative answer to the second question.




References

[ErHa60] Erd\H{o}s, P. and Hajnal, A., Some remarks on set theory. VIII . Michigan Math. J. (1960), 187-191.

[Gl62] G\l adysz, S., Bemerkungen \"uber die {U}nabh\"{a}ngigkeit der {P}unkte in {B}ezug auf mengenwertige {F}unktionen . Acta Math. Acad. Sci. Hungar. (1962), 199--201.

[He72] Hechler, S. H., On two problems in combinatorial set theory . Bull. Acad. Polon. Sci. S\'{e}r. Sci. Math. Astronom. Phys. (1972), 429-431.

[NPS87] Newelski, Ludomir and Pawlikowski, Janusz and Seredy\'nski, Witold, Infinite free set for small measure set mappings . Proc. Amer. Math. Soc. (1987), 335--339.

\noindent\rule{\linewidth}{0.4pt}


%% ============================================================
\section{Problem \#503}
\label{prob:503}

\noindent\textbf{Status:} OPEN \quad|\quad \textbf{Prize:} none \quad|\quad \textbf{Updated:} 2025-08-31

\noindent\textbf{Tags:} geometry, distances

\medskip

\noindent\textbf{Statement:}

What is the size of the largest $A\subseteq \mathbb{R}^d$ such that every three points from $A$ determine an isosceles triangle? That is, for any three points $x,y,z$ from $A$, at least two of the distances $\lvert x-y\rvert,\lvert y-z\rvert,\lvert x-z\rvert$ are equal.




When $d=2$ the answer is $6$ (due to Kelly \cite{ErKe47} - an alternative proof is given by Kov\'{a}cs \cite{Ko24c}). When $d=3$ the answer is $8$ (due to Croft \cite{Cr62}). The best upper bound known in general is due to Blokhuis \cite{Bl84} who showed that\[\lvert A\rvert \leq \binom{d+2}{2}.\]Alweiss has observed a lower bound of $\binom{d+1}{2}$ follows from considering the subset of $\mathbb{R}^{d+1}$ formed of all vectors $e_i+e_j$ where $e_i,e_j$ are distinct coordinate vectors. This set can be viewed as a subset of some $\mathbb{R}^d$, and is easily checked to have the required property. Weisenberg observed in the comments that an additional point can be added to Alweiss' construction, giving a lower bound of $\binom{d+1}{2}+1$. The fact that the truth for $d=3$ is $8$ suggests that neither of these bounds is the truth. See also [1088] for a generalisation.




References

[Bl84] Blokhuis, A., Few-distance sets . (1984), iv+70.

[Cr62] Croft, H. T., $9$-point and $7$-point configurations in $3$-space . Proc. London Math. Soc. (3) (1962), 400-424.

[ErKe47] Erd\H{o}s, Paul and Kelly, L. M., Elementary Problems and Solutions: Solutions: E735 . Amer. Math. Monthly (1947), 227-229.

[Ko24c] Z. Kov\'{a}cs, A note on Erd\H{o}s's mysterious remark . arXiv:2412.05190 (2024).

\noindent\rule{\linewidth}{0.4pt}


%% ============================================================
\section{Problem \#507}
\label{prob:507}

\noindent\textbf{Status:} OPEN \quad|\quad \textbf{Prize:} none \quad|\quad \textbf{Updated:} 2025-08-31

\noindent\textbf{Tags:} geometry

\medskip
\noindent\textit{Heilbronn's triangle problem}


\noindent\textbf{Statement:}

Let $\alpha(n)$ be such that every set of $n$ points in the unit disk contains three points which determine a triangle of area at most $\alpha(n)$. Estimate $\alpha(n)$.




Heilbronn's triangle problem . It is trivial that $\alpha(n) \ll 1/n$. Erd\H{o}s observed that $\alpha(n)\gg 1/n^2$. The current best bounds are\[\frac{\log n}{n^2}\ll \alpha(n) \ll \frac{1}{n^{7/6+o(1)}}.\]The lower bound is due to Koml\'{o}s, Pintz, and Szemer\'{e}di \cite{KPS82}. The upper bound is due to Cohen, Pohoata, and Zakharov \cite{CPZ24} (improving on their earlier work \cite{CPZ23} which itself improves an exponent of $8/7$ due to Koml\'{o}s, Pintz, and Szemer\'{e}di \cite{KPS81}). This problem is Problem 77 on Green's open problems list .




References

[CPZ23] Cohen, A. and Pohata, C. and Zakharov, D., A new upper bound for the Heilbronn triangle problem . arXiv:2305.18253 (2023).

[CPZ24] Cohen, A. and Pohata, C. and Zakharov, D., Lower bounds for incidences . arXiv:2409.07658 (2024).

[KPS81] Koml\'{o}s, J\'{a}nos and Pintz, J\'{a}nos and Szemer\'{e}di, Endre, On Heilbronn's triangle problem . J. London Math. Soc. (2) (1981), 385-396.

[KPS82] Koml\'{o}s, J\'{a}nos and Pintz, J\'{a}nos and Szemer\'{e}di, Endre, A lower bound for Heilbronn's problem . J. London Math. Soc. (2) (1982), 13-24.

\noindent\rule{\linewidth}{0.4pt}


%% ============================================================
\section{Problem \#508}
\label{prob:508}

\noindent\textbf{Status:} OPEN \quad|\quad \textbf{Prize:} none \quad|\quad \textbf{Updated:} 2025-08-31

\noindent\textbf{Tags:} geometry, ramsey theory

\medskip
\noindent\textit{Hadwiger-Nelson problem}


\noindent\textbf{Statement:}

What is the chromatic number of the plane? That is, what is the smallest number of colours required to colour $\mathbb{R}^2$ such that no two points of the same colour are distance $1$ apart?




The Hadwiger-Nelson problem . Let $\chi$ be the chromatic number of the plane. An equilateral triangle trivially shows that $\chi\geq 3$. There are several small graphs that show $\chi\geq 4$ (in particular the Moser spindle and Golomb graph). The best bounds currently known are\[5 \leq \chi \leq 7.\]The lower bound is due to de Grey \cite{dG18}. The upper bound can be seen by colouring the plane by tesselating by hexagons with diameter slightly less than $1$. Matolcsi, Ruzsa, Varga, and Zs\'{a}mboki have proved that the fractional chromatic number of the plane is at least $4$. Croft \cite{Cr67} has proved it is at most $4.359\cdots$. See also [704] , [705] , and [706] . The independence number of a finite unit distance graph is the topic of [1070] .




References

[Cr67] H. T. Croft, Incidence incidents . Eureka (1967), 22-26.

[dG18] de Grey, Aubrey D. N. J., The chromatic number of the plane is at least 5 . Geombinatorics (2018), 18-31.

\noindent\rule{\linewidth}{0.4pt}


%% ============================================================
\section{Problem \#509}
\label{prob:509}

\noindent\textbf{Status:} OPEN \quad|\quad \textbf{Prize:} none \quad|\quad \textbf{Updated:} 2025-08-31

\noindent\textbf{Tags:} analysis

\medskip

\noindent\textbf{Statement:}

Let $f(z)\in\mathbb{C}[z]$ be a monic non-constant polynomial. Can the set\[\{ z\in \mathbb{C} : \lvert f(z)\rvert \leq 1\}\]be covered by a set of circles the sum of whose radii is $\leq 2$?




Cartan proved this is true with $2$ replaced by $2e$, which was improved to $2.59$ by Pommerenke \cite{Po61}. Pommerenke \cite{Po59} proved that $2$ is achievable if the set is connected (see [1046] ). The generalisation of this to higher dimensions was asked by Erd\H{o}s as Problem 4.23 in \cite{Ha74}.




References

[Ha74] Hayman, W. K., Research problems in function theory: new problems . (1974), 155--180.

[Po59] Pommerenke, Ch., On some problems by Erd\H{o}s, Herzog and Piranian . Michigan Math. J. (1959), 221-225.

[Po61] Pommerenke, Ch., On metric properties of complex polynomials . Michigan Math. J. (1961), 97-115.

\noindent\rule{\linewidth}{0.4pt}


%% ============================================================
\section{Problem \#510}
\label{prob:510}

\noindent\textbf{Status:} OPEN \quad|\quad \textbf{Prize:} none \quad|\quad \textbf{Updated:} 2025-08-31

\noindent\textbf{Tags:} analysis

\medskip
\noindent\textit{Chowla's cosine problem}


\noindent\textbf{Statement:}

If $A\subset \mathbb{Z}$ is a finite set of size $N$ then is there some absolute constant $c&#62;0$ and $\theta$ such that\[\sum_{n\in A}\cos(n\theta) &#60; -cN^{1/2}?\]




Chowla's cosine problem. Ruzsa \cite{Ru04} (improving on an earlier result of Bourgain \cite{Bo86}), proved an upper bound of\[-\exp(O(\sqrt{\log N})).\]Polynomial bounds were proved independently by Bedert \cite{Be25c} and Jin, Milojevi\'{c}, Tomon, and Zhang \cite{JMTZ25}. The best bound follows from the method of Bedert \cite{Be25c}, which proved the existence of some $c&gt;0$ such that, for all $A$ of size $N$,\[\sum_{n\in A}\cos(n\theta) &lt; -cN^{1/7}.\]The example $A=B-B$, where $B$ is a Sidon set, shows that $N^{1/2}$ would be the best possible here. This problem is Problem 81 on Green's open problems list . This is related to [256] .




References

[Be25c] B. Bedert, Polynomial bounds for the Chowla Cosine Problem . arXiv:2509.05260 (2025).

[Bo86] Bourgain, J., Sur le minimum d'une somme de cosinus . Acta Arith. (1986), 381-389.

[JMTZ25] Z. Jin, A. Milojevi\'{c}, I. Tomon, and S. Zhang, From small eigenvalues to large cuts, and Chowla's cosine problem . arXiv:2509.03490 (2025).

[Ru04] Ruzsa, Imre Z., Negative values of cosine sums . Acta Arith. (2004), 179-186.

\noindent\rule{\linewidth}{0.4pt}


%% ============================================================
\section{Problem \#513}
\label{prob:513}

\noindent\textbf{Status:} OPEN \quad|\quad \textbf{Prize:} none \quad|\quad \textbf{Updated:} 2025-08-31

\noindent\textbf{Tags:} analysis

\medskip

\noindent\textbf{Statement:}

Let $f=\sum_{n=0}^\infty a_nz^n$ be a transcendental entire function. What is the greatest possible value of\[\liminf_{r\to \infty} \frac{\max_n\lvert a_nr^n\rvert}{\max_{\lvert z\rvert=r}\lvert f(z)\rvert}?\]




Let $B$ be the supremum of all such values. It is trivial that $B\in [1/2,1]$. K\"{o}v\'{a}ri (unpublished) observed that it must be $&gt;1/2$. Gray and Shah \cite{GrSh63} give an argument of Clunie which shows $B\leq 2/\pi$. Clunie and Hayman \cite{ClHa64} improved both bounds to\[4/7&lt; B \leq 2/\pi-c\]for some absolute constant $c&gt;0$. He and Tang \cite{HeTa26} have improved the lower bound to\[B&gt; 0.5850724\](note that $4/7\approx 0.57143$). This was further improved slightly to $0.5850788$ by GPT (as prompted by Sothanaphan). Note that $2/\pi \approx 0.63662$. See also [227] .




References

[ClHa64] Clunie, J. and Hayman, W. K., The maximum term of a power series . J. Analyse Math. (1964), 143-186.

[GrSh63] Gray, Alfred and Shah, S. M., A note on entire functions and a conjecture of Erd\H{o}s . Bull. Amer. Math. Soc. (1963), 573-577.

[HeTa26] Y. He and Q. Tang, Generalizing the Clunie--Hayman construction in an Erd\H{o}s maximum-term problem . arXiv:2602.12217 (2026).

\noindent\rule{\linewidth}{0.4pt}


%% ============================================================
\section{Problem \#514}
\label{prob:514}

\noindent\textbf{Status:} OPEN \quad|\quad \textbf{Prize:} none \quad|\quad \textbf{Updated:} 2025-08-31

\noindent\textbf{Tags:} analysis

\medskip

\noindent\textbf{Statement:}

Let $f(z)$ be an entire transcendental function. Does there exist a path $L$ so that, for every $n$,\[\lvert f(z)/z^n\rvert \to \infty\]as $z\to \infty$ along $L$? Can the length of this path be estimated in terms of $M(r)=\max_{\lvert z\rvert=r}\lvert f(z)\rvert$? Does there exist a path along which $\lvert f(z)\rvert$ tends to $\infty$ faster than a fixed function of $M(r)$ (such that $M(r)^\epsilon$)?




Boas (unpublished) has proved the first part, that such a path must exist.

\noindent\rule{\linewidth}{0.4pt}


%% ============================================================
\section{Problem \#517}
\label{prob:517}

\noindent\textbf{Status:} OPEN \quad|\quad \textbf{Prize:} none \quad|\quad \textbf{Updated:} 2025-08-31

\noindent\textbf{Tags:} analysis

\medskip

\noindent\textbf{Statement:}

Let $f(z)=\sum_{k=1}^\infty a_kz^{n_k}$ be an entire function (with $a_k\neq 0$ for all $k\geq 1$). Is it true that if $n_k/k\to \infty$ then $f(z)$ assumes every value infinitely often?




A conjecture of Fej\'{e}r and P\'{o}lya. Fej\'{e}r \cite{Fe08} proved that if $\sum\frac{1}{n_k}&#60;\infty$ then $f(z)$ assumes every value at least once, and Biernacki \cite{Bi28} proved that if $\sum\frac{1}{n_k}&#60;\infty$ then $f(z)$ assumes every value infinitely often. P\'{o}lya \cite{Po29} proved that if $f$ has finite order then $f(z)$ assumes every value infinitely often under the assumption that $\limsup (n_{k+1}-n_k)=\infty$.




References

[Bi28] Biernacki, Mi\'{e}cislas, Sur les \'{e}quations alg\'{e}briques contenant des param\'{e}tres arbitraires . (1928), 145.

[Fe08] Fej\'{e}r, Leopold, \"{U}ber die Wurzel vom kleinsten absoluten Betrage einer algebraischen Gleichung . Math. Ann. (1908), 413-423.

[Po29] P\'olya, G., Untersuchungen \"uber {L}\"ucken und Singularit\"{a}ten von {P}otenzreihen . Math. Z. (1929), 549--640.

\noindent\rule{\linewidth}{0.4pt}


%% ============================================================
\section{Problem \#520}
\label{prob:520}

\noindent\textbf{Status:} OPEN \quad|\quad \textbf{Prize:} none \quad|\quad \textbf{Updated:} 2025-08-31

\noindent\textbf{Tags:} number theory, probability

\medskip

\noindent\textbf{Statement:}

Let $f$ be a Rademacher multiplicative function: a random $\{-1,0,1\}$-valued multiplicative function, where for each prime $p$ we independently choose $f(p)\in \{-1,1\}$ uniformly at random, and for square-free integers $n$ we extend $f(p_1\cdots p_r)=f(p_1)\cdots f(p_r)$ (and $f(n)=0$ if $n$ is not squarefree). Does there exist some constant $c&#62;0$ such that, almost surely,\[\limsup_{N\to \infty}\frac{\sum_{m\leq N}f(m)}{\sqrt{N\log\log N}}=c?\]




Note that if we drop the multiplicative assumption, and simply assign $f(m)=\pm 1$ at random, then this statement is true (with $c=\sqrt{2}$), the law of the iterated logarithm . Wintner \cite{Wi44} proved that, almost surely,\[\sum_{m\leq N}f(m)\ll N^{1/2+o(1)},\]and Erd\H{o}s improved the right-hand side to $N^{1/2}(\log N)^{O(1)}$. Lau, Tenenbaum, and Wu \cite{LTW13} have shown that, almost surely,\[\sum_{m\leq N}f(m)\ll N^{1/2}(\log\log N)^{2+o(1)}.\]Caich \cite{Ca24b} has improved this to\[\sum_{m\leq N}f(m)\ll N^{1/2}(\log\log N)^{3/4+o(1)}.\]Harper \cite{Ha13} has shown that the sum is almost surely not $O(N^{1/2}/(\log\log N)^{5/2+o(1)})$, and conjectured that in fact Erd\H{o}s' conjecture is false, and almost surely\[\sum_{m\leq N}f(m) \ll N^{1/2}(\log\log N)^{1/4+o(1)}.\]




References

[Ca24b] R. Caich, Almost sure upper bound for random multiplicative functions . arXiv:2304.00943 (2024).

[Ha13] Harper, Adam J., Bounds on the suprema of Gaussian processes, and omega results for the sum of a random multiplicative function . Ann. Appl. Probab. (2013), 584-616.

[LTW13] Lau, Yuk-Kam and Tenenbaum, G\'{e}rald and Wu, Jie, On mean values of random multiplicative functions . Proc. Amer. Math. Soc. (2013), 409-420.

[Wi44] Wintner, Aurel, Random factorizations and Riemann's hypothesis . Duke Math. J. (1944), 267-275.

\noindent\rule{\linewidth}{0.4pt}


%% ============================================================
\section{Problem \#521}
\label{prob:521}

\noindent\textbf{Status:} OPEN \quad|\quad \textbf{Prize:} none \quad|\quad \textbf{Updated:} 2025-08-31

\noindent\textbf{Tags:} analysis, polynomials, probability

\medskip

\noindent\textbf{Statement:}

Let $(\epsilon_k)_{k\geq 0}$ be independently uniformly chosen at random from $\{-1,1\}$. If $R_n$ counts the number of real roots of $f_n(z)=\sum_{0\leq k\leq n}\epsilon_k z^k$ then is it true that, almost surely,\[\lim_{n\to \infty}\frac{R_n}{\log n}=\frac{2}{\pi}?\]




Erd\H{o}s and Offord \cite{EO56} showed that the number of real roots of a random degree $n$ polynomial with $\pm 1$ coefficients is $(\frac{2}{\pi}+o(1))\log n$. It is ambiguous in \cite{Er61} whether Erd\H{o}s intended the coefficients to be uniformly chosen from $\{-1,1\}$ or $\{0,1\}$. In the latter case, the constant $\frac{2}{\pi}$ should be $\frac{1}{\pi}$ (see the discussion in the comments). In the case of $\{-1,1\}$ Do \cite{Do24} proved that, if $R_n[-1,1]$ counts the number of roots in $[-1,1]$, then, almost surely,\[\lim_{n\to \infty}\frac{R_n[-1,1]}{\log n}=\frac{1}{\pi}.\]See also [522] .




References

[Do24] Y. Do, A strong law of large numbers for real roots of random polynomials . arXiv:2403.06353 (2024).

[EO56] Erd\H{o}s, Paul and Offord, A. C., On the number of real roots of a random algebraic equation . Proc. London Math. Soc. (3) (1956), 139-160.

[Er61] Erd\H{o}s, Paul, Some unsolved problems . Magyar Tud. Akad. Mat. Kutat\'{o} Int. K\"{o}zl. (1961), 221-254.

\noindent\rule{\linewidth}{0.4pt}


%% ============================================================
\section{Problem \#522}
\label{prob:522}

\noindent\textbf{Status:} OPEN \quad|\quad \textbf{Prize:} none \quad|\quad \textbf{Updated:} 2025-12-08

\noindent\textbf{Tags:} analysis, polynomials, probability

\medskip

\noindent\textbf{Statement:}

Let $f(z)=\sum_{0\leq k\leq n} \epsilon_k z^k$ be a random polynomial, where $\epsilon_k\in \{-1,1\}$ independently uniformly at random for $0\leq k\leq n$. Is it true that, if $R_n$ is the number of roots of $f(z)$ in $\{ z\in \mathbb{C} : \lvert z\rvert \leq 1\}$, then\[\frac{R_n}{n/2}\to 1\]almost surely?




Random polynomials with independently identically distributed coefficients are sometimes called Kac polynomials - this problem considers the case of Rademacher coefficients, i.e. independent uniform $\pm 1$ values. Erd\H{o}s and Offord \cite{EO56} showed that the number of real roots of a random degree $n$ polynomial with $\pm 1$ coefficients is $(\frac{2}{\pi}+o(1))\log n$. There is some ambiguity whether Erd\H{o}s intended the coefficients to be in $\{-1,1\}$ or $\{0,1\}$ - see the comments section. A weaker version of this was solved by Yakir \cite{Ya21}, who proved that\[\frac{R_n}{n/2}\to 1\]in probability. (This weaker claim was also asked by Erd\H{o}s, and also appears in a book of Hayman \cite{Ha67}.) More precisely,\[\lim_{n\to \infty} \mathbb{P}(\lvert R_n-n/2\rvert \geq n^{9/10}) =0.\]See also [521] .




References

[EO56] Erd\H{o}s, Paul and Offord, A. C., On the number of real roots of a random algebraic equation . Proc. London Math. Soc. (3) (1956), 139-160.

[Ha67] Hayman, W. K., Research problems in function theory . (1967), vii+56.

[Ya21] Yakir, Oren, Approximately half of the roots of a random {L}ittlewood polynomial are inside the disk . Studia Math. (2021), 227--240.

\noindent\rule{\linewidth}{0.4pt}


%% ============================================================
\section{Problem \#524}
\label{prob:524}

\noindent\textbf{Status:} OPEN \quad|\quad \textbf{Prize:} none \quad|\quad \textbf{Updated:} 2025-08-31

\noindent\textbf{Tags:} analysis, probability, polynomials

\medskip

\noindent\textbf{Statement:}

For any $t\in (0,1)$ let $t=\sum_{k=1}^\infty \epsilon_k(t)2^{-k}$ (where $\epsilon_k(t)\in \{0,1\}$). What is the correct order of magnitude (for almost all $t\in(0,1)$) for\[M_n(t)=\max_{x\in [-1,1]}\left\lvert \sum_{k\leq n}(-1)^{\epsilon_k(t)}x^k\right\rvert?\]




A problem of Salem and Zygmund \cite{SaZy54}. Chung showed that, for almost all $t$, there exist infinitely many $n$ such that\[M_n(t) \ll \left(\frac{n}{\log\log n}\right)^{1/2}.\]Erd\H{o}s (unpublished) showed that for almost all $t$ and every $\epsilon&#62;0$ we have $\lim_{n\to \infty}M_n(t)/n^{1/2-\epsilon}=\infty$.




References

[SaZy54] Salem, R. and Zygmund, A., Some properties of trigonometric series whose terms have random signs . Acta Math. (1954), 245-301.

\noindent\rule{\linewidth}{0.4pt}


%% ============================================================
\section{Problem \#528}
\label{prob:528}

\noindent\textbf{Status:} OPEN \quad|\quad \textbf{Prize:} none \quad|\quad \textbf{Updated:} 2025-08-31

\noindent\textbf{Tags:} geometry

\medskip

\noindent\textbf{Statement:}

Let $f(n,k)$ count the number of self-avoiding walks of $n$ steps (beginning at the origin) in $\mathbb{Z}^k$ (i.e. those walks which do not intersect themselves). Determine\[C_k=\lim_{n\to\infty}f(n,k)^{1/n}.\]




The constant $C_k$ is sometimes known as the connective constant. Hammersley and Morton \cite{HM54} showed that this limit exists, and it is trivial that $k\leq C_k\leq 2k-1$. Kesten \cite{Ke63} proved that $C_k=2k-1-1/2k+O(1/k^2)$, and more precise asymptotics are given by Clisby, Liang, and Slade \cite{CLS07}. Conway and Guttmann \cite{CG93} showed that $C_2\geq 2.62$ and Alm \cite{Al93} showed that $C_2\leq 2.696$. Jacobsen, Scullard, and Guttmann \cite{JSG16} have computed the first few decimal places of $C_2$, showing that\[C_2 = 2.6381585303279\cdots.\]See also [529] .




References

[Al93] Alm, Sven Erick, Upper bounds for the connective constant of self-avoiding walks . Combin. Probab. Comput. (1993), 115-136.

[CG93] Conway, A. R. and Guttmann, A. J., Lower bound on the connective constant for square lattice self-avoiding walks . J. Phys. A (1993), 3719-3724.

[CLS07] Clisby, Nathan and Liang, Richard and Slade, Gordon, Self-avoiding walk enumeration via the lace expansion . J. Phys. A (2007), 10973-11017.

[HM54] Hammersley, J. M. and Morton, K. W., Poor man's Monte Carlo . J. Roy. Statist. Soc. Ser. B (1954), 23-38; discussion 61-75.

[JSG16] Jacobsen, Jesper Lykke and Scullard, Christian R. and Guttmann, Anthony J., On the growth constant for square-lattice self-avoiding walks . J. Phys. A (2016), 494004, 18.

[Ke63] Kesten, Harry, On the number of self-avoiding walks . J. Mathematical Phys. (1963), 960-969.

\noindent\rule{\linewidth}{0.4pt}


%% ============================================================
\section{Problem \#529}
\label{prob:529}

\noindent\textbf{Status:} OPEN \quad|\quad \textbf{Prize:} none \quad|\quad \textbf{Updated:} 2025-08-31

\noindent\textbf{Tags:} geometry, probability

\medskip

\noindent\textbf{Statement:}

Let $d_k(n)$ be the expected distance from the origin after taking $n$ random steps from the origin in $\mathbb{Z}^k$ (conditional on no self intersections) - that is, a self-avoiding walk . Is it true that\[\lim_{n\to \infty}\frac{d_2(n)}{n^{1/2}}= \infty?\]Is it true that\[d_k(n)\ll n^{1/2}\]for $k\geq 3$?




Slade \cite{Sl87} proved that, for $k$ sufficiently large, $d_k(n)\sim Dn^{1/2}$ for some constant $D&gt;0$ (independent of $k$). Hara and Slade (\cite{HaSl91} and \cite{HaSl92}) proved this for all $k\geq 5$. For $k=2$ Duminil-Copin and Hammond \cite{DuHa13} have proved that $d_2(n)=o(n)$. It is now conjectured that $d_k(n)\ll n^{1/2}$ is false for $k=3$ and $k=4$, and more precisely (see for example Section 1.4 of \cite{MaSl93}) that $d_2(n)\sim Dn^{3/4}$, $d_3(n)\sim n^{\nu}$ where $\nu\approx 0.59$, and $d_4(n)\sim D(\log n)^{1/8}n^{1/2}$. Madras and Slade \cite{MaSl93} have a monograph on the topic of self-avoiding walks. See also [528] .




References

[DuHa13] Duminil-Copin, Hugo and Hammond, Alan, Self-avoiding walk is sub-ballistic . Comm. Math. Phys. (2013), 401--423.

[HaSl91] Hara, Takashi and Slade, Gordon, Critical behaviour of self-avoiding walk in five or more dimensions . Bull. Amer. Math. Soc. (N.S.) (1991), 417--423.

[HaSl92] Hara, Takashi and Slade, Gordon, Self-avoiding walk in five or more dimensions. I. The critical behaviour . Comm. Math. Phys. (1992), 101--136.

[MaSl93] Madras, Neal and Slade, Gordon, The self-avoiding walk . (1993), xiv+425.

[Sl87] Slade, Gordon, The diffusion of self-avoiding random walk in high dimensions . Comm. Math. Phys. (1987), 661--683.

\noindent\rule{\linewidth}{0.4pt}


%% ============================================================
\section{Problem \#530}
\label{prob:530}

\noindent\textbf{Status:} OPEN \quad|\quad \textbf{Prize:} none \quad|\quad \textbf{Updated:} 2025-08-31

\noindent\textbf{Tags:} number theory, sidon sets

\medskip

\noindent\textbf{Statement:}

Let $\ell(N)$ be maximal such that in any finite set $A\subset \mathbb{R}$ of size $N$ there exists a Sidon subset $S$ of size $\ell(N)$ (i.e. the only solutions to $a+b=c+d$ in $S$ are the trivial ones). Determine the order of $\ell(N)$. In particular, is it true that $\ell(N)\sim N^{1/2}$?




Originally asked by Riddell \cite{Ri69}. Erd\H{o}s noted the bounds\[N^{1/3} \ll \ell(N) \leq (1+o(1))N^{1/2}\](the upper bound following from the case $A=\{1,\ldots,N\}$). The lower bound was improved to $N^{1/2}\ll \ell(N)$ by Koml\'{o}s, Sulyok, and Szemer\'{e}di \cite{KSS75}. The correct constant is unknown, but it is likely that the upper bound is true, so that $\ell(N)\sim N^{1/2}$. In \cite{AlEr85} Alon and Erd\H{o}s make the stronger conjecture that perhaps $A$ can always be written as the union of at most $(1+o(1))N^{1/2}$ many Sidon sets. (This is easily verified for $A=\{1,\ldots,N\}$ using standard constructions of Sidon sets.) This is discussed in problem C9 of Guy's collection \cite{Gu04}. See also [1088] and [1208] for a higher-dimensional generalisation.




References

[AlEr85] Alon, Noga and Erd\H{o}s, P., An application of graph theory to additive number theory . European J. Combin. (1985), 201-203.

[Gu04] Guy, Richard K., Unsolved problems in number theory . (2004), xviii+437.

[KSS75] Koml\'{o}s, J. and Sulyok, M. and Szemeredi, E., Linear problems in combinatorial number theory . Acta Math. Acad. Sci. Hungar. (1975), 113-121.

[Ri69] Riddell, J., On sets of numbers containing no $l$ terms in arithmetic progression . Nieuw Arch. Wisk. (3) (1969), 204-209.

\noindent\rule{\linewidth}{0.4pt}


%% ============================================================
\section{Problem \#531}
\label{prob:531}

\noindent\textbf{Status:} OPEN \quad|\quad \textbf{Prize:} none \quad|\quad \textbf{Updated:} 2025-08-31

\noindent\textbf{Tags:} number theory, ramsey theory

\medskip

\noindent\textbf{Statement:}

Let $F(k)$ be the minimal $N$ such that if we two-colour $\{1,\ldots,N\}$ there is a set $A$ of size $k$ such that all subset sums $\sum_{a\in S}a$ (for $\emptyset\neq S\subseteq A$) are monochromatic. Estimate $F(k)$.




The existence of $F(k)$ was established by Sanders and Folkman, and it also follows from Rado's theorem. It is commonly known as Folkman's theorem . Erd\H{o}s and Spencer \cite{ErSp89} proved that\[F(k) \geq 2^{ck^2/\log k}\]for some constant $c&gt;0$. Balogh, Eberhrad, Narayanan, Treglown, and Wagner \cite{BENTW17} have improved this to\[F(k) \geq 2^{2^{k-1}/k}.\]




References

[BENTW17] Balogh, J\'{o}zsef and Eberhard, Sean and Narayanan, Bhargav and Treglown, Andrew and Wagner, Adam Zsolt, An improved lower bound for Folkman's theorem . Bull. Lond. Math. Soc. (2017), 745-747.

[ErSp89] Erd\H{o}s, Paul and Spencer, Joel, Monochromatic sumsets . J. Combin. Theory Ser. A (1989), 162-163.

\noindent\rule{\linewidth}{0.4pt}


%% ============================================================
\section{Problem \#535}
\label{prob:535}

\noindent\textbf{Status:} OPEN \quad|\quad \textbf{Prize:} none \quad|\quad \textbf{Updated:} 2025-08-31

\noindent\textbf{Tags:} number theory

\medskip

\noindent\textbf{Statement:}

Let $r\geq 3$, and let $f_r(N)$ denote the size of the largest subset of $\{1,\ldots,N\}$ such that no subset of size $r$ has the same pairwise greatest common divisor between all elements. Estimate $f_r(N)$.




Erd\H{o}s \cite{Er64} proved that\[f_r(N) \leq N^{\frac{3}{4}+o(1)},\]and Abbott and Hanson \cite{AbHa70} improved this exponent to $1/2$. Erd\H{o}s \cite{Er64} proved the lower bound, for any $r \geq 3$,\[f_r(N) &gt; N^{\frac{c_r}{\log\log N}}\]for some constant $c_r&gt;0$, and conjectured this should also be an upper bound. This is intimately connected with the sunflower problem [20] . Indeed, Erd\H{o}s \cite{Er64} noted that a positive solution to [20] would imply\[f_r(N) \leq N^{\frac{C_r}{\log\log N}}\]for some constant $C_r&gt;0$. Using the recent bounds of Alweiss, Lovett, Wu, and Zhang \cite{ALWZ20} for the sunflower problem yields\[f_r(N) \leq N^{C_r\frac{\log\log\log N}{\log\log N}},\]and in particular $f_r(N) \leq N^{o(1)}$. An alternative description of the details of these lower and upper bounds was given by Abbott and Hanson \cite{AbHa67}. (A similar detailed calculation is given in the comments by Hrishi.) See also [536] .




References

[ALWZ20] Alweiss, R. and Lovett, S. and Wu, K. and Zhang, J., Improved bounds for the sunflower lemma . (2020).

[AbHa67] No reference found.


[AbHa70] Abbott, H. L. and Hanson, D., An extremal problem in number theory . Bull. London Math. Soc. (1970), 324-326.

[Er64] Erd\H{o}s, P., On a problem in elementary number theory and a combinatorial problem . Math. Comp. (1964), 644-646.

\noindent\rule{\linewidth}{0.4pt}


%% ============================================================
\section{Problem \#536}
\label{prob:536}

\noindent\textbf{Status:} OPEN \quad|\quad \textbf{Prize:} none \quad|\quad \textbf{Updated:} 2025-08-31

\noindent\textbf{Tags:} number theory

\medskip

\noindent\textbf{Statement:}

Let $f(N)$ be the largest size of $A\subseteq \{1,\ldots,N\}$ with the property that there are no distinct $a,b,c\in A$ such that\[[a,b]=[b,c]=[a,c],\]where $[a,b]$ denotes the least common multiple. Estimate $f(N)$ - in particular, is it true that $f(N)=o(N)$?




The corresponding quantity is $\gg N$ if we ask for four elements with the same pairwise least common multiple, as shown by Erd\H{o}s \cite{Er62} (with a proof given in \cite{Er70}). This was also asked by Erd\H{o}s at the 1991 problem session of West Coast Number Theory. Abbott and Gardner \cite{AbGa67} proved\[f(N) \geq (1-o(1))(\log\log N)\frac{N}{\log N}.\]In the comments Weisenberg sketches a proof that, for some function $\omega(N)\to \infty$ as $N\to \infty$,\[(\log\log N)^{\omega(N)}\frac{N}{\log N}\ll f(N) \leq \left(\frac{221}{225}+o(1)\right)N.\]See also [535] , [537] , and [856] . A related combinatorial problem is asked at [857] .




References

[AbGa67] Abbott, H. L. and Gardner, B., An extremal problem in number theory . Canad. Math. Bull. (1967), 173--177.

[Er62] Erd\H{o}s, P\'{a}l, Remarks on number theory. IV. Extremal problems in number theory. I . Mat. Lapok (1962), 228-255.

[Er70] Erd\H{o}s, Paul, Some extremal problems in combinatorial number theory . Mathematical Essays Dedicated to A. J. Macintyre (1970), 123-133.

\noindent\rule{\linewidth}{0.4pt}


%% ============================================================
\section{Problem \#538}
\label{prob:538}

\noindent\textbf{Status:} OPEN \quad|\quad \textbf{Prize:} none \quad|\quad \textbf{Updated:} 2025-08-31

\noindent\textbf{Tags:} number theory

\medskip

\noindent\textbf{Statement:}

Let $r\geq 2$ and suppose that $A\subseteq\{1,\ldots,N\}$ is such that, for any $m$, there are at most $r$ solutions to $m=pa$ where $p$ is prime and $a\in A$. Give the best possible upper bound for\[\sum_{n\in A}\frac{1}{n}.\]




Erd\H{o}s observed that\[\sum_{n\in A}\frac{1}{n}\sum_{p\leq N}\frac{1}{p}\leq r\sum_{m\leq N^2}\frac{1}{m}\ll r\log N,\]and hence\[\sum_{n\in A}\frac{1}{n} \ll r\frac{\log N}{\log\log N}.\]See also [536] and [537] .

\noindent\rule{\linewidth}{0.4pt}


%% ============================================================
\section{Problem \#539}
\label{prob:539}

\noindent\textbf{Status:} OPEN \quad|\quad \textbf{Prize:} none \quad|\quad \textbf{Updated:} 2025-08-31

\noindent\textbf{Tags:} number theory

\medskip

\noindent\textbf{Statement:}

Let $h(n)$ be such that, for any set $A\subseteq \mathbb{N}$ of size $n$, the set\[\left\{ \frac{a}{(a,b)}: a,b\in A\right\}\]has size at least $h(n)$. Estimate $h(n)$.




Erd\H{o}s and Szemer\'{e}di proved that\[n^{1/2} \ll h(n) \ll n^{1-c}\]for some constant $c&#62;0$. The upper bound has been improved to $h(n)\ll n^{2/3}$ by Freiman and Lev. A proof of these bounds can be found in a paper of Granville and Roesler \cite{GrRo99}. Granville and Roesler also reformulate this problem as one in combinatorial geometry: given $A\subseteq \mathbb{Z}^d$ of size $n$ what is the minimum size possible of\[\{ \delta(\mathbf{a},\mathbf{b}) : \mathbf{a},\mathbf{b}\in A\}\]where $\delta(\mathbf{a},\mathbf{b})$ has $i$th coordinate $\max(0,a_i-b_i)$. Using this formulation that obtain improved lower bounds in fixed small dimension $d$.




References

[GrRo99] Granville, Andrew and Roesler, Friedrich, The set of differences of a given set . Amer. Math. Monthly (1999), 338--344.

\noindent\rule{\linewidth}{0.4pt}


%% ============================================================
\section{Problem \#544}
\label{prob:544}

\noindent\textbf{Status:} OPEN \quad|\quad \textbf{Prize:} none \quad|\quad \textbf{Updated:} 2025-08-31

\noindent\textbf{Tags:} graph theory, ramsey theory

\medskip

\noindent\textbf{Statement:}

Show that\[R(3,k+1)-R(3,k)\to\infty\]as $k\to \infty$. Similarly, prove or disprove that\[R(3,k+1)-R(3,k)=o(k).\]




A problem of Erd\H{o}s and S\'{o}s. It is known (see [165] ) that\[R(3,k) \asymp \frac{k^2}{\log k}.\]The OpenAI bound of [1014] implies that\[R(3,k+1)-R(3,k)\ll k^{-c}R(3,k)\]for some constant $c&gt;0$. This problem is #8 in Ramsey Theory in the graphs problem collection. See also [165] and [1014] .

\noindent\rule{\linewidth}{0.4pt}


%% ============================================================
\section{Problem \#545}
\label{prob:545}

\noindent\textbf{Status:} OPEN \quad|\quad \textbf{Prize:} none \quad|\quad \textbf{Updated:} 2025-12-02

\noindent\textbf{Tags:} graph theory, ramsey theory

\medskip

\noindent\textbf{Statement:}

Let $G$ be a graph with $m$ edges and no isolated vertices. Is the Ramsey number $R(G)$ maximised when $G$ is 'as complete as possible'? That is, if $m=\binom{n}{2}+t$ edges with $0\leq t&#60;n$ then is\[R(G)\leq R(H),\]where $H$ is the graph formed by connecting a new vertex to $t$ of the vertices of $K_n$?




A question of Erd\H{o}s and Graham. The weaker question of whether\[R(G) \leq 2^{O(m^{1/2})}\]is the subject of [546] . (This is true, and was proved by Sudakov \cite{Su11}.) LouisD in the comments has noted this fails for small $m$ (in particular for $2\leq m\leq 5$ and $7\leq m\leq 9$). This problem is #10 in Ramsey Theory in the graphs problem collection.




References

[Su11] Sudakov, Benny, A conjecture of Erd\H{o}s on graph Ramsey numbers . Adv. Math. (2011), 601-609.

\noindent\rule{\linewidth}{0.4pt}


%% ============================================================
\section{Problem \#550}
\label{prob:550}

\noindent\textbf{Status:} OPEN \quad|\quad \textbf{Prize:} none \quad|\quad \textbf{Updated:} 2025-08-31

\noindent\textbf{Tags:} graph theory, ramsey theory

\medskip

\noindent\textbf{Statement:}

Let $m_1\leq\cdots\leq m_k$ and $n$ be sufficiently large. If $T$ is a tree on $n$ vertices and $G$ is the complete multipartite graph with vertex class sizes $m_1,\ldots,m_k$ then prove that\[R(T,G)\leq (\chi(G)-1)(R(T,K_{m_1,m_2})-1)+m_1.\]




Chv\'{a}tal \cite{Ch77} proved that $R(T,K_m)=(m-1)(n-1)+1$. This problem is #16 in Ramsey Theory in the graphs problem collection.




References

[Ch77] Chv\'{a}tal, V., Tree-complete graph Ramsey numbers . J. Graph Theory (1977), 93.

\noindent\rule{\linewidth}{0.4pt}


%% ============================================================
\section{Problem \#552}
\label{prob:552}

\noindent\textbf{Status:} OPEN \quad|\quad \textbf{Prize:} none \quad|\quad \textbf{Updated:} 2025-08-31

\noindent\textbf{Tags:} graph theory, ramsey theory

\medskip

\noindent\textbf{Statement:}

Determine the Ramsey number\[R(C_4,S_n),\]where $S_n=K_{1,n}$ is the star on $n+1$ vertices. In particular, is it true that, for any $c&#62;0$, there are infinitely many $n$ such that\[R(C_4,S_n)\leq n+\sqrt{n}-c?\]




A problem of Burr, Erd\H{o}s, Faudree, Rousseau, and Schelp \cite{BEFRS89}. Erd\H{o}s often asked about $R(C_4,S_n)$ in the equivalent formulation of asking for a bound on the minimum degree of a graph which would guarantee the existence of a $C_4$ (see [85] ). It is known that\[ n+\sqrt{n}-6n^{11/40} \leq R(C_4,S_n)\leq n+\lceil\sqrt{n}\rceil+1.\]The lower bound is due to \cite{BEFRS89}, the upper bound is due to Parsons \cite{Pa75}. The lower bound of \cite{BEFRS89} is related to gaps between primes, and assuming e.g. Cramer's conjecture on gaps between primes their lower bound would be $n+\sqrt{n}-n^{o(1)}$. Erd\H{o}s offered \$100 for a proof or disproof of the second question in \cite{BEFRS89}. In \cite{Er96} Erd\H{o}s asks (an equivalent formulation of) whether $R(C_4,S_n)\geq n+\sqrt{n}-O(1)$, but says this is probably 'too optimistic'. They also ask, if $f(n)=R(C_4,S_n)$, whether $f(n+1)=f(n)$ infinitely often, and is the density of such $n$ $0$? Also, is it true that $f(n+1)\leq f(n)+2$ for all $n$? A similar question about an equivalent function is the subject of [85] . Parsons \cite{Pa75} proved that\[R(C_4,S_n)=n+\lceil\sqrt{n}\rceil\]whenever $n=q^2+1$ for a prime power $q$ and\[R(C_4,S_n)=n+\lceil\sqrt{n}\rceil+1\]whenever $n=q^2$ for a prime power $q$ (in particular both equalities occur infinitely often). This has been extended in various works, all in the cases $n=q^2\pm t$ for some $0\leq t\leq q$ and prime power $q$. We refer to the work of Parsons \cite{Pa75}, Wu, Sun, Zhang, and Radziszowski \cite{WSZR15}, and Zhang, Chen, and Cheng (\cite{ZCC17} and \cite{ZCC17b}) for a precise description. In every known case\[R(C_4,S_n)=n+\lceil\sqrt{n}\rceil+\{0,1\},\]and Zhang, Chen, and Cheng \cite{ZCC17} speculate whether this is in fact true for all $n\geq 2$ (whence the answer to the question above would be no). This problem is #19 in Ramsey Theory in the graphs problem collection.




References

[BEFRS89] Burr, S. and Erd\H{o}s, P. and Faudree, R. J. and Rousseau, C. C. and Schelp, R. H., Some complete bipartite graph-tree Ramsey numbers . Graph theory in memory of G. A. Dirac (Sandbjerg, 1985) (1989), 79-89.

[Er96] Erd\H{o}s, Paul, Some of my favourite problems on cycles and colourings . Tatra Mt. Math. Publ. (1996), 7-9.

[Pa75] Parsons, T. D., Ramsey graphs and block designs. I . Trans. Amer. Math. Soc. (1975), 33--44.

[WSZR15] Wu, Yali and Sun, Yongqi and Zhang, Rui and Radziszowski, Stanis\l aw P., Ramsey numbers of {$C_4$} versus wheels and stars . Graphs Combin. (2015), 2437--2446.

[ZCC17] Zhang, Xuemei and Chen, Yaojun and Cheng, T. C. Edwin, Some values of {R}amsey numbers for {$C_4$} versus stars . Finite Fields Appl. (2017), 73--85.

[ZCC17b] Zhang, Xuemei and Chen, Yaojun and Cheng, T. C. Edwin, Polarity graphs and {R}amsey numbers for {$C_4$} versus stars . Discrete Math. (2017), 655--660.

\noindent\rule{\linewidth}{0.4pt}


%% ============================================================
\section{Problem \#554}
\label{prob:554}

\noindent\textbf{Status:} OPEN \quad|\quad \textbf{Prize:} none \quad|\quad \textbf{Updated:} 2025-08-31

\noindent\textbf{Tags:} graph theory, ramsey theory

\medskip

\noindent\textbf{Statement:}

Let $R_k(G)$ denote the minimal $m$ such that if the edges of $K_m$ are $k$-coloured then there is a monochromatic copy of $G$. Show that\[\lim_{k\to \infty}\frac{R_k(C_{2n+1})}{R_k(K_3)}=0\]for any $n\geq 2$.




A problem of Erd\H{o}s and Graham. The problem is open even for $n=2$. Schur \cite{Sc16} proved\[C^k \ll R_k(K_3) \ll k!\]for some constant $C&gt;0$. Erd\H{o}s conjectured (see [183] ) that $R_k(K_3) \leq C^k$ for some contant $C&gt;0$. Bondy and Erd\H{o}s \cite{BoEr73} and Erd\H{o}s and Graham \cite{ErGr75} proved\[n2^k+1\leq R_k(C_{2n+1})\leq 2n(k+2)!.\]The lower bound is sharp for fixed $k$ and large enough $n$, as shown by Jenssen and Skokan \cite{JeSk21}. Day and Johnson \cite{DaJo17} have proved that, for fixed $n$,\[R_k(C_{2n+1})\geq 2n(2+c_n)^{k-1}\]for all large $k$, for some constant $c_n&gt;0$. Axenovich, Cames van Batenburg, Janzer, Michel, and Rundstr\"{o}m \cite{ACJMR25} have improved the upper bound to\[R_k(C_{2n+1}) \leq (4n-2)^k k^{k/n}+1.\]In particular, there exists some absolute constant $c&gt;0$ such that\[R_k(C_{2n+1})\leq (Cn)^kk!^{1/n}.\]This problem is #23 in Ramsey Theory in the graphs problem collection.




References

[ACJMR25] M. Axenovich, W. Cames van Batenburg, O. Janzer, L. Michel, and M. Rundstr\"{o}m, An improved upper bound for the multicolour Ramsey number of odd cycles . arXiv:2510.17981 (2025).

[BoEr73] Bondy, J. A. and Erd\H{o}s, P., Ramsey numbers for cycles in graphs . J. Combinatorial Theory Ser. B (1973), 46-54.

[DaJo17] Day, A. Nicholas and Johnson, J. Robert, Multicolour {R}amsey numbers of odd cycles . J. Combin. Theory Ser. B (2017), 56--63.

[ErGr75] Erd\H{o}s, P. and Graham, R. L., On partition theorems for finite graphs . Infinite and finite sets (Colloq., Keszthely, 1973; dedicated to P. Erd\H{o}s on his 60th birthday), Vols. I, II, III (1975), 515-527.

[JeSk21] Jenssen, Matthew and Skokan, Jozef, Exact {R}amsey numbers of odd cycles via nonlinear optimisation . Adv. Math. (2021), Paper No. 107444, 46.

[Sc16] I. Schur, \"{U}ber die kongruenz $x^m+y^m=z^m\pmod{p}$ . Jahresbericht der Deutschen Mathematiker-Vereinigung (1916), 114-117.

\noindent\rule{\linewidth}{0.4pt}


%% ============================================================
\section{Problem \#555}
\label{prob:555}

\noindent\textbf{Status:} OPEN \quad|\quad \textbf{Prize:} none \quad|\quad \textbf{Updated:} 2025-08-31

\noindent\textbf{Tags:} graph theory, ramsey theory

\medskip

\noindent\textbf{Statement:}

Let $R_k(G)$ denote the minimal $m$ such that if the edges of $K_m$ are $k$-coloured then there is a monochromatic copy of $G$. Determine the value of\[R_k(C_{2n}).\]




A problem of Erd\H{o}s and Graham. Erd\H{o}s \cite{Er81c} gives the bounds\[k^{1+\frac{1}{2n}}\ll R_k(C_{2n})\ll k^{1+\frac{1}{n-1}}.\]Chung and Graham \cite{ChGr75} showed that\[R_k(C_4)&gt;k^2-k+1\]when $k-1$ is a prime power and\[R_k(C_4)\leq k^2+k+1\]for all $k$. This problem is #24 in Ramsey Theory in the graphs problem collection.




References

[ChGr75] Chung, Fan R. K. and Graham, R. L., On multicolor Ramsey numbers for complete bipartite graphs . J. Combinatorial Theory Ser. B (1975), 164-169.

[Er81c] Erd\H{o}s, Paul, Some new problems and results in graph theory and other branches of combinatorial mathematics . Combinatorics and graph theory (1981), 9-17.

\noindent\rule{\linewidth}{0.4pt}


%% ============================================================
\section{Problem \#557}
\label{prob:557}

\noindent\textbf{Status:} OPEN \quad|\quad \textbf{Prize:} none \quad|\quad \textbf{Updated:} 2025-08-31

\noindent\textbf{Tags:} graph theory, ramsey theory

\medskip

\noindent\textbf{Statement:}

Let $R_k(G)$ denote the minimal $m$ such that if the edges of $K_m$ are $k$-coloured then there is a monochromatic copy of $G$. Is it true that\[R_k(T)\leq kn+O(1)\]for any tree $T$ on $n$ vertices?




A problem of Erd\H{o}s and Graham. Implied by [548] . This would be best possible since, for example, $R_k(S_n)\geq kn-O(k)$ if $S_n=K_{1,n-1}$ is a star on $n$ vertices. This problem is #26 in Ramsey Theory in the graphs problem collection.

\noindent\rule{\linewidth}{0.4pt}


%% ============================================================
\section{Problem \#558}
\label{prob:558}

\noindent\textbf{Status:} OPEN \quad|\quad \textbf{Prize:} none \quad|\quad \textbf{Updated:} 2025-08-31

\noindent\textbf{Tags:} graph theory, ramsey theory

\medskip

\noindent\textbf{Statement:}

Let $R_k(G)$ denote the minimal $m$ such that if the edges of $K_m$ are $k$-coloured then there is a monochromatic copy of $G$. Determine\[R_k(K_{s,t})\]where $K_{s,t}$ is the complete bipartite graph with $s$ vertices in one component and $t$ in the other.




Chung and Graham \cite{ChGr75} prove the general bounds\[(2\pi\sqrt{st})^{\frac{1}{s+t}}\left(\frac{s+t}{e^2}\right)k^{\frac{st-1}{s+t}}\leq R_k(K_{s,t})\leq (t-1)(k+k^{1/s})^s\]and determined\[R_k(K_{2,2})=(1+o(1))k^2.\]Alon, R\'{o}nyai, and Szab\'{o} \cite{ARS99} have proved that\[R_k(K_{3,3})=(1+o(1))k^3\]and that if $s\geq (t-1)!+1$ then\[R_k(K_{s,t})\asymp k^t.\]This problem is #27 in Ramsey Theory in the graphs problem collection.




References

[ARS99] Alon, Noga and R\'{o}nyai, Lajos and Szab\'{o}, Tibor, Norm-graphs: variations and applications . J. Combin. Theory Ser. B (1999), 280-290.

[ChGr75] Chung, Fan R. K. and Graham, R. L., On multicolor Ramsey numbers for complete bipartite graphs . J. Combinatorial Theory Ser. B (1975), 164-169.

\noindent\rule{\linewidth}{0.4pt}


%% ============================================================
\section{Problem \#560}
\label{prob:560}

\noindent\textbf{Status:} OPEN \quad|\quad \textbf{Prize:} none \quad|\quad \textbf{Updated:} 2025-08-31

\noindent\textbf{Tags:} graph theory, ramsey theory

\medskip

\noindent\textbf{Statement:}

Let $\hat{R}(G)$ denote the size Ramsey number, the minimal number of edges $m$ such that there is a graph $H$ with $m$ edges such that in any $2$-colouring of the edges of $H$ there is a monochromatic copy of $G$. Determine\[\hat{R}(K_{n,n}),\]where $K_{n,n}$ is the complete bipartite graph with $n$ vertices in each component.




We know that\[\frac{1}{60}n^22^n&lt;\hat{R}(K_{n,n})&lt; \frac{3}{2}n^32^n.\]The lower bound (which holds for $n\geq 6$) was proved by Erd\H{o}s and Rousseau \cite{ErRo93}. The upper bound was proved by Erd\H{o}s, Faudree, Rousseau, and Schelp \cite{EFRS78b} and Ne\v{s}et\v{r}il and R\"{o}dl \cite{NeRo78}. Conlon, Fox, and Wigderson \cite{CFW23} have proved that, for any $s\leq t$,\[\hat{R}(K_{s,t})\gg s^{2-\frac{s}{t}}t2^s,\]and prove that when $t\gg s\log s$ we have $\hat{R}(K_{s,t})\asymp s^2t2^s$. They conjecture that this should hold for all $s\leq t$, and so in particular we should have $\hat{R}(K_{n,n})\asymp n^32^n$. This problem is #29 in Ramsey Theory in the graphs problem collection.




References

[CFW23] Conlon, David and Fox, Jacob and Wigderson, Yuval, Three early problems on size Ramsey numbers . Combinatorica (2023), 743-768.

[EFRS78b] Erd\H{o}s, P. and Faudree, R. J. and Rousseau, C. C. and Schelp, R. H., The size Ramsey number . Period. Math. Hungar. (1978), 145-161.

[ErRo93] Erd\H{o}s, P. and Rousseau, C. C., The size Ramsey number of a complete bipartite graph . Discrete Math. (1993), 259-262.

[NeRo78] Ne\vSet\v{r}il, J. and R\"{o}dl, V., The structure of critical Ramsey graphs . Acta Math. Acad. Sci. Hungar. (1978), 295-300.

\noindent\rule{\linewidth}{0.4pt}


%% ============================================================
\section{Problem \#561}
\label{prob:561}

\noindent\textbf{Status:} OPEN \quad|\quad \textbf{Prize:} none \quad|\quad \textbf{Updated:} 2025-08-31

\noindent\textbf{Tags:} graph theory, ramsey theory

\medskip

\noindent\textbf{Statement:}

Let $\hat{R}(G)$ denote the size Ramsey number, the minimal number of edges $m$ such that there is a graph $H$ with $m$ edges such that in any $2$-colouring of the edges of $H$ there is a monochromatic copy of $G$. Let $F_1$ and $F_2$ be the union of stars. More precisely, let $F_1=\cup_{i\leq s} K_{1,n_i}$ and $F_2=\cup_{j\leq t} K_{1,m_j}$ with $n_1\geq \cdots \geq n_s\geq 1$ and $m_1\geq \cdots \geq m_t\geq 1$. Prove that\[\hat{R}(F_1,F_2) = \sum_{2\leq k\leq s+t}l_k\]where\[l_k=\max\{n_i+m_j-1 : i+j=k\}.\]




Burr, Erd\H{o}s, Faudree, Rousseau, and Schelp \cite{BEFRS78} proved this when all the $n_i$ are identical and all the $m_i$ are identical. Gy\H{o}ri and Schelp \cite{GySc02} proved this when $\binom{l_k}{2}&gt; \sum_{k\leq i\leq s+t}l_i$ for all $2\leq k\leq s+t$. More special cases, including all cases when $s=1$, or $s=2$ and $n_1=n_2$, or all $n_i$ and $m_j$ are odd, or all $n_i$ are an identical odd number and $m_1$ is odd, were proved by Davoodi, Javadi, Kamranian, and Raeisi \cite{DJKR25}. This problem is #30 in Ramsey Theory in the graphs problem collection.




References

[BEFRS78] Burr, S. A. and Erd\H{o}s, P. and Faudree, R. J. and Rousseau, C. C. and Schelp, R. H., Ramsey-minimal graphs for multiple copies . Nederl. Akad. Wetensch. Indag. Math. (1978), 187-195.

[DJKR25] Davoodi, Akbar and Javadi, Ramin and Kamranian, Azam and Raeisi, Ghaffar, On a conjecture of {E}rd\H{o}s on size {R}amsey number of star forests . Ars Math. Contemp. (2025), Paper No. 9, 10.

[GySc02] Gy\H ori, Ervin and Schelp, R. H., Two-edge colorings of graphs with bounded degree in both colors . Discrete Math. (2002), 105--110.

\noindent\rule{\linewidth}{0.4pt}


%% ============================================================
\section{Problem \#562}
\label{prob:562}

\noindent\textbf{Status:} OPEN \quad|\quad \textbf{Prize:} none \quad|\quad \textbf{Updated:} 2025-08-31

\noindent\textbf{Tags:} graph theory, ramsey theory, hypergraphs

\medskip

\noindent\textbf{Statement:}

Let $R_r(n)$ denote the $r$-uniform hypergraph Ramsey number: the minimal $m$ such that if we $2$-colour all edges of the complete $r$-uniform hypergraph on $m$ vertices then there must be some monochromatic copy of the complete $r$-uniform hypergraph on $n$ vertices. Prove that, for $r\geq 3$,\[\log_{r-1} R_r(n) \asymp_r n,\]where $\log_{r-1}$ denotes the $(r-1)$-fold iterated logarithm. That is, does $R_r(n)$ grow like\[2^{2^{\cdots n}}\]where the tower of exponentials has height $r-1$?




A problem of Erd\H{o}s, Hajnal, and Rado \cite{EHR65}. A generalisation of [564] . This problem is #38 in Ramsey Theory in the graphs problem collection.




References

[EHR65] Erd\H{o}s, P. and Hajnal, A. and Rado, R., Partition relations for cardinal numbers . Acta Math. Acad. Sci. Hungar. (1965), 93-196.

\noindent\rule{\linewidth}{0.4pt}


%% ============================================================
\section{Problem \#563}
\label{prob:563}

\noindent\textbf{Status:} OPEN \quad|\quad \textbf{Prize:} none \quad|\quad \textbf{Updated:} 2025-08-31

\noindent\textbf{Tags:} graph theory, ramsey theory, hypergraphs

\medskip

\noindent\textbf{Statement:}

Let $F(n,\alpha)$ denote the smallest $m$ such that there exists a $2$-colouring of the edges of $K_n$ so that every $X\subseteq [n]$ with $\lvert X\rvert\geq m$ contains more than $\alpha \binom{\lvert X\rvert}{2}$ many edges of each colour. Prove that, for every $0\leq \alpha&#60; 1/2$,\[F(n,\alpha)\sim c_\alpha\log n\]for some constant $c_\alpha$ depending only on $\alpha$.




It is easy to show via the probabilistic method that, for every $0\leq \alpha&lt;1/2$,\[F(n,\alpha)\asymp_\alpha \log n.\]Note that when $\alpha=0$ this is just asking for a $2$-colouring of the edges of $K_n$ which contains no monochromatic clique of size $m$, and hence we recover the classical Ramsey numbers. See also [161] for a generalisation to hypergraphs. This problem is #39 in Ramsey Theory in the graphs problem collection.

\noindent\rule{\linewidth}{0.4pt}


%% ============================================================
\section{Problem \#564}
\label{prob:564}

\noindent\textbf{Status:} OPEN \quad|\quad \textbf{Prize:} \$500 \quad|\quad \textbf{Updated:} 2025-08-31

\noindent\textbf{Tags:} graph theory, ramsey theory, hypergraphs

\medskip

\noindent\textbf{Statement:}

Let $R_3(n)$ be the minimal $m$ such that if the edges of the $3$-uniform hypergraph on $m$ vertices are $2$-coloured then there is a monochromatic copy of the complete $3$-uniform hypergraph on $n$ vertices. Is there some constant $c&#62;0$ such that\[R_3(n) \geq 2^{2^{cn}}?\]




A special case of [562] . A problem of Erd\H{o}s, Hajnal, and Rado \cite{EHR65}, who prove the bounds\[2^{cn^2}&lt; R_3(n)&lt; 2^{2^{n}}\]for some constant $c&gt;0$. Erd\H{o}s, Hajnal, M\'{a}t\'{e}, and Rado \cite{EHMR84} have proved a doubly exponential lower bound for the corresponding problem with $4$ colours. This problem is #37 in Ramsey Theory in the graphs problem collection.




References

[EHMR84] Erd\H{o}s, Paul and Hajnal, Andr\'{a}s and M\'{a}t\'{e}, Attila and Rado, Richard, Combinatorial set theory: partition relations for cardinals . (1984), 347.

[EHR65] Erd\H{o}s, P. and Hajnal, A. and Rado, R., Partition relations for cardinal numbers . Acta Math. Acad. Sci. Hungar. (1965), 93-196.

\noindent\rule{\linewidth}{0.4pt}


%% ============================================================
\section{Problem \#566}
\label{prob:566}

\noindent\textbf{Status:} OPEN \quad|\quad \textbf{Prize:} none \quad|\quad \textbf{Updated:} 2025-08-31

\noindent\textbf{Tags:} graph theory, ramsey theory

\medskip

\noindent\textbf{Statement:}

Let $G$ be such that any subgraph on $k$ vertices has at most $2k-3$ edges. Is it true that, if $H$ has $m$ edges and no isolated vertices, then\[R(G,H)\ll m?\]




In other words, is $G$ Ramsey size linear? This fails for a graph $G$ with $n$ vertices and $2n-2$ edges (for example with $H=K_n$). Erd\H{o}s, Faudree, Rousseau, and Schelp \cite{EFRS93} have shown that any graph $G$ with $n$ vertices and at most $n+1$ edges is Ramsey size linear. Implies [567] . This problem is #31 in Ramsey Theory in the graphs problem collection.




References

[EFRS93] Erd\H{o}s, Paul and Faudree, R. J. and Rousseau, C. C. and Schelp, R. H., Ramsey size linear graphs . Combin. Probab. Comput. (1993), 389-399.

\noindent\rule{\linewidth}{0.4pt}


%% ============================================================
\section{Problem \#567}
\label{prob:567}

\noindent\textbf{Status:} OPEN \quad|\quad \textbf{Prize:} none \quad|\quad \textbf{Updated:} 2025-08-31

\noindent\textbf{Tags:} graph theory, ramsey theory

\medskip

\noindent\textbf{Statement:}

Let $G$ be either $Q_3$ or $K_{3,3}$ or $H_5$ (the last formed by adding two vertex-disjoint chords to $C_5$). Is it true that, if $H$ has $m$ edges and no isolated vertices, then\[R(G,H)\ll m?\]




In other words, is $G$ Ramsey size linear? A special case of [566] . In \cite{Er95} Erd\H{o}s specifically asks about the case $G=K_{3,3}$. The graph $H_5$ can also be described as $K_4^*$, obtained from $K_4$ by subdividing one edge. ($K_4$ itself is not Ramsey size linear, since $R(4,n)\gg n^{3-o(1)}$, see [166] .) Brada\'{c}, Gishboliner, and Sudakov \cite{BGS23} have shown that every subdivision of $K_4$ on at least $6$ vertices is Ramsey size linear, and also that $R(H_5,H) \ll m$ whenever $H$ is a bipartite graph with $m$ edges and no isolated vertices. This problem is #32 in Ramsey Theory in the graphs problem collection.




References

[BGS23] Brada\'C, D. and Gishboliner, L. and Sudakov, B., On Ramsey size-linear graphs and related questions . arXiv:2202.10388 (2023).

[Er95] Erd\H{o}s, Paul, Some of my favourite problems in number theory, combinatorics, and geometry . Resenhas (1995), 165-186.

\noindent\rule{\linewidth}{0.4pt}


%% ============================================================
\section{Problem \#568}
\label{prob:568}

\noindent\textbf{Status:} OPEN \quad|\quad \textbf{Prize:} none \quad|\quad \textbf{Updated:} 2025-08-31

\noindent\textbf{Tags:} graph theory, ramsey theory

\medskip

\noindent\textbf{Statement:}

Let $G$ be a graph such that $R(G,T_n)\ll n$ for any tree $T_n$ on $n$ vertices and $R(G,K_n)\ll n^2$. Is it true that, for any $H$ with $m$ edges and no isolated vertices,\[R(G,H)\ll m?\]




In other words, is $G$ Ramsey size linear? This problem is #33 in Ramsey Theory in the graphs problem collection.

\noindent\rule{\linewidth}{0.4pt}


%% ============================================================
\section{Problem \#569}
\label{prob:569}

\noindent\textbf{Status:} OPEN \quad|\quad \textbf{Prize:} none \quad|\quad \textbf{Updated:} 2025-08-31

\noindent\textbf{Tags:} graph theory, ramsey theory

\medskip

\noindent\textbf{Statement:}

Let $k\geq 1$. What is the best possible $c_k$ such that\[R(C_{2k+1},H)\leq c_k m\]for any graph $H$ on $m$ edges without isolated vertices?




This problem is #34 in Ramsey Theory in the graphs problem collection.

\noindent\rule{\linewidth}{0.4pt}


%% ============================================================
\section{Problem \#571}
\label{prob:571}

\noindent\textbf{Status:} OPEN \quad|\quad \textbf{Prize:} none \quad|\quad \textbf{Updated:} 2025-08-31

\noindent\textbf{Tags:} graph theory, turan number

\medskip

\noindent\textbf{Statement:}

Show that for any rational $\alpha \in [1,2)$ there exists a bipartite graph $G$ such that\[\mathrm{ex}(n;G)\asymp n^{\alpha}.\]




A problem of Erd\H{o}s and Simonovits. In \cite{Er78} Erd\H{o}s wrote 'I am not entirely sure that a trivial counterexample can not be found'. Bukh and Conlon \cite{BuCo18} proved that this holds if we weaken asking for the extremal number of a single graph to asking for the extremal number of a finite family of graphs. A rational $\alpha\in [1,2)$ for which this holds is known as a Tur\'{a}n exponent. Known Tur\'{a}n exponents are: {UL} {LI} $\frac{3}{2}-\frac{1}{2s}$ for $s\geq 2$ (Conlon, Janzer, and Lee \cite{CJL21}).{/LI} {LI} $\frac{4}{3}-\frac{1}{3s}$ and $\frac{5}{4}-\frac{1}{4s}$ for $s\geq 2$ (Jiang and Qiu \cite{JiQi20}).{/LI} {LI} $2-\frac{a}{b}$ for $\lfloor b/a\rfloor^3 \leq a\leq \frac{b}{\lfloor b/a\rfloor+1}+1$ (Jiang, Jiang, and Ma \cite{JJM20}).{/LI} {LI} $2-\frac{a}{b}$ with $b&gt;a\geq 1$ and $b\equiv \pm 1\pmod{a}$ (Kang, Kim, and Liu \cite{KKL21}).{/LI} {LI} $1+a/b$ with $b&gt;a^2$ (Jiang and Qiu \cite{JiQi23}),{/LI} {LI} $2-\frac{2}{2b+1}$ for $b\geq 2$ or $7/5$ (Jiang, Ma, and Yepremyan \cite{JMY22}).{/LI} {LI} $2-a/b$ with $b\geq (a-1)^2$ (Conlon and Janzer \cite{CoJa22}).{/LI} {/UL} See also [713] . This problem is #45 in Extremal Graph Theory in the graphs problem collection.




References

[BuCo18] Bukh, Boris and Conlon, David, Rational exponents in extremal graph theory . J. Eur. Math. Soc. (JEMS) (2018), 1747-1757.

[CJL21] Conlon, David and Janzer, Oliver and Lee, Joonkyung, More on the extremal number of subdivisions . Combinatorica (2021), 465-494.

[CoJa22] Conlon, David and Janzer, Oliver, Rational exponents near two . Adv. Comb. (2022), Paper No. 9, 10.

[Er78] Erd\H{o}s, Paul, Problems and results in combinatorial analysis and combinatorial number theory . Proceedings of the Ninth Southeastern Conference on Combinatorics, Graph Theory, and Computing (Florida Atlantic Univ., Boca Raton, Fla., 1978) (1978), 29-40.

[JJM20] Jiang, Tao and Jiang, Zilin and Ma, Jie, Negligible obstructions and Tur\'{a}n exponents . arXiv:2007.02975 (2020).

[JMY22] Jiang, Tao and Ma, Jie and Yepremyan, Liana, On Tur\'{a}n exponents of bipartite graphs . Combin. Probab. Comput. (2022), 333-344.

[JiQi20] Jiang, Tao and Qiu, Yu, Tur\'{a}n numbers of bipartite subdivisions . SIAM J. Discrete Math. (2020), 556-570.

[JiQi23] Jiang, Tao and Qiu, Yu, Many Tur\'{a}n exponents via subdivisions . Combin. Probab. Comput. (2023), 134-150.

[KKL21] Kang, Dong Yeap and Kim, Jaehoon and Liu, Hong, On the rational Tur\'{a}n exponents conjecture . J. Combin. Theory Ser. B (2021), 149-172.

\noindent\rule{\linewidth}{0.4pt}


%% ============================================================
\section{Problem \#572}
\label{prob:572}

\noindent\textbf{Status:} OPEN \quad|\quad \textbf{Prize:} none \quad|\quad \textbf{Updated:} 2025-08-31

\noindent\textbf{Tags:} graph theory, turan number

\medskip

\noindent\textbf{Statement:}

Show that for $k\geq 3$\[\mathrm{ex}(n;C_{2k})\gg n^{1+\frac{1}{k}}.\]




It is easy to see that $\mathrm{ex}(n;C_{2k+1})=\lfloor n^2/4\rfloor$ for any $k\geq 1$ (and $n&gt;2k+1$) (since no bipartite graph contains an odd cycle). Erd\H{o}s and Klein \cite{Er38} proved $\mathrm{ex}(n;C_4)\asymp n^{3/2}$. Erd\H{o}s \cite{Er64c} and Bondy and Simonovits \cite{BoSi74} showed that\[\mathrm{ex}(n;C_{2k})\ll kn^{1+\frac{1}{k}}.\]Benson \cite{Be66} has proved this conjecture for $k=3$ and $k=5$. Lazebnik, Ustimenko, and Woldar \cite{LUW95} have shown that, for arbitrary $k\geq 3$,\[\mathrm{ex}(n;C_{2k})\gg n^{1+\frac{2}{3k-3+\nu}},\]where $\nu=0$ if $k$ is odd and $\nu=1$ if $k$ is even. See \cite{LUW99} for further history and references. See also [765] . This problem is #46 in Extremal Graph Theory in the graphs problem collection.




References

[Be66] Benson, Clark T., Minimal regular graphs of girths eight and twelve . Canadian J. Math. (1966), 1091-1094.

[BoSi74] Bondy, J. A. and Simonovits, M., Cycles of even length in graphs . J. Combinatorial Theory Ser. B (1974), 97-105.

[Er38] P. Erd\H{o}s, On sequences of integers no one of which divides the product of two others and on related problems . Tomsk. Gos. Univ. Ucen Zap. (1938), 74-82.

[Er64c] Erd\H{o}s, P., Extremal problems in graph theory . Theory of Graphs and its Applications (Proc. Sympos. Smolenice, 1963) (1964), 29-36.

[LUW95] Lazebnik, F. and Ustimenko, V. A. and Woldar, A. J., A new series of dense graphs of high girth . Bull. Amer. Math. Soc. (N.S.) (1995), 73-79.

[LUW99] Lazebnik, Felix and Ustimenko, Vasiliy A. and Woldar, Andrew J., Polarities and $2k$-cycle-free graphs . Discrete Math. (1999), 503-513.

\noindent\rule{\linewidth}{0.4pt}


%% ============================================================
\section{Problem \#573}
\label{prob:573}

\noindent\textbf{Status:} OPEN \quad|\quad \textbf{Prize:} none \quad|\quad \textbf{Updated:} 2025-08-31

\noindent\textbf{Tags:} graph theory, turan number

\medskip

\noindent\textbf{Statement:}

Is it true that\[\mathrm{ex}(n;\{C_3,C_4\})\sim (n/2)^{3/2}?\]




A problem of Erd\H{o}s and Simonovits, who proved that\[\mathrm{ex}(n;\{C_4,C_5\})=(n/2)^{3/2}+O(n).\]K\"{o}v\'{a}ri, S\'{o}s, and Tur\'{a}n \cite{KST54} proved that the extremal number of edges for containing either $C_4$ or an odd cycle of any length is $\sim (n/2)^{3/2}$. This problem is therefore asking whether the threshold is the same if we just forbid odd cycles of length $3$. See also [574] for the general case, and [765] for $\mathrm{ex}(n;C_4)$. This problem is #48 in Extremal Graph Theory in the graphs problem collection.




References

[KST54] K\"{o}vari, T. and S\'{o}s, V. T. and Tur\'{a}n, P., On a problem of K. Zarankiewicz . Colloq. Math. (1954), 50-57.

\noindent\rule{\linewidth}{0.4pt}


%% ============================================================
\section{Problem \#575}
\label{prob:575}

\noindent\textbf{Status:} OPEN \quad|\quad \textbf{Prize:} none \quad|\quad \textbf{Updated:} 2025-08-31

\noindent\textbf{Tags:} graph theory, turan number

\medskip

\noindent\textbf{Statement:}

If $\mathcal{F}$ is a finite set of finite graphs then $\mathrm{ex}(n;\mathcal{F})$ is the maximum number of edges a graph on $n$ vertices can have without containing any subgraphs from $\mathcal{F}$. Note that it is trivial that $\mathrm{ex}(n;\mathcal{F})\leq \mathrm{ex}(n;G)$ for every $G\in\mathcal{F}$. Is it true that, for every $\mathcal{F}$, if there is a bipartite graph in $\mathcal{F}$ then there exists some bipartite $G\in\mathcal{F}$ such that\[\mathrm{ex}(n;G)\ll_{\mathcal{F}}\mathrm{ex}(n;\mathcal{F})?\]




A problem of Erd\H{o}s and Simonovits. See also [180] . This problem is #51 in Extremal Graph Theory in the graphs problem collection.

\noindent\rule{\linewidth}{0.4pt}


%% ============================================================
\section{Problem \#576}
\label{prob:576}

\noindent\textbf{Status:} OPEN \quad|\quad \textbf{Prize:} none \quad|\quad \textbf{Updated:} 2025-08-31

\noindent\textbf{Tags:} graph theory, turan number

\medskip

\noindent\textbf{Statement:}

Let $Q_k$ be the $k$-dimensional hypercube graph (so that $Q_k$ has $2^k$ vertices and $k2^{k-1}$ edges). Determine the behaviour of\[\mathrm{ex}(n;Q_k).\]




Erd\H{o}s and Simonovits \cite{ErSi70} proved that\[(\tfrac{1}{2}+o(1))n^{3/2}\leq \mathrm{ex}(n;Q_3) \ll n^{8/5}.\](In \cite{ErSi70} they mention that Erd\H{o}s had originally conjectured that $ \mathrm{ex}(n;Q_3)\gg n^{5/3}$.) Erd\H{o}s and Simonovits also proved that, if $G$ is the graph $Q_3$ with a missing edge, then $\mathrm{ex}(n;G)\asymp n^{3/2}$. In \cite{Er74c}, \cite{Er81}, and \cite{Er93} Erd\H{o}s asked whether it is $\mathrm{ex}(n;Q_3)\asymp n^{8/5}$. A theorem of Sudakov and Tomon \cite{SuTo22} implies\[\mathrm{ex}(n;Q_k)=o(n^{2-\frac{1}{k}}).\]Janzer and Sudakov \cite{JaSu22} have improved this to\[\mathrm{ex}(n;Q_k)\ll_k n^{2-\frac{1}{k-1}+\frac{1}{(k-1)2^{k-1}}}.\]See also [1035] . This problem is #52 in Extremal Graph Theory in the graphs problem collection.




References

[Er74c] Erd\H{o}s, Paul, Extremal problems on graphs and hypergraphs . (1974), 75-84.

[Er81] Erd\H{o}s, P., On the combinatorial problems which I would most like to see solved . Combinatorica (1981), 25-42.

[Er93] Erd\H{o}s, Paul, Some of my favorite solved and unsolved problems in graph theory . Quaestiones Math. (1993), 333-350.

[ErSi70] Erd\H{o}s, P. and Simonovits, M., Some extremal problems in graph theory . Combinatorial theory and its applications, I-III (Proc. Colloq., Balatonf\"{u}red, 1969) (1970), 377-390.

[JaSu22] Janzer, O. and Sudakov, B., On the Tur\'{a}n number of the hypercube . arXiv:2211.02015 (2024).

[SuTo22] Sudakov, Benny and Tomon, Istv\'{a}n, The extremal number of tight cycles . Int. Math. Res. Not. IMRN (2022), 9663-9684.

\noindent\rule{\linewidth}{0.4pt}


%% ============================================================
\section{Problem \#579}
\label{prob:579}

\noindent\textbf{Status:} OPEN \quad|\quad \textbf{Prize:} none \quad|\quad \textbf{Updated:} 2025-08-31

\noindent\textbf{Tags:} graph theory, turan number

\medskip

\noindent\textbf{Statement:}

Let $\delta&#62;0$. If $n$ is sufficiently large and $G$ is a graph on $n$ vertices with no $K_{2,2,2}$ and at least $\delta n^2$ edges then $G$ contains an independent set of size $\gg_\delta n$.




A problem of Erd\H{o}s, Hajnal, S\'{o}s, and Szemer\'{e}di, who could prove this is true for $\delta&gt;1/8$. See also [533] and the entry in the graphs problem collection .

\noindent\rule{\linewidth}{0.4pt}


%% ============================================================
\section{Problem \#584}
\label{prob:584}

\noindent\textbf{Status:} OPEN \quad|\quad \textbf{Prize:} none \quad|\quad \textbf{Updated:} 2025-08-31

\noindent\textbf{Tags:} graph theory, cycles

\medskip

\noindent\textbf{Statement:}

Let $G$ be a graph with $n$ vertices and $\delta n^{2}$ edges. Are there subgraphs $H_1,H_2\subseteq G$ such that {UL} {LI}$H_1$ has $\gg \delta^3n^2$ edges and every two edges in $H_1$ are contained in a cycle of length at most $6$, and furthermore if two edges share a vertex they are on a cycle of length $4$, and {LI}$H_2$ has $\gg \delta^2n^2$ edges and every two edges in $H_2$ are contained in a cycle of length at most $8$. {/UL}




A problem of Erd\H{o}s, Duke, and R\"{o}dl. Duke and Erd\H{o}s \cite{DuEr82} proved the first if $n$ is sufficiently large depending on $\delta$. The real challenge is to prove this when $\delta=n^{-c}$ for some $c&gt;0$. Duke, Erd\H{o}s, and R\"{o}dl \cite{DER84} proved the first statement with a $\delta^5$ in place of a $\delta^3$. Fox and Sudakov \cite{FoSu08b} have proved the second statement when $\delta &gt;n^{-1/5}$. See also the entry in the graphs problem collection .




References

[DER84] Duke, Richard and Erd\H{o}s, Paul and R\"{o}dl, Vojt\vEch, More results on subgraphs with many short cycles . Proceedings of the fifteenth Southeastern conference on combinatorics, graph theory and computing (Baton Rouge, La., 1984) (1984), 295-300.

[DuEr82] Duke, Richard and Erd\H{o}s, Paul, Subgraphs in which each pair of edges lies in a short common cycle . Proceedings of the thirteenth Southeastern conference on combinatorics, graph theory and computing (Boca Raton, Fla., 1982) (1982), 253-260.

[FoSu08b] Fox, Jacob and Sudakov, Benny, On a problem of Duke-Erd\H{o}s-R\"{o}dl on cycle-connected subgraphs . J. Combin. Theory Ser. B (2008), 1056-1062.

\noindent\rule{\linewidth}{0.4pt}


%% ============================================================
\section{Problem \#585}
\label{prob:585}

\noindent\textbf{Status:} OPEN \quad|\quad \textbf{Prize:} none \quad|\quad \textbf{Updated:} 2025-08-31

\noindent\textbf{Tags:} graph theory, cycles

\medskip

\noindent\textbf{Statement:}

What is the maximum number of edges that a graph on $n$ vertices can have if it does not contain two edge-disjoint cycles with the same vertex set?




Pyber, R\"{o}dl, and Szemer\'{e}di \cite{PRS95} constructed such a graph with $\gg n\log\log n$ edges. Chakraborti, Janzer, Methuku, and Montgomery \cite{CJMM24} have shown that such a graph can have at most $n(\log n)^{O(1)}$ many edges. Indeed, they prove that there exists a constant $C&#62;0$ such that for any $k\geq 2$ there is a $c_k$ such that if a graph has $n$ vertices and at least $c_kn(\log n)^{C}$ many edges then it contains $k$ pairwise edge-disjoint cycles with the same vertex set.




References

[CJMM24] Chakraborti, D. and Janzer, O. and Methuku, A. and Montgomery, R., Edge-disjoint cycles with the same vertex set . arXiv:2404.07190 (2024).

[PRS95] Pyber, L. and R\"{o}dl, V. and Szemer\'{e}di, E., Dense subgraphs without 3-regular subgraphs . Journal of Combinatorial Theory, Series B (1995), 41-54.

\noindent\rule{\linewidth}{0.4pt}


%% ============================================================
\section{Problem \#588}
\label{prob:588}

\noindent\textbf{Status:} OPEN \quad|\quad \textbf{Prize:} \$100 \quad|\quad \textbf{Updated:} 2025-08-31

\noindent\textbf{Tags:} geometry

\medskip

\noindent\textbf{Statement:}

Let $f_k(n)$ be minimal such that if $n$ points in $\mathbb{R}^2$ have no $k+1$ points on a line then there must be at most $f_k(n)$ many lines containing at least $k$ points. Is it true that\[f_k(n)=o(n^2)\]for $k\geq 4$?




A generalisation of [101] (which asks about $k=4$). The restriction to $k\geq 4$ is necessary since Sylvester has shown that $f_3(n)= n^2/6+O(n)$. (See also Burr, Gr\"{u}nbaum, and Sloane \cite{BGS74} and F\"{u}redi and Pal\'{a}sti \cite{FuPa84} for constructions which show that $f_3(n)\geq(1/6+o(1))n^2$.) For $k\geq 4$, K\'{a}rteszi \cite{Ka63} proved\[f_k(n)\gg_k n\log n\](resolving a conjecture of Erd\H{o}s that $f_k(n)/n\to \infty$). Gr\"{u}nbaum \cite{Gr76} proved\[f_k(n) \gg_k n^{1+\frac{1}{k-2}}.\]Erd\H{o}s speculated this may be the correct order of magnitude, but Solymosi and Stojakovi\'{c} \cite{SoSt13} give a construction which shows\[f_k(n)\gg_k n^{2-O_k(1/\sqrt{\log n})}\]




References

[BGS74] Burr, Stefan A. and Gr\"{u}nbaum, Branko and Sloane, N. J. A., The orchard problem . Geometriae Dedicata (1974), 397-424.

[FuPa84] F\"{u}redi, Z. and Pal\'{a}sti, I., Arrangements of lines with a large number of triangles . Proc. Amer. Math. Soc. (1984), 561-566.

[Gr76] Gr\"{u}nbaum, Branko, New views on some old questions of combinatorial geometry . Colloquio Internazionale sulle Teorie Combinatorie (Roma, 1973), Tomo I (1976), 451-468.

[Ka63] F. K\'{a}rteszi, Sylvester egy t\'{e}tel\'{e}r\H{o}l \'{e}s Erd\H{o}s egy sejt\'{e}s\'{e}r\H{o}l . Matematikai Lapok (1963), 3-10.

[SoSt13] Solymosi, J\'{o}zsef and Stojakovi\'C, Milo\vS, Many collinear {$k$}-tuples with no {$k+1$} collinear points . Discrete Comput. Geom. (2013), 811-820.

\noindent\rule{\linewidth}{0.4pt}


%% ============================================================
\section{Problem \#589}
\label{prob:589}

\noindent\textbf{Status:} OPEN \quad|\quad \textbf{Prize:} none \quad|\quad \textbf{Updated:} 2025-08-31

\noindent\textbf{Tags:} geometry

\medskip

\noindent\textbf{Statement:}

Let $g(n)$ be maximal such that in any set of $n$ points in $\mathbb{R}^2$ with no four points on a line there exists a subset on $g(n)$ points with no three points on a line. Estimate $g(n)$.




The trivial greedy algorithm gives $g(n)\gg n^{1/2}$. A similar question can be asked for a set with no $k$ points on a line, searching for a subset with no $l$ points on a line, for any $3\leq l&lt;k$. Erd\H{o}s thought that $g(n) \gg n$, but in fact $g(n)=o(n)$, which follows from the density Hales-Jewett theorem proved by Furstenberg and Katznelson \cite{FuKa91} (see [185] ). F\"{u}eredi \cite{Fu91b} proved\[n^{1/2}\log n\ll g(n)=o(n).\]Balogh and Solymosi \cite{BaSo18} improved the upper bound to\[g(n) \ll n^{5/6+o(1)}.\]




References

[BaSo18] Balogh, J\'{o}zsef and Solymosi, J\'{o}zsef, On the number of points in general position in the plane . Discrete Anal. (2018), Paper No. 16, 20.

[Fu91b] F\"{u}redi, Zolt\'{a}n, Maximal independent subsets in Steiner systems and in planar sets . SIAM J. Discrete Math. (1991), 196--199.

[FuKa91] Furstenberg, H. and Katznelson, Y., A density version of the Hales-Jewett Theorem . Journal d'Analyse Math\'{e}matique (1991), 64-119.

\noindent\rule{\linewidth}{0.4pt}


%% ============================================================
\section{Problem \#592}
\label{prob:592}

\noindent\textbf{Status:} OPEN \quad|\quad \textbf{Prize:} \$1000 \quad|\quad \textbf{Updated:} 2025-08-31

\noindent\textbf{Tags:} set theory, ramsey theory

\medskip

\noindent\textbf{Statement:}

Determine which countable ordinals $\beta$ have the property that, if $\alpha=\omega^{^\beta}$, then in any red/blue colouring of the edges of $K_\alpha$ there is either a red $K_\alpha$ or a blue $K_3$.




Such $\alpha$ (in other words, those such that $\alpha \to (\alpha,3)^2$) are called partition ordinals. {UL} {LI}Specker \cite{Sp57} proved this holds for $\beta=2$ and not for $3\leq \beta &lt;\omega$.{/LI} {LI} Chang \cite{Ch72} proved this holds for $\beta=\omega$.{/LI} {LI} Galvin and Larson \cite{GaLa74} have shown that if $\beta\geq 3$ has this property then $\beta$ must be 'additively indecomposable', so that in particular $\beta=\omega^\gamma$ for some countable ordinal $\gamma$. Galvin and Larson conjecture that every $\beta\geq 3$ of this form has this property.{/LI} {LI}Schipperus \cite{Sc10} have proved this is holds if $\beta=\omega^\gamma$ in which $\gamma$ is a countable ordinal which is the sum of one or two indecomposable ordinals, and this fails to hold if $\gamma$ is the sum of four or more indecomposable ordinals.{/LI} {/UL} The remaining open case appears to be when $\gamma$ is the sum of three indecomposable ordinals. The case $\beta=\omega$ is the subject of [590] , and $\beta=\omega^2$ is the subject of [591] . See also [118] , and [1169] for the case $\alpha=\omega_1^2$.




References

[Ch72] Chang, C. C., A partition theorem for the complete graph on {$\omega\sp{\omega }$} . J. Combinatorial Theory Ser. A (1972), 396-452.

[GaLa74] Galvin, Fred and Larson, Jean, Pinning countable ordinals . Fund. Math. (1974/75), 357-361.

[Sc10] Schipperus, Rene, Countable partition ordinals . Ann. Pure Appl. Logic (2010), 1195-1215.

[Sp57] Specker, Ernst, Teilmengen von Mengen mit Relationen . Comment. Math. Helv. (1957), 302-314.

\noindent\rule{\linewidth}{0.4pt}


%% ============================================================
\section{Problem \#593}
\label{prob:593}

\noindent\textbf{Status:} OPEN \quad|\quad \textbf{Prize:} \$500 \quad|\quad \textbf{Updated:} 2025-08-31

\noindent\textbf{Tags:} set theory, graph theory, hypergraphs, chromatic number

\medskip

\noindent\textbf{Statement:}

Characterize those finite 3-uniform hypergraphs which appear in every 3-uniform hypergraph of chromatic number $&#62;\aleph_0$.




Similar problems were investigated by Erd\H{o}s, Galvin, and Hajnal \cite{EGH75}. Erd\H{o}s claims that for graphs the problem is completely solved: a graph of chromatic number $\geq \aleph_1$ must contain all finite bipartite graphs but need not contain any fixed odd cycle.




References

[EGH75] Erd\H{o}s, P. and Galvin, F. and Hajnal, A., On set-systems having large chromatic number and not containing prescribed subsystems . (1975), 425--513.

\noindent\rule{\linewidth}{0.4pt}


%% ============================================================
\section{Problem \#595}
\label{prob:595}

\noindent\textbf{Status:} OPEN \quad|\quad \textbf{Prize:} \$250 \quad|\quad \textbf{Updated:} 2025-08-31

\noindent\textbf{Tags:} graph theory, set theory

\medskip

\noindent\textbf{Statement:}

Is there an infinite graph $G$ which contains no $K_4$ and is not the union of countably many triangle-free graphs?




A problem of Erd\H{o}s and Hajnal. Folkman \cite{Fo70} and Ne\v{s}et\v{r}il and R\"{o}dl \cite{NeRo75} have proved that for every $n\geq 1$ there is a graph $G$ which contains no $K_4$ and is not the union of $n$ triangle-free graphs. See also [582] and [596] .




References

[Fo70] Folkman, Jon, Graphs with monochromatic complete subgraphs in every edge coloring . SIAM J. Appl. Math. (1970), 19-24.

[NeRo75] Ne\u set\u ril, Jaroslav and R\"odl, Vojt\v ech, Type theory of partition properties of graphs . (1975), 405-412.

\noindent\rule{\linewidth}{0.4pt}


%% ============================================================
\section{Problem \#596}
\label{prob:596}

\noindent\textbf{Status:} OPEN \quad|\quad \textbf{Prize:} none \quad|\quad \textbf{Updated:} 2025-08-31

\noindent\textbf{Tags:} graph theory, ramsey theory, set theory

\medskip

\noindent\textbf{Statement:}

For which graphs $G_1,G_2$ is it true that {UL} {LI} for every $n\geq 1$ there is a graph $H$ without a $G_1$ but if the edges of $H$ are $n$-coloured then there is a monochromatic copy of $G_2$, and yet{/LI} {LI} for every graph $H$ without a $G_1$ there is an $\aleph_0$-colouring of the edges of $H$ without a monochromatic $G_2$. {/UL}




Erd\H{o}s and Hajnal originally conjectured that there are no such $G_1,G_2$, but in fact $G_1=C_4$ and $G_2=C_6$ is an example. Indeed, for this pair Ne\v{s}et\v{r}il and R\"{o}dl established the first property and Erd\H{o}s and Hajnal the second (in fact every $C_4$-free graph is a countable union of trees). Whether this is true for $G_1=K_4$ and $G_2=K_3$ is the content of [595] .

\noindent\rule{\linewidth}{0.4pt}


%% ============================================================
\section{Problem \#597}
\label{prob:597}

\noindent\textbf{Status:} OPEN \quad|\quad \textbf{Prize:} none \quad|\quad \textbf{Updated:} 2025-08-31

\noindent\textbf{Tags:} graph theory, ramsey theory, set theory

\medskip

\noindent\textbf{Statement:}

Let $G$ be a graph on at most $\aleph_1$ vertices which contains no $K_4$ and no $K_{\aleph_0,\aleph_0}$ (the complete bipartite graph with $\aleph_0$ vertices in each class). Is it true that\[\omega_1^2 \to (\omega_1\omega, G)^2?\]What about finite $G$?




Erd\H{o}s and Hajnal proved that $\omega_1^2 \to (\omega_1\omega,3)^2$. Erd\H{o}s originally asked this with just the assumption that $G$ is $K_4$-free, but Baumgartner proved that $\omega_1^2 \not\to (\omega_1\omega, K_{\aleph_0,\aleph_0})^2$.

\noindent\rule{\linewidth}{0.4pt}


%% ============================================================
\section{Problem \#598}
\label{prob:598}

\noindent\textbf{Status:} OPEN \quad|\quad \textbf{Prize:} none \quad|\quad \textbf{Updated:} 2025-08-31

\noindent\textbf{Tags:} set theory, ramsey theory

\medskip

\noindent\textbf{Statement:}

Let $m$ be an infinite cardinal and $\kappa$ be the successor cardinal of $2^{\aleph_0}$. Can one colour the countable subsets of $m$ using $\kappa$ many colours so that every $X\subseteq m$ with $\lvert X\rvert=\kappa$ contains subsets of all possible colours?

\noindent\rule{\linewidth}{0.4pt}


%% ============================================================
\section{Problem \#600}
\label{prob:600}

\noindent\textbf{Status:} OPEN \quad|\quad \textbf{Prize:} none \quad|\quad \textbf{Updated:} 2025-08-31

\noindent\textbf{Tags:} graph theory

\medskip

\noindent\textbf{Statement:}

Let $e(n,r)$ be minimal such that every graph on $n$ vertices with at least $e(n,r)$ edges, each edge contained in at least one triangle, must have an edge contained in at least $r$ triangles. Let $r\geq 2$. Is it true that\[e(n,r+1)-e(n,r)\to \infty\]as $n\to \infty$? Is it true that\[\frac{e(n,r+1)}{e(n,r)}\to 1\]as $n\to \infty$?




Ruzsa and Szemer\'{e}di \cite{RuSz78} proved that $e(n,r)=o(n^2)$ for any fixed $r$. See also [80] .




References

[RuSz78] Ruzsa, I. Z. and Szemer\'{e}di, E., Triple systems with no six points carrying three triangles . Combinatorics (Proc. Fifth Hungarian Colloq., Keszthely, 1976), Vol. II (1978), 939-945.

\noindent\rule{\linewidth}{0.4pt}


%% ============================================================
\section{Problem \#601}
\label{prob:601}

\noindent\textbf{Status:} OPEN \quad|\quad \textbf{Prize:} \$500 \quad|\quad \textbf{Updated:} 2025-08-31

\noindent\textbf{Tags:} graph theory, set theory

\medskip

\noindent\textbf{Statement:}

For which limit ordinals $\alpha$ is it true that if $G$ is a graph with vertex set $\alpha$ then $G$ must have either an infinite path or independent set on a set of vertices with order type $\alpha$?




A problem of Erd\H{o}s, Hajnal, and Milner \cite{EHM70}, who proved this is true for $\alpha &lt; \omega_1^{\omega+2}$. In \cite{Er82e} Erd\H{o}s offers \$250 for showing what happens when $\alpha=\omega_1^{\omega+2}$ and \$500 for settling the general case. Larson \cite{La90} proved this is true for all $\alpha&lt;2^{\aleph_0}$ assuming Martin's axiom .




References

[EHM70] Erd\H{o}s, P. and Hajnal, A. and Milner, E. C., Set mappings and polarized partition relations . Combinatorial theory and its applications, I-III (Proc. Colloq., Balatonf\"{u}red, 1969) (1970), 327-363.

[Er82e] Erd\H{o}s, Paul, Some of my favourite problems which recently have been solved . (1982), 59--79.

[La90] Larson, Jean A., Martin's axiom and ordinal graphs: large independent sets or infinite paths . Ann. Pure Appl. Logic (1990), 31-39.

\noindent\rule{\linewidth}{0.4pt}


%% ============================================================
\section{Problem \#602}
\label{prob:602}

\noindent\textbf{Status:} OPEN \quad|\quad \textbf{Prize:} none \quad|\quad \textbf{Updated:} 2025-08-31

\noindent\textbf{Tags:} combinatorics, set theory

\medskip
\noindent\textit{Property B}


\noindent\textbf{Statement:}

Let $(A_i)$ be a family of sets with $\lvert A_i\rvert=\aleph_0$ for all $i$, such that for any $i\neq j$ we have $\lvert A_i\cap A_j\rvert$ finite and $\neq 1$. Is there a $2$-colouring of $\cup A_i$ such that no $A_i$ is monochromatic?




A problem of Komj\'{a}th. The existence of such a $2$-colouring is sometimes known as Property B.

\noindent\rule{\linewidth}{0.4pt}


%% ============================================================
\section{Problem \#603}
\label{prob:603}

\noindent\textbf{Status:} OPEN \quad|\quad \textbf{Prize:} none \quad|\quad \textbf{Updated:} 2025-08-31

\noindent\textbf{Tags:} combinatorics, set theory

\medskip

\noindent\textbf{Statement:}

Let $(A_i)$ be a family of countably infinite sets such that $\lvert A_i\cap A_j\rvert \neq 2$ for all $i\neq j$. Find the smallest cardinal $C$ such that $\cup A_i$ can always be coloured with at most $C$ colours so that no $A_i$ is monochromatic.




A problem of Komj\'{a}th. If instead we have $\lvert A_i\cap A_j\rvert \neq 1$ then Komj\'{a}th showed that this is possible with at most $\aleph_0$ colours.

\noindent\rule{\linewidth}{0.4pt}


%% ============================================================
\section{Problem \#604}
\label{prob:604}

\noindent\textbf{Status:} OPEN \quad|\quad \textbf{Prize:} \$500 \quad|\quad \textbf{Updated:} 2025-08-31

\noindent\textbf{Tags:} geometry, distances

\medskip
\noindent\textit{pinned distance problem}


\noindent\textbf{Statement:}

Given $n$ distinct points $A\subset\mathbb{R}^2$ must there be a point $x\in A$ such that\[\#\{ d(x,y) : y \in A\} \gg n^{1-o(1)}?\]Or even $\gg n/\sqrt{\log n}$?




The pinned distance problem, a stronger form of [89] . The example of an integer grid show that $n/\sqrt{\log n}$ would be best possible. It may be true that there are $\gg n$ many such points (Hunter has noted in the comments that this trivially follows from the existence of a single point), or that this is true on average - for example, if $d(x)$ counts the number of distinct distances from $x$ then in \cite{Er75f} Erd\H{o}s conjectured\[\sum_{x\in A}d(x) \gg \frac{n^2}{\sqrt{\log n}},\]where $A\subset \mathbb{R}^2$ is any set of $n$ points. In \cite{Er97e} Erd\H{o}s offers \$500 for a solution to this problem, but it is unclear whether he intended this for proving the existence of a single such point or for $\gg n$ many such points. In \cite{Er97e} Erd\H{o}s wrote that he initially 'overconjectured' and thought that the answer to this problem is the same as for the number of distinct distances between all pairs (see [89] ), but this was disproved by Harborth. It could be true that the answers are the same up to an additive factor of $n^{o(1)}$. The best known bound is\[\gg n^{c-o(1)},\]due to Katz and Tardos \cite{KaTa04}, where\[c=\frac{48-14e}{55-16e}=0.864137\cdots.\]




References

[Er75f] Erd\H{o}s, Paul, On some problems of elementary and combinatorial geometry . Ann. Mat. Pura Appl. (4) (1975), 99-108.

[Er97e] Erd\H{o}s, Paul, Some of my favourite unsolved problems . Math. Japon. (1997), 527-537.

[KaTa04] Katz, Nets Hawk and Tardos, G\'{a}bor, A new entropy inequality for the Erd\H{o}s distance problem . Towards a theory of geometric graphs (2004), 119-126.

\noindent\rule{\linewidth}{0.4pt}


%% ============================================================
\section{Problem \#609}
\label{prob:609}

\noindent\textbf{Status:} OPEN \quad|\quad \textbf{Prize:} none \quad|\quad \textbf{Updated:} 2025-08-31

\noindent\textbf{Tags:} graph theory, ramsey theory

\medskip

\noindent\textbf{Statement:}

Let $f(n)$ be the minimal $m$ such that if the edges of $K_{2^n+1}$ are coloured with $n$ colours then there must be a monochromatic odd cycle of length at most $m$. Estimate $f(n)$.




A problem of Erd\H{o}s and Graham. The edges of $K_{2^n}$ can be $n$-coloured to avoid odd cycles of any length. It can be shown that $C_5$ and $C_7$ can be avoided for large $n$. Chung \cite{Ch97} asked whether $f(n)\to \infty$ as $n\to \infty$. Day and Johnson \cite{DaJo17} proved this is true, and that\[f(n)\geq 2^{c\sqrt{\log n}}\]for some constant $c&gt;0$. The trivial upper bound is $2^n$. Gir\~{a}o and Hunter \cite{GiHu24} have proved that\[f(n) \ll \frac{2^n}{n^{1-o(1)}}.\]Janzer and Yip \cite{JaYi25} have improved this to\[f(n) \ll n^{3/2}2^{n/2}.\]See also the entry in the graphs problem collection .




References

[Ch97] Chung, F. R. K., Open problems of {P}aul Erd\H{o}s in graph theory . J. Graph Theory (1997), 3--36.

[DaJo17] Day, A. Nicholas and Johnson, J. Robert, Multicolour {R}amsey numbers of odd cycles . J. Combin. Theory Ser. B (2017), 56--63.

[GiHu24] A. Gir\~Ao and Z. Hunter, Monochromatic odd cycles in edge-coloured complete graphs . arXiv:2412.07708 (2024).

[JaYi25] O. Janzer and F. Yip, Short monochromatic odd cycles . arXiv:2506.14910 (2025).

\noindent\rule{\linewidth}{0.4pt}


%% ============================================================
\section{Problem \#611}
\label{prob:611}

\noindent\textbf{Status:} OPEN \quad|\quad \textbf{Prize:} none \quad|\quad \textbf{Updated:} 2025-08-31

\noindent\textbf{Tags:} graph theory

\medskip

\noindent\textbf{Statement:}

For a graph $G$ let $\tau(G)$ denote the minimal number of vertices that include at least one from each maximal clique of $G$ (sometimes called the clique transversal number). Is it true that if all maximal cliques in $G$ have at least $cn$ vertices then $\tau(G)=o_c(n)$? Similarly, estimate for $c&#62;0$ the minimal $k_c(n)$ such that if every maximal clique in $G$ has at least $k_c(n)$ vertices then $\tau(G)&#60;(1-c)n$.




A problem of Erd\H{o}s, Gallai, and Tuza \cite{EGT92}, who proved for the latter question that $k_c(n) \geq n^{c'/\log\log n}$ for some $c'&gt;0$, and that if every clique has size least $k$ then $\tau(G) \leq n-(kn)^{1/2}$. Bollob\'{a}s and Erd\H{o}s proved that if every maximal clique has at least $n+3-2\sqrt{n}$ vertices then $\tau(G)=1$ (and this threshold is best possible). See also [610] and the entry in the graphs problem collection .




References

[EGT92] Erd\H{o}s, Paul and Gallai, Tibor and Tuza, Zsolt, Covering the cliques of a graph with vertices . Discrete Math. (1992), 279-289.

\noindent\rule{\linewidth}{0.4pt}


%% ============================================================
\section{Problem \#612}
\label{prob:612}

\noindent\textbf{Status:} OPEN \quad|\quad \textbf{Prize:} none \quad|\quad \textbf{Updated:} 2025-08-31

\noindent\textbf{Tags:} graph theory

\medskip

\noindent\textbf{Statement:}

Let $G$ be a connected graph with $n$ vertices, minimum degree $d$, and diameter $D$. Show if that $G$ contains no $K_{2r}$ and $(r-1)(3r+2)\mid d$ then\[D\leq \frac{2(r-1)(3r+2)}{2r^2-1}\frac{n}{d}+O(1),\]and if $G$ contains no $K_{2r+1}$ and $3r-1 \mid d$ then\[D\leq \frac{3r-1}{r}\frac{n}{d}+O(1).\]




A problem of Erd\H{o}s, Pach, Pollack, and Tuza \cite{EPPT89}, who gave constructions showing that the above bounds would be sharp, and proved the case $2r+1=3$. It is known (see \cite{EPPT89} for example) that any connected graph on $n$ vertices with minimum degree $d$ has diameter\[D\leq 3\frac{n}{d+1}+O(1).\]This was disproven for the case of $K_{2r}$-free graphs with $r\geq 2$ by Czabarka, Singgih, and Sz\'{e}kely \cite{CSS21}, who constructed arbitrarily large connected graphs on $n$ vertices which contain no $K_{2r}$ and have minimum degree $d$, and diameter\[\frac{6r-5}{(2r-1)d+2r-3}n+O(1),\]which contradicts the above conjecture for each fixed $r$ as $d\to \infty$. They suggest the amended conjecture, which no longer divides into two cases, that if $G$ is a connected graph on $n$ vertices with minimum degree $d$ which contains no $K_{k+1}$ then the diameter of $G$ is at most\[(3-\tfrac{2}{k})\frac{n}{d}+O(1).\]This bound is known under the weaker assumption that $G$ is $k$-colourable when $k=3$ and $k=4$, shown by Czabarka, Dankelmann, and Sz\'{e}kely \cite{CDS09} and Czabarka, Smith, and Sz\'{e}kely \cite{CSS23}. Cambie and Jooken \cite{CaJo25} have given an example that shows the diameter for $K_4$-free graphs with minimum degree $16$ is at least $\frac{31}{216}n+O(1)$, giving another counterexample to the original conjecture. See also the entry in the graphs problem collection .




References

[CDS09] Czabarka, \'{e}. and Dankelmann, P. and Sz\'{e}kely, L. A., Diameter of 4-colourable graphs . European J. Combin. (2009), 1082--1089.

[CSS21] Czabarka, \'{e}va and Singgih, Inne and Sz\'{e}kely, L\'{a}szl\'{o}{} A., Counterexamples to a conjecture of {E}rd\H{o}s, {P}ach, {P}ollack and Tuza . J. Combin. Theory Ser. B (2021), 38--45.

[CSS23] Czabarka, \'{e}va and Smith, Stephen J. and Sz\'{e}kely, L\'{a}szl\'{o}, Maximum diameter of 3- and 4-colorable graphs . J. Graph Theory (2023), 262--270.

[CaJo25] S. Cambie and J. Jooken, Sharp results for the Erd\H{o}s, Pach, Pollack and Tuza problem . arXiv:2502.08626 (2025).

[EPPT89] Erd\H{o}s, Paul and Pach, J\'{a}nos and Pollack, Richard and Tuza, Zsolt, Radius, diameter, and minimum degree . J. Combin. Theory Ser. B (1989), 73-79.

\noindent\rule{\linewidth}{0.4pt}


%% ============================================================
\section{Problem \#614}
\label{prob:614}

\noindent\textbf{Status:} OPEN \quad|\quad \textbf{Prize:} none \quad|\quad \textbf{Updated:} 2025-08-31

\noindent\textbf{Tags:} graph theory

\medskip

\noindent\textbf{Statement:}

Let $f(n,k)$ be minimal such that there is a graph with $n$ vertices and $f(n,k)$ edges where every set of $k+2$ vertices induces a subgraph with maximum degree at least $k$. Determine $f(n,k)$.




See also the entry in the graphs problem collection .

\noindent\rule{\linewidth}{0.4pt}


%% ============================================================
\section{Problem \#616}
\label{prob:616}

\noindent\textbf{Status:} OPEN \quad|\quad \textbf{Prize:} none \quad|\quad \textbf{Updated:} 2025-08-31

\noindent\textbf{Tags:} graph theory

\medskip

\noindent\textbf{Statement:}

Let $r\geq 3$. For an $r$-uniform hypergraph $G$ let $\tau(G)$ denote the covering number (or transversal number), the minimum size of a set of vertices which includes at least one from each edge in $G$. Determine the best possible $t$ such that, if $G$ is an $r$-uniform hypergraph $G$ where every subgraph $G'$ on at most $3r-3$ vertices has $\tau(G')\leq 1$, we have $\tau(G)\leq t$.




Erd\H{o}s, Hajnal, and Tuza \cite{EHT91} proved that this $t$ satisfies\[\frac{3}{16}r+\frac{7}{8}\leq t \leq \frac{1}{5}r.\]




References

[EHT91] Erd\H{o}s, Paul and Hajnal, Andr\'{a}s and Tuza, Zsolt, Local constraints ensuring small representing sets . J. Combin. Theory Ser. A (1991), 78-84.

\noindent\rule{\linewidth}{0.4pt}


%% ============================================================
\section{Problem \#619}
\label{prob:619}

\noindent\textbf{Status:} OPEN \quad|\quad \textbf{Prize:} none \quad|\quad \textbf{Updated:} 2025-08-31

\noindent\textbf{Tags:} graph theory

\medskip

\noindent\textbf{Statement:}

For a triangle-free graph $G$ let $h_r(G)$ be the smallest number of edges that need to be added to $G$ so that it has diameter $r$ (while preserving the property of being triangle-free). Is it true that there exists a constant $c&#62;0$ such that if $G$ is a connected graph on $n$ vertices then $h_4(G)&#60;(1-c)n$?




A problem of Erd\H{o}s, Gy\'{a}rf\'{a}s, and Ruszink\'{o} \cite{EGR98} who proved that $h_3(G)\leq n$ and $h_5(G) \leq \frac{n-1}{2}$ and there exist connected graphs $G$ on $n$ vertices with $h_3(G)\geq n-c$ for some constant $c&gt;0$. If we omit the condition that the graph must remain triangle-free then Alon, Gy\'{a}rf\'{a}s, and Ruszink\'{o} \cite{AGR00} have proved that adding $n/2$ edges always suffices to obtain diameter at most $4$. See also [134] and [618] .




References

[AGR00] Alon, Noga and Gy\'{a}rf\'{a}s, Andr\'{a}s and Ruszink\'{o}, Mikl\'{o}s, Decreasing the diameter of bounded degree graphs . J. Graph Theory (2000), 161--172.

[EGR98] Erd\H{o}s, Paul and Gy\'{a}rf\'{a}s, Andr\'{a}s and Ruszink\'{o}, Mikl\'{o}s, How to decrease the diameter of triangle-free graphs . Combinatorica (1998), 493-501.

\noindent\rule{\linewidth}{0.4pt}


%% ============================================================
\section{Problem \#620}
\label{prob:620}

\noindent\textbf{Status:} OPEN \quad|\quad \textbf{Prize:} none \quad|\quad \textbf{Updated:} 2025-08-31

\noindent\textbf{Tags:} graph theory

\medskip
\noindent\textit{Erdős-Rogers problem}


\noindent\textbf{Statement:}

If $G$ is a graph on $n$ vertices without a $K_4$ then how large a triangle-free induced subgraph must $G$ contain?




This was first asked by Erd\H{o}s and Rogers \cite{ErRo62}, and is generally known as the Erd\H{o}s-Rogers problem. Let $f(n)$ be such that every such graph contains a triangle-free subgraph with at least $f(n)$ vertices. It is now known that $f(n)=n^{1/2+o(1)}$. Bollob\'{a}s and Hind \cite{BoHi91} proved\[n^{1/2} \ll f(n) \ll n^{7/10+o(1)}.\]Krivelevich \cite{Kr94} improved this to\[n^{1/2}(\log\log n)^{1/2} \ll f(n) \ll n^{2/3}(\log n)^{1/3}.\]Wolfovitz \cite{Wo13} proved\[f(n) \ll n^{1/2}(\log n)^{120}.\]The best bounds currently known are\[n^{1/2}\frac{(\log n)^{1/2}}{\log\log n}\ll f(n) \ll n^{1/2}\log n.\]The lower bound follows from results of Shearer \cite{Sh95}, and the upper bound was proved by Mubayi and Verstraete \cite{MuVe24}.




References

[BoHi91] Bollob\'{a}s, B. and Hind, H. R., Graphs without large triangle free subgraphs . Discrete Math. (1991), 119-131.

[ErRo62] Erd\H{o}s, P. and Rogers, C. A., The construction of certain graphs . Canadian J. Math. (1962), 702-707.

[Kr94] Krivelevich, Michael, {$K^s$}-free graphs without large {$K^r$}-free subgraphs . Combin. Probab. Comput. (1994), 349-354.

[MuVe24] D. Mubayi and J. Verstraete, On the order of Erd\H{o}s-Rogers functions . arXiv:2401.02548 (2024).

[Sh95] Shearer, James B., On the independence number of sparse graphs . Random Structures Algorithms (1995), 269--271.

[Wo13] Wolfovitz, Guy, {$K_4$}-free graphs without large induced triangle-free subgraphs . Combinatorica (2013), 623-631.

\noindent\rule{\linewidth}{0.4pt}


%% ============================================================
\section{Problem \#623}
\label{prob:623}

\noindent\textbf{Status:} OPEN \quad|\quad \textbf{Prize:} none \quad|\quad \textbf{Updated:} 2025-08-31

\noindent\textbf{Tags:} set theory

\medskip

\noindent\textbf{Statement:}

Let $X$ be a set of cardinality $\aleph_\omega$ and $f$ be a function from the finite subsets of $X$ to $X$ such that $f(A)\not\in A$ for all $A$. Must there exist an infinite $Y\subseteq X$ that is independent - that is, for all finite $B\subset Y$ we have $f(B)\not\in Y$?




A problem of Erd\H{o}s and Hajnal \cite{ErHa58}, who proved that if $\lvert X\rvert &#60;\aleph_\omega$ then the answer is no. Erd\H{o}s suggests in \cite{Er99} that this problem is 'perhaps undecidable'.




References

[Er99] Erd\H{o}s, Paul, A selection of problems and results in combinatorics . Combin. Probab. Comput. (1999), 1-6.

[ErHa58] Erd\H{o}s, P. and Hajnal, A., On the structure of set mappings . Acta Math. Acad. Sci. Hungar. (1958), 111-133.

\noindent\rule{\linewidth}{0.4pt}


%% ============================================================
\section{Problem \#624}
\label{prob:624}

\noindent\textbf{Status:} OPEN \quad|\quad \textbf{Prize:} none \quad|\quad \textbf{Updated:} 2025-08-31

\noindent\textbf{Tags:} combinatorics

\medskip

\noindent\textbf{Statement:}

Let $X$ be a finite set of size $n$ and $H(n)$ be such that there is a function $f:\{A : A\subseteq X\}\to X$ so that for every $Y\subseteq X$ with $\lvert Y\rvert \geq H(n)$ we have\[\{ f(A) : A\subseteq Y\}=X.\]Prove that\[H(n)-\log_2 n \to \infty.\]




A problem of Erd\H{o}s and Hajnal \cite{ErHa68} who proved that\[\log_2 n \leq H(n) &lt; \log_2n +(3+o(1))\log_2\log_2n.\]Erd\H{o}s said that even the weaker statement that for $n=2^k$ we have $H(n)\geq k+1$ is open, but Alon has provided the following simple proof: by the pigeonhole principle there are $\frac{n-1}{2}$ subsets $A_i$ of size $2$ such that $f(A_i)$ is the same. Any set $Y$ of size $k$ containing at least $k/2$ of them can have at most\[2^k-\lfloor k/2\rfloor+1&lt; 2^k=n\]distinct elements in the union of the images of $f(A)$ for $A\subseteq Y$. For this weaker statement, Erd\H{o}s and Gy\'{a}rf\'{a}s conjectured the stronger form that if $\lvert X\rvert=2^k$ then, for any $f:\{A : A\subseteq X\}\to X$, there must exist some $Y\subset X$ of size $k$ such that\[\#\{ f(A) : A\subseteq Y\}&lt; 2^k-k^C\]for every $C$ (with $k$ sufficiently large depending on $C$). This was proved by Alon (personal communication), who proved the stronger version that there exists some absolute constant $c&gt;0$ such that, if $k$ is large enough, there must exist some $Y\subset X$ of size $k$ such that\[\#\{ f(A) : A\subseteq Y\}&lt;(1-c)2^k.\]Alon also proved that, provided $k$ is large enough, if $\lvert X\rvert=2^k$ there exists some $f:\{A: A\subseteq X\}\to X$ such that, if $Y\subset X$ with $\lvert Y\rvert=k$, then\[\#\{ f(A) : A\subseteq Y\}&gt;\tfrac{1}{4}2^k.\]




References

[ErHa68] Erd\H{o}s, P. and Hajnal, A., On a combinatorial problem . Mat. Lapok (1968), 345-348.

\noindent\rule{\linewidth}{0.4pt}


%% ============================================================
\section{Problem \#625}
\label{prob:625}

\noindent\textbf{Status:} OPEN \quad|\quad \textbf{Prize:} \$1000 \quad|\quad \textbf{Updated:} 2025-08-31

\noindent\textbf{Tags:} graph theory, chromatic number

\medskip

\noindent\textbf{Statement:}

The cochromatic number of $G$, denoted by $\zeta(G)$, is the minimum number of colours needed to colour the vertices of $G$ such that each colour class induces either a complete graph or empty graph. Let $\chi(G)$ denote the chromatic number. If $G$ is a random graph with $n$ vertices and each edge included independently with probability $1/2$ then is it true that almost surely\[\chi(G) - \zeta(G) \to \infty\]as $n\to \infty$?




A problem of Erd\H{o}s and Gimbel (see also \cite{Gi16}). At a conference on random graphs in Poznan, Poland (most likely in 1989) Erd\H{o}s offered \$100 for a proof that this is true, and \$1000 for a proof that this is false (although later told Gimbel that \$1000 was perhaps too much). It is known that almost surely\[\frac{n}{2\log_2n}\leq \zeta(G)\leq \chi(G)\leq (1+o(1))\frac{n}{2\log_2n}.\](The final upper bound is due to Bollob\'{a}s \cite{Bo88}. The first inequality follows from the fact that almost surely $G$ has clique number and independence number $&lt; 2\log_2n$.) Heckel \cite{He24} and, independently, Steiner \cite{St24b} have shown that it is not the case that $\chi(G)-\zeta(G)$ is bounded with high probability, and in fact if $\chi(G)-\zeta(G) \leq f(n)$ with high probability then $f(n)\geq n^{1/2-o(1)}$ along an infinite sequence of $n$. Heckel conjectures that, with high probability,\[\chi(G)-\zeta(G) \asymp \frac{n}{(\log n)^3}.\]Heckel \cite{He24c} further proved that, for any $\epsilon&gt;0$, we have\[\chi(G) -\zeta(G) \geq n^{1-\epsilon}\]for roughly $95\%$ of all $n$.




References

[Bo88] Bollob\'{a}s, B., The chromatic number of random graphs . Combinatorica (1988), 49-55.

[Gi16] J. Gimbel, Some of my favorite coloring problems for graphs and digraphs . Graph Theory: Favorite conjectures and open problems (2016), 95-108.

[He24] A. Heckel, On a question of Erd\H{o}s and Gimbel on the cochromatic number . arXiv:2408.13839 (2024).

[He24c] A. Heckel, The difference between the chromatic and the cochromatic number of a random graph . arXiv:2409.17614 (2024).

[St24b] R. Steiner, On the difference between the chromatic and cochromatic number . arXiv:2408.02400 (2024).

\noindent\rule{\linewidth}{0.4pt}


%% ============================================================
\section{Problem \#626}
\label{prob:626}

\noindent\textbf{Status:} OPEN \quad|\quad \textbf{Prize:} none \quad|\quad \textbf{Updated:} 2025-08-31

\noindent\textbf{Tags:} graph theory, chromatic number, cycles

\medskip

\noindent\textbf{Statement:}

Let $k\geq 4$ and $g_k(n)$ denote the largest $m$ such that there is a graph on $n$ vertices with chromatic number $k$ and girth $&#62;m$ (i.e. contains no cycle of length $\leq m$). Does\[\lim_{n\to \infty}\frac{g_k(n)}{\log n}\]exist? Conversely, if $h^{(m)}(n)$ is the maximal chromatic number of a graph on $n$ vertices with girth $&#62;m$ then does\[\lim_{n\to \infty}\frac{\log h^{(m)}(n)}{\log n}\]exist, and what is its value?




It is known that\[\frac{1}{4\log k}\log n\leq g_k(n) \leq \frac{2}{\log(k-2)}\log n+1,\]the lower bound due to Kostochka \cite{Ko88} and the upper bound to Erd\H{o}s \cite{Er59b}. Erd\H{o}s \cite{Er59b} proved that\[\lim_{n\to \infty}\frac{\log h^{(m)}(n)}{\log n}\gg \frac{1}{m}\]and, for odd $m$,\[\lim_{n\to \infty}\frac{\log h^{(m)}(n)}{\log n}\leq \frac{2}{m+1},\]and conjectured this is sharp. He had no good guess for the value of the limit for even $m$, other that it should lie in $[\frac{2}{m+2},\frac{2}{m}]$, but could not prove this even for $m=4$. See also the entry in the graphs problem collection .




References

[Er59b] Erd\H{o}s, P., Graph theory and probability . Canadian J. Math. (1959), 34-38.

[Ko88] Kostochka, A. V., Upper bounds on the chromatic number of graphs . Trudy Inst. Mat. (Novosibirsk) (1988), 204-226, 265.

\noindent\rule{\linewidth}{0.4pt}


%% ============================================================
\section{Problem \#627}
\label{prob:627}

\noindent\textbf{Status:} OPEN \quad|\quad \textbf{Prize:} none \quad|\quad \textbf{Updated:} 2025-08-31

\noindent\textbf{Tags:} graph theory, chromatic number

\medskip

\noindent\textbf{Statement:}

Let $\omega(G)$ denote the clique number of $G$ and $\chi(G)$ the chromatic number. If $f(n)$ is the maximum value of $\chi(G)/\omega(G)$, as $G$ ranges over all graphs on $n$ vertices, then does\[\lim_{n\to\infty}\frac{f(n)}{n/(\log_2n)^2}\]exist?




Tutte and Zykov \cite{Zy52} independently proved that for every $k$ there is a graph with $\omega(G)=2$ and $\chi(G)=k$. Erd\H{o}s \cite{Er61d} proved that for every $n$ there is a graph on $n$ vertices with $\omega(G)=2$ and $\chi(G)\gg n^{1/2}/\log n$, whence $f(n) \gg n^{1/2}/\log n$. Erd\H{o}s \cite{Er67c} proved that\[f(n) \asymp \frac{n}{(\log_2n)^2}\]and that the limit in question, if it exists, must be in $[1/4,4]$. (The upper bound he states as $1$, but \cite{AFM25} note the method actually gives $4$.) Araujo, Filipe, and Miyazaki \cite{AFM25} prove that if $\lim\frac{\log R(k)}{k}$ exists and is equal to $C$ (see [77] ), and also $R(s,t)\leq R(k)$ for all $st\leq k^2$, then the limit in this question also exists and is equal to $C^2$. Using this connection they improve the upper bound from $4$ to $\approx 3.7$. See also the entry in the graphs problem collection .




References

[AFM25] I. Araujo, R. Filipe, and R. Miyazaki, A note on the maximum ratio between chromatic number and clique number . arXiv:252.16062 (2025).

[Er61d] Erd\H{o}s, P., Graph theory and probability. II . Canadian J. Math. (1961), 346-352.

[Er67c] Erd\H{o}s, P., Some remarks on chromatic graphs . Colloq. Math. (1967), 253-256.

[Zy52] Zykov, A. A., On some properties of linear complexes . Amer. Math. Soc. Translation (1952), 33.

\noindent\rule{\linewidth}{0.4pt}


%% ============================================================
\section{Problem \#629}
\label{prob:629}

\noindent\textbf{Status:} OPEN \quad|\quad \textbf{Prize:} none \quad|\quad \textbf{Updated:} 2025-08-31

\noindent\textbf{Tags:} graph theory, chromatic number

\medskip

\noindent\textbf{Statement:}

The list chromatic number $\chi_L(G)$ is defined to be the minimal $k$ such that for any assignment of a list of $k$ colours to each vertex of $G$ (perhaps different lists for different vertices) a colouring of each vertex by a colour on its list can be chosen such that adjacent vertices receive distinct colours. Determine the minimal number of vertices $n(k)$ of a bipartite graph $G$ such that $\chi_L(G)&#62;k$.




A problem of Erd\H{o}s, Rubin, and Taylor \cite{ERT80}, who proved that\[2^{k-1}&lt;n(k) &lt;k^22^{k+2}.\]They also prove that if $m(k)$ is the size of the smallest family of $k$-sets without property B (i.e. the smallest number of $k$-sets in a graph with chromatic number $3$) then $m(k)\leq n(k)\leq m(k+1)$. Bounds on $m(k)$ are the subject of [901] . The lower bounds on $m(k)$ by Radhakrishnan and Srinivasan \cite{RaSr00} imply that\[2^k \left(\frac{k}{\log k}\right)^{1/2}\ll n(k).\]Erd\H{o}s, Rubin, and Taylor \cite{ERT80} proved $n(2)=6$ and Hanson, MacGillivray, and Toft \cite{HMT96} proved $n(3)=14$ and\[n(k) \leq kn(k-2)+2^k.\]See also the entry in the graphs problem collection .




References

[ERT80] Erd\H{o}s, Paul and Rubin, Arthur L. and Taylor, Herbert, Choosability in graphs . (1980), 125-157.

[HMT96] Hanson, Denis and MacGillivray, Gary and Toft, Bjarne, Choosability of bipartite graphs . Ars Combin. (1996), 183-192.

[RaSr00] Radhakrishnan, Jaikumar and Srinivasan, Aravind, Improved bounds and algorithms for hypergraph {$2$}-coloring . Random Structures Algorithms (2000), 4--32.

\noindent\rule{\linewidth}{0.4pt}


%% ============================================================
\section{Problem \#634}
\label{prob:634}

\noindent\textbf{Status:} OPEN \quad|\quad \textbf{Prize:} \$25 \quad|\quad \textbf{Updated:} 2025-08-31

\noindent\textbf{Tags:} geometry

\medskip

\noindent\textbf{Statement:}

Find all $n$ such that there is at least one triangle which can be cut into $n$ congruent triangles.




Erd\H{o}s' question was reported by Soifer \cite{So09c}. It is easy to see that all square numbers have this property (in fact for square numbers any triangle will do). Soifer \cite{So09c} has shown that numbers of the form $2n^2,3n^2,6n^2,n^2+m^2$ also have this property. Beeson has shown (see the slides below) that $7$ and $11$ do not have this property. It is possible that any prime of the form $4n+3$ does not have this property. In particular, it is not known if $19$ has this property (i.e. are there $19$ congruent triangles which can be assembled into a triangle?). For more on this problem see these slides from a talk by Michael Beeson. As a demonstration of this problem we include {IMAGE=634Triangle,a picture} of a cutting of an equilateral triangle into $27$ congruent triangles from these slides. Soifer proved \cite{So09} that if we relax congruence to similarity then every triangle can be cut into $N$ similar triangles when $N\neq 2,3,5$. If one requires the smaller triangles to be similar to the larger triangle then the only possible values of $N$ are $n^2,n^2+m^2,3n^2$, proved by Snover, Waiveris, and Williams \cite{SWW91}. Zhang \cite{Zh25}, among other results, has proved that for any integers $a \geq b$, if\[n\geq 3\left\lceil \frac{a^2+b^2+ab-a-b}{ab}\right\rceil\]then $n^2ab$ has this property (indeed, they explicitly show that an equilateral triangle can be tiled with $n^2ab$ many triangles of side lengths $a,b,\sqrt{a^2+b^2+2+ab}$). See also [633] .




References

[SWW91] Snover, S. and Waiveris, C. and Williams, J., Rep-tiling for triangles . Discrete Math. (1991), 193-200.

[So09] Soifer, Alexander, How Does One Cut a Triangle? I . (2009), 15-23.

[So09c] Soifer, Alexander, Is there anything beyond the solution? . (2009), 47-50.

[Zh25] Y. Zhang, Tiling Triangles With $2\pi/3$ Angles . arXiv:2512.22696 (2025).

\noindent\rule{\linewidth}{0.4pt}


%% ============================================================
\section{Problem \#635}
\label{prob:635}

\noindent\textbf{Status:} OPEN \quad|\quad \textbf{Prize:} none \quad|\quad \textbf{Updated:} 2026-01-30

\noindent\textbf{Tags:} number theory

\medskip

\noindent\textbf{Statement:}

Let $t\geq 1$ and $A\subseteq \{1,\ldots,N\}$ be such that whenever $a,b\in A$ with $b-a\geq t$ we have $b-a\nmid b$. How large can $\lvert A\rvert$ be? Is it true that\[\lvert A\rvert \leq \left(\frac{1}{2}+o_t(1)\right)N?\]




Asked by Erd\H{o}s in a letter to Ruzsa in around 1980. Erd\H{o}s observes that when $t=1$ the maximum possible is\[\lvert A\rvert=\left\lfloor\frac{N+1}{2}\right\rfloor,\]achieved by taking $A$ to be all odd numbers in $\{1,\ldots,N\}$. He also observes that when $t=2$ there exists such an $A$ with\[\lvert A\rvert \geq \frac{N}{2}+c\log N\]for some constant $c&#62;0$: take $A$ to be the union of all odd numbers together with numbers of the shape $2^k$ with $k$ odd. The second question has been resolved in the affirmative by ChatGPT-5.2 (prompted by Leeham); Tao has since observed that a positive answer also follows fairly quickly from an inequality due to Elliott \cite{El79}.




References

[El79] Elliott, P. D. T. A., Probabilistic number theory. I . (1979), xxii+359+xxxiii pp. (2 plates).

\noindent\rule{\linewidth}{0.4pt}


%% ============================================================
\section{Problem \#638}
\label{prob:638}

\noindent\textbf{Status:} OPEN \quad|\quad \textbf{Prize:} none \quad|\quad \textbf{Updated:} 2025-08-31

\noindent\textbf{Tags:} graph theory, ramsey theory

\medskip

\noindent\textbf{Statement:}

Let $S$ be a family of finite graphs such that for every $n$ there is some $G_n\in S$ such that if the edges of $G_n$ are coloured with $n$ colours then there is a monochromatic triangle. Is it true that for every infinite cardinal $\aleph$ there is a graph $G$ of which every finite subgraph is in $S$ and if the edges of $G$ are coloured with $\aleph$ many colours then there is a monochromatic triangle.




Erd\H{o}s writes 'if the answer is affirmative many extensions and generalisations will be possible'. Barreto notes in the comments that (as is often the case with problems of this kind) presumably $S$ is intended to also be closed under taking subgraphs, or else a sparse family of complete graphs gives a trivial counterexample.

\noindent\rule{\linewidth}{0.4pt}


%% ============================================================
\section{Problem \#640}
\label{prob:640}

\noindent\textbf{Status:} OPEN \quad|\quad \textbf{Prize:} none \quad|\quad \textbf{Updated:} 2025-08-31

\noindent\textbf{Tags:} graph theory, chromatic number

\medskip

\noindent\textbf{Statement:}

Let $k\geq 3$. Does there exist some $f(k)$ such that if a graph $G$ has chromatic number $\geq f(k)$ then $G$ must contain some odd cycle whose vertices span a graph of chromatic number $\geq k$?




A problem of Erd\H{o}s and Hajnal. This is trivial when $k=3$, since any non-bipartite graph contains an odd cycle, and all odd cycles have chromatic number $3$. Steiner has observed in the comments that this is equivalent to the problem in which we replace 'odd cycle' with 'path'.

\noindent\rule{\linewidth}{0.4pt}


%% ============================================================
\section{Problem \#642}
\label{prob:642}

\noindent\textbf{Status:} OPEN \quad|\quad \textbf{Prize:} none \quad|\quad \textbf{Updated:} 2025-08-31

\noindent\textbf{Tags:} graph theory, cycles

\medskip

\noindent\textbf{Statement:}

Let $f(n)$ be the maximal number of edges in a graph on $n$ vertices such that all cycles have more vertices than chords. Is it true that $f(n)\ll n$?




A problem of Hamburger and Szegedy. A chord is an edge between two vertices of the cycle which are not consecutive in the cycle. Chen, Erd\H{o}s, and Staton \cite{CES96} proved $f(n) \ll n^{3/2}$. Dragani\'{c}, Methuku, Munh\'{a} Correia, and Sudakov \cite{DMMS24} have improved this to\[f(n) \ll n(\log n)^8.\]




References

[CES96] Chen, Guantao and Erd\H{o}s, Paul and Staton, William, Proof of a conjecture of {B}ollob\'{a}s on nested cycles . J. Combin. Theory Ser. B (1996), 38--43.

[DMMS24] Dragani\'c, Nemanja and Methuku, Abhishek and Munh\'{a}{} Correia, David and Sudakov, Benny, Cycles with many chords . Random Structures Algorithms (2024), 3--16.

\noindent\rule{\linewidth}{0.4pt}


%% ============================================================
\section{Problem \#643}
\label{prob:643}

\noindent\textbf{Status:} OPEN \quad|\quad \textbf{Prize:} none \quad|\quad \textbf{Updated:} 2025-08-31

\noindent\textbf{Tags:} graph theory, hypergraphs

\medskip

\noindent\textbf{Statement:}

Let $f(n;t)$ be minimal such that if a $t$-uniform hypergraph on $n$ vertices contains at least $f(n;t)$ edges then there must be four edges $A,B,C,D$ such that\[A\cup B= C\cup D\]and\[A\cap B=C\cap D=\emptyset.\]Estimate $f(n;t)$ - in particular, is it true that for $t\geq 3$\[f(n;t)=(1+o(1))\binom{n}{t-1}?\]




For $t=2$ this is asking for the maximal number of edges on a graph which contains no $C_4$, and so $f(n;2)=(1/2+o(1))n^{3/2}$. F\"{u}redi \cite{Fu84} proved that $f(n;3) \ll n^2$ and $f(n;3) &#62; \binom{n}{2}$ for infinitely many $n$. Pikhurko and Verstra\"{e}te \cite{PiVe09} have proved $f(n;3)\leq \frac{13}{9}\binom{n}{2}$ for all $n$. More generally, F\"{u}redi \cite{Fu84} proved that\[\binom{n-1}{t-1}+\left\lfloor\frac{n-1}{t}\right\rfloor\leq f(n;t) &#60; \frac{7}{2}\binom{n}{t-1},\]and conjectured the lower bound is sharp for $t\geq 4$. Pikhurko and Verstra\"{e}te \cite{PiVe09} have proved that\[1 \leq \limsup_{n\to \infty} \frac{f(n;t)}{\binom{n}{t-1}}\leq \min\left(\frac{7}{4},1+\frac{2}{\sqrt{t}}\right)\]for all $t\geq 3$. F\"{u}redi \cite{Fu84} proved that $f(n;3)/\binom{n}{2}$ converges as $n\to \infty$, but the existence of the limit for $t\geq 4$ is unknown.




References

[Fu84] F\"uredi, Z., Hypergraphs in which all disjoint pairs have distinct unions . Combinatorica (1984), 161--168.

[PiVe09] Pikhurko, Oleg and Verstra\"{e}te, Jacques, The maximum size of hypergraphs without generalized 4-cycles . J. Combin. Theory Ser. A (2009), 637--649.

\noindent\rule{\linewidth}{0.4pt}


%% ============================================================
\section{Problem \#644}
\label{prob:644}

\noindent\textbf{Status:} OPEN \quad|\quad \textbf{Prize:} none \quad|\quad \textbf{Updated:} 2025-08-31

\noindent\textbf{Tags:} combinatorics

\medskip

\noindent\textbf{Statement:}

Let $f(k,r)$ be minimal such that if $A_1,A_2,\ldots$ is a family of sets, all of size $k$, such that for every collection of $r$ of the $A_is$ there is some pair $\{x,y\}$ which intersects all of the $A_j$, then there is some set of size $f(k,r)$ which intersects all of the sets $A_i$. Is it true that\[f(k,7)=(1+o(1))\frac{3}{4}k?\]Is it true that for any $r\geq 3$ there exists some constant $c_r$ such that\[f(k,r)=(1+o(1))c_rk?\]




A problem of Erd\H{o}s, Fon-Der-Flaass, Kostochka, and Tuza \cite{EFKT92}, who proved that $f(k,3)=2k$ and $f(k,4)=\lfloor 3k/2\rfloor$ and $f(k,5)=\lfloor 5k/4\rfloor$, and further that $f(k,6)=k$.




References

[EFKT92] Erd\H{o}s, P. and Fon-Der-Flaass, D. and Kostochka, A. V. and Tuza, Zs., Small transversals in uniform hypergraphs . Siberian Adv. Math. (1992), 82-88.

\noindent\rule{\linewidth}{0.4pt}


%% ============================================================
\section{Problem \#653}
\label{prob:653}

\noindent\textbf{Status:} OPEN \quad|\quad \textbf{Prize:} none \quad|\quad \textbf{Updated:} 2025-08-31

\noindent\textbf{Tags:} geometry, distances

\medskip

\noindent\textbf{Statement:}

Let $x_1,\ldots,x_n\in \mathbb{R}^2$ and let $R(x_i)=\#\{ \lvert x_j-x_i\rvert : j\neq i\}$, where the points are ordered such that\[R(x_1)\leq \cdots \leq R(x_n).\]Let $g(n)$ be the maximum number of distinct values the $R(x_i)$ can take. Is it true that $g(n) \geq (1-o(1))n$?




Erd\H{o}s and Fishburn proved $g(n)&#62;\frac{3}{8}n$ and Csizmadia proved $g(n)&#62;\frac{7}{10}n$. Both groups proved $g(n) &lt; n-cn^{2/3}$ for some constant $c&gt;0$.

\noindent\rule{\linewidth}{0.4pt}


%% ============================================================
\section{Problem \#654}
\label{prob:654}

\noindent\textbf{Status:} OPEN \quad|\quad \textbf{Prize:} none \quad|\quad \textbf{Updated:} 2025-08-31

\noindent\textbf{Tags:} geometry, distances

\medskip

\noindent\textbf{Statement:}

Let $f(n)$ be such that, given any $x_1,\ldots,x_n\in \mathbb{R}^2$ with no four points on a circle, there exists some $x_i$ with at least $f(n)$ many distinct distances to other $x_j$. Estimate $f(n)$ - in particular, is it true that\[f(n)&#62;(1-o(1))n?\]Or at least\[f(n) &#62; (1/3+c)n\]for some $c&#62;0$, for all large $n$?




It is trivial that $f(n) \geq (n-1)/3$. In \cite{Er97e} Erd\H{o}s asks whether $f(n)&#62;(1-o(1))n$, but says this is 'perhaps too optimistic but I am sure that [the trivial bound] can be significantly improved'. In \cite{Er87b} and \cite{ErPa90} Erd\H{o}s and Pach ask this under the additional assumption that there are no three points on a line (so that the points are in general position), although they only ask the weaker question whether there is a lower bound of the shape $(\tfrac{1}{3}+c)n$ for some constant $c&#62;0$. They suggest the lower bound $(1-o(1))n$ is true under the assumption that any circle around a point $x_i$ contains at most $2$ other $x_j$. The strongest form of this conjecture was disproved by Aletheia \cite{Fe26}, which constructed a set of $n$ points in $\mathbb{R}^2$, no four on a circle, such that there are at most $\frac{3}{4}n$ distinct distances from any single point (although this construction has all the points on the union of two lines, so does not help with the stronger version with the points in general position).




References

[Er87b] Erd\H{o}s, P., Some combinatorial and metric problems in geometry . Intuitive geometry (Si\'{o}fok, 1985) (1987), 167-177.

[Er97e] Erd\H{o}s, Paul, Some of my favourite unsolved problems . Math. Japon. (1997), 527-537.

[ErPa90] Erd\H{o}s, P. and Pach, J., Variations on the theme of repeated distances . Combinatorica (1990), 261--269.

[Fe26] T. Feng et al, Semi-Autonomous Mathematics Discovery with Gemini: A Case Study on the Erd\H{o}s Problems . arXiv:2601.22401 (2026).

\noindent\rule{\linewidth}{0.4pt}


%% ============================================================
\section{Problem \#655}
\label{prob:655}

\noindent\textbf{Status:} OPEN \quad|\quad \textbf{Prize:} none \quad|\quad \textbf{Updated:} 2025-08-31

\noindent\textbf{Tags:} geometry, distances

\medskip
\noindent\textit{ambiguous statement}


\noindent\textbf{Statement:}

Let $x_1,\ldots,x_n\in \mathbb{R}^2$ be such that no circle whose centre is one of the $x_i$ contains three other points. Are there at least\[(1+c)\frac{n}{2}\]distinct distances determined between the $x_i$, for some constant $c&#62;0$ and all $n$ sufficiently large?




A problem of Erd\H{o}s and Pach. It is easy to see that this assumption implies that there are at least $\frac{n-1}{2}$ distinct distances determined by every point. Zach Hunter has observed that taking $n$ points equally spaced on a circle disproves this conjecture. In the spirit of related conjectures of Erd\H{o}s and others, presumably some kind of assumption that the points are in general position (e.g. no three on a line and no four on a circle) was intended.

\noindent\rule{\linewidth}{0.4pt}


%% ============================================================
\section{Problem \#657}
\label{prob:657}

\noindent\textbf{Status:} OPEN \quad|\quad \textbf{Prize:} none \quad|\quad \textbf{Updated:} 2025-08-31

\noindent\textbf{Tags:} geometry, distances

\medskip

\noindent\textbf{Statement:}

Is it true that if $A\subset \mathbb{R}^2$ is a set of $n$ points such that every subset of $3$ points determines $3$ distinct distances (i.e. $A$ has no isosceles triangles) then $A$ must determine at least $f(n)n$ distinct distances, for some $f(n)\to \infty$?




In \cite{Er73} Erd\H{o}s attributes this problem (more generally in $\mathbb{R}^k$) to himself and Davies. In \cite{Er97e} he does not mention Davis, but says this problem was investigated by himself, F\"{u}redi, Ruzsa, and Pach. In \cite{Er73} Erd\H{o}s says it is not even known in $\mathbb{R}$ whether $f(n)\to \infty$. Sarosh Adenwalla has observed that this is equivalent to minimising the number of distinct differences in a set $A\subset \mathbb{R}$ of size $n$ without three-term arithmetic progressions. Dumitrescu \cite{Du08} proved that, in these terms,\[(\log n)^c \leq f(n) \leq 2^{O(\sqrt{\log n})}\]for some constant $c&gt;0$. Hunter observed in the comments that a result of Ruzsa coupled with standard tools of additive combinatorics (with details given by Alfaiz and Tang) allow recent progress on the size of subsets without three-term arithmetic progression (see \cite{BlSi23} which improves slightly on the bounds due to Kelley and Meka \cite{KeMe23}) yield\[2^{c(\log n)^{1/9}}\leq f(n)\]for some constant $c&gt;0$. Straus has observed that if $2^k\geq n$ then there exist $n$ points in $\mathbb{R}^k$ which contain no isosceles triangle and determine at most $n-1$ distances. See also [135] .




References

[BlSi23] T. F. Bloom and O. Sisask, An improvement to the Kelley-Meka bounds on three-term arithmetic progressions . arXiv:2309.02353 (2023).

[Du08] Dumitrescu, Adrian, On distinct distances and {$\lambda$}-free point sets . Discrete Math. (2008), 6533--6538.

[Er73] Erd\H{o}s, P., Problems and results on combinatorial number theory . A survey of combinatorial theory (Proc. Internat. Sympos., Colorado State Univ., Fort Collins, Colo., 1971) (1973), 117-138.

[Er97e] Erd\H{o}s, Paul, Some of my favourite unsolved problems . Math. Japon. (1997), 527-537.

[KeMe23] Kelley, Z. and Meka, R., Strong Bounds for 3-Progressions . arXiv:2302.05537 (2023).

\noindent\rule{\linewidth}{0.4pt}


%% ============================================================
\section{Problem \#660}
\label{prob:660}

\noindent\textbf{Status:} OPEN \quad|\quad \textbf{Prize:} none \quad|\quad \textbf{Updated:} 2025-08-31

\noindent\textbf{Tags:} geometry, distances, convex

\medskip
\noindent\textit{ambiguous statement}


\noindent\textbf{Statement:}

Let $x_1,\ldots,x_n\in \mathbb{R}^3$ be the vertices of a convex polyhedron. Are there at least\[(1-o(1))\frac{n}{2}\]many distinct distances between the $x_i$?




For the similar problem in $\mathbb{R}^2$ there are always at least $n/2$ distances, as proved by Altman \cite{Al63} (see [93] ). In \cite{Er75f} Erd\H{o}s claims that Altman proved that the vertices determine $\gg n$ many distinct distances, but gives no reference.




References

[Al63] Altman, E., On a problem of P. Erd\H{o}s . Amer. Math. Monthly (1963), 148-157.

[Er75f] Erd\H{o}s, Paul, On some problems of elementary and combinatorial geometry . Ann. Mat. Pura Appl. (4) (1975), 99-108.

\noindent\rule{\linewidth}{0.4pt}


%% ============================================================
\section{Problem \#661}
\label{prob:661}

\noindent\textbf{Status:} OPEN \quad|\quad \textbf{Prize:} \$50 \quad|\quad \textbf{Updated:} 2025-08-31

\noindent\textbf{Tags:} geometry, distances

\medskip

\noindent\textbf{Statement:}

Are there, for all large $n$, some points $x_1,\ldots,x_n,y_1,\ldots,y_n\in \mathbb{R}^2$ such that the number of distinct distances $d(x_i,y_j)$ is\[o\left(\frac{n}{\sqrt{\log n}}\right)?\]




One can also ask this for points in $\mathbb{R}^3$. In $\mathbb{R}^4$ Lenz observed that there are $x_1,\ldots,x_n,y_1,\ldots,y_n\in \mathbb{R}^4$ such that $d(x_i,y_j)=1$ for all $i,j$, taking the points on two orthogonal circles. More generally, if $F(2n)$ is the minimal number of such distances, and $f(2n)$ is minimal number of distinct distances between any $2n$ points in $\mathbb{R}^2$, then is $F =o(f)$? See also [89] .

\noindent\rule{\linewidth}{0.4pt}


%% ============================================================
\section{Problem \#662}
\label{prob:662}

\noindent\textbf{Status:} OPEN \quad|\quad \textbf{Prize:} none \quad|\quad \textbf{Updated:} 2025-08-31

\noindent\textbf{Tags:} geometry, distances

\medskip
\noindent\textit{ambiguous statement}


\noindent\textbf{Statement:}

Consider the triangular lattice with minimal distance between two points $1$. Denote by $f(t)$ the number of distances from any points $\leq t$. For example $f(1)=6$, $f(\sqrt{3})=12$, and $f(3)=18$. Let $x_1,\ldots,x_n\in \mathbb{R}^2$ be such that $d(x_i,x_j)\geq 1$ for all $i\neq j$. Is it true that, provided $n$ is sufficiently large depending on $t$, the number of distances $d(x_i,x_j)\leq t$ is less than or equal to $f(t)$ with equality perhaps only for the triangular lattice? In particular, is it true that the number of distances $\leq \sqrt{3}-\epsilon$ is less than $1$?




A problem of Erd\H{o}s, Lov\'{a}sz, and Vesztergombi. This is essentially verbatim the problem description in \cite{Er97e}, but this does not make sense as written; there must be at least one typo. Suggestions about what this problem intends are welcome. Erd\H{o}s also goes on to write 'Perhaps the following stronger conjecture holds: Let $t_1&#60;t_2&#60;\cdots$ be the set of distances occurring in the triangular lattice. $t_1=1$ $t_2=\sqrt{3}$ $t_3=3$ $t_4=5$ etc. Is it true that there is an $\epsilon_n$ so that for every set $y_1,\ldots,$ with $d(y_i,y_j)\geq 1$ the number of distances $d(y_i,y_j)&#60;t_n$ is less than $f(t_n)$?' Again, this is nonsense interpreted literally; I am not sure what Erd\H{o}s intended.




References

[Er97e] Erd\H{o}s, Paul, Some of my favourite unsolved problems . Math. Japon. (1997), 527-537.

\noindent\rule{\linewidth}{0.4pt}


%% ============================================================
\section{Problem \#663}
\label{prob:663}

\noindent\textbf{Status:} OPEN \quad|\quad \textbf{Prize:} none \quad|\quad \textbf{Updated:} 2025-08-31

\noindent\textbf{Tags:} number theory

\medskip

\noindent\textbf{Statement:}

Let $k\geq 2$ and $q(n,k)$ denote the least prime which does not divide $\prod_{1\leq i\leq k}(n+i)$. Is it true that, if $k$ is fixed and $n$ is sufficiently large, we have\[q(n,k)&#60;(1+o(1))\log n?\]




A problem of Erd\H{o}s and Pomerance. The bound $q(n,k)&lt;(1+o(1))k\log n$ is easy. It may be true this improved bound holds even up to $k=o(\log n)$. A heuristic argument in favour of this is provided by Tao in the comments. See also [457] .

\noindent\rule{\linewidth}{0.4pt}


%% ============================================================
\section{Problem \#665}
\label{prob:665}

\noindent\textbf{Status:} OPEN \quad|\quad \textbf{Prize:} none \quad|\quad \textbf{Updated:} 2025-08-31

\noindent\textbf{Tags:} combinatorics

\medskip

\noindent\textbf{Statement:}

A pairwise balanced design for $\{1,\ldots,n\}$ is a collection of sets $A_1,\ldots,A_m\subseteq \{1,\ldots,n\}$ such that $2\leq \lvert A_i\rvert &lt;n$ and every pair of distinct elements $x,y\in \{1,\ldots,n\}$ is contained in exactly one $A_i$. Is there a constant $C&gt;0$ and, for all large $n$, a pairwise balanced design such that\[\lvert A_i\rvert &#62; n^{1/2}-C\]for all $1\leq i\leq m$?




A problem of Erd\H{o}s and Larson \cite{ErLa82}. In general, as Erd\H{o}s asks in \cite{Er97f}, find the slowest growing function $h$ such that, for all large $n$, there exists a pairwise balanced design with\[\lvert A_i\rvert &gt; n^{1/2}-h(n)\]for all $1\leq i\leq m$. The problem above asks whether $h(n)\ll 1$. Erd\H{o}s and Larson prove that $h(n) \ll n^{1/2-c}$ for some constant $c&gt;0$, and note this can be improved to $h(n)\ll (\log n)^2$ assuming Cramer-type bounds on the difference between consecutive primes. Shrikhande and Singhi \cite{ShSi85} have proved that the answer is no conditional on the conjecture that the order of every projective plane is a prime power (see [723] ), by proving that every pairwise balanced design on $n$ points in which each block is of size $\geq n^{1/2}-c$ can be embedded in a projective plane of order $n+i$ for some $i\leq c+2$, if $n$ is sufficiently large. In general, if $H(n)$ is the largest prime gap $\leq n$, then the above reuslts show that, assuming the prime power conjecture, $H(n)\asymp h(n)$.




References

[Er97f] Erd\H{o}s, Paul, Some unsolved problems . Combinatorics, geometry and probability (Cambridge, 1993) (1997), 1-10.

[ErLa82] Erd\H{o}s, P. and Larson, J., On pairwise balanced block designs with the sizes of blocks as uniform as possible . Annals of Discrete Mathematics (1982), 129-134.

[ShSi85] S. S. Shrikhande and N. M. Singhi, On a problem of Erd\H{o}s and Larson . Combinatorica (1985), 351-358.

\noindent\rule{\linewidth}{0.4pt}


%% ============================================================
\section{Problem \#667}
\label{prob:667}

\noindent\textbf{Status:} OPEN \quad|\quad \textbf{Prize:} none \quad|\quad \textbf{Updated:} 2025-08-31

\noindent\textbf{Tags:} graph theory, ramsey theory

\medskip

\noindent\textbf{Statement:}

Let $p,q\geq 1$ be fixed integers. We define $H(n)=H(N;p,q)$ to be the largest $m$ such that any graph on $n$ vertices where every set of $p$ vertices spans at least $q$ edges must contain a complete graph on $m$ vertices. Is\[c(p,q)=\liminf \frac{\log H(n)}{\log n}\]a strictly increasing function of $q$ for $1\leq q\leq \binom{p-1}{2}+1$?




A problem of Erd\H{o}s, Faudree, Rousseau, and Schelp. When $q=1$ this corresponds exactly to the classical Ramsey problem, and hence for example\[\frac{1}{p-1}\leq c(p,1) \leq \frac{2}{p+1}.\]It is easy to see that if $q=\binom{p-1}{2}+1$ then $c(p,q)=1$. Erd\H{o}s, Faudree, Rousseau, and Schelp have shown that $c(p,\binom{p-1}{2})\leq 1/2$.

\noindent\rule{\linewidth}{0.4pt}


%% ============================================================
\section{Problem \#668}
\label{prob:668}

\noindent\textbf{Status:} OPEN \quad|\quad \textbf{Prize:} none \quad|\quad \textbf{Updated:} 2025-08-31

\noindent\textbf{Tags:} geometry, distances

\medskip

\noindent\textbf{Statement:}

Is it true that the number of incongruent sets of $n$ points in $\mathbb{R}^2$ which maximise the number of unit distances tends to infinity as $n\to\infty$? Is it always $&#62;1$ for $n&#62;3$?




In fact this is $=1$ also for $n=4$, the unique example given by two equilateral triangles joined by an edge. Computational evidence of Engel, Hammond-Lee, Su, Varga, and Zs\'{a}mboki \cite{EHSVZ25} and Alexeev, Mixon, and Parshall \cite{AMP25} suggests that this count is $=1$ for various other $5\leq n\leq 21$ (although these calculations were checking only up to graph isomorphism, rather than congruency). The actual maximal number of unit distances is the subject of [90] .




References

[AMP25] B. Alexeev, D. Mixon, and H. Parshall, The Erd\H{o}s unit distance problem for small point sets . arXiv:2412.11914 (2025).

[EHSVZ25] P. Engel, O. Hammond-Lee, Y. Su, D. Varga, and P. Zs\'{a}mboki, Diverse beam search to find densest-known planar unit distance graphs . arXiv:2406.15317 (2025).

\noindent\rule{\linewidth}{0.4pt}


%% ============================================================
\section{Problem \#669}
\label{prob:669}

\noindent\textbf{Status:} OPEN \quad|\quad \textbf{Prize:} none \quad|\quad \textbf{Updated:} 2025-08-31

\noindent\textbf{Tags:} geometry

\medskip
\noindent\textit{orchard problems}


\noindent\textbf{Statement:}

Let $F_k(n)$ be minimal such that for any $n$ points in $\mathbb{R}^2$ there exist at most $F_k(n)$ many distinct lines passing through at least $k$ of the points, and $f_k(n)$ similarly but with lines passing through exactly $k$ points. Estimate $f_k(n)$ and $F_k(n)$ - in particular, determine $\lim F_k(n)/n^2$ and $\lim f_k(n)/n^2$.




Trivially $f_k(n)\leq F_k(n)$ and $f_2(n)=F_2(n)=\binom{n}{2}$. The problem with $k=3$ is the classical 'Orchard problem' of Sylvester. Burr, Gr\"{u}nbaum, and Sloane \cite{BGS74} have proved that\[f_3(n)=\frac{n^2}{6}-O(n)\]and\[F_3(n)=\frac{n^2}{6}-O(n).\]There is a trivial upper bound of $F_k(n) \leq \binom{n}{2}/\binom{k}{2}$, and hence\[\lim F_k(n)/n^2 \leq \frac{1}{k(k-1)}.\]See also [101] .




References

[BGS74] Burr, Stefan A. and Gr\"{u}nbaum, Branko and Sloane, N. J. A., The orchard problem . Geometriae Dedicata (1974), 397-424.

\noindent\rule{\linewidth}{0.4pt}


%% ============================================================
\section{Problem \#670}
\label{prob:670}

\noindent\textbf{Status:} OPEN \quad|\quad \textbf{Prize:} none \quad|\quad \textbf{Updated:} 2026-04-16

\noindent\textbf{Tags:} geometry, distances

\medskip

\noindent\textbf{Statement:}

Let $A\subseteq \mathbb{R}^d$ be a set of $n$ points such that all pairwise distances differ by at least $1$. Is the diameter of $A$ at least $(1+o(1))n^2$?




The lower bound of $\binom{n}{2}$ for the diameter is trivial. Erd\H{o}s \cite{Er97f} proved the claim when $d=1$. The quantifiers are slightly ambiguous, but likely Erd\H{o}s had in mind the regime where $d$ is fixed and the $o(1)$ term $\to 0$ as $n\to \infty$ (perhaps at a rate which depends on $d$). This was disproved if $n$ grows with $d$ by Ho \cite{Ho26}: there exist infinitely many $n$ and, with $d=n^2-n$, a set $A\subset \mathbb{R}^d$ such that all pairwise distances differ by at least $1$, with diameter at most\[\left(1-\frac{1}{\pi^2}+o(1)\right)n^2\approx 0.898n^2.\]




References

[Er97f] Erd\H{o}s, Paul, Some unsolved problems . Combinatorics, geometry and probability (Cambridge, 1993) (1997), 1-10.

[Ho26] B. S. Ho, Erd\H{o}s's diameter conjecture for separated distances fails in high dimensions . arXiv:2604.15305 (2026).

\noindent\rule{\linewidth}{0.4pt}


%% ============================================================
\section{Problem \#671}
\label{prob:671}

\noindent\textbf{Status:} OPEN \quad|\quad \textbf{Prize:} \$250 \quad|\quad \textbf{Updated:} 2025-08-31

\noindent\textbf{Tags:} analysis

\medskip

\noindent\textbf{Statement:}

Given $a_{i}^n\in [-1,1]$ for all $1\leq i\leq n&#60;\infty$ we define $p_{i}^n$ as the unique polynomial of degree $n-1$ such that $p_{i}^n(a_{i}^n)=1$ and $p_{i}^n(a_{i'}^n)=0$ if $1\leq i'\leq n$ with $i\neq i'$. We similarly define\[\mathcal{L}^nf(x) = \sum_{1\leq i\leq n}f(a_i^n)p_i^n(x),\]the unique polynomial of degree $n-1$ which agrees with $f$ on $a_i^n$ for $1\leq i\leq n$ (that is, the sequence of Lagrange interpolation polynomials). Is there such a sequence of $a_i^n$ such that for every continuous $f:[-1,1]\to \mathbb{R}$ there exists some $x\in [-1,1]$ where\[\limsup_{n\to \infty} \sum_{1\leq i\leq n}\lvert p_{i}^n(x)\rvert=\infty\]and yet\[\mathcal{L}^nf(x) \to f(x)?\]Is there such a sequence such that\[\limsup_{n\to \infty} \sum_{1\leq i\leq n}\lvert p_{i}^n(x)\rvert=\infty\]for every $x\in [-1,1]$ and yet for every continuous $f:[-1,1]\to \mathbb{R}$ there exists $x\in [-1,1]$ with\[\mathcal{L}^nf(x) \to f(x)?\]




Bernstein \cite{Be31} proved that for any choice of $a_i^n$ there exists $x_0\in [-1,1]$ such that\[\limsup_{n\to \infty} \sum_{1\leq i\leq n}\lvert p_{i}^n(x)\rvert=\infty.\]Erd\H{o}s and V\'{e}rtesi \cite{ErVe80} proved that for any choice of $a_i^n$ there exists a continuous $f:[-1,1]\to \mathbb{R}$ such that\[\limsup_{n\to \infty} \lvert \mathcal{L}^nf(x)\rvert=\infty\]for almost all $x\in [-1,1]$.




References

[Be31] S. Bernstein, Sur la limitation des valeurs d'un polynome $P_n(x)$ de degr\'{e} n sur tout un segment par ses valeurs en $(n+1)$ points du segment . Izv. Akad. Nauk. SSSR (1931), 1025-1050.

[ErVe80] Erd\H{o}s, P. and V\'{e}rtesi, P., On the almost everywhere divergence of Lagrange interpolatory polynomials for arbitrary system of nodes . Acta Math. Acad. Sci. Hungar. (1980), 71-89.

\noindent\rule{\linewidth}{0.4pt}


%% ============================================================
\section{Problem \#675}
\label{prob:675}

\noindent\textbf{Status:} OPEN \quad|\quad \textbf{Prize:} none \quad|\quad \textbf{Updated:} 2025-08-31

\noindent\textbf{Tags:} number theory

\medskip

\noindent\textbf{Statement:}

We say that $A\subset \mathbb{N}$ has the translation property if, for every $n$, there exists some integer $t_n\geq 1$ such that, for all $1\leq a\leq n$,\[a\in A\quad\textrm{ if and only if }\quad a+t_n\in A.\]{UL} {LI}Does the set of the sums of two squares have the translation property?{/LI} {LI}If we partition all primes into $P\sqcup Q$, such that each set contains $\gg x/\log x$ many primes $\leq x$ for all large $x$, then can the set of integers only divisible by primes from $P$ have the translation property?{/LI} {LI}If $A$ is the set of squarefree numbers then how fast does the minimal such $t_n$ grow? Is it true that $t_n&#62;\exp(n^c)$ for some constant $c&#62;0$?{/LI} {/UL}




Elementary sieve theory implies that the set of squarefree numbers has the translation property. More generally, Brun's sieve can be used to prove that if $B\subseteq \mathbb{N}$ is a set of pairwise coprime integers with $\sum_{b&#60;x}\frac{1}{b}=o(\log\log x)$ then $A=\{ n: b\nmid n\textrm{ for all }b\in A\}$ has the translation property. Erd\H{o}s did not know what happens if the condition on $\sum_{b&#60;x}\frac{1}{b}$ is weakened or dropped altogether.

\noindent\rule{\linewidth}{0.4pt}


%% ============================================================
\section{Problem \#676}
\label{prob:676}

\noindent\textbf{Status:} OPEN \quad|\quad \textbf{Prize:} none \quad|\quad \textbf{Updated:} 2025-08-31

\noindent\textbf{Tags:} number theory

\medskip

\noindent\textbf{Statement:}

Is every sufficiently large integer of the form\[ap^2+b\]for some prime $p$ and integer $a\geq 1$ and $0\leq b&#60;p$?




The sieve of Eratosthenes implies that almost all integers are of this form, and the Brun-Selberg sieve implies the number of exceptions in $[1,x]$ is $\ll x/(\log x)^c$ for some constant $c&#62;0$. Erd\H{o}s \cite{Er79} believed it is 'rather unlikely' that all large integers are of this form. What if the condition that $p$ is prime is omitted? Selfridge and Wagstaff made a 'preliminary computer search' and suggested that there are infinitely many $n$ not of this form even without the condition that $p$ is prime. It should be true that the number of exceptions in $[1,x]$ is $&#60;x^c$ for some constant $c&#60;1$. Most generally, given some infinite set $A\subseteq \mathbb{N}$ and function $f:A\to \mathbb{N}$ one can ask for sufficient conditions on $A$ and $f$ that guarantee every large number (or almost all numbers) can be written as\[am^2+b\]for some $m\in A$ and $a\geq 1$ and $0\leq b&#60;f(m)$. In another direction, one can ask what is the minimal $c_n$ such that $n$ can be written as $n=ap^2+b$ with $0\leq b&#60;c_np$ for some $p\leq \sqrt{n}$. This problem asks whether $c_n\leq 1$ eventually, but in \cite{Er79d} Erd\H{o}s suggests that in fact $\limsup c_n=\infty$. Is it true that $c_n&#60;n^{o(1)}$?




References

[Er79] Erd\H{o}s, Paul, Some unconventional problems in number theory . Math. Mag. (1979), 67-70.

[Er79d] Erd\H{o}s, P., Some unconventional problems in number theory . Acta Math. Acad. Sci. Hungar. (1979), 71-80.

\noindent\rule{\linewidth}{0.4pt}


%% ============================================================
\section{Problem \#677}
\label{prob:677}

\noindent\textbf{Status:} OPEN \quad|\quad \textbf{Prize:} none \quad|\quad \textbf{Updated:} 2025-08-31

\noindent\textbf{Tags:} number theory

\medskip

\noindent\textbf{Statement:}

Let $M(n,k)=[n+1,\ldots,n+k]$ be the least common multiple of $\{n+1,\ldots,n+k\}$. Is it true that for all $m\geq n+k$\[M(n,k) \neq M(m,k)?\]




The Thue-Siegel theorem implies that, for fixed $k$, there are only finitely many $m,n$ such that $m\geq n+k$ and $M(n,k)=M(m,k)$. In general, how many solutions does $M(n,k)=M(m,l)$ have when $m\geq n+k$ and $l&gt;1$? Erd\H{o}s expects very few (and none when $l\geq k$). The only solutions Erd\H{o}s knew were $M(4,3)=M(13,2)$ and $M(3,4)=M(19,2)$. In \cite{Er79d} Erd\H{o}s conjectures the stronger fact that (aside from a finite number of exceptions) if $k&gt;2$ and $m\geq n+k$ then $\prod_{i\leq k}(n+i)$ and $\prod_{i\leq k}(m+i)$ cannot have the same set of prime factors. See also [678] , [686] , and [850] . This is discussed in problem B35 of Guy's collection \cite{Gu04}.




References

[Er79d] Erd\H{o}s, P., Some unconventional problems in number theory . Acta Math. Acad. Sci. Hungar. (1979), 71-80.

[Gu04] Guy, Richard K., Unsolved problems in number theory . (2004), xviii+437.

\noindent\rule{\linewidth}{0.4pt}


%% ============================================================
\section{Problem \#679}
\label{prob:679}

\noindent\textbf{Status:} OPEN \quad|\quad \textbf{Prize:} none \quad|\quad \textbf{Updated:} 2025-08-31

\noindent\textbf{Tags:} number theory

\medskip

\noindent\textbf{Statement:}

Let $\epsilon&#62;0$ and $\omega(n)$ count the number of distinct prime factors of $n$. Are there infinitely many values of $n$ such that\[\omega(n-k) &#60; (1+\epsilon)\frac{\log k}{\log\log k}\]for all $k&#60;n$ which are sufficiently large depending on $\epsilon$ only? Can one show the stronger version with\[\omega(n-k) &#60; \frac{\log k}{\log\log k}+O(1)\]is false?




One can ask similar questions for $\Omega$, the number of prime factors with multiplicity, where $\log k/\log\log k$ is replaced by $\log_2k$. In the comments DottedCalculator has disproved the second stronger version, proving that in fact for all large $n$ there exists $k&lt;n$ such that\[\omega(n-k)\geq \frac{\log k}{\log\log k} + c\frac{\log k}{(\log\log k)^2}\]for some constant $c&gt;0$. Lau \cite{La26} has proved that there exists a constant $C&gt;0$ such that, for infinitely many $n$,\[\omega(n-k) \leq \Omega(n-k)\leq C\log k\]for all $1&lt;k&lt;n$. Lau conjectures that, for $\omega$, this can only be improved in the constant factor, and thus the original question has a negative answer. See also [248] , [413] , and [1203] .




References

[La26] C. F. Lau, On the number of prime factors of consecutive integers . arXiv:2604.15042 (2026).

\noindent\rule{\linewidth}{0.4pt}


%% ============================================================
\section{Problem \#680}
\label{prob:680}

\noindent\textbf{Status:} OPEN \quad|\quad \textbf{Prize:} none \quad|\quad \textbf{Updated:} 2025-08-31

\noindent\textbf{Tags:} number theory, primes

\medskip

\noindent\textbf{Statement:}

Is it true that, for all sufficiently large $n$, there exists some $k$ such that\[p(n+k)&#62;k^2+1,\]where $p(m)$ denotes the least prime factor of $m$? Can one prove this is false if we replace $k^2+1$ by $e^{(1+\epsilon)\sqrt{k}}+C_\epsilon$, for all $\epsilon&#62;0$, where $C_\epsilon&#62;0$ is some constant?




This follows from 'plausible assumptions on the distribution of primes' (as does the question with $k^2$ replaced by $k^d$ for any $d$); the challenge is to prove this unconditionally. Erd\H{o}s observed that Cramer's conjecture\[\limsup_{k\to \infty} \frac{p_{k+1}-p_k}{(\log k)^2}=1\]implies that for all $\epsilon&gt;0$ and all sufficiently large $n$ there exists some $k$ such that\[p(n+k)&gt;e^{(1-\epsilon)\sqrt{k}}.\]There is now evidence, however, that Cramer's conjecture is false; a more refined heuristic by Granville \cite{Gr95} suggests this $\limsup$ is $2e^{-\gamma}\approx 1.119\cdots$, and so perhaps the $1+\epsilon$ in the second question should be replaced by $2e^{-\gamma}+\epsilon$. See also [681] and [682] .




References

[Gr95] Granville, Andrew, Harald {C}ram\'{e}r and the distribution of prime numbers . Scand. Actuar. J. (1995), 12--28.

\noindent\rule{\linewidth}{0.4pt}


%% ============================================================
\section{Problem \#681}
\label{prob:681}

\noindent\textbf{Status:} OPEN \quad|\quad \textbf{Prize:} none \quad|\quad \textbf{Updated:} 2025-08-31

\noindent\textbf{Tags:} number theory, primes

\medskip

\noindent\textbf{Statement:}

Is it true that for all large $n$ there exists $k$ such that $n+k$ is composite and\[p(n+k)&#62;k^2,\]where $p(m)$ is the least prime factor of $m$?




Related to questions of Erd\H{o}s, Eggleton, and Selfridge. This may be true with $k^2$ replaced by $k^d$ for any $d$. See also [680] and [682] .

\noindent\rule{\linewidth}{0.4pt}


%% ============================================================
\section{Problem \#683}
\label{prob:683}

\noindent\textbf{Status:} OPEN \quad|\quad \textbf{Prize:} none \quad|\quad \textbf{Updated:} 2025-09-04

\noindent\textbf{Tags:} number theory, primes, binomial coefficients

\medskip

\noindent\textbf{Statement:}

Is it true that for every $1\leq k\leq n$ the largest prime divisor of $\binom{n}{k}$, say $P(\binom{n}{k})$, satisfies\[P\left(\binom{n}{k}\right)\geq \min(n-k+1, k^{1+c})\]for some constant $c&#62;0$?




A theorem of Sylvester and Schur (see \cite{Er34}) states that $P(\binom{n}{k})&gt;k$ if $k\leq n/2$. Erd\H{o}s \cite{Er55d} proved that there exists some $c&gt;0$ such that, whenever $k\leq n/2$,\[P\left(\binom{n}{k}\right)\gg k\log k.\]Erd\H{o}s \cite{Er79d} writes it 'seems certain' that this holds for every $c&gt;0$, with only a finite number of exceptions (depending on $c$). Standard heuristics on prime gaps suggest that the largest prime divisor of $\binom{n}{k}$ is, for $k\leq n/2$, in fact\[&gt;e^{c\sqrt{k}}\]for some constant $c&gt;0$. This is essentially equivalent to [961] .




References

[Er34] Erd\H{o}s, Paul, A Theorem of Sylvester and Schur . J. London Math. Soc. (1934), 282--288.

[Er55d] Erd\H{o}s, P., On consecutive integers . Nieuw Arch. Wisk. (3) (1955), 124--128.

[Er79d] Erd\H{o}s, P., Some unconventional problems in number theory . Acta Math. Acad. Sci. Hungar. (1979), 71-80.

\noindent\rule{\linewidth}{0.4pt}


%% ============================================================
\section{Problem \#684}
\label{prob:684}

\noindent\textbf{Status:} OPEN \quad|\quad \textbf{Prize:} none \quad|\quad \textbf{Updated:} 2025-08-31

\noindent\textbf{Tags:} number theory, primes, binomial coefficients

\medskip

\noindent\textbf{Statement:}

For $0\leq k\leq n$ write\[\binom{n}{k} = uv\]where the only primes dividing $u$ are in $[2,k]$ and the only primes dividing $v$ are in $(k,n]$. Let $f(n)$ be the smallest $k$ such that $u&#62;n^2$. Give bounds for $f(n)$.




A classical theorem of Mahler states that for any $\epsilon&gt;0$ and integers $k$ and $l$ then, writing\[(n+1)\cdots (n+k) = ab\]where the only primes dividing $a$ are $\leq l$ and the only primes dividing $b$ are $&gt;l$, we have $a &lt; n^{1+\epsilon}$ for all sufficiently large (depending on $\epsilon,k,l$) $n$. Mahler's theorem implies $f(n)\to \infty$ as $n\to \infty$, but is ineffective, and so gives no bounds on the growth of $f(n)$. One can similarly ask for estimates on the smallest integer $f(n,k)$ such that if $m$ is the factor of $\binom{n}{k}$ containing all primes $\leq f(n,k)$ then $m &gt; n^2$. Tang and ChatGPT have proved that\[f(n)\leq n^{30/43+o(1)}.\]The same proof would prove $f(n) \leq n^{2/3+o(1)}$ assuming the Riemann Hypothesis (or Density Hypothesis). A heuristic given by Sothanaphan and ChatGPT in the comments suggests that, at least for most $n$, $f(n)\sim 2\log n$. An internal OpenAI model \cite{APSSV26} has given an elementary argument which shows $f(n) \ll (\log n)^2$, and constructs arbitrarily large $n$ such that $f(n)\geq (1/2-o(1))\log n$.




References

[APSSV26] B. Alexeev, M. Putterman, M. Sawhney, M. Sellke, and G. Valiant, Short proofs in combinatorics and number theory . arXiv:2603.29961 (2026).

\noindent\rule{\linewidth}{0.4pt}


%% ============================================================
\section{Problem \#685}
\label{prob:685}

\noindent\textbf{Status:} OPEN \quad|\quad \textbf{Prize:} none \quad|\quad \textbf{Updated:} 2025-08-31

\noindent\textbf{Tags:} number theory, primes, binomial coefficients

\medskip

\noindent\textbf{Statement:}

Let $\epsilon&#62;0$ and $n$ be large depending on $\epsilon$. Is it true that for all $n^\epsilon&#60;k\leq n^{1-\epsilon}$ the number of distinct prime divisors of $\binom{n}{k}$ is\[(1+o(1))k\sum_{k&#60;p&#60;n}\frac{1}{p}?\]Or perhaps even when $k \geq (\log n)^c$?




It is trivial that the number of prime factors is\[&#62;\frac{\log \binom{n}{k}}{\log n},\]and this inequality becomes (asymptotic) equality if $k&#62;n^{1-o(1)}$.

\noindent\rule{\linewidth}{0.4pt}


%% ============================================================
\section{Problem \#686}
\label{prob:686}

\noindent\textbf{Status:} OPEN \quad|\quad \textbf{Prize:} none \quad|\quad \textbf{Updated:} 2025-08-31

\noindent\textbf{Tags:} number theory

\medskip

\noindent\textbf{Statement:}

Can every integer $N\geq 2$ be written as\[N=\frac{\prod_{1\leq i\leq k}(m+i)}{\prod_{1\leq i\leq k}(n+i)}\]for some $k\geq 2$ and $m\geq n+k$?




If $n$ and $k$ are fixed then can one say anything about the set of integers so represented? See also [388] and [677] .

\noindent\rule{\linewidth}{0.4pt}


%% ============================================================
\section{Problem \#687}
\label{prob:687}

\noindent\textbf{Status:} OPEN \quad|\quad \textbf{Prize:} \$1000 \quad|\quad \textbf{Updated:} 2025-08-31

\noindent\textbf{Tags:} number theory

\medskip

\noindent\textbf{Statement:}

Let $Y(x)$ be the maximal $y$ such that there exists a choice of congruence classes $a_p$ for all primes $p\leq x$ such that every integer in $[1,y]$ is congruent to at least one of the $a_p\pmod{p}$. Give good estimates for $Y(x)$. In particular, can one prove that $Y(x)=o(x^2)$ or even $Y(x)\ll x^{1+o(1)}$?




This function (associated with Jacobsthal) is closely related to the problem of gaps between primes (see [4] ). The best known upper bound is due to Iwaniec \cite{Iw78},\[Y(x) \ll x^2.\]The best lower bound is due to Ford, Green, Konyagin, Maynard, and Tao \cite{FGKMT18},\[Y(x) \gg x\frac{\log x\log\log\log x}{\log\log x},\]improving on a previous bound of Rankin \cite{Ra38}. Maier and Pomerance have conjectured that $Y(x)\ll x(\log x)^{2+o(1)}$. In \cite{Er80} he writes 'It is not clear who first formulated this problem - probably many of us did it independently. I offer the maximum of \$1000 dollars and $1/2$ my total savings for clearing up of this problem.' In \cite{Er80} Erd\H{o}s also asks about a weaker variant in which all except $o(y/\log y)$ of the integers in $[1,y]$ are congruent to at least one of the $a_p\pmod{p}$, and in particular asks if the answer is very different. See also [688] and [689] . A more general Jacobsthal function is the focus of [970] .




References

[Er80] Erd\H{o}s, Paul, A survey of problems in combinatorial number theory . Ann. Discrete Math. (1980), 89-115.

[FGKMT18] Ford, Kevin and Green, Ben and Konyagin, Sergei and Maynard, James and Tao, Terence, Long gaps between primes . J. Amer. Math. Soc. (2018), 65-105.

[Iw78] Iwaniec, Henryk, On the problem of {J}acobsthal . Demonstratio Math. (1978), 225--231.

[Ra38] Rankin, R. A., The Difference between Consecutive Prime Numbers . J. London Math. Soc. (1938), 242-247.

\noindent\rule{\linewidth}{0.4pt}


%% ============================================================
\section{Problem \#688}
\label{prob:688}

\noindent\textbf{Status:} OPEN \quad|\quad \textbf{Prize:} none \quad|\quad \textbf{Updated:} 2025-08-31

\noindent\textbf{Tags:} number theory

\medskip

\noindent\textbf{Statement:}

Define $\epsilon_n$ to be maximal such that there exists some choice of congruence class $a_p$ for all primes $n^{\epsilon_n}&#60;p\leq n$ such that every integer in $[1,n]$ satisfies at least one of the congruences $\equiv a_p\pmod{p}$. Estimate $\epsilon_n$ - in particular is it true that $\epsilon_n=o(1)$?




Erd\H{o}s could prove\[\epsilon_n \gg \frac{\log\log\log n}{\log\log n}.\]See also [687] , [689] , and [1200] .

\noindent\rule{\linewidth}{0.4pt}


%% ============================================================
\section{Problem \#689}
\label{prob:689}

\noindent\textbf{Status:} OPEN \quad|\quad \textbf{Prize:} none \quad|\quad \textbf{Updated:} 2025-08-31

\noindent\textbf{Tags:} number theory

\medskip

\noindent\textbf{Statement:}

Let $n$ be sufficiently large. Is there some choice of congruence class $a_p$ for all primes $2\leq p\leq n$ such that every integer in $[1,n]$ satisfies at least two of the congruences $\equiv a_p\pmod{p}$?




One can ask a similar question replacing $2$ by any fixed integer $r$ (provided $n$ is sufficiently large depending on $r$). See also [687] and [688] . If we replace primes with all integers then this is [1205] . This problem (with $2$ replaced by $10$) is Problem 45 on Green's open problems list .

\noindent\rule{\linewidth}{0.4pt}


%% ============================================================
\section{Problem \#691}
\label{prob:691}

\noindent\textbf{Status:} OPEN \quad|\quad \textbf{Prize:} none \quad|\quad \textbf{Updated:} 2025-08-31

\noindent\textbf{Tags:} number theory

\medskip

\noindent\textbf{Statement:}

Given $A\subseteq \mathbb{N}$ let $M_A=\{ n \geq 1 : a\mid n\textrm{ for some }a\in A\}$ be the set of multiples of $A$. Find a necessary and sufficient condition on $A$ for $M_A$ to have density $1$.




A sequence $A$ for which $M_A$ has density $1$ is called a Behrend sequence. If $A$ is a set of prime numbers (or, more generally, a set of pairwise coprime integers without $1$) then a necessary and sufficient condition is that $\sum_{p\in A}\frac{1}{p}=\infty$. The general situation is more complicated. For example suppose $A$ is the union of $(n_k,(1+\eta_k)n_k)\cap \mathbb{Z}$ where $1\leq n_1&lt;n_2&lt;\cdots$ is a lacunary sequence. (This construction is sometimes called a block sequence.) If $\sum \eta_k&lt;\infty$ then the density of $M_A$ exists and is $&lt;1$. If $\eta_k=1/k$, so $\sum \eta_k=\infty$, then the density exists and is $&lt;1$. Erd\H{o}s writes it 'seems certain' that there is some threshold $\alpha\in (0,1)$ such that, if $\eta_k=k^{-\beta}$, then the density of $M_A$ is $1$ if $\beta &lt;\alpha$ and the density is $&lt;1$ if $\beta &gt;\alpha$. Tenenbaum notes in \cite{Te96} that this is certainly not true as written since if the $n_j$ grow sufficiently quickly then this sequence is never Behrend, for any choice of $\eta_k$. He then writes 'we understand from subsequent discussions with Erd\H{o}s that he had actually in mind a two-sided condition on' $n_{j+1}/n_j$. Tenenbaum \cite{Te96} proves this conjecture: if there are constants $1&lt;C_1&lt;C_2$ such that $C_1&lt;n_{i+1}/n_i&lt;C_2$ for all $i$ and $\eta_k=k^{-\beta}$ then $A$ is Behrend if $\beta&lt;\log 2$ and not Behrend if $\beta&gt;\log 2$.




References

[Te96] Tenenbaum, G., On block {B}ehrend sequences . Math. Proc. Cambridge Philos. Soc. (1996), 355--367.

\noindent\rule{\linewidth}{0.4pt}


%% ============================================================
\section{Problem \#693}
\label{prob:693}

\noindent\textbf{Status:} OPEN \quad|\quad \textbf{Prize:} none \quad|\quad \textbf{Updated:} 2025-08-31

\noindent\textbf{Tags:} number theory, divisors

\medskip

\noindent\textbf{Statement:}

Let $k\geq 2$ and $n$ be sufficiently large depending on $k$. Let $A=\{a_1&#60;a_2&#60;\cdots \}$ be the set of those integers in $[n,n^k]$ which have a divisor in $(n,2n)$. Estimate\[\max_{i} a_{i+1}-a_i.\]Is this $\leq (\log n)^{O(1)}$?




See also [446] .

\noindent\rule{\linewidth}{0.4pt}


%% ============================================================
\section{Problem \#695}
\label{prob:695}

\noindent\textbf{Status:} OPEN \quad|\quad \textbf{Prize:} none \quad|\quad \textbf{Updated:} 2025-08-31

\noindent\textbf{Tags:} number theory

\medskip

\noindent\textbf{Statement:}

Let $p_1&#60;p_2&#60;\cdots$ be a sequence of primes such that $p_{i+1}\equiv 1\pmod{p_i}$. Is it true that\[\lim_k p_k^{1/k}=\infty?\]Does there exist such a sequence with\[p_k \leq \exp(k(\log k)^{1+o(1)})?\]




Such a sequence is sometimes called a prime chain. If we take the obvious 'greedy' chain with $2=p_1$ and $p_{i+1}$ is the smallest prime $\equiv 1\pmod{p_i}$ then Linnik's theorem implies that this sequence grows like\[p_k \leq e^{e^{O(k)}}.\]It is conjectured that, for any prime $p$, there is a prime $p'\leq p(\log p)^{O(1)}$ which is congruent to $1\pmod{p}$, which would imply this sequence grows like\[p_k\leq \exp(k(\log k)^{1+o(1)}).\]An extensive study of the growth of finite prime chains was carried out by Ford, Konyagin, and Luca \cite{FKL10}. See also [696] .




References

[FKL10] Ford, Kevin and Konyagin, Sergei V. and Luca, Florian, Prime chains and {P}ratt trees . Geom. Funct. Anal. (2010), 1231--1258.

\noindent\rule{\linewidth}{0.4pt}


%% ============================================================
\section{Problem \#700}
\label{prob:700}

\noindent\textbf{Status:} OPEN \quad|\quad \textbf{Prize:} none \quad|\quad \textbf{Updated:} 2025-08-31

\noindent\textbf{Tags:} number theory, binomial coefficients

\medskip

\noindent\textbf{Statement:}

Let\[f(n)=\min_{1&lt;k\leq n/2}\textrm{gcd}\left(n,\binom{n}{k}\right).\]{UL} {LI}Characterise those composite $n$ such that $f(n)=n/P(n)$, where $P(n)$ is the largest prime dividing $n$.{/LI} {LI}Are there infinitely many composite $n$ such that $f(n)&gt;n^{1/2}$?{/LI} {LI} Is it true that, for every composite $n$,\[f(n) \ll_A \frac{n}{(\log n)^A}\]for every $A&#62;0$?{/LI} {/UL}




A problem of Erd\H{o}s and Szekeres. It is easy to see that $f(n)\leq n/P(n)$ for composite $n$, since if $j=p^k$ where $p^k\mid n$ and $p^{k+1}\nmid n$ then $\textrm{gcd}\left(n,\binom{n}{j}\right)=n/p^k$. This implies\[f(n) \leq (1+o(1))\frac{n}{\log n}.\]It is known that $f(n)=n/P(n)$ when $n$ is the product of two primes. Another example is $n=30$. For the second problem, it is easy to see that for any $n$ we have $f(n)\geq p(n)$, where $p(n)$ is the smallest prime dividing $n$, and hence there are infinitely many $n$ (those $=p^2)$ such that $f(n)\geq n^{1/2}$.

\noindent\rule{\linewidth}{0.4pt}


%% ============================================================
\section{Problem \#701}
\label{prob:701}

\noindent\textbf{Status:} OPEN \quad|\quad \textbf{Prize:} none \quad|\quad \textbf{Updated:} 2025-08-31

\noindent\textbf{Tags:} combinatorics, intersecting family

\medskip

\noindent\textbf{Statement:}

Let $\mathcal{F}$ be a family of sets closed under taking subsets (i.e. if $B\subseteq A\in\mathcal{F}$ then $B\in \mathcal{F}$). There exists some element $x$ such that whenever $\mathcal{F}'\subseteq \mathcal{F}$ is an intersecting subfamily we have\[\lvert \mathcal{F}'\rvert \leq \lvert \{ A\in \mathcal{F} : x\in A\}\rvert.\]




A problem of Chv\'{a}tal \cite{Ch74}, who proved it replacing the closed under subsets condition with the (stronger) condition that, assuming all sets in $\mathcal{F}$ are subsets of $\{1,\ldots,n\}$, whenever $A\in \mathcal{F}$ and there is an injection $f:B\to A$ such that $x\leq f(x)$ for all $x\in B$, then $B\in \mathcal{F}$. Sterboul \cite{St74} proved this when, letting $\mathcal{G}$ be the maximal sets (under inclusion) in $\mathcal{F}$, all sets in $\mathcal{G}$ have the same size, $\lvert A\cap B\rvert\leq 1$ for all $A\neq B\in \mathcal{G}$, and at least two sets in $\mathcal{G}$ have non-empty intersection. Frankl and Kupavskii \cite{FrKu23} have proved this when $\mathcal{F}$ has covering number $2$. Borg \cite{Bo11} has proposed a weighted generalisation of this conjecture, which he proves under certain additional assumptions.




References

[Bo11] Borg, Peter, On Chv\'{a}tal's conjecture and a conjecture on families of signed sets . European J. Combin. (2011), 140-145.

[Ch74] Chv\'{a}tal, V., Intersecting families of edges in hypergraphs having the hereditary property . (1974), 61-66.

[FrKu23] Frankl, Peter and Kupavskii, Andrey, Perfect matchings in down-sets . Discrete Math. (2023), Paper No. 113323, 7.

[St74] Sterboul, F., Sur une conjecture de V. Chv\'{a}tal . (1974), 152-164.

\noindent\rule{\linewidth}{0.4pt}


%% ============================================================
\section{Problem \#704}
\label{prob:704}

\noindent\textbf{Status:} OPEN \quad|\quad \textbf{Prize:} none \quad|\quad \textbf{Updated:} 2025-08-31

\noindent\textbf{Tags:} graph theory, geometry, chromatic number

\medskip

\noindent\textbf{Statement:}

Let $G_n$ be the unit distance graph in $\mathbb{R}^n$, with two vertices joined by an edge if and only if the distance between them is $1$. Estimate the chromatic number $\chi(G_n)$. Does it grow exponentially in $n$? Does\[\lim_{n\to \infty}\chi(G_n)^{1/n}\]exist?




A generalisation of the Hadwiger-Nelson problem (which addresses $n=2$). Frankl and Wilson \cite{FrWi81} proved exponential growth:\[\chi(G_n) \geq (1+o(1))1.2^n.\]This was improved to\[\chi(G_n)\geq (1.239\cdots+o(1))^n\]by Raigorodsky \cite{Ra00}. The trivial colouring (by tiling with cubes) gives\[\chi(G_n) \leq (2+\sqrt{n})^n.\]Larman and Rogers \cite{LaRo72} improved this to\[\chi(G_n) \leq (3+o(1))^n,\]and conjecture the truth may be $(2^{3/2}+o(1))^n$. (Note that $2^{3/2}\approx 2.828$.) Prosanov \cite{Pr20} has given an alternative proof of this upper bound. See also [508] , [705] , and [706] .




References

[FrWi81] Frankl, P. and Wilson, R. M., Intersection theorems with geometric consequences . Combinatorica (1981), 357-368.

[LaRo72] Larman, D. G. and Rogers, C. A., The realization of distances within sets in Euclidean space . Mathematika (1972), 1-24.

[Pr20] Prosanov, Roman, A new proof of the Larman-Rogers upper bound for the chromatic number of the Euclidean space . Discrete Appl. Math. (2020), 115-120.

[Ra00] Ra\u igorodski\u i, A. M., On the chromatic number of a space . Uspekhi Mat. Nauk (2000), 147--148.

\noindent\rule{\linewidth}{0.4pt}


%% ============================================================
\section{Problem \#706}
\label{prob:706}

\noindent\textbf{Status:} OPEN \quad|\quad \textbf{Prize:} none \quad|\quad \textbf{Updated:} 2025-08-31

\noindent\textbf{Tags:} graph theory, chromatic number

\medskip

\noindent\textbf{Statement:}

Let $L(r)$ be such that if $G$ is a graph formed by taking a finite set of points $P$ in $\mathbb{R}^2$ and some set $A\subset (0,\infty)$ of size $r$, where the vertex set is $P$ and there is an edge between two points if and only if their distance is a member of $A$, then $\chi(G)\leq L(r)$. Estimate $L(r)$. In particular, is it true that $L(r)\leq r^{O(1)}$?




The case $r=1$ is the Hadwiger-Nelson problem , for which it is known that $5\leq L(1)\leq 7$. See also [508] , [704] , and [705] .

\noindent\rule{\linewidth}{0.4pt}


%% ============================================================
\section{Problem \#708}
\label{prob:708}

\noindent\textbf{Status:} OPEN \quad|\quad \textbf{Prize:} \$100 \quad|\quad \textbf{Updated:} 2025-08-31

\noindent\textbf{Tags:} number theory

\medskip

\noindent\textbf{Statement:}

Let $g(n)$ be minimal such that for any $A\subseteq [2,\infty)\cap \mathbb{N}$ with $\lvert A\rvert =n$ and any set $I$ of $\max(A)$ consecutive integers there exists some $B\subseteq I$ with $\lvert B\rvert=g(n)$ such that\[\prod_{a\in A} a \mid \prod_{b\in B}b.\]Is it true that\[g(n) \leq (2+o(1))n?\]Or perhaps even $g(n)\leq 2n$?




A problem of Erd\H{o}s and Sur\'{a}nyi \cite{ErSu59}, who proved that $g(n) \geq (2-o(1))n$, and that $g(3)=4$. Their lower bound construction takes $A$ as the set of $p_ip_j$ for $i\neq j$, where $p_1&lt;\cdots &lt;p_\ell$ is some set of primes such that $2p_1^2&gt;p_\ell^2$. Gallai was the first to consider problems of this type, and observed that $g(2)=2$ and $g(3)\geq 4$. In \cite{Er92c} Erd\H{o}s offers '100 dollars or 1000 rupees', whichever is more, for a proof or disproof. (In 1992 1000 rupees was worth approximately \$38.60.) Erd\H{o}s and Sur\'{a}nyi similarly asked what is the smallest $c_n\geq 1$ such that in any interval $I\subset [0,\infty)$ of length $c_n\max(A)$ there exists some $B\subseteq I\cap \mathbb{N}$ with $\lvert B\rvert=n$ such that\[\prod_{a\in A} a \mid \prod_{b\in B}b.\]They prove $c_2=1$ and $c_3=\sqrt{2}$, but have no good upper or lower bounds in general. See also [709] .




References

[Er92c] Erd\H{o}s, P., Some of my forgotten problems in number theory . Hardy-Ramanujan J. (1992), 34-50.

[ErSu59] Erd\H{o}s, P\'{a}l and Sur\'{a}nyi, J\'{a}nos, Bemerkungen zu einer Aufgabe eines mathematischen {W}ettbewerbs . Mat. Lapok (1959), 39-48.

\noindent\rule{\linewidth}{0.4pt}


%% ============================================================
\section{Problem \#709}
\label{prob:709}

\noindent\textbf{Status:} OPEN \quad|\quad \textbf{Prize:} none \quad|\quad \textbf{Updated:} 2025-08-31

\noindent\textbf{Tags:} number theory

\medskip

\noindent\textbf{Statement:}

Let $f(n)$ be minimal such that, for any $A=\{a_1,\ldots,a_n\}\subseteq [2,\infty)\cap\mathbb{N}$ of size $n$, in any interval $I$ of $f(n)\max(A)$ consecutive integers there exist distinct $x_1,\ldots,x_n\in I$ such that $a_i\mid x_i$. Obtain good bounds for $f(n)$, or even an asymptotic formula.




A problem of Erd\H{o}s and Sur\'{a}nyi \cite{ErSu59}, who proved\[(\log n)^c \ll f(n) \ll n^{1/2}\]for some constant $c&gt;0$. The lower bound here can be improved, using van Doorn's lower bound for [711] from \cite{vD26}, to\[\frac{\log n}{\log\log n}\ll f(n).\]See also [708] .




References

[ErSu59] Erd\H{o}s, P\'{a}l and Sur\'{a}nyi, J\'{a}nos, Bemerkungen zu einer Aufgabe eines mathematischen {W}ettbewerbs . Mat. Lapok (1959), 39-48.

[vD26] W. van Doorn, On the length of an interval that contains distinct multiples of the first $n$ positive integers . Integers (2026), #A7.

\noindent\rule{\linewidth}{0.4pt}


%% ============================================================
\section{Problem \#710}
\label{prob:710}

\noindent\textbf{Status:} OPEN \quad|\quad \textbf{Prize:} ₹2000 \quad|\quad \textbf{Updated:} 2025-08-31

\noindent\textbf{Tags:} number theory

\medskip

\noindent\textbf{Statement:}

Let $f(n)$ be minimal such that in $(n,n+f(n))$ there exist distinct integers $a_1,\ldots,a_n$ such that $k\mid a_k$ for all $1\leq k\leq n$. Obtain an asymptotic formula for $f(n)$.




A problem of Erd\H{o}s and Pomerance \cite{ErPo80}, who proved\[(2/\sqrt{e}+o(1))n\left(\frac{\log n}{\log\log n}\right)^{1/2}\leq f(n)\leq (1.7398\cdots+o(1))n(\log n)^{1/2}.\]In \cite{Er92c} Erd\H{o}s offered 2000 rupees for an asymptotic formula; for uniform comparison across prizes I have converted this using the 1992 exchange rates. See also [711] .




References

[Er92c] Erd\H{o}s, P., Some of my forgotten problems in number theory . Hardy-Ramanujan J. (1992), 34-50.

[ErPo80] P. Erd\H{o}s and C. Pomerance, Matching the natural numbers up to $n$ with distinct multiples of another interval . Indigationes Math. (1980), 147-151.

\noindent\rule{\linewidth}{0.4pt}


%% ============================================================
\section{Problem \#711}
\label{prob:711}

\noindent\textbf{Status:} OPEN \quad|\quad \textbf{Prize:} ₹1000 \quad|\quad \textbf{Updated:} 2025-08-31

\noindent\textbf{Tags:} number theory

\medskip

\noindent\textbf{Statement:}

Let $f(n,m)$ be minimal such that in $(m,m+f(n,m))$ there exist distinct integers $a_1,\ldots,a_n$ such that $k\mid a_k$ for all $1\leq k\leq n$. Prove that\[\max_m f(n,m) \leq n^{1+o(1)}\]and that\[\max_m (f(n,m)-f(n,n))\to \infty.\]




A problem of Erd\H{o}s and Pomerance \cite{ErPo80}, who proved that\[\max_m f(n,m) \ll n^{3/2}\]and\[n\left(\frac{\log n}{\log\log n}\right)^{1/2} \ll f(n,n)\ll n(\log n)^{1/2}.\]In \cite{Er92c} Erd\H{o}s offered 1000 rupees for a proof of either; for uniform comparison across prizes I have converted this using the 1992 exchange rates. van Doorn \cite{vD26} has provided an affirmative answer to the second question, proving that, for all large $n$, there exists $m=m(n)$ such that\[f(n,m)-f(n,n) \gg \frac{\log n}{\log\log n}n.\]See also [710] .




References

[Er92c] Erd\H{o}s, P., Some of my forgotten problems in number theory . Hardy-Ramanujan J. (1992), 34-50.

[ErPo80] P. Erd\H{o}s and C. Pomerance, Matching the natural numbers up to $n$ with distinct multiples of another interval . Indigationes Math. (1980), 147-151.

[vD26] W. van Doorn, On the length of an interval that contains distinct multiples of the first $n$ positive integers . Integers (2026), #A7.

\noindent\rule{\linewidth}{0.4pt}


%% ============================================================
\section{Problem \#712}
\label{prob:712}

\noindent\textbf{Status:} OPEN \quad|\quad \textbf{Prize:} \$500 \quad|\quad \textbf{Updated:} 2025-08-31

\noindent\textbf{Tags:} graph theory, turan number, hypergraphs

\medskip

\noindent\textbf{Statement:}

Determine, for any $k&#62;r&#62;2$, the value of\[\frac{\mathrm{ex}_r(n,K_k^r)}{\binom{n}{r}},\]where $\mathrm{ex}_r(n,K_k^r)$ is the largest number of $r$-edges which can placed on $n$ vertices so that there exists no set of $k$ vertices which is covered by all $\binom{k}{r}$ possible $r$-edges.




Tur\'{a n proved} that, when $r=2$, this limit is\[\frac{1}{2}\left(1-\frac{1}{k-1}\right).\]Erd\H{o}s \cite{Er81} offered \$500 for the determination of this value for any fixed $k&gt;r&gt;2$, and \$1000 for 'clearing up the whole set of problems'. See also [500] for the case $r=3$ and $k=4$.




References

[Er81] Erd\H{o}s, P., On the combinatorial problems which I would most like to see solved . Combinatorica (1981), 25-42.

\noindent\rule{\linewidth}{0.4pt}


%% ============================================================
\section{Problem \#713}
\label{prob:713}

\noindent\textbf{Status:} OPEN \quad|\quad \textbf{Prize:} \$500 \quad|\quad \textbf{Updated:} 2025-08-31

\noindent\textbf{Tags:} graph theory, turan number

\medskip

\noindent\textbf{Statement:}

Is it true that, for every bipartite graph $G$, there exists some $\alpha\in [1,2)$ and $c&#62;0$ such that\[\mathrm{ex}(n;G)\sim cn^\alpha?\]Must $\alpha$ be rational?




A problem of Erd\H{o}s and Simonovits. Erd\H{o}s sometimes asked this in the weaker version with just\[\mathrm{ex}(n;G)\asymp n^{\alpha}.\]Erd\H{o}s \cite{Er67d} had initially conjectured that, for any bipartite graph $G$, $\mathrm{ex}(n;G)\sim cn^{\alpha}$ for some constant $c&gt;0$ and $\alpha$ of the shape $1+\frac{1}{k}$ or $2-\frac{1}{k}$ for some integer $k\geq 2$. This was disproved by Erd\H{o}s and Simonovits \cite{ErSi70}. The analogous statement is not true for hypergraphs, as shown by Frankl and F\"{u}redi \cite{FrFu87}, who proved that if $G$ is the $5$-uniform hypergraph on $8$ vertices with edges $\{12346,12457,12358\}$ then $\mathrm{ex}(n;G)=o(n^5)$ but $\mathrm{ex}(n;G)\neq O(n^c)$ for any $c&lt;5$. A simplified proof was given by F\"{u}redi and Gerbner \cite{FuGe21}, who extended it to a counterexample for all $k\geq 5$. It remains open whether it is true for $k=3$ and $k=4$ (though F\"{u}redi and Gerbner conjecture it is not). See also [571] .




References

[Er67d] Erd\H{o}s, P., Some recent results on extremal problems in graph theory. {R}esults . (1967), 117--123 (English); pp. 124--130 (French).

[ErSi70] Erd\H{o}s, P. and Simonovits, M., Some extremal problems in graph theory . Combinatorial theory and its applications, I-III (Proc. Colloq., Balatonf\"{u}red, 1969) (1970), 377-390.

[FrFu87] Frankl, P. and F\"uredi, Z., Exact solution of some Tur\'{a}n-type problems . J. Combin. Theory Ser. A (1987), 226--262.

[FuGe21] F\"uredi, Zolt\'{a}n and Gerbner, D\'{a}niel, Hypergraphs without exponents . J. Combin. Theory Ser. A (2021), Paper No. 105517, 9.

\noindent\rule{\linewidth}{0.4pt}


%% ============================================================
\section{Problem \#714}
\label{prob:714}

\noindent\textbf{Status:} OPEN \quad|\quad \textbf{Prize:} none \quad|\quad \textbf{Updated:} 2025-08-31

\noindent\textbf{Tags:} graph theory, turan number

\medskip

\noindent\textbf{Statement:}

Is it true that\[\mathrm{ex}(n; K_{r,r}) \gg n^{2-1/r}?\]




K\"{o}v\'{a}ri, S\'{o}s, and Tur\'{a}n \cite{KST54} proved\[\mathrm{ex}(n; K_{r,r}) \ll n^{2-1/r}\]for all $r\geq 2$. Brown \cite{Br66} and, independently, Erd\H{o}s, R\'{e}nyi, and S\'{o}s \cite{ERS66}, proved the conjectured lower bound when $r=3$. When $r=2$ it is known that\[\mathrm{ex}(n;K_{2,2})=\left(\frac{1}{2}+o(1)\right)n^{3/2}\](see [768] , since $K_{2,2}=C_4$). See also [147] . The generalisation to hypergraphs is [1158] .




References

[Br66] Brown, W. G., On graphs that do not contain a Thomsen graph . Canad. Math. Bull. (1966), 281-285.

[ERS66] Erd\H{o}s, P. and R\'{e}nyi, A. and S\'os, V. T., On a problem of graph theory . Studia Sci. Math. Hungar. (1966), 215--235.

[KST54] K\"{o}vari, T. and S\'{o}s, V. T. and Tur\'{a}n, P., On a problem of K. Zarankiewicz . Colloq. Math. (1954), 50-57.

\noindent\rule{\linewidth}{0.4pt}


%% ============================================================
\section{Problem \#719}
\label{prob:719}

\noindent\textbf{Status:} OPEN \quad|\quad \textbf{Prize:} none \quad|\quad \textbf{Updated:} 2025-08-31

\noindent\textbf{Tags:} graph theory, hypergraphs

\medskip

\noindent\textbf{Statement:}

Let $\mathrm{ex}_r(n;K_{r+1}^r)$ be the maximum number of $r$-edges that can be placed on $n$ vertices without forming a $K_{r+1}^r$ (the $r$-uniform complete graph on $r+1$ vertices). Is every $r$-hypergraph $G$ on $n$ vertices the union of at most $\mathrm{ex}_{r}(n;K_{r+1}^r)$ many copies of $K_r^r$ and $K_{r+1}^r$, no two of which share a $K_r^r$?




A conjecture of Erd\H{o}s and Sauer.

\noindent\rule{\linewidth}{0.4pt}


%% ============================================================
\section{Problem \#724}
\label{prob:724}

\noindent\textbf{Status:} OPEN \quad|\quad \textbf{Prize:} none \quad|\quad \textbf{Updated:} 2025-08-31

\noindent\textbf{Tags:} combinatorics

\medskip

\noindent\textbf{Statement:}

Let $f(n)$ be the maximum number of mutually orthogonal Latin squares of order $n$. Is it true that\[f(n) \gg n^{1/2}?\]




Euler conjectured that $f(n)=1$ when $n\equiv 2\pmod{4}$, but this was disproved by Bose, Parker, and Shrikhande \cite{BPS60} who proved $f(n)\geq 2$ for $n\geq 7$. Chowla, Erd\H{o}s, and Straus \cite{CES60} proved $f(n) \gg n^{1/91}$. Wilson \cite{Wi74} proved $f(n) \gg n^{1/17}$. Beth \cite{Be83c} proved $f(n) \gg n^{1/14.8}$. The sequence of $f(n)$ is A001438 in the OEIS.




References

[BPS60] Bose, R. C. and Shrikhande, S. S. and Parker, E. T., Further results on the construction of mutually orthogonal Latin squares and the falsity of Euler's conjecture . Canadian J. Math. (1960), 189-203.

[Be83c] Beth, Thomas, Eine Bemerkung zur Absch\"{a}tzung der Anzahl orthogonaler lateinischer Quadrate mittels Siebverfahren . Abh. Math. Sem. Univ. Hamburg (1983), 284-288.

[CES60] Chowla, S. and Erd\H{o}s, P. and Straus, E. G., On the maximal number of pairwise orthogonal Latin squares of a given order . Canadian J. Math. (1960), 204-208.

[Wi74] Wilson, Richard M., Concerning the number of mutually orthogonal Latin squares . Discrete Math. (1974), 181-198.

\noindent\rule{\linewidth}{0.4pt}


%% ============================================================
\section{Problem \#725}
\label{prob:725}

\noindent\textbf{Status:} OPEN \quad|\quad \textbf{Prize:} none \quad|\quad \textbf{Updated:} 2025-08-31

\noindent\textbf{Tags:} combinatorics

\medskip

\noindent\textbf{Statement:}

Give an asymptotic formula for the number of $k\times n$ Latin rectangles .




Erd\H{o}s and Kaplansky \cite{ErKa46} proved the count is\[\sim e^{-\binom{k}{2}}(n!)^k\]when $k=o((\log n)^{3/2-\epsilon})$. Yamamoto \cite{Ya51} extended this to $k\leq n^{1/3-o(1)}$. The count of such Latin rectangles is A001009 in the OEIS.




References

[ErKa46] Erd\H{o}s, Paul and Kaplansky, Irving, The asymptotic number of Latin rectangles . Amer. J. Math. (1946), 230-236.

[Ya51] Yamamoto, Koichi, On the asymptotic number of Latin rectangles . Jpn. J. Math. (1951), 113-119.

\noindent\rule{\linewidth}{0.4pt}


%% ============================================================
\section{Problem \#726}
\label{prob:726}

\noindent\textbf{Status:} OPEN \quad|\quad \textbf{Prize:} none \quad|\quad \textbf{Updated:} 2025-08-31

\noindent\textbf{Tags:} number theory

\medskip

\noindent\textbf{Statement:}

As $n\to \infty$ ranges over integers\[\sum_{p\leq n}1_{n\in (p/2,p)\pmod{p}}\frac{1}{p}\sim \frac{\log\log n}{2}.\]




A conjecture of Erd\H{o}s, Graham, Ruzsa, and Straus \cite{EGRS75}. For comparison the classical estimate of Mertens states that\[\sum_{p\leq n}\frac{1}{p}\sim \log\log n.\]By $n\in (p/2,p)\pmod{p}$ we mean $n\equiv r\pmod{p}$ for some integer $r$ with $p/2&#60;r&#60;p$.




References

[EGRS75] Erd\H{o}s, P. and Graham, R. L. and Ruzsa, I. Z. and Straus, E. G., On the prime factors of $(\sp{2n}\sb{n})$ . Math. Comp. (1975), 83-92.

\noindent\rule{\linewidth}{0.4pt}


%% ============================================================
\section{Problem \#727}
\label{prob:727}

\noindent\textbf{Status:} OPEN \quad|\quad \textbf{Prize:} none \quad|\quad \textbf{Updated:} 2025-08-31

\noindent\textbf{Tags:} number theory, factorials

\medskip

\noindent\textbf{Statement:}

Let $k\geq 2$. Does\[(n+k)!^2 \mid (2n)!\]for infinitely many $n$?




A conjecture of Erd\H{o}s, Graham, Ruzsa, and Straus \cite{EGRS75}. It is open even for $k=2$. Balakran \cite{Ba29} proved this holds for $k=1$ - that is, $(n+1)^2\mid \binom{2n}{n}$ for infinitely many $n$. It is a classical fact that $(n+1)\mid \binom{2n}{n}$ for all $n$ (see Catalan numbers ). Erd\H{o}s, Graham, Ruzsa, and Straus observe that the method of Balakran can be further used to prove that there are infinitely many $n$ such that\[(n+k)!(n+1)! \mid (2n)!\](in fact this holds whenever $k&lt;c \log n$ for some small constant $c&gt;0$). Erd\H{o}s \cite{Er68c} proved that if $a!b!\mid n!$ then $a+b\leq n+O(\log n)$.




References

[Ba29] H. Balakran, On the values of $n$ which make $(2n)!/(n+1)!(n+1)!$ an integer . J. Indian Math. Soc. (1929), 97-100.

[EGRS75] Erd\H{o}s, P. and Graham, R. L. and Ruzsa, I. Z. and Straus, E. G., On the prime factors of $(\sp{2n}\sb{n})$ . Math. Comp. (1975), 83-92.

[Er68c] P. Erd\H{o}s, Aufgabe 557 . Elemente Math. (1968), 111-113.

\noindent\rule{\linewidth}{0.4pt}


%% ============================================================
\section{Problem \#730}
\label{prob:730}

\noindent\textbf{Status:} OPEN \quad|\quad \textbf{Prize:} none \quad|\quad \textbf{Updated:} 2025-08-31

\noindent\textbf{Tags:} number theory, binomial coefficients, base representations

\medskip

\noindent\textbf{Statement:}

Are there infinitely many pairs of integers $n\neq m$ such that $\binom{2n}{n}$ and $\binom{2m}{m}$ have the same set of prime divisors?




A problem of Erd\H{o}s, Graham, Ruzsa, and Straus \cite{EGRS75}, who believed there is 'no doubt' that the answer is yes. For example $(87,88)$ and $(607,608)$. Those $n$ such that there exists some suitable $m&gt;n$ are listed as A129515 in the OEIS. A triple of such $n$ for which $\binom{2n}{n}$ all share the same set of prime divisors is $(10003,10004,10005)$. It is not known whether there are such pairs of the shape $(n,n+k)$ for every $k\geq 1$.




References

[EGRS75] Erd\H{o}s, P. and Graham, R. L. and Ruzsa, I. Z. and Straus, E. G., On the prime factors of $(\sp{2n}\sb{n})$ . Math. Comp. (1975), 83-92.

\noindent\rule{\linewidth}{0.4pt}


%% ============================================================
\section{Problem \#731}
\label{prob:731}

\noindent\textbf{Status:} OPEN \quad|\quad \textbf{Prize:} none \quad|\quad \textbf{Updated:} 2025-08-31

\noindent\textbf{Tags:} number theory, binomial coefficients

\medskip

\noindent\textbf{Statement:}

Find some reasonable function $f(n)$ such that, for almost all integers $n$, the least integer $m$ such that $m\nmid \binom{2n}{n}$ satisfies\[m\sim f(n).\]




A problem of Erd\H{o}s, Graham, Ruzsa, and Straus \cite{EGRS75}, who say it is 'not hard to show that', for almost all $n$, the minimal such $m$ satisfies\[m=\exp((\log n)^{1/2+o(1)}).\]




References

[EGRS75] Erd\H{o}s, P. and Graham, R. L. and Ruzsa, I. Z. and Straus, E. G., On the prime factors of $(\sp{2n}\sb{n})$ . Math. Comp. (1975), 83-92.

\noindent\rule{\linewidth}{0.4pt}


%% ============================================================
\section{Problem \#734}
\label{prob:734}

\noindent\textbf{Status:} OPEN \quad|\quad \textbf{Prize:} none \quad|\quad \textbf{Updated:} 2025-08-31

\noindent\textbf{Tags:} combinatorics

\medskip

\noindent\textbf{Statement:}

Find, for all large $n$, a non-trivial pairwise balanced block design $A_1,\ldots,A_m\subseteq \{1,\ldots,n\}$ such that, for all $t$, there are $O(n^{1/2})$ many $i$ such that $\lvert A_i\rvert=t$.




$A_1,\ldots,A_m$ is a pairwise balanced block design if every pair in $\{1,\ldots,n\}$ is contained in exactly one of the $A_i$. Erd\H{o}s \cite{Er81} writes 'this will be probably not be very difficult to prove but so far I was not successful'. Erd\H{o}s and de Bruijn \cite{dBEr48} proved that if $A_1,\ldots,A_m\subseteq \{1,\ldots,n\}$ is a pairwise balanced block design then $m\geq n$, and this implies there must be some $t$ such that there are $\gg n^{1/2}$ many $t$ with $\lvert A_i\rvert=t$.




References

[Er81] Erd\H{o}s, P., On the combinatorial problems which I would most like to see solved . Combinatorica (1981), 25-42.

[dBEr48] de Bruijn, N. G. and Erd\H{o}s, P., On a combinatorial problem . Nederl. Akad. Wetensch., Proc. (1948), 1277--1279 = Indagationes Math. 10, 421--423.

\noindent\rule{\linewidth}{0.4pt}


%% ============================================================
\section{Problem \#738}
\label{prob:738}

\noindent\textbf{Status:} OPEN \quad|\quad \textbf{Prize:} none \quad|\quad \textbf{Updated:} 2025-08-31

\noindent\textbf{Tags:} graph theory, chromatic number

\medskip

\noindent\textbf{Statement:}

If $G$ has infinite chromatic number and is triangle-free (contains no $K_3$) then must $G$ contain every tree as an induced subgraph?




A conjecture of Gy\'{a}rf\'{a}s.

\noindent\rule{\linewidth}{0.4pt}


%% ============================================================
\section{Problem \#740}
\label{prob:740}

\noindent\textbf{Status:} OPEN \quad|\quad \textbf{Prize:} none \quad|\quad \textbf{Updated:} 2025-08-31

\noindent\textbf{Tags:} graph theory, chromatic number

\medskip

\noindent\textbf{Statement:}

Let $\mathfrak{m}$ be an infinite cardinal and $G$ be a graph with chromatic number $\mathfrak{m}$. Let $r\geq 1$. Must $G$ contain a subgraph of chromatic number $\mathfrak{m}$ which does not contain any odd cycle of length $\leq r$?




A question of Erd\H{o}s and Hajnal. R\"{o}dl proved this is true if $\mathfrak{m}=\aleph_0$ and $r=3$ (see [108] for the finitary version). More generally, Erd\H{o}s and Hajnal asked must there exist (for every cardinal $\mathfrak{m}$ and integer $r$) some $f_r(\mathfrak{m})$ such that every graph with chromatic number $\geq f_r(\mathfrak{m})$ contains a subgraph with chromatic number $\mathfrak{m}$ with no odd cycle of length $\leq r$? Erd\H{o}s \cite{Er95d} claimed that even the $r=3$ case of this is open: must every graph with sufficiently large chromatic number contain a triangle free graph with chromatic number $\mathfrak{m}$? In \cite{Er81} Erd\H{o}s also asks the same question but with girth (i.e. the subgraph does not contain any cycle at all of length $\leq C$).




References

[Er81] Erd\H{o}s, P., On the combinatorial problems which I would most like to see solved . Combinatorica (1981), 25-42.

[Er95d] Erd\H{o}s, Paul, On some problems in combinatorial set theory . Publ. Inst. Math. (Beograd) (N.S.) (1995), 61-65.

\noindent\rule{\linewidth}{0.4pt}


%% ============================================================
\section{Problem \#749}
\label{prob:749}

\noindent\textbf{Status:} OPEN \quad|\quad \textbf{Prize:} none \quad|\quad \textbf{Updated:} 2025-08-31

\noindent\textbf{Tags:} additive combinatorics

\medskip

\noindent\textbf{Statement:}

Let $\epsilon&#62;0$. Does there exist $A\subseteq \mathbb{N}$ such that the lower density of $A+A$ is at least $1-\epsilon$ and yet $1_A\ast 1_A(n) \ll_\epsilon 1$ for all $n$?




It is unclear from \cite{Er94b} whether Erd\H{o}s believed the answer to this should be yes or no. He also asked a similar question for upper density, for which he seemed more certain the answer should be yes. He wrote that both he and S\'{a}rk\"{o}zy believed that if $1_A\ast 1-A(n)\leq C$ for all $n$ then the upper density of $A+A$ is at most $1-\epsilon_C$, where $\epsilon_C&gt;0$ may depend on $C$. The upper density variant has been resolved by Bhalla , with the assistance of GPT 5.4. Bhalla constructs, for any $\epsilon&gt;0$, a set $A$ for which the upper density of $A+A$ is at least $1-\epsilon$ and yet $1_A\ast 1_A(n)\ll \epsilon^{-1}$ for all $n$. See also [28] .




References

[Er94b] Erd\H{o}s, Paul, Some problems in number theory, combinatorics and combinatorial geometry . Math. Pannon. (1994), 261-269.

\noindent\rule{\linewidth}{0.4pt}


%% ============================================================
\section{Problem \#750}
\label{prob:750}

\noindent\textbf{Status:} OPEN \quad|\quad \textbf{Prize:} none \quad|\quad \textbf{Updated:} 2025-08-31

\noindent\textbf{Tags:} graph theory, chromatic number

\medskip

\noindent\textbf{Statement:}

Let $f(m)$ be some function such that $f(m)\to \infty$ as $m\to \infty$. Does there exist a graph $G$ of infinite chromatic number such that every subgraph on $m$ vertices contains an independent set of size at least $\frac{m}{2}-f(m)$?




In \cite{Er69b} Erd\H{o}s conjectures this for $f(m)=\epsilon m$ for any fixed $\epsilon&gt;0$. This follows from a result of Erd\H{o}s, Hajnal, and Szemer\'{e}di \cite{EHS82}, as described by msellke in the comments. In \cite{ErHa67b} Erd\H{o}s and Hajnal prove this for $f(m)\geq cm$ for all $c&gt;1/4$. See also [75] .




References

[EHS82] Erd\H{o}s, P. and Hajnal, A. and Szemer\'{e}di, E., On almost bipartite large chromatic graphs . Theory and practice of combinatorics (1982), 117-123.

[Er69b] Erd\H{o}s, P., Problems and results in chromatic graph theory . Proof Techniques in Graph Theory (Proc. Second Ann Arbor Graph Theory Conf., Ann Arbor, Mich., 1968) (1969), 27-35.

[ErHa67b] Erd\H{o}s, P. and Hajnal, Andr\'{a}s, On chromatic graphs . Mat. Lapok (1967), 1--4.

\noindent\rule{\linewidth}{0.4pt}


%% ============================================================
\section{Problem \#757}
\label{prob:757}

\noindent\textbf{Status:} OPEN \quad|\quad \textbf{Prize:} none \quad|\quad \textbf{Updated:} 2025-08-31

\noindent\textbf{Tags:} geometry, distances, sidon sets

\medskip

\noindent\textbf{Statement:}

Let $A\subset \mathbb{R}$ be a set of size $n$ such that every subset $B\subseteq A$ with $\lvert B\rvert =4$ has $\lvert B-B\rvert\geq 11$. Find the best constant $c&#62;0$ such that $A$ must always contain a Sidon set of size $\geq cn$.




For comparison, note that if $B$ were a Sidon set then $\lvert B-B\rvert=13$, so this condition is saying that at most one difference is 'missing' from $B-B$. Equivalently, one can view $A$ as a set such that every four points determine at least five distinct distances, and ask for a subset with all distances distinct. Without loss of generality, one can assume $A\subset \mathbb{N}$. Erd\H{o}s and S\'{o}s proved that $c\geq 1/2$. Gy\'{a}rf\'{a}s and Lehel \cite{GyLe95} proved\[\frac{1}{2}+\frac{1}{141\cdot 76}\leq c\leq\frac{3}{5}.\](The example proving the upper bound is the set of the first $n$ Fibonacci numbers.) Ma and Tang \cite{MaTa26} have improved these bounds to\[\frac{9}{17}\leq c\leq \frac{4}{7}.\]




References

[GyLe95] Gy\'{a}rf\'{a}s, Andr\'{a}s and Lehel, Jen\H{o}, Linear sets with five distinct differences among any four elements . J. Combin. Theory Ser. B (1995), 108-118.

[MaTa26] J. Ma and Q. Tang, Largest Sidon subsets in weak Sidon sets . arXiv:2602.23282 (2026).

\noindent\rule{\linewidth}{0.4pt}


%% ============================================================
\section{Problem \#761}
\label{prob:761}

\noindent\textbf{Status:} OPEN \quad|\quad \textbf{Prize:} none \quad|\quad \textbf{Updated:} 2025-08-31

\noindent\textbf{Tags:} graph theory, chromatic number

\medskip

\noindent\textbf{Statement:}

The cochromatic number of $G$, denoted by $\zeta(G)$, is the minimum number of colours needed to colour the vertices of $G$ such that each colour class induces either a complete graph or empty graph. The dichromatic number of $G$, denoted by $\delta(G)$, is the minimum number $k$ of colours required such that, in any orientation of the edges of $G$, there is a $k$-colouring of the vertices of $G$ such that there are no monochromatic oriented cycles. Must a graph with large chromatic number have large dichromatic number? Must a graph with large cochromatic number contain a graph with large dichromatic number?




The first question is due to Erd\H{o}s and Neumann-Lara. The second question is due to Erd\H{o}s and Gimbel. A positive answer to the second question implies a positive answer to the first via the bound mentioned in [760] .

\noindent\rule{\linewidth}{0.4pt}


%% ============================================================
\section{Problem \#766}
\label{prob:766}

\noindent\textbf{Status:} OPEN \quad|\quad \textbf{Prize:} none \quad|\quad \textbf{Updated:} 2025-08-31

\noindent\textbf{Tags:} graph theory, turan number

\medskip

\noindent\textbf{Statement:}

Let $f(n;k,l)=\min \mathrm{ex}(n;G)$, where $G$ ranges over all graphs with $k$ vertices and $l$ edges. Give good estimates for $f(n;k,l)$ in the range $k&#60;l\leq k^2/4$. For fixed $k$ and large $n$ is $f(n;k,l)$ a strictly monotone function of $l$?




Dirac and Erd\H{o}s proved independently that when $l=\lfloor k^2/4\rfloor+1$\[f(n;k,l)\leq \lfloor n^2/4\rfloor+1.\]

\noindent\rule{\linewidth}{0.4pt}


%% ============================================================
\section{Problem \#768}
\label{prob:768}

\noindent\textbf{Status:} OPEN \quad|\quad \textbf{Prize:} none \quad|\quad \textbf{Updated:} 2025-08-31

\noindent\textbf{Tags:} number theory

\medskip

\noindent\textbf{Statement:}

Let $A\subset\mathbb{N}$ be the set of $n$ such that for every prime $p\mid n$ there exists some $d\mid n$ with $d&#62;1$ such that $d\equiv 1\pmod{p}$. Is it true that there exists some constant $c&#62;0$ such that for all large $N$\[\frac{\lvert A\cap [1,N]\rvert}{N}=\exp(-(c+o(1))\sqrt{\log N}\log\log N).\]




Erd\H{o}s could prove that there exists some constant $c&#62;0$ such that for all large $N$\[\exp(-c\sqrt{\log N}\log\log N)\leq \frac{\lvert A\cap [1,N]\rvert}{N}\]and\[\frac{\lvert A\cap [1,N]\rvert}{N}\leq \exp(-(1+o(1))\sqrt{\log N\log\log N}).\]Erd\H{o}s asked about this because $\lvert A\cap [1,N]\rvert$ provides an upper bound for the number of integers $n\leq N$ for which there is a non-cyclic simple group of order $n$.

\noindent\rule{\linewidth}{0.4pt}


%% ============================================================
\section{Problem \#769}
\label{prob:769}

\noindent\textbf{Status:} OPEN \quad|\quad \textbf{Prize:} none \quad|\quad \textbf{Updated:} 2025-08-31

\noindent\textbf{Tags:} number theory, geometry

\medskip

\noindent\textbf{Statement:}

Let $c(n)$ be minimal such that if $k\geq c(n)$ then the $n$-dimensional unit cube can be decomposed into $k$ homothetic $n$-dimensional cubes. Give good bounds for $c(n)$ - in particular, is it true that $c(n) \gg n^n$?




A problem first investigated by Hadwiger, who proved the lower bound\[c(n) \geq 2^n+2^{n-1}.\]It is easy to see that $c(2)=6$. Meier conjectured $c(3)=48$. Burgess and Erd\H{o}s \cite{Er74b} proved\[c(n) \ll n^{n+1}.\]Erd\H{o}s wrote 'I am certain that if $n+1$ is a prime then $c(n)&#62;n^n$.' Hudelson \cite{Hu98} proved that if $(2^n-1,3^n-1)=1$ then $c(n) &#60; 6^n$, and in general $c(n) \ll (2n)^{n-1}$. Connor and Marmorino \cite{CoMa18} proved that\[c(n) \geq 2^{n+1}-1\]for all $n\geq 3$,\[c(n) \leq 1.8n^{n+1}\]if $n+1$ is prime, and\[c(n) \leq e^2n^n\]otherwise.




References

[CoMa18] Connor, Peter and Marmorino, Phillip, Decomposing cubes into smaller cubes . J. Geom. (2018), Paper No. 19, 11.

[Er74b] Erd\H{o}s, P., Remarks on some problems in number theory . Math. Balkanica (1974), 197-202.

[Hu98] Hudelson, Matthew, Dissecting {$d$}-cubes into smaller {$d$}-cubes . J. Combin. Theory Ser. A (1998), 190--200.

\noindent\rule{\linewidth}{0.4pt}


%% ============================================================
\section{Problem \#770}
\label{prob:770}

\noindent\textbf{Status:} OPEN \quad|\quad \textbf{Prize:} none \quad|\quad \textbf{Updated:} 2025-08-31

\noindent\textbf{Tags:} number theory

\medskip

\noindent\textbf{Statement:}

Let $h(n)$ be minimal such that $2^n-1,3^n-1,\ldots,h(n)^n-1$ are mutually coprime. Does, for every prime $p$, the density $\delta_p$ of integers with $h(n)=p$ exist? Does $\liminf h(n)=\infty$? Is it true that if $p$ is the greatest prime such that $p-1\mid n$ and $p&#62;n^\epsilon$ then $h(n)=p$?




It is easy to see that $h(n)=n+1$ if and only if $n+1$ is prime, and that $h(n)$ is unbounded for odd $n$. It is probably true that $h(n)=3$ for infinitely many $n$. See also [820] .

\noindent\rule{\linewidth}{0.4pt}


%% ============================================================
\section{Problem \#773}
\label{prob:773}

\noindent\textbf{Status:} OPEN \quad|\quad \textbf{Prize:} none \quad|\quad \textbf{Updated:} 2025-08-31

\noindent\textbf{Tags:} number theory, sidon sets, squares

\medskip

\noindent\textbf{Statement:}

What is the size of the largest Sidon subset $A\subseteq\{1,2^2,\ldots,N^2\}$? Is it $N^{1-o(1)}$?




A question of Alon and Erd\H{o}s \cite{AlEr85}, who proved $\lvert A\rvert \geq N^{2/3-o(1)}$ is possible (via a random subset), and observed that\[\lvert A\rvert \ll \frac{N}{(\log N)^{1/4}},\]since (as shown by Landau) the density of the sums of two squares decays like $(\log N)^{-1/2}$. This upper bound was improved to\[\lvert A\rvert \leq N^{1-\frac{c}{\log\log N}}\]for some constant $c&gt;0$ by Croot, Mao, and Yip (see the comments). The lower bound was improved to\[\lvert A\rvert \gg N^{2/3}\]by Lefmann and Thiele \cite{LeTh95}. In \cite{Er80} Erd\H{o}s further defines $g(A)$ to be the maximal size of a Sidon set of $\{a^2 : a\in A\}$, and asks whether $g(A)\geq g(\{1,\ldots,N\})$ where $N=\lvert A\rvert$. (See [530] for the linear analogue.) He further asks the infinite analogue: is there an infinite set $A\subset \mathbb{N}$ such that $\lvert A\cap [1,N]\rvert \geq N^{1-o(1)}$ for all large $N$ and $\{a^2 : a \in A\}$ is a Sidon set? See [1206] for a similar question for higher powers.




References

[AlEr85] Alon, Noga and Erd\H{o}s, P., An application of graph theory to additive number theory . European J. Combin. (1985), 201-203.

[Er80] Erd\H{o}s, Paul, A survey of problems in combinatorial number theory . Ann. Discrete Math. (1980), 89-115.

[LeTh95] Lefmann, Hanno and Thiele, Torsten, Point sets with distinct distances . Combinatorica (1995), 379--408.

\noindent\rule{\linewidth}{0.4pt}


%% ============================================================
\section{Problem \#774}
\label{prob:774}

\noindent\textbf{Status:} OPEN \quad|\quad \textbf{Prize:} none \quad|\quad \textbf{Updated:} 2025-08-31

\noindent\textbf{Tags:} number theory

\medskip

\noindent\textbf{Statement:}

We call $A\subset \mathbb{N}$ dissociated if $\sum_{n\in X}n\neq \sum_{m\in Y}m$ for all finite $X,Y\subset A$ with $X\neq Y$. Let $A\subset \mathbb{N}$ be an infinite set. We call $A$ proportionately dissociated if every finite $B\subset A$ contains a dissociated set of size $\gg \lvert B\rvert$. Is every proportionately dissociated set the union of a finite number of dissociated sets?




This question appears in a paper of Alon and Erd\H{o}s \cite{AlEr85}, although the general topic was first considered by Pisier \cite{Pi83}, who observed that the converse holds, and proved that being proportionately dissociated is equivalent to being a 'Sidon set' in the harmonic analysis sense; that is, whenever $f:A\to \mathbb{C}$ there exists some $\theta\in [0,1]$ such that\[\| f\|_1 \ll \left\lvert\sum_{n\in A} f(n)e(n\theta)\right\rvert,\]where $e(x)=e^{2\pi ix}$. Alon and Erd\H{o}s write that it 'seems unlikely that [this] is also sufficient'. They also point out the same question can be asked replacing dissociated with Sidon (in the additive combinatorial sense) (see [328] ). This latter question was resolved in the negative by Ne\v{s}et\v{r}il, R\"{o}dl, and Sales \cite{NRS24}.




References

[AlEr85] Alon, Noga and Erd\H{o}s, P., An application of graph theory to additive number theory . European J. Combin. (1985), 201-203.

[NRS24] Ne\v set\v ril, Jaroslav and R\"odl, Vojt\v ech and Sales, Marcelo, On {P}isier type theorems . Combinatorica (2024), 1211--1232.

[Pi83] Pisier, Gilles, Arithmetic characterizations of Sidon sets . Bull. Amer. Math. Soc. (N.S.) (1983), 87-89.

\noindent\rule{\linewidth}{0.4pt}


%% ============================================================
\section{Problem \#776}
\label{prob:776}

\noindent\textbf{Status:} OPEN \quad|\quad \textbf{Prize:} none \quad|\quad \textbf{Updated:} 2025-08-31

\noindent\textbf{Tags:} combinatorics

\medskip

\noindent\textbf{Statement:}

Let $r\geq 2$ and $A_1,\ldots,A_m\subseteq \{1,\ldots,n\}$ be such that $A_i\not\subseteq A_j$ for all $i\neq j$ and for any $t$ if there exists some $i$ with $\lvert A_i\rvert=t$ then there must exist at least $r$ sets of that size. How large must $n$ be (as a function of $r$) to ensure that there is such a family which achieves $n-3$ distinct sizes of sets?




A problem of Erd\H{o}s and Trotter. For $r=1$ and $n&#62;3$ the maximum possible is $n-2$. For $r&#62;1$ and $n$ sufficiently large $n-3$ is achievable, but $n-2$ is never achievable. Let $n_0(r)$ be such that whenever $n&#62;n_0(r)$ there exists such a family with $n-3$ distinct set sizes. He and Tang \cite{HeTa26b} (with ChatGPT) have proved that $n_0(2)=3$, $n_0(3)=8$, and for all $r\geq 4$\[2r+2\leq n_0(r)\leq 2r+2\log_2r+O(\log\log r),\]and hence $n_0(r)\sim 2r$ as $r\to \infty$.




References

[HeTa26b] Y. He and Q. Tang, An Erd\H{o}s--Trotter problem on antichains with multiplicity r on each occurring level . arXiv:2602.09803 (2026).

\noindent\rule{\linewidth}{0.4pt}


%% ============================================================
\section{Problem \#778}
\label{prob:778}

\noindent\textbf{Status:} OPEN \quad|\quad \textbf{Prize:} none \quad|\quad \textbf{Updated:} 2025-08-31

\noindent\textbf{Tags:} graph theory

\medskip

\noindent\textbf{Statement:}

Alice and Bob play a game on the edges of $K_n$, alternating colouring edges by red (Alice) and blue (Bob). Alice goes first, and wins if at the end the largest red clique is larger than any of the blue cliques. Does Bob have a winning strategy for $n\geq 3$? (Erd\H{o}s believed the answer is yes.) If we change the game so that Bob colours two edges after each edge that Alice colours, but now require Bob's largest clique to be strictly larger than Alice's, then does Bob have a winning strategy for $n&#62;3$? Finally, consider the game when Alice wins if the maximum degree of the red subgraph is larger than the maximum degree of the blue subgraph. Who wins?




Malekshahian and Spiro \cite{MaSp24} have proved that, for the first game, the set of $n$ for which Bob wins has density at least $3/4$ - in fact they prove that if Alice wins at $n$ then Bob wins at $n+1,n+2,n+3$. Similarly, for the third game they prove that the set of $n$ for which Bob wins has density at least $2/3$, and prove the stronger statement that if Alice wins at $n$ then Bob wins at $n+1,n+2$.




References

[MaSp24] Malekshahian, A. and Spiro, S., On a clique-building game of Erd\H{o}s . arXiv:2410.18304 (2024).

\noindent\rule{\linewidth}{0.4pt}


%% ============================================================
\section{Problem \#782}
\label{prob:782}

\noindent\textbf{Status:} OPEN \quad|\quad \textbf{Prize:} none \quad|\quad \textbf{Updated:} 2025-08-31

\noindent\textbf{Tags:} number theory

\medskip

\noindent\textbf{Statement:}

Do the squares contain arbitrarily long quasi-progressions? That is, does there exist some constant $C&#62;0$ such that, for any $k$, the squares contain a sequence $x_1,\ldots,x_k$ where, for some $d$ and all $1\leq i&#60;k$,\[x_i+d\leq x_{i+1}\leq x_i+d+C.\]Do the squares contain arbitrarily large cubes\[a+\left\{ \sum_i \epsilon_ib_i : \epsilon_i\in \{0,1\}\right\}?\]




A question of Brown, Erd\H{o}s, and Freedman \cite{BEF90}. It is a classical fact that the squares do not contain arithmetic progressions of length $4$. An affirmative answer to the first question implies an affirmative answer to the second. Solymosi \cite{So07} conjectured the answer to the second question is no. Cilleruelo and Granville \cite{CiGr07} have observed that the answer to the second question is no conditional on the Bombieri-Lang conjecture .




References

[BEF90] Brown, T. C. and Erd\H{o}s, P. and Freedman, A. R., Quasi-progressions and descending waves . J. Combin. Theory Ser. A (1990), 81-95.

[CiGr07] Cilleruelo, Javier and Granville, Andrew, Lattice points on circles, squares in arithmetic progressions and sumsets of squares . (2007), 241-262.

[So07] Solymosi, J\'{o}zsef, Elementary additive combinatorics . (2007), 29-38.

\noindent\rule{\linewidth}{0.4pt}


%% ============================================================
\section{Problem \#786}
\label{prob:786}

\noindent\textbf{Status:} OPEN \quad|\quad \textbf{Prize:} none \quad|\quad \textbf{Updated:} 2025-08-31

\noindent\textbf{Tags:} number theory

\medskip

\noindent\textbf{Statement:}

Let $\epsilon&#62;0$. Is there some set $A\subset \mathbb{N}$ of density $&#62;1-\epsilon$ such that $a_1\cdots a_r=b_1\cdots b_s$ with $a_i,b_j\in A$ can only hold when $r=s$? Similarly, can one always find a set $A\subset\{1,\ldots,N\}$ with this property of size $\geq (1-o(1))N$?




An example of such a set with density $1/4$ is given by the integers $\equiv 2\pmod{4}$. Selfridge constructed such a set with density $1/e-\epsilon$ for any $\epsilon&gt;0$: let $p_1&lt;\cdots&lt;p_k$ be a sequence of large consecutive primes such that\[\sum_{i=1}^k\frac{1}{p_i}&lt;1&lt;\sum_{i=1}^{k+1}\frac{1}{p_i},\]and let $A$ be those integers divisible by exactly one of $p_1,\ldots,p_k$. For the second question the set of integers with a prime factor $&gt;N^{1/2}$ give an example of a set with size $\geq (\log 2)N$. Erd\H{o}s could improve this constant slightly. Such an improvement is given by Tao in the comments, taking $A$ to be the set of all integers with exactly one prime factor $&gt;N^{\frac{1}{1+\sqrt{e}}}$, which has size $\approx 0.8285 N$. In \cite{Er65} Erd\H{o}s reports that Ruzsa proved the maximal size of such an $A$ is $\leq (1-c)N$ for some constant $c&gt;0$ for large $N$, but the proof 'is not yet published'. It is presumably the following. For any such $A$ we can define an associated additive function $f$ by letting $f(a_1\cdots a_r)=r$, and then $A=\{ n: f(n)=1\}$. A result of Erd\H{o}s, Ruzsa, and S\'{a}rkz\"{o}zy \cite{ERS73} then implies that there exists a constant $c&gt;0$ such that, for all sufficiently large $N$, $\lvert A\rvert\leq (1-c)N$. In \cite{ERS73} they indicate that their method should be able to obtain this with $c=1/10$. In the comments Tao explains how a result of Granville and Soundararajan \cite{GrSo01} implies\[\lvert A\rvert\leq (1-c+o(1))N\]for an explicit constant $c\approx 0.1715$, and this is the best possible (see also [121] ). Theorem 2 from \cite{ERS73} also implies a negative answer to the first question, in that the density is at most $1/2$. The above all assumes that repetition is allowed amongst the $a_i$ and $b_i$ in this condition. It is ambiguous in the original sources cited whether this was intended or not. If we insist the products are of distinct elements then this is a weaker condition on $A$, and the above questions are still open. The formulation in \cite{Er80} is clear that repetitions are not allowed, but also claims that Ruzsa had shown the answer to both questions is no in this case (and that Ruzsa had shown that the upper density of $A$ is $&lt;1/e$ if $A$ is infinite with this property); it is possible that Erd\H{o}s was mistakenly thinking of the above results, in which repetitions are allowed. See also [421] and [795] .




References

[ERS73] Erd\H{o}s, P. and Ruzsa, Jr., I. and S\'{a}rk\"ozi, A., On the number of solutions of {$f(n)=a$} for additive functions . Acta Arith. (1973), 1--9.

[Er65] Erd\H{o}s, P., Extremal problems in number theory . Proc. Sympos. Pure Math., Vol. VIII (1965), 181-189.

[Er80] Erd\H{o}s, Paul, A survey of problems in combinatorial number theory . Ann. Discrete Math. (1980), 89-115.

[GrSo01] Granville, Andrew and Soundararajan, K., The spectrum of multiplicative functions . Ann. of Math. (2) (2001), 407-470.

\noindent\rule{\linewidth}{0.4pt}


%% ============================================================
\section{Problem \#787}
\label{prob:787}

\noindent\textbf{Status:} OPEN \quad|\quad \textbf{Prize:} none \quad|\quad \textbf{Updated:} 2025-08-31

\noindent\textbf{Tags:} additive combinatorics

\medskip

\noindent\textbf{Statement:}

Let $g(n)$ be maximal such that given any set $A\subset \mathbb{R}$ with $\lvert A\rvert=n$ there exists some $B\subseteq A$ of size $\lvert B\rvert\geq g(n)$ such that $b_1+b_2\not\in A$ for all $b_1\neq b_2\in B$. Estimate $g(n)$.




This function was considered by Erd\H{o}s and Moser. Choi observed that, without loss of generality, one can assume that $A\subset \mathbb{Z}$. Klarner proved $g(n) \gg \log n$ (indeed, a greedy construction suffices). Choi \cite{Ch71} proved $g(n) \ll n^{2/5+o(1)}$. The current best bounds known are\[(\log n)^{1+c} \ll g(n) \ll \exp(\sqrt{\log n})\]for some constant $c&#62;0$, the lower bound due to Sanders \cite{Sa21} and the upper bound due to Ruzsa \cite{Ru05}. Beker \cite{Be25} has proved\[(\log n)^{1+\tfrac{1}{68}+o(1)} \ll g(n).\]




References

[Be25] A. Beker, The Erd\H{o}s-Moser sum-free set problem via improved bounds for $k$-configurations . arXiv:2501.10203 (2025).

[Ch71] Choi, S. L. G., On a combinatorial problem in number theory . Proc. London Math. Soc. (3) (1971), 629-642.

[Ru05] Ruzsa, Imre Z., Sum-avoiding subsets . Ramanujan J. (2005), 77-82.

[Sa21] Sanders, Tom, The Erd\H{o}s-Moser sum-free set problem . Canad. J. Math. (2021), 63-107.

\noindent\rule{\linewidth}{0.4pt}


%% ============================================================
\section{Problem \#788}
\label{prob:788}

\noindent\textbf{Status:} OPEN \quad|\quad \textbf{Prize:} none \quad|\quad \textbf{Updated:} 2025-08-31

\noindent\textbf{Tags:} additive combinatorics

\medskip

\noindent\textbf{Statement:}

Let $f(n)$ be maximal such that if $B\subset (2n,4n)\cap \mathbb{N}$ there exists some $C\subset (n,2n)\cap \mathbb{N}$ such that $c_1+c_2\not\in B$ for all $c_1\neq c_2\in C$ and $\lvert C\rvert+\lvert B\rvert \geq f(n)$. Estimate $f(n)$. In particular is it true that $f(n)\leq n^{1/2+o(1)}$?




A conjecture of Choi \cite{Ch71}, who proved $f(n) \ll n^{3/4}$. Adenwalla in the comments has provided a simple construction that proves $f(n) \gg n^{1/2}$. Hunter in the comments has sketched an argument that gives $f(n) \ll n^{2/3+o(1)}$. The bound\[f(n) \ll (n\log n)^{2/3}\]was proved by Baltz, Schoen, and Srivastav \cite{BSS00}. In general, the argument given by Hunter shows that, if a random Cayley graph with probability $p$ has (almost surely) independence number $\ll p^{-c-o(1)}$, then $f(n) \leq n^{\frac{c}{c+1}+o(1)}$. The work of Alon and Pham \cite{AlPh25} on random Cayley graphs therefore implies that\[f(n) \leq n^{3/5+o(1)},\]and the conjectured bound of $\leq p^{-1+o(1)}$ for the independence number would yield $f(n)\leq n^{1/2+o(1)}$.




References

[AlPh25] N. Alon and H. T. Pham, Random Cayley graphs and random subsets . arXiv:2509.02561 (2025).

[BSS00] Baltz, Andreas and Schoen, Tomasz and Srivastav, Anand, Probabilistic construction of small strongly sum-free sets via large Sidon sets . Colloq. Math. (2000), 171--176.

[Ch71] Choi, S. L. G., On a combinatorial problem in number theory . Proc. London Math. Soc. (3) (1971), 629-642.

\noindent\rule{\linewidth}{0.4pt}


%% ============================================================
\section{Problem \#789}
\label{prob:789}

\noindent\textbf{Status:} OPEN \quad|\quad \textbf{Prize:} none \quad|\quad \textbf{Updated:} 2025-08-31

\noindent\textbf{Tags:} additive combinatorics

\medskip

\noindent\textbf{Statement:}

Let $h(n)$ be maximal such that if $A\subseteq \mathbb{Z}$ with $\lvert A\rvert=n$ then there is $B\subseteq A$ with $\lvert B\rvert \geq h(n)$ such that if $a_1+\cdots+a_r=b_1+\cdots+b_s$ with $a_i,b_i\in B$ then $r=s$. Estimate $h(n)$.




Erd\H{o}s \cite{Er62c} proved $h(n) \ll n^{5/6}$. Straus \cite{St66} proved $h(n) \ll n^{1/2}$. Erd\H{o}s noted the bound $h(n)\gg n^{1/3}$, taking\[B=\{ a: \{ \alpha a\} \in n^{-1/3}+\tfrac{1}{2} (-n^{-2/3},n^{-2/3})\}\]for a random $\alpha\in [0,1]$. \cite{Er62c} and Choi \cite{Ch74b} improved this to $h(n) \gg (n\log n)^{1/3}$. See also [186] and [874] .




References

[Ch74b] Choi, S. L. G., On an extremal problem in number theory . J. Number Theory (1974), 105--111.

[Er62c] Erd\H{o}s, P\'{a}l, Some remarks on number theory. {III} . Mat. Lapok (1962), 28--38.

[St66] Straus, E. G., On a problem in combinatorial number theory . J. Math. Sci. (1966), 77--80.

\noindent\rule{\linewidth}{0.4pt}


%% ============================================================
\section{Problem \#790}
\label{prob:790}

\noindent\textbf{Status:} OPEN \quad|\quad \textbf{Prize:} none \quad|\quad \textbf{Updated:} 2025-08-31

\noindent\textbf{Tags:} additive combinatorics

\medskip

\noindent\textbf{Statement:}

Let $l(n)$ be maximal such that if $A\subset\mathbb{Z}$ with $\lvert A\rvert=n$ then there exists a sum-free $B\subseteq A$ with $\lvert B\rvert \geq l(n)$ - that is, $B$ is such that there are no solutions to\[a_1=a_2+\cdots+a_r\]with $a_i\in B$ all distinct. Estimate $l(n)$. In particular, is it true that $l(n)n^{-1/2}\to \infty$? Is it true that $l(n)&lt; n^{1-c}$ for some $c&gt;0$?




Erd\H{o}s observed that $l(n)\geq (n/2)^{1/2}$, which Choi improved to $l(n)&gt;(1+c)n^{1/2}$ for some $c&gt;0$. Erd\H{o}s \cite{Er73} thought he could prove $l(n)=o(n)$ but had 'difficulties in reconstructing [his] proof'. (In \cite{Er65} he wrote 'by complicated arguments we can show $l(n)=o(n)$'.) Choi, Koml\'{o}s, and Szemer\'{e}di \cite{CKS75} proved\[\left(\frac{\log n}{\log\log n}n\right)^{1/2}\ll l(n) \ll \frac{n}{\log n}.\]They further conjecture that $l(n)\geq n^{1-o(1)}$. See also [876] .




References

[CKS75] Choi, S. L. G. and Koml\'os, J. and Szemer\'{e}di, E., On sum-free subsequences . Trans. Amer. Math. Soc. (1975), 307--313.

[Er65] Erd\H{o}s, P., Extremal problems in number theory . Proc. Sympos. Pure Math., Vol. VIII (1965), 181-189.

[Er73] Erd\H{o}s, P., Problems and results on combinatorial number theory . A survey of combinatorial theory (Proc. Internat. Sympos., Colorado State Univ., Fort Collins, Colo., 1971) (1973), 117-138.

\noindent\rule{\linewidth}{0.4pt}


%% ============================================================
\section{Problem \#791}
\label{prob:791}

\noindent\textbf{Status:} OPEN \quad|\quad \textbf{Prize:} none \quad|\quad \textbf{Updated:} 2025-08-31

\noindent\textbf{Tags:} additive combinatorics

\medskip

\noindent\textbf{Statement:}

Let $g(n)$ be minimal such that there exists $A\subseteq \{0,\ldots,n\}$ of size $g(n)$ with $\{0,\ldots,n\}\subseteq A+A$. Estimate $g(n)$. In particular is it true that $g(n)\sim 2n^{1/2}$?




Such a set is often called a finite additive $2$-basis. A problem of Rohrbach, who proved in \cite{Ro37}\[(2+c)n \leq g(n)^2 \leq 4n\]for some small constant $c&#62;0$. The current best-known bounds are\[(2.181\cdots+o(1))n\leq g(n)^2 \leq (3.458\cdots+o(1))n.\]The lower bound is due to Yu \cite{Yu15}, and the upper bound is due to Kohonen \cite{Ko17}. (The disproof of $g(n)\sim 2n^{1/2}$ was accomplished by Mrose \cite{Mr79}, who gave a construction implying $g(n)^2 \leq \frac{7}{2}n$.)




References

[Ko17] Kohonen, Jukka, An improved lower bound for finite additive 2-bases . J. Number Theory (2017), 518--524.

[Mr79] Mrose, Arnulf, Untere Schranken f\"ur die {R}eichweiten von {E}xtremalbasen fester {O}rdnung . Abh. Math. Sem. Univ. Hamburg (1979), 118--124.

[Ro37] Rohrbach, Hans, Ein {B}eitrag zur additiven {Z}ahlentheorie . Math. Z. (1937), 1--30.

[Yu15] Yu, Gang, A new upper bound for finite additive {$h$}-bases . J. Number Theory (2015), 95--104.

\noindent\rule{\linewidth}{0.4pt}


%% ============================================================
\section{Problem \#792}
\label{prob:792}

\noindent\textbf{Status:} OPEN \quad|\quad \textbf{Prize:} none \quad|\quad \textbf{Updated:} 2025-08-31

\noindent\textbf{Tags:} additive combinatorics

\medskip

\noindent\textbf{Statement:}

Let $f(n)$ be maximal such that in any $A\subset \mathbb{Z}$ with $\lvert A\rvert=n$ there exists some sum-free subset $B\subseteq A$ with $\lvert B\rvert \geq f(n)$, so that there are no solutions to\[a+b=c\]with $a,b,c\in B$. Estimate $f(n)$.




Erd\H{o}s \cite{Er65} gave a simple proof that shows $f(n) \geq n/3$. Alon and Kleitman \cite{AlKl90} improved this to $f(n)\geq \frac{n+1}{3}$, and Bourgain \cite{Bo97} further improved this to $\frac{n+2}{3}$. The best lower bound known is\[f(n)\geq \frac{n}{3}+c\log\log n\]for some constant $c&gt;0$, due to Bedert \cite{Be25b}. The best upper bound known is\[f(n) \leq \frac{n}{3}+o(n),\]due to Eberhard, Green, and Manners \cite{EGM14}. This problem is Problem 1 on Green's open problems list .




References

[AlKl90] Alon, N. and Kleitman, D. J., Sum-free subsets . (1990), 13--26.

[Be25b] B. Bedert, Large sum-free subsets of sets of integers via $L^1$-estimates for trigonometric sums . arXiv:2502.08624 (2025).

[Bo97] Bourgain, Jean, Estimates related to sumfree subsets of sets of integers . Israel J. Math. (1997), 71-92.

[EGM14] Eberhard, Sean and Green, Ben and Manners, Freddie, Sets of integers with no large sum-free subset . Ann. of Math. (2) (2014), 621-652.

[Er65] Erd\H{o}s, P., Extremal problems in number theory . Proc. Sympos. Pure Math., Vol. VIII (1965), 181-189.

\noindent\rule{\linewidth}{0.4pt}


%% ============================================================
\section{Problem \#793}
\label{prob:793}

\noindent\textbf{Status:} OPEN \quad|\quad \textbf{Prize:} none \quad|\quad \textbf{Updated:} 2025-08-31

\noindent\textbf{Tags:} number theory

\medskip

\noindent\textbf{Statement:}

Let $F(n)$ be the maximum possible size of a subset $A\subseteq\{1,\ldots,n\}$ such that $a\nmid bc$ whenever $a,b,c\in A$ with $a\neq b$ and $a\neq c$. Is there a constant $C$ such that\[F(n)=\pi(n)+(C+o(1))n^{2/3}(\log n)^{-2}?\]




Erd\H{o}s \cite{Er38} proved there exist constants $0&lt;c_1\leq c_2$ such that\[\pi(n)+c_1n^{2/3}(\log n)^{-2}\leq F(n) \leq \pi(n)+c_2n^{2/3}(\log n)^{-2}.\]Erd\H{o}s \cite{Er69} gave a simple proof that $F(n) \leq \pi(n)+n^{2/3}$: define a graph with vertex set the union of those integers in $[1,n^{2/3}]$ with all primes $p\in (n^{2/3},n]$. We have an edge $u\sim v$ if and only if $uv\in A$. It is easy to see that every $m\leq n$ can be written as $uv$ where $u\leq n^{2/3}$ and $v$ is either prime or $\leq n^{2/3}$, and hence there are $\geq \lvert A\rvert$ many edges. This graph contains no path of length $3$ and hence must be a tree and have fewer edges than vertices, and we are done. This can be improved to give the upper bound mentioned by using a subset of integers in $[1,n^{2/3}]$. More generally, one can ask for such an asymptotic for the size of sets such that no $a\in A$ divides the product of $r$ distinct other elements of $A$, with the exponent $2/3$ replaced by $\frac{2}{r+1}$. For further discussion and references concerning this generalisation to $r\geq 3$ see the comment by Wouter van Doorn. See also [425] .




References

[Er38] P. Erd\H{o}s, On sequences of integers no one of which divides the product of two others and on related problems . Tomsk. Gos. Univ. Ucen Zap. (1938), 74-82.

[Er69] Erd\H{o}s, Paul, Some applications of graph theory to number theory . The Many Facets of Graph Theory (Proc. Conf., Western Mich. Univ., Kalamazoo, Mich., 1968) (1969), 77-82.

\noindent\rule{\linewidth}{0.4pt}


%% ============================================================
\section{Problem \#796}
\label{prob:796}

\noindent\textbf{Status:} OPEN \quad|\quad \textbf{Prize:} none \quad|\quad \textbf{Updated:} 2025-08-31

\noindent\textbf{Tags:} number theory

\medskip

\noindent\textbf{Statement:}

Let $k\geq 2$ and let $g_k(n)$ be the largest possible size of $A\subseteq \{1,\ldots,n\}$ such that every $m$ has $&#60;k$ solutions to $m=a_1a_2$ with $a_1&#60;a_2\in A$. Is it true that\[g_3(n)=\frac{\log\log n}{\log n}n+(c+o(1))\frac{n}{\log n}\]for some constant $c$?




Erd\H{o}s \cite{Er64d} proved that if $2^{r-1}&lt;k\leq 2^r$ then\[g_k(n) \sim \frac{(\log\log n)^{r-1}}{(r-1)!\log n}n\](which is the asymptotic count of those integers $\leq n$ with $r$ distinct prime factors). In particular the asymptotics of $g_k(n)$ are known; in \cite{Er69b} Erd\H{o}s discussed the second order terms, and this question is implicit (especially when compared to the explicit question [425] he asked on several occasions). In \cite{Er69} this question actually appears with a denominator of $(\log n)^2$ in the second term. Furthermore, for $k=3$ he claims he could prove the existence of some $0&lt;c_1\leq c_2$ such that\[\frac{\log\log n}{\log n}n+c_1\frac{n}{(\log n)^2}\leq g_3(n)\leq \frac{\log\log n}{\log n}n+c_2\frac{n}{(\log n)^2}.\]This is strange, since in \cite{Er64d} both the upper and lower bound techniques that he used in fact prove\[\frac{\log\log n}{\log n}n+c_1\frac{n}{\log n}\leq g_3(n)\leq \frac{\log\log n}{\log n}n+c_2\frac{n}{\log n}\]for some constants $0&lt;c_1\leq c_2$. My best guess is that the denominator of $(\log n)^2$ in \cite{Er69} is just an unfortunate repeated typo, and $\log n$ was intended in both the reported bounds and the main question. This correction is thanks to Tang, who noted independently (see the comments) the improved lower bound given above (indeed with an improvement in the constant $c_1$). The special case $k=2$ is the subject of [425] .




References

[Er64d] P. Erd\H{o}s, On the multiplicative representation of integers . Israel Journal of Mathematics (1964).

[Er69] Erd\H{o}s, Paul, Some applications of graph theory to number theory . The Many Facets of Graph Theory (Proc. Conf., Western Mich. Univ., Kalamazoo, Mich., 1968) (1969), 77-82.

[Er69b] Erd\H{o}s, P., Problems and results in chromatic graph theory . Proof Techniques in Graph Theory (Proc. Second Ann Arbor Graph Theory Conf., Ann Arbor, Mich., 1968) (1969), 27-35.

\noindent\rule{\linewidth}{0.4pt}


%% ============================================================
\section{Problem \#802}
\label{prob:802}

\noindent\textbf{Status:} OPEN \quad|\quad \textbf{Prize:} none \quad|\quad \textbf{Updated:} 2025-08-31

\noindent\textbf{Tags:} graph theory

\medskip

\noindent\textbf{Statement:}

Is it true that any $K_r$-free graph on $n$ vertices with average degree $t$ contains an independent set on\[\gg_r \frac{\log t}{t}n\]many vertices?




A conjecture of Ajtai, Erd\H{o}s, Koml\'{o}s, and Szemer\'{e}di \cite{AEKS81}, who proved that there must exist an independent set on\[\gg_r \frac{\log\log(t+1)}{t}n\]many vertices. Shearer \cite{Sh95} improved this to\[\gg_r \frac{\log t}{\log\log(t+1)t}n.\]Ajtai, Koml\'{o}s, and Szemer\'{e}di \cite{AKS80} proved the conjectured bound when $r=3$. Alon \cite{Al96b} proved the conjectured bound, but replacing the $K_r$-free assumption with the stronger assumption that the induced graph on every vertex neighbourhood has chromatic number $\leq r-2$.




References

[AEKS81] Ajtai, M. and Erd\H{o}s, P. and Koml\'{o}s, J. and Szemer\'{e}di, E., On Tur\'{a}n's theorem for sparse graphs . Combinatorica (1981), 313-317.

[AKS80] Ajtai, Mikl\'{o}s and Koml\'{o}s, J\'{a}nos and Szemer\'{e}di, Endre, A note on Ramsey numbers . J. Combin. Theory Ser. A (1980), 354-360.

[Al96b] Alon, Noga, Independence numbers of locally sparse graphs and a Ramsey type problem . Random Structures Algorithms (1996), 271-278.

[Sh95] Shearer, James B., On the independence number of sparse graphs . Random Structures Algorithms (1995), 269--271.

\noindent\rule{\linewidth}{0.4pt}


%% ============================================================
\section{Problem \#805}
\label{prob:805}

\noindent\textbf{Status:} OPEN \quad|\quad \textbf{Prize:} none \quad|\quad \textbf{Updated:} 2025-08-31

\noindent\textbf{Tags:} graph theory

\medskip

\noindent\textbf{Statement:}

For which functions $g(n)$ with $n&#62;g(n)\geq (\log n)^2$ is there a graph on $n$ vertices in which every induced subgraph on $g(n)$ vertices contains a clique of size $\geq \log n$ and an independent set of size $\geq \log n$? In particular, is there such a graph for $g(n)=(\log n)^3$?




A problem of Erd\H{o}s and Hajnal, who thought that there is no such graph for $g(n)=(\log n)^3$. Alon and Sudakov \cite{AlSu07} proved that there is no such graph with\[g(n)=\frac{c}{\log\log n}(\log n)^3\]for some constant $c&gt;0$. Alon, Buci\'{c}, and Sudakov \cite{ABS21} construct such a graph with\[g(n)\leq 2^{2^{(\log\log n)^{1/2+o(1)}}}.\]See also [804] .




References

[ABS21] Alon, Noga and Buci\'c, Matija and Sudakov, Benny, Large cliques and independent sets all over the place . Proc. Amer. Math. Soc. (2021), 3145-3157.

[AlSu07] Alon, Noga and Sudakov, Benny, On graphs with subgraphs having large independence numbers . J. Graph Theory (2007), 149-157.

\noindent\rule{\linewidth}{0.4pt}


%% ============================================================
\section{Problem \#809}
\label{prob:809}

\noindent\textbf{Status:} OPEN \quad|\quad \textbf{Prize:} none \quad|\quad \textbf{Updated:} 2025-08-31

\noindent\textbf{Tags:} graph theory, ramsey theory

\medskip

\noindent\textbf{Statement:}

Define the anti-Ramsey number $\chi_S(n,e,G)$ as the smallest $r$ such that there is a graph with $n$ vertices and $e$ edges with an $r$-colouring of its edges in which every copy of $G$ has entirely distinct edge colours. Is it true that, for all $k\geq 3$,\[\chi_S(n, \lfloor n^2/4\rfloor+1,C_{2k+1})\sim n^2/8?\]




A problem of Burr, Erd\H{o}s, Graham, and S\'{o}s \cite{BEGS89} who proved that\[\chi_S(n, \lfloor n^2/4\rfloor+1,C_{2k+1})\gg_k n^2.\]This question was solved in the affirmative for all $k\geq 4$ by Buci\'{c}, Chen, and Ma \cite{BCM26}. The situation for $C_3$ and $C_5$ are quite different - it is easy to see that\[\chi_S(n, \lfloor n^2/4\rfloor+1,C_{3})=3\]and Erd\H{o}s and Simonovits proved (as reported in \cite{BEGS89}) that\[\chi_S(n, \lfloor n^2/4\rfloor+1,C_{5})=\lfloor n/2\rfloor+3\]for all large $n$. See also [810] .




References

[BCM26] M. Buci\'{c}, K. Chen, and J. Ma, On a maximal anti-Ramsey conjeture of Burr, Erd\H{o}s, Graham, and S\'{o}s . arXiv:2603.18952 (2026).

[BEGS89] Burr, S. A. and Erd\H{o}s, P. and Graham, R. L. and S\'os, V. T., Maximal anti-{R}amsey graphs and the strong chromatic number . J. Graph Theory (1989), 263--282.

\noindent\rule{\linewidth}{0.4pt}


%% ============================================================
\section{Problem \#810}
\label{prob:810}

\noindent\textbf{Status:} OPEN \quad|\quad \textbf{Prize:} none \quad|\quad \textbf{Updated:} 2025-08-31

\noindent\textbf{Tags:} graph theory, ramsey theory

\medskip

\noindent\textbf{Statement:}

Does there exist some $\epsilon&#62;0$ such that, for all sufficiently large $n$, there exists a graph $G$ on $n$ vertices with at least $\epsilon n^2$ many edges such that the edges can be coloured with $n$ colours so that every $C_4$ receives $4$ distinct colours?




A problem of Burr, Erd\H{o}s, Graham, and S\'{o}s \cite{BEGS89}, who believed the answer is no. Equivalently, if $\chi_S(n,e,G)$ is the anti-Ramsey number, the smallest $r$ such that there is a graph with $n$ vertices and $e$ edges with an $r$-colouring of its edges in which every copy of $G$ has entirely distinct edge colours, then this asks whether there exists an $\epsilon&gt;0$ such that\[\chi_S(n,\epsilon n^2,C_4)\leq n\]for all large $n$. In \cite{BEGS89} they prove there does not exist such $\epsilon$ if we replace $C_4$ with $P_4$. In \cite{BEGS89} they ask the stronger question as to whether, for fixed $\epsilon&gt;0$,\[\chi_S(n,\epsilon n^2,G)/n\to \infty\]as $n\to \infty$ for all connected bipartite $G$ which is not a star. This was proved for all such $G$ except complete bipartite graphs by S\'{a}rk\"{o}zy and Selkow \cite{SaSe06}. It remains open in particular for $C_4$. In \cite{BEGS89} they also observe that, for some constant $c&gt;0$,\[\chi_S(n, cg(n;7,4),C_4)\leq n,\]where $g(n;7,4)$ is the maximum number of edges in a $3$-uniform hypergraph on $n$ vertices in which no $7$ vertices contain $4$ edges. It is unknown whether (but likely) that $g(n;7,4)=o(n^2)$ (see [1178] ). See also [809] .




References

[BEGS89] Burr, S. A. and Erd\H{o}s, P. and Graham, R. L. and S\'os, V. T., Maximal anti-{R}amsey graphs and the strong chromatic number . J. Graph Theory (1989), 263--282.

[SaSe06] S\'{a}rk\"ozy, G\'{a}bor N. and Selkow, Stanley, On an anti-{R}amsey problem of {B}urr, {E}rd\H os, {G}raham, and T. S\'os . J. Graph Theory (2006), 147--156.

\noindent\rule{\linewidth}{0.4pt}


%% ============================================================
\section{Problem \#811}
\label{prob:811}

\noindent\textbf{Status:} OPEN \quad|\quad \textbf{Prize:} none \quad|\quad \textbf{Updated:} 2025-08-31

\noindent\textbf{Tags:} graph theory, ramsey theory

\medskip

\noindent\textbf{Statement:}

Suppose $n\equiv 1\pmod{m}$. We say that an edge-colouring of $K_n$ using $m$ colours is balanced if every vertex sees exactly $\lfloor n/m\rfloor$ many edges of each colours. For which graphs $G$ is it true that, if $m=e(G)$, for all large $n\equiv 1\pmod{m}$, every balanced edge-colouring of $K_n$ with $m$ colours contains a rainbow copy of $G$? (That is, a subgraph isomorphic to $G$ where each edge receives a different colour.)




In \cite{Er91} Erd\H{o}s credits this problem to himself, Pyber, and Tuza. This problem was explored in a paper of Erd\H{o}s and Tuza \cite{ErTu93}. In \cite{Er96} Erd\H{o}s seems to suggest that this might be true for every graph $G$, and specifically asks specific challenge posed in \cite{Er91} and \cite{Er96} is whether, in any balanced edge-colouring of $K_{6n+1}$ by $6$ colours there must exist a rainbow $C_6$ and $K_4$. In general, one can ask for a quantitative version, defining $d_G(n)$ to be minimal (if it exists) such that if $n$ is sufficiently large and the edges of $K_n$ are coloured with $e(G)$ many colours such that the minimum degree of each colour class is $\geq d_G(n)$ then there is a rainbow copy of $G$. Erd\H{o}s and Tuza \cite{ErTu93} proved that\[\lfloor n/6\rfloor \leq d_{C_4}(n) \leq \left(\frac{1}{4}-c\right)n\]for some constant $c&#62;0$. Axenovich and Clemen \cite{AxCl24} have proved that there exist infinitely many graphs without this property. In particular, they show that for any odd $\ell \geq 3$ and $m=\lfloor \sqrt{\ell}+3.5\rfloor$ there exist arbitrarily large $n$ such that $K_n$ has a balanced edge-colouring using $\ell$ colours which contains no rainbow $K_m$. They conjecture that $K_m$ lacks this property for all $m\geq 4$. Clemen and Wagner \cite{ClWa23} proved that $K_4$ does lack this property.




References

[AxCl24] Axenovich, Maria and Clemen, Felix C., Rainbow subgraphs in edge-colored complete graphs: answering two questions by {E}rd\H{o}s and Tuza . J. Graph Theory (2024), 57--66.

[ClWa23] Clemen, Felix Christian and Wagner, Adam Zsolt, Balanced edge-colorings avoiding rainbow cliques of size four . Electron. J. Combin. (2023), Paper No. 3.17, 3.

[Er91] Erd\H{o}s, P., Problems and results in combinatorial analysis and combinatorial number theory . Graph theory, combinatorics, and applications, Vol. 1 (Kalamazoo, MI, 1988) (1991), 397-406.

[Er96] Erd\H{o}s, Paul, Some of my favourite problems on cycles and colourings . Tatra Mt. Math. Publ. (1996), 7-9.

[ErTu93] Erd\H{o}s, Paul and Tuza, Zsolt, Rainbow subgraphs in edge-colorings of complete graphs . (1993), 81--88.

\noindent\rule{\linewidth}{0.4pt}


%% ============================================================
\section{Problem \#812}
\label{prob:812}

\noindent\textbf{Status:} OPEN \quad|\quad \textbf{Prize:} none \quad|\quad \textbf{Updated:} 2025-08-31

\noindent\textbf{Tags:} graph theory, ramsey theory

\medskip

\noindent\textbf{Statement:}

Is it true that\[\frac{R(n+1)}{R(n)}\geq 1+c\]for some constant $c&#62;0$, for all large $n$? Is it true that\[R(n+1)-R(n) \gg n^2?\]




Burr, Erd\H{o}s, Faudree, and Schelp \cite{BEFS89} proved that\[R(n+1)-R(n) \geq 4n-8\]for all $n\geq 2$. The lower bound of [165] implies that\[R(n+2)-R(n) \gg n^{2-o(1)}.\]




References

[BEFS89] Burr, S. A. and Erd\H{o}s, P. and Faudree, R. J. and Schelp, R. H., On the difference between consecutive {R}amsey numbers . Utilitas Math. (1989), 115--118.

\noindent\rule{\linewidth}{0.4pt}


%% ============================================================
\section{Problem \#813}
\label{prob:813}

\noindent\textbf{Status:} OPEN \quad|\quad \textbf{Prize:} none \quad|\quad \textbf{Updated:} 2025-08-31

\noindent\textbf{Tags:} graph theory

\medskip

\noindent\textbf{Statement:}

Let $h(n)$ be minimal such that every graph on $n$ vertices where every set of $7$ vertices contains a triangle (a copy of $K_3$) must contain a clique on at least $h(n)$ vertices. Estimate $h(n)$ - in particular, do there exist constants $c_1,c_2&#62;0$ such that\[n^{1/3+c_1}\ll h(n) \ll n^{1/2-c_2}?\]




A problem of Erd\H{o}s and Hajnal, who could prove that\[n^{1/3}\ll h(n) \ll n^{1/2}.\]Buci\'{c} and Sudakov \cite{BuSu23} have proved\[h(n) \gg n^{5/12-o(1)}.\]




References

[BuSu23] M. Buci\'C and B. Sudakov, Large independent sets from local considerations . arXiv:2007.03667 (2023).

\noindent\rule{\linewidth}{0.4pt}


%% ============================================================
\section{Problem \#817}
\label{prob:817}

\noindent\textbf{Status:} OPEN \quad|\quad \textbf{Prize:} none \quad|\quad \textbf{Updated:} 2025-08-31

\noindent\textbf{Tags:} additive combinatorics

\medskip

\noindent\textbf{Statement:}

Let $k\geq 3$ and define $g_k(n)$ to be the minimal $N$ such that $\{1,\ldots,N\}$ contains some $A$ of size $\lvert A\rvert=n$ such that\[\langle A\rangle = \left\{\sum_{a\in A}\epsilon_aa: \epsilon_a\in \{0,1\}\right\}\]contains no non-trivial $k$-term arithmetic progression. Estimate $g_k(n)$. In particular, is it true that\[g_3(n) \gg 3^n?\]




A problem of Erd\H{o}s and S\'{a}rk\"{o}zy who proved\[g_3(n) \gg \frac{3^n}{n^{O(1)}}.\]

\noindent\rule{\linewidth}{0.4pt}


%% ============================================================
\section{Problem \#819}
\label{prob:819}

\noindent\textbf{Status:} OPEN \quad|\quad \textbf{Prize:} none \quad|\quad \textbf{Updated:} 2025-08-31

\noindent\textbf{Tags:} additive combinatorics

\medskip

\noindent\textbf{Statement:}

Let $f(N)$ be maximal such that there exists $A\subseteq \{1,\ldots,N\}$ with $\lvert A\rvert=\lfloor N^{1/2}\rfloor$ such that $\lvert (A+A)\cap [1,N]\rvert=f(N)$. Estimate $f(N)$.




Erd\H{o}s and Freud \cite{ErFr91} proved\[\left(\frac{3}{8}-o(1)\right)N \leq f(N) \leq \left(\frac{1}{2}+o(1)\right)N,\]and note that it is closely connected to the size of the largest quasi-Sidon set (see [840] ).




References

[ErFr91] Erd\H{o}s, P. and Freud, R., On sums of a Sidon-sequence . J. Number Theory (1991), 196--205.

\noindent\rule{\linewidth}{0.4pt}


%% ============================================================
\section{Problem \#820}
\label{prob:820}

\noindent\textbf{Status:} OPEN \quad|\quad \textbf{Prize:} none \quad|\quad \textbf{Updated:} 2025-08-31

\noindent\textbf{Tags:} number theory

\medskip

\noindent\textbf{Statement:}

Let $H(n)$ be the smallest integer $l$ such that there exist $k&lt;l$ with $(k^n-1,l^n-1)=1$. Is it true that $H(n)=3$ infinitely often? (That is, $(2^n-1,3^n-1)=1$ infinitely often?) Estimate $H(n)$. Is it true that there exists some constant $c&gt;0$ such that, for all $\epsilon&#62;0$,\[H(n) &#62; \exp(n^{(c-\epsilon)/\log\log n})\]for infinitely many $n$ and\[H(n) &#60; \exp(n^{(c+\epsilon)/\log\log n})\]for all large enough $n$? Does a similar upper bound hold for the smallest $k$ such that $(k^n-1,2^n-1)=1$?




Erd\H{o}s \cite{Er74b} proved that there exists a constant $c&gt;0$ such that\[H(n) &gt; \exp(n^{c/(\log\log n)^2})\]for infinitely many $n$. van Doorn in the comments sketches a proof of the lower bound: that there exists some constant $c&gt;0$ and infinitely many $n$ such that\[H(n) &gt; \exp(n^{c/\log\log n}).\]The sequence $H(n)$ for $1\leq n\leq 10$ is\[3,3,3,6,3,18,3,6,3,12.\]The sequence of $n$ for which $(2^n-1,3^n-1)=1$ is A263647 in the OEIS. See also [770] .




References

[Er74b] Erd\H{o}s, P., Remarks on some problems in number theory . Math. Balkanica (1974), 197-202.

\noindent\rule{\linewidth}{0.4pt}


%% ============================================================
\section{Problem \#821}
\label{prob:821}

\noindent\textbf{Status:} OPEN \quad|\quad \textbf{Prize:} none \quad|\quad \textbf{Updated:} 2025-08-31

\noindent\textbf{Tags:} number theory

\medskip

\noindent\textbf{Statement:}

Let $g(n)$ count the number of $m$ such that $\phi(m)=n$. Is it true that, for every $\epsilon&#62;0$, there exist infinitely many $n$ such that\[g(n) &#62; n^{1-\epsilon}?\]




Pillai proved that $\limsup g(n)=\infty$ and Erd\H{o}s \cite{Er35b} proved that there exists some constant $c&gt;0$ such that $g(n) &gt;n^c$ for infinitely many $n$. This conjecture would follow if we knew that, for every $\epsilon&gt;0$, there are $\gg_\epsilon \frac{x}{\log x}$ many primes $p&lt;x$ such that all prime factors of $p-1$ are $&lt;p^\epsilon$. The best known bound is that there are infinitely many $n$ such that\[g(n) &gt; n^{0.71568\cdots},\]obtained by Lichtman \cite{Li22} as a consequence of proving that there are $\geq \frac{x}{(\log x)^{O(1)}}$ many primes $p\leq x$ such that all prime factors of $p-1$ are $\leq x^{0.2843\cdots}$ (which improves a number of previous exponents, most recently Baker and Harman \cite{BaHa98}). The average size of $g(n)$ was investigated by Luca and Pollack \cite{LuPo11}. See also [416] .




References

[BaHa98] Baker, R. C. and Harman, G., Shifted primes without large prime factors . Acta Arith. (1998), 331--361.

[Er35b] Erd\H{o}s, P., On the normal number of prime factors of $p-1$ and some related problems concerning Euler's $\varphi$-function . Quart. J. Math. (1935), 205-213.

[Li22] J. D. Lichtman, Primes in arithmetic progressions to large moduli and shifted primes without large prime factors . arXiv:2211.09641 (2022).

[LuPo11] Luca, Florian and Pollack, Paul, An arithmetic function arising from {C}armichael's conjecture . J. Th\'{e}or. Nombres Bordeaux (2011), 697--714.

\noindent\rule{\linewidth}{0.4pt}


%% ============================================================
\section{Problem \#824}
\label{prob:824}

\noindent\textbf{Status:} OPEN \quad|\quad \textbf{Prize:} none \quad|\quad \textbf{Updated:} 2025-08-31

\noindent\textbf{Tags:} number theory

\medskip

\noindent\textbf{Statement:}

Let $h(x)$ count the number of integers $1\leq a&lt;b&lt;x$ such that $(a,b)=1$ and $\sigma(a)=\sigma(b)$, where $\sigma$ is the sum of divisors function. Is it true that $h(x)&gt;x^{2-o(1)}$?




Erd\H{o}s \cite{Er74b} proved that $\limsup h(x)/x= \infty$, and claimed a similar proof for this problem. A complete proof that $h(x)/x\to \infty$ was provided by Pollack and Pomerance \cite{PoPo16}. A similar question can be asked if we replace the condition $(a,b)=1$ with the condition that $a$ and $b$ are squarefree. Weisenberg suggests another variant, with the condition that there are no proper factors $u\mid a$ and $v\mid b$ such that $\sigma(u)=\sigma(v)$ and $(u,a/u)=(v,b/v)=1$, which is the weakest restriction one can impose that is still strong enough to eliminate trivial duplicates.




References

[Er74b] Erd\H{o}s, P., Remarks on some problems in number theory . Math. Balkanica (1974), 197-202.

[PoPo16] Pollack, Paul and Pomerance, Carl, Some problems of Erd\H{o}s on the sum-of-divisors function . Trans. Amer. Math. Soc. Ser. B (2016), 1-26.

\noindent\rule{\linewidth}{0.4pt}


%% ============================================================
\section{Problem \#826}
\label{prob:826}

\noindent\textbf{Status:} OPEN \quad|\quad \textbf{Prize:} none \quad|\quad \textbf{Updated:} 2025-08-31

\noindent\textbf{Tags:} number theory

\medskip

\noindent\textbf{Statement:}

Are there infinitely many $n$ such that, for all $k\geq 1$,\[\tau(n+k)\ll k?\]




A stronger form of [248] . Lau \cite{La26} has proved that there exists an absolute constant $C$ such that, for infinitely many $n$,\[\tau(n+k) \ll k^C\]for all $k\geq 1$.




References

[La26] C. F. Lau, On the number of prime factors of consecutive integers . arXiv:2604.15042 (2026).

\noindent\rule{\linewidth}{0.4pt}


%% ============================================================
\section{Problem \#827}
\label{prob:827}

\noindent\textbf{Status:} OPEN \quad|\quad \textbf{Prize:} none \quad|\quad \textbf{Updated:} 2025-08-31

\noindent\textbf{Tags:} geometry

\medskip

\noindent\textbf{Statement:}

Let $n_k$ be minimal such that if $n_k$ points in $\mathbb{R}^2$ are in general position then there exists a subset of $k$ points such that all $\binom{k}{3}$ triples determine circles of different radii. Determine $n_k$.




In \cite{Er75h} Erd\H{o}s asks whether $n_k$ exists. In \cite{Er78c} he gave a simple argument which proves that it does, and in fact\[n_k \leq k+2\binom{k-1}{2}\binom{k-1}{3},\]but this argument is incorrect, as explained by Martinez and Rold\'{a}n-Pensado \cite{MaRo15}. Martinez and Rold\'{a}n-Pensado give a corrected argument that proves $n_k\ll k^9$. A simple probabilistic argument, given in the comments by FlaredRain, improves this to $n_k\ll k^5$.




References

[Er75h] Erd\H{o}s, P., Some problems on elementary geometry . Austral. Math. Soc. Gaz. (1975), 2-3.

[Er78c] Erd\H{o}s, P., Some more problems on elementary geometry . Austral. Math. Soc. Gaz. (1978), 52-54.

[MaRo15] Mart\'Inez, L. and Rold\'{a}n-Pensado, E., Points defining triangles with distinct circumradii . Acta Math. Hungar. (2015), 136--141.

\noindent\rule{\linewidth}{0.4pt}


%% ============================================================
\section{Problem \#828}
\label{prob:828}

\noindent\textbf{Status:} OPEN \quad|\quad \textbf{Prize:} none \quad|\quad \textbf{Updated:} 2025-08-31

\noindent\textbf{Tags:} number theory

\medskip

\noindent\textbf{Statement:}

Is it true that, for any $a\in\mathbb{Z}$, there are infinitely many $n$ such that\[\phi(n) \mid n+a?\]




A conjecture of Graham. Lehmer has conjectured that $\phi(n)\mid n-1$ if and only if $n$ is prime. It is an easy exercise to show that $\phi(n) \mid n$ if and only if $n=2^a3^b$. This is discussed in problem B37 of Guy's collection \cite{Gu04}.




References

[Gu04] Guy, Richard K., Unsolved problems in number theory . (2004), xviii+437.

\noindent\rule{\linewidth}{0.4pt}


%% ============================================================
\section{Problem \#829}
\label{prob:829}

\noindent\textbf{Status:} OPEN \quad|\quad \textbf{Prize:} none \quad|\quad \textbf{Updated:} 2025-08-31

\noindent\textbf{Tags:} number theory

\medskip

\noindent\textbf{Statement:}

Let $A\subset\mathbb{N}$ be the set of cubes. Is it true that\[1_A\ast 1_A(n) \ll (\log n)^{O(1)}?\]




Mordell proved that\[\limsup_{n\to \infty} 1_A\ast 1_A(n)=\infty\]and Mahler \cite{Ma35b} proved\[1_A\ast 1_A(n) \gg (\log n)^{1/4}\]for infinitely many $n$. Stewart \cite{St08} improved this to\[1_A\ast 1_A(n) \gg (\log n)^{11/13}.\]




References

[Ma35b] Mahler, Kurt, On the Lattice Points on Curves of Genus 1 . Proc. London Math. Soc. (2) (1935), 431-466.

[St08] Stewart, Cameron L., Cubic Thue equations with many solutions . Int. Math. Res. Not. IMRN (2008), Art. ID rnn040, 11.

\noindent\rule{\linewidth}{0.4pt}


%% ============================================================
\section{Problem \#830}
\label{prob:830}

\noindent\textbf{Status:} OPEN \quad|\quad \textbf{Prize:} none \quad|\quad \textbf{Updated:} 2025-08-31

\noindent\textbf{Tags:} number theory

\medskip

\noindent\textbf{Statement:}

We say that $a,b\in \mathbb{N}$ are an amicable pair if $\sigma(a)=\sigma(b)=a+b$. Are there infinitely many amicable pairs? If $A(x)$ counts the number of amicable $1\leq a\leq b\leq x$ then is it true that\[A(x)&gt;x^{1-o(1)}?\]




For example $220$ and $284$. Erd\H{o}s \cite{Er55b} proved that $A(x)=o(x)$, and Pomerance \cite{Po81} improved this to\[A(x) \leq x \exp(-(\log x)^{1/3})\]and later \cite{Po15} to\[A(x) \leq x \exp(-(\tfrac{1}{2}+o(1))(\log x\log\log x)^{1/2}).\]This is problem B4 in Guy's collection \cite{Gu04}.




References

[Er55b] Erd\H{o}s, P., On amicable numbers . Publ. Math. Debrecen (1955), 108-111.

[Gu04] Guy, Richard K., Unsolved problems in number theory . (2004), xviii+437.

[Po15] Pomerance, Carl, On amicable numbers . (2015), 321-327.

[Po81] Pomerance, Carl, On the distribution of amicable numbers. {II} . J. Reine Angew. Math. (1981), 183-188.

\noindent\rule{\linewidth}{0.4pt}


%% ============================================================
\section{Problem \#831}
\label{prob:831}

\noindent\textbf{Status:} OPEN \quad|\quad \textbf{Prize:} none \quad|\quad \textbf{Updated:} 2025-08-31

\noindent\textbf{Tags:} geometry

\medskip

\noindent\textbf{Statement:}

Let $h(n)$ be maximal such that in any $n$ points in $\mathbb{R}^2$ (with no three on a line and no four on a circle) there are at least $h(n)$ many circles of different radii passing through three points. Estimate $h(n)$.




See also [104] and [506] .

\noindent\rule{\linewidth}{0.4pt}


%% ============================================================
\section{Problem \#836}
\label{prob:836}

\noindent\textbf{Status:} OPEN \quad|\quad \textbf{Prize:} none \quad|\quad \textbf{Updated:} 2025-08-31

\noindent\textbf{Tags:} graph theory, hypergraphs, chromatic number

\medskip

\noindent\textbf{Statement:}

Let $r\geq 2$ and $G$ be a $r$-uniform hypergraph with chromatic number $3$ (that is, there is a $3$-colouring of the vertices of $G$ such that no edge is monochromatic). Suppose any two edges of $G$ have a non-empty intersection. Must $G$ contain $O(r^2)$ many vertices? Must there be two edges which meet in $\gg r$ many vertices?




A problem of Erd\H{o}s and Shelah. The Fano geometry gives an example where there are no two edges which meet in $r-1$ vertices. Are there any other examples? Erd\H{o}s and Lov\'{a}sz \cite{ErLo75} proved that there must be two edges which meet in $\gg \frac{r}{\log r}$ many vertices. Alon has provided the following counterexample to the first question: as vertices take two sets $X$ and $Y$ of sizes $2r-2$ and $\frac{1}{2}\binom{2r-2}{r-1}$ respectively, where $Y$ corresponds to all partitions of $X$ into two equal parts. The edges are all subsets of $X$ of size $r$, and also all sets consisting of a subset of $X$ of size $r-1$ together with the unique element of $Y$ corresponding to the induced partition of $X$. This hypergraph is intersecting, its chromatic number is $3$, and it has $\asymp 4^r/\sqrt{r}$ many vertices.




References

[ErLo75] Erd\H{o}s, P. and Lov\'{a}sz, L., Problems and results on {$3$}-chromatic hypergraphs and some related questions . (1975), 609--627.

\noindent\rule{\linewidth}{0.4pt}


%% ============================================================
\section{Problem \#837}
\label{prob:837}

\noindent\textbf{Status:} OPEN \quad|\quad \textbf{Prize:} none \quad|\quad \textbf{Updated:} 2025-08-31

\noindent\textbf{Tags:} graph theory, hypergraphs

\medskip

\noindent\textbf{Statement:}

Let $k\geq 2$ and $A_k\subseteq [0,1]$ be the set of $\alpha$ such that there exists some $\beta(\alpha)&#62;\alpha$ with the property that, if $G_1,G_2,\ldots$ is a sequence of $k$-uniform hypergraphs with\[\liminf \frac{e(G_n)}{\binom{\lvert G_n\rvert}{k}} &#62;\alpha\]then there exist subgraphs $H_n\subseteq G_n$ such that $\lvert H_n\rvert \to \infty$ and\[\liminf \frac{e(H_n)}{\binom{\lvert H_n\rvert}{k}} &#62;\beta,\]and further that this property does not necessarily hold if $&#62;\alpha$ is replaced by $\geq \alpha$. What is $A_3$?




A problem of Erd\H{o}s and Simonovits. It is known that\[A_2 = \left\{ 1-\frac{1}{k} : k\geq 1\right\}.\]

\noindent\rule{\linewidth}{0.4pt}


%% ============================================================
\section{Problem \#838}
\label{prob:838}

\noindent\textbf{Status:} OPEN \quad|\quad \textbf{Prize:} none \quad|\quad \textbf{Updated:} 2025-08-31

\noindent\textbf{Tags:} geometry, convex

\medskip

\noindent\textbf{Statement:}

Let $f(n)$ be maximal such that any $n$ points in $\mathbb{R}^2$, with no three on a line, determine at least $f(n)$ different convex subsets. Estimate $f(n)$ - in particular, does there exist a constant $c$ such that\[\lim \frac{\log f(n)}{(\log n)^2}=c?\]




A question of Erd\H{o}s and Hammer. Erd\H{o}s proved in \cite{Er78c} that there exist constants $c_1,c_2&gt;0$ such that\[n^{c_1\log n}&lt;f(n)&lt; n^{c_2\log n}.\]See also [107] .




References

[Er78c] Erd\H{o}s, P., Some more problems on elementary geometry . Austral. Math. Soc. Gaz. (1978), 52-54.

\noindent\rule{\linewidth}{0.4pt}


%% ============================================================
\section{Problem \#839}
\label{prob:839}

\noindent\textbf{Status:} OPEN \quad|\quad \textbf{Prize:} none \quad|\quad \textbf{Updated:} 2025-08-31

\noindent\textbf{Tags:} number theory

\medskip

\noindent\textbf{Statement:}

Let $1\leq a_1&#60;a_2&#60;\cdots$ be a sequence of integers such that no $a_i$ is the sum of consecutive $a_j$ for $j&#60;i$. Is it true that\[\limsup \frac{a_n}{n}=\infty?\]Or even\[\lim \frac{1}{\log x}\sum_{a_n&#60;x}\frac{1}{a_n}=0?\]




Erd\H{o}s writes that it is easy to see that $\liminf a_n/n&lt;\infty$ is possible, and that one can have\[\sum_{a_n&lt; x}\frac{1}{a_n}\gg \log\log x.\]The upper density of such a sequence can be $1/2$, but Erd\H{o}s thought it probably could not be $&gt;1/2$. In fact this is false - Freud \cite{Fr93} constructed a sequence with upper density $19/36$. See also [359] and [867] .




References

[Fr93] R. Freud, Adding numbers - on a problem of P. Erd\H{o}s . James Cook Mathematical Notes (1993), 6199-6202.

\noindent\rule{\linewidth}{0.4pt}


%% ============================================================
\section{Problem \#840}
\label{prob:840}

\noindent\textbf{Status:} OPEN \quad|\quad \textbf{Prize:} none \quad|\quad \textbf{Updated:} 2025-08-31

\noindent\textbf{Tags:} additive combinatorics, sidon sets

\medskip

\noindent\textbf{Statement:}

Let $f(N)$ be the size of the largest quasi-Sidon subset $A\subset\{1,\ldots,N\}$, where we say that $A$ is quasi-Sidon if\[\lvert A+A\rvert=(1+o(1))\binom{\lvert A\rvert}{2}.\]How does $f(N)$ grow?




Considered by Erd\H{o}s and Freud \cite{ErFr91}, who proved\[\left(\frac{2}{\sqrt{3}}+o(1)\right)N^{1/2} \leq f(N) \leq \left(2+o(1)\right)N^{1/2}.\](Although both bounds were already given by Erd\H{o}s in \cite{Er81h}.) Note that $2/\sqrt{3}=1.15\cdots$. The lower bound is taking a genuine Sidon set $B\subset [1,N/3]$ of size $\sim N^{1/2}/\sqrt{3}$ and taking the union with $\{N-b : b\in B\}$. The upper bound was improved by Pikhurko \cite{Pi06} to\[f(N) \leq \left(\left(\frac{1}{4}+\frac{1}{(\pi+2)^2}\right)^{-1/2}+o(1)\right)N^{1/2}\](the constant here is $=1.863\cdots$). The analogous question with $A-A$ in place of $A+A$ is simpler, and there the maximal size is $\sim N^{1/2}$, as proved by Cilleruelo . See also [30] , [819] , and [864] .




References

[Er81h] Erd\H{o}s, P., Some problems and results on additive and multiplicative number theory . Analytic number theory (Philadelphia, Pa., 1980) (1981), 171-182.

[ErFr91] Erd\H{o}s, P. and Freud, R., On sums of a Sidon-sequence . J. Number Theory (1991), 196--205.

[Pi06] Pikhurko, Oleg, Dense edge-magic graphs and thin additive bases . Discrete Math. (2006), 2097--2107.

\noindent\rule{\linewidth}{0.4pt}


%% ============================================================
\section{Problem \#849}
\label{prob:849}

\noindent\textbf{Status:} OPEN \quad|\quad \textbf{Prize:} none \quad|\quad \textbf{Updated:} 2025-08-31

\noindent\textbf{Tags:} number theory, binomial coefficients

\medskip
\noindent\textit{Singmaster's conjecture}


\noindent\textbf{Statement:}

Is it true that, for every integer $t\geq 1$, there is some integer $a$ such that\[\binom{n}{k}=a\](with $1\leq k\leq n/2$) has exactly $t$ solutions?




Erd\H{o}s \cite{Er96b} credits this to himself and Gordon 'many years ago', but it is more commonly known as Singmaster's conjecture . For $t=3$ one could take $a=120$, and for $t=4$ one could take $a=3003$. There are no known examples for $t\geq 5$. Both Erd\H{o}s and Singmaster believed the answer to this question is no, and in fact that there exists an absolute upper bound on the number of solutions. Matomäki, Radziwill, Shao, Tao, and Teräväinen \cite{MRSTT22} have proved that there are always at most two solutions if we restrict $k$ to\[k\geq \exp((\log n)^{2/3+\epsilon}),\]assuming $a$ is sufficiently large depending on $\epsilon&gt;0$.




References

[Er96b] Erd\H{o}s, Paul, Some problems I presented or planned to present in my short talk . Analytic number theory, Vol. 1 (Allerton Park, IL, 1995) (1996), 333-335.

[MRSTT22] Matom\"{a}ki, Kaisa and Radziwi\l\l, Maksym and Shao, Xuancheng and Tao, Terence and Ter\"{a}v\"{a}inen, Joni, Singmaster's conjecture in the interior of {P}ascal's triangle . Q. J. Math. (2022), 1137--1177.

\noindent\rule{\linewidth}{0.4pt}


%% ============================================================
\section{Problem \#850}
\label{prob:850}

\noindent\textbf{Status:} OPEN \quad|\quad \textbf{Prize:} none \quad|\quad \textbf{Updated:} 2025-08-31

\noindent\textbf{Tags:} number theory, primes

\medskip
\noindent\textit{Erdős-Woods conjecture}


\noindent\textbf{Statement:}

Can there exist two distinct integers $x$ and $y$ such that $x,y$ have the same prime factors, $x+1,y+1$ have the same prime factors, and $x+2,y+2$ also have the same prime factors?




This is sometimes known as the Erd\H{o}s-Woods conjecture. For just $x,y$ and $x+1,y+1$ one can take\[x=2(2^r-1)\]and\[y = x(x+2).\]Erd\H{o}s also asked whether there are any other examples. Makowski \cite{Ma68} observed that $x=75$ and $y=1215$ is another example, since\[75 = 3\cdot 5^2 \textrm{ and }1215 = 3^5\cdot 5\]while\[76 = 2^2\cdot 19\textrm{ and }1216 = 2^6\cdot 19.\](This example was also found independently by Matthew Bolan, and by Dubickas, who posed it as part of the 2024 team selection test in Lithuania.) No other examples are known. This sequence is listed as A343101 at the OEIS. Shorey and Tijdeman \cite{ShTi16} have shown that, assuming a strong form of the ABC conjecture due to Baker, then the answer to the original problem is no. See also [677] . The case of $x,y$ and $x+1,y+1$ appeared as Problem 1 in the Third Benelux Mathematical Olympiad 2011. This problem is discussed in problem B19 of Guy's collection \cite{Gu04}.




References

[Gu04] Guy, Richard K., Unsolved problems in number theory . (2004), xviii+437.

[Ma68] Makowski, Andrzej, On a problem of {E}rd\H{o}s . Enseign. Math. (2) (1968), 193.

[ShTi16] Shorey, Tarlok N. and Tijdeman, Rob, Arithmetic properties of blocks of consecutive integers . (2016), 455--471.

\noindent\rule{\linewidth}{0.4pt}


%% ============================================================
\section{Problem \#852}
\label{prob:852}

\noindent\textbf{Status:} OPEN \quad|\quad \textbf{Prize:} none \quad|\quad \textbf{Updated:} 2025-08-31

\noindent\textbf{Tags:} number theory, primes

\medskip

\noindent\textbf{Statement:}

Let $d_n=p_{n+1}-p_n$, where $p_n$ is the $n$th prime. Let $h(x)$ be maximal such that for some $n&lt;x$ the numbers $d_n,d_{n+1},\ldots,d_{n+h(x)-1}$ are all distinct. Estimate $h(x)$. In particular, is it true that\[h(x) &gt;(\log x)^c\]for some constant $c&#62;0$, and\[h(x)=o(\log x)?\]




Brun's sieve implies $h(x) \to \infty$ as $x\to \infty$.

\noindent\rule{\linewidth}{0.4pt}


%% ============================================================
\section{Problem \#853}
\label{prob:853}

\noindent\textbf{Status:} OPEN \quad|\quad \textbf{Prize:} none \quad|\quad \textbf{Updated:} 2025-08-31

\noindent\textbf{Tags:} number theory, primes

\medskip

\noindent\textbf{Statement:}

Let $d_n=p_{n+1}-p_n$, where $p_n$ is the $n$th prime. Let $r(x)$ be the smallest even integer $t$ such that $d_n=t$ has no solutions for $n\leq x$. Is it true that $r(x)\to \infty$? Or even $r(x)/\log x \to \infty$?




In \cite{Er85c} Erd\H{o}s omits the condition that $t$ be even, but this is clearly necessary.




References

[Er85c] Erd\H{o}s, P., On some of my problems in number theory I would most like to see solved . Number theory (Ootacamund, 1984) (1985), 74-84.

\noindent\rule{\linewidth}{0.4pt}


%% ============================================================
\section{Problem \#854}
\label{prob:854}

\noindent\textbf{Status:} OPEN \quad|\quad \textbf{Prize:} none \quad|\quad \textbf{Updated:} 2025-08-31

\noindent\textbf{Tags:} number theory

\medskip

\noindent\textbf{Statement:}

Let $n_k$ denote the $k$th primorial, i.e. the product of the first $k$ primes. If $1=a_1&#60;a_2&#60;\cdots a_{\phi(n_k)}=n_k-1$ is the sequence of integers coprime to $n_k$, then estimate the smallest even integer not of the form $a_{i+1}-a_i$. Are there\[\gg \max_i (a_{i+1}-a_i)\]many even integers of the form $a_{j+1}-a_j$?




This was asked by Erd\H{o}s in Oberwolfach (most likely in 1986). Clearly all differences $a_{i+1}-a_i$ are even. Erd\H{o}s first thought that (for large enough $k$) all even $t\leq \max(a_{i+1}-a_i)$ can be written as $t=a_{j+1}-a_j$ for some $j$, but in \cite{Ob1} writes 'perhaps this is false', and reports some computations of Lacampagne and Selfridge that this fails for $n_k=2\cdot 3\cdot 5\cdot 7\cdot 11\cdot 13$ which 'show some doubt on [his] conjecture', and says it could fail for all or infinitely many $k$. In \cite{Ob1} he also asks about the set of $j$ for which $a_{j+1}-a_j=\max(a_{i+1}-a_i)$, and in particular asks for estimates on the number of such $j$ and the minimal value of such a $j$.




References

[Ob1] P. Erd\H{o}s, Oberwolfach Mathematical Problems, Volume 1 . Mathematisches Forschungsinstitut Oberwolfach (Various).

\noindent\rule{\linewidth}{0.4pt}


%% ============================================================
\section{Problem \#855}
\label{prob:855}

\noindent\textbf{Status:} OPEN \quad|\quad \textbf{Prize:} none \quad|\quad \textbf{Updated:} 2026-03-14

\noindent\textbf{Tags:} number theory, primes

\medskip
\noindent\textit{second Hardy-Littlewood conjecture}


\noindent\textbf{Statement:}

If $\pi(x)$ counts the number of primes in $[1,x]$ then is it true that (for large $x$ and $y$)\[\pi(x+y) \leq \pi(x)+\pi(y)?\]




Commonly known as the second Hardy-Littlewood conjecture . In \cite{Er85c} Erd\H{o}s describes it as 'an old conjecture of mine which was probably already stated by Hardy and Littlewood'. This is probably false, since Hensley and Richards \cite{HeRi73} have shown that this is false assuming the Hardy-Littlewood prime tuples conjecture . Indeed, assuming this conjecture, they prove that for every large $y$ there are infinitely many $x$ such that\[\pi(x+y)&gt;\pi(x)+\pi(y)+(\log 2-o(1))\frac{y}{(\log y)^2},\]where the $o(1)$ term tends to $0$ as $y\to \infty$. Erd\H{o}s \cite{Er85c} reports Straus as remarking that the 'correct way' of stating this conjecture would have been\[\pi(x+y) \leq \pi(x)+2\pi(y/2).\]Clark and Jarvis \cite{ClJa01} have shown this is also incompatible with the prime tuples conjecture. In \cite{Er85c} Erd\H{o}s conjectures the weaker result (which in particular follows from the conjecture of Straus) that\[\pi(x+y) \leq \pi(x)+\pi(y)+O\left(\frac{y}{(\log y)^2}\right),\]which the Hensley and Richards result shows (conditionally) would be best possible. Richards conjectured that this is false. Erd\H{o}s and Richards further conjectured that the original inequality is true almost always - that is, the set of $x$ such that $\pi(x+y)\leq \pi(x)+\pi(y)$ for all $y&lt;x$ has density $1$. They could only prove that this set has positive lower density. They also conjectured that for every $x$ the inequality $\pi(x+y)\leq \pi(x)+\pi(y)$ is true provided $y \gg (\log x)^C$ for some large constant $C&gt;0$. In \cite{Er80} Erd\H{o}s conjectures that\[\pi(x+y)\leq \pi(x)+O\left(\frac{y}{\log x}\right)\]for every $x\geq y \gg (\log x)^C$ for some large constant $C&gt;0$. Hardy and Littlewood proved\[\pi(x+y) \leq \pi(x)+O(\pi(y)).\]The best known in this direction is a result of Montgomery and Vaughan \cite{MoVa73}, which shows\[\pi(x+y) \leq \pi(x)+2\frac{y}{\log y}.\]This is discussed in problem A9 of Guy's collection \cite{Gu04}.




References

[ClJa01] Clark, David A. and Jarvis, Norman C., Dense admissible sequences . Math. Comp. (2001), 1713--1718.

[Er80] Erd\H{o}s, Paul, A survey of problems in combinatorial number theory . Ann. Discrete Math. (1980), 89-115.

[Er85c] Erd\H{o}s, P., On some of my problems in number theory I would most like to see solved . Number theory (Ootacamund, 1984) (1985), 74-84.

[Gu04] Guy, Richard K., Unsolved problems in number theory . (2004), xviii+437.

[HeRi73] Hensley, Douglas and Richards, Ian, On the incompatibility of two conjectures concerning primes . (1973), 123--127.

[MoVa73] Montgomery, H. L. and Vaughan, R. C., The large sieve . Mathematika (1973), 119--134.

\noindent\rule{\linewidth}{0.4pt}


%% ============================================================
\section{Problem \#856}
\label{prob:856}

\noindent\textbf{Status:} OPEN \quad|\quad \textbf{Prize:} none \quad|\quad \textbf{Updated:} 2025-08-31

\noindent\textbf{Tags:} number theory

\medskip

\noindent\textbf{Statement:}

Let $k\geq 3$ and $f_k(N)$ be the maximum value of $\sum_{n\in A}\frac{1}{n}$, where $A$ ranges over all subsets of $\{1,\ldots,N\}$ which contain no subset of size $k$ with the same pairwise least common multiple. Estimate $f_k(N)$.




Erd\H{o}s \cite{Er70} notes that\[f_k(N) \ll \frac{\log N}{\log\log N}.\]Indeed, let $A$ be such a set. This in particular implies that, for every $t$, there are $&lt;k$ solutions to $t=ap$ with $a\in A$ and $p$ prime, whence\[\sum_{n\in A}\frac{1}{n}\sum_{p&lt;N}\frac{1}{p}&lt; k \sum_{t&lt;N^2}\frac{1}{t} \ll \log N,\]and the bound follows since $\sum_{p&lt;N}\frac{1}{p}\gg \log\log N$. Improved bounds have been given by Tang and Zhang \cite{TaZh25b}, who proved bounds of the shape\[(\log N)^{b_k-o(1)}\leq f_k(N)\leq (\log N)^{c_k+o(1)}\]for some constants $0&lt;b_k\leq c_k\leq 1$, and in particular\[(\log N)^{0.438}\leq f_3(N)\leq (\log N)^{0.889},\]say, for all large $N$. The exponents here are related to progress in the sunflower conjecture [857] , to which this problem is closely related. For example, the exponents $c_k$ are $&lt;1$ (and so the upper bound above is non-trivial) if and only if [857] holds for $k$-sunflowers. The analogous question with natural density in place of logarithmic density (that is, we measure $\lvert A\rvert$ in place of $\sum_{n\in A}\frac{1}{n}$) is the subject of [536] . In particular Erd\H{o}s \cite{Er70} has constructed $A\subseteq \{1,\ldots,N\}$ with $\lvert A\rvert \gg N$ where no four have the same pairwise least common multiple, and hence the interest of the natural density problem is the $k=3$ case. A related combinatorial problem is asked at [857] .




References

[Er70] Erd\H{o}s, Paul, Some extremal problems in combinatorial number theory . Mathematical Essays Dedicated to A. J. Macintyre (1970), 123-133.

[TaZh25b] Q. Tang and S. Zhang, Harmonic LCM patterns and sunflower-free capacity . arXiv:2512.20055 (2025).

\noindent\rule{\linewidth}{0.4pt}


%% ============================================================
\section{Problem \#857}
\label{prob:857}

\noindent\textbf{Status:} OPEN \quad|\quad \textbf{Prize:} none \quad|\quad \textbf{Updated:} 2025-08-31

\noindent\textbf{Tags:} combinatorics

\medskip
\noindent\textit{weak sunflower problem}


\noindent\textbf{Statement:}

Let $m=m(n,k)$ be minimal such that in any collection of sets $A_1,\ldots,A_m\subseteq \{1,\ldots,n\}$ there must exist a sunflower of size $k$ - that is, some collection of $k$ of the $A_i$ which pairwise have the same intersection. Estimate $m(n,k)$, or even better, give an asymptotic formula.




Related to [536] and [856] . In \cite{Er70} Erd\H{o}s asks this in the equivalent formulation with intersection replaced by union. This is sometimes known as the weak sunflower problem (see [20] for the strong sunflower problem). When $k=3$ this is strongly connected to the cap set problem (finding the maximal size of subsets of $\mathbb{F}_3^n$ with no three-term arithmetic progressions), as observed by Alon, Shpilka, and Umans \cite{ASU13}). Naslund and Sawin \cite{NaSa17} have proved that\[m(n,3) \leq (3/2^{2/3})^{(1+o(1))n}.\]




References

[ASU13] Alon, Noga and Shpilka, Amir and Umans, Christopher, On sunflowers and matrix multiplication . Comput. Complexity (2013), 219--243.

[Er70] Erd\H{o}s, Paul, Some extremal problems in combinatorial number theory . Mathematical Essays Dedicated to A. J. Macintyre (1970), 123-133.

[NaSa17] Naslund, Eric and Sawin, Will, Upper bounds for sunflower-free sets . Forum Math. Sigma (2017), Paper No. e15, 10.

\noindent\rule{\linewidth}{0.4pt}


%% ============================================================
\section{Problem \#859}
\label{prob:859}

\noindent\textbf{Status:} OPEN \quad|\quad \textbf{Prize:} none \quad|\quad \textbf{Updated:} 2025-08-31

\noindent\textbf{Tags:} number theory, divisors

\medskip

\noindent\textbf{Statement:}

Let $t\geq 1$ and let $d_t$ be the density of the set of integers $n\in\mathbb{N}$ for which $t$ can be represented as the sum of distinct divisors of $n$. Do there exist constants $c_1,c_2&#62;0$ such that\[d_t \sim \frac{c_1}{(\log t)^{c_2}}\]as $t\to \infty$?




Erd\H{o}s \cite{Er70} proved that $d_t$ always exists, and that there exist some constants $c_3,c_4&#62;0$ such that\[\frac{1}{(\log t)^{c_3}} &#60; d_t &#60; \frac{1}{(\log t)^{c_4}}.\]




References

[Er70] Erd\H{o}s, Paul, Some extremal problems in combinatorial number theory . Mathematical Essays Dedicated to A. J. Macintyre (1970), 123-133.

\noindent\rule{\linewidth}{0.4pt}


%% ============================================================
\section{Problem \#860}
\label{prob:860}

\noindent\textbf{Status:} OPEN \quad|\quad \textbf{Prize:} none \quad|\quad \textbf{Updated:} 2025-08-31

\noindent\textbf{Tags:} number theory, primes

\medskip

\noindent\textbf{Statement:}

Let $h(n)$ be such that, for any $m\geq 1$, in the interval $(m,m+h(n))$ there exist distinct integers $a_i$ for $1\leq i\leq \pi(n)$ such that $p_i\mid a_i$, where $p_i$ denotes the $i$th prime. Estimate $h(n)$.




A problem of Erd\H{o}s and Pomerance \cite{ErPo80}, who proved that\[h(n) \ll \frac{n^{3/2}}{(\log n)^{1/2}}.\]Erd\H{o}s and Selfridge proved $h(n)&gt;(3-o(1))n$, and Ruzsa proved $h(n)/n\to \infty$. This is discussed in problem B32 of Guy's collection \cite{Gu04}. See also [375] .




References

[ErPo80] P. Erd\H{o}s and C. Pomerance, Matching the natural numbers up to $n$ with distinct multiples of another interval . Indigationes Math. (1980), 147-151.

[Gu04] Guy, Richard K., Unsolved problems in number theory . (2004), xviii+437.

\noindent\rule{\linewidth}{0.4pt}


%% ============================================================
\section{Problem \#864}
\label{prob:864}

\noindent\textbf{Status:} OPEN \quad|\quad \textbf{Prize:} none \quad|\quad \textbf{Updated:} 2025-08-31

\noindent\textbf{Tags:} number theory, sidon sets, additive combinatorics

\medskip

\noindent\textbf{Statement:}

Let $A\subseteq \{1,\ldots N\}$ be a set such that there exists at most one $n$ with more than one solution to $n=a+b$ (with $a\leq b\in A$). Estimate the maximal possible size of $\lvert A\rvert$ - in particular, is it true that\[\lvert A\rvert \leq (1+o(1))\frac{2}{\sqrt{3}}N^{1/2}?\]




A problem of Erd\H{o}s and Freud, who prove that\[\lvert A\rvert \geq (1+o(1))\frac{2}{\sqrt{3}}N^{1/2}.\]This is shown by taking a genuine Sidon set $B\subset [1,N/3]$ of size $\sim N^{1/2}/\sqrt{3}$ and taking the union with $\{N-b : b\in B\}$. For the analogous question with $n=a-b$ they prove that $\lvert A\rvert\sim N^{1/2}$. This is a weaker form of [840] .

\noindent\rule{\linewidth}{0.4pt}


%% ============================================================
\section{Problem \#865}
\label{prob:865}

\noindent\textbf{Status:} OPEN \quad|\quad \textbf{Prize:} none \quad|\quad \textbf{Updated:} 2025-08-31

\noindent\textbf{Tags:} number theory, additive combinatorics

\medskip

\noindent\textbf{Statement:}

There exists a constant $C&#62;0$ such that, for all large $N$, if $A\subseteq \{1,\ldots,N\}$ has size at least $\frac{5}{8}N+C$ then there are distinct $a,b,c\in A$ such that $a+b,a+c,b+c\in A$.




A problem of Erd\H{o}s and S\'{o}s (also earlier considered by Choi, Erd\H{o}s, and Szemer\'{e}di \cite{CES75}, but Erd\H{o}s had forgotten this). Taking all integers in $[N/8,N/4]$ and $[N/2,N]$ shows that $\frac{5}{8}$ would be best possible here. It is a classical folklore fact that if $A\subseteq \{1,\ldots,2N\}$ has size $\geq N+2$ then there are distinct $a,b\in A$ such that $a+b\in A$, which establishes the $k=2$ case. In general, one can define $f_k(N)$ to be minimal such that if $A\subseteq \{1,\ldots,N\}$ has size at least $f_k(N)$ then there are $k$ distinct $a_i\in A$ such that all $\binom{k}{2}$ pairwise sums are elements of $A$. Erd\H{o}s and S\'{o}s conjectured that\[f_k(N)\sim \frac{1}{2}\left(1+\sum_{1\leq r\leq k-2}\frac{1}{4^r}\right) N,\]and a similar example shows that this would be best possible. Choi, Erd\H{o}s, and Szemer\'{e}di \cite{CES75} have proved that, for all $k\geq 3$, there exists $\epsilon_k&#62;0$ such that (for large enough $N$)\[f_k(N)\leq \left(\frac{2}{3}-\epsilon_k\right)N.\]




References

[CES75] Choi, S. L. G. and Erd\H{o}s, P. and Szemer\'{e}di, E., Some additive and multiplicative problems in number theory . Acta Arith. (1975), 37--50.

\noindent\rule{\linewidth}{0.4pt}


%% ============================================================
\section{Problem \#866}
\label{prob:866}

\noindent\textbf{Status:} OPEN \quad|\quad \textbf{Prize:} none \quad|\quad \textbf{Updated:} 2025-08-31

\noindent\textbf{Tags:} number theory, additive combinatorics

\medskip

\noindent\textbf{Statement:}

Let $k\geq 3$ and $g_k(N)$ be minimal such that if $A\subseteq \{1,\ldots,2N\}$ has $\lvert A\rvert \geq N+g_k(N)$ then there exist integers $b_1,\ldots,b_k$ such that all $\binom{k}{2}$ pairwise sums are in $A$ (but the $b_i$ themselves need not be in $A$). Estimate $g_k(N)$.




A problem of Choi, Erd\H{o}s, and Szemer\'{e}di. It is clear that, for the set of odd numbers in $\{1,\ldots,2N\}$, no such $b_i$ exist, whence $g_k(N)\geq 0$ always. Choi, Erd\H{o}s, and Szemer\'{e}di proved that $g_3(N)=2$ and $g_4(N) \ll 1$. van Doorn has shown that $g_4(N)\leq 2032$. Choi, Erd\H{o}s, and Szemer\'{e}di also proved that\[g_5(N)\asymp \log N\]and\[g_6(N)\asymp N^{1/2}.\]In general they proved that\[g_k(N) \ll_k N^{1-2^{-k}}\]and for every $\epsilon&gt;0$ if $k$ is sufficiently large then\[g_k(N) &gt; N^{1-\epsilon}.\]As an example, taking $A$ to be the set of all odd integers and the powers of $2$ shows that $g_5(N)\gg \log N$.

\noindent\rule{\linewidth}{0.4pt}


%% ============================================================
\section{Problem \#870}
\label{prob:870}

\noindent\textbf{Status:} OPEN \quad|\quad \textbf{Prize:} none \quad|\quad \textbf{Updated:} 2025-08-31

\noindent\textbf{Tags:} number theory, additive basis

\medskip

\noindent\textbf{Statement:}

Let $k\geq 3$ and $A$ be an additive basis of order $k$. Does there exist a constant $c=c(k)&#62;0$ such that if $r(n)\geq c\log n$ for all large $n$ then $A$ must contain a minimal basis of order $k$? (Here $r(n)$ counts the number of representations of $n$ as the sum of at most $k$ elements from $A$.)




A question of Erd\H{o}s and Nathanson \cite{ErNa79}, who proved that this is true for $k=2$ if $1_A\ast 1_A(n) &gt; (\log \frac{4}{3})^{-1}\log n$ for all large $n$. Härtter \cite{Ha56} and Nathanson \cite{Na74} proved that there exist additive bases which do not contain any minimal additive bases. See also [868] .




References

[ErNa79] Erd\H{o}s, Paul and Nathanson, Melvyn B., Systems of distinct representatives and minimal bases in additive number theory . (1979), 89--107.

[Ha56] H\"{a}rtter, Erich, Ein Beitrag zur Theorie der {M}inimalbasen . J. Reine Angew. Math. (1956), 170--204.

[Na74] Nathanson, Melvyn B., Minimal bases and maximal nonbases in additive number theory . J. Number Theory (1974), 324--333.

\noindent\rule{\linewidth}{0.4pt}


%% ============================================================
\section{Problem \#872}
\label{prob:872}

\noindent\textbf{Status:} OPEN \quad|\quad \textbf{Prize:} none \quad|\quad \textbf{Updated:} 2025-08-31

\noindent\textbf{Tags:} number theory, primitive sets

\medskip

\noindent\textbf{Statement:}

Consider the two-player game in which players alternately choose integers from $\{2,3,\ldots,n\}$ to be included in some set $A$ (the same set for both players) such that no $a\mid b$ for $a\neq b\in A$. The game ends when no legal move is possible. One player wants the game to last as long as possible, the other wants the game to end quickly. How long can the game be guaranteed to last for? At least $\epsilon n$ moves? (For $\epsilon&#62;0$ and $n$ sufficiently large.) At least $(1-\epsilon)\frac{n}{2}$ moves?




A number theoretic variant of a combinatorial game of Hajnal, in which players alternately add edges to a graph while keeping it triangle-free. This graph theoretic game must trivially end in at most $n^2/4$ moves, and F\"{u}redi and Seress \cite{FuSe91} proved that it can be guaranteed to last for $\gg n\log n$ moves. Bir\'{o}, Horn, and Wildstrom \cite{BPW16} proved that it must end in at most $(\frac{26}{121}+o(1))n^2&lt;0.215n^2$ moves. This type of game is known as a saturation game. Erd\H{o}s does not specify which player goes first, which may result in different answers. GPT-5.2 Pro, prompted by Price, has shown that the final question has a negative answer: assuming the 'Prolonger' goes first, the 'Shortener' can guarantee that it ends in at most $(\frac{23}{48}+o(1))n$ moves. This constant has been refined, see the comment section for the latest. Since all primes in $(n/2,n]$ must eventually be chosen, the game must last for at least $\gg \frac{n}{\log n}$ moves.




References

[BPW16] Bir\'{o}, Csaba and Horn, Paul and Wildstrom, D. Jacob, An upper bound on the extremal version of Hajnal's triangle-free game . Discrete Appl. Math. (2016), 20--28.

[FuSe91] F\"{u}redi, Zolt\'{a}n and Reimer, Dave and Seress, \'{A}kos, Hajnal's triangle-free game and extremal graph problems . Congr. Numer. (1991), 123--128.

\noindent\rule{\linewidth}{0.4pt}


%% ============================================================
\section{Problem \#873}
\label{prob:873}

\noindent\textbf{Status:} OPEN \quad|\quad \textbf{Prize:} none \quad|\quad \textbf{Updated:} 2025-08-31

\noindent\textbf{Tags:} number theory

\medskip

\noindent\textbf{Statement:}

Let $A=\{a_1&lt;a_2&lt;\cdots\}\subseteq \mathbb{N}$ and let $F(A,X,k)$ count the number of $i$ such that\[[a_i,a_{i+1},\ldots,a_{i+k-1}] &lt; X,\]where the left-hand side is the least common multiple. Is it true that, for every $\epsilon &gt;0$, there exists some $k$ such that\[F(A,X,k)&#60;X^\epsilon?\]




A problem of Erd\H{o}s and Szemer\'{e}di, who proved that for every $A$\[F(A,X,3) \ll X^{1/3}\log X,\]and there is an $A$ such that\[F(A,X,3) \gg X^{1/3}\log X\]for infinitely many $X$. There may be a sequence for which this holds for every $X$.

\noindent\rule{\linewidth}{0.4pt}


%% ============================================================
\section{Problem \#875}
\label{prob:875}

\noindent\textbf{Status:} OPEN \quad|\quad \textbf{Prize:} none \quad|\quad \textbf{Updated:} 2025-08-31

\noindent\textbf{Tags:} additive combinatorics

\medskip

\noindent\textbf{Statement:}

Let $A=\{a_1&#60;a_2&#60;\cdots\}\subset \mathbb{N}$ be an infinite set such that the sets\[S_r = \{ a_1+\cdots +a_r : a_1&#60;\cdots&#60;a_r\in A\}\]are disjoint for distinct $r\geq 1$. How fast can such a sequence grow? How small can $a_{n+1}-a_n$ be? In particular, for which $c$ is it possible that $a_{n+1}-a_n\leq n^{c}$?




A problem of Deshouillers and Erd\H{o}s (an infinite version of [874] ). Such sets are sometimes called admissible. Erd\H{o}s writes 'it [is not] completely trivial to find such a sequence for which $a_{n+1}/a_n\to 1$'. It is not clear from this whether Deshouillers and Erd\H{o}s knew of such a sequence.

\noindent\rule{\linewidth}{0.4pt}


%% ============================================================
\section{Problem \#876}
\label{prob:876}

\noindent\textbf{Status:} OPEN \quad|\quad \textbf{Prize:} none \quad|\quad \textbf{Updated:} 2025-08-31

\noindent\textbf{Tags:} additive combinatorics

\medskip

\noindent\textbf{Statement:}

Let $A=\{a_1&#60;a_2&#60;\cdots\}\subset \mathbb{N}$ be an infinite sum-free set - that is, there are no solutions to\[a=b_1+\cdots+b_r\]with $b_1&#60;\cdots&#60;b_r&#60;a\in A$. How small can $a_{n+1}-a_n$ be? Is it possible that $a_{n+1}-a_n&#60;n$?




Erd\H{o}s \cite{Er98} writes that Graham 'recently proved' that there is such a sequence for which $a_{n+1}-a_n&lt;n^{1+o(1)}$, and that Melfi proved a somewhat weaker result. Erd\H{o}s \cite{Er62c} proved that a sum-free set has density zero. Deshouillers, Erd\H{o}s, and Melfi \cite{DEM99} constructed a sum-free set that grows like $a_n\sim n^{3+o(1)}$. Luczak and Schoen \cite{LuSc00} have proved that, for all large $N$,\[\lvert A\cap [1,N]\rvert\ll (N\log N)^{1/2},\]and that there exists a sum-free set $B$ such that\[\lvert B\cap [1,N]\rvert \gg \frac{N^{1/2}}{(\log N)^{1/2+o(1)}}\]for all large $N$. In \cite{Er75b} and \cite{Er77c} Erd\H{o}s asks to determine the maximum possible value of $\sum_{n\in A}\frac{1}{n}$. Erd\H{o}s had proved this is $&lt;100$, and Sullivan had shown that this is $&lt;4$, and Sullivan conjectured the maximum is slightly larger than $2$. See also [790] .




References

[DEM99] Deshouillers, Jean-Marc and Erd\H{o}s, Paul and Melfi, Giuseppe, On a question about sum-free sequences . Discrete Math. (1999), 49--54.

[Er62c] Erd\H{o}s, P\'{a}l, Some remarks on number theory. {III} . Mat. Lapok (1962), 28--38.

[Er75b] Erd\H{o}s, Paul, Problems and results in combinatorial number theory . Journ\'{e}es Arithm\'{e}tiques de Bordeaux (Conf., Univ. Bordeaux, Bordeaux, 1974) (1975), 295-310.

[Er77c] Erd\H{o}s, Paul, Problems and results on combinatorial number theory. III . Number theory day (Proc. Conf., Rockefeller Univ., New York, 1976) (1977), 43-72.

[Er98] Erd\H{o}s, Paul, Some of my new and almost new problems and results in combinatorial number theory . Number theory (Eger, 1996) (1998), 169-180.

[LuSc00] \L uczak, Tomasz and Schoen, Tomasz, On the maximal density of sum-free sets . Acta Arith. (2000), 225--229.

\noindent\rule{\linewidth}{0.4pt}


%% ============================================================
\section{Problem \#878}
\label{prob:878}

\noindent\textbf{Status:} OPEN \quad|\quad \textbf{Prize:} none \quad|\quad \textbf{Updated:} 2025-08-31

\noindent\textbf{Tags:} number theory

\medskip

\noindent\textbf{Statement:}

If $n=\prod_{1\leq i\leq t} p_i^{k_i}$ is the factorisation of $n$ into distinct primes then let\[f(n)=\sum p_i^{\ell_i},\]where $\ell_i$ is chosen such that $n\in [p_i^{\ell_i},p_i^{\ell_i+1})$. Furthermore, let\[F(n)=\max \sum_{i} a_i\]where the maximum is taken over all distinct $a_1,\ldots,a_k\leq n$ such that $(a_i,a_j)=1$ for $i\neq j$ and all prime factors of each $a_i$ are prime factors of $n$. Is it true that, for almost all $n$,\[f(n)=o(n\log\log n)\]and\[F(n) \gg n\log\log n?\]Is it true that\[\max_{n\leq x}f(n)\sim \frac{x\log x}{\log\log x}?\]Is it true that (for all $x$, or perhaps just for all large $x$)\[\max_{n\leq x}f(n)=\max_{n\leq x}F(n)?\]Find an asymptotic formula for the number of $n&#60;x$ such that $f(n)=F(n)$. Find an asymptotic formula for\[H(x)=\sum_{n&#60;x}\frac{f(n)}{n}.\]Is it true that\[H(x) \ll x\log\log\log\log x?\]




Erd\H{o}s \cite{Er84e} proved that\[\max_{n\leq x}f(n)\sim \frac{x\log x}{\log\log x}\]for a sequence of $x\to \infty$. It is trivial that $f(n)\leq F(n)$ for all $n$. It may be true that, for almost all $n$,\[F(n)\sim \frac{1}{2}n\log\log n.\]Erd\H{o}s notes that $f(n)/n$ 'almost behaves as a conventional additive function', but unusually $f(n)/n$ does not have a mean value - indeed,\[\limsup \frac{1}{x}\sum_{n&lt;x}\frac{f(n)}{n}=\infty\]but\[\liminf \frac{1}{x}\sum_{n&lt;x}\frac{f(n)}{n}&lt;\infty.\]Erd\H{o}s \cite{Er84e} proved that\[x\log\log\log\log x\ll H(x) \ll x\log\log\log x.\]Barreto has noted in the comments that this upper bound for $H(x)$ implies $f(n)=o(n\log\log n)$ for almost all $n$. Barreto has also noted that $\max_{n\leq x}f(n)=\max_{n\leq x}F(n)$ is false at $x=210$. See also [879] .




References

[Er84e] Erd\H{o}s, P., On two unconventional number theoretic functions and on some related problems . (1984), 113--121.

\noindent\rule{\linewidth}{0.4pt}


%% ============================================================
\section{Problem \#879}
\label{prob:879}

\noindent\textbf{Status:} OPEN \quad|\quad \textbf{Prize:} none \quad|\quad \textbf{Updated:} 2025-08-31

\noindent\textbf{Tags:} number theory

\medskip

\noindent\textbf{Statement:}

Call a set $S\subseteq \{1,\ldots,n\}$ admissible if $(a,b)=1$ for all $a\neq b\in S$. Let\[G(n) = \max_{S\subseteq \{1,\ldots,n\}} \sum_{a\in S}a\]and\[H(n)=\sum_{p&lt;n}p+ n\pi(n^{1/2}).\]Is it true that\[G(n) &gt;H(n)-n^{1+o(1)}?\]Is it true that, for every $k\geq 2$, if $n$ is sufficiently large then the admissible set which maximises $G(n)$ contains at least one integer with at least $k$ prime factors?




Erd\H{o}s and Van Lint proved that\[H(n)-n^{3/2-o(1)}&lt;G(n)&lt;H(n)\]and\[\frac{H(n)-G(n)}{n}\to \infty.\]They proved that $G(n)&gt;H(n)-n^{1+o(1)}$ assuming 'plausible (but hopeless) assumptions about the distribution of primes'. They also prove the second claim when $k=2$. See also [878] .

\noindent\rule{\linewidth}{0.4pt}


%% ============================================================
\section{Problem \#881}
\label{prob:881}

\noindent\textbf{Status:} OPEN \quad|\quad \textbf{Prize:} none \quad|\quad \textbf{Updated:} 2025-08-31

\noindent\textbf{Tags:} number theory, additive basis

\medskip

\noindent\textbf{Statement:}

Let $A\subset\mathbb{N}$ be an additive basis of order $k$ which is minimal, in the sense that if $B\subset A$ is any infinite set then $A\backslash B$ is not a basis of order $k$. Must there exist an infinite $B\subset A$ such that $A\backslash B$ is a basis of order $k+1$?

\noindent\rule{\linewidth}{0.4pt}


%% ============================================================
\section{Problem \#883}
\label{prob:883}

\noindent\textbf{Status:} OPEN \quad|\quad \textbf{Prize:} none \quad|\quad \textbf{Updated:} 2025-08-31

\noindent\textbf{Tags:} number theory, graph theory

\medskip

\noindent\textbf{Statement:}

For $A\subseteq \{1,\ldots,n\}$ let $G(A)$ be the graph with vertex set $A$, where two integers are joined by an edge if they are coprime. Is it true that if\[\lvert A\rvert &#62;\lfloor\tfrac{n}{2}\rfloor+\lfloor\tfrac{n}{3}\rfloor-\lfloor\tfrac{n}{6}\rfloor\]then $G(A)$ contains all odd cycles of length $\leq \frac{n}{3}+1$? Is it true that, for every $\ell\geq 1$, if $n$ is sufficiently large and\[\lvert A\rvert &#62;\lfloor\tfrac{n}{2}\rfloor+\lfloor\tfrac{n}{3}\rfloor-\lfloor\tfrac{n}{6}\rfloor\]then $G(A)$ must contain a complete $(1,\ell,\ell)$ triparite graph on $2\ell+1$ vertices?




A problem of Erd\H{o}s and S\'{a}rk\H{o}zy \cite{ErSa97}, who prove that if\[\lvert A\rvert &#62;\lfloor\tfrac{n}{2}\rfloor+\lfloor\tfrac{n}{3}\rfloor-\lfloor\tfrac{n}{6}\rfloor\]then $G(A)$ contains all odd cycles of length $\leq cn$ for some constant $c&#62;0$. This threshold is the best possible, since one could take $A$ to be the set of $m\leq n$ which are divisible by either $2$ or $3$, in which case $G(A)$ contains no triangles. The second question was solved by S\'{a}rk\"{o}zy \cite{Sa99} who proved that, for large $n$, if $\lvert A\rvert$ exceeds the given threshold then $G(A)$ contains a complete $(1,\ell,\ell)$ triparite graph with\[\ell \gg \frac{\log n}{\log\log n}.\]




References

[ErSa97] Erd\H{o}s, Paul and Sarkozy, Gabor N., On cycles in the coprime graph of integers . Electron. J. Combin. (1997), Research Paper 8, approx. 11.

[Sa99] S\'{a}rk\"ozy, G\'{a}bor N., Complete tripartite subgraphs in the coprime graph of integers . Discrete Math. (1999), 227--238.

\noindent\rule{\linewidth}{0.4pt}


%% ============================================================
\section{Problem \#885}
\label{prob:885}

\noindent\textbf{Status:} OPEN \quad|\quad \textbf{Prize:} none \quad|\quad \textbf{Updated:} 2025-08-31

\noindent\textbf{Tags:} number theory, divisors

\medskip

\noindent\textbf{Statement:}

For integer $n\geq 1$ we define the factor difference set of $n$ by\[D(n) = \{\lvert a-b\rvert : n=ab\}.\]Is it true that, for every $k\geq 1$, there exist integers $N_1&#60;\cdots&#60;N_k$ such that\[\lvert \cap_i D(N_i)\rvert \geq k?\]




A question of Erd\H{o}s and Rosenfeld \cite{ErRo97}, who proved this is true for $k=2$. Jim\'{e}nez-Urroz \cite{Ji99} proved this for $k=3$ and Bremner \cite{Br19} proved this for $k=4$.




References

[Br19] Bremner, Andrew, On a problem of Erd\H{o}s related to common factor differences . Int. J. Number Theory (2019), 1059--1068.

[ErRo97] Erd\H{o}s, Paul and Rosenfeld, Moshe, The factor-difference set of integers . Acta Arith. (1997), 353--359.

[Ji99] Jim\'{e}nez-Urroz, Jorge, A note on a conjecture of Erd\H{o}s and {R}osenfeld . J. Number Theory (1999), 140--143.

\noindent\rule{\linewidth}{0.4pt}


%% ============================================================
\section{Problem \#886}
\label{prob:886}

\noindent\textbf{Status:} OPEN \quad|\quad \textbf{Prize:} none \quad|\quad \textbf{Updated:} 2025-08-31

\noindent\textbf{Tags:} number theory, divisors

\medskip

\noindent\textbf{Statement:}

Let $\epsilon&#62;0$. Is it true that, for all large $n$, the number of divisors of $n$ in $(n^{1/2},n^{1/2}+n^{1/2-\epsilon})$ is $O_\epsilon(1)$?




Erd\H{o}s attributes this conjecture to Ruzsa. Erd\H{o}s and Rosenfeld \cite{ErRo97} proved that there are infinitely many $n$ such that there are four divisors of $n$ in $(n^{1/2},n^{1/2}+16n^{1/4})$. They also proved that, for any constant $C&gt;0$, all large $n$ have at most $1+C^2$ many divisors in\[[n^{1/2}, n^{1/2}+Cn^{1/4}].\]See also [887] .




References

[ErRo97] Erd\H{o}s, Paul and Rosenfeld, Moshe, The factor-difference set of integers . Acta Arith. (1997), 353--359.

\noindent\rule{\linewidth}{0.4pt}


%% ============================================================
\section{Problem \#887}
\label{prob:887}

\noindent\textbf{Status:} OPEN \quad|\quad \textbf{Prize:} none \quad|\quad \textbf{Updated:} 2025-08-31

\noindent\textbf{Tags:} number theory, divisors

\medskip

\noindent\textbf{Statement:}

Is there an absolute constant $K$ such that, for every $C&#62;0$, if $n$ is sufficiently large then $n$ has at most $K$ divisors in $(n^{1/2},n^{1/2}+C n^{1/4})$.




A question of Erd\H{o}s and Rosenfeld \cite{ErRo97}, who proved that there are infinitely many $n$ with $4$ divisors in $(n^{1/2},n^{1/2}+n^{1/4})$, and ask whether $4$ is best possible here. They also proved that, if $n$ is large enough depending on $C$, then $n$ has at most $1+C^2$ divisors in $(n^{1/2},n^{1/2}+C n^{1/4})$. Chan \cite{Ch14} has resolved this when $n$ is a square, proving that if $n$ is a square then $n$ has at most $5$ divisors in\[[ n^{1/2}-n^{1/4}(\log n)^{1/7}, n^{1/2}+n^{1/4}(\log n)^{1/7}].\]In \cite{Ch15} Chan further proves that if $n=(N-a)(N-b)$ for some $0\leq a\leq b\leq \exp((\log n)^{2/7})$ then $n$ has at most $18$ divisors in\[[ n^{1/2}-n^{1/4}(\log n)^{1/14}, n^{1/2}+n^{1/4}(\log n)^{1/14}].\]See also [886] .




References

[Ch14] Chan, Tsz Ho, Factors of a perfect square . Acta Arith. (2014), 141--143.

[Ch15] Chan, Tsz Ho, Factors of almost squares and lattice points on circles . Int. J. Number Theory (2015), 1701--1708.

[ErRo97] Erd\H{o}s, Paul and Rosenfeld, Moshe, The factor-difference set of integers . Acta Arith. (1997), 353--359.

\noindent\rule{\linewidth}{0.4pt}


%% ============================================================
\section{Problem \#889}
\label{prob:889}

\noindent\textbf{Status:} OPEN \quad|\quad \textbf{Prize:} none \quad|\quad \textbf{Updated:} 2025-08-31

\noindent\textbf{Tags:} number theory

\medskip

\noindent\textbf{Statement:}

For $k\geq 0$ and $n\geq 1$ let $v(n,k)$ count the prime factors of $n+k$ which do not divide $n+i$ for $0\leq i&lt;k$. Equivalently, $v(n,k)$ counts the number of prime factors of $n+k$ which are $&gt;k$. Is it true that\[v_0(n)=\max_{k\geq 0}v(n,k)\to \infty\]as $n\to \infty$?




A question of Erd\H{o}s and Selfridge \cite{ErSe67}, who could only show that $v_0(n)\geq 2$ for $n\geq 17$. More generally, they conjecture that\[v_l(n)=\max_{k\geq l}v(n,k)\to \infty\]as $n\to \infty$, for every fixed $l$, but could not even prove that $v_1(n)\geq 2$ for all large $n$. This is problem B27 of Guy's collection \cite{Gu04}.




References

[ErSe67] Erd\H{o}s, P. and Selfridge, J. L., Some problems on the prime factors of consecutive integers . Illinois J. Math. (1967), 428--430.

[Gu04] Guy, Richard K., Unsolved problems in number theory . (2004), xviii+437.

\noindent\rule{\linewidth}{0.4pt}


%% ============================================================
\section{Problem \#890}
\label{prob:890}

\noindent\textbf{Status:} OPEN \quad|\quad \textbf{Prize:} none \quad|\quad \textbf{Updated:} 2025-08-31

\noindent\textbf{Tags:} number theory, primes

\medskip

\noindent\textbf{Statement:}

If $\omega_k(n)$ counts the number of distinct prime factors of $n$ which are $&#62;k$, then is it true that, for every $k\geq 1$,\[\liminf_{n\to \infty}\sum_{0\leq i&#60;k}\omega_k(n+i)\leq k?\]Is it true that\[\limsup_{n\to \infty}\left(\sum_{0\leq i&#60;k}\omega(n+i)\right) \frac{\log\log n}{\log n}=1,\]where $\omega$ counts the number of distinct prime factors without restriction?




A question of Erd\H{o}s and Selfridge \cite{ErSe67}, who observe that\[\liminf_{n\to \infty}\sum_{0\leq i&lt;k}\omega_k(n+i)\geq k-1\]for every $k$. This follows from P\'{o}lya's theorem that the set of $k$-smooth integers has unbounded gaps - indeed, provided $n$ is large, all but at most one of $n,n+1,\ldots,n+k-1$ has a prime factor $&gt;k$ by P\'{o}lya's theorem. It is a classical fact that\[\limsup_{n\to \infty}\omega(n)\frac{\log\log n}{\log n}=1.\](There appears to be an error in \cite{ErSe67}, who ask the first question with $\omega_k$ replaced by $\omega$, have a bound of $k+\pi(k)$ rather than $k$, but as explained by Meza and Tao in the comments this formulation is obviously false, and the version given here is most likely what was intended.)




References

[ErSe67] Erd\H{o}s, P. and Selfridge, J. L., Some problems on the prime factors of consecutive integers . Illinois J. Math. (1967), 428--430.

\noindent\rule{\linewidth}{0.4pt}


%% ============================================================
\section{Problem \#891}
\label{prob:891}

\noindent\textbf{Status:} OPEN \quad|\quad \textbf{Prize:} none \quad|\quad \textbf{Updated:} 2025-08-31

\noindent\textbf{Tags:} number theory

\medskip

\noindent\textbf{Statement:}

Let $2=p_1&lt;p_2&lt;\cdots$ be the primes and $k\geq 2$. Is it true that, for all sufficiently large $n$, there must exist an integer in $[n,n+p_1\cdots p_k)$ with $&gt;k$ many prime factors?




Schinzel deduced from P\'{o}lya's theorem \cite{Po18} (that the sequence of $k$-smooth integers has unbounded gaps) that this is true with $p_1\cdots p_k$ replaced by $p_1\cdots p_{k-1}p_{k+1}$. This is unknown even for $k=2$ - that is, is it true that in every interval of $6$ (sufficiently large) consecutive integers there must exist one with at least $3$ prime factors? Weisenberg has observed that Dickson's conjecture implies the answer is no if we replace $p_1\cdots p_k$ with $p_1\cdots p_k-1$. Indeed, let $L_k$ be the lowest common multiple of all integers at most $p_1\cdots p_k$. By Dickson's conjecture there are infinitely many $n'$ such that $\frac{L_k}{m}n'+1$ is prime for all $1\leq m&lt;p_1\cdots p_k$. It follows that, if $n=L_kn'+1$, then all integers in $[n,n+p_1\cdots p_k-1)$ have at most $k$ prime factors.




References

[Po18] P\'{o}lya, Georg, Zur arithmetischen {U}ntersuchung der {P}olynome . Math. Z. (1918), 143--148.

\noindent\rule{\linewidth}{0.4pt}


%% ============================================================
\section{Problem \#892}
\label{prob:892}

\noindent\textbf{Status:} OPEN \quad|\quad \textbf{Prize:} none \quad|\quad \textbf{Updated:} 2025-08-31

\noindent\textbf{Tags:} number theory, primitive sets

\medskip

\noindent\textbf{Statement:}

Is there a necessary and sufficient condition for a sequence of integers $b_1&#60;b_2&#60;\cdots$ that ensures there exists a primitive sequence $a_1&#60;a_2&#60;\cdots$ (i.e. no element divides another) with $a_n \ll b_n$ for all $n$? In particular, is this always possible if there are no non-trivial solutions to $(b_i,b_j)=b_k$? Similarly, find necessary and sufficient conditions on a sequence $n_1&#60;n_2&#60;\cdots$ that ensure there exists a primitive set $A$ such that\[\lvert A\cap [1,2^{n_i}]\rvert \gg 2^{n_i}\]for every $i$.




A problem of Erd\H{o}s, S\'{a}rk\"{o}zi, and Szemer\'{e}di \cite{ESS68}. It is known that\[\sum \frac{1}{b_n\log b_n}&lt;\infty\]and\[\sum_{b_n&lt;x}\frac{1}{b_n} =o\left(\frac{\log x}{\sqrt{\log\log x}}\right)\]are both necessary. (The former is due to Erd\H{o}s \cite{Er35}, the latter to Erd\H{o}s, S\'{a}rk\"{o}zy, and Szemer\'{e}di \cite{ESS67}.) One can ask a similar question for sequences of real numbers, as in [143] . In \cite{Er80} Erd\H{o}s suggests the first question is 'difficult and perhaps has no reasonable solution', and perhaps the the final question is more reasonable.




References

[ESS67] Erd\H{o}s, P. and S\'{a}rk\"ozy, A. and Szemer\'{e}di, E., On a theorem of Behrend . J. Austral. Math. Soc. (1967), 9--16.

[ESS68] Erd\H{o}s, P. and S\'{a}rk\"ozi, A. and Szemer\'{e}di, E., On the solvability of certain equations in sequences of positive upper logarithmic density . J. London Math. Soc. (1968), 71--78.

[Er35] Erd\H{o}s, Paul, Note on Sequences of Integers No One of Which is Divisible By Any Other . J. London Math. Soc. (1935), 126-128.

[Er80] Erd\H{o}s, Paul, A survey of problems in combinatorial number theory . Ann. Discrete Math. (1980), 89-115.

\noindent\rule{\linewidth}{0.4pt}


%% ============================================================
\section{Problem \#893}
\label{prob:893}

\noindent\textbf{Status:} OPEN \quad|\quad \textbf{Prize:} none \quad|\quad \textbf{Updated:} 2025-08-31

\noindent\textbf{Tags:} number theory, divisors

\medskip

\noindent\textbf{Statement:}

If $\tau(n)$ counts the divisors of $n$ then let\[f(n)=\sum_{1\leq k\leq n}\tau(2^k-1).\]Does $f(2n)/f(n)$ tend to a limit?




Erd\H{o}s \cite{Er98} says that 'probably there is no simple asymptotic formula for $f(n)$ since $f(n)$ increases too fast'. Kova\v{c} and Luca \cite{KoLu25} (building on a heuristic independently found by Cambie (personal communication)) have shown that there is no finite limit, in that\[\limsup_{n\to \infty}\frac{f(2n)}{f(n)}=\infty,\]and provide both theoretical and numerical evidence that suggests $\lim \frac{f(2n)}{f(n)}=\infty$.




References

[Er98] Erd\H{o}s, Paul, Some of my new and almost new problems and results in combinatorial number theory . Number theory (Eger, 1996) (1998), 169-180.

[KoLu25] V. Kova\v{c} and F. Luca, On the number of divisors of Mersenne numbers . arXiv:2506.04883 (2025).

\noindent\rule{\linewidth}{0.4pt}


%% ============================================================
\section{Problem \#901}
\label{prob:901}

\noindent\textbf{Status:} OPEN \quad|\quad \textbf{Prize:} none \quad|\quad \textbf{Updated:} 2025-08-31

\noindent\textbf{Tags:} combinatorics, hypergraphs

\medskip

\noindent\textbf{Statement:}

Let $m(n)$ be minimal such that there is an $n$-uniform hypergraph with $m(n)$ edges which is $3$-chromatic. Estimate $m(n)$.




In other words, the hypergraph does not have Property B . Property B means that there is a set $S$ which intersects all edges and yet does not contain any edge. It is known that $m(2)=3$, $m(3)=7$, and $m(4)=23$. Erd\H{o}s proved\[2^n \ll m(n) \ll n^2 2^n\](the lower bound in \cite{Er63b} and the upper bound in \cite{Er64e}). Erd\H{o}s conjectured that $m(n)/2^n\to \infty$, which was proved by Beck \cite{Be77}, who proved $m(n)\gg (\log n)2^n$, and later \cite{Be78} improved this to\[n^{1/3-o(1)}2^n \ll m(n).\]Radhakrishnan and Srinivasan \cite{RaSr00} improved this to\[\sqrt{\frac{n}{\log n}}2^n \ll m(n).\]Pluhar \cite{Pl09} gave a very short proof that $m(n) \gg n^{1/4}2^n$. In \cite{ErLo75} Erd\H{o}s and Lov\'{a}sz speculate that $n2^n$ is the correct order of magnitude for $m(n)$.




References

[Be77] Beck, J., On a combinatorial problem of {P}. {E}rd\H{o}s and {L}. {L}ov\'{a}sz . Discrete Math. (1977), 127--131.

[Be78] Beck, J., On {$3$}-chromatic hypergraphs . Discrete Math. (1978), 127--137.

[Er63b] Erd\H{o}s, P., On a combinatorial problem . Nordisk Mat. Tidskr. (1963), 5--10, 40.

[Er64e] Erd\H{o}s, P., On a combinatorial problem. {II} . Acta Math. Acad. Sci. Hungar. (1964), 445--447.

[ErLo75] Erd\H{o}s, P. and Lov\'{a}sz, L., Problems and results on {$3$}-chromatic hypergraphs and some related questions . (1975), 609--627.

[Pl09] Pluh\'{a}r, Andr\'{a}s, Greedy colorings of uniform hypergraphs . Random Structures Algorithms (2009), 216--221.

[RaSr00] Radhakrishnan, Jaikumar and Srinivasan, Aravind, Improved bounds and algorithms for hypergraph {$2$}-coloring . Random Structures Algorithms (2000), 4--32.

\noindent\rule{\linewidth}{0.4pt}


%% ============================================================
\section{Problem \#902}
\label{prob:902}

\noindent\textbf{Status:} OPEN \quad|\quad \textbf{Prize:} none \quad|\quad \textbf{Updated:} 2025-08-31

\noindent\textbf{Tags:} graph theory

\medskip

\noindent\textbf{Statement:}

Let $f(n)$ be minimal such that there is a tournament (a complete directed graph) on $f(n)$ vertices such that every set of $n$ vertices is dominated by at least one other vertex. Estimate $f(n)$.




Sch\"{u}tte asked Erd\H{o}s this in the early 1960s. It is easy to check that $f(1)=3$ and $f(2)=7$. Erd\H{o}s \cite{Er63c} proved\[2^{n+1}-1 \leq f(n) \ll n^22^n.\]Szekeres and Szekeres \cite{SzSz65} proved that $f(3)=19$ and\[n2^n \ll f(n).\]




References

[Er63c] Erd\H{o}s, P., On a problem in graph theory . Math. Gaz. (1963), 220--223.

[SzSz65] Szekeres, E. and Szekeres, G., On a problem of Sch\"{u}tte and {E}rd\H{o}s . Math. Gaz. (1965), 290--293.

\noindent\rule{\linewidth}{0.4pt}


%% ============================================================
\section{Problem \#906}
\label{prob:906}

\noindent\textbf{Status:} OPEN \quad|\quad \textbf{Prize:} none \quad|\quad \textbf{Updated:} 2025-08-31

\noindent\textbf{Tags:} analysis, iterated functions

\medskip

\noindent\textbf{Statement:}

Is there an entire non-zero function $f:\mathbb{C}\to \mathbb{C}$ such that, for any infinite sequence $n_1&#60;n_2&#60;\cdots$, the set\[\{ z: f^{(n_k)}(z)=0 \textrm{ for some }k\geq 1\}\]is everywhere dense?




Erd\H{o}s \cite{Er82e} writes that this was solved in the affirmative 'more than ten years ago', but gives no reference or indication who solved it. From context he seems to attribute this to Barth and Schneider \cite{BaSc72}, but this paper contains no such result. Tang points out that the problem is trivial if we take $f$ to be a polynomial, so presumably it is intended the function $f$ is transcendental.




References

[BaSc72] Barth, K. F. and Schneider, W. J., On a problem of Erd\H{o}s concerning the zeros of the derivatives of an entire function . Proc. Amer. Math. Soc. (1972), 229--232.

[Er82e] Erd\H{o}s, Paul, Some of my favourite problems which recently have been solved . (1982), 59--79.

\noindent\rule{\linewidth}{0.4pt}


%% ============================================================
\section{Problem \#911}
\label{prob:911}

\noindent\textbf{Status:} OPEN \quad|\quad \textbf{Prize:} none \quad|\quad \textbf{Updated:} 2025-08-31

\noindent\textbf{Tags:} graph theory, ramsey theory

\medskip

\noindent\textbf{Statement:}

Let $\hat{R}(G)$ denote the size Ramsey number, the minimal number of edges $m$ such that there is a graph $H$ with $m$ edges that is Ramsey for $G$. Is there a function $f$ such that $f(x)/x\to \infty$ as $x\to \infty$ such that, for all large $C$, if $G$ is a graph with $n$ vertices and $e\geq Cn$ edges then\[\hat{R}(G) &#62; f(C) e?\]

\noindent\rule{\linewidth}{0.4pt}


%% ============================================================
\section{Problem \#912}
\label{prob:912}

\noindent\textbf{Status:} OPEN \quad|\quad \textbf{Prize:} none \quad|\quad \textbf{Updated:} 2025-08-31

\noindent\textbf{Tags:} number theory, factorials

\medskip

\noindent\textbf{Statement:}

If\[n! = \prod_i p_i^{k_i}\]is the factorisation into distinct primes then let $h(n)$ count the number of distinct exponents $k_i$. Prove that there exists some $c&#62;0$ such that\[h(n) \sim c \left(\frac{n}{\log n}\right)^{1/2}\]as $n\to \infty$.




A problem of Erd\H{o}s and Selfridge, who proved (see \cite{Er82c})\[h(n) \asymp \left(\frac{n}{\log n}\right)^{1/2}.\]A heuristic of Tao using the Cram\'{e}r model for the primes (detailed in the comments) suggests this is true with\[c=\sqrt{2\pi}=2.506\cdots.\]




References

[Er82c] Erd\H{o}s, P., Miscellaneous problems in number theory . Congr. Numer. (1982), 25-45.

\noindent\rule{\linewidth}{0.4pt}


%% ============================================================
\section{Problem \#913}
\label{prob:913}

\noindent\textbf{Status:} OPEN \quad|\quad \textbf{Prize:} none \quad|\quad \textbf{Updated:} 2025-08-31

\noindent\textbf{Tags:} number theory

\medskip

\noindent\textbf{Statement:}

Are there infinitely many $n$ such that if\[n(n+1) = \prod_i p_i^{k_i}\]is the factorisation into distinct primes then all exponents $k_i$ are distinct?




It is likely that there are infinitely many primes $p$ such that $8p^2-1$ is also prime, in which case this is true with exponents $\{1,2,3\}$, letting $n=8p^2-1$. This problem has been formalised in Lean as part of the Google DeepMind Formal Conjectures project .

\noindent\rule{\linewidth}{0.4pt}


%% ============================================================
\section{Problem \#917}
\label{prob:917}

\noindent\textbf{Status:} OPEN \quad|\quad \textbf{Prize:} none \quad|\quad \textbf{Updated:} 2025-08-31

\noindent\textbf{Tags:} graph theory, chromatic number

\medskip

\noindent\textbf{Statement:}

Let $k\geq 4$ and $f_k(n)$ be the largest number of edges in a graph on $n$ vertices which has chromatic number $k$ and is critical (i.e. deleting any edge reduces the chromatic number). Is it true that\[f_k(n) \gg_k n^2?\]Is it true that\[f_6(n)\sim n^2/4?\]More generally, is it true that, for $k\geq 6$,\[f_k(n) \sim \frac{1}{2}\left(1-\frac{1}{\lfloor k/3\rfloor}\right)n^2?\]




Erd\H{o}s \cite{Er93} wrote 'I learned of this definition from Dirac in 1949 and immediately asked whether $f_k(n)=o(n^2)$. To my great surprise Dirac constructed a $6$ critical graph on $n$ vertices with more than $\frac{n^2}{4}$ edges.' In fact Dirac \cite{Di52} proved\[f_6(4n+2) \geq 4n^2+8n+3,\]as witnessed by taking two disjoint copies of $C_{2n+1}$ and adding all edges between them. Erd\H{o}s \cite{Er69b} observed that Dirac's construction generalises to show that, if $3\mid k$, there are infinitely many values of $n$ (those of the shape $mk/3$ where $m$ is odd) such that\[f_k(n) \geq \frac{1}{2}\left(1-\frac{1}{k/3}\right)n^2 + n.\]Toft \cite{To70} proved that $f_k(n)\gg_k n^2$ for $k\geq 4$. Constructions of Stiebitz \cite{St87} show that, for $k\geq 6$, there exist infinitely many values of $n$ such that\[f_k(n) \geq \frac{1}{2}\left(1-\frac{1}{\lfloor k/3\rfloor+\delta_k}\right)n^2\]where $\delta_k=0$ if $k\equiv 0\pmod{3}$, $\delta_k=1/7$ if $k\equiv 1\pmod{3}$, and $\delta_k\equiv 24/69$ if $k\equiv 2\pmod{3}$, which disproves Erd\H{o}s' conjectured asympotic for $k\not\equiv 0\pmod{3}$. Stiebitz also proved the general upper bound\[f_k(n) &lt; \mathrm{ex}(n;K_{k-1})\sim \frac{1}{2}\left(1-\frac{1}{k-2}\right)n^2\]for large $n$. Luo, Ma, and Yang \cite{LMY23} have improved this upper bound to\[f_k(n) \leq \frac{1}{2}\left(1-\frac{1}{k-2}-\frac{1}{36(k-1)^2}+o(1)\right)n^2\]See also [944] and [1032] .




References

[Di52] Dirac, G. A., A property of {$4$}-chromatic graphs and some remarks on critical graphs . J. London Math. Soc. (1952), 85-92.

[Er69b] Erd\H{o}s, P., Problems and results in chromatic graph theory . Proof Techniques in Graph Theory (Proc. Second Ann Arbor Graph Theory Conf., Ann Arbor, Mich., 1968) (1969), 27-35.

[Er93] Erd\H{o}s, Paul, Some of my favorite solved and unsolved problems in graph theory . Quaestiones Math. (1993), 333-350.

[LMY23] Luo, Cong and Ma, Jie and Yang, Tianchi, On the maximum number of edges in {$k$}-critical graphs . Combin. Probab. Comput. (2023), 900--911.

[St87] Stiebitz, M., Subgraphs of colour-critical graphs . Combinatorica (1987), 303--312.

[To70] Toft, B., On the maximal number of edges of critical {$k$}-chromatic graphs . Studia Sci. Math. Hungar. (1970), 461--470.

\noindent\rule{\linewidth}{0.4pt}


%% ============================================================
\section{Problem \#918}
\label{prob:918}

\noindent\textbf{Status:} OPEN \quad|\quad \textbf{Prize:} none \quad|\quad \textbf{Updated:} 2025-08-31

\noindent\textbf{Tags:} graph theory, chromatic number

\medskip

\noindent\textbf{Statement:}

Is there a graph with $\aleph_2$ vertices and chromatic number $\aleph_2$ such that every subgraph on $\aleph_1$ vertices has chromatic number $\leq\aleph_0$? Is there a graph with $\aleph_{\omega+1}$ vertices and chromatic number $\aleph_1$ such that every subgraph on $\aleph_\omega$ vertices has chromatic number $\leq\aleph_0$?




A question of Erd\H{o}s and Hajnal \cite{ErHa68b}, who proved that for every finite $k$ there is a graph with chromatic number $\aleph_1$ where each subgraph on less than $\aleph_k$ vertices has chromatic number $\leq \aleph_0$. In \cite{Er69b} it is asked with chromatic number $=\aleph_0$, but in the comments louisd observes this is (assuming subgraph and not induced subgraph was intended) trivially impossible, and hence presumably the problem was intended as written here (which is how it is posed in \cite{ErHa68b}).




References

[Er69b] Erd\H{o}s, P., Problems and results in chromatic graph theory . Proof Techniques in Graph Theory (Proc. Second Ann Arbor Graph Theory Conf., Ann Arbor, Mich., 1968) (1969), 27-35.

[ErHa68b] Erd\H{o}s, P. and Hajnal, A., On chromatic number of infinite graphs . (1968), 83--98.

\noindent\rule{\linewidth}{0.4pt}


%% ============================================================
\section{Problem \#919}
\label{prob:919}

\noindent\textbf{Status:} OPEN \quad|\quad \textbf{Prize:} none \quad|\quad \textbf{Updated:} 2025-08-31

\noindent\textbf{Tags:} graph theory, chromatic number

\medskip

\noindent\textbf{Statement:}

Is there a graph $G$ with vertex set $\omega_2^2$ and chromatic number $\aleph_2$ such that every subgraph whose vertices have a lesser type has chromatic number $\leq \aleph_0$? What if instead we ask for $G$ to have chromatic number $\aleph_1$?




This question was inspired by a theorem of Babai, that if $G$ is a graph on a well-ordered set with chromatic number $\geq \aleph_0$ there is a subgraph on vertices with order-type $\omega$ with chromatic number $\aleph_0$. Erd\H{o}s and Hajnal showed this does not generalise to higher cardinals - they (see \cite{Er69b}) constructed a set on $\omega_1^2$ with chromatic number $\aleph_1$ such that every strictly smaller subgraph has chromatic number $\leq \aleph_0$ as follows: the vertices of $G$ are the pairs $(x_\alpha,y_\beta)$ for $1\leq \alpha,\beta &#60;\omega_1$, ordered lexicographically. We connect $(x_{\alpha_1},y_{\beta_1})$ and $(x_{\alpha_2},y_{\beta_2})$ if and only if $\alpha_1&#60;\alpha_2$ and $\beta_1&#60;\beta_2$. A similar construction produces a graph on $\omega_2^2$ with chromatic number $\aleph_2$ such that every smaller subgraph has chromatic number $\leq \aleph_1$.




References

[Er69b] Erd\H{o}s, P., Problems and results in chromatic graph theory . Proof Techniques in Graph Theory (Proc. Second Ann Arbor Graph Theory Conf., Ann Arbor, Mich., 1968) (1969), 27-35.

\noindent\rule{\linewidth}{0.4pt}


%% ============================================================
\section{Problem \#920}
\label{prob:920}

\noindent\textbf{Status:} OPEN \quad|\quad \textbf{Prize:} none \quad|\quad \textbf{Updated:} 2025-08-31

\noindent\textbf{Tags:} graph theory, chromatic number

\medskip

\noindent\textbf{Statement:}

Let $f_k(n)$ be the maximum possible chromatic number of a graph with $n$ vertices which contains no $K_k$. Is it true that, for $k\geq 4$,\[f_k(n) \gg \frac{n^{1-\frac{1}{k-1}}}{(\log n)^{c_k}}\]for some constant $c_k&#62;0$?




Graver and Yackel \cite{GrYa68} proved that\[f_k(n) \ll \left(n\frac{\log\log n}{\log n}\right)^{1-\frac{1}{k-1}}.\]It is known that $f_3(n)\asymp (n/\log n)^{1/2}$ (see [1104] ). The lower bound $R(4,m) \gg m^3/(\log m)^4$ of Mattheus and Verstraete \cite{MaVe23} (see [166] ) implies\[f_4(n) \gg \frac{n^{2/3}}{(\log n)^{4/3}}.\]A positive answer to this question would follow from [986] . The known bounds for that problem imply, for $k\geq 5$,\[f_k(n) \gg \frac{n^{1-\frac{1}{k-2}}}{(\log n)^{c_k}}.\]See [1104] (and also [1013] ) for the case $k=3$.




References

[GrYa68] Graver, Jack E. and Yackel, James, Some graph theoretic results associated with Ramsey's theorem . J. Combinatorial Theory (1968), 125--175.

[MaVe23] Mattheus, S. and Verstraete, J., The asymptotics of $r(4,t)$ . arXiv:2306.04007 (2023).

\noindent\rule{\linewidth}{0.4pt}


%% ============================================================
\section{Problem \#928}
\label{prob:928}

\noindent\textbf{Status:} OPEN \quad|\quad \textbf{Prize:} none \quad|\quad \textbf{Updated:} 2025-09-04

\noindent\textbf{Tags:} number theory

\medskip

\noindent\textbf{Statement:}

Let $\alpha,\beta\in (0,1)$ and let $P(n)$ denote the largest prime divisor of $n$. Does the density of integers $n$ such that $P(n)&#60;n^{\alpha}$ and $P(n+1)&#60;(n+1)^\beta$ exist?




Dickman \cite{Di30} showed the density of smooth $n$, with largest prime factor $&lt;n^\alpha$, is $\rho(1/\alpha)$ where $\rho$ is the Dickman function . Erd\H{o}s also asked whether infinitely many such $n$ even exist, but Meza has observed that this follows immediately from Schinzel's result \cite{Sc67b} that for infinitely many $n$ the largest prime factor of $n(n+1)$ is at most $n^{O(1/\log\log n)}$. Erd\H{o}s asked whether the events $P(n)&lt;n^\alpha$ and $P(n+1)&lt;(n+1)^\beta$ are independent, in the sense that the density of $n$ satisfying both conditions is equal to $\rho(1/\alpha)\rho(1/\beta)$. Ter\"{a}v\"{a}inen \cite{Te18} has proved the logarithmic density exists and is equal to $\rho(1/\alpha)\rho(1/\beta)$. Wang \cite{Wa21} has proved the density is $\rho(1/\alpha)\rho(1/\beta)$ assuming the Elliott-Halberstam conjecture for friable integers. In \cite{Er76e} Erd\H{o}s conjectured that the density of the set of $n$ such that all of $n,\ldots,n+r-1$ are $n^{\alpha}$ smooth (i.e. $P(n+i)\leq n^{\alpha}$ for $0\leq i&lt;r$) is $\rho(1/\alpha)^r$. See also [370] .




References

[Di30] K. Dickman, On the frequency of numbers containing prime factors of a certain relative magnitude . Ark. Mat. Astr. Fys. (1930), 1-14.

[Er76e] Erd\H{o}s, P., Problems and results on consecutive integers . Publ. Math. Debrecen (1976), 271-282.

[Sc67b] Schinzel, A., On two theorems of Gelfond and some of their applications . Acta Arith. (1967/68), 177-236.

[Te18] Ter\"{a}v\"{a}inen, Joni, On binary correlations of multiplicative functions . Forum Math. Sigma (2018), Paper No. e10, 41.

[Wa21] Wang, Zhiwei, Three conjectures on {$P^+(n)$} and {$P^+(n+1)$} hold under the {E}lliott-{H}alberstam conjecture for friable integers . J. Number Theory (2021), 1--11.

\noindent\rule{\linewidth}{0.4pt}


%% ============================================================
\section{Problem \#929}
\label{prob:929}

\noindent\textbf{Status:} OPEN \quad|\quad \textbf{Prize:} none \quad|\quad \textbf{Updated:} 2025-08-31

\noindent\textbf{Tags:} number theory

\medskip

\noindent\textbf{Statement:}

Let $k\geq 2$ be large and let $S(k)$ be the minimal $x$ such that there is a positive density set of $n$ where\[n+1,n+2,\ldots,n+k\]are all divisible by primes $\leq x$. Estimate $S(k)$ - in particular, is it true that $S(k)\geq k^{1-o(1)}$?




It follows from Rosser's sieve that $S(k)&gt; k^{1/2-o(1)}$. It is trivial that $S(k)\leq k+1$ since, for example, one can take $n\equiv 1\pmod{(k+1)!}$. The best bound on large gaps between primes due to Ford, Green, Konyagin, Maynard, and Tao \cite{FGKMT18} (see [4] ) implies\[S(k) \ll k \frac{\log\log\log k}{\log\log k\log\log\log\log k}.\]




References

[FGKMT18] Ford, Kevin and Green, Ben and Konyagin, Sergei and Maynard, James and Tao, Terence, Long gaps between primes . J. Amer. Math. Soc. (2018), 65-105.

\noindent\rule{\linewidth}{0.4pt}


%% ============================================================
\section{Problem \#930}
\label{prob:930}

\noindent\textbf{Status:} OPEN \quad|\quad \textbf{Prize:} none \quad|\quad \textbf{Updated:} 2025-08-31

\noindent\textbf{Tags:} number theory

\medskip

\noindent\textbf{Statement:}

Is it true that, for every $r$, there is a $k$ such that if $I_1,\ldots,I_r$ are disjoint intervals of consecutive integers, all of length at least $k$, then\[\prod_{1\leq i\leq r}\prod_{m\in I_i}m\]is not a perfect power?




Erd\H{o}s and Selfridge \cite{ErSe75} proved that the product of consecutive integers is never a power (establishing the case $r=1$). The condition that the intervals be large in terms of $r$ is necessary for $r=2$ - see the constructions in [363] . See also [363] for the case of squares.




References

[ErSe75] Erd\H{o}s, P. and Selfridge, J. L., The product of consecutive integers is never a power . Illinois J. Math. (1975), 292-301.

\noindent\rule{\linewidth}{0.4pt}


%% ============================================================
\section{Problem \#931}
\label{prob:931}

\noindent\textbf{Status:} OPEN \quad|\quad \textbf{Prize:} none \quad|\quad \textbf{Updated:} 2025-08-31

\noindent\textbf{Tags:} number theory

\medskip

\noindent\textbf{Statement:}

Let $k_1\geq k_2\geq 3$. Are there only finitely many $n_2\geq n_1+k_1$ such that\[\prod_{1\leq i\leq k_1}(n_1+i)\textrm{ and }\prod_{1\leq j\leq k_2}(n_2+j)\]have the same prime factors?




Tijdeman gave the example\[19,20,21,22\textrm{ and }54,55,56,57.\]Erd\H{o}s \cite{Er76d} was unsure of this conjecture, and thought perhaps if the two products have the same prime factors then $n_2&gt;2(n_1+k_1)$. It is not clear but it is possible that he meant to ask this question also permitting finitely many counterexamples. Indeed, without this caveat it is false - AlphaProof has found the counterexample\[10! = 2^8\cdot 3^4\cdot 5^2\cdot 7\]and\[14\cdot 15\cdot 16 = 2^5\cdot 3\cdot 5\cdot 7,\]so that $n_1=0$, $k_1=10$, $n_2=13$, and $k_2=3$. See also [388] . This is discussed in problem B35 of Guy's collection \cite{Gu04}.




References

[Er76d] Erd\H{o}s, P., Problems and results on number theoretic properties of consecutive integers and related questions . Proceedings of the Fifth Manitoba Conference on Numerical Mathematics (Univ. Manitoba, Winnipeg, Man., 1975) (1976), 25-44.

[Gu04] Guy, Richard K., Unsolved problems in number theory . (2004), xviii+437.

\noindent\rule{\linewidth}{0.4pt}


%% ============================================================
\section{Problem \#932}
\label{prob:932}

\noindent\textbf{Status:} OPEN \quad|\quad \textbf{Prize:} none \quad|\quad \textbf{Updated:} 2025-08-31

\noindent\textbf{Tags:} number theory

\medskip

\noindent\textbf{Statement:}

Let $p_k$ denote the $k$th prime. For infinitely many $r$ there are at least two integers $p_r&#60;n&#60;p_{r+1}$ all of whose prime factors are $&#60;p_{r+1}-p_r$.




Erd\H{o}s thought this was true but that there are very few such $r$. He could show that the density of $r$ such that at least one such $n$ exist is $0$. This problem has been formalised in Lean as part of the Google DeepMind Formal Conjectures project .

\noindent\rule{\linewidth}{0.4pt}


%% ============================================================
\section{Problem \#933}
\label{prob:933}

\noindent\textbf{Status:} OPEN \quad|\quad \textbf{Prize:} none \quad|\quad \textbf{Updated:} 2025-08-31

\noindent\textbf{Tags:} number theory

\medskip

\noindent\textbf{Statement:}

If $n(n+1)=2^k3^lm$, where $(m,6)=1$, then is it true that\[\limsup_{n\to \infty} \frac{2^k3^l}{n\log n}=\infty?\]




Mahler proved (a more general result that implies in particular) that\[2^k3^l&lt;n^{1+o(1)}.\]Erd\H{o}s \cite{Er76d} wrote 'it is easy to see' that for infinitely many $n$\[2^k3^l&gt;n\log n.\]Steinerberger has noted a simple proof of this fact follows from taking $n=2^{3^r}$ for any integer $r\geq 1$, when $k=3^r$ and $l=r+1$.




References

[Er76d] Erd\H{o}s, P., Problems and results on number theoretic properties of consecutive integers and related questions . Proceedings of the Fifth Manitoba Conference on Numerical Mathematics (Univ. Manitoba, Winnipeg, Man., 1975) (1976), 25-44.

\noindent\rule{\linewidth}{0.4pt}


%% ============================================================
\section{Problem \#934}
\label{prob:934}

\noindent\textbf{Status:} OPEN \quad|\quad \textbf{Prize:} none \quad|\quad \textbf{Updated:} 2025-08-31

\noindent\textbf{Tags:} graph theory

\medskip

\noindent\textbf{Statement:}

Let $h_t(d)$ be minimal such that every graph $G$ with $h_t(d)$ edges and maximal degree $\leq d$ contains two edges whose shortest path between them has length $\geq t$. Estimate $h_t(d)$.




A problem of Erd\H{o}s and Ne\v{s}et\v{r}il. Erd\H{o}s \cite{Er88} wrote 'This problem seems to be interesting only if there is a nice expression for $h_t(d)$.' It is easy to see that $h_t(d)\leq 2d^t$ always and $h_1(d)=d+1$. Erd\H{o}s and Ne\v{s}et\v{r}il and Bermond, Bond, Paoli, and Peyrat \cite{BBPP83} independently conjectured that $h_2(d) \leq \tfrac{5}{4}d^2+1$, with equality for even $d$ (see [149] ). This was proved by Chung, Gy\'{a}rf\'{a}s, Tuza, and Trotter \cite{CGTT90}. Cambie, Cames van Batenburg, de Joannis de Verclos, and Kang \cite{CCJK22} conjectured that\[h_3(d) \leq d^3-d^2+d+2,\]with equality if and only if $d=p^k+1$ for some prime power $p^k$, and proved that $h_3(3)=23$. They also conjecture that, for all $t\geq 3$, $h_t(d)\geq (1-o(1))d^t$ for infinitely many $d$ and $h_t(d)\leq (1+o(1))d^t$ for all $d$ (where the $o(1)$ term $\to 0$ as $d\to \infty$). The same authors prove that, if $t$ is large, then there are infinitely many $d$ such that $h_t(d) \geq 0.629^td^t$, and that for all $t\geq 1$ we have\[h_t(d) \leq \tfrac{3}{2}d^t+1.\]




References

[BBPP83] Bermond, J.-C. and Bond, J. and Paoli, M. and Peyrat, C., Graphs and interconnection networks: diameter and vulnerability . (1983), 1--30.

[CCJK22] Cambie, Stijn and Cames van Batenburg, Wouter and de Joannis de Verclos, R\'{e}mi and Kang, Ross J., Maximizing line subgraphs of diameter at most {$t$} . SIAM J. Discrete Math. (2022), 939--950.

[CGTT90] Chung, F. R. K. and Gy\'{a}rf\'{a}s, A. and Tuza, Z. and Trotter, W. T., The maximum number of edges in {$2K_2$}-free graphs of bounded degree . Discrete Math. (1990), 129--135.

[Er88] Erd\H{o}s, P, Problems and results in combinatorial analysis and graph theory . Discrete Math. (1988), 81-92.

\noindent\rule{\linewidth}{0.4pt}


%% ============================================================
\section{Problem \#935}
\label{prob:935}

\noindent\textbf{Status:} OPEN \quad|\quad \textbf{Prize:} none \quad|\quad \textbf{Updated:} 2025-09-04

\noindent\textbf{Tags:} number theory

\medskip

\noindent\textbf{Statement:}

For any integer $n=\prod p^{k_p}$ let $Q_2(n)$ be the powerful part of $n$, so that\[Q_2(n) = \prod_{\substack{p\\ k_p\geq 2}}p^{k_p}.\]Is it true that, for every $\epsilon&#62;0$ and $\ell\geq 1$, if $n$ is sufficiently large then\[Q_2(n(n+1)\cdots(n+\ell))&#60;n^{2+\epsilon}?\]If $\ell\geq 2$ then is\[\limsup_{n\to \infty}\frac{Q_2(n(n+1)\cdots(n+\ell))}{n^2}\]infinite? If $\ell\geq 2$ then is\[\lim_{n\to \infty}\frac{Q_2(n(n+1)\cdots(n+\ell))}{n^{\ell+1}}=0?\]




Erd\H{o}s \cite{Er76d} writes that if this is true then it 'seems very difficult to prove'. A result of Mahler implies, for every $\ell\geq 1$,\[\limsup_{n\to \infty}\frac{Q_2(n(n+1)\cdots(n+\ell))}{n^2}\geq 1.\]All these questions can be asked replacing $Q_2$ by $Q_r$ for $r&gt;2$, only keeping those prime powers with exponent $\geq r$. The second part of this problem is essentially the same (up to constants) as [367] . The construction given there by van Doorn (also given by Aletheia \cite{Fe26}) proves the second part in the affirmative, since a construction via solutions to the Pell equation $x^2-8y^2=1$ implies\[\limsup_{n\to \infty}\frac{Q_2(n(n+1)(n+2))}{n^2}=\infty.\]In \cite{Fe26} it is similarly noted that the ABC conjecture implies a positive answer to the third question.




References

[Er76d] Erd\H{o}s, P., Problems and results on number theoretic properties of consecutive integers and related questions . Proceedings of the Fifth Manitoba Conference on Numerical Mathematics (Univ. Manitoba, Winnipeg, Man., 1975) (1976), 25-44.

[Fe26] T. Feng et al, Semi-Autonomous Mathematics Discovery with Gemini: A Case Study on the Erd\H{o}s Problems . arXiv:2601.22401 (2026).

\noindent\rule{\linewidth}{0.4pt}


%% ============================================================
\section{Problem \#936}
\label{prob:936}

\noindent\textbf{Status:} OPEN \quad|\quad \textbf{Prize:} none \quad|\quad \textbf{Updated:} 2025-08-31

\noindent\textbf{Tags:} number theory

\medskip

\noindent\textbf{Statement:}

Are\[2^n\pm 1\]and\[n!\pm 1\]powerful (i.e. if $p\mid m$ then $p^2\mid m$) for only finitely many $n$?




Cushing and Pascoe \cite{CuPa16} have shown the answer to the second question is yes assuming the abc conjecture - in fact, for any fixed $k\geq 0$, there are only finitely many $n$ and powerful $x$ such that $\lvert x-n!\rvert \leq k$. CrowdMath \cite{Cr20} has shown that the answer to the first question is yes, again assuming the abc conjecture.




References

[Cr20] P. A. CrowdMath, Applications of the abc conjecture to powerful numbers . arXiv:2005.07321 (2020).

[CuPa16] D. Cushing and J. E. Pascoe, Powerful numbers and the ABC-conjecture . arXiv:1611.01192 (2016).

\noindent\rule{\linewidth}{0.4pt}


%% ============================================================
\section{Problem \#938}
\label{prob:938}

\noindent\textbf{Status:} OPEN \quad|\quad \textbf{Prize:} none \quad|\quad \textbf{Updated:} 2025-08-31

\noindent\textbf{Tags:} number theory

\medskip

\noindent\textbf{Statement:}

Let $A=\{n_1&#60;n_2&#60;\cdots\}$ be the sequence of powerful numbers (if $p\mid n$ then $p^2\mid n$). Are there only finitely many three-term progressions of consecutive terms $n_k,n_{k+1},n_{k+2}$?




Erd\H{o}s also conjectured (see [364] ) that there are no triples of powerful numbers of the shape $n,n+1,n+2$.

\noindent\rule{\linewidth}{0.4pt}


%% ============================================================
\section{Problem \#939}
\label{prob:939}

\noindent\textbf{Status:} OPEN \quad|\quad \textbf{Prize:} none \quad|\quad \textbf{Updated:} 2025-08-31

\noindent\textbf{Tags:} number theory

\medskip

\noindent\textbf{Statement:}

Let $r\geq 2$. An $r$-powerful number $n$ is one such that if $p\mid n$ then $p^r\mid n$. If $r\geq 4$ then can the sum of $r-2$ coprime $r$-powerful numbers ever be itself $r$-powerful? Are there at most finitely many such solutions? Are there infinitely many triples of coprime $3$-powerful numbers $a,b,c$ such that $a+b=c$?




The answer to the third question is yes: Nitaj \cite{Ni95} has proved that there are infinitely many triples of coprime $3$-powerful numbers $a,b,c$ such that $a+b=c$, such as\[2^3\cdot 3^5\cdot 73^3+271^3 = 919^3.\]In Nitaj's construction at least two of $a,b,c$ are perfect cubes. Cohn \cite{Co98} constructed infinitely many such triples, none of which are perfect cubes. An alternative construction was given by Walsh \cite{Wa24}. Euler had conjectured that the sum of $k-1$ many $k$th powers is never a $k$th power, but this is false for $k=5$, as Lander and Parkin \cite{LaPa67} found\[27^5+84^5+110^5+133^5=144^5.\]Cambie has found several examples of the sum of $r-2$ coprime $r$-powerful numbers being itself $r$-powerful. For example when $r=5$\[3^761^5=2^83^{10}5^7+2^{12}23^6+11^513^5.\](Kitamura has given another example in the comments.) Cambie has also found solutions when $r=7$ or $r=8$ (the latter even with the sum of $5$ $8$-powerful numbers being $8$-powerful). GPT-5.5. Pro (prompted by Price) has given a simple construction in the comments that shows for $r\geq 6$ there are infinitely many $r$-powerful numbers which are the sum of $(r-2)$ (in fact $\lceil r/2\rceil+1$) many coprime $r$-powerful numbers.




References

[Co98] Cohn, J. H. E., A conjecture of {E}rd\H{o}s on {$3$}-powerful numbers . Math. Comp. (1998), 439--440.

[LaPa67] Lander, L. J. and Parkin, T. R., A counterexample to {E}uler's sum of powers conjecture . Math. Comp. (1967), 101--103.

[Ni95] Nitaj, Abderrahmane, On a conjecture of {E}rd\H{o}s on {$3$}-powerful numbers . Bull. London Math. Soc. (1995), 317--318.

[Wa24] P. Walsh, A question of Erd\H{o}s on 3-powerful numbers and an elliptic curve analogue of the Ankeny-Artin-Chowla conjecture . arXiv:2404.039701 (2024).

\noindent\rule{\linewidth}{0.4pt}


%% ============================================================
\section{Problem \#940}
\label{prob:940}

\noindent\textbf{Status:} OPEN \quad|\quad \textbf{Prize:} none \quad|\quad \textbf{Updated:} 2025-08-31

\noindent\textbf{Tags:} number theory

\medskip

\noindent\textbf{Statement:}

Let $r\geq 3$. A number $n$ is $r$-powerful if for every prime $p$ which divides $n$ we have $p^r\mid n$. Are there infinitely many integers which are not the sum of at most $r$ many $r$-powerful numbers? Does the set of integers which are the sum of at most $r$ $r$-powerful numbers have density $0$?




Erd\H{o}s \cite{Er76d} claims that the claim that the set has density $0$ is 'easy' for $r=2$ (a potential 'easy argument' is given in the comments by Tao; this was first proved in the literature by Baker and Br\"{u}dern \cite{BaBr94}). For $r=3$ it is not even known if those integers which are the sum of at most three cubes has density $0$. In the Oberwolfach problem book this is recorded in 1986 as a problem of Erd\H{o}s and Ivi\'{c}. In \cite{Er76d} Erd\H{o}s claims that 'a simple counting argument' implies that there are infinitely many integers which are not the sum of at most $r$ many $r$-powerful numbers, but Schinzel pointed out he made a mistake. Heath-Brown \cite{He88} has proved that all large numbers are the sum of at most three $2$-powerful numbers, see [941] . See also [1081] for a more refined question concerning the case $r=2$, and [1107] for the case of $r+1$ summands.




References

[BaBr94] Baker, R. C. and Br\"udern, J., On sums of two squarefull numbers . Math. Proc. Cambridge Philos. Soc. (1994), 1--5.

[Er76d] Erd\H{o}s, P., Problems and results on number theoretic properties of consecutive integers and related questions . Proceedings of the Fifth Manitoba Conference on Numerical Mathematics (Univ. Manitoba, Winnipeg, Man., 1975) (1976), 25-44.

[He88] Heath-Brown, D. R., Ternary quadratic forms and sums of three square-full numbers . (1988), 137--163.

\noindent\rule{\linewidth}{0.4pt}


%% ============================================================
\section{Problem \#942}
\label{prob:942}

\noindent\textbf{Status:} OPEN \quad|\quad \textbf{Prize:} none \quad|\quad \textbf{Updated:} 2025-08-31

\noindent\textbf{Tags:} number theory

\medskip

\noindent\textbf{Statement:}

Let $h(n)$ count the number of powerful (if $p\mid m$ then $p^2\mid m$) integers in $[n^2,(n+1)^2)$. Estimate $h(n)$. In particular is there some constant $c&#62;0$ such that\[h(n) &lt; (\log n)^{c+o(1)}\]and, for infinitely many $n$,\[h(n) &gt;(\log n)^{c-o(1)}?\]




Erd\H{o}s writes it is not hard to prove that $\limsup h(n)=\infty$, and that the density $\delta_l$ of integers for which $h(n)=l$ exists and $\sum \delta_l=1$. A proof that $h(n)$ is unbounded is provided by van Doorn in the comments. De Koninck and Luca \cite{DeLu04} have proved, for infinitely many $n$,\[h(n) \gg \left(\frac{\log n}{\log\log n}\right)^{1/3}.\]They also give the density ($\approx 0.275$) of those $n$ such that $h(n)=1$.




References

[DeLu04] De Koninck, Jean-Marie and Luca, Florian, Sur la proximit\'{e} des nombres puissants . Acta Arith. (2004), 149--157.

\noindent\rule{\linewidth}{0.4pt}


%% ============================================================
\section{Problem \#943}
\label{prob:943}

\noindent\textbf{Status:} OPEN \quad|\quad \textbf{Prize:} none \quad|\quad \textbf{Updated:} 2025-08-31

\noindent\textbf{Tags:} number theory

\medskip

\noindent\textbf{Statement:}

Let $A$ be the set of powerful numbers (if $p\mid n$ then $p^2\mid n$). Is it true that\[1_A\ast 1_A(n)=n^{o(1)}\]for every $n$?

\noindent\rule{\linewidth}{0.4pt}


%% ============================================================
\section{Problem \#944}
\label{prob:944}

\noindent\textbf{Status:} OPEN \quad|\quad \textbf{Prize:} none \quad|\quad \textbf{Updated:} 2025-08-31

\noindent\textbf{Tags:} graph theory, chromatic number

\medskip

\noindent\textbf{Statement:}

A critical vertex, edge, or set of edges, is one whose deletion lowers the chromatic number. Let $k\geq 4$ and $r\geq 1$. Must there exist a graph $G$ with chromatic number $k$ such that every vertex is critical, yet every critical set of edges has size $&#62;r$?




A graph $G$ with chromatic number $k$ in which every vertex is critical is called $k$-vertex-critical. This was conjectured by Dirac in 1970 for $k\geq 4$ and $r=1$. Dirac's conjecture was proved, for $k=5$, by Brown \cite{Br92}. Lattanzio \cite{La02} proved there exist such graphs for all $k$ such that $k-1$ is not prime. Independently, Jensen \cite{Je02} gave an alternative construction for all $k\geq 5$. The case $k=4$ and $r=1$ remains open. Martinsson and Steiner \cite{MaSt25} proved this is true for every $r\geq 1$ if $k$ is sufficiently large, depending on $r$. Skottova and Steiner \cite{SkSt25} have improved this, proving that such graphs exist for all $k\geq 5$ and $r\geq 1$. The only remaining open case is $k=4$ (even the case $k=4$ and $r=1$ remains open). Erd\H{o}s also asked a stronger quantitative form of this question: let $f_k(n)$ denote the largest $r\geq 1$ such that there exists a $k$-vertex-critical graph on $n$ vertices such that no set of at most $r$ edges is critical. He then asks whether $f_k(n)\to \infty$ as $n\to \infty$. Skottova and Steiner \cite{SkSt25} have proved this for $k\geq 5$, establishing the bounds\[n^{1/3}\ll_k f_k(n) \ll_k \frac{n}{(\log n)^C}\]for all $k\geq 5$, where $C&gt;0$ is an absolute constant. This is Problem 91 in the graph problems collection. See also [917] and [1032] .




References

[Br92] Brown, Jason I., A vertex critical graph without critical edges . Discrete Math. (1992), 99--101.

[Je02] Jensen, Tommy R., Dense critical and vertex-critical graphs . Discrete Math. (2002), 63--84.

[La02] Lattanzio, John J., A note on a conjecture of {D}irac . Discrete Math. (2002), 323--330.

[MaSt25] Martinsson, Anders and Steiner, Raphael, Vertex-critical graphs far from edge-criticality . Combin. Probab. Comput. (2025), 151--157.

[SkSt25] E. Skottova and R. Steiner, Critical edge sets in vertex-critical graphs . arXiv:2508.08703 (2025).

\noindent\rule{\linewidth}{0.4pt}


%% ============================================================
\section{Problem \#945}
\label{prob:945}

\noindent\textbf{Status:} OPEN \quad|\quad \textbf{Prize:} none \quad|\quad \textbf{Updated:} 2025-08-31

\noindent\textbf{Tags:} number theory, divisors

\medskip

\noindent\textbf{Statement:}

Let $F(x)$ be the maximal $k$ such that there exist $n+1,\ldots,n+k\leq x$ with $\tau(n+1),\ldots,\tau(n+k)$ all distinct (where $\tau(m)$ counts the divisors of $m$). Estimate $F(x)$. In particular, is it true that\[F(x) \leq (\log x)^{O(1)}?\]In other words, is there a constant $C&#62;0$ such that, for all large $x$, every interval $[x,x+(\log x)^C]$ contains two integers with the same number of divisors?




A problem of Erd\H{o}s and Mirsky \cite{ErMi52}, who proved that\[\frac{(\log x)^{1/2}}{\log\log x}\ll F(x) \ll \exp\left(O\left(\frac{(\log x)^{1/2}}{\log\log x}\right)\right).\]Erd\H{o}s \cite{Er85e} claimed that the lower bound could be improved via their method 'with some more work' to $(\log x)^{1-o(1)}$. Beker has improved the upper bound to\[F(x) \ll \exp\left(O\left((\log x)^{1/3+o(1)}\right)\right).\]Cambie has observed that Cram\'{e r's conjecture} implies that $F(x) \ll (\log x)^2$, and furthermore if every interval in $[x,2x]$ of length $\gg \log x$ contains a squarefree number (see [208] ) then every interval of length $\gg (\log x)^2$ contains two numbers with the same number of divisors, whence\[F(x) \ll (\log x)^2.\]See [1004] for the analogous problem with the Euler totient function. This problem is discussed in problem B18 of Guy's collection \cite{Gu04}.




References

[Er85e] Erd\H{o}s, P., Some problems and results in number theory . Number theory and combinatorics. Japan 1984 (Tokyo, Okayama and Kyoto, 1984) (1985), 65-87.

[ErMi52] Erd\H{o}s, P. and Mirsky, L., The distribution of values of the divisor function {$d(n)$} . Proc. London Math. Soc. (3) (1952), 257--271.

[Gu04] Guy, Richard K., Unsolved problems in number theory . (2004), xviii+437.

\noindent\rule{\linewidth}{0.4pt}


%% ============================================================
\section{Problem \#948}
\label{prob:948}

\noindent\textbf{Status:} OPEN \quad|\quad \textbf{Prize:} none \quad|\quad \textbf{Updated:} 2025-08-31

\noindent\textbf{Tags:} number theory, ramsey theory

\medskip

\noindent\textbf{Statement:}

Is there a function $f(n)$ and a $k$ such that in any $k$-colouring of the integers there exists a sequence $a_1&#60;\cdots$ such that $a_n&#60;f(n)$ for infinitely many $n$ and the set\[\left\{ \sum_{i\in S}a_i : \textrm{finite }S\right\}\]does not contain all colours?




Erd\H{o}s initially asked whether this is possible with the set being monochromatic, but Galvin showed that this is not always possible, considering the two colouring where, writing $n=2^km$ with $m$ odd, we colour $n$ red if $m\geq F(k)$ and blue if $m&lt;F(k)$ (for some sufficiently quickly growing $F$). In other words, the answer to this question is no when $k=2$. The original question is open even in the case of $\aleph_0$-many colours. This is asked by Erd\H{o}s and Galvin in \cite{ErGa91}, where they note that they do not even know whether this is possible with $k=3$. In \cite{ErGa91} they do prove that, for all $k\geq 2$, in any $k$-colouring of the integers there is a sequence such that $a_n&lt;2^{2^{O(n)}}$ for all $n$ and\[\left\{ \sum_{i\in I}a_i : \textrm{finite intervals } I\right\}\]is coloured with only two colours. This is asking about a variant of Hindman's theorem (see [532] ).




References

[ErGa91] Erd\H{o}s, Paul and Galvin, Fred, Some {R}amsey-type theorems . Discrete Math. (1991), 261--269.

\noindent\rule{\linewidth}{0.4pt}


%% ============================================================
\section{Problem \#949}
\label{prob:949}

\noindent\textbf{Status:} OPEN \quad|\quad \textbf{Prize:} none \quad|\quad \textbf{Updated:} 2025-08-31

\noindent\textbf{Tags:} ramsey theory

\medskip

\noindent\textbf{Statement:}

Let $S\subset \mathbb{R}$ be a set containing no solutions to $a+b=c$. Must there be a set $A\subseteq \mathbb{R}\backslash S$ of cardinality continuum such that $A+A\subseteq \mathbb{R}\backslash S$?




Erd\H{o}s suggests that if the answer is no one could consider the variant where we assume that $S$ is Sidon (i.e. all $a+b$ with $a,b\in S$ are distinct, aside from the trivial coincidences). In the comments Dillies gives a positive proof of this, found by AlphaProof: in other words, if $S\subset \mathbb{R}$ is a Sidon set then there exists $A\subseteq \mathbb{R}\backslash S$ of cardinality continuum such that $A+A\subseteq \mathbb{R}\backslash S$.

\noindent\rule{\linewidth}{0.4pt}


%% ============================================================
\section{Problem \#950}
\label{prob:950}

\noindent\textbf{Status:} OPEN \quad|\quad \textbf{Prize:} none \quad|\quad \textbf{Updated:} 2025-08-31

\noindent\textbf{Tags:} number theory, primes

\medskip

\noindent\textbf{Statement:}

Let\[f(n) = \sum_{p&#60;n}\frac{1}{n-p}.\]Is it true that\[\liminf f(n)=1\]and\[\limsup f(n)=\infty?\]Is it true that $f(n)=o(\log\log n)$ for all $n$?




This function was considered by de Bruijn, Erd\H{o}s, and Tur\'{a}n, who showed that\[\sum_{n&lt;x}f(n)\sim \sum_{n&lt;x}f(n)^2\sim x.\]They give no proofs, but a proof of the (harder) second claim is given by Gorodetsky here . The existence of some $c&gt;0$ such that there are $\gg n^c/\log n$ many primes in $[n,n+n^c]$ implies that $\liminf f(n)&gt;0$. Erd\H{o}s writes that a 'weaker conjecture which is perhaps not quite inaccessible' is that, for every $\epsilon&gt;0$, if $x$ is sufficiently large there exists $y&lt;x$ such that\[\pi(x)&lt; \pi(y)+\epsilon \pi(x-y).\](Compare this to [855] .) He notes that if\[\pi(x)&lt; \pi(y)+O\left(\frac{x-y}{\log x}\right)\]for all $y&lt;x-(\log x)^C$ for some constant $C&gt;0$ then $f(n)\ll \log\log\log n$. The study of $f(p)$ is even harder, and Erd\H{o}s could not prove that\[\sum_{p&lt;x}f(p)^2\sim \pi(x).\]See also [1210] .

\noindent\rule{\linewidth}{0.4pt}


%% ============================================================
\section{Problem \#951}
\label{prob:951}

\noindent\textbf{Status:} OPEN \quad|\quad \textbf{Prize:} none \quad|\quad \textbf{Updated:} 2025-08-31

\noindent\textbf{Tags:} number theory

\medskip

\noindent\textbf{Statement:}

Let $1&#60;a_1&#60;\cdots$ be a sequence of real numbers such that\[\left\lvert \prod_i a_i^{k_i}-\prod_j a_j^{\ell_j}\right\rvert \geq 1\]for every distinct pair of non-negative finitely supported integer tuples $k_i,\ell_j\geq 0$. Is it true that\[\#\{ a_i \leq x\} \leq \pi(x)?\]




Erd\H{o}s \cite{Er77c} said this question was asked 'during [his] lecture at Queens College [by] one member of the audience (perhaps S. Shapiro)'. (Although in \cite{Er80} he seems sure it was Shapiro, but also recalled that he had himself asked this in \cite{Er69}.) In \cite{Er80} he further asks whether equality holds if and only if the $a_i$ are the set of primes. Such a sequence of $a_i$ is sometimes called a set of Beurling prime numbers (and the sequence of products called the associated generalised integers). Beurling conjectured that if the number of reals in $[1,x]$ of the form $\prod a_i^{k_i}$ is $x+o(\log x)$ then the $a_i$ must be the sequence of primes. It is unclear whether this is intended to hold for all $x$ or just all sufficiently large $x$. The former question can be disproved by finite calculation (since any finite sequence of $a_i$ can be extended to a suitable infinite sequence via a greedy algorithm). A finite counterexample for $x=10$ was found by ChatGPT-5.2 Pro (prompted by Leeham).




References

[Er69] Erd\H{o}s, Paul, Some applications of graph theory to number theory . The Many Facets of Graph Theory (Proc. Conf., Western Mich. Univ., Kalamazoo, Mich., 1968) (1969), 77-82.

[Er77c] Erd\H{o}s, Paul, Problems and results on combinatorial number theory. III . Number theory day (Proc. Conf., Rockefeller Univ., New York, 1976) (1977), 43-72.

[Er80] Erd\H{o}s, Paul, A survey of problems in combinatorial number theory . Ann. Discrete Math. (1980), 89-115.

\noindent\rule{\linewidth}{0.4pt}


%% ============================================================
\section{Problem \#952}
\label{prob:952}

\noindent\textbf{Status:} OPEN \quad|\quad \textbf{Prize:} none \quad|\quad \textbf{Updated:} 2025-08-31

\noindent\textbf{Tags:} number theory

\medskip
\noindent\textit{Gaussian moat problem}


\noindent\textbf{Statement:}

Is there an infinite sequence of distinct Gaussian primes $x_1,x_2,\ldots$ such that\[\lvert x_{n+1}-x_n\rvert \ll 1?\]




The Gaussian moat problem . This is not actually a problem of Erd\H{o}s, but has been erroneously attributed to him in the past. In \cite{Er77c} Erd\H{o}s recalls: 'The conjecture was told me by Motzkin at the Pasadena number theory meeting 1963 November and it was apparently raised by Basil Gordon and Motzkin. I naturally liked it very much and told it right away to many people, naturally attributing it to Motzkin, but this was later forgotten. Thus the problem is returned to its rightful owners.' In \cite{Er80} Erd\H{o}s writes 'the answer is almost certainly negative'.




References

[Er77c] Erd\H{o}s, Paul, Problems and results on combinatorial number theory. III . Number theory day (Proc. Conf., Rockefeller Univ., New York, 1976) (1977), 43-72.

[Er80] Erd\H{o}s, Paul, A survey of problems in combinatorial number theory . Ann. Discrete Math. (1980), 89-115.

\noindent\rule{\linewidth}{0.4pt}


%% ============================================================
\section{Problem \#953}
\label{prob:953}

\noindent\textbf{Status:} OPEN \quad|\quad \textbf{Prize:} none \quad|\quad \textbf{Updated:} 2025-08-31

\noindent\textbf{Tags:} geometry, distances

\medskip

\noindent\textbf{Statement:}

Let $A\subset \{ x\in \mathbb{R}^2 : \lvert x\rvert &#60;r\}$ be a measurable set with no integer distances, that is, such that $\lvert a-b\rvert \not\in \mathbb{Z}$ for any distinct $a,b\in A$. How large can the measure of $A$ be?




A problem of Erd\H{o}s and S\'{a}rk\"{o}zi. Erd\H{o}s \cite{Er77c} writes that 'S\'{a}rk\"{o}zi has the sharpest results, but nothing has been published yet'. The trivial upper bound is $O(r)$. Koizumi and Kovac have observed in the comments that S\'{a}rk\"{o}zy's lower bound in [466] can be adapted to give a lower bound of $\gg_\epsilon r^{1/2-\epsilon}$ for all $\epsilon&gt;0$. See also [465] for a similar problem (concerning upper bounds) and [466] for a similar problem (concerning lower bounds).




References

[Er77c] Erd\H{o}s, Paul, Problems and results on combinatorial number theory. III . Number theory day (Proc. Conf., Rockefeller Univ., New York, 1976) (1977), 43-72.

\noindent\rule{\linewidth}{0.4pt}


%% ============================================================
\section{Problem \#954}
\label{prob:954}

\noindent\textbf{Status:} OPEN \quad|\quad \textbf{Prize:} none \quad|\quad \textbf{Updated:} 2025-08-31

\noindent\textbf{Tags:} number theory

\medskip

\noindent\textbf{Statement:}

Let $0=a_0&#60;a_1&#60;a_2&#60;\cdots$ be the sequence of integers defined by $a_0=0$ and $a_1=1$, and $a_{k+1}$ is the smallest integer $n$ for which the number of solutions to $a_i+a_j \leq n$ (with $0\leq i\leq j\leq k$ and $j\geq 1$) is $&#60;n$. Is the number of solutions to $a_i+a_j \leq x$ equal to $x+O(x^{1/4+o(1)})$?




This sequence was constructed by Rosen. Note that the number of solutions to $a_i+a_j\leq x$ is always at least $x$ by construction, but Erd\H{o}s and Rosen could not even prove whether it is at most $(1+o(1))x$. The sequence begins\[0,1,3,5,9,13,17,24,31,38,45,\ldots.\]

\noindent\rule{\linewidth}{0.4pt}


%% ============================================================
\section{Problem \#955}
\label{prob:955}

\noindent\textbf{Status:} OPEN \quad|\quad \textbf{Prize:} none \quad|\quad \textbf{Updated:} 2025-08-31

\noindent\textbf{Tags:} number theory

\medskip

\noindent\textbf{Statement:}

Let\[s(n)=\sigma(n)-n=\sum_{\substack{d\mid n\\ d&#60;n}}d\]be the sum of proper divisors function. If $A\subset \mathbb{N}$ has density $0$ then $s^{-1}(A)$ must also have density $0$.




A conjecture of Erd\H{o}s, Granville, Pomerance, and Spiro \cite{EGPS90}. It is possible for $s(A)$ to have positive density even if $A$ has zero density (for example taking $A$ to be the product of two distinct primes). Erd\H{o}s \cite{Er73b} proved that there are sets $A$ of positive density such that $s^{-1}(A)$ is empty. Pollack \cite{Po14b} has shown that this is true if $A$ is the set of primes. Troupe \cite{Tr15} has shown that this is true if $A$ is the set of integers with unusually many prime factors. Troupe \cite{Tr20} has also shown this is true if $A$ is the set of integers which are the sum of two squares. Pollack, Pomerance, and Thompson \cite{PPT18} prove that if $\epsilon(x)=o(1)$ and $A\subset \mathbb{N}$ has size at most $x^{1/2+\epsilon(x)}$ then\[\#\{ n\leq x: s(n)\in A\} =o(x)\]as $x\to \infty$. It follows that (using $s(n)\ll n\log\log n$) if $A$ grows like $\lvert A\cap [1,x]\rvert\leq x^{1/2+o(1)}$ then $s^{-1}(A)$ has density $0$. Alanen calls $k$ such that $s(n)=k$ has no solutions untouchable. These are discussed in problem B10 of Guy's collection \cite{Gu04}.




References

[EGPS90] Erd\H{o}s, P. and Granville, A. and Pomerance, C. and Spiro, C., On the normal behavior of the iterates of some arithmetic functions . Analytic number theory (Allerton Park, IL, 1989) (1990), 165-204.

[Er73b] Erd\H{o}s, P., \"{U}ber die Zahlen der Form $\sigma (n)-n$ und $n-\phi(n)$ . Elem. Math. (1973), 83-86.

[Gu04] Guy, Richard K., Unsolved problems in number theory . (2004), xviii+437.

[PPT18] Pollack, Paul and Pomerance, Carl and Thompson, Lola, Divisor-sum fibers . Mathematika (2018), 330--342.

[Po14b] Pollack, Paul, Some arithmetic properties of the sum of proper divisors and the sum of prime divisors . Illinois J. Math. (2014), 125--147.

[Tr15] Troupe, Lee, On the number of prime factors of values of the sum-of-proper-divisors function . J. Number Theory (2015), 120--135.

[Tr20] Troupe, Lee, Divisor sums representable as the sum of two squares . Proc. Amer. Math. Soc. (2020), 4189--4202.

\noindent\rule{\linewidth}{0.4pt}


%% ============================================================
\section{Problem \#956}
\label{prob:956}

\noindent\textbf{Status:} OPEN \quad|\quad \textbf{Prize:} none \quad|\quad \textbf{Updated:} 2025-08-31

\noindent\textbf{Tags:} geometry, distances, convex

\medskip

\noindent\textbf{Statement:}

If $C,D\subseteq \mathbb{R}^2$ then the distance between $C$ and $D$ is defined by\[\delta(C,D)=\inf_{\substack{c\in C\\ d\in D}}\| c-d\|.\]Let $h(n)$ be the maximal number of unit distances between disjoint convex translates. That is, the maximal $m$ such that there is a compact convex set $C\subset \mathbb{R}^2$ and a set $X$ of size $n$ such that all $(C+x)_{x\in X}$ are disjoint and there are $m$ pairs $x_1,x_2\in X$ such that\[\delta(C+x_1,C+x_2)=1.\]Determine $h(n)$ - in particular, prove that there exists a constant $c&#62;0$ such that $h(n)&#62;n^{1+c}$ for all large $n$.




A problem of Erd\H{o}s and Pach \cite{ErPa90}, who proved that $h(n) \ll n^{4/3}$. They also consider the related function where we consider $n$ disjoint convex sets (not necessarily translates), for which they give an upper bound of $\ll n^{7/5}$. It is trivial that $h(n)\geq f(n)$, where $f(n)$ is the maximal number of unit distances determined by $n$ points in $\mathbb{R}^2$ (see [90] ).




References

[ErPa90] Erd\H{o}s, P. and Pach, J., Variations on the theme of repeated distances . Combinatorica (1990), 261--269.

\noindent\rule{\linewidth}{0.4pt}


%% ============================================================
\section{Problem \#959}
\label{prob:959}

\noindent\textbf{Status:} OPEN \quad|\quad \textbf{Prize:} none \quad|\quad \textbf{Updated:} 2025-08-31

\noindent\textbf{Tags:} geometry, distances

\medskip

\noindent\textbf{Statement:}

Let $A\subset \mathbb{R}^2$ be a set of size $n$ and let $\{d_1,\ldots,d_k\}$ be the set of distinct distances determined by $A$. Let $f(d)$ be the number of times the distance $d$ is determined, and suppose the $d_i$ are ordered such that\[f(d_1)\geq f(d_2)\geq \cdots \geq f(d_k).\]Estimate\[\max (f(d_1)-f(d_2)),\]where the maximum is taken over all $A$ of size $n$.




More generally, one can ask about\[\max (f(d_r)-f(d_{r+1})).\]Clemen, Dumitrescu, and Liu \cite{CDL25}, have shown that\[\max (f(d_1)-f(d_2))\gg n\log n.\]More generally, for any $1\leq k\leq \log n$, there exists a set $A$ of $n$ points such that\[f(d_r)-f(d_{r+1})\gg \frac{n\log n}{r}.\]They conjecture that $n\log n$ can be improved to $n^{1+c/\log\log n}$ for some constant $c&#62;0$.




References

[CDL25] F. Clemen, A. Dumitrescu, and D. Liu, On multiplicities of interpoint distances . arXiv:2505.04283 (2025).

\noindent\rule{\linewidth}{0.4pt}


%% ============================================================
\section{Problem \#961}
\label{prob:961}

\noindent\textbf{Status:} OPEN \quad|\quad \textbf{Prize:} none \quad|\quad \textbf{Updated:} 2025-08-31

\noindent\textbf{Tags:} number theory

\medskip

\noindent\textbf{Statement:}

Let $f(k)$ be the minimal $n$ such that every set of $n$ consecutive integers $&#62;k$ contains an integer divisible by a prime $&#62;k$. Estimate $f(k)$.




In other words, how large can a consecutive set of $k$-smooth integers be? Sylvester and Schur (see \cite{Er34}) proved $f(k)\leq k$ and Erd\H{o}s \cite{Er55d} proved\[f(k)&lt;3\frac{k}{\log k}.\]Jutila \cite{Ju74} and Ramachandra, and Shorey \cite{RaSh73} proved\[f(k) \ll \frac{\log\log\log k}{\log \log k}\frac{k}{\log k}.\]It is likely that $f(k) \ll (\log k)^{O(1)}$. This is essentially equivalent to [683] .




References

[Er34] Erd\H{o}s, Paul, A Theorem of Sylvester and Schur . J. London Math. Soc. (1934), 282--288.

[Er55d] Erd\H{o}s, P., On consecutive integers . Nieuw Arch. Wisk. (3) (1955), 124--128.

[Ju74] Jutila, Matti, On numbers with a large prime factor. {II} . J. Indian Math. Soc. (N.S.) (1974), 125--130.

[RaSh73] Ramachandra, K. and Shorey, T. N., On gaps between numbers with a large prime factor . Acta Arith. (1973), 99--111.

\noindent\rule{\linewidth}{0.4pt}


%% ============================================================
\section{Problem \#962}
\label{prob:962}

\noindent\textbf{Status:} OPEN \quad|\quad \textbf{Prize:} none \quad|\quad \textbf{Updated:} 2025-08-31

\noindent\textbf{Tags:} number theory

\medskip

\noindent\textbf{Statement:}

Let $k(n)$ be the maximal $k$ such that there exists $m\leq n$ such that each of the integers\[m+1,\ldots,m+k\]are divisible by at least one prime $&#62;k$. Estimate $k(n)$ - in particular, is it true that\[\log k(n) \leq (\log n)^{1/2+o(1)}?\]




Erd\H{o}s \cite{Er65} wrote it is 'not hard to prove' that\[\log k(n)\geq (\log n)^{1/2-o(1)}\]and it 'seems likely' that $k(n)=o(n^\epsilon)$, but had no non-trivial upper bound for $k(n)$. In \cite{Er76e} he gives an argument which proves\[\log k(n) \gg \sqrt{\log n\log\log n},\]and thought that this bound was 'fairly sharp'. Tao in the comments has given a simple argument proving $k(n) \leq (1+o(1))n^{1/2}$. In \cite{Er76e} Erd\H{o}s reports he could prove that\[k(n) \leq \exp(-(\log n)^c)n^{1/2}\]for some constant $c&gt;0$, but could not prove that $k(n) \leq n^{1/2-c}$ for a constant $c&gt;0$, which he said 'seems a ridiculously weak result'. Tang has proved a lower bound of\[\log k(n)\geq \left(\frac{1}{\sqrt{2}}-o(1)\right)\sqrt{\log n\log\log n}.\]




References

[Er65] Erd\H{o}s, P., Extremal problems in number theory . Proc. Sympos. Pure Math., Vol. VIII (1965), 181-189.

[Er76e] Erd\H{o}s, P., Problems and results on consecutive integers . Publ. Math. Debrecen (1976), 271-282.

\noindent\rule{\linewidth}{0.4pt}


%% ============================================================
\section{Problem \#963}
\label{prob:963}

\noindent\textbf{Status:} OPEN \quad|\quad \textbf{Prize:} none \quad|\quad \textbf{Updated:} 2025-08-31

\noindent\textbf{Tags:} number theory

\medskip

\noindent\textbf{Statement:}

Let $f(n)$ be the maximal $k$ such that in any set $A\subset \mathbb{R}$ of size $n$ there is a subset $B\subseteq A$ of size $\lvert B\rvert\geq k$ which is dissociated that is, the sums $\sum_{b\in S}b$ are distinct for all $S\subseteq B$. Estimate $f(n)$ - in particular, is it true that\[f(n)\geq \lfloor \log_2 n\rfloor?\]




Erd\H{o}s noted that the greedy algorithm showed $f(n)\geq \lfloor \log_3 n\rfloor$.

\noindent\rule{\linewidth}{0.4pt}


%% ============================================================
\section{Problem \#968}
\label{prob:968}

\noindent\textbf{Status:} OPEN \quad|\quad \textbf{Prize:} none \quad|\quad \textbf{Updated:} 2025-08-31

\noindent\textbf{Tags:} number theory

\medskip

\noindent\textbf{Statement:}

Let $u_n=p_n/n$, where $p_n$ is the $n$th prime. Does the set of $n$ such that $u_n&#60;u_{n+1}$ have positive density?




Erd\H{o}s and Prachar \cite{ErPr61} proved that\[\sum_{p_n&lt;x} \lvert u_{n+1}-u_n\rvert \asymp (\log x)^2,\]and that the set of $n$ such that $u_n&gt;u_{n+1}$ has positive density. As with many of Erd\H{o}s' questions, by 'positive density' here he most likely meant positive lower density - in other words, there exists a constant $c&#62;0$ such that for all large $x$ the number of such $n\leq x$ is at least $cx$. Erd\H{o}s also asks whether\[u_n&lt;u_{n+1}&lt;u_{n+2}\]or\[u_n&gt;u_{n+1}&#62;u_{n+2}\]have infinitely many solutions.




References

[ErPr61] Erd\H{o}s, P. and Prachar, K., S\"{a}tze und {P}robleme \"{u}ber {$p\sb{k}/k$} . Abh. Math. Sem. Univ. Hamburg (1961/62), 251--256.

\noindent\rule{\linewidth}{0.4pt}


%% ============================================================
\section{Problem \#969}
\label{prob:969}

\noindent\textbf{Status:} OPEN \quad|\quad \textbf{Prize:} none \quad|\quad \textbf{Updated:} 2025-08-31

\noindent\textbf{Tags:} number theory

\medskip

\noindent\textbf{Statement:}

Let $Q(x)$ count the number of squarefree integers in $[1,x]$. Determine the order of magnitude in the error term in the asymptotic\[Q(x)=\frac{6}{\pi^2}x+E(x).\]




It is elementary to prove $E(x)\ll x^{1/2}$, and the prime number theorem implies $o(x^{1/2})$. The best known unconditional upper bound is of the shape $x^{1/2-o(1)}$, due to Walfisz \cite{Wa63}. Evelyn and Linfoot \cite{EvLi31} proved that\[E(x) \gg x^{1/4},\]and this is likely the true order of magnitude. The Riemann Hypothesis would follow from $E(x)\ll x^{1/4}$. The true order of magnitude is unknown even assuming the Riemann Hypothesis. Conditional on this assumption, the best known upper bound is\[E(x)\ll x^{\frac{11}{35}+o(1)},\]due to Liu \cite{Li16}.




References

[EvLi31] Evelyn, C. J. A. and Linfoot, E. H., On a problem in the additive theory of numbers . Ann. of Math. (2) (1931), 261--270.

[Li16] Liu, H.-Q., On the distribution of squarefree numbers . J. Number Theory (2016), 202--222.

[Wa63] Walfisz, Arnold, Weylsche {E}xponentialsummen in der neueren {Z}ahlentheorie . (1963), 231.

\noindent\rule{\linewidth}{0.4pt}


%% ============================================================
\section{Problem \#970}
\label{prob:970}

\noindent\textbf{Status:} OPEN \quad|\quad \textbf{Prize:} none \quad|\quad \textbf{Updated:} 2025-08-31

\noindent\textbf{Tags:} number theory

\medskip
\noindent\textit{Jacobsthal's function}


\noindent\textbf{Statement:}

Let $h(k)$ be Jacobsthal's function, defined to as the minimal $m$ such that, if $n$ has at most $k$ prime factors, then in any set of $m$ consecutive integers there exists an integer coprime to $n$. Determine the order of magnitude of $h(k)$. In particular, is it true that\[h(k) \ll k^2?\]




That $h(k)\ll k^2$ is a conjecture of Jacobsthal. Iwaniec \cite{Iw78} proved\[h(k) \ll (k\log k)^2.\]The best lower bound known is\[h(k) \gg \frac{(\log k)(\log\log\log k)}{(\log\log k)^2}k,\]due to Ford, Green, Konyagin, Maynard, and Tao \cite{FGKMT18}. This is a more general form of the function considered in [687] .




References

[FGKMT18] Ford, Kevin and Green, Ben and Konyagin, Sergei and Maynard, James and Tao, Terence, Long gaps between primes . J. Amer. Math. Soc. (2018), 65-105.

[Iw78] Iwaniec, Henryk, On the problem of {J}acobsthal . Demonstratio Math. (1978), 225--231.

\noindent\rule{\linewidth}{0.4pt}


%% ============================================================
\section{Problem \#971}
\label{prob:971}

\noindent\textbf{Status:} OPEN \quad|\quad \textbf{Prize:} none \quad|\quad \textbf{Updated:} 2025-08-31

\noindent\textbf{Tags:} number theory

\medskip

\noindent\textbf{Statement:}

Let $p(a,d)$ be the least prime congruent to $a\pmod{d}$. Does there exist a constant $c&#62;0$ such that, for all large $d$,\[p(a,d) &#62; (1+c)\phi(d)\log d\]for $\gg \phi(d)$ many values of $a$?




Erd\H{o}s \cite{Er49c} could prove this is true for an infinite sequence of $d$. He also proved that, for any $\epsilon&#62;0$,\[p(a,d)&#60; \epsilon \phi(d)\log d\]for $\gg_\epsilon \phi(d)$ many values of $a$.




References

[Er49c] Erd\H{o}s, P., On some applications of {B}run's method . Acta Univ. Szeged. Sect. Sci. Math. (1949), 57--63.

\noindent\rule{\linewidth}{0.4pt}


%% ============================================================
\section{Problem \#972}
\label{prob:972}

\noindent\textbf{Status:} OPEN \quad|\quad \textbf{Prize:} none \quad|\quad \textbf{Updated:} 2025-08-31

\noindent\textbf{Tags:} number theory

\medskip

\noindent\textbf{Statement:}

Let $\alpha&#62;1$ be irrational. Are there infinitely many primes $p$ such that $\lfloor p\alpha\rfloor$ is also prime?




Vinogradov \cite{Vi48} proved that the sequence $\{p\alpha\}$ is uniformly distributed for every irrational $\alpha$, and hence there are infinitely many primes $p$ of the shape $p=\lfloor n\alpha\rfloor$ for every irrational $\alpha&#62;1$. Indeed, this occurs if and only if\[\frac{p}{\alpha}\leq n&lt;\frac{p+1}{\alpha},\]which is true if and only if $\{p\alpha^{-1}\}&gt;1-\alpha^{-1}$, which happens infinitely often by the uniform distribution of $\{p\alpha^{-1}\}$.




References

[Vi48] Vinogradov, I. M., On an estimate of trigonometric sums with prime numbers . Izv. Akad. Nauk SSSR Ser. Mat. (1948), 225--248.

\noindent\rule{\linewidth}{0.4pt}


%% ============================================================
\section{Problem \#973}
\label{prob:973}

\noindent\textbf{Status:} OPEN \quad|\quad \textbf{Prize:} none \quad|\quad \textbf{Updated:} 2025-08-31

\noindent\textbf{Tags:} analysis

\medskip

\noindent\textbf{Statement:}

Does there exist a constant $C&#62;1$ such that, for every $n\geq 2$, there exists a sequence $z_i\in \mathbb{C}$ with $z_1=1$ and $\lvert z_i\rvert \geq 1$ for all $1\leq i\leq n$ with\[\max_{2\leq k\leq n+1}\left\lvert \sum_{1\leq i\leq n}z_i^k\right\rvert &#60; C^{-n}?\]




This is Problem 7.3 in \cite{Ha74}, where it is attributed to Erd\H{o}s. Erd\H{o}s proved (as described on p.35 of \cite{Tu84b}) that such a sequence does exist with $\lvert z_i\rvert\leq 1$. Indeed, Erd\H{o}s' construction gives a value of $C\approx 1.32$. In \cite{Er92f} (a different) Erd\H{o}s refines this analysis, proving that if\[M_2=\min_{z_j} \max_{2\leq k\leq n+1} \left\lvert \sum_{1\leq j\leq n}z_j^k\right\rvert,\]where the minimum is take over all $z_j\in \mathbb{C}$ with $\max \lvert z_j\rvert=1$, then\[(1.746)^{-n} &lt; M_2 &lt; (1.745)^{-n}.\]Tang notes in the comments that Theorem 6.1 of \cite{Tu84b} implies that, if $\lvert z_i\rvert \geq 1$ for all $i$, then\[\max_{2\leq k\leq n+1}\left\lvert \sum_{1\leq i\leq n}z_i^k\right\rvert \geq (2e)^{-(1+o(1))n}.\]See also [519] .




References

[Er92f] Erd\H{o}s, L., On some problems of {P}. Tur\'{a}n concerning power sums of complex numbers . Acta Math. Hungar. (1992), 11--24.

[Ha74] Hayman, W. K., Research problems in function theory: new problems . (1974), 155--180.

[Tu84b] Tur\'{a}n, Paul, On a new method of analysis and its applications . (1984), xvi+584.

\noindent\rule{\linewidth}{0.4pt}


%% ============================================================
\section{Problem \#975}
\label{prob:975}

\noindent\textbf{Status:} OPEN \quad|\quad \textbf{Prize:} none \quad|\quad \textbf{Updated:} 2025-08-31

\noindent\textbf{Tags:} number theory, divisors

\medskip

\noindent\textbf{Statement:}

Let $f\in \mathbb{Z}[x]$ be an irreducible non-constant polynomial such that $f(n)\geq 1$ for all large $n\in\mathbb{N}$. Does there exist a constant $c=c(f)&#62;0$ such that\[\sum_{n\leq X} \tau(f(n))\sim cX\log X,\]where $\tau$ is the divisor function?




Van der Corput \cite{Va39} proved that\[\sum_{n\leq X} \tau(f(n))\gg_f X\log X.\]Erd\H{o}s \cite{Er52b} proved using elementary methods that\[\sum_{n\leq X} \tau(f(n))\ll_f X\log X.\]Such an asymptotic formula is known whenever $f$ is an irreducible quadratic, as proved by Hooley \cite{Ho63}. The form of $c$ depends on $f$ in a complicated fashion (see the work of McKee \cite{Mc95}, \cite{Mc97}, and \cite{Mc99} for expressions for various types of quadratic $f$). For example,\[\sum_{n\leq x}\tau(n^2+1)=\frac{3}{\pi}x\log x+O(x).\]Tao has a blog post on this topic.




References

[Er52b] Erd\H{o}s, P., On the sum {$\sum^x_{k=1} d(f(k))$} . J. London Math. Soc. (1952), 7--15.

[Ho63] Hooley, Christopher, On the number of divisors of a quadratic polynomial . Acta Math. (1963), 97--114.

[Mc95] McKee, James, On the average number of divisors of quadratic polynomials . Math. Proc. Cambridge Philos. Soc. (1995), 389--392.

[Mc97] McKee, James, A note on the number of divisors of quadratic polynomials . (1997), 275--281.

[Mc99] McKee, James, The average number of divisors of an irreducible quadratic polynomial . Math. Proc. Cambridge Philos. Soc. (1999), 17--22.

[Va39] van der Corput, J. G., Une in\'{e}galit\'{e}{} relative au nombre des diviseurs . Nederl. Akad. Wetensch., Proc. (1939), 547--553.

\noindent\rule{\linewidth}{0.4pt}


%% ============================================================
\section{Problem \#976}
\label{prob:976}

\noindent\textbf{Status:} OPEN \quad|\quad \textbf{Prize:} none \quad|\quad \textbf{Updated:} 2025-08-31

\noindent\textbf{Tags:} number theory

\medskip

\noindent\textbf{Statement:}

Let $f\in \mathbb{Z}[x]$ be an irreducible polynomial of degree $d\geq 2$. Let $F_f(n)$ be maximal such that there exists $1\leq m\leq n$ with $f(m)$ is divisible by a prime $\geq F_f(n)$. Equivalently, $F_f(n)$ is the greatest prime divisor of\[\prod_{1\leq m\leq n}f(m).\]Estimate $F_f(n)$. In particular, is it true that $F_f(n)\gg n^{1+c}$ for some constant $c&#62;0$? Or even $\gg n^d$?




Nagell \cite{Na22} and Ricci \cite{Ri34} proved that\[F_f(n) \gg n\log n,\]which Erd\H{o}s \cite{Er52c} improved to\[F_f(n) \gg n(\log n)^{\log\log\log n}.\]In \cite{Er65b} he claimed a proof of\[F_f(n) \gg n\exp((\log n)^c)\]for some constant $c&#62;0$, but said he had never published the proof, which was 'fairly complicated'. This seems to have been flawed, since Erd\H{o}s and Schinzel \cite{ErSc90} later published a weaker bound. A proof of the stronger bound above was finally provided by Tenenbaum \cite{Te90}.




References

[Er52c] Erd\H{o}s, P., On the greatest prime factor of {$\prod^x_{k=1}f(k)$} . J. London Math. Soc. (1952), 379--384.

[Er65b] Erd\H{o}s, Paul, Some recent advances and current problems in number theory . Lectures on Modern Mathematics, Vol. III (1965), 196-244.

[ErSc90] Erd\H{o}s, P. and Schinzel, A., On the greatest prime factor of {$\prod^x_{k=1}f(k)$} . Acta Arith. (1990), 191--200.

[Na22] Nagell, . Abhandlungen aus dem Math. Seminar Hamburg (1922), 179-194.

[Ri34] G. Ricci, . Annali de Mat. (1934), 295-303.

[Te90] Tenenbaum, G\'{e}rald, Sur une question d'{E}rd\H{o}s et Schinzel . (1990), 405--443.

\noindent\rule{\linewidth}{0.4pt}


%% ============================================================
\section{Problem \#978}
\label{prob:978}

\noindent\textbf{Status:} OPEN \quad|\quad \textbf{Prize:} none \quad|\quad \textbf{Updated:} 2025-08-31

\noindent\textbf{Tags:} number theory

\medskip

\noindent\textbf{Statement:}

Let $f\in \mathbb{Z}[x]$ be an irreducible polynomial of degree $k&#62;2$ (and suppose that $k\neq 2^l$ for any $l\geq 1$) such that the leading coefficient of $f$ is positive. Does the set of integers $n\geq 1$ for which $f(n)$ is $(k-1)$-power-free have positive density? If $k&#62;3$, and for all primes $p$ there exists $n$ such that $p^{k-2}\nmid f(n)$, then are there infinitely many $n$ for which $f(n)$ is $(k-2)$-power-free? In particular, does\[n^4+2\]represent infinitely many squarefree numbers?




Erd\H{o}s \cite{Er53} proved there are infinitely many $n$ for which $f(n)$ is $(k-1)$-power-free, except when $k=2^l$ and $2^{k-1}\mid f(n)$ for all $n$. This latter exception is possible, for example if\[f(x)=k!\left(\binom{x}{k}+1\right).\]Hooley \cite{Ho67} settled the first question, in fact providing a precise asymptotic for the number of such $n\leq x$. Heath-Brown \cite{He06} proved the answer to the second question is yes when $k\geq 10$, and Browning \cite{Br11} extended this to $k\geq 9$ (in fact establishing an asymptotic formula for the number of such $n$). (This is under the assumption that there does not exist a prime $p$ such that $p^{k-2}\mid f(n)$ for all $n$, which Erd\H{o}s did not state explicitly, but is an obvious necessary condition, and is discussed explicitly in \cite{He06} and \cite{Br11}.) In \cite{Er65b} Erd\H{o}s mentions the question of whether $2^n\pm 1$ represents infinitely many $k$th power-free integers, or $n!\pm 1$, but that these are 'intractable at present'. (See also [936] .)




References

[Br11] Browning, T. D., Power-free values of polynomials . Arch. Math. (Basel) (2011), 139--150.

[Er53] Erd\H{o}s, P., Arithmetical properties of polynomials . J. London Math. Soc. (1953), 416--425.

[Er65b] Erd\H{o}s, Paul, Some recent advances and current problems in number theory . Lectures on Modern Mathematics, Vol. III (1965), 196-244.

[He06] Heath-Brown, D. R., Counting rational points on algebraic varieties . (2006), 51--95.

[Ho67] Hooley, C., On the power free values of polynomials . Mathematika (1967), 21--26.

\noindent\rule{\linewidth}{0.4pt}


%% ============================================================
\section{Problem \#979}
\label{prob:979}

\noindent\textbf{Status:} OPEN \quad|\quad \textbf{Prize:} none \quad|\quad \textbf{Updated:} 2025-08-31

\noindent\textbf{Tags:} number theory

\medskip

\noindent\textbf{Statement:}

Let $k\geq 2$, and let $f_k(n)$ count the number of solutions to\[n=p_1^k+\cdots+p_k^k,\]where the $p_i$ are prime numbers. Is it true that $\limsup f_k(n)=\infty$?




Erd\H{o}s \cite{Er37b} proved this is true when $k=2$, and also when $k=3$ (but this proof appears to be unpublished).




References

[Er37b] Erd\H{o}s, Paul, On the Sum and {D}ifference of Squares of {P}rimes . J. London Math. Soc. (1937), 133--136.

\noindent\rule{\linewidth}{0.4pt}


%% ============================================================
\section{Problem \#983}
\label{prob:983}

\noindent\textbf{Status:} OPEN \quad|\quad \textbf{Prize:} none \quad|\quad \textbf{Updated:} 2025-08-31

\noindent\textbf{Tags:} number theory

\medskip

\noindent\textbf{Statement:}

Let $n\geq 2$ and $\pi(n)&lt;k\leq n$. Let $f(k,n)$ be the smallest integer $r$ such that in any $A\subseteq \{1,\ldots,n\}$ of size $\lvert A\rvert=k$ there exist primes $p_1,\ldots,p_r$ such that $&gt;r$ many $a\in A$ are only divisible by primes from $\{p_1,\ldots,p_r\}$. Is it true that\[2\pi(n^{1/2})-f(\pi(n)+1,n)\to \infty\]as $n\to \infty$? In general, estimate $f(k,n)$, particularly when $\pi(n)+1&#60;k=o(n)$.




It is trivial that $f(k,n)\leq \pi(n)$. Erd\H{o}s and Straus \cite{Er70b} proved that\[f(\pi(n)+1,n)=2\pi(n^{1/2})+o_A\left(\frac{n^{1/2}}{(\log n)^A}\right)\]for any $A&#62;0$ and, for any constant $1&#62;c&#62;0$,\[f(cn,n)=\log\log n+(c_1+o(1))\sqrt{2\log\log n},\]where\[c=\frac{1}{\sqrt{2\pi}}\int_{-\infty}^{c_1}e^{-x^2/2}\mathrm{d}x.\]




References

[Er70b] Erd\H{o}s, P., Some applications of graph theory to number theory . Proc. Second Chapel Hill Conf. on Combinatorial Mathematics and its Applications (Univ. North Carolina, Chapel Hill, N.C., 1970) (1970), 136-145.

\noindent\rule{\linewidth}{0.4pt}


%% ============================================================
\section{Problem \#985}
\label{prob:985}

\noindent\textbf{Status:} OPEN \quad|\quad \textbf{Prize:} none \quad|\quad \textbf{Updated:} 2025-08-31

\noindent\textbf{Tags:} number theory

\medskip

\noindent\textbf{Statement:}

Is it true that, for every prime $p$, there is a prime $q&#60;p$ which is a primitive root modulo $p$?




Artin conjectured that $2$ is a primitive root for infinitely many primes $p$, which Hooley \cite{Ho67b} proved assuming the Generalised Riemann Hypothesis. Heath-Brown \cite{He86b} proved that at least one of $2$, $3$, or $5$ is a primitive root for infinitely many primes $p$.




References

[He86b] Heath-Brown, D. R., Artin's conjecture for primitive roots . Quart. J. Math. Oxford Ser. (2) (1986), 27--38.

[Ho67b] Hooley, Christopher, On {A}rtin's conjecture . J. Reine Angew. Math. (1967), 209--220.

\noindent\rule{\linewidth}{0.4pt}


%% ============================================================
\section{Problem \#986}
\label{prob:986}

\noindent\textbf{Status:} OPEN \quad|\quad \textbf{Prize:} none \quad|\quad \textbf{Updated:} 2025-08-31

\noindent\textbf{Tags:} graph theory, ramsey theory

\medskip

\noindent\textbf{Statement:}

For any fixed $s\geq 3$,\[R(s,k) \gg \frac{k^{s-1}}{(\log k)^c}\]for some constant $c=c(s)&#62;0$.




According to Chung and Graham \cite{ChGr98} this was first conjectured by Erd\H{o}s in 1947. Spencer \cite{Sp77} proved this for $s=3$ and Mattheus and Verstraete \cite{MaVe23} proved this for $s=4$. The best general bounds available for $s\geq 5$ are\[\frac{k^{s-2}}{(\log k)^{2s-6}}\ll_k R(s,k) \ll_s \frac{k^{s-1}}{(\log k)^{s-2}}.\]The lower bound was proved by Brada\v{c} \cite{Br26}, improving on earlier work of Bohman and Keevash \cite{BoKe10}. The upper bound was proved by Ajtai, Koml\'{o}s, and Szemer\'{e}di \cite{AKS80}. Li, Rousseau, and Zang \cite{LRZ01} have shown that $\ll_s$ in the upper bound can be improved to $\leq (1+o(1))$. The special case $s=3$ is the topic of [165] and $s=4$ is the topic of [166] . This problem is #6 in Ramsey Theory in the graphs problem collection. See also [920] .




References

[AKS80] Ajtai, Mikl\'{o}s and Koml\'{o}s, J\'{a}nos and Szemer\'{e}di, Endre, A note on Ramsey numbers . J. Combin. Theory Ser. A (1980), 354-360.

[BoKe10] Bohman, Tom and Keevash, Peter, The early evolution of the {$H$}-free process . Invent. Math. (2010), 291--336.

[Br26] D. Brada\v{c}, Nearly tight exponents for off-diagonal Ramsey numbers . arXiv:2605.28793 (2026).

[ChGr98] F. Chung and R. Graham, Erd\H{o}s on Graphs . His legacy of unsolved problems, A K Peters Ltd, Wellesley, MA (1998).

[LRZ01] Li, Yusheng and Rousseau, Cecil C. and Zang, Wenan, Asymptotic upper bounds for {R}amsey functions . Graphs Combin. (2001), 123--128.

[MaVe23] Mattheus, S. and Verstraete, J., The asymptotics of $r(4,t)$ . arXiv:2306.04007 (2023).

[Sp77] Spencer, J., Asymptotic lower bounds for Ramsey functions . Discrete Math. (1977), 69-76.

\noindent\rule{\linewidth}{0.4pt}


%% ============================================================
\section{Problem \#995}
\label{prob:995}

\noindent\textbf{Status:} OPEN \quad|\quad \textbf{Prize:} none \quad|\quad \textbf{Updated:} 2025-09-07

\noindent\textbf{Tags:} analysis, discrepancy

\medskip

\noindent\textbf{Statement:}

Let $n_1&#60;n_2&#60;\cdots$ be a lacunary sequence of integers and $f\in L^2([0,1])$. Estimate the growth of, for almost all $\alpha$,\[\sum_{1\leq k\leq N}f(\{ \alpha n_k\}).\]For example, is it true that, for almost all $\alpha$,\[\sum_{1\leq k\leq N}f(\{ \alpha n_k\})=o(N\sqrt{\log\log N})?\]




Erd\H{o}s \cite{Er49d} constructed a lacunary sequence and $f\in L^2([0,1])$ such that, for every $\epsilon&#62;0$, for almost all $\alpha$\[\limsup_{N\to \infty}\frac{1}{N(\log\log N)^{\frac{1}{2}-\epsilon}}\sum_{1\leq k\leq N}f(\{\alpha n_k\})=\infty.\]Erd\H{o}s also proved that, for every lacunary sequence and $f\in L^2$, for every $\epsilon&#62;0$, for almost all $\alpha$,\[\sum_{1\leq k\leq N}\sum_{1\leq k\leq N}f(\{\alpha n_k\})=o( N(\log N)^{\frac{1}{2}+\epsilon}).\]Erd\H{o}s \cite{Er64b} thought that his lower bound was closer to the truth.




References

[Er49d] Erd\H{o}s, P., On the strong law of large numbers . Trans. Amer. Math. Soc. (1949), 51--56.

[Er64b] Erd\H{o}s, P., Problems and results on diophantine approximations . Compositio Math. (1964), 52-65.

\noindent\rule{\linewidth}{0.4pt}


%% ============================================================
\section{Problem \#996}
\label{prob:996}

\noindent\textbf{Status:} OPEN \quad|\quad \textbf{Prize:} none \quad|\quad \textbf{Updated:} 2025-09-07

\noindent\textbf{Tags:} analysis

\medskip

\noindent\textbf{Statement:}

Let $n_1&lt;n_2&lt;\cdots$ be a lacunary sequence of integers, and let $f\in L^2([0,1])$. Let $f_n$ be the $n$th partial sum of the Fourier series of $f(x)$. Is there an absolute constant $C&gt;0$ such that, if\[\| f-f_n\|_2 \ll \frac{1}{(\log\log\log n)^{C}}\]then\[\lim_{N\to\infty}\frac{1}{N}\sum_{k\leq N}f(\{\alpha n_k\})=\int_0^1 f(x)\mathrm{d}x\]for almost every $\alpha$?




Raikov proved the conclusion always holds (for every $f\in L^2([0,1])$, with no assumption on $\| f-f_n\|_2$) if $n_k=a^k$ for some integer $a\geq 2$. Erd\H{o}s \cite{Er64b} also asked whether this is true for $n_k=\lfloor a^k\rfloor$ for some $a&#62;1$. Kac, Salem, and Zygmund \cite{KSZ48} proved that the conclusion holds if\[\| f-f_n\|_2 \ll \frac{1}{(\log n)^{c}}\]for some constant $c&#62;1$. Erd\H{o}s \cite{Er49d} proved that the conclusion holds if\[\| f-f_n\|_2 \ll \frac{1}{(\log\log n)^{c}}\]for some constant $c&#62;1$. Matsuyama \cite{Ma66} improved this to $c&#62;1/2$. In \cite{Er64b} Erd\H{o}s asked whether the conclusion holds for all bounded functions $f$ and lacunary sequences $n_k$.




References

[Er49d] Erd\H{o}s, P., On the strong law of large numbers . Trans. Amer. Math. Soc. (1949), 51--56.

[Er64b] Erd\H{o}s, P., Problems and results on diophantine approximations . Compositio Math. (1964), 52-65.

[KSZ48] Kac, M. and Salem, R. and Zygmund, A., A gap theorem . Trans. Amer. Math. Soc. (1948), 235--243.

[Ma66] Matsuyama, Noboru, On the strong law of large numbers . Tohoku Math. J. (2) (1966), 259--269.

\noindent\rule{\linewidth}{0.4pt}


%% ============================================================
\section{Problem \#1002}
\label{prob:1002}

\noindent\textbf{Status:} OPEN \quad|\quad \textbf{Prize:} none \quad|\quad \textbf{Updated:} 2025-09-07

\noindent\textbf{Tags:} analysis, diophantine approximation

\medskip

\noindent\textbf{Statement:}

For any $0&#60;\alpha&#60;1$, let\[f(\alpha,n)=\frac{1}{\log n}\sum_{1\leq k\leq n}(\tfrac{1}{2}-\{ \alpha k\}).\]Does $f(\alpha,n)$ have an asymptotic distribution function? In other words, is there a non-decreasing function $g$ such that $g(-\infty)=0$, $g(\infty)=1$, and\[\lim_{n\to \infty}\lvert \{ \alpha\in (0,1): f(\alpha,n)\leq c\}\rvert=g(c)?\]




Kesten \cite{Ke60} proved that if\[f(\alpha,\beta,n)=\frac{1}{\log n}\sum_{1\leq k\leq n}(\tfrac{1}{2}-\{\beta+\alpha k\})\]then $f(\alpha,\beta,n)$ has asymptotic distribution function\[g(c)=\frac{1}{\pi}\int_{-\infty}^{\rho c}\frac{1}{1+t^2}\mathrm{d}t,\]where $\rho&#62;0$ is an explicit constant.




References

[Ke60] Kesten, Harry, Uniform distribution {${\rm mod}\,1$} . Ann. of Math. (2) (1960), 445--471.

\noindent\rule{\linewidth}{0.4pt}


%% ============================================================
\section{Problem \#1003}
\label{prob:1003}

\noindent\textbf{Status:} OPEN \quad|\quad \textbf{Prize:} none \quad|\quad \textbf{Updated:} 2025-09-08

\noindent\textbf{Tags:} number theory

\medskip

\noindent\textbf{Statement:}

Are there infinitely many solutions to $\phi(n)=\phi(n+1)$, where $\phi$ is the Euler totient function?




Erd\H{o}s \cite{Er85e} says that, presumably, for every $k\geq 1$ the equation\[\phi(n)=\phi(n+1)=\cdots=\phi(n+k)\]has infinitely many solutions. Erd\H{o}s, Pomerance, and S\'{a}rk\"{o}zy \cite{EPS87} proved that the number of $n\leq x$ with $\phi(n)=\phi(n+1)$ is at most\[\frac{x}{\exp((\log x)^{1/3})}.\]See [946] for the analogous question with the divisor function, and [415] for a question about more general inequality patterns among consecutive values of $\phi$.




References

[EPS87] Erd\H{o}s, Paul and Pomerance, Carl and S\'{a}rk\"ozy, Andr\'{a}s, On locally repeated values of certain arithmetic functions. {III} . Proc. Amer. Math. Soc. (1987), 1--7.

[Er85e] Erd\H{o}s, P., Some problems and results in number theory . Number theory and combinatorics. Japan 1984 (Tokyo, Okayama and Kyoto, 1984) (1985), 65-87.

\noindent\rule{\linewidth}{0.4pt}


%% ============================================================
\section{Problem \#1004}
\label{prob:1004}

\noindent\textbf{Status:} OPEN \quad|\quad \textbf{Prize:} none \quad|\quad \textbf{Updated:} 2025-09-07

\noindent\textbf{Tags:} number theory

\medskip

\noindent\textbf{Statement:}

Let $c&#62;0$. If $x$ is sufficiently large then does there exist $n\leq x$ such that the values of $\phi(n+k)$ are all distinct for $1\leq k\leq (\log x)^c$, where $\phi$ is the Euler totient function?




Erd\H{o}s, Pomerance, and S\'{a}rk\"{o}zy \cite{EPS87} proved that if $\phi(n+k)$ are all distinct for $1\leq k\leq K$ then\[K \leq \frac{n}{\exp(c(\log n)^{1/3})}\]for some constant $c&gt;0$. See [945] for the analogous problem with the divisor function.




References

[EPS87] Erd\H{o}s, Paul and Pomerance, Carl and S\'{a}rk\"ozy, Andr\'{a}s, On locally repeated values of certain arithmetic functions. {III} . Proc. Amer. Math. Soc. (1987), 1--7.

\noindent\rule{\linewidth}{0.4pt}


%% ============================================================
\section{Problem \#1005}
\label{prob:1005}

\noindent\textbf{Status:} OPEN \quad|\quad \textbf{Prize:} none \quad|\quad \textbf{Updated:} 2025-09-09

\noindent\textbf{Tags:} number theory

\medskip

\noindent\textbf{Statement:}

Let $\frac{a_1}{b_1},\frac{a_2}{b_2},\ldots$ be the Farey fractions of order $n\geq 4$. Let $f(n)$ be the largest integer such that if $1\leq k&lt;l\leq k+f(n)$ then $\frac{a_k}{b_k}$ and $\frac{a_l}{b_l}$ are similarly ordered - in other words,\[(a_k-a_l)(b_k-b_l)\geq 0.\]Estimate $f(n)$ - in particular, is there a constant $c&gt;0$ such that $f(n)=(c+o(1))n$ for all large $n$?




The function $f(n)$ was first considered by Mayer \cite{Ma42}, who proved $f(n)\to \infty$ as $n\to \infty$. Erd\H{o}s \cite{Er43} proved $f(n)\gg n$. van Doorn \cite{vD25b} has proved that\[\left(\frac{1}{12}-o(1)\right)n\leq f(n) \leq \frac{1}{4}n+O(1),\]and conjectures that the upper bound is optimal.




References

[Er43] Erd\H{o}s, P., A note on {F}arey series . Quart. J. Math. Oxford Ser. (1943), 82--85.

[Ma42] Mayer, A. E., A mean value theorem concerning {F}arey series . Quart. J. Math. Oxford Ser. (1942), 48--57.

[vD25b] W. van Doorn, Improved bounds for the Mayer-Erd\H{o}s phenomenon on similarly ordered Farey fractions . arXiv:2509.00121 (2025).

\noindent\rule{\linewidth}{0.4pt}


%% ============================================================
\section{Problem \#1011}
\label{prob:1011}

\noindent\textbf{Status:} OPEN \quad|\quad \textbf{Prize:} none \quad|\quad \textbf{Updated:} 2025-09-10

\noindent\textbf{Tags:} graph theory

\medskip

\noindent\textbf{Statement:}

Let $f_r(n)$ be minimal such that every graph on $n$ vertices with $\geq f_r(n)$ edges and chromatic number $\geq r$ contains a triangle. Determine $f_r(n)$.




Tur\'{a}n's theorem implies $f_2(n)=\lfloor n^2/4\rfloor+1$. Erd\H{o}s and Gallai \cite{Er62d} proved $f_3(n)=\lfloor \frac{1}{4}(n-1)^2\rfloor+2$. Simonovits showed in his PhD thesis (see the discussion on p. 358 of \cite{Si74}) that\[f_r(n)=\frac{n^2}{4}-\frac{g(r)}{2}{n}+O(1),\]where $g(r)$ is the largest $m$ such that, for any triangle-free graph with chromatic number $\geq r$, at least $m$ vertices of $G$ need to be removed to obtain a bipartite graph. Simonovits \cite{Si74} notes\[\frac{\log r}{\log\log r}r^2 \ll g(r) \ll (\log r)^2r^2.\]Hunter in the comments has noted that other results imply $g(r)\asymp r^2\log r$ - in fact\[(1/2-o(1))r^2\log r\leq g(r)\leq (2+o(1))r^2\log r.\]The lower bound follows from work of Davies and Illingworth \cite{DaIl22} (see [1104] ). The upper bound follows from work of Hefty, Horn, King, and Pfender \cite{HHKP25} on $R(3,k)$. Ren, Wang, Wang, and Yang \cite{RWWY24} showed that, for $n\geq 150$,\[f_4(n)=\left\lfloor\frac{(n-3)^2}{4}\right\rfloor+6.\]




References

[DaIl22] Davies, Ewan and Illingworth, Freddie, The {$\chi$}-{R}amsey problem for triangle-free graphs . SIAM J. Discrete Math. (2022), 1124--1134.

[Er62d] Erd\H{o}s, P., On a theorem of {R}ademacher-Tur\'{a}n . Illinois J. Math. (1962), 122--127.

[HHKP25] Z. Hefty, P. Horn, D. King, and F. Pfender, Improving $R(3,k)$ in just two bites . arXiv:2510.19718 (2025).

[RWWY24] S. Ren, J. Wang, S. Wang, and W. Yang, Extremal triangle-free graphs with chromatic number at least four . arXiv:2404.07486 (2024).

[Si74] Simonovits, M., Extermal graph problems with symmetrical extremal graphs. {A}dditional chromatic conditions . Discrete Math. (1974), 349--376.

\noindent\rule{\linewidth}{0.4pt}


%% ============================================================
\section{Problem \#1013}
\label{prob:1013}

\noindent\textbf{Status:} OPEN \quad|\quad \textbf{Prize:} none \quad|\quad \textbf{Updated:} 2025-09-10

\noindent\textbf{Tags:} graph theory

\medskip

\noindent\textbf{Statement:}

Let $h_3(k)$ be the minimal $n$ such that there exists a triangle-free graph on $n$ vertices with chromatic number $k$. Find an asymptotic for $h_3(k)$, and also prove\[\lim_{k\to \infty}\frac{h_3(k+1)}{h_3(k)}=1.\]




The function $h_3(k)$ is dual to the function $f(n)$ considered in [1104] , in that $h_3(k)= n$ if and only if $n$ is minimal such that $f(n)=k$. Graver and Yackel \cite{GrYa68} proved\[h_3(k)\gg \frac{\log k}{\log\log k}k^2.\]The bounds for $f(n)$ from [1104] imply\[\left(\frac{1}{2}-o(1)\right)k^2\log k\leq h_3(k) \leq (1+o(1))k^2\log k.\]See also [920] for a generalisation to $K_r$-free graphs.




References

[GrYa68] Graver, Jack E. and Yackel, James, Some graph theoretic results associated with Ramsey's theorem . J. Combinatorial Theory (1968), 125--175.

\noindent\rule{\linewidth}{0.4pt}


%% ============================================================
\section{Problem \#1016}
\label{prob:1016}

\noindent\textbf{Status:} OPEN \quad|\quad \textbf{Prize:} none \quad|\quad \textbf{Updated:} 2025-09-10

\noindent\textbf{Tags:} graph theory, cycles

\medskip
\noindent\textit{pancyclic graphs}


\noindent\textbf{Statement:}

Let $h(n)$ be minimal such that there is a graph on $n$ vertices with $n+h(n)$ edges which contains a cycle on $k$ vertices, for all $3\leq k\leq n$. Estimate $h(n)$. In particular, is it true that\[h(n) \geq \log_2n+\log_*n-O(1),\]where $\log_*n$ is the iterated logarithmic function?




Such graphs are called pancyclic. A problem of Bondy \cite{Bo71}, who claimed a proof (without details) of\[\log_2(n-1)-1\leq h(n) \leq \log_2n+\log_*n+O(1).\]Erd\H{o}s \cite{Er71} believed the upper bound is closer to the truth, but could not even prove $h(n)-\log_2n\to \infty$. A proof of the above lower bound is provided by Griffin \cite{Gr13}. The first published proof of the upper bound appears to be in Chapter 4.5 of George, Khodkar, and Wallis \cite{GKW16}.




References

[Bo71] Bondy, J. A., Pancyclic graphs. I . J. Combinatorial Theory Ser. B (1971), 80--84.

[Er71] Erd\H{o}s, P., Some unsolved problems in graph theory and combinatorial analysis . Combinatorial Mathematics and its Applications (Proc. Conf., Oxford, 1969) (1971), 97-109.

[GKW16] George, John C. and Khodkar, Abdollah and Wallis, W. D., Pancyclic and bipancyclic graphs . (2016), xii+108.

[Gr13] S. Griffin, Minimal Pancyclicity . arXiv:1312.0274 (2013).

\noindent\rule{\linewidth}{0.4pt}


%% ============================================================
\section{Problem \#1017}
\label{prob:1017}

\noindent\textbf{Status:} OPEN \quad|\quad \textbf{Prize:} none \quad|\quad \textbf{Updated:} 2025-09-12

\noindent\textbf{Tags:} graph theory

\medskip

\noindent\textbf{Statement:}

Let $f(n,k)$ be such that every graph on $n$ vertices and $k$ edges can be partitioned into at most $f(n,k)$ edge-disjoint complete graphs. Estimate $f(n,k)$ for $k&#62;n^2/4$.




The function $f(n,k)$ is sometimes called the clique partition number. Erd\H{o}s, Goodman, and P\'{o}sa \cite{EGP66} proved that $f(n,k)\leq n^2/4$ for all $k$ (and in fact the complete graphs can be taken to be edges and triangles), which is best possible in general, as witnessed for example by a complete bipartite graph. In \cite{Er71} Erd\H{o} asks vaguely whether this result can be 'sharpened' for $k&gt;n^2/4$. Lov\'{a}sz \cite{Lo68} proved that every graph on $n$ vertices and $k$ edges is the union of $\binom{n}{2}-k+t$ complete graphs, where $t$ is maximal such that $t^2-t\leq \binom{n}{2}-k$, but without the assumption that the complete graphs are edge disjoint. Lov\'{a}sz's result is sharp in many cases. If $k&gt;n^2/4$ and the graph contains no $K_4$ then this is equivalent to finding the minimum number of edge disjoint triangles. This special case was also asked about by Erd\H{o}s. A complete answer was provided by Gy\"{o}ri and Keszegh \cite{GyKe17}, who proved that every $K_4$-free graph with $n$ vertices and $\lfloor n^2/4\rfloor+m$ edges contins $m$ pairwise edge disjoint triangles. See also [184] for an analogous problem decomposing into edges and cycles, and [583] for decomposing into paths. The clique partition problem for chordal graphs is the subject of [81] .




References

[EGP66] Erd\H{o}s, Paul and Goodman, A. W. and P\'{o}sa, Lajos, The representation of a graph by set intersections . Canadian J. Math. (1966), 106-112.

[Er71] Erd\H{o}s, P., Some unsolved problems in graph theory and combinatorial analysis . Combinatorial Mathematics and its Applications (Proc. Conf., Oxford, 1969) (1971), 97-109.

[GyKe17] Gy\H{o}ri, Ervin and Keszegh, Bal\'{a}zs, On the number of edge-disjoint triangles in {$K_4$}-free graphs . Combinatorica (2017), 1113--1124.

[Lo68] Lov\'{a}sz, L., On covering of graphs . Theory of Graphs (Proc. Colloq., Tihany, 1966) (1968), 231-236.

\noindent\rule{\linewidth}{0.4pt}


%% ============================================================
\section{Problem \#1029}
\label{prob:1029}

\noindent\textbf{Status:} OPEN \quad|\quad \textbf{Prize:} \$100 \quad|\quad \textbf{Updated:} 2025-09-13

\noindent\textbf{Tags:} graph theory, ramsey theory

\medskip

\noindent\textbf{Statement:}

If $R(k)$ is the Ramsey number for $K_k$, the minimal $n$ such that every $2$-colouring of the edges of $K_n$ contains a monochromatic copy of $K_k$, then\[\frac{R(k)}{k2^{k/2}}\to \infty.\]




In \cite{Er93} Erd\H{o}s offers \$100 for a proof of this and \$1000 for a disproof, but says 'this last offer is to some extent phoney: I am sure that [this] is true (but I have been wrong before).' Erd\H{o}s and Szekeres \cite{ErSz35} proved\[k2^{k/2} \ll R(k) \leq \binom{2k-1}{k-1}.\]One of the first applications of the probabilistic method pioneered by Erd\H{o}s gives\[R(k) \geq (1+o(1))\frac{1}{\sqrt{2}e}k2^{k/2},\]which Spencer \cite{Sp75} improved by a factor of $2$ to\[R(k) \geq (1+o(1))\frac{\sqrt{2}}{e}k2^{k/2}.\]See also [77] for a more general problem concerning $\lim R(k)^{1/k}$, and discussion of upper bounds for $R(k)$.




References

[Er93] Erd\H{o}s, Paul, Some of my favorite solved and unsolved problems in graph theory . Quaestiones Math. (1993), 333-350.

[ErSz35] Erd\H{o}s, P. and Szekeres, G., A combinatorial problem in geometry . Compos. Math. (1935), 463-470.

[Sp75] Spencer, Joel, Ramsey's theorem---a new lower bound . J. Combinatorial Theory Ser. A (1975), 108--115.

\noindent\rule{\linewidth}{0.4pt}


%% ============================================================
\section{Problem \#1030}
\label{prob:1030}

\noindent\textbf{Status:} OPEN \quad|\quad \textbf{Prize:} none \quad|\quad \textbf{Updated:} 2025-09-13

\noindent\textbf{Tags:} graph theory, ramsey theory

\medskip

\noindent\textbf{Statement:}

Let $R(k,l)$ be the usual Ramsey number: the smallest $n$ such that if the edges of $K_n$ are coloured red and blue then there exists either a red $K_k$ or a blue $K_l$. Prove the existence of some $c&#62;0$ such that\[\lim_{k\to \infty}\frac{R(k+1,k)}{R(k,k)}&#62; 1+c.\]




A problem of Erd\H{o}s and S\'{o}s, who could not even prove whether $R(k+1,k)-R(k,k)&gt;k^c$ for any $c&gt;1$. It is trivial that $R(k+1,k)-R(k,k)\geq k-2$. Burr, Erd\H{o}s, Faudree, and Schelp \cite{BEFS89} proved\[R(k+1,k)-R(k,k)\geq 2k-5.\]See also [544] for a similar question concerning $R(3,k)$, and [1014] for the general off-diagonal case.




References

[BEFS89] Burr, S. A. and Erd\H{o}s, P. and Faudree, R. J. and Schelp, R. H., On the difference between consecutive {R}amsey numbers . Utilitas Math. (1989), 115--118.

\noindent\rule{\linewidth}{0.4pt}


%% ============================================================
\section{Problem \#1032}
\label{prob:1032}

\noindent\textbf{Status:} OPEN \quad|\quad \textbf{Prize:} none \quad|\quad \textbf{Updated:} 2025-09-13

\noindent\textbf{Tags:} graph theory, chromatic number

\medskip

\noindent\textbf{Statement:}

We say that a graph is $4$-chromatic critical if it has chromatic number $4$, and removing any edge decreases the chromatic number to $3$. Is there, for arbitrarily large $n$, a $4$-chromatic critical graph on $n$ vertices with minimum degree $\gg n$?




In \cite{Er93} Erd\H{o}s said he asked this 'more than 20 years ago'. Dirac gave an example of a $6$-chromatic critical graph with minimum degree $&gt;n/2$. This problem is also open for $5$-chromatic critical graphs. Simonovits \cite{Si72} and Toft \cite{To72} independently constructed $4$-chromatic critical graphs with minimum degree $\gg n^{1/3}$. Toft conjectured that a $4$-chromatic critical graph on $n$ vertices has at least $(\frac{5}{3}+o(1))n$ vertices, and has examples to show this would be the best possible. See also [917] and [944] .




References

[Er93] Erd\H{o}s, Paul, Some of my favorite solved and unsolved problems in graph theory . Quaestiones Math. (1993), 333-350.

[Si72] Simonovits, M., On colour-critical graphs . Studia Sci. Math. Hungar. (1972), 67--81.

[To72] Toft, B., Two theorems on critical {$4$}-chromatic graphs . Studia Sci. Math. Hungar. (1972), 83--89.

\noindent\rule{\linewidth}{0.4pt}


%% ============================================================
\section{Problem \#1033}
\label{prob:1033}

\noindent\textbf{Status:} OPEN \quad|\quad \textbf{Prize:} none \quad|\quad \textbf{Updated:} 2025-12-12

\noindent\textbf{Tags:} graph theory

\medskip

\noindent\textbf{Statement:}

Let $h(n)$ be such that every graph on $n$ vertices with $&#62;n^2/4$ many edges contains a triangle whose vertices have degrees summing to at least $h(n)$. Estimate $h(n)$. In particular, is it true that\[h(n)\geq (2(\sqrt{3}-1)-o(1))n?\]




A conjecture of Bollob\'{a}s and Erd\H{o}s. In \cite{Er82e} it is asked whether $h(n)\geq \frac{3}{2}n$. It is now known that\[2(\sqrt{3}-1)n +O(1)\geq h(n) \geq \frac{21}{16}n.\]The upper bound is due to Erd\H{o}s and Laskar \cite{ErLa85}, and the lower bound is due to Fan \cite{Fa88}. (Note that $2(\sqrt{3}-1)\approx 1.464$ and $21/16=1.3125$.) The upper bound is not made explicit in \cite{ErLa85}, which is concerned with chordal subgraphs (i.e. subgraphs which do not contain any induced cycle on $&gt;3$ vertices). The connection to this problem is that taking a triangle together with all incident edges forms a chordal subgraph of $G$. The construction of the upper bound is made more explicit in \cite{Fa88}. In brief, let $m=\lfloor n^2/4\rfloor+1$ and take a complete bipartite graph between $k$ and $l$ vertices where $k+l=n$ and $k=cn$. If $c$ is chosen such that $m-kl\leq k^2/4$ then we can add $m-kl$ edges between the $k$ vertices creating no triangles, while having maximum degree (amongst these vertices) $\lceil 2\frac{m-kl}{k}\rceil$. This creates a graph $G$ on $n$ vertices with $m$ edges. Every triangle uses exactly one vertex from the side with $l$ vertices and two vertices from the side with $k$ vertices, and hence the sum of degrees is at most\[k+2\left\lceil 2\frac{m-kl}{k}\right\rceil+2l=(c^{-1}+3c-2)n+O(1).\]This is minimised at $c=1/\sqrt{3}$, where it is $2(\sqrt{3}-1)n+O(1)$ (and a short calculation verifies that with this choice of $c$ we do indeed have $m-kl\leq k^2/4$). More generally, let $\Delta_r(n,m)$ be such that every graph $G$ with $n$ vertices and $m$ edges contains a clique on $r$ vertices, say $x_1,\ldots,x_r$, such that\[d(x_1)+\cdots+d(x_r)\geq \Delta_r(n,m).\]Let $r\geq 2$ and let $t_r(n)$ be the Tur\'{a}n number (the maximal number of edges in a graph on $n$ vertices with no $K_{r+1}$). One can ask more generally about the behaviour of $\Delta_r(n,m)$ for $m&gt;t_{r-1}(n)$ (when $m\leq t_{r-1}(n)$ this is trivial since there may be no clique on $r$ vertices at all). The behaviour of $\Delta_r(n,m)$ in the regime $m\geq t_r(n)$ is the subject of [904] . Erd\H{o}s (see \cite{Fa92}) proved that, for every $\epsilon&gt;0$, there exists $\delta&gt;0$ such that if $t_{r-1}(n)&lt;m&lt;t_r(n)-\delta n^2$ then\[\Delta_r(n,m)\leq (1-\epsilon)\frac{2rm}{n}.\]Bollob\'{a}s and Nikiforov \cite{BoNi05} proved that, for every $\epsilon&gt;0$, there exists $\delta&gt;0$ such that if $m&gt;t_r(n)-\delta n^2$ then\[\Delta_r(n,m)\geq (1-\epsilon)\frac{2rm}{n}.\]




References

[BoNi05] Bollob\'{a}s, B\'ela and Nikiforov, Vladimir, The sum of degrees in cliques . Electron. J. Combin. (2005), Note 21, 10.

[Er82e] Erd\H{o}s, Paul, Some of my favourite problems which recently have been solved . (1982), 59--79.

[ErLa85] Erd\H{o}s, Paul and Laskar, Renu, A note on the size of a chordal subgraph . Congr. Numer. (1985), 81--86.

[Fa88] Fan, Genghua, Degree sum for a triangle in a graph . J. Graph Theory (1988), 249--263.

[Fa92] Faudree, Ralph J., Complete subgraphs with large degree sums . J. Graph Theory (1992), 327--334.

\noindent\rule{\linewidth}{0.4pt}


%% ============================================================
\section{Problem \#1035}
\label{prob:1035}

\noindent\textbf{Status:} OPEN \quad|\quad \textbf{Prize:} none \quad|\quad \textbf{Updated:} 2025-12-26

\noindent\textbf{Tags:} graph theory

\medskip

\noindent\textbf{Statement:}

Is there a constant $c&#62;0$ such that every graph on $2^n$ vertices with minimum degree $&#62;(1-c)2^n$ contains the $n$-dimensional hypercube $Q_n$?




Erd\H{o}s \cite{Er93} says 'if the conjecture is false, two related problems could be asked': {UL} {LI}Determine or estimate the smallest $m&gt;2^n$ such that every graph on $m$ vertices with minimum degree $&gt;(1-c)2^n$ contains a $Q_n$, and {/LI} {LI}For which $u_n$ is it true that every graph on $2^n$ vertices with minimum degree $&gt;2^n-u_n$ contains a $Q_n$.{/LI} {/UL} See also [576] for the extremal number of edges that guarantee a $Q_n$.




References

[Er93] Erd\H{o}s, Paul, Some of my favorite solved and unsolved problems in graph theory . Quaestiones Math. (1993), 333-350.

\noindent\rule{\linewidth}{0.4pt}


%% ============================================================
\section{Problem \#1038}
\label{prob:1038}

\noindent\textbf{Status:} OPEN \quad|\quad \textbf{Prize:} none \quad|\quad \textbf{Updated:} 2025-09-15

\noindent\textbf{Tags:} analysis

\medskip

\noindent\textbf{Statement:}

Determine the infimum and supremum of\[\lvert \{ x\in \mathbb{R} : \lvert f(x)\rvert &#60; 1\}\rvert\]as $f\in \mathbb{R}[x]$ ranges over all non-constant monic polynomials, all of whose roots are real and in the interval $[-1,1]$.




A problem of Erd\H{o}s, Herzog, and Piranian \cite{EHP58}, who proved that the measure of the set in question is always at most $2\sqrt{2}$ under the assumption that all the roots are in $\{-1,1\}$, and conjecture this is the best possible upper bound. They also note that the infimum of the set in question is less than $2$, as witnessed by $f(x)=(x+1)(x-1)^m$ for $m\geq 3$. They further note that if the roots are restricted to $[-2,2]$ then the infimum is zero, as witnessed by a small perturbation of the Chebyshev polynomials. They further conjectured that, if the roots are restricted to $[-2,2]$, then\[\lvert \{ x\in \mathbb{R} : \lvert f(x)\rvert &lt; 1\}\rvert\geq n^{-c}\]for an absolute constant $c&gt;0$. This was proved by Pommerenke \cite{Po61}, who in fact showed that this set must contain an interval of width $\gg n^{-4}$. The current best known bounds (see the discussion in the comments) are\[1.519\approx 2^{4/3}-1\leq \inf \leq 1.835\cdots\]and\[\sup = 2\sqrt{2}\approx 2.828.\]




References

[EHP58] Erd\H{o}s, P. and Herzog, F. and Piranian, G., Metric properties of polynomials . J. Analyse Math. (1958), 125-148.

[Po61] Pommerenke, Ch., On metric properties of complex polynomials . Michigan Math. J. (1961), 97-115.

\noindent\rule{\linewidth}{0.4pt}


%% ============================================================
\section{Problem \#1039}
\label{prob:1039}

\noindent\textbf{Status:} OPEN \quad|\quad \textbf{Prize:} none \quad|\quad \textbf{Updated:} 2025-09-15

\noindent\textbf{Tags:} analysis

\medskip

\noindent\textbf{Statement:}

Let $f(z)=\prod_{i=1}^n(z-z_i)\in \mathbb{C}[z]$ with $\lvert z_i\rvert \leq 1$ for all $i$. Let $\rho(f)$ be the radius of the largest disc which is contained in $\{z: \lvert f(z)\rvert&#60; 1\}$. Determine the behaviour of $\rho(f)$. In particular, is it always true that $\rho(f)\gg 1/n$?




A problem of Erd\H{o}s, Herzog, and Piranian, who note that $f(z)=z^n-1$ has $\rho(f) \leq \frac{\pi/2}{n}$. Pommerenke \cite{Po61} proved that\[\rho(f) \geq \frac{1}{2en^2}.\]Krishnapur, Lundberg, and Ramachandran \cite{KLR25} proved\[\rho(f) \gg \frac{1}{n\sqrt{\log n}}.\]




References

[KLR25] M. Krishnapur, E. Lundberg, and K. Ramachandran, On the area of polynomial lemniscates . arXiv:2503.18270 (2025).

[Po61] Pommerenke, Ch., On metric properties of complex polynomials . Michigan Math. J. (1961), 97-115.

\noindent\rule{\linewidth}{0.4pt}


%% ============================================================
\section{Problem \#1040}
\label{prob:1040}

\noindent\textbf{Status:} OPEN \quad|\quad \textbf{Prize:} none \quad|\quad \textbf{Updated:} 2025-09-15

\noindent\textbf{Tags:} analysis

\medskip

\noindent\textbf{Statement:}

Let $F\subseteq \mathbb{C}$ be a closed infinite set, and let $\mu(F)$ be the infimum of\[\lvert \{ z: \lvert f(z)\rvert &#60; 1\}\rvert,\]as $f$ ranges over all polynomials of the shape $\prod (z-z_i)$ with $z_i\in F$. Is $\mu(F)$ determined by the transfinite diameter of $F$? In particular, is $\mu(F)=0$ whenever the transfinite diameter of $F$ is $\geq 1$?




A problem of Erd\H{o}s, Herzog, and Piranian \cite{EHP58}, who show that the answer is yes if $F$ is a line segment or disc, and that if the transfinite diameter is $&#60;1$ then $\{ z: \lvert f(z)\rvert &#60; 1\}$ always contains a disc of radius $\gg_F 1$. Erd\H{o}s and Netanyahu \cite{ErNe73} proved that if $F$ is also bounded and connected, with transfinite diameter $0&#60;c&#60;1$, then $\{ z: \lvert f(z)\rvert &#60; 1\}$ always contains a disc of radius $\gg_c 1$. The transfinite diameter of $F$, also known as the logarithmic capacity, is defined by\[\rho(F)=\lim_{n\to \infty}\sup_{z_1,\ldots,z_n\in F}\left(\prod_{i&#60;j}\lvert z_i-z_j\rvert\right)^{1/\binom{n}{2}}.\]Aletheia \cite{Fe26} has shown that $\mu(F)$ is not determined by the transfinite diameter of $F$, by producing two distinct closed infinite sets $F_1$ and $F_2$, both of which have transfinite diameter $0$, and yet $\mu(F_1)\geq \pi/4$ while $\mu(F_2)$ can be made arbitrarily close to $0$.




References

[EHP58] Erd\H{o}s, P. and Herzog, F. and Piranian, G., Metric properties of polynomials . J. Analyse Math. (1958), 125-148.

[ErNe73] Erd\H{o}s, P. and Netanyahu, E., A remark on polynomials and the transfinite diameter . Israel J. Math. (1973), 23--25.

[Fe26] T. Feng et al, Semi-Autonomous Mathematics Discovery with Gemini: A Case Study on the Erd\H{o}s Problems . arXiv:2601.22401 (2026).

\noindent\rule{\linewidth}{0.4pt}


%% ============================================================
\section{Problem \#1045}
\label{prob:1045}

\noindent\textbf{Status:} OPEN \quad|\quad \textbf{Prize:} none \quad|\quad \textbf{Updated:} 2026-03-14

\noindent\textbf{Tags:} analysis

\medskip

\noindent\textbf{Statement:}

Let $z_1,\ldots,z_n\in \mathbb{C}$ with $\lvert z_i-z_j\rvert\leq 2$ for all $i,j$, and\[\Delta(z_1,\ldots,z_n)=\prod_{i\neq j}\lvert z_i-z_j\rvert.\]What is the maximum possible value of $\Delta$? Is it maximised by taking the $z_i$ to be the vertices of a regular polygon?




A problem of Erd\H{o}s, Herzog, and Piranian \cite{EHP58}, who proved that, for any monic polynomial $f$, if $\{ z: \lvert f(z)\rvert &lt;1\}$ is connected and $f$ has roots $z_1,\ldots,z_n$ then $\prod_{i\neq j}\lvert z_i-z_j\rvert &lt;n^n$. The value of $\Delta$ when the $z_i$ are the vertices of a regular polygon is $n^n$ when $n$ is even and\[\cos(\pi/2n)^{-n(n-1)}n^n \sim e^{\pi^2/8}n^n\]when $n$ is odd. Pommerenke \cite{Po61} proved that $\Delta \leq 2^{O(n)}n^n$ for all $z_i$ with $\lvert z_i-z_j\rvert \leq 2$. Hu and Tang (see the comments) found examples when $n=4$ and $n=6$ that show that vertices of a regular polygon do not maximise $\Delta$. Cambie (also in the comments) showed that, in general, the vertices of a regular polygon are not a maximiser for all even $n\geq 4$. There is a lot of discussion of this problem in the comments (see also the papers of Sothanaphan \cite{So25} and Cambie, Decadt, Dong, Hu, and Tang \cite{CDDHT26}). It is now known that, for even $n$,\[\liminf \frac{\max \Delta}{n^n}\geq C\]for some $C&gt;0$. The best value of $C$ available is $C\approx 1.268$, proved in \cite{CDDHT26}. If we restrict to $n$ divisible by $6$ this can be improved to $C\approx 1.304$. It remains possible that the regular polygon is a maximiser for odd $n$; this is explicitly conjectured in \cite{CDDHT26}. It is not yet known whether, for odd $n$,\[\lim \frac{\max \Delta}{n^n}=e^{\pi^2/8}\approx 3.433.\]The paper \cite{CDDHT26} also contains a number of structural results about the form of an extremal set of points.




References

[CDDHT26] S. Cambie, A. Decadt, Y. Dong, T. Hu, and Q. Tang, On the maximum product of distances of diameter 2 point sets . arXiv:2603.07088 (2026).

[EHP58] Erd\H{o}s, P. and Herzog, F. and Piranian, G., Metric properties of polynomials . J. Analyse Math. (1958), 125-148.

[Po61] Pommerenke, Ch., On metric properties of complex polynomials . Michigan Math. J. (1961), 97-115.

[So25] N. Sothanaphan, An improved lower bound to Erdos' problem concerning products of distances for fixed diameter . arXiv:2512.14251 (2025).

\noindent\rule{\linewidth}{0.4pt}


%% ============================================================
\section{Problem \#1049}
\label{prob:1049}

\noindent\textbf{Status:} OPEN \quad|\quad \textbf{Prize:} none \quad|\quad \textbf{Updated:} 2025-09-28

\noindent\textbf{Tags:} irrationality

\medskip

\noindent\textbf{Statement:}

Let $t&#62;1$ be a rational number. Is\[\sum_{n=1}^\infty\frac{1}{t^n-1}=\sum_{n=1}^\infty \frac{\tau(n)}{t^n}\]irrational, where $\tau(n)$ counts the divisors of $n$?




A conjecture of Chowla. Erd\H{o}s \cite{Er48} proved that this is true if $t\geq 2$ is an integer.




References

[Er48] Erd\H{o}s, P., On arithmetical properties of Lambert series . J. Indian Math. Soc. (N.S.) (1948), 63-66.

\noindent\rule{\linewidth}{0.4pt}


%% ============================================================
\section{Problem \#1052}
\label{prob:1052}

\noindent\textbf{Status:} OPEN \quad|\quad \textbf{Prize:} \$10 \quad|\quad \textbf{Updated:} 2025-09-28

\noindent\textbf{Tags:} number theory

\medskip
\noindent\textit{unitary perfect numbers}


\noindent\textbf{Statement:}

A unitary divisor of $n$ is $d\mid n$ such that $(d,n/d)=1$. A number $n\geq 1$ is a unitary perfect number if it is the sum of its unitary divisors (aside from $n$ itself). Are there only finite many unitary perfect numbers?




Guy \cite{Gu04} reports that Carlitz, Erd\H{o}s, and Subbarao offer \$10 for settling this question, and that Subbarao offers 10 cents for each new example. There are no odd unitary perfect numbers. There are five known unitary perfect numbers ( A002827 in the OEIS):\[6, 60, 90, 87360, 146361946186458562560000.\]This is problem B3 in Guy's collection \cite{Gu04}.




References

[Gu04] Guy, Richard K., Unsolved problems in number theory . (2004), xviii+437.

\noindent\rule{\linewidth}{0.4pt}


%% ============================================================
\section{Problem \#1053}
\label{prob:1053}

\noindent\textbf{Status:} OPEN \quad|\quad \textbf{Prize:} none \quad|\quad \textbf{Updated:} 2025-09-28

\noindent\textbf{Tags:} number theory

\medskip
\noindent\textit{multiply perfect numbers}


\noindent\textbf{Statement:}

Call a number $k$-perfect if $\sigma(n)=kn$, where $\sigma(n)$ is the sum of the divisors of $n$. Must $k=o(\log\log n)$?




A question of Erd\H{o}s, as reported in problem B2 of Guy's collection \cite{Gu04}. Guy further writes 'It has even been suggested that there may be only finitely many $k$-perfect numbers with $k\geq 3$.' The largest $k$ for which a $k$-perfect number has been found is $k=11$ - see this page for more information. These are known as multiply perfect numbers . When $k=2$ this is the definition of a perfect number .




References

[Gu04] Guy, Richard K., Unsolved problems in number theory . (2004), xviii+437.

\noindent\rule{\linewidth}{0.4pt}


%% ============================================================
\section{Problem \#1054}
\label{prob:1054}

\noindent\textbf{Status:} OPEN \quad|\quad \textbf{Prize:} none \quad|\quad \textbf{Updated:} 2025-09-28

\noindent\textbf{Tags:} number theory, divisors

\medskip

\noindent\textbf{Statement:}

Let $f(n)$ be the minimal integer $m$ such that $n$ is the sum of the $k$ smallest divisors of $m$ for some $k\geq 1$. Is it true that $f(n)=o(n)$? Or is this true only for almost all $n$, and $\limsup f(n)/n=\infty$?




A question of Erd\H{o}s reported in problem B2 of Guy's collection \cite{Gu04}. The function $f(n)$ is undefined for $n=2$ and $n=5$, but is likely well-defined for all $n\geq 6$ (which would follow from a strong form of Goldbach's conjecture). The sequence of values of $f(n)$ is given by A167485 in the OEIS. See also [468] . The strong claim that $f(n)=o(n)$ was disproved by Tao in the comments to [468] , in which he proves that the upper density of $\{ n : f(n)\leq \delta n\}$ is $\ll \delta^2$.




References

[Gu04] Guy, Richard K., Unsolved problems in number theory . (2004), xviii+437.

\noindent\rule{\linewidth}{0.4pt}


%% ============================================================
\section{Problem \#1055}
\label{prob:1055}

\noindent\textbf{Status:} OPEN \quad|\quad \textbf{Prize:} none \quad|\quad \textbf{Updated:} 2025-09-28

\noindent\textbf{Tags:} number theory, primes

\medskip

\noindent\textbf{Statement:}

A prime $p$ is in class $1$ if the only prime divisors of $p+1$ are $2$ or $3$. In general, a prime $p$ is in class $r$ if every prime factor of $p+1$ is in some class $\leq r-1$, with equality for at least one prime factor. Are there infinitely many primes in each class? If $p_r$ is the least prime in class $r$, then how does $p_r^{1/r}$ behave?




A classification due to Erd\H{o}s and Selfridge. It is easy to prove that the number of primes $\leq n$ in class $r$ is at most $n^{o(1)}$. The sequence $p_r$ begins $2,13,37,73,1021$ ( A005113 in the OEIS). Erd\H{o}s thought $p_r^{1/r}\to \infty$, while Selfridge thought it quite likely to be bounded. A similar question can be asked replacing $p+1$ with $p-1$. This is problem A18 in Guy's collection \cite{Gu04}.




References

[Gu04] Guy, Richard K., Unsolved problems in number theory . (2004), xviii+437.

\noindent\rule{\linewidth}{0.4pt}


%% ============================================================
\section{Problem \#1056}
\label{prob:1056}

\noindent\textbf{Status:} OPEN \quad|\quad \textbf{Prize:} none \quad|\quad \textbf{Updated:} 2025-09-28

\noindent\textbf{Tags:} number theory

\medskip

\noindent\textbf{Statement:}

Let $k\geq 2$. Does there exist a prime $p$ and consecutive intervals $I_1,\ldots,I_k$ such that\[\prod_{n\in I_i}n \equiv 1\pmod{p}\]for all $1\leq i\leq k$?




This is problem A15 in Guy's collection \cite{Gu04}, where he reports that in a letter in 1979 Erd\H{o}s observed that\[3\cdot 4\equiv 5\cdot 6\cdot 7\equiv 1\pmod{11},\]establishing the case $k=2$. Makowski \cite{Ma83} found, for $k=3$,\[2\cdot 3\cdot 4\cdot 5\equiv 6\cdot 7\cdot 8\cdot 9\cdot 10\cdot 11\equiv 12\cdot 13\cdot 14\cdot 15\equiv 1\pmod{17}.\]Noll and Simmons asked, more generally, whether there are solutions to $q_1!\equiv\cdots \equiv q_k!\pmod{p}$ for arbitrarily large $k$ (with $q_1&#60;\cdots&#60;q_k$).




References

[Gu04] Guy, Richard K., Unsolved problems in number theory . (2004), xviii+437.

[Ma83] M\polhk akowski, Andrzej, On a number-theoretical problem of {E}rd\H{o}s . Elem. Math. (1983), 101--102.

\noindent\rule{\linewidth}{0.4pt}


%% ============================================================
\section{Problem \#1057}
\label{prob:1057}

\noindent\textbf{Status:} OPEN \quad|\quad \textbf{Prize:} none \quad|\quad \textbf{Updated:} 2025-09-28

\noindent\textbf{Tags:} number theory

\medskip
\noindent\textit{Carmichael numbers}


\noindent\textbf{Statement:}

Let $C(x)$ count the number of Carmichael numbers in the interval $[1,x]$. Is it true that $C(x)=x^{1-o(1)}$?




Erd\H{o}s \cite{Er56c} proved\[C(x) &lt; x \exp\left(-c \frac{\log x\log\log\log x}{\log\log x}\right)\]for some constant $c&gt;0$. Pomerance \cite{Po89} gave a heuristic suggesting that this is the true order of growth, and in fact\[C(x)= x \exp\left(-(1+o(1))\frac{\log x\log\log\log x}{\log\log x}\right).\]Alford, Granville, and Pomerance \cite{AGP94} proved that $C(x)\to \infty$, and in fact $C(x)&#62;x^{2/7}$ for large $x$. The lower bound\[C(x)&#62; x^{0.33336704}\]was proved by Harman \cite{Ha08}. This exponent was improved to $0.3389$ by Lichtman \cite{Li22}. Korselt observed that $n$ being a Carmichael number is equivalent to $n$ being squarefree and $p-1\mid n-1$ for all primes $p\mid n$. This is discussed in problem A13 of Guy's collection \cite{Gu04}.




References

[AGP94] Alford, W. R. and Granville, Andrew and Pomerance, Carl, There are infinitely many {C}armichael numbers . Ann. of Math. (2) (1994), 703--722.

[Er56c] Erd\H{o}s, P., On pseudoprimes and {C}armichael numbers . Publ. Math. Debrecen (1956), 201--206.

[Gu04] Guy, Richard K., Unsolved problems in number theory . (2004), xviii+437.

[Ha08] Harman, Glyn, Watt's mean value theorem and {C}armichael numbers . Int. J. Number Theory (2008), 241--248.

[Li22] J. D. Lichtman, Primes in arithmetic progressions to large moduli and shifted primes without large prime factors . arXiv:2211.09641 (2022).

[Po89] Pomerance, Carl, Two methods in elementary analytic number theory . (1989), 135--161.

\noindent\rule{\linewidth}{0.4pt}


%% ============================================================
\section{Problem \#1059}
\label{prob:1059}

\noindent\textbf{Status:} OPEN \quad|\quad \textbf{Prize:} none \quad|\quad \textbf{Updated:} 2025-09-28

\noindent\textbf{Tags:} number theory, primes

\medskip

\noindent\textbf{Statement:}

Are there infinitely many primes $p$ such that $p-k!$ is composite for each $k$ such that $1\leq k!&#60;p$?




A question of Erd\H{o}s reported in problem A2 of Guy's collection \cite{Gu04}. Examples are $p=101$ and $p=211$. Erd\H{o}s suggested it may be easier to show that there are infinitely many $n$ such that, if $l!&lt;n\leq (l+1)!$, then all the prime factors of $n$ are $&gt;l$, and all the numbers $n-k!$ are composite for $1\leq k\leq l$.




References

[Gu04] Guy, Richard K., Unsolved problems in number theory . (2004), xviii+437.

\noindent\rule{\linewidth}{0.4pt}


%% ============================================================
\section{Problem \#1060}
\label{prob:1060}

\noindent\textbf{Status:} OPEN \quad|\quad \textbf{Prize:} none \quad|\quad \textbf{Updated:} 2025-09-28

\noindent\textbf{Tags:} number theory

\medskip

\noindent\textbf{Statement:}

Let $f(n)$ count the number of solutions to $k\sigma(k)=n$, where $\sigma(k)$ is the sum of divisors of $k$. Is it true that $f(n)\leq n^{o(\frac{1}{\log\log n})}$? Perhaps even $\leq (\log n)^{O(1)}$?




This is discussed in problem B11 of Guy's collection \cite{Gu04}.




References

[Gu04] Guy, Richard K., Unsolved problems in number theory . (2004), xviii+437.

\noindent\rule{\linewidth}{0.4pt}


%% ============================================================
\section{Problem \#1061}
\label{prob:1061}

\noindent\textbf{Status:} OPEN \quad|\quad \textbf{Prize:} none \quad|\quad \textbf{Updated:} 2025-09-28

\noindent\textbf{Tags:} number theory

\medskip

\noindent\textbf{Statement:}

How many solutions are there to\[\sigma(a)+\sigma(b)=\sigma(a+b)\]with $a+b\leq x$, where $\sigma$ is the sum of divisors function? Is it $\sim cx$ for some constant $c&#62;0$?




A question of Erd\H{o}s reported in problem B15 of Guy's collection \cite{Gu04}.




References

[Gu04] Guy, Richard K., Unsolved problems in number theory . (2004), xviii+437.

\noindent\rule{\linewidth}{0.4pt}


%% ============================================================
\section{Problem \#1062}
\label{prob:1062}

\noindent\textbf{Status:} OPEN \quad|\quad \textbf{Prize:} none \quad|\quad \textbf{Updated:} 2025-09-28

\noindent\textbf{Tags:} number theory

\medskip

\noindent\textbf{Statement:}

Let $f(n)$ be the size of the largest subset $A\subseteq \{1,\ldots,n\}$ such that there are no three distinct elements $a,b,c\in A$ such that $a\mid b$ and $a\mid c$. How large can $f(n)$ be? Is $\lim f(n)/n$ irrational?




The example $[m+1,3m+2]$ shows that $f(n)\geq\lceil \frac{2}{3}n\rceil$. Lebensold \cite{Le76} has shown that, for large $n$,\[0.6725 n \leq f(n) \leq 0.6736 n.\]This is problem B24 in Guy's collection \cite{Gu04}.




References

[Gu04] Guy, Richard K., Unsolved problems in number theory . (2004), xviii+437.

[Le76] Lebensold, Kenneth, A divisibility problem . Studies in Appl. Math. (1976/77), 291--294.

\noindent\rule{\linewidth}{0.4pt}


%% ============================================================
\section{Problem \#1063}
\label{prob:1063}

\noindent\textbf{Status:} OPEN \quad|\quad \textbf{Prize:} none \quad|\quad \textbf{Updated:} 2025-10-01

\noindent\textbf{Tags:} number theory

\medskip

\noindent\textbf{Statement:}

Let $k\geq 2$ and define $n_k\geq 2k$ to be the least value of $n$ such that $n-i$ divides $\binom{n}{k}$ for all but one $0\leq i&#60;k$. Estimate $n_k$.




A problem of Erd\H{o}s and Selfridge posed in \cite{ErSe83}. Erd\H{o}s and Selfridge noted (and a proof can be found in \cite{Mo85}) that if $n\geq 2k$ then there must exist at least one $0\leq i&#60;k$ such that $n-i$ does not divide $\binom{n}{k}$. We have $n_2=4$, $n_3=6$, $n_4=9$, and $n_5=12$. Monier \cite{Mo85} observed that $n_k\leq k!$ for $k\geq 3$, since $\binom{k!}{k}$ is divisible by $k!-i$ for $1\leq i&#60;k$. Cambie observes in the comments that this can be improved to\[n_k\leq k[2,3,\ldots,k-1]\leq e^{(1+o(1))k},\]where $[\cdots]$ is the least common multiple. This is discussed in problem B31 of Guy's collection \cite{Gu04}.




References

[ErSe83] Erdos, P. and Selfridge, J. L., Problem 6447 . Amer. Math. Monthly (1983), 710.

[Gu04] Guy, Richard K., Unsolved problems in number theory . (2004), xviii+437.

[Mo85] Monier, Jean-Marie, Problems and Solutions: Solutions of {A}dvanced {P}roblems: 6447 . Amer. Math. Monthly (1985), 435--436.

\noindent\rule{\linewidth}{0.4pt}


%% ============================================================
\section{Problem \#1065}
\label{prob:1065}

\noindent\textbf{Status:} OPEN \quad|\quad \textbf{Prize:} none \quad|\quad \textbf{Updated:} 2025-10-01

\noindent\textbf{Tags:} number theory

\medskip

\noindent\textbf{Statement:}

Are there infinitely many primes $p$ such that $p=2^kq+1$ for some prime $q$ and $k\geq 0$? Or $p=2^k3^lq+1$?




This is mentioned in problem B46 of Guy's collection \cite{Gu04}.




References

[Gu04] Guy, Richard K., Unsolved problems in number theory . (2004), xviii+437.

\noindent\rule{\linewidth}{0.4pt}


%% ============================================================
\section{Problem \#1066}
\label{prob:1066}

\noindent\textbf{Status:} OPEN \quad|\quad \textbf{Prize:} none \quad|\quad \textbf{Updated:} 2025-10-01

\noindent\textbf{Tags:} graph theory, planar graphs

\medskip

\noindent\textbf{Statement:}

Let $G$ be a graph given by $n$ points in $\mathbb{R}^2$, where any two distinct points are at least distance $1$ apart, and we draw an edge between two points if they are distance $1$ apart. Let $g(n)$ be maximal such that any such graph always has an independent set on at least $g(n)$ vertices. Estimate $g(n)$, or perhaps $\lim \frac{g(n)}{n}$.




Such graphs are always planar. Erd\H{o}s initially thought that $g(n)=n/3$, but Chung and Graham, and independently Pach, gave a construction that shows $g(n)\leq \frac{6}{19}n$. Pach and Toth \cite{PaTo96} improved this to $g(n)\leq \frac{5}{16}n$. Pollack \cite{Po85} noted that the Four colour theorem implies $g(n)\geq n/4$, since the graph is planar. Pollack reports that Pach observed that this in for unit distance graphs the four colour theorem can be proved by a simple induction. This lower bound has been improved to $\frac{9}{35}n$ by Csizmadia \cite{Cs98} and then $\frac{8}{31}n$ by Swanepoel \cite{Sw02}. The current record bounds are therefore\[\frac{8}{31}n \approx 0.258n \leq g(n) \leq 0.3125n=\frac{5}{16}n.\]Pollack \cite{Po85} also reports a letter from Erd\H{o}s which poses the more general problem of, given $n$ points in $\mathbb{R}^d$ with minimum distance $1$, let $g_d(n)$ be maximal such that there always exist at least $g_d(n)$ many points which have minimum distance $&gt;1$. Is it true that $g_d(n) \gg n/d$ in general? The upper bound $g_d(n) \ll n/d$ is trivial, considering widely spaced unit simplices. See [1070] for the general estimate of independence number of unit distance graphs.




References

[Cs98] Csizmadia, G., On the independence number of minimum distance graphs . Discrete Comput. Geom. (1998), 179--187.

[PaTo96] Pach, J\'{a}nos and T\'oth, G\'{e}za, On the independence number of coin graphs . Geombinatorics (1996), 30--33.

[Po85] Pollack, R., Increasing the minimum distance of a set of points . J. Combin. Theory Ser. A (1985), 450.

[Sw02] Swanepoel, Konrad J., Independence numbers of planar contact graphs . Discrete Comput. Geom. (2002), 649--670.

\noindent\rule{\linewidth}{0.4pt}


%% ============================================================
\section{Problem \#1068}
\label{prob:1068}

\noindent\textbf{Status:} OPEN \quad|\quad \textbf{Prize:} none \quad|\quad \textbf{Updated:} 2025-10-01

\noindent\textbf{Tags:} graph theory, set theory, chromatic number

\medskip

\noindent\textbf{Statement:}

Does every graph with chromatic number $\aleph_1$ contain a countable subgraph which is infinitely vertex-connected?




This is described in \cite{BoPi24} as a 'version of the Erd\H{o}s-Hajnal problem' (which is [1067] ), but it does not seem to appear in \cite{ErHa66}. We say a graph is infinitely (vertex) connected if any two vertices are connected by infinitely many pairwise vertex-disjoint paths. Soukup \cite{So15} constructed a graph with uncountable chromatic number in which every uncountable set is finitely vertex-connected. A simpler construction was given by Bowler and Pitz \cite{BoPi24}. See also [1067] .




References

[BoPi24] N. Bowler and M. Pitz, A note on uncountably chromatic graphs . arXiv:2402.05984 (2024).

[ErHa66] Erd\H{o}s, P. and Hajnal, A., On chromatic number of graphs and set-systems . Acta Math. Acad. Sci. Hungar. (1966), 61-99.

[So15] Soukup, D\'{a}niel T., Trees, ladders and graphs . J. Combin. Theory Ser. B (2015), 96--116.

\noindent\rule{\linewidth}{0.4pt}


%% ============================================================
\section{Problem \#1070}
\label{prob:1070}

\noindent\textbf{Status:} OPEN \quad|\quad \textbf{Prize:} none \quad|\quad \textbf{Updated:} 2025-10-05

\noindent\textbf{Tags:} geometry

\medskip

\noindent\textbf{Statement:}

Let $f(n)$ be maximal such that, given any $n$ points in $\mathbb{R}^2$, there exist $f(n)$ points such that no two are distance $1$ apart. Estimate $f(n)$. In particular, is it true that $f(n)\geq n/4$?




In other words, estimate the minimal independence number of a unit distance graph with $n$ vertices. If $\omega$ is the independence number and $\chi$ is the chromatic number then $\omega \chi\geq n$, and hence $f(n)\geq n/\chi$, where $\chi$ is the answer to the Hadwiger-Nelson problem [508] . The Moser spindle shows $f(n)\leq \frac{2}{7}n\approx 0.285n$. Larman and Rogers \cite{LaRo72} noted that if $m_1$ is the supremum of the upper densities of measurable subsets of $\mathbb{R}^2$ which have no unit distance pairs then\[f(n)\geq m_1n.\]Croft \cite{Cr67} gave the best-known lower bound of $m_1\geq 0.22936$ and hence\[0.22936n \leq f(n) \leq \frac{2}{7}n\approx 0.285n.\]Ambrus, Csisz\'{a}rik, Matolcsi, Varga, and Zs\'{a}mboki \cite{ACMVZ23} have proved that $m_1\leq 0.247$, and hence this approach cannot achieve $f(n)\geq n/4$. See [232] for more on $m_1$. Matolcsi, Ruzsa, Varga, and Zs\'{a}mboki \cite{MRVZ23} have improved the upper bound to\[f(n) \leq \left(\frac{1}{4}+o(1)\right)n.\]They conjecture that $m_1=0.22936\cdots$ (the lower bound of Croft mentioned above) and $f(n)=(1/4+o(1))n$. If we also insist that no two points are distance $&lt;1$ apart then this is problem becomes [1066] .




References

[ACMVZ23] G. Ambrus, A. Csisz\'{a}rik, M. Matolcsi, D. Varga, and P. Zs\'{a}mboki, The density of planar sets avoiding unit distances . arXiv:2207.14179 (2023).

[Cr67] H. T. Croft, Incidence incidents . Eureka (1967), 22-26.

[LaRo72] Larman, D. G. and Rogers, C. A., The realization of distances within sets in Euclidean space . Mathematika (1972), 1-24.

[MRVZ23] M. Matolcsi, I. Z. Ruzsa, D. Varga, and P. Zs\'{a}mboki, The fractional chromatic number of the plane is at least $4$ . arXiv:2311.10069 (2023).

\noindent\rule{\linewidth}{0.4pt}


%% ============================================================
\section{Problem \#1072}
\label{prob:1072}

\noindent\textbf{Status:} OPEN \quad|\quad \textbf{Prize:} none \quad|\quad \textbf{Updated:} 2025-10-05

\noindent\textbf{Tags:} number theory

\medskip

\noindent\textbf{Statement:}

For any prime $p$, let $f(p)$ be the least integer such that $f(p)!+1\equiv 0\pmod{p}$. Is it true that there are infinitely many $p$ for which $f(p)=p-1$? Is it true that $f(p)/p\to 0$ for almost all $p$?




Questions formulated by Erd\H{o}s, Hardy, and Subbarao \cite{HaSu02}, who believed that the number of $p\leq x$ for which $f(p)=p-1$ is $o(x/\log x)$. These are mentioned in problem A2 of Guy's collection.




References

[HaSu02] Hardy, G. E. and Subbarao, M. V., A modified problem of Pillai and some related questions . Amer. Math. Monthly (2002), 554--559.

\noindent\rule{\linewidth}{0.4pt}


%% ============================================================
\section{Problem \#1073}
\label{prob:1073}

\noindent\textbf{Status:} OPEN \quad|\quad \textbf{Prize:} none \quad|\quad \textbf{Updated:} 2025-10-05

\noindent\textbf{Tags:} number theory

\medskip

\noindent\textbf{Statement:}

Let $A(x)$ count the number of composite $u&#60;x$ such that $n!+1\equiv 0\pmod{u}$ for some $n$. Is it true that $A(x)\leq x^{o(1)}$?




A question of Erd\H{o}s raised in discussions with Hardy and Subbarao \cite{HaSu02}. The sequence of such $u$ begins\[25,121,169,437,\ldots\]and is A256519 in the OEIS. Wilson's theorem states that $u$ is prime if and only if $(u-1)!+1\equiv 0\pmod{u}$. This is mentioned in problem A2 of Guy's collection.




References

[HaSu02] Hardy, G. E. and Subbarao, M. V., A modified problem of Pillai and some related questions . Amer. Math. Monthly (2002), 554--559.

\noindent\rule{\linewidth}{0.4pt}


%% ============================================================
\section{Problem \#1074}
\label{prob:1074}

\noindent\textbf{Status:} OPEN \quad|\quad \textbf{Prize:} none \quad|\quad \textbf{Updated:} 2025-10-05

\noindent\textbf{Tags:} number theory

\medskip

\noindent\textbf{Statement:}

Let $S$ be the set of all $m\geq 1$ such that there exists a prime $p\not\equiv 1\pmod{m}$ such that $m!+1\equiv 0\pmod{p}$. Does\[\lim \frac{\lvert S\cap [1,x]\rvert}{x}\]exist? What is it? Similarly, if $P$ is the set of all primes $p$ such that there exists an $m$ with $p\not\equiv 1\pmod{m}$ such that $m!+1\equiv 0\pmod{p}$, then does\[\lim \frac{\lvert P\cap [1,x]\rvert}{\pi(x)}\]exist? What is it?




Questions raised by Erd\H{o}s, Hardy, and Subbarao, who called the set $S$ 'EHS numbers' and the set $P$ 'Pillai primes', and proved that both $S$ and $P$ are infinite. Pillai \cite{Pi30} raised the question of whether there exist any primes in $P$. This was answered by Chowla, who noted that, for example,\[14!+1\equiv 18!+1\equiv 0\pmod{23}.\]The sequence $S$ begins\[8,9,13,14,15,16,17,\ldots\]and is A064164 in the OEIS. The sequence $P$ begins\[23,29,59,61,67,71,\ldots\]and is A063980 in the OEIS. Regarding the first question, Hardy and Subbarao computed all EHS numbers up to $2^{10}$, and write '...if this trend conditions we expect [the limit] to be around $0.5$, if it exists. The frequency with which the EHS numbers occur - most often in long sequences of consecutive integers - makes us believe that their asymptotic density exists and is unity. Erd\H{o}s, though initially hesitant, later agreed with this view.' Regarding the second question, they write '[from the data] it would appear that if the limit exists, it is perhaps between $0.5$ and $0.6$. But then there seems to be no reason why the ratio should not tend to $1$, even though very slowly and certainly not monotonically.' This is discussed in problem A2 of Guy's collection \cite{Gu04}.




References

[Gu04] Guy, Richard K., Unsolved problems in number theory . (2004), xviii+437.

[Pi30] S. S. Pillai, Question 1490 . J. Indian Math. Soc. (1930), 230.

\noindent\rule{\linewidth}{0.4pt}


%% ============================================================
\section{Problem \#1075}
\label{prob:1075}

\noindent\textbf{Status:} OPEN \quad|\quad \textbf{Prize:} none \quad|\quad \textbf{Updated:} 2025-10-05

\noindent\textbf{Tags:} hypergraphs

\medskip

\noindent\textbf{Statement:}

Let $r\geq 3$. There exists $c_r&#62;r^{-r}$ such that, for any $\epsilon&#62;0$, if $n$ is sufficiently large, the following holds. Any $r$-uniform hypergraph on $n$ vertices with at least $(1+\epsilon)(n/r)^r$ many edges contains a subgraph on $m$ vertices with at least $c_rm^r$ edges, where $m=m(n)\to \infty$ as $n\to \infty$.




Erd\H{o}s \cite{Er64f} proved that this is true with $c_r=r^{-r}$ whenever the graph has at least $\epsilon n^r$ many edges.




References

[Er64f] Erd\H{o}s, P., On extremal problems of graphs and generalized graphs . Israel J. Math. (1964), 183--190.

\noindent\rule{\linewidth}{0.4pt}


%% ============================================================
\section{Problem \#1083}
\label{prob:1083}

\noindent\textbf{Status:} OPEN \quad|\quad \textbf{Prize:} none \quad|\quad \textbf{Updated:} 2025-10-17

\noindent\textbf{Tags:} geometry, distances

\medskip

\noindent\textbf{Statement:}

Let $d\geq 3$, and let $f_d(n)$ be the minimal $m$ such that every set of $n$ points in $\mathbb{R}^d$ determines at least $m$ distinct distances. Estimate $f_d(n)$ - in particular, is it true that\[f_d(n)=n^{\frac{2}{d}-o(1)}?\]




A generalisation of the distinct distance problem [89] to higher dimensions. Erd\H{o}s \cite{Er46b} proved\[n^{1/d}\ll_d f_d(n)\ll_d n^{2/d},\]the upper bound construction being given by a set of lattice points. {UL} {LI} Clarkson, Edelsbrunner, Gubias, Sharir, and Welzl \cite{CEGSW90} proved $f_3(n)\gg n^{1/2}$.{/LI} {LI}Aronov, Pach, Sharir, and Tardos \cite{APST04} proved $f_d(n)\gg n^{\frac{1}{d-90/77}-o(1)}$ for any $d\geq 3$ (for example, $f_3(n)\gg n^{0.546}$).{/LI} {LI}Solymosi and Vu \cite{SoVu08} proved $f_3(n) \gg n^{3/5}$ and\[ f_d(n)\gg_d n^{\frac{2}{d}-\frac{c}{d^2}}\]for all $d\geq 4$ for some constant $c&gt;0$. (The result in their paper for $d=3$ is slightly weaker than stated here, but uses as a black box the bound for distinct distances in $2$ dimensions; we have recorded the consequence of combining their method with the work of Guth and Katz on [89] .){/LI} {/UL} The function $f_d(n)$ is essentially the inverse of the function $g_d(n)$ considered in [1089] - with our definitions, $g_d(n)&gt;m$ if and only if $f_d(m)&lt;n$. The emphasis in this problem is, however, on the behaviour as $d$ is fixed and $n\to \infty$.




References

[APST04] Aronov, Boris and Pach, J\'{a}nos and Sharir, Micha and Tardos, G\'{a}bor, Distinct distances in three and higher dimensions . Combin. Probab. Comput. (2004), 283--293.

[CEGSW90] Clarkson, Kenneth L. and Edelsbrunner, Herbert and Guibas, Leonidas J. and Sharir, Micha and Welzl, Emo, Combinatorial complexity bounds for arrangements of curves and spheres . Discrete Comput. Geom. (1990), 99--160.

[Er46b] Erd\H{o}s, P., On sets of distances of {$n$} points . Amer. Math. Monthly (1946), 248--250.

[SoVu08] Solymosi, J\'ozsef and Vu, Van H., Near optimal bounds for the {E}rd\H{o}s distinct distances problem in high dimensions . Combinatorica (2008), 113--125.

\noindent\rule{\linewidth}{0.4pt}


%% ============================================================
\section{Problem \#1084}
\label{prob:1084}

\noindent\textbf{Status:} OPEN \quad|\quad \textbf{Prize:} none \quad|\quad \textbf{Updated:} 2025-10-17

\noindent\textbf{Tags:} geometry, distances

\medskip
\noindent\textit{contact number problem}


\noindent\textbf{Statement:}

Let $f_d(n)$ be minimal such that in any collection of $n$ points in $\mathbb{R}^d$, all of distance at least $1$ apart, there are at most $f_d(n)$ many pairs of points which are distance $1$ apart. Estimate $f_d(n)$.




This is sometimes known as the contact number problem. It is easy to see that $f_1(n)=n-1$ and $f_2(n)&lt;3n$ (since there can be at most $6$ points of distance $1$ from any point). Erd\H{o}s \cite{Er46b} showed\[f_2(n)&lt;3n-cn^{1/2}\]for some constant $c&gt;0$, which the triangular lattice shows is the best possible up to the value of $c$. In \cite{Er75f} he speculated that the triangular lattice is exactly the best possible, and in particular\[f_2(3n^2+3n+1)=9n^2+3n.\]Harborth \cite{Ha74b} proved this, and more generally\[f_2(n)=\lfloor 3n-\sqrt{12n-3}\rfloor\]for all $n\geq 2$. In \cite{Er75f} he claims the existence of $c_1,c_2&gt;0$ such that\[6n-c_1n^{2/3}&lt; f_3(n) &lt; 6n-c_2n^{2/3}.\]An upper bound of\[f_3(n) &lt; 6n-0.926n^{2/3}\]for all $n\geq 2$ was proved by Bezdek and Reid \cite{BeRe13}. In general, it is known that\[(d-o(1))n \leq f_d(n) \leq 2^{O(d)}n,\]the lower bound coming from points arranged in an integer grid and the upper bound from the fact that $2^{O(d)}$ many non-intersecting congruent balls can touch a fixed ball (the kissing number problem). A recent survey on contact numbers for sphere packings is by Bezdek and Khan \cite{BeKh18}. See [223] for the analogous problem with maximal distance $1$.




References

[BeKh18] Bezdek, K\'{a}roly and Khan, Muhammad A., Contact numbers for sphere packings . (2018), 25--47.

[BeRe13] Bezdek, K\'{a}roly and Reid, Samuel, Contact graphs of unit sphere packings revisited . J. Geom. (2013), 57--83.

[Er46b] Erd\H{o}s, P., On sets of distances of {$n$} points . Amer. Math. Monthly (1946), 248--250.

[Er75f] Erd\H{o}s, Paul, On some problems of elementary and combinatorial geometry . Ann. Mat. Pura Appl. (4) (1975), 99-108.

[Ha74b] Harborth, Heiko, L\"{o}sung zu Problem 664A . Elem. Math. (1974), 14-15.

\noindent\rule{\linewidth}{0.4pt}


%% ============================================================
\section{Problem \#1085}
\label{prob:1085}

\noindent\textbf{Status:} OPEN \quad|\quad \textbf{Prize:} none \quad|\quad \textbf{Updated:} 2025-10-17

\noindent\textbf{Tags:} geometry, distances

\medskip

\noindent\textbf{Statement:}

Let $f_d(n)$ be minimal such that, in any set of $n$ points in $\mathbb{R}^d$, there exist at most $f_d(n)$ pairs of points which distance $1$ apart. Estimate $f_d(n)$.




The most difficult cases are $d=2$ and $d=3$. When $d=2$ this is the unit distance problem [90] , and the best known bounds are\[n^{1+c}&lt; f_2(n) \ll n^{4/3}\]for some small constant $c&gt;0$, the lower bound by an internal model at OpenAI, at least for infinitely many $n$ (see [90] ) and the upper bound by Spencer, Szemer\'{e}di, and Trotter \cite{SST84}. When $d=3$ the best known bounds are\[n^{4/3}\log\log n \ll f_3(n) \ll n^{3/2}\beta(n)\]where $\beta(n)$ is a very slowly growing function, the lower bound by Erd\H{o}s \cite{Er60b} and the upper bound by Clarkson, Edelsbrunner, Guibas, Sharir, and Welzl \cite{CEGSW90}. A construction of Lenz (taking points on orthogonal circles) shows that, for $d\geq 4$,\[f_d(n)\geq \frac{p-1}{2p}n^2-O(1)\]with $p=\lfloor d/2\rfloor$. Erd\H{o}s \cite{Er60b} showed that the Erd\H{o}s-Stone theorem implies\[f_d(n) \leq \left(\frac{p-1}{2p}+o(1)\right)n^2\]for $d\geq 4$. Erd\H{o}s \cite{Er67e} determined $f_d(n)$ up to $O(1)$ for all even $d\geq 4$. Brass \cite{Br97} determined $f_4(n)$ exactly. Swanepoel \cite{Sw09} determined $f_d(n)$ exactly for even $d\geq 6$. For odd $d\geq 5$ Erd\H{o}s and Pach \cite{ErPa90} proved that there exist constants $c_1(d),c_2(d)&gt;0$ such that\[\frac{p-1}{2p}n^2 +c_1n^{4/3}\leq f_d(n) \leq \frac{p-1}{2p}n^2 +c_2n^{4/3}.\]




References

[Br97] Brass, P., On the maximum number of unit distances among {$n$} points in dimension four . (1997), 277--290.

[CEGSW90] Clarkson, Kenneth L. and Edelsbrunner, Herbert and Guibas, Leonidas J. and Sharir, Micha and Welzl, Emo, Combinatorial complexity bounds for arrangements of curves and spheres . Discrete Comput. Geom. (1990), 99--160.

[Er60b] Erd\H{o}s, P., On sets of distances of {$n$} points in {E}uclidean space . Magyar Tud. Akad. Mat. Kutat\'o{} Int. K\"ozl. (1960), 165--169.

[Er67e] Erd\H{o}s, P., On some applications of graph theory to geometry . Canadian J. Math. (1967), 968--971.

[ErPa90] Erd\H{o}s, P. and Pach, J., Variations on the theme of repeated distances . Combinatorica (1990), 261--269.

[SST84] Spencer, J. and Szemer\'{e}di, E. and Trotter, Jr., W., Unit distances in the Euclidean plane . Graph theory and combinatorics (Cambridge, 1983) (1984), 293-303.

[Sw09] Swanepoel, Konrad J., Unit distances and diameters in {E}uclidean spaces . Discrete Comput. Geom. (2009), 1--27.

\noindent\rule{\linewidth}{0.4pt}


%% ============================================================
\section{Problem \#1086}
\label{prob:1086}

\noindent\textbf{Status:} OPEN \quad|\quad \textbf{Prize:} none \quad|\quad \textbf{Updated:} 2025-10-17

\noindent\textbf{Tags:} geometry, distances

\medskip

\noindent\textbf{Statement:}

Let $g(n)$ be minimal such that any set of $n$ points in $\mathbb{R}^2$ contains the vertices of at most $g(n)$ many triangles with the same area. Estimate $g(n)$.




Equivalently, how many triangles of area $1$ can a set of $n$ points in $\mathbb{R}^2$ determine? Erd\H{o}s and Purdy attribute this question to Oppenheim. Erd\H{o}s and Purdy \cite{ErPu71} proved\[n^2\log\log n \ll g(n) \ll n^{5/2},\]and believed the lower bound to be closer to the truth. The upper bound has been improved a number of times - by Pach and Sharir \cite{PaSh92}, Dumitrescu, Sharir, and T\'{o}th \cite{DST09}, Apfelbaum and Sharir \cite{ApSh10}, and Apfaulbaum \cite{Ap13}. The best known bound is\[g(n) \ll n^{20/9}\]by Raz and Sharir \cite{RaSh17}. Erd\H{o}s and Purdy also ask a similar question about the higher-dimensional generalisations - more generally, let $g_d^{r}(n)$ be minimal such that any set of $n$ points in $\mathbb{R}^d$ contains the vertices of at most $g_d^{r}(n)$ many $r$-dimensional simplices with the same volume. Erd\H{o}s and Purdy \cite{ErPu71} proved $g_3^2(n) \ll n^{8/3}$, and Dumitrescu, Sharir, and T\'{o}th \cite{DST09} improved this to $g_3^2(n) \ll n^{2.4286}$. Erd\H{o}s and Purdy \cite{ErPu71} proved $g_6^2(n)\gg n^3$. Purdy \cite{Pu74} proved\[g_4^2(n)\leq g^2_5(n) \ll n^{3-c}\]for some constant $c&gt;0$. An observation of Oppenheim (using a construction of Lenz) detailed in \cite{ErPu71} shows that\[g_{2k+2}^k(n)\geq \left(\frac{1}{(k+1)^{k+1}}+o(1)\right)n^{k+1}\]and Erd\H{o}s and Purdy conjecture this is the best possible. See also [90] and [755] .




References

[Ap13] R. Apfelbaum, Geometric Incidences and Repeated Configurations . Ph.D. Dissertation, School of Computer Science, Tel Aviv University (2013).

[ApSh10] Apfelbaum, Roel and Sharir, Micha, An improved bound on the number of unit area triangles . Discrete Comput. Geom. (2010), 753--761.

[DST09] Dumitrescu, Adrian and Sharir, Micha and T\'oth, Csaba D., Extremal problems on triangle areas in two and three dimensions . J. Combin. Theory Ser. A (2009), 1177--1198.

[ErPu71] Erd\H{o}s, Paul and Purdy, George, Some extremal problems in geometry . J. Combinatorial Theory Ser. A (1971), 246--252.

[PaSh92] Pach, J\'{a}nos and Sharir, Micha, Repeated angles in the plane and related problems . J. Combin. Theory Ser. A (1992), 12--22.

[Pu74] Purdy, George, Some extremal problems in geometry . Discrete Math. (1974), 305--315.

[RaSh17] Raz, Orit E. and Sharir, Micha, The number of unit-area triangles in the plane: theme and variation . Combinatorica (2017), 1221--1240.

\noindent\rule{\linewidth}{0.4pt}


%% ============================================================
\section{Problem \#1087}
\label{prob:1087}

\noindent\textbf{Status:} OPEN \quad|\quad \textbf{Prize:} none \quad|\quad \textbf{Updated:} 2025-10-17

\noindent\textbf{Tags:} geometry, distances

\medskip

\noindent\textbf{Statement:}

Let $f(n)$ be minimal such that every set of $n$ points in $\mathbb{R}^2$ contains at most $f(n)$ many sets of four points which are 'degenerate' in the sense that some pair are the same distance apart. Estimate $f(n)$ - in particular, is it true that $f(n)\leq n^{3+o(1)}$?




A question of Erd\H{o}s and Purdy \cite{ErPu71}, who proved\[n^3\log n \ll f(n) \ll n^{7/2}.\]




References

[ErPu71] Erd\H{o}s, Paul and Purdy, George, Some extremal problems in geometry . J. Combinatorial Theory Ser. A (1971), 246--252.

\noindent\rule{\linewidth}{0.4pt}


%% ============================================================
\section{Problem \#1088}
\label{prob:1088}

\noindent\textbf{Status:} OPEN \quad|\quad \textbf{Prize:} none \quad|\quad \textbf{Updated:} 2025-10-17

\noindent\textbf{Tags:} geometry

\medskip

\noindent\textbf{Statement:}

Let $f_d(n)$ be the minimal $m$ such that any set of $m$ points in $\mathbb{R}^d$ contains a set of $n$ points such that any two determined distances are distinct. Estimate $f_d(n)$. In particular, is it true that, for fixed $n\geq 3$,\[f_d(n)=2^{o(d)}?\]




It is easy to prove that $f_d(n) \leq n^{O_d(1)}$. Erd\H{o}s \cite{Er75f} claimed that he and Straus proved $f_d(n)\leq c_n^d$ for some constant $c_n&gt;0$. When $d=1$ this is the subject of [530] , and $f_1(n)\asymp n^2$. When $n=3$ this is the subject of [503] . Erd\H{o}s could prove $f_2(3)=7$ and Croft \cite{Cr62} proved $f_3(3)=9$. The results described at [503] demonstrate that $f_d(3)=d^2/2+O(d)$. The behaviour of $f_d(n)$ for fixed $d$ as $n\to \infty$ is the subject of [1208] .




References

[Cr62] Croft, H. T., $9$-point and $7$-point configurations in $3$-space . Proc. London Math. Soc. (3) (1962), 400-424.

[Er75f] Erd\H{o}s, Paul, On some problems of elementary and combinatorial geometry . Ann. Mat. Pura Appl. (4) (1975), 99-108.

\noindent\rule{\linewidth}{0.4pt}


%% ============================================================
\section{Problem \#1093}
\label{prob:1093}

\noindent\textbf{Status:} OPEN \quad|\quad \textbf{Prize:} none \quad|\quad \textbf{Updated:} 2025-10-18

\noindent\textbf{Tags:} number theory, binomial coefficients

\medskip

\noindent\textbf{Statement:}

For $n\geq 2k$ we define the deficiency of $\binom{n}{k}$ as follows. If $\binom{n}{k}$ is divisible by a prime $p\leq k$ then the deficiency is undefined. Otherwise, the deficiency is the number of $0\leq i&lt;k$ such that $n-i$ is $k$-smooth, that is, divisible only by primes $\leq k$. Are there infinitely many binomial coefficients with deficiency $1$? Are there only finitely many with deficiency $&gt;1$?




A problem of Erd\H{o}s, Lacampagne, and Selfridge \cite{ELS88}, that was also asked in the 1986 problem session of West Coast Number Theory (as reported here ). In \cite{ELS93} they prove that if the deficiency exists and is $\geq 1$ then $n\ll 2^k\sqrt{k}$. The following examples are either from \cite{ELS88} or here . The following have deficiency $1$ (there are $58$ examples with $n\leq 10^5$):\[\binom{7}{3},\binom{13}{4},\binom{14}{4},\binom{23}{5},\binom{62}{6},\binom{94}{10},\binom{95}{10}.\]The examples which follow are the only known examples with deficiency $&gt;1$. The following have deficiency $2$:\[\binom{44}{8},\binom{74}{10},\binom{174}{12},\binom{239}{14},\binom{5179}{27},\binom{8413}{28},\binom{8414}{28},\binom{96622}{42}.\]The following have deficiency $3$:\[\binom{46}{10},\binom{47}{10},\binom{241}{16},\binom{2105}{25},\binom{1119}{27},\binom{6459}{33}.\]The following has deficiency $4$:\[\binom{47}{11}.\]The following has deficiency $9$:\[\binom{284}{28}.\]See also [384] and [1094] . Barreto in the comments has given a positive answer to the second question, conditional on two (strong) conjectures.




References

[ELS88] Erd\H{o}s, P. and Lacampagne, C. B. and Selfridge, J. L., Prime factors of binomial coefficients and related problems . Acta Arith. (1988), 507--523.

[ELS93] Erd\H{o}s, P. and Lacampagne, C. B. and Selfridge, J. L., Estimates of the least prime factor of a binomial coefficient . Math. Comp. (1993), 215--224.

\noindent\rule{\linewidth}{0.4pt}


%% ============================================================
\section{Problem \#1094}
\label{prob:1094}

\noindent\textbf{Status:} OPEN \quad|\quad \textbf{Prize:} none \quad|\quad \textbf{Updated:} 2025-10-18

\noindent\textbf{Tags:} number theory, binomial coefficients

\medskip

\noindent\textbf{Statement:}

For all $n\geq 2k$ the least prime factor of $\binom{n}{k}$ is $\leq \max(n/k,k)$, with only finitely many exceptions.




A stronger form of [384] that appears in a paper of Erd\H{o}s, Lacampagne, and Selfridge \cite{ELS88}. Erd\H{o}s observed that the least prime factor is always $\leq n/k$ provided $n$ is sufficiently large depending on $k$. Selfridge \cite{Se77} further conjectured that this always happens if $n\geq k^2-1$, except $\binom{62}{6}$. The threshold $g(k)$ below which $\binom{n}{k}$ is guaranteed to be divisible by a prime $\leq k$ is the subject of [1095] . More precisely, in \cite{ELS88} they conjecture that if $n\geq 2k$ then the least prime factor of $\binom{n}{k}$ is $\leq \max(n/k,k)$ with the following $14$ exceptions:\[\binom{7}{3},\binom{13}{4},\binom{23}{5},\binom{14}{4},\binom{44}{8},\binom{46}{10},\binom{47}{10},\]\[\binom{47}{11},\binom{62}{6},\binom{74}{10},\binom{94}{10},\binom{95}{10},\binom{241}{16},\binom{284}{28}.\]They also suggest the stronger conjecture that, with a finite number of exceptions, the least prime factor is $\leq \max(n/k,\sqrt{k})$, or perhaps even $\leq \max(n/k,O(\log k))$. Indeed, in \cite{ELS93} they provide some further computational evidence, and point out it is consistent with what they know that in fact this holds with $\leq \max(n/k,13)$, with only $12$ exceptions. Discussed in problem B31 and B33 of Guy's collection \cite{Gu04} - there Guy credits Selfridge with the conjecture that if $n&gt; 17.125k$ then $\binom{n}{k}$ has a prime factor $p\leq n/k$. This is related to [1093] , in that the only counterexamples to this conjecture can occur from $\binom{n}{k}$ with deficiency $\geq 1$. There is an interesting discussion about this problem on MathOverflow .




References

[ELS88] Erd\H{o}s, P. and Lacampagne, C. B. and Selfridge, J. L., Prime factors of binomial coefficients and related problems . Acta Arith. (1988), 507--523.

[ELS93] Erd\H{o}s, P. and Lacampagne, C. B. and Selfridge, J. L., Estimates of the least prime factor of a binomial coefficient . Math. Comp. (1993), 215--224.

[Gu04] Guy, Richard K., Unsolved problems in number theory . (2004), xviii+437.

[Se77] J. L. Selfridge, Some problems on the prime factors of consecutive integers . Notices Amer. Math. Soc. (1977), A456-457.

\noindent\rule{\linewidth}{0.4pt}


%% ============================================================
\section{Problem \#1095}
\label{prob:1095}

\noindent\textbf{Status:} OPEN \quad|\quad \textbf{Prize:} none \quad|\quad \textbf{Updated:} 2025-10-18

\noindent\textbf{Tags:} number theory, binomial coefficients

\medskip

\noindent\textbf{Statement:}

Let $g(k)&#62;k+1$ be the smallest $n$ such that all prime factors of $\binom{n}{k}$ are $&#62;k$. Estimate $g(k)$.




A question of Ecklund, Erd\H{o}s, and Selfridge \cite{EES74}, who proved\[k^{1+c}&lt;g(k)\leq \exp((1+o(1))k)\]for some constant $c&gt;0$, and conjectured $g(k)&lt;L_k=[1,\ldots,k]$, the least common multiple of all integers $\leq k$, for all large $k$. In \cite{EES74} they further conjecture that\[\limsup \frac{g(k+1)}{g(k)}=\infty\]and\[\liminf \frac{g(k+1)}{g(k)}=0.\]The lower bound was improved by Erd\H{o}s, Lacampagne, and Selfridge \cite{ELS93} and Granville and Ramar\'{e} \cite{GrRa96}. The current record is\[g(k) \gg \exp(c(\log k)^2)\]for some $c&gt;0$, due to Konyagin \cite{Ko99b}. Erd\H{o}s, Lacampagne, and Selfridge \cite{ELS93} write 'it is clear to every right-thinking person' that $g(k)\geq\exp(c\frac{k}{\log k})$ for some constant $c&gt;0$. Sorenson, Sorenson, and Webster \cite{SSW20} give heuristic evidence that\[\log g(k) \asymp \frac{k}{\log k}.\]See also [1094] .




References

[EES74] Ecklund, Jr., E. F. and Erd\H{o}s, P. and Selfridge, J. L., A new function associated with the prime factors of {$(\sp{n}\sb{k})$} . Math. Comp. (1974), 647--649.

[ELS93] Erd\H{o}s, P. and Lacampagne, C. B. and Selfridge, J. L., Estimates of the least prime factor of a binomial coefficient . Math. Comp. (1993), 215--224.

[GrRa96] Granville, Andrew and Ramar\'{e}, Olivier, Explicit bounds on exponential sums and the scarcity of squarefree binomial coefficients . Mathematika (1996), 73--107.

[Ko99b] Konyagin, S. V., Estimates of the least prime factor of a binomial coefficient . Mathematika (1999), 41--55.

[SSW20] Sorenson, Brianna and Sorenson, Jonathan and Webster, Jonathan, An algorithm and estimates for the {E}rd\H{o}s-Selfridge function . (2020), 371--385.

\noindent\rule{\linewidth}{0.4pt}


%% ============================================================
\section{Problem \#1097}
\label{prob:1097}

\noindent\textbf{Status:} OPEN \quad|\quad \textbf{Prize:} none \quad|\quad \textbf{Updated:} 2025-10-18

\noindent\textbf{Tags:} number theory, additive combinatorics

\medskip

\noindent\textbf{Statement:}

Let $A$ be a set of $n$ integers. How many distinct $d$ can occur as the common difference of a three-term arithmetic progression in $A$? In particular, are there always $O(n^{3/2})$ many such $d$?




A problem Erd\H{o}s posed in the 1989 problem session of Great Western Number Theory. He states that Erd\H{o}s and Ruzsa gave an explicit construction which achieved $n^{1+c}$ for some $c&gt;0$, and Erd\H{o}s and Spencer gave a probabilistic proof which achieved $n^{3/2}$, and speculated this may be the best possible. In the comment section, Chan has noticed that this problem is exactly equivalent to a sums-differences question of Bourgain \cite{Bo99}, introduced as an arithmetic path towards the Kakeya conjecture: find the smallest $c\in [1,2]$ such that, for any finite sets of integers $A$ and $B$ and $G\subseteq A\times B$ we have\[\lvert A\overset{G}{-}B\rvert \ll \max(\lvert A\rvert,\lvert B\rvert, \lvert A\overset{G}{+}B\rvert)^c\](where, for example, $A\overset{G}{+}B$ denotes the set of $a+b$ with $(a,b)\in G$). This is equivalent in the sense that the greatest exponent $c$ achievable for the main problem here is equal to the smallest constant achievable for the sums-differences question. The current best bounds known are thus\[1.77898\cdots \leq c \leq 11/6 \approx 1.833.\]The upper bound is due to Katz and Tao \cite{KaTa99}. The lower bound is due to Lemm \cite{Le15} (with a very small improvement found by AlphaEvolve \cite{GGTW25}). This resolves the second part of the question negatively, but the general question of the correct order of magnitude remains open.




References

[Bo99] Bourgain, J., On the dimension of {K}akeya sets and related maximal inequalities . Geom. Funct. Anal. (1999), 256--282.

[GGTW25] B. Georgiev, J. G\'{o}mez-Serrano, T. Tao, and A. Wagner, Mathematical exploration and discovery at scale . arXiv:2511.02864 (2025).

[KaTa99] Katz, Nets Hawk and Tao, Terence, Bounds on arithmetic projections, and applications to the {K}akeya conjecture . Math. Res. Lett. (1999), 625--630.

[Le15] Lemm, Marius, New counterexamples for sums-differences . Proc. Amer. Math. Soc. (2015), 3863--3868.

\noindent\rule{\linewidth}{0.4pt}


%% ============================================================
\section{Problem \#1100}
\label{prob:1100}

\noindent\textbf{Status:} OPEN \quad|\quad \textbf{Prize:} none \quad|\quad \textbf{Updated:} 2025-10-19

\noindent\textbf{Tags:} number theory, divisors

\medskip

\noindent\textbf{Statement:}

If $1=d_1&#60;\cdots&#60;d_{\tau(n)}=n$ are the divisors of $n$, then let $\tau_\perp(n)$ count the number of $i$ for which $(d_i,d_{i+1})=1$. Is it true that $\tau_\perp(n)/\omega(n)\to \infty$ for almost all $n$? Is it true that\[\tau_\perp(n)&#60; \exp((\log n)^{o(1)})\]for all $n$? Let\[g(k) = \max_{\omega(n)=k}\tau_\perp(n),\]where $\omega(n)$ counts the number of distinct prime divisors of $n$, and $n$ is restricted to squarefree integers. Determine the growth of $g(k)$.




The function $\tau_\perp(n)$ was considered by Erd\H{o}s and Hall \cite{ErHa78}. It is trivial that $\tau_\perp(n)\geq \omega(n)$ (with equality for infinitely many $n$). Erd\H{o}s and Hall prove, for all $\epsilon&#62;0$ and sufficiently large $x$,\[\max_{n&lt;x} \tau_\perp(n) &gt; \exp((\log\log x)^{2-\epsilon}).\]Erd\H{o}s and Simonovits (see \cite{Er81h}) proved\[(2^{1/2}+o(1))^k &lt; g(k) &lt; (2-c)^k\]for some constant $c&gt;0$.




References

[Er81h] Erd\H{o}s, P., Some problems and results on additive and multiplicative number theory . Analytic number theory (Philadelphia, Pa., 1980) (1981), 171-182.

[ErHa78] Erd\H{o}s, P. and Hall, R. R., On some unconventional problems on the divisors of integers . J. Austral. Math. Soc. Ser. A (1978), 479--485.

\noindent\rule{\linewidth}{0.4pt}


%% ============================================================
\section{Problem \#1101}
\label{prob:1101}

\noindent\textbf{Status:} OPEN \quad|\quad \textbf{Prize:} none \quad|\quad \textbf{Updated:} 2025-10-19

\noindent\textbf{Tags:} number theory

\medskip

\noindent\textbf{Statement:}

If $u=\{u_1&lt;u_2&lt;\cdots\}$ is a sequence of integers such that $(u_i,u_j)=1$ for all $i\neq j$ and $\sum \frac{1}{u_i}&lt;\infty$ then let $\{a_1&lt;a_2&lt;\cdots\}$ be the sequence of integers which are not divisible by any of the $u_i$. For any $x$ define $t_x$ by\[u_1\cdots u_{t_x}\leq x&lt; u_1\cdots u_{t_x}u_{t_x+1}.\]We call such a sequence $u_i$ good if, for all $\epsilon&gt;0$, if $x$ is sufficiently large then\[\max_{a_k&#60;x} (a_{k+1}-a_k) &#60; (1+\epsilon)t_x \prod_{i}\left(1-\frac{1}{u_i}\right)^{-1}.\]Is there a good sequence such that $u_n&#60; n^{O(1)}$? Is there a good sequence such that $u_n\leq e^{o(n)}$?




Erd\H{o}s \cite{Er81h} believed the answer to the first question is no and the second question is yes. He proved the existence of some good sequence (in which all the $u_i$ are primes). An easy sieve argument proves that we always have, for any sequence $u$ with those properties,\[\max_{a_k&lt;x} (a_{k+1}-a_k)&gt; (1+o(1))t_x \prod_{i}\left(1-\frac{1}{u_i}\right)^{-1}.\]The strong form of [208] is asking whether if $u_i=p_i^2$, the sequence of prime squares, is good.




References

[Er81h] Erd\H{o}s, P., Some problems and results on additive and multiplicative number theory . Analytic number theory (Philadelphia, Pa., 1980) (1981), 171-182.

\noindent\rule{\linewidth}{0.4pt}


%% ============================================================
\section{Problem \#1103}
\label{prob:1103}

\noindent\textbf{Status:} OPEN \quad|\quad \textbf{Prize:} none \quad|\quad \textbf{Updated:} 2025-10-19

\noindent\textbf{Tags:} number theory

\medskip

\noindent\textbf{Statement:}

Let $A$ be an infinite sequence of integers such that every $n\in A+A$ is squarefree. How fast must $A$ grow?




Erd\H{o}s notes there exists such a sequence which grows exponentially, but does not expect such a sequence of polynomial growth. In \cite{Er81h} he asked whether there is an infinite sequence of integers $A$ such that, for every $a\in A$ and prime $p$, if\[a\equiv t\pmod{p^2}\]then $1\leq t&lt;p^2/2$. He noted that such a sequence has the property that every $n\in A+A$ is squarefree. He wrote 'I am doubtful if such a sequence exists. I formulated this problem while writing these lines and must ask the indulgence of the reader if it turns out to be trivial.' Indeed, there are trivially at most finitely many such $a$, since there cannot be any primes $p\in (a^{1/2},(2a)^{1/2}]$, but there exist primes in $(x,\sqrt{2}x)$ for all large $x$. If $A=\{a_1&lt;a_2&lt;\cdots\}$ is such a sequence then van Doorn and Tao \cite{vDTa25} have shown that $a_j &gt; 0.24j^{4/3}$ for all $j$, and further that there exists such a sequence (furthermore with squarefree terms) such that\[a_j &lt; \exp(5j/\log j)\]for all large $j$. A superior lower bound of $a_j \gg j^{15/11-o(1)}$ had earlier been found by Konyagin \cite{Ko04} when considering the finite case [1109] . They also obtain further results for the generalisation from squarefree to $k$-free integers, and also replacing $A+A$ with $A\cup (A+A)\cup(A+A+A)$. See also [1109] for the finite analogue of this problem.




References

[Er81h] Erd\H{o}s, P., Some problems and results on additive and multiplicative number theory . Analytic number theory (Philadelphia, Pa., 1980) (1981), 171-182.

[Ko04] Konyagin, S. V., Problems of the set of square-free numbers . Izv. Ross. Akad. Nauk Ser. Mat. (2004), 63--90.

[vDTa25] W. van Doorn and T. Tao, Growth rates of sequences governed by the squarefree properties of its translates . arXiv:2512.01087 (2025).

\noindent\rule{\linewidth}{0.4pt}


%% ============================================================
\section{Problem \#1104}
\label{prob:1104}

\noindent\textbf{Status:} OPEN \quad|\quad \textbf{Prize:} none \quad|\quad \textbf{Updated:} 2025-10-26

\noindent\textbf{Tags:} graph theory, chromatic number

\medskip

\noindent\textbf{Statement:}

Let $f(n)$ be the maximum possible chromatic number of a triangle-free graph on $n$ vertices. Estimate $f(n)$.




The best bounds available are\[(1-o(1))(n/\log n)^{1/2}\leq f(n) \leq (2+o(1))(n/\log n)^{1/2}.\]The upper bound is due to Davies and Illingworth \cite{DaIl22}, the lower bound follows from a construction of Hefty, Horn, King, and Pfender \cite{HHKP25}. One can ask a similar question for the maximum possible chromatic number of a triangle-free graph on $m$ edges. Let this be $g(m)$. Davies and Illingworth \cite{DaIl22} prove\[g(m) \leq (3^{5/3}+o(1))\left(\frac{m}{(\log m)^2}\right)^{1/3}.\]Kim \cite{Ki95} gave a construction which implies $g(m) \gg (m/(\log m)^2)^{1/3}$. The function $f(n)$ is the inverse to the function $h_3(k)$ considered in [1013] . A generalisation of $f(n)$ is considered in [920] .




References

[DaIl22] Davies, Ewan and Illingworth, Freddie, The {$\chi$}-{R}amsey problem for triangle-free graphs . SIAM J. Discrete Math. (2022), 1124--1134.

[HHKP25] Z. Hefty, P. Horn, D. King, and F. Pfender, Improving $R(3,k)$ in just two bites . arXiv:2510.19718 (2025).

[Ki95] Kim, J. H., The Ramsey number $R(3,t)$ has order of magnitude $t^2/\log t$ . Random Structures and Algorithms (1995), 173-207.

\noindent\rule{\linewidth}{0.4pt}


%% ============================================================
\section{Problem \#1106}
\label{prob:1106}

\noindent\textbf{Status:} OPEN \quad|\quad \textbf{Prize:} none \quad|\quad \textbf{Updated:} 2025-11-17

\noindent\textbf{Tags:} number theory

\medskip

\noindent\textbf{Statement:}

Let $p(n)$ denote the partition function of $n$ and let $F(n)$ count the number of distinct prime factors of\[\prod_{1\leq k\leq n}p(k).\]Does $F(n)\to \infty$ with $n$? Is $F(n)&gt;n$ for all sufficiently large $n$?




Asked by Erd\H{o}s at Oberwolfach in 1986. Schinzel noted in the Oberwolfach problem book that $F(n)\to \infty$ follows from the asymptotic formula for $p(n)$ and a result of Tijdeman \cite{Ti73}. This is not obvious; details are given in a paper of Erd\H{o}s and Ivi\'{c} (see page 69 of \cite{ErIv90}). Schinzel and Wirsing \cite{ScWi87} have proved $F(n) \gg \log n$. Ono \cite{On00} has proved that every prime divides $p(n)$ for some $n\geq 1$ (indeed this holds, for any fixed prime, for a positive density set of $n$).




References

[ErIv90] Erd\H{o}s, Paul and Ivi\'c, Aleksandar, The distribution of values of a certain class of arithmetic functions at consecutive integers . (1990), 45--91.

[On00] Ono, Ken, Distribution of the partition function modulo {$m$} . Ann. of Math. (2) (2000), 293--307.

[ScWi87] Schinzel, A. and Wirsing, E., Multiplicative properties of the partition function . Proc. Indian Acad. Sci. Math. Sci. (1987), 297--303.

[Ti73] Tijdeman, R., On integers with many small prime factors . Compositio Math. (1973), 319--330.

\noindent\rule{\linewidth}{0.4pt}


%% ============================================================
\section{Problem \#1107}
\label{prob:1107}

\noindent\textbf{Status:} OPEN \quad|\quad \textbf{Prize:} none \quad|\quad \textbf{Updated:} 2025-11-17

\noindent\textbf{Tags:} number theory, powerful

\medskip

\noindent\textbf{Statement:}

Let $r\geq 2$. A number $n$ is $r$-powerful if for every prime $p$ which divides $n$ we have $p^r\mid n$. Is every large integer the sum of at most $r+1$ many $r$-powerful numbers?




Given in the 1986 Oberwolfach problem book as a problem of Erd\H{o}s and Ivi\'{c}. This is true when $r=2$, as proved by Heath-Brown \cite{He88} (see [941] ). See [940] for the problem of which integers are the sum of at most $r$ many $r$-powerful numbers.




References

[He88] Heath-Brown, D. R., Ternary quadratic forms and sums of three square-full numbers . (1988), 137--163.

\noindent\rule{\linewidth}{0.4pt}


%% ============================================================
\section{Problem \#1108}
\label{prob:1108}

\noindent\textbf{Status:} OPEN \quad|\quad \textbf{Prize:} none \quad|\quad \textbf{Updated:} 2025-11-17

\noindent\textbf{Tags:} number theory, factorials

\medskip

\noindent\textbf{Statement:}

Let\[A = \left\{ \sum_{n\in S}n! : S\subset \mathbb{N}\textrm{ finite}\right\}.\]If $k\geq 2$, then does $A$ contain only finitely many $k$th powers? Does it contain only finitely many powerful numbers?




Asked by Erd\H{o}s at Oberwolfach in 1988. It is open even whether there are infinitely many squares of the form $1+n!$ (see [398] ). This was motivated in part by a problem of Mahler which he discussed with Erd\H{o}s a few days before his death in 1988: if $k\geq 5$ and\[A_k= \left\{ \sum_{n\in S}k^n : S\subset \mathbb{N}\textrm{ finite}\right\}\]then does $A_k$ contain only finitely many squares? Mahler showed that there are infinitely many squares in $A_k$ for $k\leq 4$, and found only one square for $k\geq 5$, namely\[1+7+7^2+7^3=400.\]Brindza and Erd\H{o}s \cite{BrEr91} proved that, for any $r$, if $n_1!+\cdots+n_r!$ is powerful then $n_1\ll_r 1$.




References

[BrEr91] Brindza, B. and Erd\H{o}s, P., On some {D}iophantine problems involving powers and factorials . J. Austral. Math. Soc. Ser. A (1991), 1--7.

\noindent\rule{\linewidth}{0.4pt}


%% ============================================================
\section{Problem \#1109}
\label{prob:1109}

\noindent\textbf{Status:} OPEN \quad|\quad \textbf{Prize:} none \quad|\quad \textbf{Updated:} 2025-12-03

\noindent\textbf{Tags:} number theory

\medskip

\noindent\textbf{Statement:}

Let $f(N)$ be the size of the largest subset $A\subseteq \{1,\ldots,N\}$ such that every $n\in A+A$ is squarefree. Estimate $f(N)$. In particular, is it true that $f(N)\leq N^{o(1)}$, or even $f(N) \leq (\log N)^{O(1)}$?




First studied by Erd\H{o}s and S\'{a}rk\"{o}zy \cite{ErSa87}, who proved\[\log N \ll f(N) \ll N^{3/4}\log N,\]and guessed the lower bound is nearer the truth. S\'{a}rk\"{o}zy \cite{Sa92c} extended this to consider the case of $A+B$ and also looking for sumsets which are $k$-power-free. Gyarmati \cite{Gy01} gave an alternative proof of $f(N)\gg \log N$, and also gave new bounds for the case of $A+B$. Konyagin \cite{Ko04} improved this to\[ \log\log N(\log N)^2\ll f(N) \ll N^{11/15+o(1)}.\]The infinite analogue of this problem is [1103] . (In particular upper bounds for this $f(N)$ directly imply lower bounds for the size of the $a_j$ considered there.)




References

[ErSa87] Erd\H{o}s, P. and S\'{a}rk\"ozy, A., On divisibility properties of integers of the form {$a+a'$} . Acta Math. Hungar. (1987), 117--122.

[Gy01] Gyarmati, Katalin, On divisibility properties of integers of the form {$ab+1$} . Period. Math. Hungar. (2001), 71--79.

[Ko04] Konyagin, S. V., Problems of the set of square-free numbers . Izv. Ross. Akad. Nauk Ser. Mat. (2004), 63--90.

[Sa92c] S\'{a}rk\"ozy, G. N., On a problem of {P}. {E}rd\H{o}s . Acta Math. Hungar. (1992), 271--282.

\noindent\rule{\linewidth}{0.4pt}


%% ============================================================
\section{Problem \#1110}
\label{prob:1110}

\noindent\textbf{Status:} OPEN \quad|\quad \textbf{Prize:} none \quad|\quad \textbf{Updated:} 2025-12-07

\noindent\textbf{Tags:} number theory

\medskip

\noindent\textbf{Statement:}

Let $p&#62;q\geq 2$ be two coprime integers. We call $n$ representable if it is the sum of integers of the form $p^kq^l$, none of which divide each other. If $\{p,q\}\neq \{2,3\}$ then what can be said about the density of non-representable numbers? Are there infinitely many coprime non-representable numbers?




A problem of Erd\H{o}s and Lewin \cite{ErLe96}, who proved that there are finitely many non-representable numbers if and only if $\{p,q\}=\{2,3\}$. Indeed, in \cite{Er92b} Erd\H{o}s wrote 'last year I made the following silly conjecture': every integer $n$ can be written as the sum of distinct integers of the form $2^k3^l$, none of which divide any other. He wrote 'I mistakenly thought that this was a nice and difficult conjecture but Jansen and several others found a simple proof by induction.' This simple proof is as follows: one proves the stronger fact that such a representation always exists, and moreover if $n$ is even then all the summands can be taken to be even: if $n=2m$ we are done applying the inductive hypothesis to $m$. Otherwise if $n$ is odd then let $3^k$ be the largest power of $3$ which is $\leq n$ and apply the inductive hypothesis to $n-3^k$ (which is even). Yu and Chen \cite{YuCh22} prove that the set of representable numbers has density zero whenever $q&gt;3$ or $q=3$ and $p&gt;6$ or $q=2$ and $p&gt;10$. They also prove that there are infinitely many coprime non-representable numbers if $q&gt;3$ or $q=3$ and $p\neq 5$ or $q=2$ and $p\not\in \{3,5,9\}$. Erd\H{o}s and Lewin \cite{ErLe96} also asked whether all large integers $n$ can be written as a sum of $2^k3^l$, none of which divide another, each of which is $&gt;f(n)$ for some $f(n)\to \infty$. Let $f(n)$ be the fastest growing such $f(n)$. Yu and Chen \cite{YuCh22} proved\[\frac{n}{(\log n)^{\log_23}}\ll f(n) \ll \frac{n}{\log n}.\]Yang and Zhao \cite{YaZh25} improved the lower bound to $f(n)\gg n/\log n$. van Doorn has observed in the comments that a result of Blecksmith, McCallum, and Selfridge \cite{BMS98} already implies\[f(n)\sim \frac{\log 2\log 3}{2}\frac{n}{\log n}.\]The case of three powers is the subject of [123] , and see also [845] for more on the case $\{p,q\}=\{2,3\}$. The problem [246] addresses the topic without the non-divisibility condition.




References

[BMS98] Blecksmith, Richard and McCallum, Michael and Selfridge, J. L., {$3$}-smooth representations of integers . Amer. Math. Monthly (1998), 529--543.

[Er92b] Erd\H{o}s, Paul, Some of my favourite problems in various branches of combinatorics . Matematiche (Catania) (1992), 231-240.

[ErLe96] Erd\H{o}s, P. and Lewin, Mordechai, $d$-complete sequences of integers . Math. Comp. (1996), 837-840.

[YaZh25] Yang, Quan-Hui and Zhao, Lilu, A conjecture of {Y}u and {C}hen related to the {E}rd\H os-{L}ewin theorem . Acta Arith. (2025), 277--286.

[YuCh22] Yu, Wang-Xing and Chen, Yong-Gao, On a conjecture of {E}rd\H{o}s and {L}ewin . J. Number Theory (2022), 763--778.

\noindent\rule{\linewidth}{0.4pt}


%% ============================================================
\section{Problem \#1111}
\label{prob:1111}

\noindent\textbf{Status:} OPEN \quad|\quad \textbf{Prize:} none \quad|\quad \textbf{Updated:} 2025-12-07

\noindent\textbf{Tags:} graph theory

\medskip

\noindent\textbf{Statement:}

If $G$ is a finite graph and $A,B$ are disjoint sets of vertices then we call $A,B$ anticomplete if there are no edges between $A$ and $B$. If $t,c\geq 1$ then there exists $d\geq 1$ such that if $\chi(G)\geq d$ and $\omega(G)&#60;t$ then there are anticomplete sets $A,B$ with $\chi(A)\geq \chi(B)\geq c$.




A problem of El Zahar and Erd\H{o}s \cite{ElEr85}, who show that it suffices to consider the case $t\leq c$. Let $d(t,c)$ be the minimal such $d$. El Zahar and Erd\H{o}s note that a result of Wagon \cite{Wa80b} implies $d(t,2)\leq \binom{t}{2}+1$ (and in fact $d(t+1,2)\leq d(t,2)+t$). We also have $t(2,2)=2$ and $t(3,2)=4$ and $t(4,2)=5$. El Zahar and Erd\H{o}s proved $d(3,3)\leq 8$ and\[d(t,3) \leq 2\binom{t-1}{3}+7\binom{t-1}{2}+t\]for $t&#62;3$. Nguyen, Scott, and Seymour \cite{NSS24} prove that if $t,c\geq 1$ then there exists $d\geq 1$ such that if $\chi(G)\geq d$ and $\omega(G)&#60;t$ then there are anticomplete sets $A,B$ with $\chi(B)\geq c$ and such that the minimum degree of the induced graph on $A$ is at least $c$.




References

[ElEr85] El-Zahar, M. and Erd\H{o}s, P., On the existence of two nonneighboring subgraphs in a graph . Combinatorica (1985), 295--300.

[NSS24] Nguyen, Tung and Scott, Alex and Seymour, Paul, On a problem of {E}l-{Z}ahar and {E}rd\H{o}s . J. Combin. Theory Ser. B (2024), 211--222.

[Wa80b] Wagon, Stanley, A bound on the chromatic number of graphs without certain induced subgraphs . J. Combin. Theory Ser. B (1980), 345--346.

\noindent\rule{\linewidth}{0.4pt}


%% ============================================================
\section{Problem \#1112}
\label{prob:1112}

\noindent\textbf{Status:} OPEN \quad|\quad \textbf{Prize:} none \quad|\quad \textbf{Updated:} 2025-12-28

\noindent\textbf{Tags:} additive combinatorics

\medskip

\noindent\textbf{Statement:}

Let $1\leq d_1&#60;d_2$ and $k\geq 3$. Does there exist an integer $r$ such that if $B=\{b_1&#60;\cdots\}$ is a lacunary sequence of positive integers with $b_{i+1}\geq rb_i$ then there exists a sequence of positive integers $A=\{a_1&#60;\cdots\}$ such that\[d_1\leq a_{i+1}-a_i\leq d_2\]for all $i\geq 1$ and $(kA)\cap B=\emptyset$, where $kA$ is the $k$-fold sumset?




Erd\H{o}s and Graham \cite{ErGr80} noted that if $B=\{b_1&#60;b_2&#60;\cdots\}$ with $b_1\geq 5$ and $b_{i+1}\geq 2b_i$ then there is a set $A=\{a_1&#60;a_2&#60;\cdots\}$ with $2\leq a_{k+1}-a_k\leq 3$ for all $k$ such that $(A+A)\cap B=\emptyset$. Bollob\'{a}s, Hegyv\'{a}ri, and Jin \cite{BHJ97} showed that such an $A$ must exist if $b_{i+1}\geq 2b_i-O(1)$, and that this is best possible. Erd\H{o}s and Graham go on to say 'whether such behavior can hold for $A+A+A$ (or more summands) is not known'. The above question is my best interpretation of what they intended. Bollob\'{a}s, Hegyv\'{a}ri, and Jin \cite{BHJ97} provide a negative answer in that, for any sequence of integers $1\leq r_1&#60;r_2&#60;\cdots$, there is a $B$ as above with $b_{i+1}\geq r_ib_i$ such that $(A+A+A)\cap B\neq\emptyset$ for any $A$ with $2\leq a_{i+1}-a_i\leq 3$. They define, more generally, $r_k(d_1,d_2)$ as the smallest $r$ (if it exists) such that if $b_{i+1}\geq rb_i$ then there exists $A$ with $d_1\leq a_{i+1}-a_i\leq d_2$ such that $(kA)\cap B=\emptyset$, where $kA$ is the $k$-fold sumset. It follows from the above results that $r_2(2,3)=2$ and that $r_3(2,3)$ does not exist. Chen \cite{Ch00} proved that $r_2(a,b)\leq 2$ for any integers $a&#60;b$ with $b\neq 2a$, and that $r_2(a,2a)\geq 2$ for all integers $a$. The more general question of existence of $r_k(a,b)$ for $k\geq 3$ remains open. Some further technical non-existence results are given by Tang and Yang \cite{TaYa21}. [Note the stated problem is a generous interpretation of a very ambiguous remark in \cite{ErGr80}, so it might be more appropriate to call this a problem 'inspired by Erd\H{o}s and Graham'.]




References

[BHJ97] Bollob\'{a}s, B\'{e}la and Hegyv\'{a}ri, Norbert and Jin, Guoping, On a problem of {E}rd\H{o}s and {G}raham . Discrete Math. (1997), 253--257.

[Ch00] Chen, Yong-Gao, On sums and intersects of sequences . Discrete Math. (2000), 351--354.

[ErGr80] Erd\H{o}s, P. and Graham, R., Old and new problems and results in combinatorial number theory . Monographies de L'Enseignement Mathematique (1980).

[TaYa21] Tang, Min and Yang, Quan-Hui, On a problem of {E}rd\H{o}s and {G}raham . Publ. Math. Debrecen (2021), 485--493.

\noindent\rule{\linewidth}{0.4pt}


%% ============================================================
\section{Problem \#1113}
\label{prob:1113}

\noindent\textbf{Status:} OPEN \quad|\quad \textbf{Prize:} none \quad|\quad \textbf{Updated:} 2025-12-28

\noindent\textbf{Tags:} number theory, covering systems

\medskip
\noindent\textit{Sierpinski numbers}


\noindent\textbf{Statement:}

A positive odd integer $m$ such that none of $2^km+1$ are prime for $k\geq 0$ is called a Sierpinski number . We say that a set of primes $P$ is a covering set for $m$ if every $2^km+1$ is divisible by some $p\in P$. Are there Sierpinski numbers with no finite covering set of primes?




Sierpinski \cite{Si60} proved that there are infinitely many Sierpinski numbers, using covering systems to construct suitable covering sets for any $m$ satisfying a certain congruence. This establishes that there is a positive density set of such $m$. The smallest Sierpinski number is believed to be $78557$, which was found by Selfridge. Erd\H{o}s and Graham \cite{ErGr80} asked whether there are Sierpinski numbers for which a covering system is not 'responsible', for which the best interpretation seems to be the above question. This is formulated precisely in problem F13 of Guy's collection \cite{Gu04}. Erd\H{o}s and Graham thought the answer is yes (in that there are such Sierpinski numbers), since otherwise this would imply there are infinitely many Fermat primes. There is now further evidence with a concrete example: an argument of Izotov \cite{Iz95}, given in more detail by Filaseta, Finch, and Kozek \cite{FFK08}, suggests that $m=734110615000775^4$ is a Sierpinski number without a covering set. (Izotov proved that this $m$ is indeed a Sierpinski number.) Filaseta, Finch, and Kozek \cite{FFK08} give a revised conjecture, suggesting that every Sierpinski number is either a perfect power or else has a finite covering set of primes. They also prove that for every $l\geq 1$ there is an $m$ such that $2^km^i+1$ is composite for all $1\leq i\leq l$ and $k\geq 0$. See also [203] , and [276] for another problem in which the question is whether covering systems are always responsible.




References

[ErGr80] Erd\H{o}s, P. and Graham, R., Old and new problems and results in combinatorial number theory . Monographies de L'Enseignement Mathematique (1980).

[FFK08] Filaseta, Michael and Finch, Carrie and Kozek, Mark, On powers associated with Sierpi\'nski numbers, {R}iesel numbers and {P}olignac's conjecture . J. Number Theory (2008), 1916--1940.

[Gu04] Guy, Richard K., Unsolved problems in number theory . (2004), xviii+437.

[Iz95] Izotov, Anatoly S., A note on Sierpi\'nski numbers . Fibonacci Quart. (1995), 206--207.

[Si60] Sierpi\'nski, W., Sur un probl\`eme concernant les nombres {$k\cdot 2\sp{n}+1$} . Elem. Math. (1960), 73--74.

\noindent\rule{\linewidth}{0.4pt}


%% ============================================================
\section{Problem \#1117}
\label{prob:1117}

\noindent\textbf{Status:} OPEN \quad|\quad \textbf{Prize:} none \quad|\quad \textbf{Updated:} 2025-12-29

\noindent\textbf{Tags:} analysis

\medskip

\noindent\textbf{Statement:}

Let $f(z)$ be an entire function which is not a monomial. Let $\nu(r)$ count the number of $z$ with $\lvert z\rvert=r$ such that $\lvert f(z)\rvert=\max_{\lvert z\rvert=r}\lvert f(z)\rvert$. (This is a finite quantity if $f$ is not a monomial.) Is it possible for\[\limsup \nu(r)=\infty?\]Is it possible for\[\liminf \nu(r)=\infty?\]




This is Problem 2.16 in \cite{Ha74}, where it is attributed to Erd\H{o}s. The answer to the first question is yes, as shown by Herzog and Piranian \cite{HePi68}. The second question is still open, although an 'approximate' affirmative answer is given by Gl\"{u}cksam and Pardo-Sim\'{o}n \cite{GlPa24}.




References

[GlPa24] Gl\"ucksam, Adi and Pardo-Sim\'on, Leticia, An approximate solution to {E}rd\H{o}s' maximum modulus points problem . J. Math. Anal. Appl. (2024), Paper No. 127768, 20.

[Ha74] Hayman, W. K., Research problems in function theory: new problems . (1974), 155--180.

[HePi68] Herzog, F. and Piranian, G., The counting function for points of maximum modulus . (1968), 240--243.

\noindent\rule{\linewidth}{0.4pt}


%% ============================================================
\section{Problem \#1120}
\label{prob:1120}

\noindent\textbf{Status:} OPEN \quad|\quad \textbf{Prize:} none \quad|\quad \textbf{Updated:} 2025-12-29

\noindent\textbf{Tags:} analysis

\medskip

\noindent\textbf{Statement:}

Let $f\in \mathbb{C}[z]$ be a monic polynomial of degree $n$, all of whose roots satisfy $\lvert z\rvert\leq 1$. Let\[E= \{ z : \lvert f(z)\rvert \leq 1\}.\]What is the shortest length of a path in $E$ joining $z=0$ to $\lvert z\rvert =1$?




This is Problem 4.22 in \cite{Ha74}, where it is attributed to Erd\H{o}s. In \cite{Ha74} it is reported that Clunie and Netanyahu (personal communication) showed that a path always exists which joins $z=0$ to $\lvert z\rvert=1$ in $A$. Erd\H{o}s wrote 'presumably this tends to infinity with $n$, but not too fast'. The trivial lower bound for the length of this path is $1$, which is achieved for $f(z)=z^n$. The interesting side of this question is what the worst case behaviour is (as a function of $n$). See also [1041] .




References

[Ha74] Hayman, W. K., Research problems in function theory: new problems . (1974), 155--180.

\noindent\rule{\linewidth}{0.4pt}


%% ============================================================
\section{Problem \#1122}
\label{prob:1122}

\noindent\textbf{Status:} OPEN \quad|\quad \textbf{Prize:} none \quad|\quad \textbf{Updated:} 2025-12-30

\noindent\textbf{Tags:} number theory

\medskip

\noindent\textbf{Statement:}

Let $f:\mathbb{N}\to \mathbb{R}$ be an additive function (i.e. $f(ab)=f(a)+f(b)$ whenever $(a,b)=1$). Let\[A=\{ n \geq 1: f(n+1)&#60; f(n)\}.\]If $\lvert A\cap [1,X]\rvert =o(X)$ then must $f(n)=c\log n$ for some $c\in \mathbb{R}$?




Erd\H{o}s proved that $f(n)=c\log n$ for some $c\in\mathbb{R}$ if $A$ is empty, or if $f(n+1)-f(n)=o(1)$. Partial progress was made by Mangerel \cite{Ma22}, who proved that this is true if\[\lvert A\cap [1,X]\rvert \ll \frac{X}{(\log X)^{2+c}}\]for some $c&gt;0$, and if $f(p)$ does not have very large values (in a certain technical sense). See also [491] .




References

[Ma22] Mangerel, Alexander P., Additive functions in short intervals, gaps and a conjecture of {E}rd\H{o}s . Ramanujan J. (2022), 1023--1090.

\noindent\rule{\linewidth}{0.4pt}


%% ============================================================
\section{Problem \#1131}
\label{prob:1131}

\noindent\textbf{Status:} OPEN \quad|\quad \textbf{Prize:} none \quad|\quad \textbf{Updated:} 2026-01-01

\noindent\textbf{Tags:} analysis, polynomials

\medskip

\noindent\textbf{Statement:}

For $x_1,\ldots,x_n\in [-1,1]$ let\[l_k(x)=\frac{\prod_{i\neq k}(x-x_i)}{\prod_{i\neq k}(x_k-x_i)},\]which are such that $l_k(x_k)=1$ and $l_k(x_i)=0$ for $i\neq k$. What is the minimal value of\[I(x_1,\ldots,x_n)=\int_{-1}^1 \sum_k \lvert l_k(x)\rvert^2\mathrm{d}x?\]In particular, is it true that\[\min I =2-(1+o(1))\frac{1}{n}?\]




Erd\H{o}s first conjectured this minimum was achieved by taking the $x_i$ to be the roots of the integral of the Legendre polynomial, since Fejer \cite{Fe32} had earlier shown these to be minimisers of\[\max_{x\in [-1,1]}\sum_k \lvert l_k(x)\rvert^2.\]This was disproved by Szabados \cite{Sz66} for every $n&#62;3$. Erd\H{o}s, Szabados, Varma, and V\'{e}rtesi \cite{ESVV94} proved that\[2-O\left(\frac{(\log n)^2}{n}\right)\leq \min I\leq 2-\frac{2}{2n-1}\]where the upper bound is witnessed by the roots of the integral of the Legendre polynomial as above.




References

[ESVV94] Erd\H{o}s, P. and Szabados, J. and Varma, A. K. and V\'{e}rtesi, P., On an interpolation theoretical extremal problem . Studia Sci. Math. Hungar. (1994), 55--60.

[Fe32] Fej\'{e}r, Leopold, Bestimmung derjenigen {A}bszissen eines Intervalles, f\"ur welche die {Q}uadratsumme der {G}rundfunktionen der {L}agrangeschen Interpolation im Intervalle ein {M}\"oglichst kleines {M}aximum {B}esitzt . Ann. Scuola Norm. Super. Pisa Cl. Sci. (2) (1932), 263--276.

[Sz66] Szabados, J., On a problem of {P}. {E}rd\H{o}s . Acta Math. Acad. Sci. Hungar. (1966), 155--157.

\noindent\rule{\linewidth}{0.4pt}


%% ============================================================
\section{Problem \#1132}
\label{prob:1132}

\noindent\textbf{Status:} OPEN \quad|\quad \textbf{Prize:} none \quad|\quad \textbf{Updated:} 2026-01-01

\noindent\textbf{Tags:} analysis, polynomials

\medskip

\noindent\textbf{Statement:}

For $x_1,\ldots,x_n\in [-1,1]$ let\[l_k(x)=\frac{\prod_{i\neq k}(x-x_i)}{\prod_{i\neq k}(x_k-x_i)},\]which are such that $l_k(x_k)=1$ and $l_k(x_i)=0$ for $i\neq k$. Let $x_1,x_2,\ldots\in [-1,1]$ be an infinite sequence, and let\[L_n(x) = \sum_{1\leq k\leq n}\lvert l_k(x)\rvert,\]where each $l_k(x)$ is defined above with respect to $x_1,\ldots,x_n$. Must there exist $x\in (-1,1)$ such that\[L_n(x) &#62;\frac{2}{\pi}\log n-O(1)\]for infinitely many $n$? Is it true that\[\limsup_{n\to \infty}\frac{L_n(x)}{\log n}\geq \frac{2}{\pi}\]for almost all $x\in (-1,1)$?




A result of Bernstein \cite{Be31} implies that the set of $x\in(-1,1)$ for which\[\limsup_{n\to \infty}\frac{L_n(x)}{\log n}\geq \frac{2}{\pi}\]is everywhere dense. Erd\H{o}s \cite{Er61c} proved that, for any fixed $x_1,\ldots,x_n\in [-1,1]$,\[\max_{x\in [-1,1]}\sum_{1\leq k\leq n}\lvert l_k(x)\rvert&gt;\frac{2}{\pi}\log n-O(1).\]Tao \cite{Ta26b} has proved that for any function $\omega(n)$ which tends to infinity with $n$ there exists a dense set of $x\in (-1,1)$ such that\[L_n(x)\geq \frac{2}{\pi}\log n-\omega(n)\]for infinitely many $n$. Tao also notes that it is unclear in the question of Erd\H{o}s whether the constant in the $O(1)$ term may depend on $x$. See also [1129] for more on $L_n(x)$, and also [1153] .




References

[Be31] S. Bernstein, Sur la limitation des valeurs d'un polynome $P_n(x)$ de degr\'{e} n sur tout un segment par ses valeurs en $(n+1)$ points du segment . Izv. Akad. Nauk. SSSR (1931), 1025-1050.

[Er61c] Erd\H{o}s, P., Problems and results on the theory of interpolation. II . Acta Math. Acad. Sci. Hungar. (1961), 235-244.

[Ta26b] T. Tao, Local Bernstein theory, and lower bounds for Lebesgue constants . arXiv:2603.21453 (2026).

\noindent\rule{\linewidth}{0.4pt}


%% ============================================================
\section{Problem \#1133}
\label{prob:1133}

\noindent\textbf{Status:} OPEN \quad|\quad \textbf{Prize:} none \quad|\quad \textbf{Updated:} 2026-01-01

\noindent\textbf{Tags:} analysis, polynomials

\medskip

\noindent\textbf{Statement:}

Let $C&#62;0$. There exists $\epsilon&#62;0$ such that if $n$ is sufficiently large the following holds. For any $x_1,\ldots,x_n\in [-1,1]$ there exist $y_1,\ldots,y_n\in [-1,1]$ such that, if $P$ is a polynomial of degree $m&lt;(1+\epsilon)n$ with $P(x_i)=y_i$ for at least $(1-\epsilon)n$ many $1\leq i\leq n$, then\[\max_{x\in [-1,1]}\lvert P(x)\rvert &gt;C.\]




Erd\H{o}s proved that, for any $C&#62;0$, there exists $\epsilon&#62;0$ such that if $n$ is sufficiently large and $m=\lfloor (1+\epsilon)n\rfloor$ then for any $x_1,\ldots,x_m\in [-1,1]$ there is a polynomial $P$ of degree $n$ such that $\lvert P(x_i)\rvert\leq 1$ for $1\leq i\leq m$ and\[\max_{x\in [-1,1]}\lvert P(x)\rvert&#62;C.\]The conjectured statement would also imply this, but Erd\H{o}s in \cite{Er67} says he could not even prove it for $m=n$.




References

[Er67] Erd\H{o}s, P., Problems and results on the convergence and divergence properties of the Lagrange interpolation polynomials and some extremal problems . Mathematica (Cluj) (1967), 65-73.

\noindent\rule{\linewidth}{0.4pt}


%% ============================================================
\section{Problem \#1135}
\label{prob:1135}

\noindent\textbf{Status:} OPEN \quad|\quad \textbf{Prize:} \$500 \quad|\quad \textbf{Updated:} 2026-01-11

\noindent\textbf{Tags:} number theory

\medskip
\noindent\textit{Collatz conjecture}


\noindent\textbf{Statement:}

Define $f:\mathbb{N}\to \mathbb{N}$ by $f(n)=n/2$ if $n$ is even and $f(n)=\frac{3n+1}{2}$ if $n$ is odd. Given any integer $m\geq 1$ does there exist $k\geq 1$ such that $f^{(k)}(m)=1$?




The infamous Collatz conjecture . For a detailed discussion of the history and theory surrounding this problem we refer to the overview by Lagarias \cite{La10}. This is not a problem due to Erd\H{o}s; it was first devised by Collatz before 1952. Erd\H{o}s referred to this problem on several occasions as 'hopeless'. As Lagarias \cite{La16} notes, the closest Erd\H{o}s ever came to working on problems of this nature is the theorem described in the remarks to [1134] . It is often claimed that Erd\H{o}s offered \$500 for a solution to this problem; this claim originated in a survey article by Lagarias \cite{La85}. Lagarias reported, in personal communication, that this came from a conversation he had with Erd\H{o}s and Graham around 1983, in which Graham asked Erd\H{o}s to make an estimate of what value Erd\H{o}s would put the problem on his prize scale, to which Erd\H{o}s replied \$500. Therefore, strictly speaking, Erd\H{o}s never offered \$500 specifically as a prize, but we include this prize value here for comparing those problems which Erd\H{o}s rated as 'prize problems'. This is Problem E16 in Guy's collection \cite{Gu04}, in which Guy quotes Erd\H{o}s as saying "Mathematics may not be ready for such problems".




References

[Gu04] Guy, Richard K., Unsolved problems in number theory . (2004), xviii+437.

[La10] Lagarias, Jeffrey C., The {$3x+1$} problem: an overview . (2010), 3--29.

[La16] Lagarias, Jeffrey C., Erd\H{o}s, {K}larner, and the {$3x+1$} problem . Amer. Math. Monthly (2016), 753--776.

[La85] Lagarias, Jeffrey C., The {$3x+1$} problem and its generalizations . Amer. Math. Monthly (1985), 3--23.

\noindent\rule{\linewidth}{0.4pt}


%% ============================================================
\section{Problem \#1137}
\label{prob:1137}

\noindent\textbf{Status:} OPEN \quad|\quad \textbf{Prize:} none \quad|\quad \textbf{Updated:} 2026-01-23

\noindent\textbf{Tags:} number theory, primes

\medskip

\noindent\textbf{Statement:}

Let $d_n=p_{n+1}-p_n$, where $p_n$ denotes the $n$th prime. Is it true that\[\frac{\max_{n&#60;x}d_{n}d_{n-1}}{(\max_{n&#60;x}d_n)^2}\to 0\]as $x\to \infty$?

\noindent\rule{\linewidth}{0.4pt}


%% ============================================================
\section{Problem \#1139}
\label{prob:1139}

\noindent\textbf{Status:} OPEN \quad|\quad \textbf{Prize:} none \quad|\quad \textbf{Updated:} 2026-01-23

\noindent\textbf{Tags:} number theory, primes

\medskip

\noindent\textbf{Statement:}

Let $1\leq u_1&#60;u_2&#60;\cdots$ be the sequence of integers with at most $2$ prime factors. Is it true that\[\limsup \frac{u_{k+1}-u_k}{\log k}=\infty?\]

\noindent\rule{\linewidth}{0.4pt}


%% ============================================================
\section{Problem \#1142}
\label{prob:1142}

\noindent\textbf{Status:} OPEN \quad|\quad \textbf{Prize:} none \quad|\quad \textbf{Updated:} 2026-01-23

\noindent\textbf{Tags:} number theory, primes

\medskip

\noindent\textbf{Statement:}

Are there infinitely many $n$ (or any $n&#62;105$) such that $n-2^k$ is prime for all $1&#60;2^k&#60;n$?




The only known such $n$ are\[4,7,15,21,45,75,105.\]This is A039669 in the OEIS. Mientka and Weitzenkamp \cite{MiWe69} have proved there are no other such $n\leq 2^{44}$. Vaughan \cite{Va73} has proved that the number of $n\leq N$ such that $n-2^k$ is prime for all $1&lt;2^k&lt;n$ is\[&lt; \exp\left(-c\frac{\log \log \log N}{\log\log N}\log N\right)N\]for some constant $c&gt;0$. This is discussed in problem A19 of Guy's collection \cite{Gu04}. There is also further discussion on the Prime Puzzles website . Erd\H{o}s made the stronger conjecture (see [236] ) that that number of $1&lt;2^k&lt;n$ for which $n-2^k$ is prime is $o(\log n)$.




References

[Gu04] Guy, Richard K., Unsolved problems in number theory . (2004), xviii+437.

[MiWe69] Mientka, Walter E. and Weitzenkamp, Roger C., On {$f$}-plentiful numbers . J. Combinatorial Theory (1969), 374--377.

[Va73] Vaughan, R. C., Some applications of {M}ontgomery's sieve . J. Number Theory (1973), 64--79.

\noindent\rule{\linewidth}{0.4pt}


%% ============================================================
\section{Problem \#1143}
\label{prob:1143}

\noindent\textbf{Status:} OPEN \quad|\quad \textbf{Prize:} none \quad|\quad \textbf{Updated:} 2026-01-23

\noindent\textbf{Tags:} number theory, primes

\medskip

\noindent\textbf{Statement:}

Let $p_1&lt;\cdots&lt;p_u$ be primes and let $k\geq 1$. Let $F_k(p_1,\ldots,p_u)$ be such that every interval of $k$ positive integers contains at least $F_k(p_1,\ldots,p_u)$ multiples of at least one of the $p_i$. Estimate $F_k(p_1,\ldots,p_u)$, particularly in the range $k=\alpha p_u$ for constant $\alpha&gt;2$.




In \cite{Va99} it is reported that Erd\H{o}s and Selfridge found 'the exact bound' when $2&lt;\alpha&lt;3$, and that 'if $\alpha&gt;3$ then very little is known'. No reference is given, and I cannot find a relevant paper of Erd\H{o}s and Selfridge. See also [970] .




References

[Va99] Various, Some of Paul's favorite problems . Booklet produced for the conference "Paul Erd\H{o}s and his mathematics", Budapest, July 1999 (1999).

\noindent\rule{\linewidth}{0.4pt}


%% ============================================================
\section{Problem \#1144}
\label{prob:1144}

\noindent\textbf{Status:} OPEN \quad|\quad \textbf{Prize:} none \quad|\quad \textbf{Updated:} 2026-01-23

\noindent\textbf{Tags:} number theory, probability

\medskip

\noindent\textbf{Statement:}

Let $f$ be a random completely multiplicative function, where for each prime $p$ we independently choose $f(p)\in \{-1,1\}$ uniformly at random. Is it true that\[\limsup_{N\to \infty}\frac{\sum_{m\leq N}f(m)}{\sqrt{N}}=\infty\]with probability $1$?




This model of a random multiplicative function is sometimes called a Rademacher function, although this is sometimes reserved for a merely multiplicative function (which is $0$ on non-squarefree integers). See [520] for the partial sums of this alternative model. It should also be compared to another popular model of random completely multiplicative functions, Steinhaus functions, which have $f(p)$ uniformly distributed over the unit circle. Atherfold \cite{At25} has proved that, almost surely,\[\sum_{m\leq N}f(m)\ll N^{1/2}(\log N)^{1+o(1)}.\]




References

[At25] C. Atherfold, Almost sure bounds for weighted sums of Rademacher random multiplicative functions . arXiv:2501.11076 (2025).

\noindent\rule{\linewidth}{0.4pt}


%% ============================================================
\section{Problem \#1145}
\label{prob:1145}

\noindent\textbf{Status:} OPEN \quad|\quad \textbf{Prize:} none \quad|\quad \textbf{Updated:} 2026-01-23

\noindent\textbf{Tags:} additive combinatorics, additive basis

\medskip

\noindent\textbf{Statement:}

Let $A=\{1\leq a_1&#60;a_2&#60;\cdots\}$ and $B=\{1\leq b_1&#60;b_2&#60;\cdots\}$ be sets of integers with $a_n/b_n\to 1$. If $A+B$ contains all sufficiently large positive integers then is it true that $\limsup 1_A\ast 1_B(n)=\infty$?




A conjecture of Erd\H{o}s and S\'{a}rk\"{o}zy. Some condition relating $A$ and $B$ is necessary since, for example, if $A$ is the set of all integers with only even binary digits and $B$ is the set of all integers with only odd binary digits then $1_A\ast 1_B(n)=1$ for all $n$. This is a stronger form of [28] . See also [331] .

\noindent\rule{\linewidth}{0.4pt}


%% ============================================================
\section{Problem \#1146}
\label{prob:1146}

\noindent\textbf{Status:} OPEN \quad|\quad \textbf{Prize:} none \quad|\quad \textbf{Updated:} 2026-01-23

\noindent\textbf{Tags:} number theory

\medskip

\noindent\textbf{Statement:}

We say that $A\subset \mathbb{N}$ is an essential component if $d_s(A+B)&gt;d_s(B)$ for every $B\subset \mathbb{N}$ with $0&lt;d_s(B)&lt;1$ where $d_s$ is the Schnirelmann density . Is $B=\{2^m3^n : m,n\geq 0\}$ an essential component?




In \cite{Ru99} Ruzsa states "The simplest set with a chance to be an essential component is the collection of numbers in the form $2^m3^n$ and Erd\H{o}s often asked whether it is an essential component or not; I do not even have a plausible guess." See also [37] .




References

[Ru99] Ruzsa, I., Erd\H{o}s and the Integers . Journal of Number Theory (1999), 115-163.

\noindent\rule{\linewidth}{0.4pt}


%% ============================================================
\section{Problem \#1150}
\label{prob:1150}

\noindent\textbf{Status:} OPEN \quad|\quad \textbf{Prize:} none \quad|\quad \textbf{Updated:} 2026-01-23

\noindent\textbf{Tags:} analysis, polynomials

\medskip

\noindent\textbf{Statement:}

Does there exist a constant $c&#62;0$ such that, for all large $n$ and all polynomials $P$ of degree $n$ with coefficients $\pm 1$,\[\max_{\lvert z\rvert=1}\lvert P(z)\rvert &#62; (1+c)\sqrt{n}?\]




In other words, does there exist an 'ultraflat' polynomial with coefficients $\pm 1$. The answer is yes if the coefficients can take any values on the unit circle (see [230] ). The lower bound\[\max_{\lvert z\rvert=1}\lvert P(z)\rvert\geq \sqrt{n}\]is trivial from Parseval's theorem. A weaker 'flatness' question is the subject of [228] .

\noindent\rule{\linewidth}{0.4pt}


%% ============================================================
\section{Problem \#1151}
\label{prob:1151}

\noindent\textbf{Status:} OPEN \quad|\quad \textbf{Prize:} none \quad|\quad \textbf{Updated:} 2026-01-23

\noindent\textbf{Tags:} analysis, polynomials

\medskip

\noindent\textbf{Statement:}

Given $a_1,\ldots,a_n\in [-1,1]$ let\[\mathcal{L}^nf(x) = \sum_{1\leq i\leq n}f(a_i)\ell_i(x)\]be the unique polynomial of degree $n-1$ which agrees with $f$ on $a_i$ for $1\leq i\leq n$ (that is, the Lagrange interpolation polynomial). Let $a_i$ be the set of Chebyshev nodes . Prove that, for any closed $A\subseteq [-1,1]$, there exists a continuous function $f$ such that $A$ is the set of limit points of $\mathcal{L}^nf(x)$.




This is as the problem is given in \cite{Va99}, but I am unclear exactly what is intended here - is this meant for fixed, arbitrary, $x\in [-1,1]$? Erd\H{o}s \cite{Er41} proved that, if $x=\cos(\pi p/q)$ for some odd integers $p,q\geq 1$, then there is a continuous function $f$ such that\[\lim_{n\to \infty}\mathcal{L}^nf(x)=\infty,\]where the limit is taken over the sequence of Chebyshev nodes as $n\to \infty$. In \cite{Er43} he claims (without proof) that for any closed set $A$ there exists a continuous $f$ such that the limit points of $\mathcal{L}^nf(x)$ is $A$ (for specific $x$ of this shape).




References

[Er41] Erd\H{o}s, P., On divergence properties of the {L}agrange interpolation parabolas . Ann. of Math. (2) (1941), 309--315.

[Er43] Erd\H{o}s, P., A note on {F}arey series . Quart. J. Math. Oxford Ser. (1943), 82--85.

[Va99] Various, Some of Paul's favorite problems . Booklet produced for the conference "Paul Erd\H{o}s and his mathematics", Budapest, July 1999 (1999).

\noindent\rule{\linewidth}{0.4pt}


%% ============================================================
\section{Problem \#1152}
\label{prob:1152}

\noindent\textbf{Status:} OPEN \quad|\quad \textbf{Prize:} none \quad|\quad \textbf{Updated:} 2026-01-23

\noindent\textbf{Tags:} analysis, polynomials

\medskip

\noindent\textbf{Statement:}

For $n\geq 1$ fix some sequence of $n$ distinct numbers $x_{1n},\ldots,x_{nn}\in [-1,1]$. Let $\epsilon=\epsilon(n)\to 0$. Does there always exist a continuous function $f:[-1,1]\to \mathbb{R}$ such that if $p_n$ is a sequence of polynomials, with degrees $\deg p_n&#60;(1+\epsilon(n))n$, such that $p_n(x_{kn})=f(x_{kn})$ for all $1\leq k\leq n$, then $p_n(x)\not\to f(x)$ for almost all $x\in [-1,1]$?




Erd\H{o}s, Kro\'{o}, and Szabados \cite{EKS89} proved that, if $\epsilon&#62;0$ is fixed and does not $\to 0$, then there exist sequences $x_{ij}$ such that, for any continuous function $f$, there exists a sequence of polynomials $p_n$, with degrees $\deg p_n&#60;(1+\epsilon)n$, such that $p_n(x_{kn})=f(x_{kn})$ for all $1\leq k\leq n$, and $p_n(x)\to f(x)$ uniformly for all $x\in [-1,1]$.




References

[EKS89] Erd\H{o}s, P. and Kro\'o, A. and Szabados, J., On convergent interpolatory polynomials . J. Approx. Theory (1989), 232--241.

\noindent\rule{\linewidth}{0.4pt}


%% ============================================================
\section{Problem \#1155}
\label{prob:1155}

\noindent\textbf{Status:} OPEN \quad|\quad \textbf{Prize:} none \quad|\quad \textbf{Updated:} 2026-01-23

\noindent\textbf{Tags:} graph theory

\medskip

\noindent\textbf{Statement:}

Construct a random graph on $n$ vertices in the following way: begin with the complete graph $K_n$. At each stage, choose uniformly a random triangle in the graph and delete all the edges of this triangle. Repeat until the graph is triangle-free. Describe the typical parameters and structure of such a graph. In particular, if $f(n)$ is the number of edges remaining, then is it true that\[\mathbb{E}f(n)\asymp n^{3/2}\]and that $f(n) \ll n^{3/2}$ almost surely?




A problem of Bollob\'{a}s and Erd\H{o}s, described in \cite{Va99} as 'motivated by the task of generating a random triangle-free graph'. In \cite{Bo98} it says they asked this at the 'Quo Vadis, Graph Theory?' conference in Fairbanks, Alaska, in 1990, 'while admiring the playful bears'. Grable \cite{Gr97} proved that, for every $\epsilon&#62;0$,\[\mathbb{P}(f(n)&#62;n^{7/4+\epsilon})\to 0.\]Bohman, Frieze, and Lubetzky \cite{BFL15} proved that $f(n)=n^{3/2+o(1)}$ almost surely - in other words, for every $\epsilon&#62;0$,\[\mathbb{P}(n^{3/2-\epsilon}&#60;f(n)&#60;n^{3/2+\epsilon})\to 1.\]




References

[BFL15] Bohman, Tom and Frieze, Alan and Lubetzky, Eyal, Random triangle removal . Adv. Math. (2015), 379--438.

[Bo98] Bollob\'{a}s, B\'ela, To prove and conjecture: {P}aul {E}rd\H os and his mathematics . Amer. Math. Monthly (1998), 209--237.

[Gr97] Grable, David A., On random greedy triangle packing . Electron. J. Combin. (1997), Research Paper 11, 19.

[Va99] Various, Some of Paul's favorite problems . Booklet produced for the conference "Paul Erd\H{o}s and his mathematics", Budapest, July 1999 (1999).

\noindent\rule{\linewidth}{0.4pt}


%% ============================================================
\section{Problem \#1156}
\label{prob:1156}

\noindent\textbf{Status:} OPEN \quad|\quad \textbf{Prize:} none \quad|\quad \textbf{Updated:} 2026-01-23

\noindent\textbf{Tags:} graph theory, chromatic number

\medskip

\noindent\textbf{Statement:}

Let $G$ be a random graph on $n$ vertices, in which every edge is included independently with probability $1/2$. Is there some constant $C$ such that that chromatic number $\chi(G)$ is, almost surely, concentrated on at most $C$ values? Is it true that, if $\omega(n)\to \infty$ sufficiently slowly, then for every function $f(n)$\[\mathbb{P}(\lvert\chi(G)-f(n)\rvert&#60;\omega(n))&#60;1/2\]if $n$ is sufficiently large?




Bollob\'{a}s \cite{Bo88} proved that $\chi(G) \sim \frac{n}{2\log_2n}$ with high probability. Shamir and Spencer \cite{ShSp87} proved that, for any function $\omega(n)$ such that $\omega(n)/\sqrt{n}\to \infty$, there is a function $f(n)$ such that\[\mathbb{P}(\lvert\chi(G)-f(n)\rvert&lt;\omega(n))\to 1\]as $n\to \infty$. This is proved with $\omega(n)\frac{\log n}{\sqrt{n}}\to \infty$ in Exercise 3 of Section 7.9 of Alon and Spencer \cite{AlSp16} (a proof is also given by Scott \cite{Sc17}). Heckel \cite{He21} proved that if $f$ and $\omega$ are such that\[\mathbb{P}(\lvert\chi(G)-f(n)\rvert&lt;\omega(n))\to 1\]as $n\to \infty$ then, for any $c&lt;1/4$, there are infinitely many $n$ such that $\omega(n)&gt;n^c$. This was improved to any $c&#60;1/2$ by Heckel and Riordan \cite{HeRi23}.




References

[AlSp16] Alon, Noga and Spencer, Joel H., The probabilistic method . (2016), xiv+375.

[Bo88] Bollob\'{a}s, B., The chromatic number of random graphs . Combinatorica (1988), 49-55.

[He21] Heckel, Annika, Non-concentration of the chromatic number of a random graph . J. Amer. Math. Soc. (2021), 245--260.

[HeRi23] Heckel, Annika and Riordan, Oliver, How does the chromatic number of a random graph vary? . J. Lond. Math. Soc. (2) (2023), 1769--1815.

[Sc17] A. Scott, On the concentration of the chromatic number of random graphs . arXiv:0806.0178 (2017).

[ShSp87] Shamir, E. and Spencer, J., Sharp concentration of the chromatic number on random graphs {$G_{n,p}$} . Combinatorica (1987), 121--129.

\noindent\rule{\linewidth}{0.4pt}


%% ============================================================
\section{Problem \#1157}
\label{prob:1157}

\noindent\textbf{Status:} OPEN \quad|\quad \textbf{Prize:} none \quad|\quad \textbf{Updated:} 2026-01-23

\noindent\textbf{Tags:} hypergraphs, turan number

\medskip

\noindent\textbf{Statement:}

Let $t,k,r\geq 2$. Let $\mathcal{F}$ be the family of all $r$-uniform hypergraphs with $k$ vertices and $s$ edges. Determine\[\mathrm{ex}_r(n,\mathcal{F}).\]




This is a very difficult and general question, and many partial results are known. The paper of Brown, Erd\H{o}s, and S\'{o}s \cite{BES73} contains many of them, in particular the general lower bound, for all $k&gt;r$ and $s&gt;1$,\[\mathrm{ex}_r(n,\mathcal{F})\gg_{k,s} n^{\frac{rs-k}{s-1}}.\]A general conjecture of Brown, Erd\H{o}s, and S\'{o}s is that, for all $r&gt;t\geq 2$ and $s\geq 3$,\[\mathrm{ex}_t(n,\mathcal{F})=o(n^t)\]whenever $k\geq (r-t)s+t+1$. The case $t=2$ is the subject of [1178] . The case $s=r=3$ and $k=6$ is the subject of [716] . The case of $r=3$ and $k=s+2$ is the subject of [1076] .




References

[BES73] Brown, W. G. and Erd\H{o}s, P. and S\'os, V. T., Some extremal problems on {$r$}-graphs . (1973), 53--63.

\noindent\rule{\linewidth}{0.4pt}


%% ============================================================
\section{Problem \#1158}
\label{prob:1158}

\noindent\textbf{Status:} OPEN \quad|\quad \textbf{Prize:} none \quad|\quad \textbf{Updated:} 2026-01-23

\noindent\textbf{Tags:} hypergraphs, turan number

\medskip

\noindent\textbf{Statement:}

Let $K_{t}(r)$ be the complete $t$-partite $t$-uniform hypergraph with $r$ vertices in each class. Is it true that\[\mathrm{ex}_t(n,K_t(r)) \geq n^{t-r^{1-t}-o(1)}\]for all $t,r$?




Erd\H{o}s \cite{Er64f} proved that\[n^{t-O(r^{1-t})}\leq \mathrm{ex}_t(n,K_t(r)) \ll n^{t-r^{1-t}}.\]This is only known when $t=2$ and $2\leq r\leq 3$. The case $t=2$ is [714] .




References

[Er64f] Erd\H{o}s, P., On extremal problems of graphs and generalized graphs . Israel J. Math. (1964), 183--190.

\noindent\rule{\linewidth}{0.4pt}


%% ============================================================
\section{Problem \#1159}
\label{prob:1159}

\noindent\textbf{Status:} OPEN \quad|\quad \textbf{Prize:} none \quad|\quad \textbf{Updated:} 2026-01-23

\noindent\textbf{Tags:} combinatorics

\medskip

\noindent\textbf{Statement:}

Determine whether there exists a constant $C&gt;1$ such that the following holds. Let $P$ be a finite projective plane . Must there exist a set of points $S$ such that $1\leq \lvert S\cap \ell\rvert \leq C$ for all lines $\ell$?




A set which meets all lines at once is called a blocking set. In \cite{Er81} Erd\H{o}s asks the stronger question of whether this is true for all pairwise balanced block designs. Erd\H{o}s, Silverman, and Stein \cite{ESS83} proved this is true with $\lvert S\cap\ell \rvert\ll \log n$ for all lines $\ell$ (where $n$ is the order of the projective plane). See also [664] for a stronger question. This problem is mentioned after Problem 68 on Green's open problems list .




References

[ESS83] Erd\H{o}s, P. and Silverman, R. and Stein, A., Intersection properties of families containing sets of nearly the same size . Ars Combin. (1983), 247--259.

[Er81] Erd\H{o}s, P., On the combinatorial problems which I would most like to see solved . Combinatorica (1981), 25-42.

\noindent\rule{\linewidth}{0.4pt}


%% ============================================================
\section{Problem \#1160}
\label{prob:1160}

\noindent\textbf{Status:} OPEN \quad|\quad \textbf{Prize:} none \quad|\quad \textbf{Updated:} 2026-01-23

\noindent\textbf{Tags:} group theory

\medskip

\noindent\textbf{Statement:}

Let $g(n)$ denote the number of groups of order $n$. If $n\leq 2^m$ then $g(n)\leq g(2^m)$.




This is listed as an open problem (Question 22.16) in \cite{BNV07}, which reports it as a 'quite natural conjecture, whose origin we have been unable to trace satisfactorily. We have heard it attributed at various times to various people, such as Paul Erd\H{o}s and Graham Higman.' Question 22.18 of \cite{BNV07} suggests the even stronger conjecture\[\sum_{n&#60;2^m}g(n) \leq g(2^m)\]for all sufficiently large $m$ (perhaps even as soon as $m\geq 7$). Pantelidakis \cite{Pa03} proved that the original conjecture is true if $n$ is odd and $m\geq 3619$.




References

[BNV07] Blackburn, Simon R. and Neumann, Peter M. and Venkataraman, Geetha, Enumeration of finite groups . (2007), xii+281.

[Pa03] I. Pantelidakis, On the Number of Non-isomorphic Groups of the Same Order . DPhil Thesis, University of Oxford (2003).

\noindent\rule{\linewidth}{0.4pt}


%% ============================================================
\section{Problem \#1162}
\label{prob:1162}

\noindent\textbf{Status:} OPEN \quad|\quad \textbf{Prize:} none \quad|\quad \textbf{Updated:} 2026-01-23

\noindent\textbf{Tags:} group theory

\medskip

\noindent\textbf{Statement:}

Give an asymptotic formula for the number of subgroups of $S_n$. Is there a statistical theorem on their order?




A problem of Erd\H{o}s and Tur\'{a}n. Let $f(n)$ count the number of subgroups of $S_n$. Pyber \cite{Py93} proved that\[\log f(n) \asymp n^2.\]Roney-Dougal and Tracey \cite{RoTr25} have proved that\[\log f(n)=\left(\frac{1}{16}+o(1)\right)n^2.\]




References

[Py93] Pyber, L\'{a}szl\'o, Asymptotic results for permutation groups . (1993), 197--219.

[RoTr25] C. Roney-Dougal and G. Tracey, Subgroups of symmetric groups: enumeration and asymptotic properties . arXiv:2503.05416 (2025).

\noindent\rule{\linewidth}{0.4pt}


%% ============================================================
\section{Problem \#1163}
\label{prob:1163}

\noindent\textbf{Status:} OPEN \quad|\quad \textbf{Prize:} none \quad|\quad \textbf{Updated:} 2026-01-23

\noindent\textbf{Tags:} group theory

\medskip
\noindent\textit{ambiguous statement}


\noindent\textbf{Statement:}

Describe (by statistical means) the arithmetic structure of the orders of subgroups of $S_n$.




This is given in \cite{Va99} as a problem of Erd\H{o}s and Tur\'{a}n. I have reproduced the problem verbatim; I am not entirely sure what it is asking for.




References

[Va99] Various, Some of Paul's favorite problems . Booklet produced for the conference "Paul Erd\H{o}s and his mathematics", Budapest, July 1999 (1999).

\noindent\rule{\linewidth}{0.4pt}


%% ============================================================
\section{Problem \#1167}
\label{prob:1167}

\noindent\textbf{Status:} OPEN \quad|\quad \textbf{Prize:} none \quad|\quad \textbf{Updated:} 2026-01-23

\noindent\textbf{Tags:} set theory, probability

\medskip

\noindent\textbf{Statement:}

Let $r\geq 2$ be finite and $\lambda$ be an infinite cardinal. Let $\kappa_\alpha$ be cardinals for all $\alpha&#60;\gamma$. Is it true that\[2^\lambda \to (\kappa_\alpha+1)_{\alpha&#60;\gamma}^{r+1}\]implies\[\lambda \to (\kappa_\alpha)_{\alpha&#60;\gamma}^{r}?\]Here $+$ means cardinal addition, so that $\kappa_\alpha+1=\kappa_\alpha$ if $\kappa_\alpha$ is infinite.




A problem of Erd\H{o}s, Hajnal, and Rado.

\noindent\rule{\linewidth}{0.4pt}


%% ============================================================
\section{Problem \#1168}
\label{prob:1168}

\noindent\textbf{Status:} OPEN \quad|\quad \textbf{Prize:} none \quad|\quad \textbf{Updated:} 2026-01-23

\noindent\textbf{Tags:} set theory, ramsey theory

\medskip

\noindent\textbf{Statement:}

Prove that\[\aleph_{\omega+1}\not\to (\aleph_{\omega+1}, 3,\ldots,3)_{\aleph_0}^2\]without assuming the generalised continuum hypothesis.




A problem of Erd\H{o}s, Hajnal, and Rado.

\noindent\rule{\linewidth}{0.4pt}


%% ============================================================
\section{Problem \#1170}
\label{prob:1170}

\noindent\textbf{Status:} OPEN \quad|\quad \textbf{Prize:} none \quad|\quad \textbf{Updated:} 2026-01-23

\noindent\textbf{Tags:} set theory, ramsey theory

\medskip

\noindent\textbf{Statement:}

Is it consistent that\[\omega_2\to (\alpha)_2^2\]for every $\alpha &#60;\omega_2$?




Laver \cite{La82} proved the consistency of $\omega_2\to (\omega_12+1,\alpha)^2$ for all $\alpha&#60;\omega_2$. Foreman and Hajnal \cite{FoHa03} proved the consistency of $\omega_2\to (\omega_1^2+1,\alpha)^2$ for all $\alpha&#60;\omega_2$.




References

[FoHa03] Foreman, Matthew and Hajnal, Andras, A partition relation for successors of large cardinals . Math. Ann. (2003), 583--623.

[La82] Laver, Richard, An {$(\aleph \sb{2},\,\aleph \sb{2},\,\aleph \sb{0})$}-saturated ideal on {$\omega \sb{1}$} . (1982), 173--180.

\noindent\rule{\linewidth}{0.4pt}


%% ============================================================
\section{Problem \#1171}
\label{prob:1171}

\noindent\textbf{Status:} OPEN \quad|\quad \textbf{Prize:} none \quad|\quad \textbf{Updated:} 2026-01-23

\noindent\textbf{Tags:} set theory, ramsey theory

\medskip

\noindent\textbf{Statement:}

Is it true that, for all finite $k&#60;\omega$,\[\omega_1^2\to (\omega_1\omega, 3,\ldots,3)_{k+1}^2?\]




Baumgartner \cite{Ba89b} proved that, assuming a form of Martin's axiom, that $\omega_1\omega\to (\omega_1\omega, 3)^2$.




References

[Ba89b] Baumgartner, James E., Remarks on partition ordinals . (1989), 5--17.

\noindent\rule{\linewidth}{0.4pt}


%% ============================================================
\section{Problem \#1172}
\label{prob:1172}

\noindent\textbf{Status:} OPEN \quad|\quad \textbf{Prize:} none \quad|\quad \textbf{Updated:} 2026-01-23

\noindent\textbf{Tags:} set theory, ramsey theory

\medskip

\noindent\textbf{Statement:}

Establish whether the following are true assuming the generalised continuum hypothesis:\[\omega_3 \to (\omega_2,\omega_1+2)^2,\]\[\omega_3\to (\omega_2+\omega_1,\omega_2+\omega)^2,\]\[\omega_2\to (\omega_1^{\omega+2}+2, \omega_1+2)^2.\]Establish whether the following is consistent with the generalised continuum hypothesis:\[\omega_2\to (\omega_1+\omega)_2^2,\]or even $\omega_2 \to (\xi)_2^2$ for all $\xi&#60;\omega_2$.




A problem of Erd\H{o}s and Hajnal. The Erd\H{o}s-Rado partition theorem \cite{ErRa56} states that\[(2^{\kappa})^+ \to (\kappa^++1)_\kappa^2\]for every infinite cardinal $\kappa$.




References

[ErRa56] Erd\H{o}s, P. and Rado, R., A partition calculus in set theory . Bull. Amer. Math. Soc. (1956), 427-489.

\noindent\rule{\linewidth}{0.4pt}


%% ============================================================
\section{Problem \#1173}
\label{prob:1173}

\noindent\textbf{Status:} OPEN \quad|\quad \textbf{Prize:} none \quad|\quad \textbf{Updated:} 2026-01-23

\noindent\textbf{Tags:} set theory, combinatorics

\medskip

\noindent\textbf{Statement:}

Assume the generalised continuum hypothesis. Let\[f: \omega_{\omega+1}\to [\omega_{\omega+1}]^{\leq \aleph_\omega}\]be a set mapping such that\[\lvert f(\alpha)\cap f(\beta)\rvert &#60;\aleph_\omega\]for all $\alpha\neq \beta$. Does there exist a free set of cardinality $\aleph_{\omega+1}$?




A problem of Erd\H{o}s and Hajnal.

\noindent\rule{\linewidth}{0.4pt}


%% ============================================================
\section{Problem \#1175}
\label{prob:1175}

\noindent\textbf{Status:} OPEN \quad|\quad \textbf{Prize:} none \quad|\quad \textbf{Updated:} 2026-01-23

\noindent\textbf{Tags:} set theory, chromatic number

\medskip

\noindent\textbf{Statement:}

Let $\kappa$ be an uncountable cardinal. Must there exist a cardinal $\lambda$ such that every graph with chromatic number $\lambda$ contains a triangle-free subgraph with chromatic number $\kappa$?




Shelah proved that a negative answer is consistent if $\kappa=\lambda=\aleph_1$.

\noindent\rule{\linewidth}{0.4pt}


%% ============================================================
\section{Problem \#1177}
\label{prob:1177}

\noindent\textbf{Status:} OPEN \quad|\quad \textbf{Prize:} none \quad|\quad \textbf{Updated:} 2026-01-23

\noindent\textbf{Tags:} set theory, chromatic number, hypergraphs

\medskip

\noindent\textbf{Statement:}

Let $G$ be a finite $3$-uniform hypergraph, and let $F_G(\kappa)$ denote the collection of $3$-uniform hypergraphs with chromatic number $\kappa$ not containing $G$. If $F_G(\aleph_1)$ is not empty then there exists $X\in F_G(\aleph_1)$ of cardinality at most $2^{2^{\aleph_0}}$. If both $F_G(\aleph_1)$ and $F_H(\aleph_1)$ are non-empty then $F_G(\aleph_1)\cap F_H(\aleph_1)$ is non-empty. If $\kappa,\lambda$ are uncountable cardinals and $F_G(\kappa)$ is non-empty then $F_G(\lambda)$ is non-empty.




A problem of Erd\H{o}s, Galvin, and Hajnal.

\noindent\rule{\linewidth}{0.4pt}


%% ============================================================
\section{Problem \#1178}
\label{prob:1178}

\noindent\textbf{Status:} OPEN \quad|\quad \textbf{Prize:} none \quad|\quad \textbf{Updated:} 2026-01-25

\noindent\textbf{Tags:} graph theory, hypergraphs

\medskip

\noindent\textbf{Statement:}

For $r\geq 3$ let $d_r(e)$ be the minimal $d$ such that\[\mathrm{ex}_r(n,\mathcal{F})=o(n^2),\]where $\mathcal{F}$ is the family of $r$-uniform hypergraphs on $d$ vertices with $e$ edges. Prove that\[d_r(e)=(r-2)e+3\]for all $r,e\geq 3$.




A conjecture of Brown, Erd\H{o}s, and S\'{o}s \cite{BES73}, who proved that $d_r(e)\geq (r-2)e+3$ (see also [1076] ). Ruzsa and Szemer\'{e}di \cite{RuSz78} proved $d_3(3)=6$ (see [716] ). Erd\H{o}s, Frankl, and R\"{o}dl \cite{EFR86} proved\[d_r(3)= (r-2)3+3\]for all $r\geq 3$. S\'{a}rk\"{o}zy and Selkow \cite{SaSe05} proved\[d_r(e) \leq (r-2)e+2+\lfloor \log_2 e\rfloor\]for all $r,e\geq 3$. Solymosi and Solymosi \cite{SoSo17} proved that $d_3(10)\leq 14$. Conlon, Gishboliner, Levanzov, and Shapira \cite{CGLS23} proved\[d_3(e)\leq e+O\left(\frac{\log e}{\log\log e}\right)\]for all $e\geq 3$. In \cite{Er75b} Erd\H{o}s even asks whether, if $\mathcal{F}$ is the family of $3$-uniform hypergraphs on $k$ vertices with $k-3$ edges,\[\mathrm{ex}_3(n,\mathcal{F})\asymp n r_{k-3}(n),\]where $r_{k-3}(n)$ is the maximal size of a subset of $\{1,\ldots,n\}$ that does not contain a non-trivial arithmetic progression of length $k-3$. He states that Ruzsa has proved the lower bound for $k=6,7,8$. See [1157] for the general case.




References

[BES73] Brown, W. G. and Erd\H{o}s, P. and S\'os, V. T., Some extremal problems on {$r$}-graphs . (1973), 53--63.

[CGLS23] Conlon, David and Gishboliner, Lior and Levanzov, Yevgeny and Shapira, Asaf, A new bound for the {B}rown-{E}rd\H os-S\'os problem . J. Combin. Theory Ser. B (2023), 1--35.

[EFR86] Erd\H{o}s, P. and Frankl, P. and R\"{o}dl, V., The asymptotic number of graphs not containing a fixed subgraph and a problem for hypergraphs having no exponent . Graphs Combin. (1986), 113-121.

[Er75b] Erd\H{o}s, Paul, Problems and results in combinatorial number theory . Journ\'{e}es Arithm\'{e}tiques de Bordeaux (Conf., Univ. Bordeaux, Bordeaux, 1974) (1975), 295-310.

[RuSz78] Ruzsa, I. Z. and Szemer\'{e}di, E., Triple systems with no six points carrying three triangles . Combinatorics (Proc. Fifth Hungarian Colloq., Keszthely, 1976), Vol. II (1978), 939-945.

[SaSe05] S\'{a}rk\"ozy, G\'{a}bor N. and Selkow, Stanley, An extension of the {R}uzsa-Szemer\'edi theorem . Combinatorica (2005), 77--84.

[SoSo17] Solymosi, David and Solymosi, Jozsef, Small cores in 3-uniform hypergraphs . J. Combin. Theory Ser. B (2017), 897--910.

\noindent\rule{\linewidth}{0.4pt}


%% ============================================================
\section{Problem \#1181}
\label{prob:1181}

\noindent\textbf{Status:} OPEN \quad|\quad \textbf{Prize:} none \quad|\quad \textbf{Updated:} 2026-03-07

\noindent\textbf{Tags:} number theory

\medskip

\noindent\textbf{Statement:}

Let $q(n,k)$ denote the least prime which does not divide $\prod_{1\leq i\leq k}(n+i)$. Is it true that there exists some $c&#62;0$ such that, for all large $n$,\[q(n,\log n)&#60;(1-c)(\log n)^2?\]




The upper bound $q(n,\log n)\leq (1+o(1))(\log n)^2$ follows from observing that the primorial of $q(n,k)$ is at most $\prod_{1\leq i\leq k}(n+i)$. Probabilistic heuristics described by Tao in the comments of [457] suggest that\[q(n,\log n)\ll \frac{\log\log n}{\log\log\log n}{\log n}\]for all $n$. [457] concerns lower bounds for $q(n,\log n)$ (and in particular the constructions described there proved that the above upper bound would be best possible).

\noindent\rule{\linewidth}{0.4pt}


%% ============================================================
\section{Problem \#1182}
\label{prob:1182}

\noindent\textbf{Status:} OPEN \quad|\quad \textbf{Prize:} none \quad|\quad \textbf{Updated:} 2026-03-07

\noindent\textbf{Tags:} graph theory, ramsey theory

\medskip

\noindent\textbf{Statement:}

Let $f(n)$ be maximal such that there is a connected graph $G$ with $n$ vertices and $f(n)$ edges such that\[R(K_3,G)= 2n-1.\]Let $F(n)$ be maximal such that every connected graph $G$ with $n$ vertices and $\leq F(n)$ edges has\[R(K_3,G)= 2n-1.\]Estimate $f(n)$ and $F(n)$. In particular, is it true that $F(n)/n\to \infty$?




A problem of Burr, Erd\H{o}s, Faudree, Rousseau, and Schelp. It is trivial that $f(n)\geq F(n)$, and that $R(K_3,G)\geq 2n-1$ for any connected graph $G$ on $n$ vertices. By a result of Chav\'{a}tal $R(K_3,T)=2n-1$ whenever $T$ is a tree on $n$ vertices, whence $F(n)\geq n-1$. Burr, Erd\H{o}s, Faudree, Roussea, and Schelp \cite{BEFRS80} proved\[\frac{17n+1}{15}\leq F(n)\leq \left(\frac{27}{4}+o(1)\right)n(\log n)^2,\]the lower bound for all $n\geq 4$. The upper bound was improved to\[F(n)\leq 84n\]by Brandt \cite{Br96}, who expects that for all large $n$ we have $2n&lt;F(n)&lt;6n$. (This resolves the final question in the negative.) In \cite{BEFRS80} they also prove\[n^{3/2}(\log n)^{1/2} \ll f(n) \ll n^{5/3}(\log n)^{2/3}.\]The first few values of $F(n)$ and $f(n)$ (beginning with $n=2$) are $1,2,5,7,8$ and $1,2,5,8,12$ respectively, as reported in \cite{BEFRS80}. Similar questions can be asked with $K_3$ replaced by $K_m$ (where the $2n-1$ is replaced with $(m-1)(n-1)+1$). For the corresponding generalisations of $f$ and $F$ Burr, Erd\H{o}s, Faudree, Roussea, and Schelp \cite{BEFRS80} proved\[n^{\frac{2}{m-1}}\ll F_m(n)-n\ll n^{\frac{4}{m+1}+o(1)}\]and\[n^{1+\frac{1}{m-1}}\ll f_m(n) \ll n^{1+\frac{2}{m}+o(1)}.\]




References

[BEFRS80] Burr, S. A. and Erd\H{o}s, P. and Faudree, R. J. and Rousseau, C. C. and Schelp, R. H., An extremal problem in generalized {R}amsey theory . Ars Combin. (1980), 193--203.

[Br96] S. Brandt, Expanding graphs and Ramsey numbers . Seri A Mathematik, Preprints (1996).

\noindent\rule{\linewidth}{0.4pt}


%% ============================================================
\section{Problem \#1183}
\label{prob:1183}

\noindent\textbf{Status:} OPEN \quad|\quad \textbf{Prize:} none \quad|\quad \textbf{Updated:} 2026-03-07

\noindent\textbf{Tags:} combinatorics, ramsey theory

\medskip

\noindent\textbf{Statement:}

Let $f(n)$ be maximal such that in any $2$-colouring of the subsets of $\{1,\ldots,n\}$ there is always a monochromatic family of at least $f(n)$ sets which is closed under taking unions and intersections. Estimate $f(n)$. Let $F(n)$ be defined similarly, except that we only require the family be closed under taking unions. Estimate $F(n)$. In particular, is it true that $F(n)\geq n^{\omega(n)}$ for some $\omega(n)\to \infty$ as $n\to \infty$, and $F(n)&#60;(1+o(1))^n$?




A problem of Erd\H{o}s and Ulam. It is trivial that $f(n)\geq \frac{n+1}{2}$, since there is a sequence of $n+1$ many nested subsets. This seems to be all that they knew; Erd\H{o}s \cite{Er78} wrote 'we have no plausible conjecture for the true order of magnitude of $f(n)$'. They similarly had no good guess about $F(n)$. In \cite{Er78} Erd\H{o}s wrote that if the colouring is such that all subsets of same size receive the same colour then Howorka had proved that $F(n)&#62;n^{\omega(n)}$ for some $\omega(n)\to \infty$, but gave no reference.




References

[Er78] Erd\H{o}s, Paul, Problems and results in combinatorial analysis and combinatorial number theory . Proceedings of the Ninth Southeastern Conference on Combinatorics, Graph Theory, and Computing (Florida Atlantic Univ., Boca Raton, Fla., 1978) (1978), 29-40.

\noindent\rule{\linewidth}{0.4pt}


%% ============================================================
\section{Problem \#1184}
\label{prob:1184}

\noindent\textbf{Status:} OPEN \quad|\quad \textbf{Prize:} none \quad|\quad \textbf{Updated:} 2026-04-04

\noindent\textbf{Tags:} number theory, primes

\medskip

\noindent\textbf{Statement:}

Let $f(n,k)$ count the number of $1\leq i\leq k$ such that $P(n+i)&gt;k$ (where $P(m)$ is the largest prime divisor of $m$). Is it true that, if $\alpha&gt;1$ is such that $n=k^{\alpha+o(1)}$, then\[f(n,k)=(1-\rho(\alpha)+o(1))k,\]where $\rho$ is the Dickman function ?




Erd\H{o}s \cite{Er76e} proved that, for every $\alpha&#62;1$, if $k$ is sufficiently large and $n&#62;k^\alpha-k$ then\[f(n,k) &#62; \left(1-\frac{1}{\alpha}+c_\alpha\right)k\]for some constant $c_\alpha&#62;0$, and if $1&lt;\alpha&lt;2$ and $n\leq k^\alpha-k$ then\[f(n,k) &lt; (\alpha-1+o(1))k.\]He knew of no non-trivial bounds when $\alpha \geq 2$. Ramachandra, Shorey, and Tijdeman \cite{RST75b} proved that if $n&gt;\exp(c(\log k)^2)$ for a constant $c&#62;0$ then $f(n,k)\geq k-\pi(k)$.




References

[Er76e] Erd\H{o}s, P., Problems and results on consecutive integers . Publ. Math. Debrecen (1976), 271-282.

[RST75b] Ramachandra, K. and Shorey, T. N. and Tijdeman, R., On {G}rimm's problem relating to factorisation of a block of consecutive integers. {II} . J. Reine Angew. Math. (1976), 192--201.

\noindent\rule{\linewidth}{0.4pt}


%% ============================================================
\section{Problem \#1186}
\label{prob:1186}

\noindent\textbf{Status:} OPEN \quad|\quad \textbf{Prize:} none \quad|\quad \textbf{Updated:} 2026-04-04

\noindent\textbf{Tags:} additive combinatorics, arithmetic progressions

\medskip

\noindent\textbf{Statement:}

Let $\delta_k$ be such that in any $2$-colouring of $\{1,\ldots,n\}$ there exist at least $(\delta_k+o(1))n^2$ many monochromatic $k$-term arithmetic progressions. Give reasonable bounds (or even an asymptotic formula) for $\delta_k$.




In \cite{Er80} Erd\H{o}s said 'perhaps one can get an asymptotic formula' for $\delta_k$, but this seems unlikely. Van der Waerden's theorem implies $\delta_k\gg_k 1$. The probabilistic method implies\[\delta_k \leq \frac{1}{(k-1)2^{k}}.\]Parrilo, Robertson, and Saracino \cite{PRS08} have proved\[0.0511\approx \frac{1675}{32768}\leq \delta_3 \leq \frac{117}{2192}\approx 0.0533,\]and conjectured the upper bound is correct. They also report that Graham offered \$100 for determining the value of $\delta_3$ at a conference in 1999. It is easier to study this quantity if we replace $\{1,\ldots,n\}$ with a finite field $\mathbb{F}_p$; let this analogue be denoted by $\tilde{\delta}_k$. A random colouring implies $\tilde{\delta}_k\leq 1/2^{k}$. It is straightforward to prove that $\tilde{\delta}_3=1/8$ (see, for example, \cite{Wo10}). Improving on work of Cameron, Cilleruelo, and Serra \cite{CCS07}, Wolf \cite{Wo10} proved\[\frac{1}{32}\leq \tilde{\delta}_4\leq \left(1-\frac{1}{259200}\right)\frac{1}{16}.\]These bounds were improved by Lu and Peng \cite{LuPe12} to\[\frac{7}{192}\leq \tilde{\delta}_4\leq \frac{17}{300}.\]




References

[CCS07] Cameron, Peter and Cilleruelo, Javier and Serra, Oriol, On monochromatic solutions of equations in groups . Rev. Mat. Iberoam. (2007), 385--395.

[Er80] Erd\H{o}s, Paul, A survey of problems in combinatorial number theory . Ann. Discrete Math. (1980), 89-115.

[LuPe12] Lu, Linyuan and Peng, Xing, Monochromatic 4-term arithmetic progressions in 2-colorings of {$\Bbb Z_n$} . J. Combin. Theory Ser. A (2012), 1048--1065.

[PRS08] Parrilo, Pablo A. and Robertson, Aaron and Saracino, Dan, On the asymptotic minimum number of monochromatic 3-term arithmetic progressions . J. Combin. Theory Ser. A (2008), 185--192.

[Wo10] Wolf, J., The minimum number of monochromatic 4-term progressions in {$\Bbb Z_p$} . J. Comb. (2010), 53--68.

\noindent\rule{\linewidth}{0.4pt}


%% ============================================================
\section{Problem \#1188}
\label{prob:1188}

\noindent\textbf{Status:} OPEN \quad|\quad \textbf{Prize:} none \quad|\quad \textbf{Updated:} 2026-04-04

\noindent\textbf{Tags:} number theory, covering systems

\medskip

\noindent\textbf{Statement:}

Call a set of distinct integers $1&#60;n_1&#60;\cdots&#60;n_k$ with associated congruence classes $a_i\pmod{n_i}$ a distinct covering system if every integer satisfies at least one of these congruences. A minimal distinct covering system is one such that no proper subset forms a covering system. Let $F(x)$ count the number of minimal distinct covering systems with all moduli in $[1,x]$. Estimate $F(x)$.




Erd\H{o}s \cite{Er80} asked this without the minimal assumption, but otherwise this is trivial (since e.g. any system containing some fixed covering system would work, resulting in $\gg 2^x$ many such systems). He expected $F(x)\to \infty$ as $x\to \infty$ 'very slowly'. In \cite{Er80} Erd\H{o}s also asked to estimate the number of irreducible covering sets in $[1,x]$ (that is, those sets of moduli which contain no proper sets of moduli which form a covering system, perhaps with a different choice of congruence class). If this count is $G(x)$ then trivially $G(x)\leq F(x)$. See [1189] for more on irreducible covering sets. It follows from the solution of Hough \cite{Ho15} to [2] that $F(x)\to \infty$ as $x\to \infty$ (although a lower bound of $F(x) \gg \log x$ is easily established by elementary means, as observed by van Doorn in the comments). The lower bound construction of Balister, Bollob\'{a}s, Morris, Sahasrabudhe, and Tiba \cite{BBMST24} implies that there are at least $\exp(k^{3+o(1)})$ many minimal distinct covering systems, all of whose moduli are at most $\prod_{p&lt;k}p=e^{(1+o(1))k}$, which implies\[F(x)\geq \exp((\log x)^{3-o(1)})\](in particular $F(x)$ grows faster than any polynomial, perhaps contradicting Erd\H{o}s' guess that it should grow 'very slowly', although it is possible he was thinking of the number of different moduli possible). The trivial upper bound is\[F(x) \leq \exp(O(x\log x)).\]




References

[BBMST24] Balister, Paul and Bollob\'{a}s, B\'ela and Morris, Robert and Sahasrabudhe, Julian and Tiba, Marius, The structure and number of {E}rd\H os covering systems . J. Eur. Math. Soc. (JEMS) (2024), 75--109.

[Er80] Erd\H{o}s, Paul, A survey of problems in combinatorial number theory . Ann. Discrete Math. (1980), 89-115.

[Ho15] Hough, Bob, Solution of the minimum modulus problem for covering systems . Ann. of Math. (2) (2015), 361-382.

\noindent\rule{\linewidth}{0.4pt}


%% ============================================================
\section{Problem \#1189}
\label{prob:1189}

\noindent\textbf{Status:} OPEN \quad|\quad \textbf{Prize:} none \quad|\quad \textbf{Updated:} 2026-04-04

\noindent\textbf{Tags:} number theory, covering systems

\medskip

\noindent\textbf{Statement:}

Call a set of distinct integers $1&lt;n_1&lt;\cdots&lt;n_k$ a covering set if there is a choice of $a_i\pmod{n_i}$ for $1\leq i\leq k$ such that every integer satisfies at least one of these congruences. A set is an irreducible covering set if no proper subset is a covering set. How many irreducible covering sets of size $k$ are there? What is the minimum and maximum that $n_k$ can be? Determine or estimate $\max \sum\frac{1}{n_i}$, where the maximum ranges over all irreducible covering sets of size $k$. Are there infinitely many $n$ such that the divisors of $n$ (which are $&gt;1$) form an irreducible covering set?




This is subtly different from the (more usually studied) notion of minimal covering system: an irreducible covering set corresponds to the set of moduli of a minimal covering system, but not necessarily vice versa. The divisors of $12$ (which are $&gt;1$) are an example of an irreducible covering set. Sun \cite{Su07} proved more generally that, for any odd prime $p$, the divisors (which are $&gt;1$) of $2^{p-1}p$ form an irreducible covering set, settling the final question. Simpson \cite{Si85} proved $n_k\leq 2^{k-1}$. Let $I(k)$ count the number of irreducible covering sets of size $k$. Balister, Bollob\'{a}s, Morris, Sahasrabudhe, and Tiba \cite{BBMST24} proved that\[I(k) \leq \exp\left((c+o(1))\frac{k^{3/2}}{(\log k)^{1/2}}\right)\]by proving that the right-hand side is an asymptotic for the number of minimal covering systems on $k$ moduli. It is trivial that $\sum\frac{1}{n_i}&gt;1$ for any irreducible covering set. See also [1188] .




References

[BBMST24] Balister, Paul and Bollob\'{a}s, B\'ela and Morris, Robert and Sahasrabudhe, Julian and Tiba, Marius, The structure and number of {E}rd\H os covering systems . J. Eur. Math. Soc. (JEMS) (2024), 75--109.

[Si85] Simpson, R. J., Regular coverings of the integers by arithmetic progressions . Acta Arith. (1985), 145--152.

[Su07] Sun, Zhi-Wei, On covering numbers . (2007), 443--453.

\noindent\rule{\linewidth}{0.4pt}


%% ============================================================
\section{Problem \#1191}
\label{prob:1191}

\noindent\textbf{Status:} OPEN \quad|\quad \textbf{Prize:} \$1000 \quad|\quad \textbf{Updated:} 2026-04-04

\noindent\textbf{Tags:} additive combinatorics, sidon sets

\medskip

\noindent\textbf{Statement:}

Let $A\subset\mathbb{N}$ be an infinite Sidon set. Is it true that\[\liminf_{x\to \infty} \frac{\lvert A\cap [1,x]\rvert}{x^{1/2}}(\log x)^{1/2}=0?\]Does there exist an infinite Sidon set $A$ such that\[\liminf_{x\to \infty} \frac{\lvert A\cap [1,x]\rvert}{x^{1/2}}(\log x)^c&#62;0\]for some $c&#62;0$?




Erd\H{o}s proved (see for example \cite{HaRo66}) that if $A$ is an infinite Sidon set then\[\liminf_{x\to \infty} \frac{\lvert A\cap [1,x]\rvert}{x^{1/2}}(\log x)^{1/2}\leq c\]for some constant $c&gt;0$. In \cite{Er80} he offered \$1000 'for clearing up the problems' raised by this; he may have meant for finding the optimal function $f$ such that\[\liminf_{x\to \infty} \frac{\lvert A\cap [1,x]\rvert}{x^{1/2}}f(x)=0\]for all infinite Sidon sets $A$. The second question is a stronger form of [39] . See also [729] for the behaviour of the $\limsup$.




References

[Er80] Erd\H{o}s, Paul, A survey of problems in combinatorial number theory . Ann. Discrete Math. (1980), 89-115.

[HaRo66] Halberstam, H. and Roth, K. F., Sequences. {V}ol. I . (1966), xx+291.

\noindent\rule{\linewidth}{0.4pt}


%% ============================================================
\section{Problem \#1192}
\label{prob:1192}

\noindent\textbf{Status:} OPEN \quad|\quad \textbf{Prize:} none \quad|\quad \textbf{Updated:} 2026-04-04

\noindent\textbf{Tags:} additive combinatorics, additive basis

\medskip

\noindent\textbf{Statement:}

For $A\subset \mathbb{N}$ let $f_r(n)$ count the number of solutions to $n=a_1+\cdots+a_r$ with $a_i\in A$. Does there exist, for all $r\geq 2$, a basis $A$ of order $r$ (so that $f_r(n)&#62;0$ for all large $n$) such that\[\sum_{n\leq x}f_r(n)^2 \ll x\]for all $x$?




Erd\H{o}s and R\'{e}nyi proved by the probabilistic method that there exists a set $A$ such that\[\sum_{n\leq x}f_r(n)^2 \ll x\]and\[\lvert A\cap [1,x]\rvert\gg x^{1/r}\]for all $x$. Ruzsa \cite{Ru90} proved that the answer is yes for $r=2$.




References

[Ru90] Ruzsa, Imre Z., A just basis . Monatsh. Math. (1990), 145--151.

\noindent\rule{\linewidth}{0.4pt}


%% ============================================================
\section{Problem \#1194}
\label{prob:1194}

\noindent\textbf{Status:} OPEN \quad|\quad \textbf{Prize:} none \quad|\quad \textbf{Updated:} 2026-04-04

\noindent\textbf{Tags:} additive combinatorics, additive basis, sidon sets

\medskip

\noindent\textbf{Statement:}

Let $A\subset\mathbb{N}$ be such that every integer $n\geq 1$ can be written uniquely as $a_n-b_n$ for some $a_n,b_n\in A$. How fast must $a_n/n$ increase?




Such a set is called a perfect difference set; they can be constructed, for example, by a greedy approach (as detailed by Lev \cite{Le04}). The greedy construction achieves $a_n\ll n^3$. Erd\H{o}s \cite{Er80} wrote it is 'easy to see' that\[\limsup_{n\to \infty}\frac{a_n}{n}=\infty.\]Indeed, since $A$ must be a Sidon set, if $a_n \leq f(n)n$ for all $n$, then all differences in $[1,x/f(x)]$ can be realised by differences from $A\cap [1,x]$, and hence $\lvert A\cap [1,x]\rvert^2 \gg x/f(x)$. Erd\H{o}s (see for example \cite{HaRo66}) proved that for infinitely many $x$ $\lvert A\cap [1,x]\rvert \ll (x/\log x)^{1/2}$, which is a contradiction for a suitable choice of $f$. This argument proves that $a_n\gg n\log n$ for infinitely many $n$. Cilleruelo and Nathanson \cite{CiNa08} describe a method for construction dense perfect difference sets from dense Sidon sets, allowing many of the results for the latter to transfer to the former. GPT-5.4 Pro (see the comments) has given an argument which proves that $a_n\gg n^{2-o(1)}$ infinitely often; in fact $a_n \gg n^2/f(n)$ infinitely often for any function $f$ such that $\sum \frac{1}{nf(n)}$ diverges.




References

[CiNa08] Cilleruelo, Javier and Nathanson, Melvyn B., Perfect difference sets constructed from Sidon sets . Combinatorica (2008), 401--414.

[Er80] Erd\H{o}s, Paul, A survey of problems in combinatorial number theory . Ann. Discrete Math. (1980), 89-115.

[HaRo66] Halberstam, H. and Roth, K. F., Sequences. {V}ol. I . (1966), xx+291.

[Le04] Lev, Vsevolod F., Reconstructing integer sets from their representation functions . Electron. J. Combin. (2004), Research Paper 78, 6.

\noindent\rule{\linewidth}{0.4pt}


%% ============================================================
\section{Problem \#1199}
\label{prob:1199}

\noindent\textbf{Status:} OPEN \quad|\quad \textbf{Prize:} none \quad|\quad \textbf{Updated:} 2026-04-04

\noindent\textbf{Tags:} additive combinatorics, ramsey theory

\medskip

\noindent\textbf{Statement:}

Is it true that in any $2$-colouring of $\mathbb{N}$ there exists an infinite set $A$ such that all elements of $A+A$ are the same colour?




A conjecture of Owings \cite{Ow74}. Hindman \cite{Hi79} has shown that this is false for $3$-colourings. If we do not ask for the elements $2a$ for $a\in A$ to also be the same colour the answer is yes by Hindman's theorem (see [532] ).




References

[Hi79] Hindman, Neil, Partitions and sums of integers with repetition . J. Combin. Theory Ser. A (1979), 19--32.

[Ow74] J. Owings, E2494 . Amer. Math. Monthly (1974), 902.

\noindent\rule{\linewidth}{0.4pt}


%% ============================================================
\section{Problem \#1200}
\label{prob:1200}

\noindent\textbf{Status:} OPEN \quad|\quad \textbf{Prize:} none \quad|\quad \textbf{Updated:} 2026-04-04

\noindent\textbf{Tags:} number theory, primes

\medskip

\noindent\textbf{Statement:}

There exists a constant $C$ such that for all large $x$ there is a collection of primes $p_1&#60;\ldots&#60;p_k&#60;x$ with $\sum\frac{1}{p_i}&#60;C$ together with a system of congruences $a_i\pmod{p_i}$ such that every integer $n&#60;x$ satisfies at least one of these congruences.




A conjecture of Erd\H{o}s and Ruzsa which Erd\H{o}s described in \cite{Er80} as 'surprising'. He says that if this is correct then 'very likely' the $\epsilon_n$ in [688] is $\gg 1$ - certainly proving $\epsilon_n\geq c$ would prove this conjecture (taking $P$ to be all primes in $[x^c,x]$). In \cite{ErRu80} this is asked as a question: if $p_1&lt;\cdots&lt;p_k&lt;x$ are primes with $\sum \frac{1}{p_i}\leq C$ and $a_i\pmod{p_i}$ are any residue classes then must there always be $\gg_K x$ many integers $n&lt;x$ avoiding all of them? (Of course is the answer is yes then this disproves the main conjecture.) Erd\H{o}s and Ruzsa \cite{ErRu80} proved that for any $C&gt;0$ there exists a set of primes $P$ such that $\sum_{p\in P}\frac{1}{p}\leq C$ and the number of integers $n\leq x$ divisible by at least one $p\in P$ is $\gg_C x$. See also [783] and [784] .




References

[Er80] Erd\H{o}s, Paul, A survey of problems in combinatorial number theory . Ann. Discrete Math. (1980), 89-115.

[ErRu80] Erd\H{o}s, P. and Ruzsa, I. Z., On the small sieve. I. Sifting by primes . J. Number Theory (1980), 385--394.

\noindent\rule{\linewidth}{0.4pt}


%% ============================================================
\section{Problem \#1201}
\label{prob:1201}

\noindent\textbf{Status:} OPEN \quad|\quad \textbf{Prize:} none \quad|\quad \textbf{Updated:} 2026-04-04

\noindent\textbf{Tags:} number theory, primes

\medskip

\noindent\textbf{Statement:}

Is it true that for every $\epsilon,\eta&#62;0$ there exists a $k$ such that the density of $n$ for which\[P(n(n+1)\cdots(n+k))&#62;n^{1-\epsilon}\]is at least $1-\eta$ (where $P(m)$ is the greatest prime divisor of $m$)?




Erd\H{o}s wrote he could prove this for $\epsilon=1/2$.

\noindent\rule{\linewidth}{0.4pt}


%% ============================================================
\section{Problem \#1203}
\label{prob:1203}

\noindent\textbf{Status:} OPEN \quad|\quad \textbf{Prize:} none \quad|\quad \textbf{Updated:} 2026-04-04

\noindent\textbf{Tags:} number theory

\medskip

\noindent\textbf{Statement:}

If $\omega(n)$ counts the number of distinct prime divisors of $n$ then let\[F(n)=\max_k \omega(n+k)\frac{\log\log k}{\log k}.\]Prove that $F(n)\to \infty$ as $n\to \infty$.




It is easy to prove that $F(n)\geq 1-o(1)$. See also [248] , [679] , and [890]

\noindent\rule{\linewidth}{0.4pt}


%% ============================================================
\section{Problem \#1204}
\label{prob:1204}

\noindent\textbf{Status:} OPEN \quad|\quad \textbf{Prize:} none \quad|\quad \textbf{Updated:} 2026-04-04

\noindent\textbf{Tags:} number theory

\medskip

\noindent\textbf{Statement:}

We call a sequence of integers $0\leq a_1&#60;\cdots &#60;a_k$ admissible if it is missing at least one congruence class modulo every prime $p$. Let $A(k)=\min a_k$. Estimate $A(k)$ - in particular, is it true that\[A(k)\sim k\log k?\]Estimate\[B(k)=\min \frac{a_1+\cdots+a_k}{k}.\]




Erd\H{o}s \cite{Er80} attributes this problem to Elliott. The bounds are misstated in \cite{Er80}, but it is known that\[(1/2+o(1))k\log k \leq A(k)\leq (1+o(1))k\log k.\]The upper bound Erd\H{o}s attributes to Davenport, and follows simply by taking the $k$ smallest primes $&gt;k$. The lower order terms have been improved by Hensley and Richard \cite{HeRi73} (see Section 10 of \cite{Po14c} for more details). The lower bound is originally due to Elliott \cite{El65}, but was rediscovered by the Polymath project on bounded gaps between primes (see \cite{Po14c}). An upper bound of $\pi(x+y)\leq \pi(x)+(1+o(1))\pi(y)$ (see [855] ) together with the prime tuples conjecture would imply $A(k)\geq (1+o(1))k\log k$. It is trivial that $B(k)&lt;A(k)$. Taking the first $k$ primes which are $&gt;k$ implies\[B(k)\leq (1/2+o(1))k\log k,\]and since $a_j\geq A(j)$ we have $B(k) \geq \frac{1}{k}\sum_{1\leq j\leq k}A(j)$, whence if $A(k)\geq (c-o(1))k\log k$ we have $B(k) \geq (c/2-o(1))k\log k$, and so it is likely that\[B(k)\sim(1/2+o(1))k\log k.\]In \cite{Er80} Erd\H{o}s also asks about the behaviour of $a_k$ where $a_1&lt;a_2&lt;\cdots$ is the greedy admissible sequence (i.e. $a_1=0$ and $a_{i+1}$ is the smallest integer $&gt;a_{i-1}$ such that $a_1,\ldots,a_i$ is admissible).




References

[El65] Elliott, P. D. T. A., On sequences of integers . Quart. J. Math. Oxford Ser. (2) (1965), 35--45.

[Er80] Erd\H{o}s, Paul, A survey of problems in combinatorial number theory . Ann. Discrete Math. (1980), 89-115.

[HeRi73] Hensley, Douglas and Richards, Ian, On the incompatibility of two conjectures concerning primes . (1973), 123--127.

[Po14c] Polymath, D. H. J., Variants of the Selberg sieve, and bounded intervals containing many primes . Res. Math. Sci. (2014), Art. 12, 83.

\noindent\rule{\linewidth}{0.4pt}


%% ============================================================
\section{Problem \#1206}
\label{prob:1206}

\noindent\textbf{Status:} OPEN \quad|\quad \textbf{Prize:} none \quad|\quad \textbf{Updated:} 2026-04-04

\noindent\textbf{Tags:} number theory, sidon sets

\medskip

\noindent\textbf{Statement:}

Does $\{1,2^3,\ldots,N^3\}$ contain a Sidon set of size $\gg N$? Is there an infinite set $A\subset \mathbb{N}$ of positive density such that $\{a^3 : a\in A\}$ is a Sidon set?




This question may also be asked for higher powers, and [324] suggests that for fifth powers $\{n^5 :n\geq 1\}$ may itself be a Sidon set. Gabdullin and Konyagin \cite{GaKo24} proved that there is a constant $c&gt;0$ such that $\{ n^3 : N-cN^{1/2}\leq n\leq N\}$ is a Sidon set. Garaev, Garayev, and Konyagin \cite{GGK26} have proved that the exponent $1/2$ can be improved to $4/7-o(1)$ for infinitely many $N$, and to $3/5$ for all $N$ if we replace $n^3$ with $n^4$. In \cite{Er80} Erd\H{o}s defines $g_k(A)$ to be the maximal size of a Sidon set of $\{a^k : a\in A\}$, and asks whether $g_k(A)\geq g_k(\{1,\ldots,N\})$ where $N=\lvert A\rvert$. (See [530] for the linear analogue.) See also [773] for an analogous question about squares.




References

[Er80] Erd\H{o}s, Paul, A survey of problems in combinatorial number theory . Ann. Discrete Math. (1980), 89-115.

[GGK26] M. Garaev, F. Garayev, and S. Konyagin, On Sidon sets with squares, cubes, and quartics in short intervals . arXiv:2602.08807 (2026).

[GaKo24] Gabdullin, M. R. and Konyagin, S. V., Trigonometric polynomials with frequencies in the set of cubes . Math. Notes (2024), 336--340.

\noindent\rule{\linewidth}{0.4pt}


%% ============================================================
\section{Problem \#1207}
\label{prob:1207}

\noindent\textbf{Status:} OPEN \quad|\quad \textbf{Prize:} none \quad|\quad \textbf{Updated:} 2026-04-04

\noindent\textbf{Tags:} geometry, distances

\medskip

\noindent\textbf{Statement:}

Let $P_d(n)$ be such that in any set of $n$ points in $\mathbb{R}^d$ there exist at least $P_d(n)$ many points which do not contain an isosceles triangle. Estimate $P_d(n)$ - in particular, is it true that\[P_2(n)&lt;n^{1-c}\]for some constant $c&gt;0$?




Erd\H{o}s \cite{Er80} attributes this problem to Riddell, and states it is not hard to prove that $P_d(n)&gt;n^{\epsilon_d}$ where $\epsilon_d\to 0$ as $d\to \infty$. Indeed, one can take any subset of points with all distinct distances, whence\[P_d(n)\geq f_d(n) \geq n^{\frac{1}{3d-3}-o(1)}\]with $f_d(n)$ defined as in [1208] . Erd\H{o}s suggests that perhaps the worst set of points is the set of lattice points in a sphere of suitable radius. The case $d=1$ is equivalent to finding sets with no three-term arithmetic progressions, and is the $k=3$ case of [201] . The bound of Pach and Tardos \cite{PaTa02} on the maximum number of isosceles triangles in a point set, together with a random deletion argument, shows that $P_2(n) \gg n^{0.432}$. In Section 5.3 of \cite{BMP05} they claim that the regular polygon shows that $P_2(n) \ll n^{1/2}$ but give no details. This seems to be incorrect, since an isosceles triangle from the vertices of a regular polygon correspond to three points whose angles form an arithmetic progression. This shows that $P_2(n)\ll r_3(n)$, where $r_3(n)$ is the maximum size of a subset of $\{1,\ldots,n\}$ without a non-trivial three-term arithmetic progression (see [140] ). See [1208] (and also [657] ) for a variant in which we ask for all sets of three points to have distinct distances.




References

[BMP05] Brass, Peter and Moser, William and Pach, J\'{a}nos, Research problems in discrete geometry . (2005), xii+499.

[Er80] Erd\H{o}s, Paul, A survey of problems in combinatorial number theory . Ann. Discrete Math. (1980), 89-115.

[PaTa02] Pach, J\'{a}nos and Tardos, G\'{a}bor, Isosceles triangles determined by a planar point set . Graphs Combin. (2002), 769--779.

\noindent\rule{\linewidth}{0.4pt}


%% ============================================================
\section{Problem \#1208}
\label{prob:1208}

\noindent\textbf{Status:} OPEN \quad|\quad \textbf{Prize:} none \quad|\quad \textbf{Updated:} 2026-04-04

\noindent\textbf{Tags:} geometry, distances

\medskip

\noindent\textbf{Statement:}

For $d\geq 2$ let $F_d(n)$ be minimal such that every set of $n$ points in $\mathbb{R}^d$ contains a set of $F_d(n)$ points with distinct distances. Estimate $F_d(n)$ for fixed $d$ as $n\to \infty$.




It is known that\[\frac{n^{1/3}}{(\log n)^{1/3}}\ll F_2(n) \ll \frac{n^{1/2}}{(\log n)^{1/4}}.\]The upper bound is demonstrated by the set of integer lattice points in a square grid (which have $\ll n/\sqrt{\log n}$ many distinct distances). The lower bound was proved by Charalambides \cite{Ch13}. Thiele \cite{Th95} proved $F_d(n) \gg n^{\frac{1}{3d-2}}$ for all $d\geq 3$. This was improved to\[F_d(n) \gg_d n^{\frac{1}{3d-3}}(\log n)^{1/3-\frac{2}{3d-3}}\]by Conlon, Fox, Gasarch, Harris, Ulrich, and Zbarsky \cite{CFGHUZ15}. The example of integer lattice points shows that $F_d(n)\ll n^{1/d}$ for $d\geq 2$. Erd\H{o}s and Guy \cite{ErGu70} consider the problem of determining the maximal size of such a subset of the $n^{1/d}\times\cdots \times n^{1/d}$ integer grid, showing that it is at least $n^{\frac{2}{3d}-o(1)}$ (and at most $n^{1/d}$). Lefmann and Thiele \cite{LeTh95} have improved the lower bound to $n^{\frac{2}{3d}}$. The case $d=1$ is the subject of [530] , and Koml\'{o}s, Sulyok, and Szemer\'{e}di \cite{KSS75} have proved that $F_1(n)\asymp n^{1/2}$. See also [1207] . The behaviour of $F_d(n)$ for fixed $n\geq 3$ as $d\to \infty$ is the subject of [1088] .




References

[CFGHUZ15] Conlon, David and Fox, Jacob and Gasarch, William and Harris, David G. and Ulrich, Douglas and Zbarsky, Samuel, Distinct volume subsets . SIAM J. Discrete Math. (2015), 472--480.

[Ch13] Charalambides, Marcos, A note on distinct distance subsets . J. Geom. (2013), 439--442.

[ErGu70] Erd\H{o}s, P. and Guy, R. K., Distinct distances between lattice points . Elem. Math. (1970), 121--123.

[KSS75] Koml\'{o}s, J. and Sulyok, M. and Szemeredi, E., Linear problems in combinatorial number theory . Acta Math. Acad. Sci. Hungar. (1975), 113-121.

[LeTh95] Lefmann, Hanno and Thiele, Torsten, Point sets with distinct distances . Combinatorica (1995), 379--408.

[Th95] T. Thiele, Geometric selection problems and hypergraphs . PhD thesis, Freir Universitat Berlin (1995).

\noindent\rule{\linewidth}{0.4pt}


%% ============================================================
\section{Problem \#1209}
\label{prob:1209}

\noindent\textbf{Status:} OPEN \quad|\quad \textbf{Prize:} none \quad|\quad \textbf{Updated:} 2026-04-04

\noindent\textbf{Tags:} number theory

\medskip

\noindent\textbf{Statement:}

Let $A=\{a_1&#60;a_2&#60;\cdots\}$ be a sequence of integers which tends to infinity sufficiently fast. If there is an $n$ such that all $n+a_k$ are primes then must there exist infinitely many such $n$? What if we ask for $n+a_k$ to be squarefree instead of prime? Are there $n$ such that $n+2^{2^k}$ is always a prime (or always squarefree, or infinitely often a prime, or infinitely often squarefree)?




For the first two questions Erd\H{o}s \cite{Er80} wrote 'unless I overlook a trivial way of getting a counterexample these questions are quite hopeless'. There is indeed a trivial counterexample (a variant of the construction in [429] ): define $a_1=2$ and for $k\geq 2$ let $a_k&gt;a_{k-1}$ be a prime such that $a_k+k\equiv 0\pmod{q_k}$, where $q_k$ is some prime not dividing $k$. This sequence can be made to grow arbitrarily fast. A similar construction with $\pmod{q_k^2}$ provides a counterexample to the squarefree question. ebarschkis and GPT have proved that there are no $n$ such that $n+2^{2^k}$ is always prime: let $n\geq 3$ be any odd integer. If $k$ is chosen sufficiently large, and $p=n+2^{2^{k}}$ is prime, then the multiplicative order of $2^{2^k}\pmod{p}$, say $m$ is odd, and hence if $l$ is chosen such that $2^l\equiv 1\pmod{m}$ then $p\mid n+2^{2^{k+rl}}$ for all $r\geq 1$. See also [429] and [1102] .




References

[Er80] Erd\H{o}s, Paul, A survey of problems in combinatorial number theory . Ann. Discrete Math. (1980), 89-115.

\noindent\rule{\linewidth}{0.4pt}


%% ============================================================
\section{Problem \#1210}
\label{prob:1210}

\noindent\textbf{Status:} OPEN \quad|\quad \textbf{Prize:} none \quad|\quad \textbf{Updated:} 2026-04-04

\noindent\textbf{Tags:} number theory

\medskip

\noindent\textbf{Statement:}

Let $A\subseteq [1,n)$ be a set of integers such that $(a,b)=1$ for all distinct $a,b\in A$. Is it true that\[\sum_{a\in A}\frac{1}{n-a}\leq \sum_{p&#60;n}\frac{1}{p}+O(1)?\]




In \cite{Er80} he claims he 'did not state [this] quite correctly' in \cite{Er77c}. The problem in \cite{Er77c} which Erd\H{o}s is presumably referring to states that if $n&lt;q_1&lt;\cdots&lt;q_k\leq m$ is the set of primes in $(n,m]$ then\[\sum \frac{1}{q_i-n}&lt; \sum_{p&lt;m-n}\frac{1}{p}+O(1).\]See also [460] and [950] .




References

[Er77c] Erd\H{o}s, Paul, Problems and results on combinatorial number theory. III . Number theory day (Proc. Conf., Rockefeller Univ., New York, 1976) (1977), 43-72.

[Er80] Erd\H{o}s, Paul, A survey of problems in combinatorial number theory . Ann. Discrete Math. (1980), 89-115.

\noindent\rule{\linewidth}{0.4pt}


%% ============================================================
\section{Problem \#1212}
\label{prob:1212}

\noindent\textbf{Status:} OPEN \quad|\quad \textbf{Prize:} none \quad|\quad \textbf{Updated:} 2026-04-04

\noindent\textbf{Tags:} number theory, primes

\medskip

\noindent\textbf{Statement:}

Let $G$ be the graph with vertex set those pairs $(x,y)\in \mathbb{N}^2$ with $\mathrm{gcd}(x,y)=1$, in which we join two vertices if the differ in only one coordinate, and there by $\pm 1$. Is there a path going to infinity on $G$, say $P$, such that for all $(x,y)\in P$ both $\min(x,y)&#62;1$ and at least one of $x$ or $y$ is composite?




Herzog and Stewart studied the graph $G$, which they called visible lattice points, and proved that there is only one infinite component, and conjectured that if $p$ is prime and $p\nmid a$ then $(a,p)$ always belongs to the infinite component. (This is according to \cite{Er80}; I can find no such results of Herzog and Stewart, who have a paper \cite{HeSt71} on which finite patterns appear in the set of visible lattice points, but does not address this question.) Erd\H{o}s originally asked this question with just the condition that $\min(x,y)&#62;1$ and (as he recounts in \cite{Er80}) 'foolishly offered 25 dollars for a proof. In the evening Stewart gave the simple proof: $(p_k,p_{k+1})$ can be joined to $(p_{k+1},p_{k+2})$'. Indeed, if $p_k$ is the $k$th prime then there is a path\[(p_k,p_{k+1})\to (p_k,p_{k+1}+1)\to \cdots \]\[\to (p_{k},p_{k+2})\to (p_{k}+1,p_{k+2})\to (p_{k+1},p_{k+2}).\]This is permissible provided $[p_{k+1},p_{k+2}]$ contains no multiple of $p_k$, which is true provided $p_{k+2}&lt;2p_{k}$, which is true for all $k\geq 4$. One could further demand that the path be monotone (i.e. every step increases the distance from the origin). Is there a monotone path where we change direction after a bounded number of steps? The first $50\times 50$ block of the graph $G$ restricted to those vertices $(x,y)$ with both $x,y&gt;1$ and at least one of $x$ or $y$ composite is shown {IMAGE=1212,here}.




References

[Er80] Erd\H{o}s, Paul, A survey of problems in combinatorial number theory . Ann. Discrete Math. (1980), 89-115.

[HeSt71] Herzog, Fritz and Stewart, B. M., Patterns of visible and nonvisible lattice points . Amer. Math. Monthly (1971), 487--496.

\noindent\rule{\linewidth}{0.4pt}

\end{document}
